% GENERATED FILE -- DO NOT EDIT.
%
% Produced by harness/assemble_contract.py from module.yaml and claims/*.tex.
% Edit the claim files and regenerate; `--check` fails if this file is stale.
\documentclass[11pt,a4paper]{article}
\pdftrailerid{0c488c0e399615b38612ea17e9de5b03}
%% --- 1. Mathematics (order matters: physics after braket wins \ket, \tr) ---
\usepackage{amsmath}
\usepackage{amssymb}
\usepackage{amsfonts}
\usepackage{amsthm}
\usepackage{mathtools}
\usepackage{dsfont}
\usepackage{braket}
\usepackage{physics}
%% --- 2. Page furniture ------------------------------------------------------
\usepackage{xcolor}
\usepackage{array}    % >{...} column prefixes in the wide tables of some claims
\usepackage{enumitem}
\usepackage{graphicx}
\usepackage{booktabs}
\usepackage{longtable}
\usepackage{xltabular}              % front- and back-matter tables across pages
\usepackage{needspace}
\usepackage{xurl}
\usepackage[noautomatic]{imakeidx}  % defined-terms index; latexmk runs makeindex
\makeindex[title=Index of defined terms,intoc]
\usepackage{tcolorbox}
\tcbuselibrary{skins,breakable}
%% --- 3. Links and navigation (hyperref late; bookmark, refcount after it) ---
\usepackage{hyperref}
\usepackage{bookmark}
\usepackage{refcount}
%% --- 4. Colours: every distinction also has a shape (rule, caps, size, weight)
\definecolor{nnGrey}{gray}{0.40}      % non-normative text
\definecolor{metaGrey}{gray}{0.30}    % the generated meta line
\definecolor{techGrey}{gray}{0.45}    % heading of a technical claim
\definecolor{mainRule}{RGB}{0,110,140}
\definecolor{reqAccent}{RGB}{165,70,0}
\definecolor{linkBlue}{RGB}{0,62,150}
\definecolor{extLink}{RGB}{120,40,140}  % links into another PDF
\hypersetup{colorlinks=true, linkcolor=linkBlue, citecolor=linkBlue,
  urlcolor=linkBlue, filecolor=extLink, bookmarksnumbered=true,
  bookmarksopen=true, bookmarksopenlevel=1, pdfdisplaydoctitle=true}
% A link into another PDF (\href to a .pdf, \hyperref{file}..., xr-hyper) is a
% relative URI, <file>#nameddest=<destination>, in the file-link colour:
% browser viewers (Edge, Chrome) ignore hyperref's default GoToR action.
\makeatletter
\renewcommand*{\hyper@linkfile}[3]{% anchor text, file name, destination
  \begingroup
    \ifx\\#3\\\def\CQ@target{#2}\else\def\CQ@target{#2\hyper@hash nameddest=#3}\fi
    \let\@urlcolor\@filecolor
    \hyper@linkurl{#1}{\CQ@target}%
  \endgroup}
\makeatother
%% --- 5. The mapped Unicode set (harness/contract_navigation.py, UNICODE_TEX)
\DeclareUnicodeCharacter{03B1}{\ensuremath{\alpha}}
\DeclareUnicodeCharacter{03B2}{\ensuremath{\beta}}
\DeclareUnicodeCharacter{03B3}{\ensuremath{\gamma}}
\DeclareUnicodeCharacter{03B4}{\ensuremath{\delta}}
\DeclareUnicodeCharacter{03B5}{\ensuremath{\varepsilon}}
\DeclareUnicodeCharacter{03F5}{\ensuremath{\epsilon}}
\DeclareUnicodeCharacter{03B6}{\ensuremath{\zeta}}
\DeclareUnicodeCharacter{03B7}{\ensuremath{\eta}}
\DeclareUnicodeCharacter{03B8}{\ensuremath{\theta}}
\DeclareUnicodeCharacter{03D1}{\ensuremath{\vartheta}}
\DeclareUnicodeCharacter{03B9}{\ensuremath{\iota}}
\DeclareUnicodeCharacter{03BA}{\ensuremath{\kappa}}
\DeclareUnicodeCharacter{03BB}{\ensuremath{\lambda}}
\DeclareUnicodeCharacter{03BC}{\ensuremath{\mu}}
\DeclareUnicodeCharacter{03BD}{\ensuremath{\nu}}
\DeclareUnicodeCharacter{03BE}{\ensuremath{\xi}}
\DeclareUnicodeCharacter{03C0}{\ensuremath{\pi}}
\DeclareUnicodeCharacter{03C1}{\ensuremath{\rho}}
\DeclareUnicodeCharacter{03F1}{\ensuremath{\varrho}}
\DeclareUnicodeCharacter{03C3}{\ensuremath{\sigma}}
\DeclareUnicodeCharacter{03C2}{\ensuremath{\varsigma}}
\DeclareUnicodeCharacter{03C4}{\ensuremath{\tau}}
\DeclareUnicodeCharacter{03C5}{\ensuremath{\upsilon}}
\DeclareUnicodeCharacter{03C6}{\ensuremath{\varphi}}
\DeclareUnicodeCharacter{03D5}{\ensuremath{\phi}}
\DeclareUnicodeCharacter{03C7}{\ensuremath{\chi}}
\DeclareUnicodeCharacter{03C8}{\ensuremath{\psi}}
\DeclareUnicodeCharacter{03C9}{\ensuremath{\omega}}
\DeclareUnicodeCharacter{0393}{\ensuremath{\Gamma}}
\DeclareUnicodeCharacter{0394}{\ensuremath{\Delta}}
\DeclareUnicodeCharacter{0398}{\ensuremath{\Theta}}
\DeclareUnicodeCharacter{039B}{\ensuremath{\Lambda}}
\DeclareUnicodeCharacter{039E}{\ensuremath{\Xi}}
\DeclareUnicodeCharacter{03A0}{\ensuremath{\Pi}}
\DeclareUnicodeCharacter{03A3}{\ensuremath{\Sigma}}
\DeclareUnicodeCharacter{03A5}{\ensuremath{\Upsilon}}
\DeclareUnicodeCharacter{03A6}{\ensuremath{\Phi}}
\DeclareUnicodeCharacter{03A8}{\ensuremath{\Psi}}
\DeclareUnicodeCharacter{03A9}{\ensuremath{\Omega}}
\DeclareUnicodeCharacter{00E1}{\'a}
\DeclareUnicodeCharacter{00E0}{\`a}
\DeclareUnicodeCharacter{00E2}{\^a}
\DeclareUnicodeCharacter{00E4}{\"a}
\DeclareUnicodeCharacter{00E5}{\aa{}}
\DeclareUnicodeCharacter{00E9}{\'e}
\DeclareUnicodeCharacter{00E8}{\`e}
\DeclareUnicodeCharacter{00EA}{\^e}
\DeclareUnicodeCharacter{00EB}{\"e}
\DeclareUnicodeCharacter{00ED}{\'{\i}}
\DeclareUnicodeCharacter{00EF}{\"{\i}}
\DeclareUnicodeCharacter{00F3}{\'o}
\DeclareUnicodeCharacter{00F2}{\`o}
\DeclareUnicodeCharacter{00F4}{\^o}
\DeclareUnicodeCharacter{00F6}{\"o}
\DeclareUnicodeCharacter{00F8}{\o{}}
\DeclareUnicodeCharacter{00FA}{\'u}
\DeclareUnicodeCharacter{00F9}{\`u}
\DeclareUnicodeCharacter{00FB}{\^u}
\DeclareUnicodeCharacter{00FC}{\"u}
\DeclareUnicodeCharacter{00F1}{\~n}
\DeclareUnicodeCharacter{00E7}{\c{c}}
\DeclareUnicodeCharacter{00DF}{\ss{}}
\DeclareUnicodeCharacter{00C4}{\"A}
\DeclareUnicodeCharacter{00C5}{\AA{}}
\DeclareUnicodeCharacter{00C9}{\'E}
\DeclareUnicodeCharacter{00D6}{\"O}
\DeclareUnicodeCharacter{00DC}{\"U}
\DeclareUnicodeCharacter{2013}{\textendash{}}
\DeclareUnicodeCharacter{2014}{\textemdash{}}
\DeclareUnicodeCharacter{2018}{`}
\DeclareUnicodeCharacter{2019}{'}
\DeclareUnicodeCharacter{201C}{``}
\DeclareUnicodeCharacter{201D}{''}
\DeclareUnicodeCharacter{2026}{\ldots{}}
\DeclareUnicodeCharacter{00A0}{~}
\DeclareUnicodeCharacter{00A7}{\S{}}
\DeclareUnicodeCharacter{2264}{\ensuremath{\le}}
\DeclareUnicodeCharacter{2265}{\ensuremath{\ge}}
\DeclareUnicodeCharacter{2260}{\ensuremath{\neq}}
\DeclareUnicodeCharacter{2248}{\ensuremath{\approx}}
\DeclareUnicodeCharacter{2261}{\ensuremath{\equiv}}
\DeclareUnicodeCharacter{223C}{\ensuremath{\sim}}
\DeclareUnicodeCharacter{2243}{\ensuremath{\simeq}}
\DeclareUnicodeCharacter{2245}{\ensuremath{\cong}}
\DeclareUnicodeCharacter{221D}{\ensuremath{\propto}}
\DeclareUnicodeCharacter{226A}{\ensuremath{\ll}}
\DeclareUnicodeCharacter{226B}{\ensuremath{\gg}}
\DeclareUnicodeCharacter{2192}{\ensuremath{\to}}
\DeclareUnicodeCharacter{2190}{\ensuremath{\leftarrow}}
\DeclareUnicodeCharacter{2194}{\ensuremath{\leftrightarrow}}
\DeclareUnicodeCharacter{21D2}{\ensuremath{\Rightarrow}}
\DeclareUnicodeCharacter{21D0}{\ensuremath{\Leftarrow}}
\DeclareUnicodeCharacter{21D4}{\ensuremath{\Leftrightarrow}}
\DeclareUnicodeCharacter{21A6}{\ensuremath{\mapsto}}
\DeclareUnicodeCharacter{2208}{\ensuremath{\in}}
\DeclareUnicodeCharacter{2209}{\ensuremath{\notin}}
\DeclareUnicodeCharacter{2286}{\ensuremath{\subseteq}}
\DeclareUnicodeCharacter{2282}{\ensuremath{\subset}}
\DeclareUnicodeCharacter{2287}{\ensuremath{\supseteq}}
\DeclareUnicodeCharacter{2283}{\ensuremath{\supset}}
\DeclareUnicodeCharacter{2229}{\ensuremath{\cap}}
\DeclareUnicodeCharacter{222A}{\ensuremath{\cup}}
\DeclareUnicodeCharacter{2216}{\ensuremath{\setminus}}
\DeclareUnicodeCharacter{2205}{\ensuremath{\emptyset}}
\DeclareUnicodeCharacter{2200}{\ensuremath{\forall}}
\DeclareUnicodeCharacter{2203}{\ensuremath{\exists}}
\DeclareUnicodeCharacter{2115}{\ensuremath{\mathbb{N}}}
\DeclareUnicodeCharacter{2124}{\ensuremath{\mathbb{Z}}}
\DeclareUnicodeCharacter{211A}{\ensuremath{\mathbb{Q}}}
\DeclareUnicodeCharacter{211D}{\ensuremath{\mathbb{R}}}
\DeclareUnicodeCharacter{2102}{\ensuremath{\mathbb{C}}}
\DeclareUnicodeCharacter{00B7}{\ensuremath{\cdot}}
\DeclareUnicodeCharacter{00D7}{\ensuremath{\times}}
\DeclareUnicodeCharacter{2212}{\ensuremath{-}}
\DeclareUnicodeCharacter{00B1}{\ensuremath{\pm}}
\DeclareUnicodeCharacter{2213}{\ensuremath{\mp}}
\DeclareUnicodeCharacter{221A}{\ensuremath{\surd}}
\DeclareUnicodeCharacter{2211}{\ensuremath{\sum}}
\DeclareUnicodeCharacter{220F}{\ensuremath{\prod}}
\DeclareUnicodeCharacter{222B}{\ensuremath{\int}}
\DeclareUnicodeCharacter{221E}{\ensuremath{\infty}}
\DeclareUnicodeCharacter{2202}{\ensuremath{\partial}}
\DeclareUnicodeCharacter{2207}{\ensuremath{\nabla}}
\DeclareUnicodeCharacter{2297}{\ensuremath{\otimes}}
\DeclareUnicodeCharacter{2295}{\ensuremath{\oplus}}
\DeclareUnicodeCharacter{2218}{\ensuremath{\circ}}
\DeclareUnicodeCharacter{2308}{\ensuremath{\lceil}}
\DeclareUnicodeCharacter{2309}{\ensuremath{\rceil}}
\DeclareUnicodeCharacter{230A}{\ensuremath{\lfloor}}
\DeclareUnicodeCharacter{230B}{\ensuremath{\rfloor}}
\DeclareUnicodeCharacter{27E8}{\ensuremath{\langle}}
\DeclareUnicodeCharacter{27E9}{\ensuremath{\rangle}}
\DeclareUnicodeCharacter{2016}{\ensuremath{\|}}
\DeclareUnicodeCharacter{2113}{\ensuremath{\ell}}
\DeclareUnicodeCharacter{00B0}{\ensuremath{^{\circ}}}
\DeclareUnicodeCharacter{00B2}{\ensuremath{^{2}}}
\DeclareUnicodeCharacter{00B3}{\ensuremath{^{3}}}
\DeclareUnicodeCharacter{00B9}{\ensuremath{^{1}}}
\DeclareUnicodeCharacter{2080}{\ensuremath{_{0}}}
\DeclareUnicodeCharacter{2081}{\ensuremath{_{1}}}
\DeclareUnicodeCharacter{2082}{\ensuremath{_{2}}}
\DeclareUnicodeCharacter{207F}{\ensuremath{^{n}}}
\DeclareUnicodeCharacter{2096}{\ensuremath{_{k}}}
%% --- 6. Numbering: one counter per section -----------------------
\numberwithin{equation}{section}
\setcounter{tocdepth}{1}
%% --- 7. The project macro set, then the claim environments -----------------
% spec/macros.tex, read relative to this module's directory (latexmk -cd).
% =============================================================================
%  spec/macros.tex -- the one macro set every generated contract inputs.
%
%  Two parts: part A, the macros spec/notation.md catalogues, and part B,
%  the macros local to one module, grouped by module.
%  harness/assemble_contract.py emits % =============================================================================
%  spec/macros.tex -- the one macro set every generated contract inputs.
%
%  Two parts: part A, the macros spec/notation.md catalogues, and part B,
%  the macros local to one module, grouped by module.
%  harness/assemble_contract.py emits % =============================================================================
%  spec/macros.tex -- the one macro set every generated contract inputs.
%
%  Two parts: part A, the macros spec/notation.md catalogues, and part B,
%  the macros local to one module, grouped by module.
%  harness/assemble_contract.py emits \input{../../macros} for it, so the
%  file is read from spec/modules/<id>/.
%
%  RULES
%    * A macro that spec/notation.md catalogues is a verbatim copy of it
%      (part A). Changing one is a notation change: bump spec/notation.md.
%    * Part B holds the module-local macros. A macro defined by two modules
%      with the same expansion appears once, under the first module (in the
%      order below) that defines it. Two macros with the same name and
%      DIFFERENT expansions are both kept under a module prefix, and the
%      claims write the prefixed names:
%        \Chain  -> \redChain (reductions, \mathfrak{C}: the reduction chain)
%                   \pcChain  (protocol-compiler, \mathfrak{R}: the two-segment
%                             chain; \Chsrc and \Chrec are built on it)
%        \Nmsm   -> \redNmsm  (reductions, \mathcal{N}: a POVM measurement
%                             channel)
%                   \pcNmsm   (\mathcal{N}_{\mathsf A}: Alice's measurement
%                             channel, written by
%                             transmitter/thm:certificate-sufficiency,
%                             transmitter/rem:xreg and
%                             reductions/lem:source-transfer)
%    * Deliberately omitted: a \provenance{tag}{details} macro. Provenance
%      is a header field of every claim file (spec/claim-format.md section
%      3); a claim body that calls such a macro fails to compile, which is
%      the intended signal.
%
%  REQUIRED PACKAGES, loaded by the assembler's preamble before this file:
%    amsmath, amssymb, amsfonts, mathtools, dsfont, braket, physics, xcolor
%    (braket before physics, so that physics wins for \ket, \bra, \tr, ...).
%
%  This file contains only definitions: no \usepackage, no document body.
% =============================================================================


% =============================================================================
%  PART A -- notation macros, the ones spec/notation.md catalogues, each a
%  verbatim copy of its catalogue row; grouped under numbered headings.
% =============================================================================

% -----------------------------------------------------------------------------
% 1. State spaces and operator sets (a "system label" -> spaces of operators)
% -----------------------------------------------------------------------------
\newcommand{\Hs}[1]{\mathcal{H}(#1)}        % Hilbert space of a system label
\newcommand{\Lin}[1]{\mathcal{L}(#1)}       % bounded linear operators on H(#1)
\newcommand{\Linn}[2]{\mathcal{L}(#1,#2)}   % linear maps H(#1) -> H(#2)
\newcommand{\Pos}[1]{\mathrm{Pos}(#1)}      % positive semidefinite operators on H(#1)
\newcommand{\Density}{\mathrm{D}}           % density operators: \Density(A) = D(A)
\newcommand{\Seq}[1]{S_{=}(#1)}             % normalized states (trace = 1)
\newcommand{\Sleq}[1]{S_{\leq}(#1)}         % subnormalized states (0 <= trace <= 1)
\newcommand{\id}{\mathds{1}}                % identity operator
\newcommand{\idon}[1]{\mathds{1}_{#1}}      % identity on a named system, e.g. \idon{C}

% DEPRECATED: applies \mathcal to both arguments itself, so it takes bare
% letters, not system labels. Prefer \Linn; do not use in new contracts.
\providecommand{\Lop}[2]{\mathcal{L}(\mathcal{#1}, \mathcal{#2})}


% -----------------------------------------------------------------------------
% 2. Composite systems
% -----------------------------------------------------------------------------
\newcommand{\dsum}{\oplus}                  % direct sum of system labels
% Tensor product: use \otimes directly.
% n-fold direct sum of an operator space: \Lin{A}^{\oplus n}


% -----------------------------------------------------------------------------
% 3. Events, conditioning and classical registers (for STATES)
%    Convention:  \wedge  -> subnormalized "and-event" component, e.g. \rho_{A\wedge x}
%                 |        -> normalized conditional state,        e.g. \rho_{A|x}
% (Write these directly in subscripts; no extra macros needed.)
% -----------------------------------------------------------------------------


% -----------------------------------------------------------------------------
% 4. Photon-number subspaces / direct-sum blocks
%    Convention:  operators carry the restriction "on top", in brackets;
%                 states carry it as a \wedge subscript (see Sec. 3).
% -----------------------------------------------------------------------------
\newcommand{\Pnum}[1]{P_{#1}}               % photon-number projector (system from context)
\newcommand{\Pnumon}[2]{P_{#1}^{#2}}        % with explicit system superscript, e.g. P_n^C
\newcommand{\blk}[2]{#1^{(#2)}}             % operator block, e.g. \blk{M}{\leqN} = M^{(<=N)}
\newcommand{\leqN}{\leq N}                  % use as \blk{M}{\leqN}
\newcommand{\gtN}{> N}                      % use as \blk{M}{\gtN}
\newcommand{\eqn}{= n}                      % use as \blk{M}{\eqn}


% -----------------------------------------------------------------------------
% 5. System labels used in the protocol
% -----------------------------------------------------------------------------
\newcommand{\Chan}{C}                       % quantum-channel (Eve) input to Bob
\newcommand{\Prb}{P}                        % probe input
\newcommand{\Flag}{{\mathrm{Flg}}}          % flag system (EXISTS ONLY AFTER squashing)
\newcommand{\Loss}{{\mathrm{Lss}}}          % loss system
\newcommand{\Vc}{{\mathrm{Vcm}}}            % vacuum system
\newcommand{\Nout}{N_{\mathrm{out}}}
\newcommand{\Nin}{N_{\mathrm{in}}}


% -----------------------------------------------------------------------------
% 6. Alphabets
% -----------------------------------------------------------------------------
\newcommand{\Yset}{\mathcal{Y}}             % Bob's (coarse-grained) detection alphabet
\newcommand{\Yraw}{\mathcal{Y}_{\mathrm{raw}}}% raw click patterns {0,1}^4
\newcommand{\Yscset}{\mathcal{Y}_{\singlec}}% single-click outcomes
\newcommand{\Ymcset}{\mathcal{Y}_{\mc}}     % multi-click outcomes
\newcommand{\Tset}{\mathcal{T}_{\mathrm{set}}}% private setting set (the probe settings in the optical instance)
\newcommand{\Xset}{\mathcal{X}}             % Alice's alphabet
\newcommand{\Dset}{\{H,V,+,-\}}             % polarization / detector labels
\newcommand{\pass}{\top}                    % pass sifting


% -----------------------------------------------------------------------------
% 7. Mode / outcome abbreviations
% -----------------------------------------------------------------------------
\newcommand{\on}{{\mathrm{on}}}
\newcommand{\off}{{\mathrm{off}}}
\newcommand{\dc}{{\mathrm{dc}}}             % dark count
\newcommand{\sq}{\mathrm{sq}}               % squashed
\newcommand{\mc}{\mathrm{mc}}               % multi-click
\newcommand{\nc}{\mathrm{nc}}               % no-click
\newcommand{\singlec}{\mathrm{sc}}          % single-click
\newcommand{\vac}{\mathrm{vac}}             % vacuum
\newcommand{\kout}{{k_\mathrm{out}}}
\newcommand{\kin}{{k_\mathrm{in}}}


% -----------------------------------------------------------------------------
% 8. Maps, channels and miscellaneous operators
% -----------------------------------------------------------------------------
\newcommand{\Prob}{\mathrm{Prob}}           % probability (use \Prob[\cdot] or \Prob(\cdot))
\newcommand{\lmin}{\lambda_{\min}}          % minimum eigenvalue
\newcommand{\Wflag}{W_{\Flag}}              % flag-space weight
\newcommand{\transpose}{\mathsf{T}}
\newcommand{\innerp}[2]{\left\langle #1 \vert #2 \right\rangle}
\renewcommand{\vec}[1]{\boldsymbol{#1}}     % NB: overrides the LaTeX arrow accent


% -----------------------------------------------------------------------------
% 9. Key-rate numerics
% -----------------------------------------------------------------------------
\newcommand{\cG}{\mathcal{G}}               % protocol map
\newcommand{\cZ}{\mathcal{Z}}               % pinching map
\newcommand{\Amap}{\mathcal{A}_{\mathrm{ann}}} % announcement + measurement map
\newcommand{\Amapnb}{\mathcal{A}_{\mathrm{ann}|\mathrm{nb}}}
\newcommand{\simpl}{}                       % empty suffix for simplified block maps
\newcommand{\acSet}{\mathcal{S}_{\mathrm{acc}}} % acceptance set
\newcommand{\testE}{\mathcal{E}_{\mathrm{test}}} % test-statistics map
\newcommand{\ecLeak}{\mathrm{leak}^{\mathrm{EC}}_{\mathrm{obs}}}
\newcommand{\pgen}{p_{\mathrm{gen}}}
\newcommand{\askeyr}{R_{\infty}}             % asymptotic key rate
\newcommand{\ppass}{p_{\mathrm{pass}}}       % sifting pass probability


% -----------------------------------------------------------------------------
% 10. Finite-size analysis / Renyi entropies
% -----------------------------------------------------------------------------

% --- Round types: test vs. key generation (register T, alphabet \Fset) --------
\newcommand{\tst}{\mathrm{test}}            % test round
\newcommand{\gen}{\mathrm{gen}}             % key-generation round
\newcommand{\Fset}{\mathcal{T}}             % round-type alphabet {test, gen} (not \Tset)
\newcommand{\Chat}{\widehat{C}}             % acceptance-test register
\newcommand{\Chatset}{\widehat{\mathcal{C}}}% its alphabet

% --- Basis / key extraction functions ----------------------------------------
\newcommand{\bA}{b_{A}}                     % Alice's basis-extraction function
\newcommand{\bB}{b_{B}}                     % Bob's basis-extraction function
\newcommand{\kB}{k_{B}}                     % Bob's key-extraction function

% --- Maps --------------------------------------------------------------------
\newcommand{\Mprot}{\mathcal{M}}            % single-round protocol map
\newcommand{\Phimap}{\boldsymbol{\Phi}}     % test-statistics map
\newcommand{\Pdeph}{\mathcal{P}}            % dephasing (pinching) map on \Blreg

% --- Isometries --------------------------------------------------------------
\newcommand{\Vmsm}{V_{\mathrm{msm}}}        % measurement isometry
\newcommand{\Vann}{V_{\mathrm{ann}}}        % announcement isometry
\newcommand{\VI}{V_{I}}                     % Eve's copy of the announcements
\newcommand{\VChat}{V_{\Chat}}              % acceptance-data isometry
\newcommand{\VR}{V_{R}}                     % sifting + key map -> raw key R
\newcommand{\VZ}{V_{Z}}                     % raw key -> key register Z
\newcommand{\Wprot}{W}                      % full protocol isometry

% --- Renyi entropies ---------------------------------------------------------
% SECURITY-CRITICAL: see spec/notation.md §4. Never change an arrow to
% make an inequality work out; never convert silently between the sandwiched
% (tilde) and Petz (overline) families.
\newcommand{\Hsandup}{\widetilde{H}^{\uparrow}_{\alpha}}   % sandwiched, up
\newcommand{\Hsanddown}{\widetilde{H}^{\downarrow}_{\alpha}}% sandwiched, down
\newcommand{\Hpetzup}{\overline{H}^{\uparrow}_{\alpha}}    % Petz, up
\newcommand{\Hpetzdown}{\overline{H}^{\downarrow}_{\alpha}}% Petz, down
\newcommand{\hup}{h^{\uparrow}_{\alpha}}                   % single-round rate
\newcommand{\hqkd}{h^{\uparrow,\mathrm{QKD}}_{\alpha}}     % single-round QKD rate
\newcommand{\hlow}{\underline{h}_{\alpha}}                 % lower bound on \hqkd
% --- spec/notation.md 3.5 and 4.5 -------------------------------------------
% \hlb  is the first-order linearization lower bound on \hlow: the headline
%       key rate (the `first-order` tier of foundations/def:tiers). NOT \hlow
%       itself.
% \tnum is the instance's one declared numerical tolerance. Its
%       near-collision with the declared marginal \tau_A is disclosed in
%       spec/notation.md 3.5; bare \tau is reserved and denotes neither.
\newcommand{\hlb}{\underline{h}^{\mathrm{FO}}_{\alpha}}    % first-order lower bound on \hlow
\newcommand{\tnum}{\tau_{\mathrm{num}}}                    % declared numerical tolerance


% =============================================================================
%  PART B -- module-local macros, grouped by module. Deprecated aliases,
%  defined so that a text written with them compiles, used in no new contract: \Lop,
%  \testE, \Pdeph, \pass, \ecLeak (spec/notation.md section 3.9).
% =============================================================================

% --- meat-core -----------------------------------------------------------------
\newcommand{\SecSet}{\mathbb{S}}            % sector index set of Q
\newcommand{\Qsys}{Q}                       % declared sector-carrying measured system
\newcommand{\Vset}{\mathcal{V}}             % declared joint outcome alphabet of the Q-POVM
\newcommand{\CAset}{\mathcal{C}_{A}}        % Alice announcement alphabet
\newcommand{\CBset}{\mathcal{C}_{B}}        % Bob announcement alphabet
\newcommand{\Party}{\mathfrak{P}}           % declared party partition
\newcommand{\Oset}[1]{\mathcal{O}_{#1}}     % readout alphabet of party #1
\newcommand{\rd}[1]{\pi_{#1}}               % readout map of party #1
\newcommand{\Sadm}{\mathcal{S}_{A}}         % admissible Alice subspace
\newcommand{\nomarg}{\mathsf{none}}         % "no trusted marginal declared"
\newcommand{\Zset}{\mathcal{Z}_{\perp}}     % raw-key alphabet incl. \perp
\newcommand{\Zval}{\mathcal{Z}}             % key alphabet without \perp
\newcommand{\Etil}{\widetilde{E}}           % single-round purifying register
\newcommand{\Ehat}{\widehat{E}}             % Eve's n-round register
\newcommand{\OmAT}{\Omega_{\mathrm{AT}}}    % acceptance event
\newcommand{\epa}{\varepsilon_{\mathrm{PA}}}
\newcommand{\eev}{\varepsilon_{\mathrm{EV}}}
\newcommand{\lEC}{\lambda_{\mathrm{EC}}}
\newcommand{\Decl}{\mathfrak{D}}            % protocol declaration
\newcommand{\CList}{\mathfrak{L}}           % declared constraint list
\newcommand{\PinS}{\mathcal{P}_{\SecSet}}   % sector pinching on Q
\newcommand{\Feas}{\mathcal{W}}             % feasible set of block states
\newcommand{\Chkey}{\Chatset_{\mathrm{key}}}% key-live acceptance values
\newcommand{\Kann}{\mathcal{K}}             % announcement-class alphabet
\newcommand{\hold}{\mathsf{h}}              % declared key holder (A or B)
\newcommand{\sfD}{D}                        % classical KL divergence (base 2)
\newcommand{\Dtwo}{D}                       % quantum relative entropy (base 2)
\newcommand{\Dsand}{\widetilde{D}_{\alpha}} % sandwiched Renyi divergence
\newcommand{\idmap}{\mathrm{id}}            % identity channel
\newcommand{\Aev}{\mathcal{A}}              % Eve's per-round attack channel
\newcommand{\AttSet}{\mathfrak{A}}          % licensed attack-channel set
\newcommand{\Chacc}{\Chatset_{\mathrm{acc}}}% acceptance support of \acSet
\newcommand{\epsstar}{\varepsilon^{\star}}      % distance to an exact POVM
\newcommand{\dfH}{\Delta_{\mathrm{herm}}}       % Hermiticity defect
\newcommand{\dfP}{\Delta_{\mathrm{pos}}}        % negative part
\newcommand{\dfC}{\Delta_{\mathrm{cpl}}}        % completeness defect
\newcommand{\lmoq}{g}                       % acceptance-side gradient block
\newcommand{\lmoJ}{G}                       % feasible-side gradient block
\newcommand{\bacc}{\beta_{\mathrm{acc}}}    % certified acceptance-side minimum
\newcommand{\Itset}{\mathfrak{I}}           % finite index set of iterates (C22)

% --- reductions ----------------------------------------------------------------
\newcommand{\Bsys}{B}                        % physical (pre-squashing) system
\newcommand{\Tcl}[1]{\mathcal{T}(#1)}        % trace-class operators on H(#1)
\newcommand{\Ins}{\mathcal{I}}               % generalized-instrument branch map
\newcommand{\redNmsm}{\mathcal{N}}              % POVM measurement channel
\newcommand{\Eprot}{\mathcal{E}}             % n-round protocol channel
\newcommand{\Grnd}{\mathsf{G}}               % one round map of a protocol
\newcommand{\Kfin}{\mathsf{K}}               % end-of-protocol processing
\newcommand{\Crit}{\mathsf{sec}}             % value of a security criterion
\newcommand{\Pflag}{\Pi_{\Flag}}             % flag-sector projector on Q
\newcommand{\Pleq}{\Pi_{\leqN}}              % preserved-block projector
\newcommand{\Pgt}{\Pi_{\gtN}}                % above-cutoff projector
\newcommand{\Sph}{S^{\mathrm{ph}}}           % family of physical inputs
\newcommand{\Srch}{S^{\mathrm{rch}}}         % reachable squashed states
\newcommand{\Ssq}{S^{\sq}}                   % squashed admissible set
\newcommand{\Ffss}{F^{\mathrm{FSS}}}         % flag-state-squasher target
\newcommand{\Fwp}{F^{\mathrm{WP}}}           % multi-bound WPFSS target
\newcommand{\Mwp}{M^{\mathrm{WP}}}           % renormalized above-cutoff POVM
\newcommand{\Mdec}{{M'}^{\mathrm{ph}}}       % declared physical family on B
\newcommand{\Ppvm}{P^{\Flag}}                % flag PVM
\newcommand{\Lmp}{{\Lambda_{\mathrm{mp}}}}   % measure-and-prepare channel
\newcommand{\Lwp}{{\Lambda^{\mathrm{WP}}}}   % multi-bound WPFSS channel
\newcommand{\Lfss}{{\Lambda^{\mathrm{FSS}}}} % flag-state-squasher channel
\newcommand{\lmax}{\lambda_{\max}}           % maximum eigenvalue
\newcommand{\Lnoise}{\widetilde{\Lambda}}    % flag-noise channel
\newcommand{\cg}{\varpi}                     % declared coarse-graining map
\newcommand{\MCK}{\mathsf{MC}}               % multiclick raw-pattern set
\newcommand{\redChain}{\mathfrak{C}}            % reduction chain
\newcommand{\Lnk}{\ell}                      % a single typed link
\newcommand{\dsq}{\preceq}                      % the squashing relation
\newcommand{\opn}[1]{\left\lVert #1\right\rVert}           % operator (spectral) norm
\newcommand{\mxn}[1]{\left\lVert #1\right\rVert_{\max}}    % entrywise max modulus
\newcommand{\Fld}{\mathbb{F}}                   % exact ordered substrate field
\newcommand{\Fldi}{\mathbb{F}_{i}}              % its Gaussian extension F + iF
\newcommand{\lat}[1]{\big\lfloor #1\big\rfloor_{\kappa}} % round DOWN onto the lattice

% --- protocol-compiler ---------------------------------------------------------
\newcommand{\Spec}{\mathfrak{S}}             % protocol specification (compiler input)
\newcommand{\pcChain}{\mathfrak{R}}            % reduction chain (two typed segments)
\newcommand{\Chsrc}{\pcChain_{\mathrm{src}}}   % its Alice-side source-map segment
\newcommand{\Chrec}{\pcChain_{\mathrm{rec}}}   % its receiver-side squashing segment
\newcommand{\rvirt}{\rho^{\mathrm{v}}}       % virtual (PRE-source-map) signal state
\newcommand{\Fph}{F^{\mathrm{ph}}}           % physical (pre-reduction) POVM element
\newcommand{\Vph}{\mathcal{V}^{\mathrm{ph}}} % physical outcome alphabet
\newcommand{\Tpub}{\mathcal{T}_{\mathrm{pub}}}% public-choice alphabet
\newcommand{\Uset}{\mathcal{U}}              % residual outcome alphabet
\newcommand{\idxm}{\mathsf{idx}}             % label-to-index map
\newcommand{\Asig}{A'}                       % transmitted (signal) system -- the REAL one
\newcommand{\Avirt}{A''}                     % virtual transmitted system (source descriptor of Xi_1)
\newcommand{\Abar}{\bar{A}}                  % Alice's measured register
\newcommand{\Ahat}{\hat{A}}                  % Alice's shield register
\newcommand{\pcNmsm}{\mathcal{N}_{\mathsf A}}  % Alice's measurement channel (see the RULES block)
\newcommand{\Phip}{\Phi^{+}}                 % maximally entangled vector label
\newcommand{\Smap}{\mathfrak{s}}             % the emitted sector map (fraktur: a compile-time object)

% --- receiver ------------------------------------------------------------------
\newcommand{\Modes}{\mathbb{M}}              % declared ordered mode set
\newcommand{\Inp}{\mathbb{I}}                % declared input modes
\newcommand{\InpC}{\mathbb{I}_{\Chan}}       % channel input modes
\newcommand{\InpS}{\mathbb{I}_{\Prb}}        % probe (source) input modes
\newcommand{\Env}{\mathbb{E}}                % dilation (environment) modes
\newcommand{\Port}[1]{\mathbb{P}_{#1}}       % mode set of detector port #1
\newcommand{\Emd}[1]{E_{#1}}                 % coordinate projector on modes #1
\newcommand{\Gmat}{G}                        % instrument (transfer) matrix
\newcommand{\Fock}{\mathfrak{F}}             % Fock / second-quantization functor
\newcommand{\Symn}[1]{\mathrm{Sym}^{#1}}     % symmetric tensor power
\newcommand{\normord}[1]{{:}\mspace{2mu}#1\mspace{2mu}{:}}   % normal ordering
\newcommand{\rk}{\operatorname{rank}}        % rank (physics.sty owns \rank)
\newcommand{\SCK}{\mathsf{SC}}               % single-click pattern set
\newcommand{\Recv}{\mathfrak{R}}             % receiver instrument tuple
\newcommand{\Cat}{\mathfrak{C}}              % component catalogue
\newcommand{\Grph}{\mathfrak{E}}             % declared element graph
\newcommand{\defic}{\mathbb{L}}              % loss deficiency operator
\newcommand{\Pvac}[1]{\Pi_{\vac}(#1)}        % joint no-click element on ports #1
\newcommand{\pmc}{p_{\mc}}                   % uniform high-photon mc weight
\newcommand{\Bnd}{\beta}                     % the scalar bound function beta(n)
\newcommand{\Rot}{R}                         % planar rotation matrix
\newcommand{\Wplate}{W}                      % retarder (wave-plate) matrix
\newcommand{\Bsp}{B}                         % beamsplitter core matrix

% --- statistics ----------------------------------------------------------------
\newcommand{\honest}{p^{\mathrm{hon}}}          % honest per-round distribution
\newcommand{\ecomp}{\varepsilon_{\mathrm{comp}}}% completeness budget
\newcommand{\simplex}[1]{\Delta(#1)}            % probability simplex
\newcommand{\Chatcon}{\Chatset_{\mathrm{con}}}  % constrained outcomes
\newcommand{\Ireg}[3]{\mathcal{I}_{#1}(#2,#3)}  % regularized incomplete beta
\newcommand{\Ncount}[1]{\mathsf{N}_{#1}}        % per-outcome count
\newcommand{\qlo}{\mathbf{q}^{-}}               % box lower bounds
\newcommand{\qhi}{\mathbf{q}^{+}}               % box upper bounds
 for it, so the
%  file is read from spec/modules/<id>/.
%
%  RULES
%    * A macro that spec/notation.md catalogues is a verbatim copy of it
%      (part A). Changing one is a notation change: bump spec/notation.md.
%    * Part B holds the module-local macros. A macro defined by two modules
%      with the same expansion appears once, under the first module (in the
%      order below) that defines it. Two macros with the same name and
%      DIFFERENT expansions are both kept under a module prefix, and the
%      claims write the prefixed names:
%        \Chain  -> \redChain (reductions, \mathfrak{C}: the reduction chain)
%                   \pcChain  (protocol-compiler, \mathfrak{R}: the two-segment
%                             chain; \Chsrc and \Chrec are built on it)
%        \Nmsm   -> \redNmsm  (reductions, \mathcal{N}: a POVM measurement
%                             channel)
%                   \pcNmsm   (\mathcal{N}_{\mathsf A}: Alice's measurement
%                             channel, written by
%                             transmitter/thm:certificate-sufficiency,
%                             transmitter/rem:xreg and
%                             reductions/lem:source-transfer)
%    * Deliberately omitted: a \provenance{tag}{details} macro. Provenance
%      is a header field of every claim file (spec/claim-format.md section
%      3); a claim body that calls such a macro fails to compile, which is
%      the intended signal.
%
%  REQUIRED PACKAGES, loaded by the assembler's preamble before this file:
%    amsmath, amssymb, amsfonts, mathtools, dsfont, braket, physics, xcolor
%    (braket before physics, so that physics wins for \ket, \bra, \tr, ...).
%
%  This file contains only definitions: no \usepackage, no document body.
% =============================================================================


% =============================================================================
%  PART A -- notation macros, the ones spec/notation.md catalogues, each a
%  verbatim copy of its catalogue row; grouped under numbered headings.
% =============================================================================

% -----------------------------------------------------------------------------
% 1. State spaces and operator sets (a "system label" -> spaces of operators)
% -----------------------------------------------------------------------------
\newcommand{\Hs}[1]{\mathcal{H}(#1)}        % Hilbert space of a system label
\newcommand{\Lin}[1]{\mathcal{L}(#1)}       % bounded linear operators on H(#1)
\newcommand{\Linn}[2]{\mathcal{L}(#1,#2)}   % linear maps H(#1) -> H(#2)
\newcommand{\Pos}[1]{\mathrm{Pos}(#1)}      % positive semidefinite operators on H(#1)
\newcommand{\Density}{\mathrm{D}}           % density operators: \Density(A) = D(A)
\newcommand{\Seq}[1]{S_{=}(#1)}             % normalized states (trace = 1)
\newcommand{\Sleq}[1]{S_{\leq}(#1)}         % subnormalized states (0 <= trace <= 1)
\newcommand{\id}{\mathds{1}}                % identity operator
\newcommand{\idon}[1]{\mathds{1}_{#1}}      % identity on a named system, e.g. \idon{C}

% DEPRECATED: applies \mathcal to both arguments itself, so it takes bare
% letters, not system labels. Prefer \Linn; do not use in new contracts.
\providecommand{\Lop}[2]{\mathcal{L}(\mathcal{#1}, \mathcal{#2})}


% -----------------------------------------------------------------------------
% 2. Composite systems
% -----------------------------------------------------------------------------
\newcommand{\dsum}{\oplus}                  % direct sum of system labels
% Tensor product: use \otimes directly.
% n-fold direct sum of an operator space: \Lin{A}^{\oplus n}


% -----------------------------------------------------------------------------
% 3. Events, conditioning and classical registers (for STATES)
%    Convention:  \wedge  -> subnormalized "and-event" component, e.g. \rho_{A\wedge x}
%                 |        -> normalized conditional state,        e.g. \rho_{A|x}
% (Write these directly in subscripts; no extra macros needed.)
% -----------------------------------------------------------------------------


% -----------------------------------------------------------------------------
% 4. Photon-number subspaces / direct-sum blocks
%    Convention:  operators carry the restriction "on top", in brackets;
%                 states carry it as a \wedge subscript (see Sec. 3).
% -----------------------------------------------------------------------------
\newcommand{\Pnum}[1]{P_{#1}}               % photon-number projector (system from context)
\newcommand{\Pnumon}[2]{P_{#1}^{#2}}        % with explicit system superscript, e.g. P_n^C
\newcommand{\blk}[2]{#1^{(#2)}}             % operator block, e.g. \blk{M}{\leqN} = M^{(<=N)}
\newcommand{\leqN}{\leq N}                  % use as \blk{M}{\leqN}
\newcommand{\gtN}{> N}                      % use as \blk{M}{\gtN}
\newcommand{\eqn}{= n}                      % use as \blk{M}{\eqn}


% -----------------------------------------------------------------------------
% 5. System labels used in the protocol
% -----------------------------------------------------------------------------
\newcommand{\Chan}{C}                       % quantum-channel (Eve) input to Bob
\newcommand{\Prb}{P}                        % probe input
\newcommand{\Flag}{{\mathrm{Flg}}}          % flag system (EXISTS ONLY AFTER squashing)
\newcommand{\Loss}{{\mathrm{Lss}}}          % loss system
\newcommand{\Vc}{{\mathrm{Vcm}}}            % vacuum system
\newcommand{\Nout}{N_{\mathrm{out}}}
\newcommand{\Nin}{N_{\mathrm{in}}}


% -----------------------------------------------------------------------------
% 6. Alphabets
% -----------------------------------------------------------------------------
\newcommand{\Yset}{\mathcal{Y}}             % Bob's (coarse-grained) detection alphabet
\newcommand{\Yraw}{\mathcal{Y}_{\mathrm{raw}}}% raw click patterns {0,1}^4
\newcommand{\Yscset}{\mathcal{Y}_{\singlec}}% single-click outcomes
\newcommand{\Ymcset}{\mathcal{Y}_{\mc}}     % multi-click outcomes
\newcommand{\Tset}{\mathcal{T}_{\mathrm{set}}}% private setting set (the probe settings in the optical instance)
\newcommand{\Xset}{\mathcal{X}}             % Alice's alphabet
\newcommand{\Dset}{\{H,V,+,-\}}             % polarization / detector labels
\newcommand{\pass}{\top}                    % pass sifting


% -----------------------------------------------------------------------------
% 7. Mode / outcome abbreviations
% -----------------------------------------------------------------------------
\newcommand{\on}{{\mathrm{on}}}
\newcommand{\off}{{\mathrm{off}}}
\newcommand{\dc}{{\mathrm{dc}}}             % dark count
\newcommand{\sq}{\mathrm{sq}}               % squashed
\newcommand{\mc}{\mathrm{mc}}               % multi-click
\newcommand{\nc}{\mathrm{nc}}               % no-click
\newcommand{\singlec}{\mathrm{sc}}          % single-click
\newcommand{\vac}{\mathrm{vac}}             % vacuum
\newcommand{\kout}{{k_\mathrm{out}}}
\newcommand{\kin}{{k_\mathrm{in}}}


% -----------------------------------------------------------------------------
% 8. Maps, channels and miscellaneous operators
% -----------------------------------------------------------------------------
\newcommand{\Prob}{\mathrm{Prob}}           % probability (use \Prob[\cdot] or \Prob(\cdot))
\newcommand{\lmin}{\lambda_{\min}}          % minimum eigenvalue
\newcommand{\Wflag}{W_{\Flag}}              % flag-space weight
\newcommand{\transpose}{\mathsf{T}}
\newcommand{\innerp}[2]{\left\langle #1 \vert #2 \right\rangle}
\renewcommand{\vec}[1]{\boldsymbol{#1}}     % NB: overrides the LaTeX arrow accent


% -----------------------------------------------------------------------------
% 9. Key-rate numerics
% -----------------------------------------------------------------------------
\newcommand{\cG}{\mathcal{G}}               % protocol map
\newcommand{\cZ}{\mathcal{Z}}               % pinching map
\newcommand{\Amap}{\mathcal{A}_{\mathrm{ann}}} % announcement + measurement map
\newcommand{\Amapnb}{\mathcal{A}_{\mathrm{ann}|\mathrm{nb}}}
\newcommand{\simpl}{}                       % empty suffix for simplified block maps
\newcommand{\acSet}{\mathcal{S}_{\mathrm{acc}}} % acceptance set
\newcommand{\testE}{\mathcal{E}_{\mathrm{test}}} % test-statistics map
\newcommand{\ecLeak}{\mathrm{leak}^{\mathrm{EC}}_{\mathrm{obs}}}
\newcommand{\pgen}{p_{\mathrm{gen}}}
\newcommand{\askeyr}{R_{\infty}}             % asymptotic key rate
\newcommand{\ppass}{p_{\mathrm{pass}}}       % sifting pass probability


% -----------------------------------------------------------------------------
% 10. Finite-size analysis / Renyi entropies
% -----------------------------------------------------------------------------

% --- Round types: test vs. key generation (register T, alphabet \Fset) --------
\newcommand{\tst}{\mathrm{test}}            % test round
\newcommand{\gen}{\mathrm{gen}}             % key-generation round
\newcommand{\Fset}{\mathcal{T}}             % round-type alphabet {test, gen} (not \Tset)
\newcommand{\Chat}{\widehat{C}}             % acceptance-test register
\newcommand{\Chatset}{\widehat{\mathcal{C}}}% its alphabet

% --- Basis / key extraction functions ----------------------------------------
\newcommand{\bA}{b_{A}}                     % Alice's basis-extraction function
\newcommand{\bB}{b_{B}}                     % Bob's basis-extraction function
\newcommand{\kB}{k_{B}}                     % Bob's key-extraction function

% --- Maps --------------------------------------------------------------------
\newcommand{\Mprot}{\mathcal{M}}            % single-round protocol map
\newcommand{\Phimap}{\boldsymbol{\Phi}}     % test-statistics map
\newcommand{\Pdeph}{\mathcal{P}}            % dephasing (pinching) map on \Blreg

% --- Isometries --------------------------------------------------------------
\newcommand{\Vmsm}{V_{\mathrm{msm}}}        % measurement isometry
\newcommand{\Vann}{V_{\mathrm{ann}}}        % announcement isometry
\newcommand{\VI}{V_{I}}                     % Eve's copy of the announcements
\newcommand{\VChat}{V_{\Chat}}              % acceptance-data isometry
\newcommand{\VR}{V_{R}}                     % sifting + key map -> raw key R
\newcommand{\VZ}{V_{Z}}                     % raw key -> key register Z
\newcommand{\Wprot}{W}                      % full protocol isometry

% --- Renyi entropies ---------------------------------------------------------
% SECURITY-CRITICAL: see spec/notation.md §4. Never change an arrow to
% make an inequality work out; never convert silently between the sandwiched
% (tilde) and Petz (overline) families.
\newcommand{\Hsandup}{\widetilde{H}^{\uparrow}_{\alpha}}   % sandwiched, up
\newcommand{\Hsanddown}{\widetilde{H}^{\downarrow}_{\alpha}}% sandwiched, down
\newcommand{\Hpetzup}{\overline{H}^{\uparrow}_{\alpha}}    % Petz, up
\newcommand{\Hpetzdown}{\overline{H}^{\downarrow}_{\alpha}}% Petz, down
\newcommand{\hup}{h^{\uparrow}_{\alpha}}                   % single-round rate
\newcommand{\hqkd}{h^{\uparrow,\mathrm{QKD}}_{\alpha}}     % single-round QKD rate
\newcommand{\hlow}{\underline{h}_{\alpha}}                 % lower bound on \hqkd
% --- spec/notation.md 3.5 and 4.5 -------------------------------------------
% \hlb  is the first-order linearization lower bound on \hlow: the headline
%       key rate (the `first-order` tier of foundations/def:tiers). NOT \hlow
%       itself.
% \tnum is the instance's one declared numerical tolerance. Its
%       near-collision with the declared marginal \tau_A is disclosed in
%       spec/notation.md 3.5; bare \tau is reserved and denotes neither.
\newcommand{\hlb}{\underline{h}^{\mathrm{FO}}_{\alpha}}    % first-order lower bound on \hlow
\newcommand{\tnum}{\tau_{\mathrm{num}}}                    % declared numerical tolerance


% =============================================================================
%  PART B -- module-local macros, grouped by module. Deprecated aliases,
%  defined so that a text written with them compiles, used in no new contract: \Lop,
%  \testE, \Pdeph, \pass, \ecLeak (spec/notation.md section 3.9).
% =============================================================================

% --- meat-core -----------------------------------------------------------------
\newcommand{\SecSet}{\mathbb{S}}            % sector index set of Q
\newcommand{\Qsys}{Q}                       % declared sector-carrying measured system
\newcommand{\Vset}{\mathcal{V}}             % declared joint outcome alphabet of the Q-POVM
\newcommand{\CAset}{\mathcal{C}_{A}}        % Alice announcement alphabet
\newcommand{\CBset}{\mathcal{C}_{B}}        % Bob announcement alphabet
\newcommand{\Party}{\mathfrak{P}}           % declared party partition
\newcommand{\Oset}[1]{\mathcal{O}_{#1}}     % readout alphabet of party #1
\newcommand{\rd}[1]{\pi_{#1}}               % readout map of party #1
\newcommand{\Sadm}{\mathcal{S}_{A}}         % admissible Alice subspace
\newcommand{\nomarg}{\mathsf{none}}         % "no trusted marginal declared"
\newcommand{\Zset}{\mathcal{Z}_{\perp}}     % raw-key alphabet incl. \perp
\newcommand{\Zval}{\mathcal{Z}}             % key alphabet without \perp
\newcommand{\Etil}{\widetilde{E}}           % single-round purifying register
\newcommand{\Ehat}{\widehat{E}}             % Eve's n-round register
\newcommand{\OmAT}{\Omega_{\mathrm{AT}}}    % acceptance event
\newcommand{\epa}{\varepsilon_{\mathrm{PA}}}
\newcommand{\eev}{\varepsilon_{\mathrm{EV}}}
\newcommand{\lEC}{\lambda_{\mathrm{EC}}}
\newcommand{\Decl}{\mathfrak{D}}            % protocol declaration
\newcommand{\CList}{\mathfrak{L}}           % declared constraint list
\newcommand{\PinS}{\mathcal{P}_{\SecSet}}   % sector pinching on Q
\newcommand{\Feas}{\mathcal{W}}             % feasible set of block states
\newcommand{\Chkey}{\Chatset_{\mathrm{key}}}% key-live acceptance values
\newcommand{\Kann}{\mathcal{K}}             % announcement-class alphabet
\newcommand{\hold}{\mathsf{h}}              % declared key holder (A or B)
\newcommand{\sfD}{D}                        % classical KL divergence (base 2)
\newcommand{\Dtwo}{D}                       % quantum relative entropy (base 2)
\newcommand{\Dsand}{\widetilde{D}_{\alpha}} % sandwiched Renyi divergence
\newcommand{\idmap}{\mathrm{id}}            % identity channel
\newcommand{\Aev}{\mathcal{A}}              % Eve's per-round attack channel
\newcommand{\AttSet}{\mathfrak{A}}          % licensed attack-channel set
\newcommand{\Chacc}{\Chatset_{\mathrm{acc}}}% acceptance support of \acSet
\newcommand{\epsstar}{\varepsilon^{\star}}      % distance to an exact POVM
\newcommand{\dfH}{\Delta_{\mathrm{herm}}}       % Hermiticity defect
\newcommand{\dfP}{\Delta_{\mathrm{pos}}}        % negative part
\newcommand{\dfC}{\Delta_{\mathrm{cpl}}}        % completeness defect
\newcommand{\lmoq}{g}                       % acceptance-side gradient block
\newcommand{\lmoJ}{G}                       % feasible-side gradient block
\newcommand{\bacc}{\beta_{\mathrm{acc}}}    % certified acceptance-side minimum
\newcommand{\Itset}{\mathfrak{I}}           % finite index set of iterates (C22)

% --- reductions ----------------------------------------------------------------
\newcommand{\Bsys}{B}                        % physical (pre-squashing) system
\newcommand{\Tcl}[1]{\mathcal{T}(#1)}        % trace-class operators on H(#1)
\newcommand{\Ins}{\mathcal{I}}               % generalized-instrument branch map
\newcommand{\redNmsm}{\mathcal{N}}              % POVM measurement channel
\newcommand{\Eprot}{\mathcal{E}}             % n-round protocol channel
\newcommand{\Grnd}{\mathsf{G}}               % one round map of a protocol
\newcommand{\Kfin}{\mathsf{K}}               % end-of-protocol processing
\newcommand{\Crit}{\mathsf{sec}}             % value of a security criterion
\newcommand{\Pflag}{\Pi_{\Flag}}             % flag-sector projector on Q
\newcommand{\Pleq}{\Pi_{\leqN}}              % preserved-block projector
\newcommand{\Pgt}{\Pi_{\gtN}}                % above-cutoff projector
\newcommand{\Sph}{S^{\mathrm{ph}}}           % family of physical inputs
\newcommand{\Srch}{S^{\mathrm{rch}}}         % reachable squashed states
\newcommand{\Ssq}{S^{\sq}}                   % squashed admissible set
\newcommand{\Ffss}{F^{\mathrm{FSS}}}         % flag-state-squasher target
\newcommand{\Fwp}{F^{\mathrm{WP}}}           % multi-bound WPFSS target
\newcommand{\Mwp}{M^{\mathrm{WP}}}           % renormalized above-cutoff POVM
\newcommand{\Mdec}{{M'}^{\mathrm{ph}}}       % declared physical family on B
\newcommand{\Ppvm}{P^{\Flag}}                % flag PVM
\newcommand{\Lmp}{{\Lambda_{\mathrm{mp}}}}   % measure-and-prepare channel
\newcommand{\Lwp}{{\Lambda^{\mathrm{WP}}}}   % multi-bound WPFSS channel
\newcommand{\Lfss}{{\Lambda^{\mathrm{FSS}}}} % flag-state-squasher channel
\newcommand{\lmax}{\lambda_{\max}}           % maximum eigenvalue
\newcommand{\Lnoise}{\widetilde{\Lambda}}    % flag-noise channel
\newcommand{\cg}{\varpi}                     % declared coarse-graining map
\newcommand{\MCK}{\mathsf{MC}}               % multiclick raw-pattern set
\newcommand{\redChain}{\mathfrak{C}}            % reduction chain
\newcommand{\Lnk}{\ell}                      % a single typed link
\newcommand{\dsq}{\preceq}                      % the squashing relation
\newcommand{\opn}[1]{\left\lVert #1\right\rVert}           % operator (spectral) norm
\newcommand{\mxn}[1]{\left\lVert #1\right\rVert_{\max}}    % entrywise max modulus
\newcommand{\Fld}{\mathbb{F}}                   % exact ordered substrate field
\newcommand{\Fldi}{\mathbb{F}_{i}}              % its Gaussian extension F + iF
\newcommand{\lat}[1]{\big\lfloor #1\big\rfloor_{\kappa}} % round DOWN onto the lattice

% --- protocol-compiler ---------------------------------------------------------
\newcommand{\Spec}{\mathfrak{S}}             % protocol specification (compiler input)
\newcommand{\pcChain}{\mathfrak{R}}            % reduction chain (two typed segments)
\newcommand{\Chsrc}{\pcChain_{\mathrm{src}}}   % its Alice-side source-map segment
\newcommand{\Chrec}{\pcChain_{\mathrm{rec}}}   % its receiver-side squashing segment
\newcommand{\rvirt}{\rho^{\mathrm{v}}}       % virtual (PRE-source-map) signal state
\newcommand{\Fph}{F^{\mathrm{ph}}}           % physical (pre-reduction) POVM element
\newcommand{\Vph}{\mathcal{V}^{\mathrm{ph}}} % physical outcome alphabet
\newcommand{\Tpub}{\mathcal{T}_{\mathrm{pub}}}% public-choice alphabet
\newcommand{\Uset}{\mathcal{U}}              % residual outcome alphabet
\newcommand{\idxm}{\mathsf{idx}}             % label-to-index map
\newcommand{\Asig}{A'}                       % transmitted (signal) system -- the REAL one
\newcommand{\Avirt}{A''}                     % virtual transmitted system (source descriptor of Xi_1)
\newcommand{\Abar}{\bar{A}}                  % Alice's measured register
\newcommand{\Ahat}{\hat{A}}                  % Alice's shield register
\newcommand{\pcNmsm}{\mathcal{N}_{\mathsf A}}  % Alice's measurement channel (see the RULES block)
\newcommand{\Phip}{\Phi^{+}}                 % maximally entangled vector label
\newcommand{\Smap}{\mathfrak{s}}             % the emitted sector map (fraktur: a compile-time object)

% --- receiver ------------------------------------------------------------------
\newcommand{\Modes}{\mathbb{M}}              % declared ordered mode set
\newcommand{\Inp}{\mathbb{I}}                % declared input modes
\newcommand{\InpC}{\mathbb{I}_{\Chan}}       % channel input modes
\newcommand{\InpS}{\mathbb{I}_{\Prb}}        % probe (source) input modes
\newcommand{\Env}{\mathbb{E}}                % dilation (environment) modes
\newcommand{\Port}[1]{\mathbb{P}_{#1}}       % mode set of detector port #1
\newcommand{\Emd}[1]{E_{#1}}                 % coordinate projector on modes #1
\newcommand{\Gmat}{G}                        % instrument (transfer) matrix
\newcommand{\Fock}{\mathfrak{F}}             % Fock / second-quantization functor
\newcommand{\Symn}[1]{\mathrm{Sym}^{#1}}     % symmetric tensor power
\newcommand{\normord}[1]{{:}\mspace{2mu}#1\mspace{2mu}{:}}   % normal ordering
\newcommand{\rk}{\operatorname{rank}}        % rank (physics.sty owns \rank)
\newcommand{\SCK}{\mathsf{SC}}               % single-click pattern set
\newcommand{\Recv}{\mathfrak{R}}             % receiver instrument tuple
\newcommand{\Cat}{\mathfrak{C}}              % component catalogue
\newcommand{\Grph}{\mathfrak{E}}             % declared element graph
\newcommand{\defic}{\mathbb{L}}              % loss deficiency operator
\newcommand{\Pvac}[1]{\Pi_{\vac}(#1)}        % joint no-click element on ports #1
\newcommand{\pmc}{p_{\mc}}                   % uniform high-photon mc weight
\newcommand{\Bnd}{\beta}                     % the scalar bound function beta(n)
\newcommand{\Rot}{R}                         % planar rotation matrix
\newcommand{\Wplate}{W}                      % retarder (wave-plate) matrix
\newcommand{\Bsp}{B}                         % beamsplitter core matrix

% --- statistics ----------------------------------------------------------------
\newcommand{\honest}{p^{\mathrm{hon}}}          % honest per-round distribution
\newcommand{\ecomp}{\varepsilon_{\mathrm{comp}}}% completeness budget
\newcommand{\simplex}[1]{\Delta(#1)}            % probability simplex
\newcommand{\Chatcon}{\Chatset_{\mathrm{con}}}  % constrained outcomes
\newcommand{\Ireg}[3]{\mathcal{I}_{#1}(#2,#3)}  % regularized incomplete beta
\newcommand{\Ncount}[1]{\mathsf{N}_{#1}}        % per-outcome count
\newcommand{\qlo}{\mathbf{q}^{-}}               % box lower bounds
\newcommand{\qhi}{\mathbf{q}^{+}}               % box upper bounds
 for it, so the
%  file is read from spec/modules/<id>/.
%
%  RULES
%    * A macro that spec/notation.md catalogues is a verbatim copy of it
%      (part A). Changing one is a notation change: bump spec/notation.md.
%    * Part B holds the module-local macros. A macro defined by two modules
%      with the same expansion appears once, under the first module (in the
%      order below) that defines it. Two macros with the same name and
%      DIFFERENT expansions are both kept under a module prefix, and the
%      claims write the prefixed names:
%        \Chain  -> \redChain (reductions, \mathfrak{C}: the reduction chain)
%                   \pcChain  (protocol-compiler, \mathfrak{R}: the two-segment
%                             chain; \Chsrc and \Chrec are built on it)
%        \Nmsm   -> \redNmsm  (reductions, \mathcal{N}: a POVM measurement
%                             channel)
%                   \pcNmsm   (\mathcal{N}_{\mathsf A}: Alice's measurement
%                             channel, written by
%                             transmitter/thm:certificate-sufficiency,
%                             transmitter/rem:xreg and
%                             reductions/lem:source-transfer)
%    * Deliberately omitted: a \provenance{tag}{details} macro. Provenance
%      is a header field of every claim file (spec/claim-format.md section
%      3); a claim body that calls such a macro fails to compile, which is
%      the intended signal.
%
%  REQUIRED PACKAGES, loaded by the assembler's preamble before this file:
%    amsmath, amssymb, amsfonts, mathtools, dsfont, braket, physics, xcolor
%    (braket before physics, so that physics wins for \ket, \bra, \tr, ...).
%
%  This file contains only definitions: no \usepackage, no document body.
% =============================================================================


% =============================================================================
%  PART A -- notation macros, the ones spec/notation.md catalogues, each a
%  verbatim copy of its catalogue row; grouped under numbered headings.
% =============================================================================

% -----------------------------------------------------------------------------
% 1. State spaces and operator sets (a "system label" -> spaces of operators)
% -----------------------------------------------------------------------------
\newcommand{\Hs}[1]{\mathcal{H}(#1)}        % Hilbert space of a system label
\newcommand{\Lin}[1]{\mathcal{L}(#1)}       % bounded linear operators on H(#1)
\newcommand{\Linn}[2]{\mathcal{L}(#1,#2)}   % linear maps H(#1) -> H(#2)
\newcommand{\Pos}[1]{\mathrm{Pos}(#1)}      % positive semidefinite operators on H(#1)
\newcommand{\Density}{\mathrm{D}}           % density operators: \Density(A) = D(A)
\newcommand{\Seq}[1]{S_{=}(#1)}             % normalized states (trace = 1)
\newcommand{\Sleq}[1]{S_{\leq}(#1)}         % subnormalized states (0 <= trace <= 1)
\newcommand{\id}{\mathds{1}}                % identity operator
\newcommand{\idon}[1]{\mathds{1}_{#1}}      % identity on a named system, e.g. \idon{C}

% DEPRECATED: applies \mathcal to both arguments itself, so it takes bare
% letters, not system labels. Prefer \Linn; do not use in new contracts.
\providecommand{\Lop}[2]{\mathcal{L}(\mathcal{#1}, \mathcal{#2})}


% -----------------------------------------------------------------------------
% 2. Composite systems
% -----------------------------------------------------------------------------
\newcommand{\dsum}{\oplus}                  % direct sum of system labels
% Tensor product: use \otimes directly.
% n-fold direct sum of an operator space: \Lin{A}^{\oplus n}


% -----------------------------------------------------------------------------
% 3. Events, conditioning and classical registers (for STATES)
%    Convention:  \wedge  -> subnormalized "and-event" component, e.g. \rho_{A\wedge x}
%                 |        -> normalized conditional state,        e.g. \rho_{A|x}
% (Write these directly in subscripts; no extra macros needed.)
% -----------------------------------------------------------------------------


% -----------------------------------------------------------------------------
% 4. Photon-number subspaces / direct-sum blocks
%    Convention:  operators carry the restriction "on top", in brackets;
%                 states carry it as a \wedge subscript (see Sec. 3).
% -----------------------------------------------------------------------------
\newcommand{\Pnum}[1]{P_{#1}}               % photon-number projector (system from context)
\newcommand{\Pnumon}[2]{P_{#1}^{#2}}        % with explicit system superscript, e.g. P_n^C
\newcommand{\blk}[2]{#1^{(#2)}}             % operator block, e.g. \blk{M}{\leqN} = M^{(<=N)}
\newcommand{\leqN}{\leq N}                  % use as \blk{M}{\leqN}
\newcommand{\gtN}{> N}                      % use as \blk{M}{\gtN}
\newcommand{\eqn}{= n}                      % use as \blk{M}{\eqn}


% -----------------------------------------------------------------------------
% 5. System labels used in the protocol
% -----------------------------------------------------------------------------
\newcommand{\Chan}{C}                       % quantum-channel (Eve) input to Bob
\newcommand{\Prb}{P}                        % probe input
\newcommand{\Flag}{{\mathrm{Flg}}}          % flag system (EXISTS ONLY AFTER squashing)
\newcommand{\Loss}{{\mathrm{Lss}}}          % loss system
\newcommand{\Vc}{{\mathrm{Vcm}}}            % vacuum system
\newcommand{\Nout}{N_{\mathrm{out}}}
\newcommand{\Nin}{N_{\mathrm{in}}}


% -----------------------------------------------------------------------------
% 6. Alphabets
% -----------------------------------------------------------------------------
\newcommand{\Yset}{\mathcal{Y}}             % Bob's (coarse-grained) detection alphabet
\newcommand{\Yraw}{\mathcal{Y}_{\mathrm{raw}}}% raw click patterns {0,1}^4
\newcommand{\Yscset}{\mathcal{Y}_{\singlec}}% single-click outcomes
\newcommand{\Ymcset}{\mathcal{Y}_{\mc}}     % multi-click outcomes
\newcommand{\Tset}{\mathcal{T}_{\mathrm{set}}}% private setting set (the probe settings in the optical instance)
\newcommand{\Xset}{\mathcal{X}}             % Alice's alphabet
\newcommand{\Dset}{\{H,V,+,-\}}             % polarization / detector labels
\newcommand{\pass}{\top}                    % pass sifting


% -----------------------------------------------------------------------------
% 7. Mode / outcome abbreviations
% -----------------------------------------------------------------------------
\newcommand{\on}{{\mathrm{on}}}
\newcommand{\off}{{\mathrm{off}}}
\newcommand{\dc}{{\mathrm{dc}}}             % dark count
\newcommand{\sq}{\mathrm{sq}}               % squashed
\newcommand{\mc}{\mathrm{mc}}               % multi-click
\newcommand{\nc}{\mathrm{nc}}               % no-click
\newcommand{\singlec}{\mathrm{sc}}          % single-click
\newcommand{\vac}{\mathrm{vac}}             % vacuum
\newcommand{\kout}{{k_\mathrm{out}}}
\newcommand{\kin}{{k_\mathrm{in}}}


% -----------------------------------------------------------------------------
% 8. Maps, channels and miscellaneous operators
% -----------------------------------------------------------------------------
\newcommand{\Prob}{\mathrm{Prob}}           % probability (use \Prob[\cdot] or \Prob(\cdot))
\newcommand{\lmin}{\lambda_{\min}}          % minimum eigenvalue
\newcommand{\Wflag}{W_{\Flag}}              % flag-space weight
\newcommand{\transpose}{\mathsf{T}}
\newcommand{\innerp}[2]{\left\langle #1 \vert #2 \right\rangle}
\renewcommand{\vec}[1]{\boldsymbol{#1}}     % NB: overrides the LaTeX arrow accent


% -----------------------------------------------------------------------------
% 9. Key-rate numerics
% -----------------------------------------------------------------------------
\newcommand{\cG}{\mathcal{G}}               % protocol map
\newcommand{\cZ}{\mathcal{Z}}               % pinching map
\newcommand{\Amap}{\mathcal{A}_{\mathrm{ann}}} % announcement + measurement map
\newcommand{\Amapnb}{\mathcal{A}_{\mathrm{ann}|\mathrm{nb}}}
\newcommand{\simpl}{}                       % empty suffix for simplified block maps
\newcommand{\acSet}{\mathcal{S}_{\mathrm{acc}}} % acceptance set
\newcommand{\testE}{\mathcal{E}_{\mathrm{test}}} % test-statistics map
\newcommand{\ecLeak}{\mathrm{leak}^{\mathrm{EC}}_{\mathrm{obs}}}
\newcommand{\pgen}{p_{\mathrm{gen}}}
\newcommand{\askeyr}{R_{\infty}}             % asymptotic key rate
\newcommand{\ppass}{p_{\mathrm{pass}}}       % sifting pass probability


% -----------------------------------------------------------------------------
% 10. Finite-size analysis / Renyi entropies
% -----------------------------------------------------------------------------

% --- Round types: test vs. key generation (register T, alphabet \Fset) --------
\newcommand{\tst}{\mathrm{test}}            % test round
\newcommand{\gen}{\mathrm{gen}}             % key-generation round
\newcommand{\Fset}{\mathcal{T}}             % round-type alphabet {test, gen} (not \Tset)
\newcommand{\Chat}{\widehat{C}}             % acceptance-test register
\newcommand{\Chatset}{\widehat{\mathcal{C}}}% its alphabet

% --- Basis / key extraction functions ----------------------------------------
\newcommand{\bA}{b_{A}}                     % Alice's basis-extraction function
\newcommand{\bB}{b_{B}}                     % Bob's basis-extraction function
\newcommand{\kB}{k_{B}}                     % Bob's key-extraction function

% --- Maps --------------------------------------------------------------------
\newcommand{\Mprot}{\mathcal{M}}            % single-round protocol map
\newcommand{\Phimap}{\boldsymbol{\Phi}}     % test-statistics map
\newcommand{\Pdeph}{\mathcal{P}}            % dephasing (pinching) map on \Blreg

% --- Isometries --------------------------------------------------------------
\newcommand{\Vmsm}{V_{\mathrm{msm}}}        % measurement isometry
\newcommand{\Vann}{V_{\mathrm{ann}}}        % announcement isometry
\newcommand{\VI}{V_{I}}                     % Eve's copy of the announcements
\newcommand{\VChat}{V_{\Chat}}              % acceptance-data isometry
\newcommand{\VR}{V_{R}}                     % sifting + key map -> raw key R
\newcommand{\VZ}{V_{Z}}                     % raw key -> key register Z
\newcommand{\Wprot}{W}                      % full protocol isometry

% --- Renyi entropies ---------------------------------------------------------
% SECURITY-CRITICAL: see spec/notation.md §4. Never change an arrow to
% make an inequality work out; never convert silently between the sandwiched
% (tilde) and Petz (overline) families.
\newcommand{\Hsandup}{\widetilde{H}^{\uparrow}_{\alpha}}   % sandwiched, up
\newcommand{\Hsanddown}{\widetilde{H}^{\downarrow}_{\alpha}}% sandwiched, down
\newcommand{\Hpetzup}{\overline{H}^{\uparrow}_{\alpha}}    % Petz, up
\newcommand{\Hpetzdown}{\overline{H}^{\downarrow}_{\alpha}}% Petz, down
\newcommand{\hup}{h^{\uparrow}_{\alpha}}                   % single-round rate
\newcommand{\hqkd}{h^{\uparrow,\mathrm{QKD}}_{\alpha}}     % single-round QKD rate
\newcommand{\hlow}{\underline{h}_{\alpha}}                 % lower bound on \hqkd
% --- spec/notation.md 3.5 and 4.5 -------------------------------------------
% \hlb  is the first-order linearization lower bound on \hlow: the headline
%       key rate (the `first-order` tier of foundations/def:tiers). NOT \hlow
%       itself.
% \tnum is the instance's one declared numerical tolerance. Its
%       near-collision with the declared marginal \tau_A is disclosed in
%       spec/notation.md 3.5; bare \tau is reserved and denotes neither.
\newcommand{\hlb}{\underline{h}^{\mathrm{FO}}_{\alpha}}    % first-order lower bound on \hlow
\newcommand{\tnum}{\tau_{\mathrm{num}}}                    % declared numerical tolerance


% =============================================================================
%  PART B -- module-local macros, grouped by module. Deprecated aliases,
%  defined so that a text written with them compiles, used in no new contract: \Lop,
%  \testE, \Pdeph, \pass, \ecLeak (spec/notation.md section 3.9).
% =============================================================================

% --- meat-core -----------------------------------------------------------------
\newcommand{\SecSet}{\mathbb{S}}            % sector index set of Q
\newcommand{\Qsys}{Q}                       % declared sector-carrying measured system
\newcommand{\Vset}{\mathcal{V}}             % declared joint outcome alphabet of the Q-POVM
\newcommand{\CAset}{\mathcal{C}_{A}}        % Alice announcement alphabet
\newcommand{\CBset}{\mathcal{C}_{B}}        % Bob announcement alphabet
\newcommand{\Party}{\mathfrak{P}}           % declared party partition
\newcommand{\Oset}[1]{\mathcal{O}_{#1}}     % readout alphabet of party #1
\newcommand{\rd}[1]{\pi_{#1}}               % readout map of party #1
\newcommand{\Sadm}{\mathcal{S}_{A}}         % admissible Alice subspace
\newcommand{\nomarg}{\mathsf{none}}         % "no trusted marginal declared"
\newcommand{\Zset}{\mathcal{Z}_{\perp}}     % raw-key alphabet incl. \perp
\newcommand{\Zval}{\mathcal{Z}}             % key alphabet without \perp
\newcommand{\Etil}{\widetilde{E}}           % single-round purifying register
\newcommand{\Ehat}{\widehat{E}}             % Eve's n-round register
\newcommand{\OmAT}{\Omega_{\mathrm{AT}}}    % acceptance event
\newcommand{\epa}{\varepsilon_{\mathrm{PA}}}
\newcommand{\eev}{\varepsilon_{\mathrm{EV}}}
\newcommand{\lEC}{\lambda_{\mathrm{EC}}}
\newcommand{\Decl}{\mathfrak{D}}            % protocol declaration
\newcommand{\CList}{\mathfrak{L}}           % declared constraint list
\newcommand{\PinS}{\mathcal{P}_{\SecSet}}   % sector pinching on Q
\newcommand{\Feas}{\mathcal{W}}             % feasible set of block states
\newcommand{\Chkey}{\Chatset_{\mathrm{key}}}% key-live acceptance values
\newcommand{\Kann}{\mathcal{K}}             % announcement-class alphabet
\newcommand{\hold}{\mathsf{h}}              % declared key holder (A or B)
\newcommand{\sfD}{D}                        % classical KL divergence (base 2)
\newcommand{\Dtwo}{D}                       % quantum relative entropy (base 2)
\newcommand{\Dsand}{\widetilde{D}_{\alpha}} % sandwiched Renyi divergence
\newcommand{\idmap}{\mathrm{id}}            % identity channel
\newcommand{\Aev}{\mathcal{A}}              % Eve's per-round attack channel
\newcommand{\AttSet}{\mathfrak{A}}          % licensed attack-channel set
\newcommand{\Chacc}{\Chatset_{\mathrm{acc}}}% acceptance support of \acSet
\newcommand{\epsstar}{\varepsilon^{\star}}      % distance to an exact POVM
\newcommand{\dfH}{\Delta_{\mathrm{herm}}}       % Hermiticity defect
\newcommand{\dfP}{\Delta_{\mathrm{pos}}}        % negative part
\newcommand{\dfC}{\Delta_{\mathrm{cpl}}}        % completeness defect
\newcommand{\lmoq}{g}                       % acceptance-side gradient block
\newcommand{\lmoJ}{G}                       % feasible-side gradient block
\newcommand{\bacc}{\beta_{\mathrm{acc}}}    % certified acceptance-side minimum
\newcommand{\Itset}{\mathfrak{I}}           % finite index set of iterates (C22)

% --- reductions ----------------------------------------------------------------
\newcommand{\Bsys}{B}                        % physical (pre-squashing) system
\newcommand{\Tcl}[1]{\mathcal{T}(#1)}        % trace-class operators on H(#1)
\newcommand{\Ins}{\mathcal{I}}               % generalized-instrument branch map
\newcommand{\redNmsm}{\mathcal{N}}              % POVM measurement channel
\newcommand{\Eprot}{\mathcal{E}}             % n-round protocol channel
\newcommand{\Grnd}{\mathsf{G}}               % one round map of a protocol
\newcommand{\Kfin}{\mathsf{K}}               % end-of-protocol processing
\newcommand{\Crit}{\mathsf{sec}}             % value of a security criterion
\newcommand{\Pflag}{\Pi_{\Flag}}             % flag-sector projector on Q
\newcommand{\Pleq}{\Pi_{\leqN}}              % preserved-block projector
\newcommand{\Pgt}{\Pi_{\gtN}}                % above-cutoff projector
\newcommand{\Sph}{S^{\mathrm{ph}}}           % family of physical inputs
\newcommand{\Srch}{S^{\mathrm{rch}}}         % reachable squashed states
\newcommand{\Ssq}{S^{\sq}}                   % squashed admissible set
\newcommand{\Ffss}{F^{\mathrm{FSS}}}         % flag-state-squasher target
\newcommand{\Fwp}{F^{\mathrm{WP}}}           % multi-bound WPFSS target
\newcommand{\Mwp}{M^{\mathrm{WP}}}           % renormalized above-cutoff POVM
\newcommand{\Mdec}{{M'}^{\mathrm{ph}}}       % declared physical family on B
\newcommand{\Ppvm}{P^{\Flag}}                % flag PVM
\newcommand{\Lmp}{{\Lambda_{\mathrm{mp}}}}   % measure-and-prepare channel
\newcommand{\Lwp}{{\Lambda^{\mathrm{WP}}}}   % multi-bound WPFSS channel
\newcommand{\Lfss}{{\Lambda^{\mathrm{FSS}}}} % flag-state-squasher channel
\newcommand{\lmax}{\lambda_{\max}}           % maximum eigenvalue
\newcommand{\Lnoise}{\widetilde{\Lambda}}    % flag-noise channel
\newcommand{\cg}{\varpi}                     % declared coarse-graining map
\newcommand{\MCK}{\mathsf{MC}}               % multiclick raw-pattern set
\newcommand{\redChain}{\mathfrak{C}}            % reduction chain
\newcommand{\Lnk}{\ell}                      % a single typed link
\newcommand{\dsq}{\preceq}                      % the squashing relation
\newcommand{\opn}[1]{\left\lVert #1\right\rVert}           % operator (spectral) norm
\newcommand{\mxn}[1]{\left\lVert #1\right\rVert_{\max}}    % entrywise max modulus
\newcommand{\Fld}{\mathbb{F}}                   % exact ordered substrate field
\newcommand{\Fldi}{\mathbb{F}_{i}}              % its Gaussian extension F + iF
\newcommand{\lat}[1]{\big\lfloor #1\big\rfloor_{\kappa}} % round DOWN onto the lattice

% --- protocol-compiler ---------------------------------------------------------
\newcommand{\Spec}{\mathfrak{S}}             % protocol specification (compiler input)
\newcommand{\pcChain}{\mathfrak{R}}            % reduction chain (two typed segments)
\newcommand{\Chsrc}{\pcChain_{\mathrm{src}}}   % its Alice-side source-map segment
\newcommand{\Chrec}{\pcChain_{\mathrm{rec}}}   % its receiver-side squashing segment
\newcommand{\rvirt}{\rho^{\mathrm{v}}}       % virtual (PRE-source-map) signal state
\newcommand{\Fph}{F^{\mathrm{ph}}}           % physical (pre-reduction) POVM element
\newcommand{\Vph}{\mathcal{V}^{\mathrm{ph}}} % physical outcome alphabet
\newcommand{\Tpub}{\mathcal{T}_{\mathrm{pub}}}% public-choice alphabet
\newcommand{\Uset}{\mathcal{U}}              % residual outcome alphabet
\newcommand{\idxm}{\mathsf{idx}}             % label-to-index map
\newcommand{\Asig}{A'}                       % transmitted (signal) system -- the REAL one
\newcommand{\Avirt}{A''}                     % virtual transmitted system (source descriptor of Xi_1)
\newcommand{\Abar}{\bar{A}}                  % Alice's measured register
\newcommand{\Ahat}{\hat{A}}                  % Alice's shield register
\newcommand{\pcNmsm}{\mathcal{N}_{\mathsf A}}  % Alice's measurement channel (see the RULES block)
\newcommand{\Phip}{\Phi^{+}}                 % maximally entangled vector label
\newcommand{\Smap}{\mathfrak{s}}             % the emitted sector map (fraktur: a compile-time object)

% --- receiver ------------------------------------------------------------------
\newcommand{\Modes}{\mathbb{M}}              % declared ordered mode set
\newcommand{\Inp}{\mathbb{I}}                % declared input modes
\newcommand{\InpC}{\mathbb{I}_{\Chan}}       % channel input modes
\newcommand{\InpS}{\mathbb{I}_{\Prb}}        % probe (source) input modes
\newcommand{\Env}{\mathbb{E}}                % dilation (environment) modes
\newcommand{\Port}[1]{\mathbb{P}_{#1}}       % mode set of detector port #1
\newcommand{\Emd}[1]{E_{#1}}                 % coordinate projector on modes #1
\newcommand{\Gmat}{G}                        % instrument (transfer) matrix
\newcommand{\Fock}{\mathfrak{F}}             % Fock / second-quantization functor
\newcommand{\Symn}[1]{\mathrm{Sym}^{#1}}     % symmetric tensor power
\newcommand{\normord}[1]{{:}\mspace{2mu}#1\mspace{2mu}{:}}   % normal ordering
\newcommand{\rk}{\operatorname{rank}}        % rank (physics.sty owns \rank)
\newcommand{\SCK}{\mathsf{SC}}               % single-click pattern set
\newcommand{\Recv}{\mathfrak{R}}             % receiver instrument tuple
\newcommand{\Cat}{\mathfrak{C}}              % component catalogue
\newcommand{\Grph}{\mathfrak{E}}             % declared element graph
\newcommand{\defic}{\mathbb{L}}              % loss deficiency operator
\newcommand{\Pvac}[1]{\Pi_{\vac}(#1)}        % joint no-click element on ports #1
\newcommand{\pmc}{p_{\mc}}                   % uniform high-photon mc weight
\newcommand{\Bnd}{\beta}                     % the scalar bound function beta(n)
\newcommand{\Rot}{R}                         % planar rotation matrix
\newcommand{\Wplate}{W}                      % retarder (wave-plate) matrix
\newcommand{\Bsp}{B}                         % beamsplitter core matrix

% --- statistics ----------------------------------------------------------------
\newcommand{\honest}{p^{\mathrm{hon}}}          % honest per-round distribution
\newcommand{\ecomp}{\varepsilon_{\mathrm{comp}}}% completeness budget
\newcommand{\simplex}[1]{\Delta(#1)}            % probability simplex
\newcommand{\Chatcon}{\Chatset_{\mathrm{con}}}  % constrained outcomes
\newcommand{\Ireg}[3]{\mathcal{I}_{#1}(#2,#3)}  % regularized incomplete beta
\newcommand{\Ncount}[1]{\mathsf{N}_{#1}}        % per-outcome count
\newcommand{\qlo}{\mathbf{q}^{-}}               % box lower bounds
\newcommand{\qhi}{\mathbf{q}^{+}}               % box upper bounds

\makeatletter
\newcommand*{\ClaimHeadFont}{\bfseries}   % switched by prominence
\newcommand*{\ClaimNoteFont}{\mdseries}
\newcommand*{\ClaimShortId}{}
\newcommand*{\SetShortId}[1]{\renewcommand*{\ClaimShortId}{#1}}
\newcommand*{\RestatedName}{}
% Upright bodies throughout: italic is reserved for non-normative text.
% \@endparpenalty\@M keeps the meta line on the statement's page.
\newtheoremstyle{claim}{\topsep}{\topsep}{\normalfont\@endparpenalty\@M}{}%
  {\ClaimHeadFont}{.}{.5em}{\thmname{#1}\thmnumber{ #2}\thmnote{ {\ClaimNoteFont(#3)}}}
\newtheoremstyle{claimass}{\topsep}{\topsep}{\normalfont\@endparpenalty\@M}{}%
  {\ClaimHeadFont}{.}{.5em}{\thmname{#1}\thmnumber{ #2}%
  \ifx\ClaimShortId\@empty\else\ [\textsf{\ClaimShortId}]\fi%
  \thmnote{ {\ClaimNoteFont(#3)}}}
\newtheoremstyle{claimremark}{.5\topsep}{.5\topsep}{\small\@endparpenalty\@M}{}%
  {\itshape}{.}{.5em}{\thmname{#1}\thmnumber{ #2}\thmnote{ {\upshape(#3)}}}
\newtheoremstyle{claimreq}{\topsep}{\topsep}{\normalfont\@endparpenalty\@M}{}%
  {\color{reqAccent}\scshape}{.}{.5em}%
  {\thmname{#1}\thmnumber{ #2}\thmnote{ {\upshape(#3)}}}
\makeatother
\theoremstyle{claim}
\newtheorem{theorem}{Theorem}[section]
\newtheorem{lemma}[theorem]{Lemma}
\newtheorem{proposition}[theorem]{Proposition}
\newtheorem{corollary}[theorem]{Corollary}
\newtheorem{fact}[theorem]{Fact}
\newtheorem{definition}[theorem]{Definition}
\newtheorem*{restatedclaim}{\RestatedName}
\theoremstyle{claimass}
\newtheorem{assumption}[theorem]{Assumption}
\newtheorem*{restatedassumption}{\RestatedName}
\theoremstyle{claimremark}
\newtheorem{remark}[theorem]{Remark}
\newtheorem{example}[theorem]{Example}
\theoremstyle{claimreq}
\newtheorem{requirement}[theorem]{Requirement}
\newtheorem*{restatedrequirement}{\RestatedName}
%% --- 8. Prominence: main results a margin rule and a bold heading ----------
\newtcolorbox{mainclaim}{enhanced jigsaw, breakable, blanker,
  borderline west={1.5pt}{-8pt}{mainRule}, left=0pt, right=0pt, top=0pt, bottom=0pt,
  boxsep=0pt, before skip=\topsep, after skip=\topsep, parbox=false,
  before upper={\renewcommand*{\ClaimNoteFont}{\bfseries}}}
\newenvironment{technicalclaim}{\renewcommand*{\ClaimHeadFont}{\mdseries\color{techGrey}}%
  \renewcommand*{\ClaimNoteFont}{}}{}
%% --- 9. Non-normative text, the meta line, navigation ------------------------
\makeatletter
% Tie what follows to this paragraph: no page break, no list penalty before
% the next theorem head (the \@nobreak route a section heading uses).
\newcommand{\keepwithnext}{\par\nobreak\vspace{3pt}\@nobreaktrue
  \everypar{\if@nobreak\@nobreakfalse\clubpenalty\@M\else\clubpenalty\@clubpenalty\everypar{}\fi}}
\newcommand{\claimsummary}[2]{% #1 claim id (named destination), #2 summary
  \par\addvspace{\bigskipamount}\Needspace{6\baselineskip}%
  {\small\itshape\color{nnGrey}\noindent\hypertarget{claim-#1}{}#2\par}\keepwithnext}
\newcommand{\claimanchor}[1]{% a claim without a summary: the destination alone
  \par\addvspace{\bigskipamount}\Needspace{6\baselineskip}%
  \phantomsection\hypertarget{claim-#1}{}\keepwithnext}
\newcommand{\claimleadin}[1]{{\color{nnGrey}\noindent#1\par}\keepwithnext}
\newcommand{\claimmeta}[1]{\par\nobreak\vspace{2pt}%
  {\footnotesize\color{metaGrey}\raggedright\interlinepenalty\@M\noindent#1\par}%
  \addvspace{\medskipamount}}
% After an assumption: the sentence generated from its verification field; after
% any claim with one: its plain-text reading (user_statement). Both
% non-normative, both grey.
\newcommand{\claimverification}[1]{\par\nobreak\vspace{2pt}%
  {\footnotesize\color{metaGrey}\noindent#1\par}\nobreak}
\newcommand{\claimreading}[1]{\par\nobreak\vspace{2pt}%
  {\footnotesize\color{nnGrey}\noindent\emph{Reading:}\enspace#1\par}\nobreak}
\makeatother
\newcommand{\claimbookmark}[4]{% #1 id, #2 kind word, #3 label, #4 plain title
  \bookmark[dest=claim-#1,level=2]{#2 \getrefnumber{#3} #4}}
\newcommand{\metasep}{\unskip\enspace\textperiodcentered\enspace\ignorespaces}
\newcommand{\metalabel}[1]{\emph{#1}\enspace}
\newcommand{\sectionintro}[1]{{\color{nnGrey}\noindent#1\par}\smallskip}
\newcommand{\sectionimports}[1]{{\footnotesize\color{metaGrey}\raggedright\noindent\metalabel{Imports}#1\par}\smallskip}
\newenvironment{overview}{\par\color{nnGrey}\itshape\noindent\ignorespaces}{\par}
% References: the whole "Def. 2.1" is the link; ids are the fallback text of
% every link into another PDF (a relative URI to a named destination).
\newcommand{\mref}[2]{\hyperref[#2]{#1~\ref*{#2}}}
\newcommand{\aref}[2]{\hyperref[#2]{\textsf{#1}}}
\DeclareUrlCommand{\claimid}{\urlstyle{sf}\def\UrlBreaks{\do\/\do\:\do\-\do\.\do\_}%
  \def\UrlBigBreaks{}\def\UrlNoBreaks{}}
\newcommand{\extref}[2]{\hyperref{#1}{}{claim-#2}{\claimid{#2}}}
\newcommand{\extaref}[3]{\hyperref{#1}{}{claim-#2}{\textsf{#3}}}
\newcommand{\coderef}[2]{\href{#1}{\claimid{#2}}}
% Claim-body commands (spec/claim-format.md section 4): \xref{<full id>} prints
% the kind, title and id of a claim of another contract, linked to its stable
% destination in that contract's PDF (one \DeclareXref per cited claim below);
% \term[gloss]{phrase} marks a defined phrase for the index (the gloss feeds
% the Terms table); \tref[phrase]{text} is a use of a defined term. The
% assembler rewrites both in the body into \termdef (the definition's named
% destination term-<slug>) and \termref / \exttermref (a link to it, in this
% or another contract's PDF); the fallbacks below serve text it copies
% elsewhere, such as the digest. \iface is an interface key.
\makeatletter
\newcommand{\DeclareXref}[4]{% id, kind word, title, contract PDF
  \@namedef{xref@#1}{\hyperref{#4}{}{claim-#1}{#2 ``#3'' ({\small\claimid{#1}})}}}
\newcommand{\xref}[1]{\@ifundefined{xref@#1}%
  {\textbf{[unresolved \nolinkurl{#1}]}\@latex@warning{Unresolved xref #1}}%
  {\@nameuse{xref@#1}}}
\makeatother
\DeclareUrlCommand{\iface}{\urlstyle{tt}}
\newcommand{\term}[2][]{\emph{#2}\index{#2}}
\newcommand{\tref}[2][]{#2}
\newcommand{\termdef}[2]{\hypertarget{#1}{}\emph{#2}\index{#2}}
\newcommand{\termref}[2]{\hyperlink{#1}{#2}}
\newcommand{\exttermref}[3]{\hyperref{#1}{}{#2}{#3}}
\hypersetup{pdftitle={Receiver POVM (receiver) -- contract 0.11.0}}

% Claims of other contracts cited here (\xref, [[...]]).
\DeclareXref{foundations/ass:snapped-measurement}{Assumption}{The device measures the snapped family where the run reads the snap}{../foundations/contract.pdf}
\DeclareXref{foundations/cor:least-eigenvalue-over-enclosure}{Corollary}{Least eigenvalue and rank over an enclosure}{../foundations/contract.pdf}
\DeclareXref{foundations/cor:rank-over-factor-enclosure}{Corollary}{Rank over a factor enclosure}{../foundations/contract.pdf}
\DeclareXref{foundations/def:binary64}{Definition}{Binary64 numbers, declared values and dyadic lattices}{../foundations/contract.pdf}
\DeclareXref{foundations/def:declared-data}{Definition}{Declared data and the parties to a module boundary}{../foundations/contract.pdf}
\DeclareXref{foundations/def:evaluated-off-snap}{Definition}{A reference handed beside a snap; a block evaluated off its snap}{../foundations/contract.pdf}
\DeclareXref{foundations/def:exact-island}{Definition}{The exact island}{../foundations/contract.pdf}
\DeclareXref{foundations/def:exact-reference}{Definition}{Reference of a measurement family, exact or enclosed}{../foundations/contract.pdf}
\DeclareXref{foundations/def:factor-enclosure}{Definition}{Factor enclosure of a block}{../foundations/contract.pdf}
\DeclareXref{foundations/def:round-homogeneity}{Definition}{Round homogeneity of a per-round family}{../foundations/contract.pdf}
\DeclareXref{foundations/def:sector-map}{Definition}{Sector maps and block-list operators}{../foundations/contract.pdf}
\DeclareXref{foundations/def:snap}{Definition}{Snapping declared data onto a lattice}{../foundations/contract.pdf}
\DeclareXref{foundations/def:tiers}{Definition}{The tiers of a reported key rate}{../foundations/contract.pdf}
\DeclareXref{foundations/lem:enclosure-psd}{Lemma}{Positive semidefiniteness from an entrywise enclosure}{../foundations/contract.pdf}
\DeclareXref{foundations/lem:least-eigenvalue-bound}{Lemma}{Least-eigenvalue lower bound from an approximate eigenbasis}{../foundations/contract.pdf}
\DeclareXref{foundations/lem:rational-obstruction}{Lemma}{Rationality obstruction to the domination of a lattice family}{../foundations/contract.pdf}
\DeclareXref{foundations/rem:noise-floor}{Remark}{What domination controls, and why it needs a noise floor}{../foundations/contract.pdf}
\DeclareXref{foundations/req:consumer-decides-exactly}{Requirement}{Exact decision by every consumer}{../foundations/contract.pdf}
\DeclareXref{foundations/req:rounding-direction}{Requirement}{Rounding direction of emitted bounds}{../foundations/contract.pdf}
\DeclareXref{foundations/req:snap-reference}{Requirement}{Snaps of a measurement family: kept, or moved and disclosed}{../foundations/contract.pdf}
\DeclareXref{foundations/req:support-safe-snap}{Requirement}{Support-safe snaps of a measurement family}{../foundations/contract.pdf}
\DeclareXref{foundations/req:tiers}{Requirement}{Reporting of the tiers}{../foundations/contract.pdf}
\DeclareXref{foundations/req:tolerance-scope}{Requirement}{Scope of the declared tolerance}{../foundations/contract.pdf}
\DeclareXref{foundations/req:who-snaps}{Requirement}{Snapping by the declarer only}{../foundations/contract.pdf}
\DeclareXref{meat-core/frank-wolfe/rem:fo-arithmetic}{Remark}{Exact real arithmetic in the first-order bound}{../meat-core/contract.pdf}
\DeclareXref{meat-core/frank-wolfe/req:bracket-on-reference}{Requirement}{The first-order bound evaluated on the reference companion}{../meat-core/contract.pdf}
\DeclareXref{meat-core/frank-wolfe/req:certified-evaluation}{Requirement}{The certified evaluation of the first-order bound}{../meat-core/contract.pdf}
\DeclareXref{reductions/def:joint-family}{Definition}{Joint alphabet, classes and the joint family of a private setting}{../reductions/contract.pdf}
\DeclareXref{reductions/lem:row-relaxation}{Lemma}{A rounded constraint entry with its constant relaxed}{../reductions/contract.pdf}

\begin{document}
\begin{center}
{\LARGE\bfseries Receiver POVM (receiver)}\\[4pt]
{\large Contract, version 0.11.0}
\end{center}
\begin{overview}
This contract turns a declared passive optical receiver into the click-pattern measurement it performs and into lower bounds on its multiclick weight on every photon-number block above a cutoff N; both are recorded in the receiver's artifact, which the flag-state squasher of the reductions module consumes. A receiver may be declared by its port Grams (route (i)), by an \termref{term-instrument-matrix}{instrument matrix} (route (ii)) or by a graph of catalogue components (route (iii)); the three routes converge on the port Gram family, which determines the click statistics exactly. A receiver fed by its own probe source as well as by the channel needs a submodule that constructs its \termref{term-setting-family}{setting family} and is otherwise refused. The main results of this core module are the click-pattern POVM, which follows by inclusion and exclusion and is block diagonal in the total input photon number (\hyperlink{claim-receiver/thm:click-povm}{Theorem~\ref*{thm:click-povm}}); its \termref{term-declared-coarse-graining}{declared coarse-graining} into the protocol's alphabet, under which lower bounds on disjoint raw outcomes add (\hyperlink{claim-receiver/thm:coarse-graining}{Theorem~\ref*{thm:coarse-graining}}); the weight in [0, 1] emitted for each probe setting, which bounds that setting's \termref{term-multiclick-element}{multiclick element} from below on every block of channel photon number above N (\hyperlink{claim-receiver/thm:emitted-bound}{Theorem~\ref*{thm:emitted-bound}}); and the range of the weights the artifact carries in its \termref{term-strengthening-slot}{strengthening slot}, one per outcome set and setting, which are non-negative and sum to at most one at every setting (\hyperlink{claim-receiver/prop:slotvalid}{Proposition~\ref*{prop:slotvalid}}). The weight comes from an analytic lower bound on the smallest eigenvalue of every multiclick block, made uniform over all blocks above N either by a collapse lemma, which reduces the non-negative part of the infimum over those blocks to that of the value at the first of them, or by the minimum over a finite window of blocks capped by an analytic term for the tail beyond it; an optional exact window uses lower bounds on the smallest eigenvalues of finitely many blocks. Every statement is exact unless it names a receiver tolerance or the \termref{term-eigenvalue-slack}{eigenvalue slack}: where the declaration defines an admissible receiver exactly, the weights are evaluated rigorously, in exact rational arithmetic from enclosures of the model taken on the conservative side, and are lower bounds by theorem with no premise about arithmetic (\hyperlink{claim-receiver/thm:enclosed-weight}{Theorem~\ref*{thm:enclosed-weight}}), while elsewhere they are evaluated in binary64 floating point, residuals diagnose the evaluation, and an exact window rests on the assumption that each computed eigenvalue less a declared slack does not exceed the exact one. Separately, for every declaration that defines an admissible receiver exactly, whatever its parameters, the photon-number blocks are enclosed rigorously in exact rational arithmetic, from enclosures of the square roots, cosines and sines the construction takes (\hyperlink{claim-receiver/thm:enclosed-blocks}{Theorem~\ref*{thm:enclosed-blocks}}); where the port Grams are Gaussian-rational the blocks of at most one photon are known exactly (\hyperlink{claim-receiver/thm:exact-blocks}{Theorem~\ref*{thm:exact-blocks}}), and, read through a \termref{term-port-dilation}{port dilation}, a block of a receiver without probe inputs can be given besides in enclosed factor form, which fixes its rank where the Gram matrix of its columns has a positive floor, as no entrywise box with positive radii can for a block of lower rank whose support no coordinate vectors span (\hyperlink{claim-receiver/thm:factor-blocks}{Theorem~\ref*{thm:factor-blocks}}). These enclosures are the reference a snap of the emitted blocks is certified against. The results are conditional on the device assumptions of the Assumptions section, among them a trusted, exactly characterised, passive device whose declared modes carry every path to a detector: the main results rest only on the few of these their proofs use, and the others enter the weight through the corollaries that apply \hyperlink{claim-receiver/thm:emitted-bound}{Theorem~\ref*{thm:emitted-bound}} to each setting family. Detector imperfections and the \termref{term-component-catalogue}{component catalogue} are submodules, each with main results of its own.
\end{overview}
\medskip
\noindent{\small\textcolor{nnGrey}{\itshape Grey text}, here and throughout, is non-normative: it explains, and nothing in it is assumed or proved.}

\smallskip
\noindent{\small The \href{https://github.com/fraunhofersg/certiqs-key/blob/v0.1.0/spec/modules/receiver/design.md}{design notes} (non-normative) explain why the module is built as it is and draw its processing flowchart. The dependencies between the sections are drawn in \href{graphs/overview.svg}{the section dependency matrix}, and \href{https://github.com/fraunhofersg/certiqs-key/blob/v0.1.0/spec/modules/receiver/graphs/README.md}{the graph index} opens every section's dependency graph. The \href{digest.pdf}{statements digest} holds every definition, assumption, result and requirement without the proofs. \href{https://github.com/fraunhofersg/certiqs-key/blob/v0.1.0/harness/contract-reading.md}{How to read a contract} explains this layout and its links.}

{\small\tableofcontents}

\section*{Main results}\addcontentsline{toc}{section}{Main results}
\sectionintro{The module's most important results, those its overview names as what it establishes, with every assumption each rests on, directly or through the claims it uses. Each links to its statement and proof; the statements digest holds them in full.}
\par\noindent\hyperref[thm:click-povm]{\textbf{Theorem~\ref*{thm:click-povm}}}\enspace Click-pattern POVM\par\nobreak
{\small\itshape\color{nnGrey}\noindent The click patterns of a receiver instrument form a POVM whose elements follow from the port Grams by inclusion and exclusion. Every element is block diagonal in the total input photon number, and vacuum gives only the no-click outcome.\par}\nobreak
\claimmeta{\metalabel{Rests on}\aref{A-PORT}{ass:port}, \aref{A-VAC}{ass:vac}}
\par\noindent\hyperref[thm:coarse-graining]{\textbf{Theorem~\ref*{thm:coarse-graining}}}\enspace Coarse-graining of the click-pattern POVM\par\nobreak
{\small\itshape\color{nnGrey}\noindent The emitted POVM of a declared coarse-graining is a POVM, block diagonal in the photon number, and lower bounds on disjoint raw outcomes add. A bound on a merged outcome gives nothing about the click patterns it merges.\par}\nobreak
\claimmeta{\metalabel{Rests on}\aref{A-PORT}{ass:port}, \aref{A-VAC}{ass:vac}}
\par\noindent\hyperref[thm:emitted-bound]{\textbf{Theorem~\ref*{thm:emitted-bound}}}\enspace Multiclick bounds carried by the emitted weights\par\nobreak
{\small\itshape\color{nnGrey}\noindent Each setting's emitted weight bounds its multiclick element on every photon-number block above the cutoff, and their minimum does so at every setting, under the eigenvalue-slack assumption \hyperlink{claim-receiver/ass:eig}{A-EIG} only where computed block eigenvalues enter. Each weight lies in [0, 1] and is finitely computable.\par}\nobreak
\claimmeta{\metalabel{Rests on}\aref{A-EIG}{ass:eig}, \aref{A-PORT}{ass:port}, \aref{A-VAC}{ass:vac}}
\par\noindent\hyperref[prop:slotvalid]{\textbf{Proposition~\ref*{prop:slotvalid}}}\enspace Range of the strengthening-slot weights\par\nobreak
{\small\itshape\color{nnGrey}\noindent At every setting the strengthening slot's weights (high-photon lower bounds on disjoint sets of raw outcomes) are non-negative and sum to at most one, as are their minima over the settings, and each emitted multiclick weight lies in [0, 1]. That the emitted weights are valid bounds is \hyperlink{claim-receiver/thm:emitted-bound}{Theorem~\ref*{thm:emitted-bound}}.\par}\nobreak
\claimmeta{\metalabel{Rests on}\aref{A-EIG}{ass:eig}, \aref{A-PORT}{ass:port}, \aref{A-VAC}{ass:vac}}
\par\noindent\hyperref[thm:catalogue-admissibility]{\textbf{Theorem~\ref*{thm:catalogue-admissibility}}}\enspace Admissibility of the catalogue entries\par\nobreak
{\small\itshape\color{nnGrey}\noindent Every catalogue entry maps its whole declared parameter domain to contractions on its modes, lossless exactly on the subfamily the catalogue's table lists for it. Only the generic block, an arbitrary matrix on declared modes, needs its contractivity checked at admission.\par}\nobreak
\claimmeta{\metalabel{Rests on}no assumption}
\par\noindent\hyperref[thm:efficiency]{\textbf{Theorem~\ref*{thm:efficiency}}}\enspace Detector efficiency in the port Grams\par\nobreak
{\small\itshape\color{nnGrey}\noindent Efficiency \ensuremath{\eta}\_k uniform on detector port k scales that port's Gram by \ensuremath{\eta}\_k, keeping its rank if \ensuremath{\eta}\_k \textgreater{} 0; varying within a port, it is no scalar and can lower the rank. It enters only as the compiled instrument matrix's last layer, and the loss decomposition \hyperlink{claim-receiver/lem:deficiency}{Lemma~\ref*{lem:deficiency}} checks that loss is counted once.\par}\nobreak
\claimmeta{\metalabel{Rests on}\aref{A-EXACT}{ass:exact}, \aref{A-MODES}{ass:modes}, \aref{A-NOSEL}{ass:nosel}, \aref{A-PASSIVE}{ass:passive}, \aref{A-PORT}{ass:port}, \aref{A-ROUND}{ass:round}, \aref{A-TRUST}{ass:trust}, \aref{A-VAC}{ass:vac}}
\par\noindent\hyperref[thm:dark-counts]{\textbf{Theorem~\ref*{thm:dark-counts}}}\enspace Dark counts on a photon-number block-diagonal POVM\par\nobreak
{\small\itshape\color{nnGrey}\noindent A dark-count channel keeps the photon-number block-diagonal POVM structure and never moves weight out of the multiclick set. It raises the multiclick bound p\_mc to (1 \ensuremath{-} r\_min)p\_mc + r\_min, where r\_min is the least multiclick mass over every non-multiclick column, the no-click column included.\par}\nobreak
\claimmeta{\metalabel{Rests on}\aref{A-EXACT}{ass:exact}, \aref{A-MODES}{ass:modes}, \aref{A-NOSEL}{ass:nosel}, \aref{A-PASSIVE}{ass:passive}, \aref{A-PORT}{ass:port}, \aref{A-ROUND}{ass:round}, \aref{A-THRESH}{ass:thresh}, \aref{A-TRUST}{ass:trust}, \aref{A-VAC}{ass:vac}}

\section*{Assumptions at a glance}\addcontentsline{toc}{section}{Assumptions at a glance}
\sectionintro{Every assumption this contract states or its claims assume, with its category and the claims of this contract that assume it directly. An assumption of another contract links into that contract's PDF.}
{\footnotesize\begin{xltabular}{\linewidth}{@{}>{\raggedright\arraybackslash}p{1.6cm}>{\raggedright\arraybackslash}p{1.9cm}>{\raggedright\arraybackslash}X>{\raggedright\arraybackslash}p{2.3cm}>{\raggedright\arraybackslash}p{2.8cm}@{}}
\toprule
Assumption & Category & Title & Stated as & Assumed directly by\\
\midrule
\endhead
\aref{A-EXACT}{ass:exact} & receiver & Exact declared parameters & \mref{Assumption}{ass:exact} & \mref{Lem.}{lem:deficiency}, \mref{Cor.}{cor:unconditioned-family}, \mref{Cor.}{cor:emitted-unconditioned}, \mref{Prop.}{prop:efficiency-normalization}\\
\addlinespace
\aref{A-MODES}{ass:modes} & receiver & Completeness of the declared modes & \mref{Assumption}{ass:modes} & \mref{Lem.}{lem:deficiency}, \mref{Cor.}{cor:unconditioned-family}, \mref{Cor.}{cor:emitted-unconditioned}, \mref{Rem.}{rem:out}, \mref{Prop.}{prop:efficiency-normalization}\\
\addlinespace
\aref{A-NOSEL}{ass:nosel} & receiver & Device fixed independently of the adversary & \mref{Assumption}{ass:nosel} & \mref{Lem.}{lem:deficiency}, \mref{Cor.}{cor:unconditioned-family}, \mref{Cor.}{cor:emitted-unconditioned}, \mref{Prop.}{prop:efficiency-normalization}\\
\addlinespace
\aref{A-PASSIVE}{ass:passive} & receiver & Passive linear optics & \mref{Assumption}{ass:passive} & \mref{Lem.}{lem:deficiency}, \mref{Cor.}{cor:unconditioned-family}, \mref{Cor.}{cor:emitted-unconditioned}, \mref{Rem.}{rem:out}, \mref{Prop.}{prop:efficiency-normalization}\\
\addlinespace
\aref{A-PORT}{ass:port} & receiver & Mode-insensitive port response & \mref{Assumption}{ass:port} & \mref{Lem.}{lem:vacelim}, \mref{Cor.}{cor:unconditioned-family}, \mref{Cor.}{cor:emitted-unconditioned}, \mref{Prop.}{prop:efficiency-normalization}\\
\addlinespace
\aref{A-ROUND}{ass:round} & receiver & Round homogeneity of the instrument, probe and dark counts & \mref{Assumption}{ass:round} & \mref{Lem.}{lem:deficiency}, \mref{Cor.}{cor:unconditioned-family}, \mref{Cor.}{cor:emitted-unconditioned}, \mref{Prop.}{prop:efficiency-normalization}\\
\addlinespace
\aref{A-THRESH}{ass:thresh} & receiver & Threshold detectors without memory & \mref{Assumption}{ass:thresh} & \mref{Cor.}{cor:unconditioned-family}, \mref{Cor.}{cor:emitted-unconditioned}, \mref{Rem.}{rem:out}, \mref{Rem.}{rem:corrdark}, \mref{Prop.}{prop:efficiency-normalization}\\
\addlinespace
\aref{A-TRUST}{ass:trust} & receiver & Trusted and exactly characterised receiver & \mref{Assumption}{ass:trust} & \mref{Lem.}{lem:deficiency}, \mref{Cor.}{cor:unconditioned-family}, \mref{Cor.}{cor:emitted-unconditioned}, \mref{Rem.}{rem:powervsclick}, \mref{Prop.}{prop:efficiency-normalization}\\
\addlinespace
\aref{A-VAC}{ass:vac} & receiver & Vacuum on undeclared input modes & \mref{Assumption}{ass:vac} & \mref{Lem.}{lem:vacelim}, \mref{Cor.}{cor:compiled-instrument}, \mref{Lem.}{lem:deficiency}, \mref{Cor.}{cor:unconditioned-family}, \mref{Cor.}{cor:emitted-unconditioned}, \mref{Rem.}{rem:layers}, \mref{Rem.}{rem:bsconv}, \mref{Prop.}{prop:efficiency-normalization}\\
\addlinespace
\aref{A-EIG}{ass:eig} & numerical trust & Declared slack on the computed block eigenvalues & \mref{Assumption}{ass:eig} & \mref{Rem.}{rem:exacteig}, \mref{Thm.}{thm:emitted-bound}, \mref{Cor.}{cor:emitted-unconditioned}, \mref{Thm.}{thm:exact-and-evaluated}\\
\bottomrule
\end{xltabular}}

\section*{Notation}\addcontentsline{toc}{section}{Notation}
\sectionintro{The symbols of the repository's notation catalogue (\href{../../notation.md}{spec/notation.md}) that this contract uses, in the catalogue's order, with their meaning.}
{\small\begin{xltabular}{\linewidth}{@{}>{\raggedright\arraybackslash}p{3.6cm}>{\raggedright\arraybackslash}X@{}}
\toprule
Symbol & Meaning\\
\midrule
\endhead
$\Hs{A}$ & Hilbert space of the system label $A$\\
$\Lin{A}$ & Bounded linear operators on $\mathcal{H}(A)$\\
$\id$, $\idon{C}$ & Identity operator (contextual / named)\\
$\idmap$ & Identity channel\\
$\Pnum{n}$, $\Pnumon{n}{C}$ & Photon-number projector\\
$\blk{M}{\leqN}$ & Operator block restriction\\
$\leqN$, $\gtN$, $\eqn$ & Block indices for \texttt{\textbackslash{}blk}\\
$\Vset$ & The measuring parties' declared joint outcome alphabet \textemdash{} Bob's alone under the default partition\\
$\Chan$, $\Prb$ & Channel input; probe input\\
$\Yset$ & Bob's coarse-grained detection alphabet, $\vert\mathcal{Y}\vert=6$\\
$\Tset$ & Bob's private setting set, weights $q_t$ (\extref{../reductions/contract.pdf}{reductions/def:joint-family}); in the optical instance the probe settings (register $T_{\mathrm{set}}$)\\
$\on$, $\off$, $\dc$, $\sq$, $\mc$, $\nc$, $\singlec$, $\vac$ & Mode / outcome abbreviations\\
$\lmin$ & Minimum eigenvalue\\
$\transpose$ & Transpose, $X^{\mathsf T}$\\
$\vec{v}$ & Vector (bold)\\
\bottomrule
\end{xltabular}}

\section*{Terms}\addcontentsline{toc}{section}{Terms}
\sectionintro{The phrases this contract defines, and the terms of other contracts it uses, each with a one-line gloss and the definition that fixes it. In the text, a linked term opens its definition.}
{\small\begin{xltabular}{\linewidth}{@{}>{\raggedright\arraybackslash}p{3.6cm}>{\raggedright\arraybackslash}X>{\raggedright\arraybackslash}p{2.6cm}@{}}
\toprule
Term & Meaning & Defined in\\
\midrule
\endhead
\termref{term-admissible-bound-data}{admissible bound data} & Scalars between 0 and 1, a shared \ensuremath{\gamma} and c and a \ensuremath{\beta} and d per port, with \ensuremath{\gamma} at most every \ensuremath{\beta} and c at most every d, from which the bound function is built. & \mref{Definition}{def:bounddata}\\
\addlinespace[2pt]
\termref{term-admissible-gram-family}{admissible Gram family} & A finite family of positive semidefinite operators whose sum is at most the identity; an exact condition, distinct from exact admissibility of declared data. & \mref{Definition}{def:gram}\\
\addlinespace[2pt]
\termref{term-angle-enclosure}{angle enclosure} & a pair of complex enclosures with rational centres of the cosine and the sine of an angle given as a rational or as a rational multiple of pi, exact at the multiples of pi/2 & \mref{Definition}{def:complex-enclosure}\\
\addlinespace[2pt]
\termref{term-bound-data-of-the-setting}{bound data of the setting} & The admissible bound data built from a setting's prefactors and the extreme eigenvalues of its operators B. & \mref{Definition}{def:setting-family}\\
\addlinespace[2pt]
\termref{term-bound-function}{bound function} & The function 1 + (K-1)\ensuremath{\gamma}c\ensuremath{^{n}} - \ensuremath{\Sigma}\ensuremath{_{k}} \ensuremath{\beta}\ensuremath{_{k}}d\ensuremath{_{k}}\ensuremath{^{n}} of the photon number n, built from admissible bound data. & \mref{Definition}{def:bounddata}\\
\addlinespace[2pt]
\termref{term-catalogue-entries}{catalogue entries} & The ten entry types C1 to C10 of the component catalogue, each with a support, a parameter domain and a matrix. & \mref{Definition}{def:catalogue-entries}\\
\addlinespace[2pt]
\exttermref{../foundations/contract.pdf}{term-certified}{certified} & the tier of a lower bound on the minimum of the objective of a declaration, formed with every arithmetic step from that declaration's data exact or rigorously enclosed, so that none of its hypotheses concerns arithmetic & \extref{../foundations/contract.pdf}{foundations/def:tiers}\\
\addlinespace[2pt]
\termref{term-channel-input-space}{channel input space} & The D-dimensional space whose Fock space a setting family acts on. & \mref{Definition}{def:setting-family}\\
\addlinespace[2pt]
\termref{term-channel-modes}{channel modes} & The declared input modes through which the adversary's signal enters. & \mref{Definition}{def:instrument}\\
\addlinespace[2pt]
\termref{term-compiled-instrument-matrix}{compiled instrument matrix} & The matrix D\_eta times the ordered product of a well-formed graph's embedded elements. & \mref{Definition}{def:compile}\\
\addlinespace[2pt]
\termref{term-complex-enclosure}{complex enclosure} & a Gaussian-rational centre and a non-negative rational radius that bounds the distance of a complex number from the centre; a matrix is enclosed entry by entry & \mref{Definition}{def:complex-enclosure}\\
\addlinespace[2pt]
\termref{term-component-catalogue}{component catalogue} & A finite set of entry types from which a route-(iii) receiver is assembled. & \mref{Definition}{def:catalogue}\\
\addlinespace[2pt]
\termref{term-dark-composed-family}{dark-composed family} & the click patterns of each setting composed with a dark-count channel, and their sums over the fibres of the declared coarse-graining & \mref{Definition}{def:dark-composed}\\
\addlinespace[2pt]
\termref{term-dark-count-channel}{dark-count channel} & A column-stochastic matrix on click patterns that can add clicks and never removes one. & \mref{Definition}{def:darkadm}\\
\addlinespace[2pt]
\termref{term-declared-coarse-graining}{declared coarse-graining} & A surjection from the raw click patterns onto the declared emitted alphabet, supplied as data by the protocol declaration. & \mref{Definition}{def:coarse}\\
\addlinespace[2pt]
\termref{term-declared-dark-count-route}{declared dark-count route} & The dark-count channel is a matrix R declared as such. & \mref{Definition}{def:darkadm}\\
\addlinespace[2pt]
\exttermref{../foundations/contract.pdf}{term-declared-data}{declared data} & the data an instance states and the data a producer emits as its own declaration, with the sets such data define & \extref{../foundations/contract.pdf}{foundations/def:declared-data}\\
\addlinespace[2pt]
\termref{term-declared-element-graph}{declared element graph} & A finite directed graph whose vertices are catalogue elements, each with an entry type, a parameter tuple and a support. & \mref{Definition}{def:graph}\\
\addlinespace[2pt]
\termref{term-declared-input-modes}{declared input modes} & The modes a receiver instrument declares as inputs, the channel modes and the declared source modes; every other mode carries vacuum at the input. & \mref{Definition}{def:instrument}\\
\addlinespace[2pt]
\termref{term-declared-mode-space}{declared mode space} & A finite, totally ordered set of optical mode labels together with the Hilbert space whose orthonormal basis they index. & \mref{Definition}{def:modes}\\
\addlinespace[2pt]
\exttermref{../foundations/contract.pdf}{term-declared-order}{declared order} & a total order on a finite alphabet, fixed together with the alphabet & \extref{../foundations/contract.pdf}{foundations/def:declared-order}\\
\addlinespace[2pt]
\termref{term-declared-source-modes}{declared source modes} & The declared input modes that carry the receiver's own probe state. & \mref{Definition}{def:instrument}\\
\addlinespace[2pt]
\exttermref{../foundations/contract.pdf}{term-declared-tolerance}{declared tolerance} & a binary64 number, at least zero and below one, that an instance states in its run configuration & \extref{../foundations/contract.pdf}{foundations/def:tolerance}\\
\addlinespace[2pt]
\exttermref{../foundations/contract.pdf}{term-declared-value}{declared value} & a binary64 number stated as a datum, which denotes exactly that dyadic rational and approximates nothing & \extref{../foundations/contract.pdf}{foundations/def:binary64}\\
\addlinespace[2pt]
\termref{term-detector-ports}{detector ports} & The pairwise disjoint sets of modes each watched by one threshold detector. & \mref{Definition}{def:instrument}\\
\addlinespace[2pt]
\termref{term-eigenvalue-slack}{eigenvalue slack} & The declared number subtracted from each computed block eigenvalue of an exact window. & \mref{Definition}{def:pmcex}\\
\addlinespace[2pt]
\termref{term-elements}{elements} & A vertex of a declared element graph: a catalogue component with its parameters and support, not a POVM element. & \mref{Definition}{def:graph}\\
\addlinespace[2pt]
\termref{term-emitted-povm}{emitted POVM} & The measurement the receiver emits, one element per emitted outcome, each the sum of the click-pattern elements of its fibre. & \mref{Definition}{def:coarse}\\
\addlinespace[2pt]
\termref{term-enclosed-bound-data}{enclosed bound data} & Reals that bound the bound data of a setting from the conservative side, c and \ensuremath{\gamma} from below and each d and \ensuremath{\beta} from above, all within 0 and 1. & \mref{Definition}{def:enclosed-weight}\\
\addlinespace[2pt]
\termref{term-enclosed-weight}{enclosed weight} & The per-setting weight formed on enclosed bound data and, at an exact window, on rigorous window bounds, with no computed eigenvalue and no eigenvalue slack. & \mref{Definition}{def:enclosed-weight}\\
\addlinespace[2pt]
\termref{term-entry-type}{entry type} & A kind of optical component: a support shape, a parameter domain and a rule giving a local element for each parameter tuple. & \mref{Definition}{def:catalogue}\\
\addlinespace[2pt]
\exttermref{../foundations/contract.pdf}{term-entrywise-enclosure}{entrywise enclosure} & a Gaussian-rational Hermitian centre together with symmetric non-negative rational radii that bound, entry by entry, the modulus of the difference between a matrix and that centre & \extref{../foundations/contract.pdf}{foundations/def:exact-reference}\\
\addlinespace[2pt]
\exttermref{../foundations/contract.pdf}{term-evaluated-off-its-snap}{evaluated off its snap} & said of a block of a snap in a run: no number from which the run computes its key length reads the snapped block, save the points at which those numbers are evaluated and the values other than the output length that the run fixes before its first round & \extref{../foundations/contract.pdf}{foundations/def:evaluated-off-snap}\\
\addlinespace[2pt]
\termref{term-evaluation-model}{evaluation model} & The arithmetic of this contract: its exact claims evaluated in binary64 with diagnostic residuals and no rigorous enclosure of a floating-point number beyond the eigenvalue-slack assumption, beside exact rational enclosures of the blocks and a rigorous evaluation of the weights. & \mref{Definition}{def:evaluation-model}\\
\addlinespace[2pt]
\termref{term-exact-block-bound}{exact block bound} & The smallest eigenvalue of the n-photon block of the multiclick element, a function of n. & \mref{Definition}{def:exactbnd}\\
\addlinespace[2pt]
\termref{term-exact-flag}{exact flag} & The flag of a setting, exact or general; it may read exact only where every prefactor is one by construction. & \mref{Definition}{def:setting-family}\\
\addlinespace[2pt]
\exttermref{../foundations/contract.pdf}{term-exact-island}{exact island} & the computations of the engine, its instances' declarations included, by which a consumer of data or the declarer about to declare them decides whether they satisfy the conditions their declaration states, a declarer computes and decides the reference of a measurement family, a consumer's replacement of an emitted scalar bound is decided to lie on the bound's weakening side, a producer constructs the operator data it emits, the acceptance side of the optimisation is constructed and certified, or a lower bound on a linear minimum over the feasible side of the optimisation is certified & \extref{../foundations/contract.pdf}{foundations/def:exact-island}\\
\addlinespace[2pt]
\exttermref{../foundations/contract.pdf}{term-exact-reference}{exact reference} & a reference all of whose radii are zero, so that each block is its Gaussian-rational centre & \extref{../foundations/contract.pdf}{foundations/def:exact-reference}\\
\addlinespace[2pt]
\termref{term-exact-route}{exact route} & the declarations whose declared data determine, exactly, an admissible family of port Grams with Gaussian-rational entries, with a probe of the Gaussian tier & \mref{Definition}{def:exact-route}\\
\addlinespace[2pt]
\termref{term-exact-window-depth}{exact-window depth} & The declared integer up to which the weight of a setting uses, block by block, the larger of the analytic bound and a computed block eigenvalue less the eigenvalue slack. & \mref{Definition}{def:pmcex}\\
\addlinespace[2pt]
\exttermref{../foundations/contract.pdf}{term-factor-enclosure}{factor enclosure} & a matrix box of columns together with an exact positive diagonal core, such that the matrix is the product of a matrix L the box contains, the core and the adjoint of L & \extref{../foundations/contract.pdf}{foundations/def:factor-enclosure}\\
\addlinespace[2pt]
\termref{term-factored-loss}{factored loss} & The factor t, between the largest declared efficiency and one, by which a normalized station divides every efficiency. & \mref{Definition}{def:normalized-station}\\
\addlinespace[2pt]
\termref{term-fock-space}{Fock space} & The direct sum over n of the symmetric n-photon subspaces of a finite-dimensional mode space. & \mref{Definition}{def:fock}\\
\addlinespace[2pt]
\termref{term-gaussian-tier}{Gaussian tier} & The probe declarations whose states are products of single-mode thermal states, vacuum included. & \mref{Definition}{def:probe}\\
\addlinespace[2pt]
\termref{term-generated-dark-count-route}{generated dark-count route} & The dark-count channel is the independent product form of the declared per-port probabilities. & \mref{Definition}{def:darkadm}\\
\addlinespace[2pt]
\termref{term-generic-block}{Generic block} & Entry C10: an arbitrary matrix on a declared mode subset, admitted after an explicit contractivity check. & \mref{Definition}{def:catalogue-entries}\\
\addlinespace[2pt]
\termref{term-global-embedding}{global embedding} & The matrix on the whole mode space that acts as a local element on its support and as the identity elsewhere. & \mref{Definition}{def:embed}\\
\addlinespace[2pt]
\exttermref{../foundations/contract.pdf}{term-handed-reference}{handed reference} & a reference of a snapped measurement family that its declarer passes, with its centres and radii and unaltered, to every consumer of the snap & \extref{../foundations/contract.pdf}{foundations/def:evaluated-off-snap}\\
\addlinespace[2pt]
\termref{term-hermiticity-and-positivity-defect}{Hermiticity-and-positivity defect} & The larger of the spectral norm of X minus its adjoint and the negative part of the smallest eigenvalue of the Hermitian part of X: one number, measured in the spectral norm. & \mref{Definition}{def:residuals}\\
\addlinespace[2pt]
\termref{term-independent-product-form}{independent product form} & The dark-count channel built from per-port dark-count probabilities, each port firing independently. & \mref{Definition}{def:darkadm}\\
\addlinespace[2pt]
\termref{term-instrument-matrix}{instrument matrix} & The contraction of a receiver instrument from input modes to output modes. & \mref{Definition}{def:instrument}\\
\addlinespace[2pt]
\termref{term-joint-no-click-elements}{joint no-click elements} & For a set J of ports, the sum of the POVM elements of every outcome in which no port of J clicks. & \mref{Definition}{def:setting-family}\\
\addlinespace[2pt]
\exttermref{../foundations/contract.pdf}{term-lattice-value}{lattice value} & at a lattice exponent k, a real number that is an integer multiple of two to the power \ensuremath{-}k and has modulus at most one; a complex number or a matrix is one when the real and the imaginary part of every entry are & \extref{../foundations/contract.pdf}{foundations/def:binary64}\\
\addlinespace[2pt]
\termref{term-linear-extension}{linear extension} & An enumeration of the elements of a graph that respects its reachability order. & \mref{Definition}{def:graph}\\
\addlinespace[2pt]
\termref{term-local-element}{local element} & A contraction acting on a chosen subset of the declared modes. & \mref{Definition}{def:embed}\\
\addlinespace[2pt]
\termref{term-loss-deficiency}{loss deficiency} & The identity on the declared input modes minus the sum of all port Grams. & \mref{Definition}{def:gram}\\
\addlinespace[2pt]
\termref{term-lossless-element}{lossless element} & A local element that is an isometry on its support. & \mref{Definition}{def:embed}\\
\addlinespace[2pt]
\exttermref{../foundations/contract.pdf}{term-matrix-box}{matrix box} & a matrix of Gaussian-rational centres together with non-negative rational radii that bound, entry by entry, the modulus of the difference between a rectangular complex matrix and that centre & \extref{../foundations/contract.pdf}{foundations/def:factor-enclosure}\\
\addlinespace[2pt]
\termref{term-measurement-isometry}{measurement isometry} & The isometry from the Fock space of the declared input modes that prepares vacuum elsewhere and applies the lifted unitary dilation. & \mref{Definition}{def:physmodel}\\
\addlinespace[2pt]
\termref{term-multiclick-element}{multiclick element} & The sum of the POVM elements of every outcome in which at least two ports click. & \mref{Definition}{def:setting-family}\\
\addlinespace[2pt]
\termref{term-multiclick-set}{multiclick set} & The outcomes in which at least two ports click. & \mref{Definition}{def:clickalphabet}\\
\addlinespace[2pt]
\termref{term-no-click-outcome}{no-click outcome} & The outcome in which no port clicks. & \mref{Definition}{def:clickalphabet}\\
\addlinespace[2pt]
\termref{term-normalized-station}{normalized station} & A receiver instrument with the same channel modes and ports whose port Grams are those of the receiver divided by the factored loss t. & \mref{Definition}{def:normalized-station}\\
\addlinespace[2pt]
\termref{term-occupation-diagonal}{occupation-diagonal} & Said of a state on the probe Fock space that is diagonal in its occupation basis. & \mref{Definition}{def:probe}\\
\addlinespace[2pt]
\termref{term-optical-mode}{optical mode} & A label recording a spatial path, a polarization and a reserved temporal or spectral index. & \mref{Definition}{def:modes}\\
\addlinespace[2pt]
\termref{term-outcome}{outcome} & The set of detector ports that click in one round. & \mref{Definition}{def:clickalphabet}\\
\addlinespace[2pt]
\termref{term-path}{path} & A spatial index carrying the two polarization modes H and V. & \mref{Definition}{def:catalogue-entries}\\
\addlinespace[2pt]
\termref{term-per-setting-weight}{per-setting weight} & The high-photon multiclick weight of one setting: its collapse value, its window and tail, or its exact window. & \mref{Definition}{def:pmc}\\
\addlinespace[2pt]
\termref{term-polarization-dependent-beamsplitter}{polarization-dependent beamsplitter} & A beamsplitter that declares its transmittance and reflection phase separately for the H and the V modes. & \mref{Definition}{def:catalogue-entries}\\
\addlinespace[2pt]
\termref{term-port-dilation}{port dilation} & a finite ordered mode set with disjoint ports and an isometry from the declared input modes into it whose Gram on each port is that port's Gram matrix & \mref{Definition}{def:port-dilation}\\
\addlinespace[2pt]
\termref{term-port-gram-matrices}{port Gram matrices} & The positive operators on the declared input modes, one per detector port, whose quadratic form is the squared norm of the amplitude a receiver instrument sends to that port. & \mref{Definition}{def:gram}\\
\addlinespace[2pt]
\termref{term-predicate-x}{predicate (X)} & The condition that a setting's exact flag reads exact; it selects the collapse branch of the per-setting weight. & \mref{Definition}{def:pmc}\\
\addlinespace[2pt]
\termref{term-probe-declaration}{probe declaration} & A finite alphabet of probe settings, vacuum among them, with a declared occupation-diagonal state on the declared source modes for each setting. & \mref{Definition}{def:probe}\\
\addlinespace[2pt]
\termref{term-pure-loss-channel}{pure-loss channel} & The channel on the channel system that sends every channel mode through a beamsplitter of transmittance t with vacuum in its other input and discards the reflected light. & \mref{Definition}{def:pure-loss}\\
\addlinespace[2pt]
\termref{term-raw-alphabet}{raw alphabet} & The set of all outcomes, the subsets of the K ports. & \mref{Definition}{def:clickalphabet}\\
\addlinespace[2pt]
\termref{term-receiver-instrument}{receiver instrument} & A declared mode space with channel and probe inputs, disjoint detector ports and a contractive instrument matrix. & \mref{Definition}{def:instrument}\\
\addlinespace[2pt]
\termref{term-receiver-tolerances}{receiver tolerances} & The eight declared numbers at which the receiver decides its exact conditions on declared data, bounds its residuals and spends its eigenvalue slack; not the instance's declared tolerance. & \mref{Definition}{def:tolmap}\\
\addlinespace[2pt]
\exttermref{../foundations/contract.pdf}{term-reference}{reference} & a family of Hermitian positive semidefinite block lists complete at the constant on every block, read together with an entrywise enclosure of every block of every element & \extref{../foundations/contract.pdf}{foundations/def:exact-reference}\\
\addlinespace[2pt]
\termref{term-residuals}{residuals} & The six consistency checks on the floating-point evaluation of this contract, each exactly zero in exact arithmetic. & \mref{Definition}{def:residuals}\\
\addlinespace[2pt]
\termref{term-retarder}{Retarder} & Entry C3: a wave plate with a retardance, a fast-axis angle and a scalar loss, with half-wave and quarter-wave presets. & \mref{Definition}{def:catalogue-entries}\\
\addlinespace[2pt]
\termref{term-rigorous-window-bounds}{rigorous window bounds} & Reals at most the exact block bound of a setting at each index of its exact window. & \mref{Definition}{def:enclosed-weight}\\
\addlinespace[2pt]
\exttermref{../foundations/contract.pdf}{term-rigorously-enclosed}{rigorously enclosed} & said of a computation each of whose operations returns a set of numbers, such as an interval or a ball, containing the exact result of the operation at every value its inputs' sets contain & \extref{../foundations/contract.pdf}{foundations/def:tiers}\\
\addlinespace[2pt]
\termref{term-rotator}{rotator} & Entry C4: the circular retarder, a real rotation of the two polarization modes of one path. & \mref{Definition}{def:catalogue-entries}\\
\addlinespace[2pt]
\exttermref{../foundations/contract.pdf}{term-round-homogeneous}{round-homogeneous} & one fixed list of objects describes every round, and every round behaves as they describe whatever the other rounds do & \extref{../foundations/contract.pdf}{foundations/def:round-homogeneity}\\
\addlinespace[2pt]
\exttermref{../foundations/contract.pdf}{term-sector-map}{sector map} & a block kind together with a non-empty, ordered list of block identifiers, each with a dimension & \extref{../foundations/contract.pdf}{foundations/def:sector-map}\\
\addlinespace[2pt]
\termref{term-setting-family}{setting family} & A finite set of settings, each with a photon-number block-diagonal POVM whose joint no-click elements have the prefactor form, and an exact flag. & \mref{Definition}{def:setting-family}\\
\addlinespace[2pt]
\termref{term-single-click-set}{single-click set} & The outcomes in which exactly one port clicks. & \mref{Definition}{def:clickalphabet}\\
\addlinespace[2pt]
\termref{term-square-root-enclosure}{square-root enclosure} & a complex enclosure with a rational centre of the square root of a non-negative rational, exact where that rational is the square of a rational & \mref{Definition}{def:complex-enclosure}\\
\addlinespace[2pt]
\termref{term-strengthening-slot}{strengthening slot} & The list of disjoint raw-outcome sets with per-setting high-photon lower-bound weights that the artifact carries. & \mref{Definition}{def:slot}\\
\addlinespace[2pt]
\exttermref{../foundations/contract.pdf}{term-structural-support}{structural support} & of a block of an element: the set of indices whose row or column holds an entry that is not structurally zero & \extref{../foundations/contract.pdf}{foundations/def:snap}\\
\addlinespace[2pt]
\exttermref{../foundations/contract.pdf}{term-supplier}{supplier} & a module whose own claims decide their conditions on the data it reads or emits in floating point, at slacks an instance declares to it, rather than exactly & \extref{../foundations/contract.pdf}{foundations/def:declared-data}\\
\addlinespace[2pt]
\termref{term-support}{support} & The set of modes a local element acts on. & \mref{Definition}{def:embed}\\
\addlinespace[2pt]
\termref{term-support-comparability}{support comparability} & Condition (E3): any two elements whose supports intersect are ordered by the graph. & \mref{Definition}{def:graph}\\
\addlinespace[2pt]
\termref{term-terminal-detector-efficiency-layer}{terminal detector-efficiency layer} & The diagonal matrix D\_eta applying each port's detector efficiency after every element. & \mref{Definition}{def:compile}\\
\addlinespace[2pt]
\termref{term-uniform-weight}{uniform weight} & The minimum of the per-setting weights over all settings of a setting family. & \mref{Definition}{def:pmc}\\
\addlinespace[2pt]
\termref{term-uniform-weight-of-an-entry}{uniform weight of an entry} & The minimum over the settings of a slot entry's weights; it bounds the entry at every setting. & \mref{Definition}{def:slot}\\
\addlinespace[2pt]
\termref{term-well-formed}{well formed} & Said of a declared element graph that satisfies (E1) to (E3): valid parameters, acyclicity and support comparability. & \mref{Definition}{def:graph}\\
\bottomrule
\end{xltabular}}

\clearpage
\part{Mathematics}

\section{Modes, Fock space and the mode algebra}\label{group-receiver/def:modes}
\sectionintro{A receiver acts on a finite set of declared \termref{term-optical-mode}{optical modes}, and its photon statistics live in the \termref{term-fock-space}{Fock space} over them. This section lifts a matrix on the modes to every photon-number sector and records the algebra the later proofs run on: embedding, products and adjoints preserve contractions and elements on disjoint modes commute (\hyperlink{claim-receiver/prop:mode-algebra}{Proposition~\ref*{prop:mode-algebra}}), and the lift respects composition, adjoints and projectors (\hyperlink{claim-receiver/lem:fock-functor}{Lemma~\ref*{lem:fock-functor}}), also for maps between two mode spaces, where the lift of a direct sum is a tensor product (\hyperlink{claim-receiver/lem:fock-maps}{Lemma~\ref*{lem:fock-maps}}). The lift of 1 \ensuremath{-} A is called a lifted no-click operator because, with A the sum of the port Grams of a set of ports, it is the operator for the event that no port of the set clicks, pulled back to the declared inputs (\hyperlink{claim-receiver/thm:port-grams}{Theorem~\ref*{thm:port-grams}}). Its spectrum on each photon-number block is computed from the eigenvalues of A (\hyperlink{claim-receiver/lem:fock-block-spectrum}{Lemma~\ref*{lem:fock-block-spectrum}}), which is where the analytic high-photon bound starts. The normal-ordered form is stated for comparison with the optics literature; no proof uses it.}

% claim receiver/def:modes
\claimsummary{receiver/def:modes}{A declared mode space is a finite, totally ordered set of optical modes, each labelled by a path, a polarization and a reserved temporal or spectral index that takes one declared value. Matrices on the modes follow this order; on the input modes the channel modes come first (\hyperlink{claim-receiver/def:instrument}{Definition~\ref*{def:instrument}}).}
\claimbookmark{receiver/def:modes}{Definition}{def:modes}{Declared mode space}
\claimleadin{%
The temporal or spectral index of a mode label is reserved: it takes a single
value in this contract, so mode labels carry no temporal semantics. The order compares polarizations last,
so the $H$ and $V$ modes of a path are adjacent.
}
\begin{definition}[Declared mode space]\label{def:modes}
A \termdef{term-declared-mode-space}{declared mode space} is a finite totally ordered set $\Modes$, each of
whose elements is an \termdef{term-optical-mode}{optical mode}, together with the Hilbert space
$\mathbb{C}^{\Modes}$ carrying the orthonormal basis $\{\ket{m}\}_{m\in\Modes}$
indexed in the \exttermref{../foundations/contract.pdf}{term-declared-order}{declared order}. A mode is a label $m=(p,\pi,\tau)$ recording a
spatial path $p$, a polarization $\pi\in\{H,V\}$ and a temporal/spectral index
$\tau$; the order is lexicographic in $(p,\tau,\pi)$ with $H$ before $V$. The
index $\tau$ is \emph{reserved}: it is carried in the label and in the emitted
artifact, it may take only a single \exttermref{../foundations/contract.pdf}{term-declared-value}{declared value} in this contract, and no
catalogue entry (Definition~\ref{def:catalogue-entries}) reads or writes it.
For $\Sigma\subseteq\Modes$ we write
$\Emd{\Sigma}=\sum_{m\in\Sigma}\dyad{m}$ for the coordinate projector, and
$\Modes$ is identified with $\{0,\dots,|\Modes|-1\}$ when matrix indices are
needed.
\end{definition}
\claimmeta{\metalabel{Used by}\mref{Prop.}{prop:mode-algebra}, \mref{Def.}{def:instrument}, \mref{Ass.}{ass:thresh}
\metasep newly derived
\metasep accepted}

% claim receiver/def:fock
\claimsummary{receiver/def:fock}{The Fock space over a finite-dimensional space is the direct sum of its symmetric photon-number sectors, and a linear map lifts to it sector by sector. Every photon-number block of the receiver's measurement is assembled from blocks of such lifts.}
\claimbookmark{receiver/def:fock}{Definition}{def:fock}{Fock space and the second-quantization functor}
\claimleadin{%
The lift is written for a map between two spaces as well as for an operator on
one; the two-space form is used for coordinate inclusions of mode spaces,
among them the inclusion of the \termref{term-declared-input-modes}{declared input modes} into the dilated
mode space (Definition~\ref{def:physmodel}).
}
\begin{definition}[Fock space and the second-quantization functor]
\label{def:fock}
For a finite-dimensional Hilbert space $\mathcal{K}_0$ the \termdef{term-fock-space}{Fock space} is
$\mathcal{F}(\mathcal{K}_0)=\bigoplus_{n\ge0}\Symn{n}(\mathcal{K}_0)$, where
$\Symn{n}(\mathcal{K}_0)\subseteq\mathcal{K}_0^{\otimes n}$ is the subspace
invariant under every permutation of the $n$ tensor factors and
$\Symn{0}(\mathcal{K}_0)=\mathbb{C}\ket{\vac}$. For a linear operator $X$ on
$\mathcal{K}_0$ put
\begin{equation}\label{eq:fockfunctor}
  \Fock(X)\coloneqq\bigoplus_{n\ge0}\big(X^{\otimes n}\big)\big|_{\Symn{n}},
  \qquad \Fock(X)\big|_{\Symn{0}}=1 .
\end{equation}
This is well defined because $X^{\otimes n}$ commutes with every permutation of
the factors and therefore preserves $\Symn{n}(\mathcal{K}_0)$
(Lemma~\ref{lem:fock-functor}). The same formula defines
$\Fock(X):\mathcal{F}(\mathcal{K}_0)\to\mathcal{F}(\mathcal{K}_1)$ for a linear
map $X:\mathcal{K}_0\to\mathcal{K}_1$ between two spaces, because
$X^{\otimes n}$ intertwines the permutation actions and therefore carries
$\Symn{n}(\mathcal{K}_0)$ into $\Symn{n}(\mathcal{K}_1)$. Occupation vectors
$\vec m\in\mathbb{N}_0^{\dim\mathcal{K}_0}$ with $|\vec m|=n$ index the
orthonormal basis $\{\ket{\vec m}\}$ of $\Symn{n}(\mathcal{K}_0)$ obtained by
symmetrizing and normalizing $\ket{e_1}^{\otimes m_1}\otimes\cdots$ in a fixed
orthonormal basis $\{\ket{e_j}\}$ of $\mathcal{K}_0$; $a_j$ and $a_j^\dagger$
denote the annihilation and creation operators of $\ket{e_j}$.
\end{definition}
\claimmeta{\metalabel{Used by}\mref{Lem.}{lem:fock-functor}, \mref{Lem.}{lem:fock-maps}, \mref{Def.}{def:instrument}, \mref{Def.}{def:probe}, \mref{Def.}{def:physmodel}, \mref{Lem.}{lem:vacelim}, \mref{Def.}{def:exactbnd}, \mref{Def.}{def:setting-family}, \mref{Lem.}{lem:fock-permanent}, \mref{Lem.}{lem:fock-permanent-maps}, \mref{Def.}{def:port-dilation}, \mref{Thm.}{thm:factor-blocks}, \mref{Def.}{def:dark-composed}, \mref{Prop.}{prop:noise-floor}, \mref{Def.}{def:pure-loss}
\metasep newly derived
\metasep accepted}

% claim receiver/def:embed
\claimsummary{receiver/def:embed}{A local element is a contraction on a subset of the modes, its support, and its global embedding acts as the identity on every other mode. Embedding is what lets component descriptions compose as matrices on the whole mode space.}
\claimbookmark{receiver/def:embed}{Definition}{def:embed}{Local element and its embedding}
\begin{definition}[Local element and its embedding]\label{def:embed}
A \termdef{term-local-element}{local element} on a mode subset $\Sigma\subseteq\Modes$ is a matrix
$\Gmat_\Sigma\in\mathbb{C}^{\Sigma\times\Sigma}$ with
$\Gmat_\Sigma^\dagger\Gmat_\Sigma\preceq\id_\Sigma$; $\Sigma$ is its
\termdef{term-support}{support}. Its \termdef{term-global-embedding}{global embedding} is
\begin{equation}\label{eq:embed}
  \hat\Gmat\coloneqq\Gmat_\Sigma\oplus\id_{\Modes\setminus\Sigma}
  \ \in\ \mathbb{C}^{\Modes\times\Modes},
\end{equation}
that is, $\hat\Gmat\,\Emd{\Sigma}=\Gmat_\Sigma\Emd{\Sigma}$ and
$\hat\Gmat\,(\id-\Emd{\Sigma})=\id-\Emd{\Sigma}$. A local element with
$\Gmat_\Sigma^\dagger\Gmat_\Sigma=\id_\Sigma$ is a \termdef{term-lossless-element}{lossless element}.
\end{definition}
\claimmeta{\metalabel{Used by}\mref{Prop.}{prop:mode-algebra}, \mref{Def.}{def:catalogue}, \mref{Thm.}{thm:catalogue-admissibility}, \mref{Lem.}{lem:catalogue-dilations}
\metasep newly derived
\metasep accepted}

% claim receiver/prop:mode-algebra
\claimsummary{receiver/prop:mode-algebra}{Embedding preserves contractivity, products of contractions are contractions, elements on disjoint supports commute, and embedding respects adjoints and unitarity. Graph compilation rests on these four facts.}
\claimbookmark{receiver/prop:mode-algebra}{Proposition}{prop:mode-algebra}{Algebra of embedded local elements}
\begin{proposition}[Algebra of embedded local elements]\label{prop:mode-algebra}
Let $\hat\Gmat,\hat\Gmat'$ be \termref{term-global-embedding}{global embeddings} of
\termref{term-local-element}{local elements} with supports $\Sigma,\Sigma'$.
\begin{enumerate}[label=(\roman*),leftmargin=3.0em]
\item \emph{Embedding preserves contractivity:}
  $\hat\Gmat^\dagger\hat\Gmat\preceq\id_{\Modes}$, with equality exactly when
  the element is lossless.
\item \emph{Products of contractions are contractions:} if
  $X^\dagger X\preceq\id$ and $Y^\dagger Y\preceq\id$ then
  $(XY)^\dagger(XY)\preceq\id$, with equality for all vectors exactly when
  $Y$ is an isometry and $X$ is isometric on $\operatorname{ran}Y$.
\item \emph{Disjoint supports commute:} if $\Sigma\cap\Sigma'=\varnothing$ then
  $\hat\Gmat\hat\Gmat'=\hat\Gmat'\hat\Gmat$.
\item \emph{Adjoints and unitaries:} $\hat\Gmat^\dagger$ is the embedding of
  $\Gmat_\Sigma^\dagger$, and $\hat\Gmat$ is unitary exactly when
  $\Gmat_\Sigma$ is.
\end{enumerate}
\end{proposition}
\claimmeta{\metalabel{Uses}\mref{Def.}{def:modes}, \mref{Def.}{def:embed}
\metasep \metalabel{Used by}\mref{Thm.}{thm:compilation-well-defined}, \mref{Thm.}{thm:efficiency}, \mref{Def.}{def:normalized-station}
\metasep newly derived
\metasep accepted
\metasep \metalabel{Oracles}\coderef{https://github.com/fraunhofersg/certiqs-key/blob/v0.1.0/tests/oracles/registry_receiver_povm.py}{registry_receiver_povm.py}}

\begin{proof}
(i) In the block decomposition $\mathbb{C}^{\Modes}=\mathbb{C}^{\Sigma}\oplus
\mathbb{C}^{\Modes\setminus\Sigma}$ one has
$\hat\Gmat^\dagger\hat\Gmat=(\Gmat_\Sigma^\dagger\Gmat_\Sigma)
\oplus\id_{\Modes\setminus\Sigma}$, which is $\preceq\id$ precisely when
$\Gmat_\Sigma^\dagger\Gmat_\Sigma\preceq\id_\Sigma$, and equals $\id$ precisely
when $\Gmat_\Sigma^\dagger\Gmat_\Sigma=\id_\Sigma$.

(ii) For any vector $v$, $\lVert XYv\rVert\le\lVert Yv\rVert\le\lVert v\rVert$,
using $\lVert Xw\rVert\le\lVert w\rVert$ at $w=Yv$ and then
$\lVert Yv\rVert\le\lVert v\rVert$. Hence
$v^\dagger(XY)^\dagger(XY)v\le v^\dagger v$ for all $v$, which is
$(XY)^\dagger(XY)\preceq\id$. Equality for all $v$ forces equality in both
steps for all $v$: the second gives $\lVert Yv\rVert=\lVert v\rVert$ for all
$v$, i.e.\ $Y$ isometric, and the first then gives
$\lVert Xw\rVert=\lVert w\rVert$ for every $w\in\operatorname{ran}Y$.
Conversely, if $Y$ is an isometry and $X$ is isometric on
$\operatorname{ran}Y$, then $\lVert XYv\rVert=\lVert Yv\rVert=\lVert v\rVert$
for every $v$, which is equality for all vectors.

(iii) Write $\mathbb{C}^{\Modes}=\mathbb{C}^{\Sigma}\oplus\mathbb{C}^{\Sigma'}
\oplus\mathbb{C}^{\Modes\setminus(\Sigma\cup\Sigma')}$, legitimate because
$\Sigma\cap\Sigma'=\varnothing$. In that decomposition
$\hat\Gmat=\Gmat_\Sigma\oplus\id\oplus\id$ and
$\hat\Gmat'=\id\oplus\Gmat'_{\Sigma'}\oplus\id$; both products equal
$\Gmat_\Sigma\oplus\Gmat'_{\Sigma'}\oplus\id$.

(iv) Immediate from the block form.
\end{proof}

% claim receiver/rem:overlap
\claimsummary{receiver/rem:overlap}{The contract's only commutation result, \hyperlink{claim-receiver/prop:mode-algebra}{Proposition~\ref*{prop:mode-algebra}} (iii), needs disjoint mode supports; being unrelated in a declared graph is not enough. Unrelated elements that share a mode can depend on their order, which is why a well-formed graph must order them.}
\claimbookmark{receiver/rem:overlap}{Remark}{rem:overlap}{Disjoint mode supports as the commutation hypothesis}
\begin{remark}[Disjoint mode supports as the commutation hypothesis]\label{rem:overlap}
Proposition~\ref{prop:mode-algebra}\,(iii) is the \emph{only} commutation this
contract has, and it needs disjointness of the \emph{mode supports}, not
disjointness in the wiring diagram. Two elements that are unrelated in a
declared graph may still touch a common mode, and then their product can
depend on the order. That is exactly why Definition~\ref{def:graph} carries
condition~(E3) and why Theorem~\ref{thm:compilation-well-defined} is a theorem rather than a
formality.
\end{remark}
\claimmeta{newly derived
\metasep accepted}

% claim receiver/lem:fock-functor
\claimsummary{receiver/lem:fock-functor}{The Fock lift respects composition and adjoints, preserves every photon-number sector, and takes unitaries, contractions and projectors to unitaries, contractions and projectors. The lift of a unitary implements the mode transformation on the annihilation operators.}
\claimbookmark{receiver/lem:fock-functor}{Lemma}{lem:fock-functor}{Functoriality of the Fock lift}
\claimleadin{%
The later derivations work block by block on the lift of
Definition~\ref{def:fock}; this lemma collects the properties of the lift they
use.
}
\begin{lemma}[Functoriality of the Fock lift]\label{lem:fock-functor}
Let $\mathcal{K}_0$ be finite-dimensional and let $X,Y$ be operators on
$\mathcal{K}_0$. Then, with $\Fock$ as in \eqref{eq:fockfunctor}:
\begin{enumerate}[label=(\roman*),leftmargin=3.0em]
\item $X^{\otimes n}$ preserves $\Symn{n}(\mathcal{K}_0)$, so $\Fock(X)$ is
  well defined, and it is block diagonal with respect to the photon-number
  decomposition: $\Fock(X)\Pnum{n}=\Pnum{n}\Fock(X)$ for every $n$;
\item $\Fock(XY)=\Fock(X)\Fock(Y)$, $\Fock(\id)=\id$, and
  $\Fock(X)^\dagger=\Fock(X^\dagger)$;
\item if $U$ is unitary then $\Fock(U)$ is unitary; if
  $X^\dagger X\preceq\id$ then $\lVert\Fock(X)\rVert\le1$, so $\Fock(X)$ is a
  bounded operator on $\mathcal{F}(\mathcal{K}_0)$;
\item if $P$ is an orthogonal projector then $\Fock(P)$ is the orthogonal
  projector onto $\mathcal{F}(\operatorname{ran}P)\subseteq
  \mathcal{F}(\mathcal{K}_0)$ --- equivalently, with the occupation vectors
  formed in an orthonormal basis of $\mathcal{K}_0$ adapted to
  $\operatorname{ran}P\oplus\ker P$, onto the closed span of those carrying
  no photon outside $\operatorname{ran}P$;
\item if $U$ is unitary then $\Fock(U)$ implements the mode transformation on
  the annihilation operators,
  $\Fock(U)^\dagger a_i\Fock(U)=\sum_jU_{ij}a_j$.
\end{enumerate}
\end{lemma}
\claimmeta{\metalabel{Uses}\mref{Def.}{def:fock}
\metasep \metalabel{Used by}\mref{Lem.}{lem:fock-maps}, \mref{Rem.}{rem:migdal}, \mref{Lem.}{lem:fock-block-spectrum}, \mref{Lem.}{lem:fock-monotone}, \mref{Prop.}{prop:normord}, \mref{Rem.}{rem:konig}, \mref{Rem.}{rem:whyprojector}, \mref{Def.}{def:physmodel}, \mref{Lem.}{lem:vacelim}, \mref{Thm.}{thm:port-grams}, \mref{Thm.}{thm:click-povm}, \mref{Lem.}{lem:exactbasic}, \mref{Lem.}{lem:fock-permanent}, \mref{Thm.}{thm:factor-blocks}, \mref{Ex.}{ex:bb84exact}, \mref{Prop.}{prop:noise-floor}
\metasep newly derived
\metasep implemented (effective: accepted)
\metasep \metalabel{Code}\coderef{https://github.com/fraunhofersg/certiqs-key/blob/v0.1.0/src/certiqs_key/receiver/fock.py}{certiqs_key.receiver.fock:fock_block}
\metasep \metalabel{Oracles}\coderef{https://github.com/fraunhofersg/certiqs-key/blob/v0.1.0/tests/oracles/registry_receiver_povm.py}{registry_receiver_povm.py}}

\begin{proof}
(i) For a permutation $\pi$ of $n$ factors let $\Pi_\pi$ be the corresponding
unitary on $\mathcal{K}_0^{\otimes n}$. Then $\Pi_\pi X^{\otimes n}=
X^{\otimes n}\Pi_\pi$, because both sides send
$v_1\otimes\cdots\otimes v_n$ to the tensor product of the $Xv_i$ permuted by
$\pi$. Hence $X^{\otimes n}$ commutes with the orthogonal projector
$P_{\Symn{n}}=\frac{1}{n!}\sum_\pi\Pi_\pi$ and therefore preserves its range.
Block diagonality in $n$ is the definition of $\Fock$ as a direct sum.

(ii) On $\Symn{n}$, $(XY)^{\otimes n}=X^{\otimes n}Y^{\otimes n}$ and both
factors preserve $\Symn{n}$, so the restriction of the product is the product
of the restrictions; $\id^{\otimes n}=\id$; and
$(X^{\otimes n})^\dagger=(X^\dagger)^{\otimes n}$, whose restriction to the
invariant subspace $\Symn{n}$ is the adjoint of the restriction. Summing over
$n$ gives the three identities.

(iii) With $U$ unitary, (ii) gives
$\Fock(U)^\dagger\Fock(U)=\Fock(U^\dagger U)=\id$ and likewise on the other
side. With $X^\dagger X\preceq\id$, $\lVert X\rVert\le1$, so
$\lVert X^{\otimes n}\rVert=\lVert X\rVert^n\le1$; restriction to a subspace
does not increase the norm, so every block of $\Fock(X)$ has norm $\le1$ and
the direct sum is bounded by $1$.

(iv) $P^{\otimes n}$ is an orthogonal projector on $\mathcal{K}_0^{\otimes n}$
with range $(\operatorname{ran}P)^{\otimes n}$, and it commutes with
$P_{\Symn{n}}$ by the computation in (i). The product of two commuting
orthogonal projectors is the orthogonal projector onto the intersection of
their ranges, so the restriction of $P^{\otimes n}$ to $\Symn{n}$ is the
orthogonal projector onto
$\Symn{n}(\mathcal{K}_0)\cap(\operatorname{ran}P)^{\otimes n}
=\Symn{n}(\operatorname{ran}P)$. Choosing an orthonormal basis of
$\mathcal{K}_0$ adapted to $\operatorname{ran}P\oplus\ker P$ identifies that
subspace with the span of the occupation vectors supported on
$\operatorname{ran}P$.

(v) Fix the orthonormal basis $\{\ket{e_j}\}$ used in
Definition~\ref{def:fock}. The creation operator satisfies
$a_i^\dagger\ket{\psi}=\sqrt{n+1}\,P_{\Symn{n+1}}(\ket{e_i}\otimes\ket{\psi})$
for $\ket{\psi}\in\Symn{n}$, which is the standard identification of the mode
description with the symmetric-tensor description. For unitary $U$,
$\Fock(U)P_{\Symn{n+1}}=P_{\Symn{n+1}}U^{\otimes(n+1)}$ on
$\mathcal{K}_0^{\otimes(n+1)}$ by (i), hence
\begin{align}
  \Fock(U)\,a_i^\dagger\ket{\psi}
  &=\sqrt{n+1}\;P_{\Symn{n+1}}\big(U\ket{e_i}\otimes U^{\otimes n}\ket{\psi}\big)
  \nonumber\\
  &=\sum_j U_{ji}\,\sqrt{n+1}\;P_{\Symn{n+1}}
      \big(\ket{e_j}\otimes\Fock(U)\ket{\psi}\big)
  =\Big(\sum_jU_{ji}a_j^\dagger\Big)\Fock(U)\ket{\psi},
\end{align}
using $U\ket{e_i}=\sum_jU_{ji}\ket{e_j}$. Therefore
$\Fock(U)a_i^\dagger\Fock(U)^\dagger=\sum_jU_{ji}a_j^\dagger$; taking adjoints
gives $\Fock(U)a_i\Fock(U)^\dagger=\sum_j\overline{U_{ji}}\,a_j$, and
replacing $U$ by $U^\dagger$ gives
$\Fock(U)^\dagger a_i\Fock(U)=\sum_j\overline{(U^\dagger)_{ji}}a_j
=\sum_jU_{ij}a_j$.
\end{proof}

% claim receiver/lem:fock-maps
\claimsummary{receiver/lem:fock-maps}{The Fock lift respects products and adjoints of maps between two mode spaces and takes contractions to contractions and isometries to isometries. On disjoint mode sets the occupation identification is unitary, and the lift of a direct sum is the tensor product of the lifts.}
\claimbookmark{receiver/lem:fock-maps}{Lemma}{lem:fock-maps}{The lift of maps between mode spaces and of direct sums}
\claimleadin{%
Lemma~\ref{lem:fock-functor} is stated for operators on one space. A loss in
front of a receiver is an isometry into a larger mode space, so the same
properties are needed for maps between two spaces, together with the
identification of the Fock space over two disjoint mode sets with the tensor
product of their Fock spaces.
}
\begin{technicalclaim}
\begin{lemma}[The lift of maps between mode spaces and of direct sums]\label{lem:fock-maps}
Let $\mathcal{K}_0,\mathcal{K}_1,\mathcal{K}_2$ be finite-dimensional, and let
$X\colon\mathcal{K}_0\to\mathcal{K}_1$ and $Y\colon\mathcal{K}_1\to\mathcal{K}_2$
be linear maps.
\begin{enumerate}[label=(\roman*),leftmargin=3.0em]
\item $\Fock(YX)=\Fock(Y)\Fock(X)$ and $\Fock(X)^{\dagger}=\Fock(X^{\dagger})$;
  if $X^{\dagger}X\preceq\id$ then $\lVert\Fock(X)\rVert\le1$, and if
  $X^{\dagger}X=\id$ then $\Fock(X)$ is an isometry.
\item Let $S$ and $T$ be disjoint finite mode sets. The linear map
  $\mathcal{F}(\mathbb{C}^{S})\otimes\mathcal{F}(\mathbb{C}^{T})\to
  \mathcal{F}(\mathbb{C}^{S\sqcup T})$ defined on occupation vectors
  (Definition~\ref{def:fock}) by
  $\ket{\vec m_1}\otimes\ket{\vec m_2}\mapsto\ket{\vec m_1,\vec m_2}$, the
  \emph{occupation identification},
  is unitary, and on its $n$-photon level it is
  $\ket{\varphi}\otimes\ket{\chi}\mapsto\binom{n}{p}^{1/2}P_n(\varphi\otimes\chi)$
  for $\varphi\in\Symn{p}(\mathbb{C}^{S})$, $\chi\in\Symn{q}(\mathbb{C}^{T})$,
  $p+q=n$, with $P_n$ the orthogonal projector of
  $(\mathbb{C}^{S\sqcup T})^{\otimes n}$ onto its symmetric subspace.
\item Under that identification, for operators $X_1$ on $\mathbb{C}^{S}$ and
  $X_2$ on $\mathbb{C}^{T}$,
  $\Fock(X_1\oplus X_2)=\Fock(X_1)\otimes\Fock(X_2)$.
\end{enumerate}
\end{lemma}
\end{technicalclaim}
\claimmeta{\metalabel{Uses}\mref{Def.}{def:fock}, \mref{Lem.}{lem:fock-functor}
\metasep \metalabel{Used by}\mref{Def.}{def:physmodel}, \mref{Lem.}{lem:vacelim}, \mref{Def.}{def:pure-loss}, \mref{Lem.}{lem:pure-loss-lift}
\metasep newly derived
\metasep accepted}

\begin{proof}
(i) $X^{\otimes n}$ maps $\Symn{n}(\mathcal{K}_0)$ into
$\Symn{n}(\mathcal{K}_1)$ (Definition~\ref{def:fock}), and so does
$(X^{\dagger})^{\otimes n}$ from $\Symn{n}(\mathcal{K}_1)$ back, both
intertwining the permutation actions. Then $(YX)^{\otimes n}=Y^{\otimes n}X^{\otimes n}$
restricted to $\Symn{n}(\mathcal{K}_0)$ is the product of the restrictions,
and the restriction of $(X^{\otimes n})^{\dagger}=(X^{\dagger})^{\otimes n}$ to
$\Symn{n}(\mathcal{K}_1)$ is the adjoint of the restriction of
$X^{\otimes n}$, since both restrictions map between the invariant subspaces;
summing over $n$ gives the two identities, as in
Lemma~\ref{lem:fock-functor}\,(ii). If $X^{\dagger}X\preceq\id$ then
$\lVert X\rVert\le1$, so every block $X^{\otimes n}$ has norm at most one, as
in Lemma~\ref{lem:fock-functor}\,(iii). If $X^{\dagger}X=\id$, the two
identities give $\Fock(X)^{\dagger}\Fock(X)=\Fock(X^{\dagger}X)=\Fock(\id)=\id$.

(ii) For an occupation vector $\vec m$ with $|\vec m|=n$ write
$e^{\otimes\vec m}$ for the tensor product of $m_1$ copies of the first basis
vector, then $m_2$ of the second, and so on, and $\vec m!\coloneqq\prod_jm_j!$.
$P_ne^{\otimes\vec m}$ is $\vec m!/n!$ times the sum of the $n!/\vec m!$
distinct reorderings of $e^{\otimes\vec m}$, which are orthonormal, so
$\lVert P_ne^{\otimes\vec m}\rVert^{2}=\vec m!/n!$ and the normalized
symmetrization of Definition~\ref{def:fock} is
$\ket{\vec m}=(n!/\vec m!)^{1/2}P_ne^{\otimes\vec m}$. Let $J_{p,q}$ be the
displayed map on the level $p+q=n$. Since $P_n(P_p\otimes P_q)=P_n$, the
product $P_p\otimes P_q$ averaging over a subgroup of the permutations, and a
reordering of the factors is a permutation,
\begin{equation*}
  J_{p,q}\big(\ket{\vec m_1}\otimes\ket{\vec m_2}\big)
  =\Big(\tbinom{n}{p}\tfrac{p!\,q!}{\vec m_1!\,\vec m_2!}\Big)^{1/2}
   P_n\big(e^{\otimes\vec m_1}\otimes e^{\otimes\vec m_2}\big)
  =\ket{\vec m_1,\vec m_2}.
\end{equation*}
So $\bigoplus_{p+q=n}J_{p,q}$ maps an orthonormal basis of the $n$-photon
level of the tensor product onto one of $\Symn{n}(\mathbb{C}^{S\sqcup T})$,
and the identification, their sum over $n$, is unitary.

(iii) $(X_1\oplus X_2)^{\otimes n}$ commutes with $P_n$
(Lemma~\ref{lem:fock-functor}\,(i)) and acts on
$(\mathbb{C}^{S})^{\otimes p}\otimes(\mathbb{C}^{T})^{\otimes q}$ as
$X_1^{\otimes p}\otimes X_2^{\otimes q}$, so
$(X_1\oplus X_2)^{\otimes n}J_{p,q}(\varphi\otimes\chi)
=J_{p,q}\big(X_1^{\otimes p}\varphi\otimes X_2^{\otimes q}\chi\big)$ for
$\varphi\in\Symn{p}(\mathbb{C}^{S})$ and $\chi\in\Symn{q}(\mathbb{C}^{T})$,
which is the claim on every level.
\end{proof}

% claim receiver/rem:migdal
\claimsummary{receiver/rem:migdal}{The lift of a unitary acts on n photons as its n-fold tensor power on the symmetric subspace, the standard dictionary between the mode and particle descriptions. A published statement of it corroborates the lemma; no claim depends on the citation.}
\claimbookmark{receiver/rem:migdal}{Remark}{rem:migdal}{Mode and particle descriptions in the literature}
\begin{remark}[Mode and particle descriptions in the literature]\label{rem:migdal}
Lemma~\ref{lem:fock-functor}\,(v) with (i) is the standard dictionary between the
mode (second-quantized) and particle (symmetric-tensor) descriptions of $n$
photons in $d$ modes: a mode transformation $a_i\mapsto\sum_jU_{ij}a_j$ acts on
the $n$-photon sector as $U^{\otimes n}$ restricted to $\Symn{n}$. The same
correspondence is recorded in \cite[Eqs.~(9)--(10)]{migdal2014multiphoton},
where Eq.~(9) is the mode-level transformation of the creation operators and
Eq.~(10) states that it is ``equivalent to the action of the same $U$ on each
particle'', $\ket{\psi'}_P=U^{\otimes n}\ket{\psi}_P$ in the particle
representation. That reference is corroboration only: the lemma is proved in
full, and no claim of this contract depends on the reference.
\end{remark}
\claimmeta{\metalabel{Uses}\mref{Lem.}{lem:fock-functor}
\metasep newly derived
\metasep accepted}

% claim receiver/lem:fock-block-spectrum
\claimsummary{receiver/lem:fock-block-spectrum}{On the n-photon block, the lift of 1 \ensuremath{-} A has as eigenvalues the products of n factors 1 \ensuremath{-} \ensuremath{\lambda} over the eigenvalues \ensuremath{\lambda} of A, so it lies between the n-th powers of 1 \ensuremath{-} \ensuremath{\lambda}max and 1 \ensuremath{-} \ensuremath{\lambda}min, both attained. The analytic high-photon bound of the receiver starts from these two extremes.}
\claimbookmark{receiver/lem:fock-block-spectrum}{Lemma}{lem:fock-block-spectrum}{Spectrum of a lifted no-click operator on a photon-number block}
\claimleadin{%
For $A$ the sum of the \termref{term-port-gram-matrices}{port Grams} of a set of ports,
$\Fock(\id-A)$ is the operator for the event that no port of the set clicks,
pulled back to the declared input modes (Theorem~\ref{thm:port-grams}\,(iii));
hence the name lifted no-click operator.
}
\begin{lemma}[Spectrum of a lifted no-click operator on a photon-number block]
\label{lem:fock-block-spectrum}
Let $A$ be Hermitian on $\mathcal{K}_0$ with $0\preceq A\preceq\id$ and
eigenvalues $\lambda_1,\dots,\lambda_d$ ($d=\dim\mathcal{K}_0$). Then
$\Fock(\id-A)$ is a positive semidefinite contraction, its $n$-photon block
$\Fock(\id-A)^{(n)}=\big((\id-A)^{\otimes n}\big)\big|_{\Symn{n}}$ has spectrum
\begin{equation}\label{eq:spectrum}
  \Big\{\textstyle\prod_{j=1}^{d}(1-\lambda_j)^{m_j}\ :\
  \vec m\in\mathbb{N}_0^{d},\ |\vec m|=n\Big\},
\end{equation}
and consequently
\begin{equation}\label{eq:blockbounds}
  \big(1-\lmax(A)\big)^{n}\,\id\ \preceq\ \Fock(\id-A)^{(n)}\ \preceq\
  \big(1-\lmin(A)\big)^{n}\,\id ,
\end{equation}
both bounds being attained as eigenvalues.
\end{lemma}
\claimmeta{\metalabel{Uses}\mref{Lem.}{lem:fock-functor}
\metasep \metalabel{Used by}\mref{Prop.}{prop:normord}, \mref{Thm.}{thm:coarse-graining}, \mref{Def.}{def:bounddata}, \mref{Thm.}{thm:multiclick-bound}, \mref{Lem.}{lem:exactbasic}, \mref{Rem.}{rem:rmintrap}, \mref{Prop.}{prop:noise-floor}
\metasep newly derived
\metasep implemented (effective: accepted)
\metasep \metalabel{Code}\coderef{https://github.com/fraunhofersg/certiqs-key/blob/v0.1.0/src/certiqs_key/receiver/fock.py}{certiqs_key.receiver.fock:block_spectrum_from_eigenvalues}
\metasep \metalabel{Oracles}\coderef{https://github.com/fraunhofersg/certiqs-key/blob/v0.1.0/tests/oracles/registry_receiver_povm.py}{registry_receiver_povm.py}}

\begin{proof}
Diagonalize $A=U^\dagger\Lambda U$ with $U$ unitary and
$\Lambda=\operatorname{diag}(\lambda_1,\dots,\lambda_d)$; then
$\id-A=U^\dagger(\id-\Lambda)U$ and Lemma~\ref{lem:fock-functor}\,(ii),(iii) give
$\Fock(\id-A)=\Fock(U)^\dagger\Fock(\id-\Lambda)\Fock(U)$ with $\Fock(U)$
unitary and photon-number preserving, so the two operators have the same
spectrum block by block. In the occupation basis built from the eigenbasis of
$A$, the product vector $\ket{e_{j_1}}\otimes\cdots\otimes\ket{e_{j_n}}$ is an
eigenvector of $(\id-\Lambda)^{\otimes n}$ with eigenvalue
$\prod_r(1-\lambda_{j_r})$, and that eigenvalue is unchanged by permuting the
factors; hence the symmetrized basis vector $\ket{\vec m}$ is an eigenvector of
$(\id-\Lambda)^{\otimes n}|_{\Symn{n}}$ with eigenvalue
$\prod_j(1-\lambda_j)^{m_j}$. These $\ket{\vec m}$, $|\vec m|=n$, form an
orthonormal basis of $\Symn{n}$, so \eqref{eq:spectrum} is the full spectrum.
Each factor satisfies $1-\lambda_j\in[1-\lmax(A),1-\lmin(A)]\subseteq[0,1]$,
so a product of $n$ of them lies in
$[(1-\lmax(A))^n,(1-\lmin(A))^n]$, which is \eqref{eq:blockbounds}; the
endpoints are attained by the occupation vectors putting all $n$ photons on an
eigenvector for $\lmax(A)$, respectively for $\lmin(A)$. Positivity and
contractivity follow from \eqref{eq:blockbounds} and
Lemma~\ref{lem:fock-functor}\,(iii).
\end{proof}

% claim receiver/lem:fock-monotone
\claimsummary{receiver/lem:fock-monotone}{The Fock lift preserves the Loewner order on positive semidefinite operators. A larger detection operator A therefore gives a smaller lifted no-click operator, the lift of 1 \ensuremath{-} A.}
\claimbookmark{receiver/lem:fock-monotone}{Lemma}{lem:fock-monotone}{Order monotonicity of the Fock lift}
\begin{lemma}[Order monotonicity of the Fock lift]\label{lem:fock-monotone}
If $X\succeq Y\succeq0$ on $\mathcal{K}_0$ then $\Fock(X)\succeq\Fock(Y)$.
Consequently, if $0\preceq A\preceq A'\preceq\id$ then
$\Fock(\id-A)\succeq\Fock(\id-A')$.
\end{lemma}
\claimmeta{\metalabel{Uses}\mref{Lem.}{lem:fock-functor}
\metasep \metalabel{Used by}\mref{Rem.}{rem:positivity}, \mref{Lem.}{lem:exactbasic}
\metasep newly derived
\metasep accepted
\metasep \metalabel{Oracles}\coderef{https://github.com/fraunhofersg/certiqs-key/blob/v0.1.0/tests/oracles/registry_receiver_povm.py}{registry_receiver_povm.py}}

\begin{proof}
Telescoping,
\begin{equation}\label{eq:telescope}
  X^{\otimes n}-Y^{\otimes n}
  =\sum_{j=0}^{n-1}X^{\otimes j}\otimes(X-Y)\otimes Y^{\otimes(n-1-j)} .
\end{equation}
Each summand is a tensor product of positive semidefinite operators and is
therefore positive semidefinite, so $X^{\otimes n}\succeq Y^{\otimes n}$ on
$\mathcal{K}_0^{\otimes n}$. Both $X^{\otimes n}$ and $Y^{\otimes n}$ preserve
$\Symn{n}$ by Lemma~\ref{lem:fock-functor}\,(i), so for
$\ket{\psi}\in\Symn{n}$ the difference of the restrictions is
$\bra{\psi}(X^{\otimes n}-Y^{\otimes n})\ket{\psi}\ge0$: the restriction of a
positive semidefinite operator to an invariant subspace is positive
semidefinite. Summing over $n$ gives $\Fock(X)\succeq\Fock(Y)$. The second
statement is the first with $X=\id-A$, $Y=\id-A'$, both positive semidefinite
because $A,A'\preceq\id$.
\end{proof}

% claim receiver/prop:normord
\claimsummary{receiver/prop:normord}{The lift of 1 \ensuremath{-} A equals the normally ordered exponential of minus the quadratic form of A in the mode operators, a series that is a finite sum on each photon-number block. It is the form of the optics literature; no proof of this contract uses it.}
\claimbookmark{receiver/prop:normord}{Proposition}{prop:normord}{Normal-ordered form of a lifted no-click operator}
\claimleadin{%
Remark~\ref{rem:whyprojector} says why the derivations of this contract do
without this form.
}
\begin{proposition}[Normal-ordered form of a lifted no-click operator]\label{prop:normord}
For $0\preceq A\preceq\id$ on $\mathcal{K}_0$,
\begin{equation}\label{eq:normordform}
  \Fock(\id-A)=\normord{\exp\!\big(-\vec a^\dagger A\vec a\big)}
  \coloneqq\sum_{p\ge0}\frac{(-1)^p}{p!}
  \sum_{i_1\dots i_p}\sum_{j_1\dots j_p}
  \Big(\prod_{r=1}^{p}A_{i_rj_r}\Big)
  a_{i_1}^\dagger\cdots a_{i_p}^\dagger a_{j_1}\cdots a_{j_p},
\end{equation}
the series acting on each photon-number block as a finite sum.
\end{proposition}
\claimmeta{\metalabel{Uses}\mref{Lem.}{lem:fock-functor}, \mref{Lem.}{lem:fock-block-spectrum}
\metasep \metalabel{Used by}\mref{Rem.}{rem:konig}, \mref{Rem.}{rem:whyprojector}, \mref{Rem.}{rem:normordelements}, \mref{Rem.}{rem:sperlingeta}
\metasep newly derived
\metasep accepted}

\begin{proof}
Both sides are unchanged in form under $A\mapsto UAU^\dagger$ combined with
conjugation by $\Fock(U)$: for the left side by
Lemma~\ref{lem:fock-functor}\,(ii), for the right side because
Lemma~\ref{lem:fock-functor}\,(v) maps creation operators to combinations of
creation operators and annihilation operators to combinations of annihilation
operators, so no creator crosses an annihilator and the normal-ordered form
survives with $A$ replaced by $UAU^\dagger$. It therefore suffices to prove
\eqref{eq:normordform} for diagonal $A=\Lambda$. For diagonal $\Lambda$ only
the terms with $i_r=j_r$ survive, operators of distinct modes commute, and
collecting the $p!/\vec p\,!$ tuples with a given multiplicity vector $\vec p$
(where $\vec p\,!\coloneqq\prod_jp_j!$) turns the right side into
$\prod_j\sum_{p_j\ge0}\frac{(-\lambda_j)^{p_j}}{p_j!}
a_j^{\dagger p_j}a_j^{p_j}$. On an occupation vector, the standard identity
$a_j^{\dagger p}a_j^{p}\ket{\vec m}=\frac{m_j!}{(m_j-p)!}\ket{\vec m}$ for
$p\le m_j$ and $0$ otherwise gives, by the binomial theorem,
\begin{equation}\label{eq:singlemode}
  \sum_{p\ge0}\frac{(-\lambda_j)^{p}}{p!}a_j^{\dagger p}a_j^{p}\ket{\vec m}
  =\sum_{p=0}^{m_j}\binom{m_j}{p}(-\lambda_j)^p\ket{\vec m}
  =(1-\lambda_j)^{m_j}\ket{\vec m} ,
\end{equation}
so the right side of \eqref{eq:normordform} acts on $\ket{\vec m}$ as
$\prod_j(1-\lambda_j)^{m_j}$ --- which is exactly the action of
$\Fock(\id-\Lambda)$ computed in the proof of Lemma~\ref{lem:fock-block-spectrum}.
Both sides agree on an orthonormal basis, hence are equal.
\end{proof}

% claim receiver/rem:konig
\claimsummary{receiver/rem:konig}{The single-mode step of the normal-ordered form is the published identity that the normally ordered exponential of z times the number operator is (1 + z) to the number operator, and its \ensuremath{\lambda} = 1 case is the perfect on/off detector. Both citations corroborate; neither is used.}
\claimbookmark{receiver/rem:konig}{Remark}{rem:konig}{The single-mode normal-ordering identity in the literature}
\begin{remark}[The single-mode normal-ordering identity in the literature]\label{rem:konig}
Equation \eqref{eq:singlemode} is the single-mode statement
$\normord{e^{z\hat n}}=(1+z)^{\hat n}$ at $z=-\lambda_j$, where
$\hat n=a^\dagger a$. That operator identity is
\cite[Eq.~(3.23)]{konig2021newton}, which reads
$G_f(z;\hat n)=(1+z)^{\hat n}=\normord{e^{z\hat n}}$ and is proved there
without any assumption on the state of the field. Its $\lambda=1$ case,
$\normord{e^{-\hat n}}=\dyad{\vac}$, is the no-click element of a perfect
on/off detector, recorded as \cite[Eq.~(6)]{sperling2012} in the form
$\hat\pi_0=\dyad{0}=\normord{e^{-\hat n}}$, $\hat\pi_1=\normord{\hat 1-e^{-\hat
n}}$. Both are corroboration: \eqref{eq:singlemode} is proved in
Proposition~\ref{prop:normord}, and Theorem~\ref{thm:click-povm} derives the
no-click element from the projector picture of
Lemma~\ref{lem:fock-functor}\,(iv) without using normal ordering at all.
\end{remark}
\claimmeta{\metalabel{Uses}\mref{Lem.}{lem:fock-functor}, \mref{Prop.}{prop:normord}
\metasep newly derived
\metasep accepted}

% claim receiver/rem:whyprojector
\claimsummary{receiver/rem:whyprojector}{The normal-ordered form is stated for comparison and used in no proof. Every operator identity from the port Grams to the emitted weight follows from the functoriality of the Fock lift and vacuum elimination, block by block in finite dimension, with no convergence question.}
\claimbookmark{receiver/rem:whyprojector}{Remark}{rem:whyprojector}{Derivations without normal ordering}
\begin{remark}[Derivations without normal ordering]
\label{rem:whyprojector}
Proposition~\ref{prop:normord} is stated because it is the form in which these
operators appear in the optics literature and in every numerical
implementation, and because it makes the efficiency substitution of
Theorem~\ref{thm:efficiency} transparent. It is not used in the proof of any claim.
Every operator identity of the sections from the port Grams to the emitted
weight is obtained from Lemma~\ref{lem:fock-functor}\,(i)--(iv) and from the
vacuum elimination of Lemma~\ref{lem:vacelim}, which involve only
finite-dimensional linear algebra on each photon-number block: no unbounded
operator, no operator-valued series and no convergence question arises
anywhere in
Theorems~\ref{thm:port-grams}, \ref{thm:click-povm}, \ref{thm:representation-independence}, \ref{thm:coarse-graining},
\ref{thm:multiclick-bound} and~\ref{thm:emitted-bound}. A submodule whose
closed forms need series or integrals, as conditioning on a probe state can,
says so where it uses them.
\end{remark}
\claimmeta{\metalabel{Uses}\mref{Lem.}{lem:fock-functor}, \mref{Prop.}{prop:normord}
\metasep newly derived
\metasep accepted}

\section{The receiver instrument and its physical model}\label{group-receiver/def:instrument}
\sectionintro{A \termref{term-receiver-instrument}{receiver instrument} is a \termref{term-declared-mode-space}{declared mode space}, a split of the declared inputs into \termref{term-channel-modes}{channel modes}, through which the adversary's signal enters, and probe modes, disjoint \termref{term-detector-ports}{detector ports} and a contractive \termref{term-instrument-matrix}{instrument matrix}; its port Grams are the compact data the measurement is derived from. The probe modes carry the receiver's own probe light, in a declared state for each probe setting, which the receiver chooses privately in each round. A contraction becomes a physical device by a unitary dilation with vacuum environment modes, and the detectors are commuting threshold measurements on the ports. Vacuum on every undeclared input makes the restriction to the declared input columns exact (\hyperlink{claim-receiver/lem:vacelim}{Lemma~\ref*{lem:vacelim}}); a probe mode carries its declared state and is never eliminated.}

% claim receiver/def:instrument
\claimsummary{receiver/def:instrument}{A receiver instrument is a declared mode space with its inputs split into at least one channel mode and the probe modes, disjoint detector ports and a contractive instrument matrix. Undeclared inputs carry vacuum, and a mode in no port is undetected.}
\claimbookmark{receiver/def:instrument}{Definition}{def:instrument}{Receiver instrument}
\claimleadin{%
Besides the channel, a receiver may be fed by its own probe source, such as the
light source of a detector-attack countermeasure; its light enters through the
inputs $\InpS$, whose subscript $\Prb$ is the shared notation's label for the
probe input. The definition ends by identifying the shared notation's
channel-input system $\Hs{\Chan}$ with the \termref{term-fock-space}{Fock space} over $\InpC$. Its
photon-number blocks $\Hs{\Chan_n}$, grouped at a cutoff $N$ into the blocks
below and above it, are the canonical instance of the sectors that the
platform's security analysis declares on the receiver's space
(\href{../../../spec/notation.md}{\nolinkurl{spec/notation.md}}~\S2.4).
}
\begin{definition}[Receiver instrument]\label{def:instrument}
A \termdef{term-receiver-instrument}{receiver instrument} is a tuple
\begin{equation}\label{eq:instrument}
  \Recv=\Big(\Modes;\ \Inp=\InpC\sqcup\InpS;\
  \{\Port{k}\}_{k=1}^{K};\ \Gmat\Big)
\end{equation}
in which
\begin{enumerate}[label=(\roman*),leftmargin=3.0em]
\item $\Modes$ is a \termref{term-declared-mode-space}{declared mode space} (Definition~\ref{def:modes});
\item $\Inp\subseteq\Modes$ is the set of \termdef{term-declared-input-modes}{declared input modes}, split
  into the \termdef{term-channel-modes}{channel modes} $\InpC\neq\varnothing$ (the adversary's access) and the
  \termdef{term-declared-source-modes}{declared source modes} $\InpS$ (the probe); every mode of
  $\Modes\setminus\Inp$ carries vacuum at the input
  (Assumption~\ref{ass:vac});
\item $\Port{1},\dots,\Port{K}$ are pairwise disjoint subsets of $\Modes$, the
  \termdef{term-detector-ports}{detector ports}, with $K\ge1$; a mode in no port is undetected;
\item $\Gmat\in\mathbb{C}^{\Modes\times\Modes}$ is the \termdef{term-instrument-matrix}{instrument matrix},
  indexed rows-by-outputs and columns-by-inputs, with
  $\Gmat^\dagger\Gmat\preceq\id$.
\end{enumerate}
The basis $\{\ket{m}\}_{m\in\Inp}$ of $\mathbb{C}^{\Inp}$ is ordered with the
channel modes $\InpC$ before the declared source modes $\InpS$, each part in
the \exttermref{../foundations/contract.pdf}{term-declared-order}{declared order} of $\Modes$; matrices on $\mathbb{C}^{\Inp}$ and the
columns of a matrix with column set $\Inp$ are indexed in this order, which
differs from the declared order of $\Modes$ exactly when some declared source
mode precedes some channel mode in it.
We write $\Gmat_{\Inp}\in\mathbb{C}^{\Modes\times\Inp}$ for $\Gmat$ restricted
to the columns $\Inp$, and $\Emd{J}\coloneqq\Emd{\bigcup_{k\in J}\Port{k}}$ for
$J\subseteq[K]$, where $[K]\coloneqq\{1,\dots,K\}$.
The channel-input system of the shared notation is the \termref{term-fock-space}{Fock space} over
the channel modes: $\Hs{\Chan}=\bigoplus_n\Hs{\Chan_n}$ with
$\Hs{\Chan_n}=\Symn{n}(\mathbb{C}^{\InpC})$ and photon-number projectors
$\Pnumon{n}{\Chan}$.
\end{definition}
\claimmeta{\metalabel{Uses}\mref{Def.}{def:modes}, \mref{Def.}{def:fock}
\metasep \metalabel{Used by}\mref{Def.}{def:gram}, \mref{Def.}{def:probe}, \mref{Def.}{def:physmodel}, \mref{Def.}{def:compile}, \mref{Cor.}{cor:compiled-instrument}, \mref{Lem.}{lem:deficiency}, \mref{Def.}{def:coarse}, \mref{Cor.}{cor:unconditioned-family}, \mref{Cor.}{cor:emitted-unconditioned}, \mref{Thm.}{thm:enclosed-blocks}, \mref{Def.}{def:port-dilation}, \mref{Def.}{def:exact-route}, \mref{Def.}{def:tolmap}, \mref{Lem.}{lem:catalogue-dilations}, \mref{Prop.}{prop:noise-floor}, \mref{Def.}{def:pure-loss}, \mref{Def.}{def:normalized-station}
\metasep newly derived
\metasep accepted}

% claim receiver/def:gram
\claimsummary{receiver/def:gram}{The port Gram matrix of port k is the positive operator on the declared inputs whose quadratic form is the squared norm of the amplitude reaching that port, and the loss deficiency is the identity minus their sum. These Grams are the compact data the click measurement is derived from.}
\claimbookmark{receiver/def:gram}{Definition}{def:gram}{Port Grams and the loss deficiency}
\claimleadin{%
The admissibility defined at the end is an exact condition, while a declared
family is admitted by a floating-point test, in the \termref{term-evaluation-model}{evaluation model} of
Definition~\ref{def:evaluation-model}.
Theorem~\ref{thm:exact-and-evaluated} states what this contract asserts about
the exact family and what about its floating-point evaluation.
}
\begin{definition}[Port Grams and the loss deficiency]\label{def:gram}
The \termdef{term-port-gram-matrices}{port Gram matrices} of $\Recv$ are the operators on
$\mathbb{C}^{\Inp}$
\begin{equation}\label{eq:gramdef}
  A_k\coloneqq\Gmat_{\Inp}^\dagger\,\Emd{\{k\}}\,\Gmat_{\Inp},
  \qquad k=1,\dots,K ,
\end{equation}
and for $J\subseteq[K]$ we write $A_J\coloneqq\sum_{k\in J}A_k$ and
$A_\Sigma\coloneqq A_{[K]}$. The \termdef{term-loss-deficiency}{loss deficiency} is
\begin{equation}\label{eq:deficiency}
  \defic\coloneqq\id_{\Inp}-A_\Sigma .
\end{equation}
A family $\{A_k\}_{k=1}^{K}$ of operators on a finite-dimensional space is an
\termdef{term-admissible-gram-family}{admissible Gram family} when $A_k\succeq0$ for every $k$ and
$\sum_kA_k\preceq\id$. This is an exact condition; on a \emph{declared} family
it is \emph{decided} in floating point at the
\termref{term-receiver-tolerances}{receiver tolerance}
$\varepsilon_2$, by Definition~\ref{def:tolmap}\,(i).
\end{definition}
\claimmeta{\metalabel{Uses}\mref{Def.}{def:instrument}
\metasep \metalabel{Used by}\mref{Thm.}{thm:port-grams}, \mref{Lem.}{lem:deficiency}, \mref{Thm.}{thm:representation-independence}, \mref{Cor.}{cor:unconditioned-family}, \mref{Thm.}{thm:enclosed-blocks}, \mref{Def.}{def:port-dilation}, \mref{Thm.}{thm:factor-blocks}, \mref{Def.}{def:exact-route}, \mref{Def.}{def:tolmap}, \mref{Rem.}{rem:bound-data-band}, \mref{Req.}{req:exact-route-decision}, \mref{Lem.}{lem:catalogue-dilations}, \mref{Def.}{def:normalized-station}, \mref{Prop.}{prop:efficiency-normalization}
\metasep newly derived
\metasep implemented (effective: accepted)
\metasep \metalabel{Code}\coderef{https://github.com/fraunhofersg/certiqs-key/blob/v0.1.0/src/certiqs_key/receiver/grams.py}{certiqs_key.receiver.grams:GramFamily}}

% claim receiver/def:probe
\claimsummary{receiver/def:probe}{A probe declaration is a finite alphabet of probe settings, vacuum among them, with a declared probe state per setting that is diagonal in the occupation basis. In the Gaussian tier each state is a product of thermal states given by their mean photon numbers.}
\claimbookmark{receiver/def:probe}{Definition}{def:probe}{Probe declaration}
\claimleadin{%
A probe setting is classical data of known bias that the receiver chooses
privately in each round; for a detector-attack countermeasure it fixes
whether, and how brightly, the probe light shines. The bias, the weights $q_t$ of the shared
notation, belongs to the protocol declaration, and this contract uses only the
probe state of each setting. The word tier names a class of probe states here
and is unrelated to the tiers a reported key rate carries.
}
\begin{definition}[Probe declaration]\label{def:probe}
A \termdef{term-probe-declaration}{probe declaration} is a pair $\big(\Tset,\{\rho^{(t)}\}_{t\in\Tset}\big)$
in which $\Tset$ is a finite non-empty alphabet of probe settings (including
the setting $\off$, whose state is vacuum) and, for each $t$, $\rho^{(t)}$ is a
declared state on $\mathcal{F}(\mathbb{C}^{\InpS})$ that is
\termdef{term-occupation-diagonal}{occupation-diagonal}: diagonal in the occupation basis of
$\mathcal{F}(\mathbb{C}^{\InpS})$ (diagonal in the Fock basis). In the
\termdef{term-gaussian-tier}{Gaussian tier} each $\rho^{(t)}$ is a zero-mean Gaussian state with
diagonal normally-ordered second-moment matrix
$\Gamma^{(t)}=\operatorname{diag}(\bar n^{(t)}_m)_{m\in\InpS}\succeq0$, that is,
a product of single-mode thermal states with mean photon numbers
$\bar n^{(t)}_m\ge0$; $\Gamma^{(t)}=0$ is vacuum. A state of the Gaussian tier
is occupation-diagonal by construction, a thermal state being diagonal in the
occupation basis. When $\InpS=\varnothing$ the
declaration is $\Tset=\{\off\}$ and every statement about conditioning is
vacuous.
\end{definition}
\claimmeta{\metalabel{Uses}\mref{Def.}{def:fock}, \mref{Def.}{def:instrument}
\metasep \metalabel{Used by}\mref{Rem.}{rem:lossvsprobe}, \mref{Def.}{def:exact-route}, \mref{Req.}{req:probe-diagonal-admission}, \mref{Rem.}{rem:layers}, \mref{Ex.}{ex:bb84bound}
\metasep newly derived
\metasep accepted}

% claim receiver/lem:dilation
\claimsummary{receiver/lem:dilation}{Every contractive instrument matrix is the input-to-output block of a unitary on the modes together with a copy of them as environment modes. The amplitude the matrix removes leaves through those extra outputs, which makes the matrix a passive physical device.}
\claimbookmark{receiver/lem:dilation}{Lemma}{lem:dilation}{Unitary dilation of a contraction}
\claimleadin{%
An \termref{term-instrument-matrix}{instrument matrix} is a contraction, and a contraction is not by itself
the mode transformation of a physical device: the amplitude it removes has to
leave through some output, and the input ports it does not use have to be
filled with a state. The dilation of this lemma makes it physical, and
Lemma~\ref{lem:vacelim} then shows that the dilation changes nothing:
restricting to the declared input columns is \emph{exact}, not an
approximation.
}
\begin{lemma}[Unitary dilation of a contraction]\label{lem:dilation}
Let $\Gmat$ be a contraction on $\mathbb{C}^{\Modes}$ and let
$\Env$ be a disjoint copy of $\Modes$. With the defect operators
$D_{\mathrm{out}}\coloneqq(\id-\Gmat\Gmat^\dagger)^{1/2}$ and
$D_{\mathrm{in}}\coloneqq(\id-\Gmat^\dagger\Gmat)^{1/2}$, the matrix
\begin{equation}\label{eq:dilationmatrix}
  U_{\Gmat}=\begin{pmatrix}\Gmat & D_{\mathrm{out}}\\
                           D_{\mathrm{in}} & -\Gmat^\dagger\end{pmatrix}
  \quad\text{on }\ \mathbb{C}^{\Modes}\oplus\mathbb{C}^{\Env}
\end{equation}
is unitary, and its block from the inputs $\Modes$ to the outputs $\Modes$ is
$\Gmat$.
\end{lemma}
\claimmeta{\metalabel{Used by}\mref{Def.}{def:physmodel}, \mref{Thm.}{thm:port-grams}
\metasep newly derived
\metasep accepted}

\begin{proof}
Both defect operators exist: $\Gmat^\dagger\Gmat\preceq\id$ by hypothesis, and
$\Gmat\Gmat^\dagger$ has the same nonzero eigenvalues as
$\Gmat^\dagger\Gmat$, so $\Gmat\Gmat^\dagger\preceq\id$ as well; the positive
semidefinite square root of a positive semidefinite operator is unique. Write a
singular value decomposition $\Gmat=\sum_js_j\,y_jx_j^\dagger$ with
$s_j\in[0,1]$ and orthonormal $\{x_j\},\{y_j\}$. Then
$\Gmat^\dagger\Gmat=\sum_js_j^2x_jx_j^\dagger$ and
$\Gmat\Gmat^\dagger=\sum_js_j^2y_jy_j^\dagger$, so
$D_{\mathrm{in}}x_j=\sqrt{1-s_j^2}\,x_j$ and
$D_{\mathrm{out}}y_j=\sqrt{1-s_j^2}\,y_j$, both acting as the identity on the
orthogonal complements. Hence
\begin{equation}\label{eq:intertwine}
  D_{\mathrm{out}}\Gmat=\sum_js_j\sqrt{1-s_j^2}\,y_jx_j^\dagger
  =\Gmat D_{\mathrm{in}} .
\end{equation}
Multiplying \eqref{eq:dilationmatrix} by its adjoint blockwise and using
$D_{\mathrm{out}}^\dagger=D_{\mathrm{out}}$,
$D_{\mathrm{in}}^\dagger=D_{\mathrm{in}}$,
\begin{equation}
  U_{\Gmat}U_{\Gmat}^\dagger=
  \begin{pmatrix}
    \Gmat\Gmat^\dagger+D_{\mathrm{out}}^2 &
    \Gmat D_{\mathrm{in}}-D_{\mathrm{out}}\Gmat\\
    D_{\mathrm{in}}\Gmat^\dagger-\Gmat^\dagger D_{\mathrm{out}} &
    D_{\mathrm{in}}^2+\Gmat^\dagger\Gmat
  \end{pmatrix}
  =\begin{pmatrix}\id&0\\0&\id\end{pmatrix},
\end{equation}
the diagonal blocks by the definitions of the defect operators and the
off-diagonal blocks by \eqref{eq:intertwine} and its adjoint. A square matrix
with $U U^\dagger=\id$ is unitary.
\end{proof}

% claim receiver/def:physmodel
\claimsummary{receiver/def:physmodel}{The receiver's measurement is a unitary dilation of its instrument matrix applied to the declared inputs with vacuum everywhere else, followed by commuting threshold measurements on the detector ports. The input-side click POVM is derived from this model.}
\claimbookmark{receiver/def:physmodel}{Definition}{def:physmodel}{The physical measurement model}
\begin{definition}[The physical measurement model]\label{def:physmodel}
Fix a \termref{term-receiver-instrument}{receiver instrument} $\Recv$ and a unitary extension $U$ of $\Gmat$
on $\mathbb{C}^{\Modes\sqcup\Env}$ (Lemma~\ref{lem:dilation} supplies one; any
other may be used, and by Theorem~\ref{thm:port-grams}\,(v) the choice does not
matter). Let
$\iota:\mathbb{C}^{\Inp}\hookrightarrow\mathbb{C}^{\Modes\sqcup\Env}$ be the
inclusion of the \termref{term-declared-input-modes}{declared input modes} and
$\Fock(\iota):\mathcal{F}(\mathbb{C}^{\Inp})\to
\mathcal{F}(\mathbb{C}^{\Modes\sqcup\Env})$ the induced map, which under the
factorization
$\mathcal{F}(\mathbb{C}^{\Modes\sqcup\Env})\cong
 \mathcal{F}(\mathbb{C}^{\Inp})\otimes
 \mathcal{F}(\mathbb{C}^{(\Modes\sqcup\Env)\setminus\Inp})$, the occupation
identification of Lemma~\ref{lem:fock-maps}\,(ii),
is $\ket{\psi}\mapsto\ket{\psi}\otimes\ket{\vac}$: it prepares vacuum on every
undeclared input mode and on every dilation mode
(Assumption~\ref{ass:vac}). The \termdef{term-measurement-isometry}{measurement isometry} of the receiver is
\begin{equation}\label{eq:measisometry}
  V\coloneqq\Fock(U)\,\Fock(\iota)
  \ :\ \mathcal{F}(\mathbb{C}^{\Inp})\longrightarrow
        \mathcal{F}(\mathbb{C}^{\Modes\sqcup\Env}) .
\end{equation}
The detectors measure, on the output side, the $K$ commuting binary projective
measurements $\{\Pi^{0}_{k},\id-\Pi^{0}_{k}\}$ where
$\Pi^0_k\coloneqq\Fock(\id-\Emd{\{k\}})$, with the projector $\Emd{\{k\}}$ of
$\mathbb{C}^{\Modes}$ extended by zero to $\mathbb{C}^{\Modes\sqcup\Env}$, is
the projector onto ``no photon in the modes of $\Port{k}$''
(Lemma~\ref{lem:fock-functor}\,(iv), Assumption~\ref{ass:port}). For $J\subseteq[K]$ put
$\Pi^0_J\coloneqq\prod_{k\in J}\Pi^0_k$.
\end{definition}
\claimmeta{\metalabel{Uses}\mref{Def.}{def:fock}, \mref{Lem.}{lem:fock-functor}, \mref{Lem.}{lem:fock-maps}, \mref{Def.}{def:instrument}, \mref{Lem.}{lem:dilation}
\metasep \metalabel{Used by}\mref{Lem.}{lem:vacelim}, \mref{Thm.}{thm:click-povm}, \mref{Rem.}{rem:positivity}, \mref{Prop.}{prop:noise-floor}, \mref{Prop.}{prop:efficiency-normalization}
\metasep newly derived
\metasep accepted}

% claim receiver/lem:vacelim
\claimsummary{receiver/lem:vacelim}{Compressing the lift of any operator by vacuum on the modes outside a coordinate subspace gives exactly the lift of its block on that subspace, and the measurement isometry is an isometry. Dropping the undeclared input columns is therefore exact under the vacuum assumption.}
\claimbookmark{receiver/lem:vacelim}{Lemma}{lem:vacelim}{Exactness of vacuum elimination}
\claimleadin{%
Assumption~\ref{ass:vac} is what licenses restricting $\Gmat$ to the columns
$\Inp$; the lemma shows that the restriction is exact.
}
\begin{lemma}[Exactness of vacuum elimination]\label{lem:vacelim}
Let $S\subseteq T$ be finite sets of modes and
$\iota:\mathbb{C}^{S}\hookrightarrow\mathbb{C}^{T}$ the coordinate inclusion.
Under the factorization
$\mathcal{F}(\mathbb{C}^{T})\cong\mathcal{F}(\mathbb{C}^{S})\otimes
\mathcal{F}(\mathbb{C}^{T\setminus S})$ of occupation bases, the occupation
identification of Lemma~\ref{lem:fock-maps}\,(ii), $\Fock(\iota)$ is
$\ket{\psi}\mapsto\ket{\psi}\otimes\ket{\vac}$; and for any operator $X$ on
$\mathbb{C}^{T}$,
\begin{equation}\label{eq:vacelim}
  \Fock(\iota)^\dagger\,\Fock(X)\,\Fock(\iota)
  \;=\;\Fock\big(\iota^\dagger X\iota\big),
\end{equation}
where $\iota^\dagger X\iota$ is the $S\times S$ block of $X$. In
particular $\Fock(\iota)$ is an isometry. The inclusion $\iota$ of
Definition~\ref{def:physmodel}, with $S=\Inp$ and $T=\Modes\sqcup\Env$, is
one instance, and for it $V$ of \eqref{eq:measisometry} is an isometry as
well.
\end{lemma}
\claimmeta{\metalabel{Uses}\mref{Def.}{def:fock}, \mref{Lem.}{lem:fock-functor}, \mref{Lem.}{lem:fock-maps}, \mref{Def.}{def:physmodel}
\metasep \metalabel{Assumes}\aref{A-VAC}{ass:vac}, \aref{A-PORT}{ass:port}
\metasep \metalabel{Used by}\mref{Rem.}{rem:lossvsprobe}, \mref{Thm.}{thm:port-grams}, \mref{Thm.}{thm:click-povm}, \mref{Ex.}{ex:bb84}, \mref{Ex.}{ex:bb84bound}
\metasep newly derived
\metasep accepted}

\begin{proof}
$\iota$ maps each basis vector $\ket{m}$, $m\in S$, to the basis vector
$\ket{m}$ of $\mathbb{C}^{T}$, so $\iota^{\otimes n}$ maps the symmetrized
and normalized $\ket{e_{m_1}}\otimes\cdots\otimes\ket{e_{m_n}}$ of
Definition~\ref{def:fock} to the same symmetrized and normalized vector over
$T$: the occupation vector $\ket{\vec m}$ of $S$ goes to the occupation
vector of $T$ that agrees with $\vec m$ on $S$ and is $0$ on $T\setminus S$,
which the factorization identifies with $\ket{\vec m}\otimes\ket{\vac}$, and
by linearity $\Fock(\iota)\ket{\psi}=\ket{\psi}\otimes\ket{\vac}$ for every
$\ket{\psi}$. Next, $\Fock(\iota)$ maps
$\Symn{n}(\mathbb{C}^{S})$ into $\Symn{n}(\mathbb{C}^{T})$ by
$\iota^{\otimes n}$, and distinct
$n$ land in orthogonal summands, so both sides of \eqref{eq:vacelim} are block
diagonal in $n$ and it suffices to compare blocks. For
$\ket{\psi},\ket{\phi}\in\Symn{n}(\mathbb{C}^{S})$,
\begin{equation}
  \bra{\psi}\Fock(\iota)^\dagger\Fock(X)\Fock(\iota)\ket{\phi}
  =\bra{\iota^{\otimes n}\psi}X^{\otimes n}\ket{\iota^{\otimes n}\phi}
  =\bra{\psi}(\iota^\dagger X\iota)^{\otimes n}\ket{\phi},
\end{equation}
using $(\iota^\dagger)^{\otimes n}X^{\otimes n}\iota^{\otimes n}
=(\iota^\dagger X\iota)^{\otimes n}$, which holds factor by factor. Taking
$X=\id$ gives $\Fock(\iota)^\dagger\Fock(\iota)=\Fock(\id_{S})=\id$, so
$\Fock(\iota)$ is an isometry; for the inclusion of
Definition~\ref{def:physmodel}, $\Fock(U)$ is unitary by
Lemma~\ref{lem:fock-functor}\,(iii), so $V$ is an isometry.
\end{proof}

% claim receiver/rem:lossvsprobe
\claimsummary{receiver/rem:lossvsprobe}{Vacuum elimination removes an input mode only because it carries vacuum; a declared source mode carries the probe state and stays a column of the instrument matrix. Eliminating it would replace the probe by vacuum and describe a different device, a security-relevant mismatch.}
\claimbookmark{receiver/rem:lossvsprobe}{Remark}{rem:lossvsprobe}{Loss ancillas and probe inputs}
\begin{remark}[Loss ancillas and probe inputs]
\label{rem:lossvsprobe}
Lemma~\ref{lem:vacelim} eliminates an input mode precisely because that mode
carries \emph{vacuum}. A declared source mode does not: it carries
$\rho^{(t)}$, it stays as a column of $\Gmat$, and its state is supplied
(Definition~\ref{def:probe}). The difference between ``a loss ancilla'' and ``a
probe input'' is exactly the difference between an undeclared vacuum input and
a declared-state input, and the two must not be conflated: eliminating a
probe mode would silently replace the probe by vacuum, so the derived POVM
would describe a different device than the one deployed; since the probe is
a countermeasure, that mismatch is security-relevant.
\end{remark}
\claimmeta{\metalabel{Uses}\mref{Def.}{def:probe}, \mref{Lem.}{lem:vacelim}
\metasep newly derived
\metasep accepted}

\section{Component graphs and compilation}\label{group-receiver/def:catalogue}
\sectionintro{A receiver may be declared as a directed graph of catalogue components. The graph determines an \termref{term-instrument-matrix}{instrument matrix} when it is \termref{term-well-formed}{well formed} (parameters in their domains, acyclic, and any two elements that share a mode ordered): every topological order then compiles to the same contraction (\hyperlink{claim-receiver/thm:compilation-well-defined}{Theorem~\ref*{thm:compilation-well-defined}}). A well-formed graph, with its \termref{term-declared-input-modes}{declared input modes}, ports and efficiencies, then gives a \termref{term-receiver-instrument}{receiver instrument} (\hyperlink{claim-receiver/cor:compiled-instrument}{Corollary~\ref*{cor:compiled-instrument}}). Without the ordering condition two sorts of the same graph can give different devices, which is why it is a well-formedness condition rather than a modelling remark.}

% claim receiver/def:catalogue
\claimsummary{receiver/def:catalogue}{A component catalogue is a finite set of entry types, each fixing a support shape, a parameter domain and a rule that yields a local element. The catalogue submodule declares the entries and proves that each maps its whole domain into local elements.}
\claimbookmark{receiver/def:catalogue}{Definition}{def:catalogue}{Component catalogue}
\claimleadin{%
Route (iii) declares a receiver as a graph of standard components rather than
as a matrix. This section shows when such a graph determines the
\termref{term-instrument-matrix}{instrument matrix}: it does, under a condition that is stated rather
than assumed.
}
\begin{definition}[Component catalogue]\label{def:catalogue}
A \termdef{term-component-catalogue}{component catalogue} $\Cat$ is a finite set of entry types. An
\termdef{term-entry-type}{entry type} fixes a support shape (which modes of $\Modes$ it acts on), a
parameter domain, and a rule assigning to every parameter tuple in that domain
a \termref{term-local-element}{local element} in the sense of Definition~\ref{def:embed}. The
catalogue of this contract is the one of Definition~\ref{def:catalogue-entries}
(Table~\ref{tab:catalogue}) in the component-catalogue submodule, and
Theorem~\ref{thm:catalogue-admissibility} there states that every entry type maps its whole
declared domain into local elements.
\end{definition}
\claimmeta{\metalabel{Uses}\mref{Def.}{def:embed}
\metasep \metalabel{Used by}\mref{Def.}{def:graph}
\metasep newly derived
\metasep accepted}

% claim receiver/def:graph
\claimsummary{receiver/def:graph}{A declared element graph is well formed when every element's parameters lie in their domain, the graph is acyclic, and any two elements sharing a mode are ordered by it. Support comparability keeps an arbitrary sort from supplying a physical ordering the declaration left open.}
\claimbookmark{receiver/def:graph}{Definition}{def:graph}{Declared element graph}
\begin{definition}[Declared element graph]\label{def:graph}
A \termdef{term-declared-element-graph}{declared element graph} is a finite directed graph $\Grph=(V,E)$ whose
vertices $e\in V$ are \termdef{term-elements}{elements}, each carrying a catalogue entry type, a
parameter tuple, and a support $\Sigma_e\subseteq\Modes$. Write $\preceq_\Grph$
for the reachability relation: $e\preceq_\Grph f$ when $e=f$ or there is a
directed path from $e$ to $f$. The graph is \termdef{term-well-formed}{well formed} when
\begin{enumerate}[label=(E\arabic*),leftmargin=3.0em]
\item every element's parameter tuple lies in the declared domain of its entry
  type, so that, by Definition~\ref{def:catalogue}, $\Gmat_{\Sigma_e}$ is a
  \termref{term-local-element}{local element}: defined, with
  $\Gmat_{\Sigma_e}^\dagger\Gmat_{\Sigma_e}\preceq\id$;
\item $\Grph$ is acyclic, so that $\preceq_\Grph$ is a partial order;
\item (\termdef{term-support-comparability}{support comparability}) any two elements whose supports intersect
  are comparable: $\Sigma_e\cap\Sigma_f\neq\varnothing\implies
  e\preceq_\Grph f$ or $f\preceq_\Grph e$.
\end{enumerate}
A \termdef{term-linear-extension}{linear extension} of $\Grph$ is an enumeration
$e_1,\dots,e_L$ of $V$ with $e_i\preceq_\Grph e_j\implies i\le j$.
\end{definition}
\claimmeta{\metalabel{Uses}\mref{Def.}{def:catalogue}
\metasep \metalabel{Used by}\mref{Def.}{def:compile}, \mref{Thm.}{thm:compilation-well-defined}, \mref{Cor.}{cor:compiled-instrument}, \mref{Rem.}{rem:noloops}, \mref{Def.}{def:tolmap}, \mref{Req.}{req:graph-refusal}, \mref{Thm.}{thm:catalogue-admissibility}, \mref{Lem.}{lem:catalogue-dilations}, \mref{Rem.}{rem:out}
\metasep newly derived
\metasep implemented (effective: accepted)
\metasep \metalabel{Code}\coderef{https://github.com/fraunhofersg/certiqs-key/blob/v0.1.0/src/certiqs_key/receiver/graph.py}{certiqs_key.receiver.graph:ElementGraph}}

% claim receiver/def:compile
\claimsummary{receiver/def:compile}{The compiled instrument matrix is the product of the embedded elements along a linear extension, later elements multiplying from the left, followed by the terminal detector-efficiency layer. This fixes the matrix convention of the whole receiver contract.}
\claimbookmark{receiver/def:compile}{Definition}{def:compile}{Compilation of a declared element graph}
\claimleadin{%
The product below follows the matrix convention the interface records
(\href{../../../spec/modules/receiver/interface.yaml}{\nolinkurl{spec/modules/receiver/interface.yaml}}): a factor acts on column
vectors indexed by the modes, so the element that light meets later stands
further to the left.
}
\begin{definition}[Compilation of a declared element graph]\label{def:compile}
Let $\Grph$ be a well-formed graph; it is in particular acyclic
(condition (E2) of Definition~\ref{def:graph}), so that a linear extension exists. For a
linear extension $\sigma=(e_1,\dots,e_L)$ of $\Grph$ put
\begin{equation}\label{eq:compile}
  \Gmat_{\Grph}^{\sigma}\coloneqq
  \hat\Gmat_{e_L}\,\hat\Gmat_{e_{L-1}}\cdots\hat\Gmat_{e_1}
  \ \in\ \mathbb{C}^{\Modes\times\Modes},
\end{equation}
later elements multiplying from the left. The
\termdef{term-compiled-instrument-matrix}{compiled instrument matrix} of the receiver is
$\Gmat\coloneqq D_\eta\,\Gmat^{\sigma}_{\Grph}$, where
$D_\eta=\sum_k\sqrt{\eta_k}\,\Emd{\{k\}}+\Emd{\Modes\setminus\bigcup_k
\Port{k}}$ is the \termdef{term-terminal-detector-efficiency-layer}{terminal detector-efficiency layer}, applied after every
element, and $\eta_k\in[0,1]$ is the declared efficiency of port $k$ (the
effect of that layer on the port Grams is Theorem~\ref{thm:efficiency} of the
detector-imperfections submodule). The \termref{term-detector-ports}{detector ports} $\Port{k}$ and the
port-embedding abbreviation $\Emd{\{k\}}$ written here are those of a
\termref{term-receiver-instrument}{receiver instrument} tuple: clause (iii) of
Definition~\ref{def:instrument} and its closing notation.
\end{definition}
\claimmeta{\metalabel{Uses}\mref{Def.}{def:instrument}, \mref{Def.}{def:graph}
\metasep \metalabel{Used by}\mref{Thm.}{thm:compilation-well-defined}, \mref{Cor.}{cor:compiled-instrument}, \mref{Thm.}{thm:enclosed-blocks}, \mref{Req.}{req:exact-route-decision}, \mref{Thm.}{thm:catalogue-admissibility}, \mref{Lem.}{lem:catalogue-dilations}, \mref{Rem.}{rem:layers}, \mref{Ex.}{ex:bb84}, \mref{Ex.}{ex:bb84bound}, \mref{Ex.}{ex:bb84exact}, \mref{Thm.}{thm:efficiency}, \mref{Def.}{def:normalized-station}
\metasep newly derived
\metasep implemented (effective: accepted)
\metasep \metalabel{Code}\coderef{https://github.com/fraunhofersg/certiqs-key/blob/v0.1.0/src/certiqs_key/receiver/graph.py}{certiqs_key.receiver.graph:compile_element_graph}}

% claim receiver/thm:compilation-well-defined
\claimsummary{receiver/thm:compilation-well-defined}{Every linear extension of a well-formed graph compiles to the same contraction before the efficiency layer, which is unitary exactly when every element is lossless, and well-formedness is decidable by finite inspection. A well-formed route-(iii) graph therefore determines its instrument matrix.}
\claimbookmark{receiver/thm:compilation-well-defined}{Theorem}{thm:compilation-well-defined}{Well-definedness of compilation}
\claimleadin{%
Two elements that share a mode but are not ordered by the graph can make the
compiled product depend on the sort (Remark~\ref{rem:e3}). Support
comparability rules this out,
and item (i) shows that it is enough for every sort to give the same matrix.
}
\begin{theorem}[Well-definedness of compilation]\label{thm:compilation-well-defined}
In items (i)--(iii), let $\Grph$ be a well-formed declared element graph
(Definition~\ref{def:graph}); in particular $\Grph$ is acyclic
(condition (E2) of Definition~\ref{def:graph}), so that $\preceq_\Grph$ is a partial order and
linear extensions of $\Grph$ and of each of its induced subgraphs exist.
Then:
\begin{enumerate}[label=(\roman*),leftmargin=3.0em]
\item \emph{Independence of the sort.} $\Gmat^{\sigma}_{\Grph}$ is the same
  matrix for every linear extension $\sigma$ of $\Grph$; write
  $\Gmat_{\Grph}$ for it. In particular, if $V$ is partitioned into
  \emph{layers} $L_1,\dots,L_r$ such that elements of one layer have pairwise
  disjoint supports and $e\preceq_\Grph f$ implies
  $\operatorname{layer}(e)\le\operatorname{layer}(f)$, then
  $\Gmat_\Grph=\big(\prod_{e\in L_r}\hat\Gmat_e\big)\cdots
  \big(\prod_{e\in L_1}\hat\Gmat_e\big)$ with each intra-layer product
  order-independent.
\item \emph{Admissibility by construction.}
  $\Gmat_\Grph^\dagger\Gmat_\Grph\preceq\id$, and hence also
  $\Gmat^\dagger\Gmat\preceq\id$ for the compiled instrument matrix of
  Definition~\ref{def:compile}.
\item \emph{Losslessness.} $\Gmat_\Grph$ is unitary if and only if every
  element of $\Grph$ is lossless on its support. If in addition every
  $\eta_k=1$ then $\Gmat$ is unitary.
\item \emph{Decidability.} For an arbitrary declared element graph, not
  necessarily well formed, conditions (E1)--(E3) are decidable by finite
  inspection: (E1) by evaluating the declared domain predicate of each
  element, (E2) by a topological sort, (E3) by testing
  $\Sigma_e\cap\Sigma_f\neq\varnothing$ against the transitive closure of $E$
  for each of the $\binom{|V|}{2}$ pairs.
\end{enumerate}
\end{theorem}
\claimmeta{\metalabel{Uses}\mref{Prop.}{prop:mode-algebra}, \mref{Def.}{def:graph}, \mref{Def.}{def:compile}
\metasep \metalabel{Used by}\mref{Cor.}{cor:compiled-instrument}, \mref{Req.}{req:graph-refusal}, \mref{Req.}{req:exact-route-decision}, \mref{Thm.}{thm:catalogue-admissibility}
\metasep newly derived
\metasep implemented (effective: accepted)
\metasep \metalabel{Code}\coderef{https://github.com/fraunhofersg/certiqs-key/blob/v0.1.0/src/certiqs_key/receiver/graph.py}{certiqs_key.receiver.graph:compile_element_graph}
\metasep \metalabel{Oracles}\coderef{https://github.com/fraunhofersg/certiqs-key/blob/v0.1.0/tests/oracles/registry_receiver_povm.py}{registry_receiver_povm.py}}

\begin{proof}
(i) Induction on $L=|V|$. For $L\le1$ there is nothing to prove. Let $\sigma$
and $\tau$ be linear extensions and let $e\coloneqq\sigma(1)$. Since $\sigma$
is a linear extension and $e$ is first, no element other than $e$ satisfies
$f\preceq_\Grph e$; that is, $e$ is minimal for $\preceq_\Grph$. Let $j$ be the
position of $e$ in $\tau$. Every element $f=\tau(i)$ with $i<j$ satisfies
$e\not\preceq_\Grph f$ (else $f$ would have to follow $e$ in $\tau$) and
$f\not\preceq_\Grph e$ (minimality of $e$), so $f$ and $e$ are incomparable.
By (E3), $\Sigma_f\cap\Sigma_e=\varnothing$, and by
Proposition~\ref{prop:mode-algebra}\,(iii) the embeddings $\hat\Gmat_f$ and
$\hat\Gmat_e$ commute. In the product
$\Gmat^{\tau}_{\Grph}=\hat\Gmat_{\tau(L)}\cdots\hat\Gmat_{\tau(1)}$ the factor
$\hat\Gmat_e$ sits at position $j$ from the right and may therefore be moved
past the $j-1$ factors to its right one at a time, giving
$\Gmat^{\tau}_{\Grph}=\big(\hat\Gmat_{\tau(L)}\cdots\widehat{\hat\Gmat_{e}}
\cdots\hat\Gmat_{\tau(1)}\big)\hat\Gmat_e$, where the hat denotes omission and
the remaining order is $\tau$ restricted to $V\setminus\{e\}$. The induced
subgraph on $V\setminus\{e\}$ is itself well formed: (E1) and (E2) pass to any
induced subgraph, and since $e$ is minimal no directed path between two other
elements passes through $e$, so reachability among them, and with it (E3), is
inherited. That restriction of $\tau$ is therefore a linear extension of the
induced subgraph; so is $\sigma$ restricted to $V\setminus\{e\}$, and both
compile to the same matrix by the induction hypothesis. Multiplying on the
right by $\hat\Gmat_e$ gives $\Gmat^\tau_\Grph=\Gmat^\sigma_\Grph$.

For the layered form, note that a layering as described yields a linear
extension by listing $L_1$ in an order compatible with $\preceq_\Grph$ (one
exists, the restriction of $\preceq_\Grph$ to a layer being a partial order),
then $L_2$ likewise, and so on: an element preceding another in
$\preceq_\Grph$ lies in an earlier layer or earlier within the same one. The
intra-layer products are order-independent by
Proposition~\ref{prop:mode-algebra}\,(iii), since the elements of one layer
have pairwise disjoint supports, so the layered product equals the product
along that linear extension, which is $\Gmat_\Grph$ by the first part.

(ii) Each $\hat\Gmat_e$ is a contraction by (E1) and
Proposition~\ref{prop:mode-algebra}\,(i), so the product is a contraction by
repeated application of Proposition~\ref{prop:mode-algebra}\,(ii). $D_\eta$ is
diagonal with entries in $[0,1]$, hence a contraction, and one more application
of (ii) gives $\Gmat^\dagger\Gmat\preceq\id$.

(iii) If every element is lossless then every $\hat\Gmat_e$ is unitary
(Proposition~\ref{prop:mode-algebra}\,(i) and (iv), a square isometry being
unitary), and so is their product. Conversely suppose $\Gmat_\Grph$ is unitary,
and write $X_i\coloneqq\hat\Gmat_{e_i}\cdots\hat\Gmat_{e_1}$. For every vector
$v$,
$\lVert v\rVert=\lVert X_Lv\rVert\le\lVert X_{L-1}v\rVert\le\cdots\le
\lVert X_1v\rVert\le\lVert v\rVert$, so every inequality is an equality for
every $v$. From $\lVert X_1v\rVert=\lVert v\rVert$ for all $v$, $\hat\Gmat_{e_1}$
is a square isometry, hence unitary, hence $\operatorname{ran}X_1=
\mathbb{C}^{\Modes}$; then $\lVert X_2v\rVert=\lVert X_1v\rVert$ for all $v$
says $\hat\Gmat_{e_2}$ is isometric on all of $\mathbb{C}^{\Modes}$, hence
unitary, and so on. Each $\hat\Gmat_{e_i}$ unitary forces
$\Gmat_{\Sigma_{e_i}}$ unitary, i.e.\ lossless. The statement about $D_\eta$ is
the same argument applied to the one extra factor, $D_\eta$ being unitary
exactly when every $\eta_k=1$.

(iv) Each test is a finite computation on finite data, as described.
\end{proof}

% claim receiver/cor:compiled-instrument
\claimsummary{receiver/cor:compiled-instrument}{A well-formed graph with a declared mode space, input split, port map and efficiencies yields a receiver instrument whose matrix does not depend on the sort. This joins the graph route to the object the measurement theorems are stated of.}
\claimbookmark{receiver/cor:compiled-instrument}{Corollary}{cor:compiled-instrument}{Receiver instrument from a well-formed graph declaration}
\claimleadin{%
The three input routes converge on a receiver instrument. This corollary is
the step for route (iii). Once it holds, a result this contract proves for
every receiver instrument applies to a route-(iii) declaration by
instantiation; a result proved of a receiver instrument carrying further
hypotheses applies only where the declaration meets those too.
}
\begin{corollary}[Receiver instrument from a well-formed graph declaration]
\label{cor:compiled-instrument}
Let a route-(iii) declaration supply a \termref{term-declared-mode-space}{declared mode space} $\Modes$; a
well-formed declared element graph $\Grph$ on $\Modes$
(Definition~\ref{def:graph}), in particular acyclic
(condition (E2) of Definition~\ref{def:graph}); a set
$\Inp\subseteq\Modes$ of \termref{term-declared-input-modes}{declared input modes} with a split
$\Inp=\InpC\sqcup\InpS$, $\InpC\neq\varnothing$, every mode of $\Modes\setminus\Inp$ carrying vacuum
at the input (Assumption~\ref{ass:vac}); and a port map
$\Port{1},\dots,\Port{K}$, pairwise disjoint subsets of $\Modes$ with
$K\ge1$, carrying detector efficiencies $\eta_k$ in the domain $[0,1]$
declared for them in Definition~\ref{def:compile}. Let $\sigma$ be a linear
extension of $\Grph$ and let
$\Gmat\coloneqq D_\eta\,\Gmat^{\sigma}_{\Grph}$ be the compiled instrument
matrix of Definition~\ref{def:compile}. Then
\begin{enumerate}[label=(\roman*),leftmargin=3.0em]
\item $\Gmat$ does not depend on $\sigma$, so it is a function of the
  declaration alone; and
\item the tuple
  $\Recv=\big(\Modes;\ \Inp=\InpC\sqcup\InpS;\ \{\Port{k}\}_{k=1}^{K};\
  \Gmat\big)$ is a receiver instrument in the sense of
  Definition~\ref{def:instrument}: it has the shape \eqref{eq:instrument} and
  satisfies its four clauses.
\end{enumerate}
\end{corollary}
\claimmeta{\metalabel{Uses}\mref{Def.}{def:instrument}, \mref{Def.}{def:graph}, \mref{Def.}{def:compile}, \mref{Thm.}{thm:compilation-well-defined}
\metasep \metalabel{Assumes}\aref{A-VAC}{ass:vac}
\metasep \metalabel{Used by}\mref{Ex.}{ex:bb84}
\metasep newly derived
\metasep accepted}

\begin{proof}
\emph{Clauses (i) to (iii) of Definition~\ref{def:instrument}.} Clause (i) is
the declared mode space and clause (iii) is the declared port map; that the
ports are pairwise disjoint and that $K\ge1$ are conditions on the
declaration, decided by finite inspection of it, and nothing here derives
them. Clause (ii) is the declared split $\Inp=\InpC\sqcup\InpS$ with
$\InpC\neq\varnothing$, again a condition on the declaration, together with
the condition that every mode of $\Modes\setminus\Inp$ carries vacuum at the
input, which is Assumption~\ref{ass:vac}: it is a physical condition on the
declaration that the tuple's data does not encode and nothing here derives,
so this corollary declares it in its own right, as
Definition~\ref{def:instrument}\,(ii) does. What remains is clause (iv) and
item (i), and both are read off Theorem~\ref{thm:compilation-well-defined}.

\emph{Item (i).} Because $\Grph$ is well formed it is acyclic (condition (E2)
of Definition~\ref{def:graph}), so $\preceq_\Grph$
is a partial order and a linear extension $\sigma$ of $\Grph$ exists at all.
Hence the product \eqref{eq:compile} is defined and
$\Gmat^\sigma_\Grph\in\mathbb{C}^{\Modes\times\Modes}$. By
Theorem~\ref{thm:compilation-well-defined}\,(i) it is the same matrix for every linear extension of
$\Grph$; write $\Gmat_\Grph$ for that common value. The terminal layer
$D_\eta$ is built from the port map and the efficiencies alone and contains no
element of $\Grph$ and no ordering, so $\Gmat=D_\eta\,\Gmat^\sigma_\Grph=
D_\eta\,\Gmat_\Grph$ is independent of $\sigma$ as well and depends on $\Grph$
and on the declared pairs $\{(\Port{k},\eta_k)\}_{k=1}^{K}$ alone.

\emph{Clause (iv).} The terminal layer $D_\eta=\sum_k\sqrt{\eta_k}\,\Emd{\{k\}}+
\Emd{\Modes\setminus\bigcup_k\Port{k}}$ of Definition~\ref{def:compile} is
diagonal, with entry $\sqrt{\eta_k}$ on the modes of $\Port{k}$ (well
defined as a diagonal because the ports are pairwise disjoint, so that no mode
is given two entries) and with entries in $[0,1]$ because
$\eta_k\in[0,1]$; it is therefore a contraction. The indexing of $\Gmat$ is
that of \eqref{eq:compile}: each factor $\hat\Gmat_e$ acts on column vectors
indexed by $\Modes$, so the product, and with it
$\Gmat\in\mathbb{C}^{\Modes\times\Modes}$, is indexed rows by outputs and
columns by inputs, which is what clause (iv) of
Definition~\ref{def:instrument} requires of an \termref{term-instrument-matrix}{instrument matrix}.
Finally Theorem~\ref{thm:compilation-well-defined}\,(ii) states $\Gmat^\dagger\Gmat\preceq\id$ of
exactly this compiled instrument matrix. So clause (iv) holds and $\Recv$ is a
receiver instrument.
\end{proof}

% claim receiver/rem:e3
\claimsummary{receiver/rem:e3}{An acyclic graph of a rotation and a polarizer on the same two modes, with no edge between them, compiles at generic angles to two different contractions under its two sorts. Acyclicity alone therefore does not make a declared graph determine an instrument.}
\claimbookmark{receiver/rem:e3}{Remark}{rem:e3}{Order dependence without support comparability}
\begin{remark}[Order dependence without support comparability]
\label{rem:e3}
Take $\Modes=\{m\}$ and two elements $e,f$ on $\Sigma_e=\Sigma_f=\{m\}$ with no
edge between them, $\Gmat_{\Sigma_e}=(\tfrac12)$ and
$\Gmat_{\Sigma_f}=(\tfrac12)$: here the two orders agree, and the example
looks harmless. Now take $\Modes=\{m_1,m_2\}$ with
$\Sigma_e=\Sigma_f=\Modes$, $\Gmat_{\Sigma_e}=\Rot(\theta)$ a rotation and
$\Gmat_{\Sigma_f}=\operatorname{diag}(1,0)$ a polarizer. The graph
$\Grph=(\{e,f\},\varnothing)$ is acyclic and every element is admissible, yet
$\Gmat_{\Sigma_f}\Gmat_{\Sigma_e}\neq\Gmat_{\Sigma_e}\Gmat_{\Sigma_f}$ for
$\theta\notin\frac{\pi}{2}\mathbb{Z}$: rotating and then projecting is not
projecting and then rotating. Two topological sorts of the \emph{same} declared
receiver would compile to two different instruments, and, because both are
contractions with the same declared port map, both would pass every other
validation in this contract. Without (E3), ``the declared receiver'' does not
denote.
\end{remark}
\claimmeta{newly derived
\metasep accepted}

% claim receiver/rem:edgemeaning
\claimsummary{receiver/rem:edgemeaning}{An edge fixes the order in which two elements act; moving light between modes is done by the elements' matrices. Components drawn on different arms carry no edge and have disjoint supports, so support comparability constrains only an ambiguous drawing.}
\claimbookmark{receiver/rem:edgemeaning}{Remark}{rem:edgemeaning}{Meaning of an edge}
\begin{remark}[Meaning of an edge]\label{rem:edgemeaning}
An edge $e\to f$ fixes the order in which $e$ and $f$ act: $e$ first. Moving
light from one mode to another is done by the elements' matrices, for
instance by a mode-permutation element (entry C9 of
Table~\ref{tab:catalogue}), not by the edge; between elements with disjoint
supports the order the edge fixes does not by itself change the compiled
matrix, since they commute (Proposition~\ref{prop:mode-algebra}\,(iii)); such
an edge matters only through the order it induces, by reachability, between
elements whose supports intersect. A user who
draws two components in series draws the edge; a user who draws them on
different arms of a beamsplitter draws none, and (E3) then requires their mode
supports to be disjoint, which, on different arms, they are. The condition
therefore constrains exactly the drawings that were ambiguous.
\end{remark}
\claimmeta{newly derived
\metasep accepted}

% claim receiver/rem:noloops
\claimsummary{receiver/rem:noloops}{A cavity, a ring resonator or any other feedback path is a cycle, which (E2) excludes. Compiling one would need a fixed-point solve whose contractivity must itself be certified, so such a receiver needs a new claim rather than a larger computation.}
\claimbookmark{receiver/rem:noloops}{Remark}{rem:noloops}{Feedback paths outside the acyclic scope}
\begin{remark}[Feedback paths outside the acyclic scope]\label{rem:noloops}
A cycle in $\Grph$, such as a cavity, a ring resonator or any feedback path,
is excluded by condition (E2) of Definition~\ref{def:graph}. Compiling one would require solving a fixed-point
equation on the mode graph, and the resulting matrix is a contraction only
under a spectral condition that would itself have to be certified. This
contract does neither: a declared cycle leaves no finite product to compile,
and it is refused, not solved (Requirement~\ref{req:graph-refusal}). A
receiver that needs a feedback path is a new claim, not a bigger computation.
\end{remark}
\claimmeta{\metalabel{Uses}\mref{Def.}{def:graph}
\metasep newly derived
\metasep accepted}

% claim receiver/rem:acyc-efficiency
\claimsummary{receiver/rem:acyc-efficiency}{The claims that read only the Gram substitution of the efficiency theorem involve the terminal efficiency layer alone, a function of the port map and efficiencies, and no element graph. The theorem's items on the compiled product and its record hold only for a well-formed, hence acyclic, graph.}
\claimbookmark{receiver/rem:acyc-efficiency}{Remark}{rem:acyc-efficiency}{Where the element graph enters the efficiency theorem}
\begin{remark}[Where the element graph enters the efficiency theorem]
\label{rem:acyc-efficiency}
The two remarks that use Theorem~\ref{thm:efficiency}
(Remarks~\ref{rem:powervsclick} and~\ref{rem:sperlingeta}) and
Requirement~\ref{req:efficiency-convention}, which rests on the first of them,
read only the substitution of item (i) of that theorem: they involve the
\termref{term-terminal-detector-efficiency-layer}{terminal detector-efficiency layer} $D_\eta$ of
Definition~\ref{def:compile}, which is a function of the port map and the
per-port efficiencies alone, and no graph enters their statements or proofs.
The same holds of items (i), (ii) and (v) of the theorem. Items (iii) and (iv)
are statements about the compiled product of Definition~\ref{def:compile} and
its record, which exist only for a well-formed, hence acyclic, graph
(condition (E2) of Definition~\ref{def:graph}); so are
Requirements~\ref{req:efficiency-schema} and~\ref{req:efficiency-record},
which state their engine side.
\end{remark}
\claimmeta{newly derived
\metasep accepted}

\section{Port Grams and the click-pattern POVM}\label{group-receiver/thm:port-grams}
\sectionintro{The joint no-click element of any set of ports pulls back to the lift of the identity minus the sum of the Grams of the ports in that set, whatever dilation is chosen (\hyperlink{claim-receiver/thm:port-grams}{Theorem~\ref*{thm:port-grams}}), and the \termref{term-loss-deficiency}{loss deficiency} splits exactly into loss through the dilation and light reaching unmonitored declared modes. An outcome is the set of ports that click. The main result, \hyperlink{claim-receiver/thm:click-povm}{Theorem~\ref*{thm:click-povm}}, is the click-pattern POVM: it follows by inclusion and exclusion, is positive by construction rather than by its alternating form, and is block diagonal in the total input photon number.}

% claim receiver/thm:port-grams
\claimsummary{receiver/thm:port-grams}{The port Grams form an admissible family, and every joint no-click element pulls back to the lift of the identity minus the sum of the Grams of the ports in its set. The measurement depends only on the declared input columns and on no choice of dilation.}
\claimbookmark{receiver/thm:port-grams}{Theorem}{thm:port-grams}{Port Grams and dilation independence}
\claimleadin{%
The port Grams are the compact data the whole measurement is derived from.
This theorem shows that they are admissible, that the \termref{term-joint-no-click-elements}{joint no-click
elements} are functions of them alone, and that the unitary dilation used to
define the physical model leaves no trace in the result.
}
\begin{theorem}[Port Grams and dilation independence]\label{thm:port-grams}
Let $\Recv$ be a \termref{term-receiver-instrument}{receiver instrument} with
\termref{term-port-gram-matrices}{port Grams} $\{A_k\}$ of
Definition~\ref{def:gram}.
\begin{enumerate}[label=(\roman*),leftmargin=3.0em]
\item \emph{Admissibility.} $A_k\succeq0$ for every $k$;
  $\sum_{k}A_k=\Gmat_{\Inp}^\dagger\Emd{[K]}\Gmat_{\Inp}\preceq
  \Gmat_{\Inp}^\dagger\Gmat_{\Inp}\preceq\id_{\Inp}$, so
  $\{A_k\}$ is admissible and the \termref{term-loss-deficiency}{loss deficiency}
  $\defic=\id_{\Inp}-A_\Sigma$ is positive semidefinite.
\item \emph{Rank.} $\rk A_k\le\min\big(|\Inp|,|\Port{k}|\big)$.
\item \emph{Pullback.} For every $J\subseteq[K]$,
  \begin{equation}\label{eq:pullback}
    V^\dagger\,\Pi^0_J\,V=\Fock\big(\id_{\Inp}-A_J\big).
  \end{equation}
\item \emph{Vacuum elimination is exact.} The right-hand side of
  \eqref{eq:pullback} depends on the instrument only through
  $\Gmat_{\Inp}$, the columns of $\Gmat$ indexed by the
  \termref{term-declared-input-modes}{declared input modes}. Dropping the other columns is therefore an identity, not an
  approximation.
\item \emph{Dilation independence.} \eqref{eq:pullback} holds for
  \emph{every} unitary $U$ on $\mathbb{C}^{\Modes\sqcup\Env'}$, with any
  finite mode set $\Env'$, whose block from the inputs $\Inp$ to the outputs
  $\bigcup_k\Port{k}$ equals the corresponding block of $\Gmat$. In particular
  the sub-unitary description (a contraction $\Gmat$ on the declared modes) and
  the unitary-dilation description (a unitary on system plus environment, with
  vacuum environment and the environment discarded) induce the \emph{same}
  measurement on $\mathcal{F}(\mathbb{C}^{\Inp})$.
\end{enumerate}
\end{theorem}
\claimmeta{\metalabel{Uses}\mref{Lem.}{lem:fock-functor}, \mref{Def.}{def:gram}, \mref{Lem.}{lem:dilation}, \mref{Lem.}{lem:vacelim}
\metasep \metalabel{Used by}\mref{Thm.}{thm:click-povm}, \mref{Thm.}{thm:representation-independence}, \mref{Cor.}{cor:unconditioned-family}, \mref{Rem.}{rem:bound-data-band}, \mref{Req.}{req:exact-route-decision}, \mref{Ex.}{ex:bb84}, \mref{Prop.}{prop:noise-floor}, \mref{Prop.}{prop:efficiency-normalization}
\metasep newly derived
\metasep implemented (effective: accepted)
\metasep \metalabel{Code}\coderef{https://github.com/fraunhofersg/certiqs-key/blob/v0.1.0/src/certiqs_key/receiver/grams.py}{certiqs_key.receiver.grams:check_gram_family_admissible}
\metasep \metalabel{Oracles}\coderef{https://github.com/fraunhofersg/certiqs-key/blob/v0.1.0/tests/oracles/registry_receiver_povm.py}{registry_receiver_povm.py}}

\begin{proof}
(i) $A_k=(\Emd{\{k\}}\Gmat_{\Inp})^\dagger(\Emd{\{k\}}\Gmat_{\Inp})\succeq0$
because $\Emd{\{k\}}$ is an orthogonal projector, hence idempotent and
self-adjoint. The ports are pairwise disjoint, so
$\sum_k\Emd{\{k\}}=\Emd{[K]}\preceq\id_{\Modes}$, and conjugating that
inequality by $\Gmat_{\Inp}$ gives
$\sum_kA_k\preceq\Gmat_{\Inp}^\dagger\Gmat_{\Inp}$, which is
$\preceq\id_{\Inp}$ because $\Gmat^\dagger\Gmat\preceq\id_{\Modes}$ restricts
to the principal submatrix on $\Inp$.

(ii) $A_k$ is the Gram matrix of the rows of $\Gmat_{\Inp}$ indexed by
$\Port{k}$, so $\rk A_k=\rk(\Emd{\{k\}}\Gmat_{\Inp})$, and that matrix has at
most $|\Port{k}|$ nonzero rows and $|\Inp|$ columns.

(iii) By Lemma~\ref{lem:fock-functor}\,(iv), $\Pi^0_k$ is the orthogonal projector
onto the \termref{term-fock-space}{Fock space} over the modes outside $\Port{k}$, and since the
$\Port{k}$ are pairwise disjoint the $\Pi^0_k$ are commuting projectors with
$\Pi^0_J=\Fock(\id-\Emd{J})$: indeed
$\prod_{k\in J}\Fock(\id-\Emd{\{k\}})=\Fock\big(\prod_{k\in
J}(\id-\Emd{\{k\}})\big)=\Fock(\id-\Emd{J})$ by
Lemma~\ref{lem:fock-functor}\,(ii) and disjointness. Then, using
Lemma~\ref{lem:fock-functor}\,(ii) and Lemma~\ref{lem:vacelim},
\begin{equation}
  V^\dagger\Pi^0_JV
  =\Fock(\iota)^\dagger\Fock\big(U^\dagger(\id-\Emd{J})U\big)\Fock(\iota)
  =\Fock\Big(\iota^\dagger U^\dagger(\id-\Emd{J})U\iota\Big).
\end{equation}
Now $\iota^\dagger U^\dagger(\id-\Emd{J})U\iota
=\id_{\Inp}-(\Emd{J}U\iota)^\dagger(\Emd{J}U\iota)$, and $\Emd{J}U\iota$ is the
submatrix of $U$ with rows in $\bigcup_{k\in J}\Port{k}\subseteq\Modes$ and
columns in $\Inp\subseteq\Modes$, which is $\Emd{J}\Gmat_{\Inp}$ because the
$\Modes\times\Modes$ block of $U$ is $\Gmat$. Hence
$\iota^\dagger U^\dagger(\id-\Emd{J})U\iota=\id_{\Inp}-\Gmat_{\Inp}^\dagger
\Emd{J}\Gmat_{\Inp}=\id_{\Inp}-A_J$.

(iv) Immediate from the closed form just obtained: $A_J$ is built from
$\Gmat_{\Inp}$ alone.

(v) The computation in (iii) used only two facts about $U$: that it is unitary,
and that $\Emd{J}U\iota=\Emd{J}\Gmat_{\Inp}$ for every $J$, i.e.\ that its
block from $\Inp$ to the port modes agrees with that of $\Gmat$. Any dilation
with that property therefore yields the same $V^\dagger\Pi^0_JV$ for every
$J$. It yields the same measurement too: the $\Pi^0_k$ commute, so expanding
$\prod_{k\in S}(\id-\Pi^0_k)=\sum_{T\subseteq S}(-1)^{|T|}\prod_{k\in T}\Pi^0_k$
gives
$V^\dagger\big[\prod_{k\in S}(\id-\Pi^0_k)\prod_{k\notin S}\Pi^0_k\big]V
=\sum_{T\subseteq S}(-1)^{|T|}\,V^\dagger\Pi^0_{([K]\setminus S)\cup T}V$,
a combination of those operators, for every click pattern $S$ (the elements of
the POVM of Theorem~\ref{thm:click-povm}). Conversely a contraction $\Gmat$ always
admits such a dilation (Lemma~\ref{lem:dilation}), so the two descriptions are
interchangeable.
\end{proof}

% claim receiver/lem:deficiency
\claimsummary{receiver/lem:deficiency}{The undetected fraction splits exactly into light lost through the dilation and light reaching declared modes that no port monitors, both positive semidefinite. It is the check that loss is counted exactly once.}
\claimbookmark{receiver/lem:deficiency}{Lemma}{lem:deficiency}{Decomposition of the loss deficiency}
\claimleadin{%
The identity below holds for every receiver instrument and is the correct
check that loss was accounted for exactly once.
}
\begin{lemma}[Decomposition of the loss deficiency]\label{lem:deficiency}
Let $\Recv$ be a receiver instrument with port Grams $\{A_k\}$ and let
$\Emd{\mathrm{un}}$ be the projector onto the declared modes lying in no
port. Then
\begin{equation}\label{eq:deficiencyid}
  \defic\;=\;\id_{\Inp}-A_\Sigma
  \;=\;\underbrace{\big(\id_{\Inp}
    -\Gmat_{\Inp}^\dagger\Gmat_{\Inp}\big)}_{\text{left through the
    dilation}}
  \;+\;\underbrace{\Gmat_{\Inp}^\dagger\,\Emd{\mathrm{un}}\,
    \Gmat_{\Inp}}_{\text{reached an unmonitored declared mode}} ,
\end{equation}
and both summands are positive semidefinite.
\end{lemma}
\claimmeta{\metalabel{Uses}\mref{Def.}{def:instrument}, \mref{Def.}{def:gram}
\metasep \metalabel{Assumes}\aref{A-TRUST}{ass:trust}, \aref{A-PASSIVE}{ass:passive}, \aref{A-VAC}{ass:vac}, \aref{A-MODES}{ass:modes}, \aref{A-EXACT}{ass:exact}, \aref{A-NOSEL}{ass:nosel}, \aref{A-ROUND}{ass:round}
\metasep \metalabel{Used by}\mref{Def.}{def:residuals}, \mref{Thm.}{thm:efficiency}, \mref{Rem.}{rem:wrongcheck}
\metasep newly derived
\metasep accepted
\metasep \metalabel{Oracles}\coderef{https://github.com/fraunhofersg/certiqs-key/blob/v0.1.0/tests/oracles/registry_receiver_povm.py}{registry_receiver_povm.py}}

\begin{proof}
The ports and the unmonitored declared modes partition $\Modes$, so
$\Emd{[K]}+\Emd{\mathrm{un}}=\id_{\Modes}$ and hence
$\Gmat_{\Inp}^\dagger\Gmat_{\Inp}=A_\Sigma+\Gmat_{\Inp}^\dagger
\Emd{\mathrm{un}}\Gmat_{\Inp}$. Rearranging gives \eqref{eq:deficiencyid}. The
first summand is positive semidefinite because $\Gmat$ is a contraction and the
restriction of $\id-\Gmat^\dagger\Gmat\succeq0$ to the principal submatrix on
$\Inp$ is $\id_{\Inp}-\Gmat_{\Inp}^\dagger\Gmat_{\Inp}$; the second is a
conjugated projector.
\end{proof}

% claim receiver/rem:unmonitored
\claimsummary{receiver/rem:unmonitored}{Deleting the row of an output mode that no port monitors changes no port Gram and hence no POVM element; the deleted amplitude is carried by the dilation. Which modes lie in a port is a declaration with no observable content beyond the Grams it produces.}
\claimbookmark{receiver/rem:unmonitored}{Remark}{rem:unmonitored}{Unmonitored output modes}
\begin{remark}[Unmonitored output modes]
\label{rem:unmonitored}
A mode that lies in no port is undetected: the amplitude reaching it is loss.
Nothing is lost by that convention. If a model retains an output mode that no
port monitors, deleting the corresponding row of $\Gmat$ changes no $A_J$
(the row does not appear in any $\Emd{J}\Gmat_{\Inp}$), hence changes no POVM
element; the deleted amplitude is carried by the dilation instead. The
partition into ``in some port'' and ``in no port'' is therefore a
\emph{declaration} with no observable content beyond the $A_k$ it produces.
\end{remark}
\claimmeta{\metalabel{Used by}\mref{Req.}{req:declared-port-map}
\metasep newly derived
\metasep accepted}

% claim receiver/def:clickalphabet
\claimsummary{receiver/def:clickalphabet}{A raw outcome is the set of detector ports that click, and the raw alphabet splits into the no-click outcome, the single-click set and the multiclick set. Any merging into a protocol's alphabet comes later.}
\claimbookmark{receiver/def:clickalphabet}{Definition}{def:clickalphabet}{Click-pattern alphabet}
\claimleadin{%
At $K=4$ the sets defined here are the raw, single-click and multiclick sets
of the optical-profile notation
(\href{../../../spec/notation.md\#37-optical-profile-symbols}{\nolinkurl{spec/notation.md}\texttt{\#37-optical-profile-symbols}}).
}
\begin{definition}[Click-pattern alphabet]\label{def:clickalphabet}
An \termdef{term-outcome}{outcome} is the set $S\subseteq[K]$ of ports that click; the
\termdef{term-raw-alphabet}{raw alphabet} is $2^{[K]}$, in bijection with $\{0,1\}^{K}$ by indicator
vectors. We write
\begin{equation}\label{eq:mcsc}
  \SCK\coloneqq\big\{\{k\}:k\in[K]\big\},\qquad
  \MCK\coloneqq\big\{S\subseteq[K]:|S|\ge2\big\}
\end{equation}
for the \termdef{term-single-click-set}{single-click set} and the \termdef{term-multiclick-set}{multiclick set}, so that
$2^{[K]}=\{\varnothing\}\sqcup\SCK\sqcup\MCK$, and $\varnothing$ is the
\termdef{term-no-click-outcome}{no-click outcome}.
\end{definition}
\claimmeta{\metalabel{Used by}\mref{Def.}{def:setting-family}
\metasep newly derived
\metasep accepted}

% claim receiver/thm:click-povm
\claimsummary{receiver/thm:click-povm}{The click patterns of a receiver instrument form a POVM whose elements follow from the port Grams by inclusion and exclusion. Every element is block diagonal in the total input photon number, and vacuum gives only the no-click outcome.}
\claimbookmark{receiver/thm:click-povm}{Theorem}{thm:click-povm}{Click-pattern POVM}
\claimleadin{%
The element of an outcome $S$ is the pull-back of ``every port in $S$ clicks
and no other does''. Positivity and completeness come from that projective
form, and item (ii) then gives the closed form in the port Grams.
}
\begin{mainclaim}
\begin{theorem}[Click-pattern POVM]\label{thm:click-povm}
Let $\Recv$ be a receiver instrument with \termref{term-measurement-isometry}{measurement isometry} $V$
(Definition~\ref{def:physmodel}) and port Grams $\{A_k\}$. Define
\begin{equation}\label{eq:MSdef}
  M_S\coloneqq V^\dagger\Big[\prod_{k\in S}\big(\id-\Pi^0_k\big)
  \prod_{k\notin S}\Pi^0_k\Big]V ,\qquad S\subseteq[K].
\end{equation}
Then:
\begin{enumerate}[label=(\roman*),leftmargin=3.0em]
\item \emph{$\{M_S\}$ is a POVM.} $M_S\succeq0$ for every $S$ and
  $\sum_{S\subseteq[K]}M_S=\id$. Moreover
  $\Pvac{J}\coloneqq V^\dagger\Pi^0_JV=\sum_{S\subseteq[K]\setminus J}M_S$ for
  every $J\subseteq[K]$.
\item \emph{Closed form.} $\Pvac{J}=\Fock(\id_{\Inp}-A_J)$ and
  \begin{equation}\label{eq:incexc}
    M_S=\sum_{T\subseteq S}(-1)^{|T|}\,
        \Fock\big(\id_{\Inp}-A_{([K]\setminus S)\cup T}\big).
  \end{equation}
\item \emph{Photon-number block structure.} Every $M_S$ commutes with the total
  photon number of the declared inputs, so
  $M_S=\bigoplus_{n\ge0}\blk{M}{\eqn}_S$ with
  $\blk{M}{\eqn}_S$ acting on $\Symn{n}(\mathbb{C}^{\Inp})$, and
  \begin{equation}\label{eq:blockMS}
    \blk{M}{\eqn}_S=\sum_{T\subseteq S}(-1)^{|T|}
    \Big(\big(\id-A_{([K]\setminus S)\cup T}\big)^{\otimes n}\Big)
    \Big|_{\Symn{n}} .
  \end{equation}
  In particular $\blk{M}{(=0)}_S=\delta_{S,\varnothing}$: no photon, no click.
\item \emph{The coarse elements.} With $A_\Sigma=A_{[K]}$,
  \begin{align}
    M_{\nc}&\coloneqq M_{\varnothing}=\Fock(\id-A_\Sigma),
      \label{eq:Mnc}\\
    M_{\singlec_k}&\coloneqq M_{\{k\}}
      =\Fock\big(\id-A_{[K]\setminus\{k\}}\big)-\Fock(\id-A_\Sigma),
      \label{eq:Msc}\\
    M_{\mc}&\coloneqq\sum_{S\in\MCK}M_S
      =\id+(K-1)\,\Fock(\id-A_\Sigma)
        -\sum_{k=1}^{K}\Fock\big(\id-A_{[K]\setminus\{k\}}\big).
      \label{eq:Mmc}
  \end{align}
\end{enumerate}
\end{theorem}
\end{mainclaim}
\claimmeta{\metalabel{Uses}\mref{Lem.}{lem:fock-functor}, \mref{Def.}{def:physmodel}, \mref{Lem.}{lem:vacelim}, \mref{Thm.}{thm:port-grams}
\metasep \metalabel{Used by}\mref{Rem.}{rem:positivity}, \mref{Rem.}{rem:normordelements}, \mref{Rem.}{rem:small-k}, \mref{Rem.}{rem:blockmeaning}, \mref{Thm.}{thm:representation-independence}, \mref{Def.}{def:coarse}, \mref{Thm.}{thm:coarse-graining}, \mref{Rem.}{rem:k-minus-1}, \mref{Def.}{def:exactbnd}, \mref{Cor.}{cor:unconditioned-family}, \mref{Rem.}{rem:notsharp}, \mref{Thm.}{thm:enclosed-blocks}, \mref{Thm.}{thm:factor-blocks}, \mref{Thm.}{thm:exact-blocks}, \mref{Thm.}{thm:dark-counts}, \mref{Prop.}{prop:dark-slack-bound}, \mref{Rem.}{rem:d2slack}, \mref{Prop.}{prop:noise-floor}, \mref{Prop.}{prop:efficiency-normalization}, \mref{Rem.}{rem:normalization-gain}, \extref{../reductions/contract.pdf}{reductions/req:recv-row-provenance}, \extref{../reductions/contract.pdf}{reductions/wpfss/cor:normalized-flat}, \extref{../reductions/contract.pdf}{reductions/wpfss/prop:route}, \extref{../reductions/contract.pdf}{reductions/wpfss/rem:receiver-imports}
\metasep newly derived
\metasep implemented (effective: accepted)
\metasep \metalabel{Code}\coderef{https://github.com/fraunhofersg/certiqs-key/blob/v0.1.0/src/certiqs_key/receiver/povm.py}{certiqs_key.receiver.povm:unconditioned_family}
\metasep \metalabel{Oracles}\coderef{https://github.com/fraunhofersg/certiqs-key/blob/v0.1.0/tests/oracles/registry_receiver_povm.py}{registry_receiver_povm.py}}

\begin{proof}
(i) The $\Pi^0_k$ are orthogonal projectors (Lemma~\ref{lem:fock-functor}\,(iv))
acting on the Fock spaces over pairwise disjoint mode sets, hence they commute:
$\Pi^0_k\Pi^0_j=\Fock\big((\id-\Emd{\{k\}})(\id-\Emd{\{j\}})\big)=\Pi^0_j
\Pi^0_k$ by Lemma~\ref{lem:fock-functor}\,(ii) and commutativity of the coordinate
projectors. A product of commuting orthogonal projectors is an orthogonal
projector, so the bracket in \eqref{eq:MSdef} is a product of commuting
orthogonal projectors and is therefore positive semidefinite; conjugating by
$V$ preserves that. Completeness: expanding the commuting product
$\prod_{k}\big(\Pi^0_k+(\id-\Pi^0_k)\big)=\id$ enumerates exactly the brackets
of \eqref{eq:MSdef} over $S\subseteq[K]$, so their sum is $\id$, and
$V^\dagger\id V=\id$ because $V$ is an isometry (Lemma~\ref{lem:vacelim}).
Summing \eqref{eq:MSdef} over $S\subseteq[K]\setminus J$ leaves the terms in
which every $k\in J$ contributes the factor $\Pi^0_k$, and the same expansion
of the remaining product gives $\prod_{k\in J}\Pi^0_k=\Pi^0_J$; conjugating by
$V$ gives the last identity.

(ii) $\Pvac{J}=\Fock(\id-A_J)$ is Theorem~\ref{thm:port-grams}\,(iii). Expanding the
commuting product,
$\prod_{k\in S}(\id-\Pi^0_k)=\sum_{T\subseteq S}(-1)^{|T|}\prod_{k\in T}
\Pi^0_k$, and multiplying by $\prod_{k\notin S}\Pi^0_k$ gives
$\sum_{T\subseteq S}(-1)^{|T|}\Pi^0_{([K]\setminus S)\cup T}$, since
$T\subseteq S$ is disjoint from $[K]\setminus S$. Applying
$V^\dagger(\cdot)V$ termwise and using $\Pvac{J}=\Fock(\id-A_J)$ gives
\eqref{eq:incexc}.

(iii) Each $\Fock(\cdot)$ is photon-number block diagonal by
Lemma~\ref{lem:fock-functor}\,(i), and \eqref{eq:blockMS} is \eqref{eq:incexc}
restricted to $\Symn{n}(\mathbb{C}^{\Inp})$. At $n=0$ every
$\Fock(\cdot)^{(0)}=1$, so $\blk{M}{(=0)}_S=\sum_{T\subseteq S}(-1)^{|T|}
=\sum_{j=0}^{|S|}\binom{|S|}{j}(-1)^j=(1-1)^{|S|}$, which is $1$ for
$S=\varnothing$ and $0$ otherwise.

(iv) \eqref{eq:Mnc} is \eqref{eq:incexc} at $S=\varnothing$ and \eqref{eq:Msc}
at $S=\{k\}$. For \eqref{eq:Mmc}, completeness gives
$M_{\mc}=\id-M_{\varnothing}-\sum_kM_{\{k\}}$, and substituting the previous
two lines contributes $-\Fock(\id-A_\Sigma)$ once from $M_\varnothing$ and
$+\Fock(\id-A_\Sigma)$ once from each of the $K$ single-click terms, for a net
coefficient $K-1$.
\end{proof}

% claim receiver/rem:positivity
\claimsummary{receiver/rem:positivity}{The inclusion-exclusion formula is an alternating sum and nothing in its form makes it positive. Positivity comes first, from commuting projective measurements pulled back through an isometry, and the single-click case also follows from monotonicity of the lift.}
\claimbookmark{receiver/rem:positivity}{Remark}{rem:positivity}{Source of positivity of the click-pattern elements}
\begin{remark}[Source of positivity of the click-pattern elements]
\label{rem:positivity}
Formula \eqref{eq:incexc} is an alternating sum of positive operators and
nothing about its form makes it positive; the element of a non-empty outcome
$S$ has $2^{|S|}$ terms, half of them with a minus sign. Positivity is established \emph{before} the expansion, in
part~(i) of Theorem~\ref{thm:click-povm}, from the physical construction: $K$
commuting binary projective measurements on disjoint output mode sets, pulled
back through an isometry. Definition~\ref{def:physmodel} is set up for that
order of argument. A derivation that started from \eqref{eq:incexc} would
have to prove positivity of an alternating sum of Fock functors, which is both
harder and unnecessary. For a small independent check,
$\eqref{eq:Msc}\succeq0$ follows separately from
Lemma~\ref{lem:fock-monotone}, since $A_{[K]\setminus\{k\}}\preceq A_\Sigma$ gives
$\Fock(\id-A_{[K]\setminus\{k\}})\succeq\Fock(\id-A_\Sigma)$.
\end{remark}
\claimmeta{\metalabel{Uses}\mref{Lem.}{lem:fock-monotone}, \mref{Def.}{def:physmodel}, \mref{Thm.}{thm:click-povm}
\metasep newly derived
\metasep accepted}

% claim receiver/rem:normordelements
\claimsummary{receiver/rem:normordelements}{Each joint no-click element is the normal-ordered exponential of the summed Gram A\_J of its ports, the form used in the optics literature. For one port and one input mode with Gram A\_1 = 1 this is the perfect on/off detector pair of the cited literature.}
\claimbookmark{receiver/rem:normordelements}{Remark}{rem:normordelements}{Click-pattern elements in normal-ordered form}
\begin{remark}[Click-pattern elements in normal-ordered form]\label{rem:normordelements}
By Proposition~\ref{prop:normord}, $\Pvac{J}=\normord{\exp(-\vec
a^\dagger A_J\vec a)}$, which is the form these operators take in the optics
literature; \eqref{eq:incexc} then reads
$M_S=\sum_{T\subseteq S}(-1)^{|T|}\normord{\exp(-\vec a^\dagger
A_{([K]\setminus S)\cup T}\vec a)}$. At $K=1$ and a single input mode with
$A_1=\lambda$ this is $M_\varnothing=\normord{e^{-\lambda\hat n}}$ and
$M_{\{1\}}=\normord{\id-e^{-\lambda\hat n}}$, which at $\lambda=1$ is the
perfect on/off detector pair $\hat\pi_0=\dyad{0}=\normord{e^{-\hat n}}$,
$\hat\pi_1=\normord{\hat 1-e^{-\hat n}}$ of
\cite[Eq.~(6)]{sperling2012}. The comparison is corroboration only: the
identity is Proposition~\ref{prop:normord} and the POVM property is
Theorem~\ref{thm:click-povm}\,(i).
\end{remark}
\claimmeta{\metalabel{Uses}\mref{Prop.}{prop:normord}, \mref{Thm.}{thm:click-povm}
\metasep newly derived
\metasep accepted}

% claim receiver/rem:small-k
\claimsummary{receiver/rem:small-k}{With one port the multiclick element vanishes, as the closed form confirms; with two ports each single-click element is the other port dark minus both dark. These cases check the inclusion-exclusion formula by hand.}
\claimbookmark{receiver/rem:small-k}{Remark}{rem:small-k}{Small numbers of ports}
\begin{remark}[Small numbers of ports]\label{rem:small-k}
$K=1$: $M_\varnothing=\Fock(\id-A_1)$, $M_{\{1\}}=\id-\Fock(\id-A_1)$, and
$\MCK=\varnothing$ so $M_{\mc}=0$, consistent with \eqref{eq:Mmc}, whose
right-hand side is $\id+0-\Fock(\id-A_\varnothing)=\id-\Fock(\id)=0$.
$K=2$, $S=\{1\}$: $M_{\{1\}}=\Fock(\id-A_2)-\Fock(\id-A_1-A_2)$, i.e.\ ``port
$2$ dark'' minus ``both dark''. $K=2$, $\MCK=\{\{1,2\}\}$:
$M_{\{1,2\}}=\id+\Fock(\id-A_\Sigma)-\Fock(\id-A_2)-\Fock(\id-A_1)$.
\end{remark}
\claimmeta{\metalabel{Uses}\mref{Thm.}{thm:click-povm}
\metasep newly derived
\metasep accepted}

% claim receiver/rem:blockmeaning
\claimsummary{receiver/rem:blockmeaning}{The block structure is in the total input photon number, which equals the channel photon number only without a probe; an occupation-diagonal probe state suffices to recover the channel grading. The emitted low blocks decompose the POVM and are not compressions or a squashing argument.}
\claimbookmark{receiver/rem:blockmeaning}{Remark}{rem:blockmeaning}{Scope of the photon-number block structure}
\claimleadin{%
The difference between an operator block and a compression is the third
caveat of the sector convention
(\href{../../../spec/notation.md\#24-sector-decompositions-and-block-restrictions}{\nolinkurl{spec/notation.md}\texttt{\#24-sector-decompositions-and-block-restrictions}}).
}
\begin{remark}[Scope of the photon-number block structure]\label{rem:blockmeaning}
The grading in Theorem~\ref{thm:click-povm}\,(iii) is by the \emph{total} photon
number of the declared inputs $\Inp=\InpC\sqcup\InpS$. When a probe is
declared, that is not the same as the channel photon number, and the
downstream modules need the \emph{channel} grading; recovering it is what the
occupation-diagonality of the probe states (Definition~\ref{def:probe}) is
required for, and the submodule that conditions on the probe states proves
that it suffices. When no probe is
declared ($\InpS=\varnothing$) the two gradings coincide. The emitted
low-photon blocks $\blk{M}{\leqN}_S$ are operators on
$\Hs{\Chan_{\leq N}}$ and are \emph{not} the compressions
$\Pnum{\le N}M_S\Pnum{\le N}$; emitting them is a decomposition of the
physical POVM, never a squashing argument.
\end{remark}
\claimmeta{\metalabel{Uses}\mref{Thm.}{thm:click-povm}
\metasep newly derived
\metasep accepted}

\section{Representation independence and coarse-graining}\label{group-receiver/thm:representation-independence}
\sectionintro{Two receivers induce the same click POVM exactly when their port Grams agree, every \termref{term-admissible-gram-family}{admissible Gram family} is realised by some instrument, and a one-mode-per-port description exists exactly when every Gram has rank at most one (\hyperlink{claim-receiver/thm:representation-independence}{Theorem~\ref*{thm:representation-independence}}). Phases inside a single port after all interference are invisible; a phase on an earlier element whose output paths are later recombined is not covered and can change the Grams. The main result, \hyperlink{claim-receiver/thm:coarse-graining}{Theorem~\ref*{thm:coarse-graining}}, is that under a \termref{term-declared-coarse-graining}{declared coarse-graining}, which merges raw outcomes into the protocol's alphabet, lower bounds on disjoint raw outcomes add, and a bound on a merged outcome gives nothing about its parts.}

% claim receiver/thm:representation-independence
\claimsummary{receiver/thm:representation-independence}{Two receivers have the same click-pattern POVM exactly when their port Grams agree, every admissible Gram family is realised by some instrument, and a one-mode-per-port description exists exactly when every Gram has rank at most one.}
\claimbookmark{receiver/thm:representation-independence}{Theorem}{thm:representation-independence}{Representation independence of the click-pattern POVM}
\claimleadin{%
The port Grams carry exactly the information the click statistics can
reveal, no more and no less. The theorem makes this precise in both
directions and says when a description with one \termref{term-optical-mode}{optical mode} per
detector port exists. Under the mode-insensitive port response of
Assumption~\ref{ass:port}, item (iii) is what makes $A_k$, rather than the
individual rows feeding port $k$, the right object.
}
\begin{theorem}[Representation independence of the click-pattern POVM]\label{thm:representation-independence}
Let $\Recv,\Recv'$ be \termref{term-receiver-instrument}{receiver instruments} on the same declared input space
$\mathbb{C}^{\Inp}$ with the same number of ports $K$.
\begin{enumerate}[label=(\roman*),leftmargin=3.0em]
\item \emph{Determination, in both directions.} $\Recv$ and $\Recv'$ induce the
  same click-pattern POVM if and only if $A_k=A_k'$ for every $k$. In
  particular the POVM depends on the \termref{term-instrument-matrix}{instrument matrix} only through the
  port Grams.
\item \emph{Realizability.} Every \termref{term-admissible-gram-family}{admissible family} $\{A_k\}_{k=1}^{K}$ on
  $\mathbb{C}^{\Inp}$ (Definition~\ref{def:gram}) is the port-Gram family of
  some receiver instrument, in which port $k$ collects at most
  $\rk A_k\le|\Inp|$ modes.
\item \emph{Port-local unitaries are invisible.} If
  $\Gmat'=\big(\bigoplus_kV_k\oplus V_{\mathrm{rest}}\big)\Gmat$ with each
  $V_k$ unitary on $\mathbb{C}^{\Port{k}}$ and $V_{\mathrm{rest}}$ unitary on
  the remaining modes, then $A_k'=A_k$ for every $k$.
\item \emph{One mode per port.} There is a receiver instrument with
  $|\Port{k}|=1$ for every $k$ inducing the same POVM as $\Recv$ if and only if
  $\rk A_k\le1$ for every $k$. In that case, writing
  $A_k=\lambda_ku_ku_k^\dagger$ with $\lambda_k\ge0$ and $u_k$ a unit vector,
  the single row of port $k$ is $g_k=e^{i\varphi_k}\sqrt{\lambda_k}\,
  u_k^\dagger$, unique up to the phase $\varphi_k$.
\end{enumerate}
\end{theorem}
\claimmeta{\metalabel{Uses}\mref{Def.}{def:gram}, \mref{Thm.}{thm:port-grams}, \mref{Thm.}{thm:click-povm}
\metasep \metalabel{Used by}\mref{Cor.}{cor:phases}, \mref{Rem.}{rem:rankone}, \mref{Def.}{def:exact-route}, \mref{Def.}{def:residuals}, \mref{Req.}{req:one-mode-per-port}, \mref{Ex.}{ex:bb84}, \mref{Thm.}{thm:efficiency}, \mref{Rem.}{rem:effrank}, \mref{Def.}{def:normalized-station}, \extref{../reductions/contract.pdf}{reductions/req:recv-row-provenance}
\metasep newly derived
\metasep implemented (effective: accepted)
\metasep \metalabel{Code}\coderef{https://github.com/fraunhofersg/certiqs-key/blob/v0.1.0/src/certiqs_key/receiver/grams.py}{certiqs_key.receiver.grams:realizing_instrument}
\metasep \metalabel{Oracles}\coderef{https://github.com/fraunhofersg/certiqs-key/blob/v0.1.0/tests/oracles/registry_receiver_povm.py}{registry_receiver_povm.py}}

\begin{proof}
\emph{Idea.} The \termref{term-joint-no-click-elements}{joint no-click elements} are lifts of $\id-A_J$, and the
click-pattern POVM and these elements determine each other; the one-photon
block of a lift is the matrix itself, which recovers the Grams. Realisation
takes square roots of the Grams as rows. For (iv), by (i) a
one-mode-per-port instrument with the same POVM has the same Grams, each of
rank at most one by Theorem~\ref{thm:port-grams}\,(ii); conversely, Grams of
rank at most one give single rows that realise the family.

(i) If $A_k=A_k'$ for all $k$ then $A_J=A_J'$ for all $J$, so by
Theorem~\ref{thm:port-grams}\,(iii) the two receivers have the same joint no-click
elements $\Pvac{J}=\Fock(\id-A_J)$, and by \eqref{eq:incexc} of
Theorem~\ref{thm:click-povm} the same $M_S$ for every $S$. Conversely, suppose the
POVMs agree. ``No port in $J$ clicks'' is the union of the disjoint outcomes
whose click set avoids $J$, so
\begin{equation}\label{eq:recoverpivac}
  \Pvac{J}=\sum_{S\subseteq[K]\setminus J}M_S
\end{equation}
(Theorem~\ref{thm:click-povm}\,(i)), and hence $\Fock(\id-A_J)=\Fock(\id-A_J')$.
Reading off the one-photon block, whose value is
$\Fock(\id-A_J)^{(1)}=\id-A_J$, gives $A_J=A_J'$ for all $J$, and in
particular $A_{\{k\}}=A_k=A_k'$.

(ii) Write $r_k\coloneqq\rk A_k\le|\Inp|$ and diagonalize
$A_k=U_k\Lambda_kU_k^\dagger$ over its non-zero eigenvalues, with $\Lambda_k$
the positive diagonal $r_k\times r_k$ matrix of those eigenvalues and
$U_k\in\mathbb{C}^{\Inp\times r_k}$ with orthonormal columns. Put
$C_k\coloneqq\Lambda_k^{1/2}U_k^\dagger\in\mathbb{C}^{r_k\times\Inp}$, so
that $C_k^\dagger C_k=U_k\Lambda_kU_k^\dagger=A_k$ with $r_k$ rows. Let
$\Modes'$ consist of the $|\Inp|$ input modes, with the split of $\Inp$ that
$\Recv$ declares, together with $K$ blocks of $r_1,\dots,r_K$ output modes,
let $\Port{k}'$ be the $k$-th output block, and let $\Gmat'$ have the rows of
$C_k$ in the block $\Port{k}'$ and zeros elsewhere. Then
$\Emd{\{k\}}\Gmat'_{\Inp}=C_k$, so $A_k'=C_k^\dagger C_k=A_k$, and
$\Gmat'^\dagger_{\Inp}\Gmat'_{\Inp}=\sum_kA_k\preceq\id$, so $\Gmat'$ is a
contraction and extends to a contraction on all of $\mathbb{C}^{\Modes'}$ by
padding with zeros. Port $k$ collects $r_k=\rk A_k$ modes.

(iii) The ports are pairwise disjoint, so
$\Emd{\{k\}}\big(\bigoplus_jV_j\oplus V_{\mathrm{rest}}\big)
=V_k\Emd{\{k\}}$ as maps into $\mathbb{C}^{\Port{k}}$. Hence
$A_k'=\Gmat_{\Inp}^\dagger\Emd{\{k\}}V_k^\dagger V_k\Emd{\{k\}}\Gmat_{\Inp}
=\Gmat_{\Inp}^\dagger\Emd{\{k\}}\Gmat_{\Inp}=A_k$.

(iv) If a one-mode-per-port instrument induces the same POVM then by (i) its
Grams are the $A_k$, and by Theorem~\ref{thm:port-grams}\,(ii) each has rank at most
$|\Port{k}|=1$. Conversely, if $\rk A_k\le1$, write
$A_k=\lambda_ku_ku_k^\dagger$ (possible for a rank-$\le1$ positive semidefinite
matrix, with $\lambda_k=0$ allowed) and set
$g_k\coloneqq\sqrt{\lambda_k}u_k^\dagger$. Stacking the $g_k$ as $K$ single
rows gives a matrix with $\sum_kg_k^\dagger g_k=\sum_kA_k\preceq\id$, hence a
contraction, whose port Grams are the $A_k$; by (i) it induces the same POVM.
For uniqueness, let $g$ be a row vector with $g^\dagger g=\lambda uu^\dagger$.
Taking traces, $\lVert g\rVert^2=\lambda$. The range of $g^\dagger g$ is spanned
by $g^\dagger$ when $g\neq0$, and equals $\operatorname{span}\{u\}$, so
$g^\dagger=cu$ for some $c\in\mathbb{C}$ with $|c|^2=\lambda$, i.e.\
$g=\bar c\,u^\dagger$ with $\bar c=e^{i\varphi}\sqrt\lambda$. If $\lambda=0$
then $g=0$ and the statement is vacuous.
\end{proof}

% claim receiver/cor:phases
\claimsummary{receiver/cor:phases}{A unitary acting inside a single detector port after all interference changes no click statistic. A phase inserted earlier, on an element whose output paths are later recombined, is not covered by this invariance.}
\claimbookmark{receiver/cor:phases}{Corollary}{cor:phases}{Unobservability of per-port phases}
\claimleadin{%
``Global phases do not matter'' holds here only with a stated hypothesis:
the unitary acts inside one port after all interference. That is why every
entry of the \termref{term-component-catalogue}{component catalogue} carries its phases, rather than dropping
those that look global on the element they belong to
(\href{../../../spec/modules/receiver/design.md}{\nolinkurl{spec/modules/receiver/design.md}}).
}
\begin{corollary}[Unobservability of per-port phases]
\label{cor:phases}
\begin{enumerate}[label=(\roman*),leftmargin=3.0em]
\item Multiplying the output row of a port by a phase, or more generally
  applying any unitary \emph{inside} a single port's mode set, changes no
  click statistic (Theorem~\ref{thm:representation-independence}\,(iii)).
\item \emph{The hypothesis matters}: the unitary must act inside one port's
  modes and after all interference has taken place. A phase that is global
  on an element in the \emph{middle} of the receiver can become relative once
  that element's output paths are recombined with other paths, and it can
  then change the Grams.
\end{enumerate}
\end{corollary}
\claimmeta{\metalabel{Uses}\mref{Thm.}{thm:representation-independence}
\metasep \metalabel{Used by}\mref{Ex.}{ex:bb84}
\metasep newly derived
\metasep accepted}

\begin{proof}
Item (i) is Theorem~\ref{thm:representation-independence}\,(iii) applied to the unitary in question,
which is port-local by hypothesis. For item (ii), a unitary that acts on an
element in the middle of the receiver is not of the form
$\bigoplus_kV_k\oplus V_{\mathrm{rest}}$ on the output modes once its paths
are recombined, so part~(iii) does not apply to it; and it can change the
Grams. Take two modes $a,b$ with $\Inp=\InpC=\{a\}$, $b$ in vacuum, and
ports $\{a\}$ and $\{b\}$; apply the balanced beamsplitter
$H=2^{-1/2}\big(\begin{smallmatrix}1&1\\1&-1\end{smallmatrix}\big)$ to
$(a,b)$, then the phase $e^{i\theta}$ on $a$, a global phase on that
one-mode element, then $H$ again. The instrument matrix is
$\Gmat=H\operatorname{diag}(e^{i\theta},1)H$, whose input column is
$\tfrac12\big(e^{i\theta}+1,\ e^{i\theta}-1\big)^{\transpose}$, so
$A_1=\big|\tfrac12(e^{i\theta}+1)\big|^2=\cos^2(\theta/2)$ depends on
$\theta$.
\end{proof}

% claim receiver/cor:extinction
\claimsummary{receiver/cor:extinction}{At two input modes a nonzero port Gram has rank one exactly when the port extinguishes some polarization completely. At more input modes one null direction bounds the rank only by the number of input modes minus one.}
\claimbookmark{receiver/cor:extinction}{Corollary}{cor:extinction}{Extinction criterion for rank one at two input modes}
\claimleadin{%
The rank criterion of the one-mode-per-port description has a laboratory
reading when there are two input modes, such as two polarizations.
}
\begin{corollary}[Extinction criterion for rank one at two input modes]\label{cor:extinction}
Suppose $|\Inp|=2$. Then $\rk A_k\le1$ holds if and only if either $A_k=0$ or
there is a unit vector $v\in\mathbb{C}^{2}$ with $v^\dagger A_kv=0$, that is,
a polarization the port extinguishes completely. For $|\Inp|>2$ the two
conditions are \emph{not} equivalent: the existence of one null direction of a
positive semidefinite $A_k$ says only $\rk A_k\le|\Inp|-1$.
\end{corollary}
\claimmeta{newly derived
\metasep accepted}

\begin{proof}
For $A_k\succeq0$ and $v$ a unit vector, $v^\dagger A_kv=0$ iff
$A_k^{1/2}v=0$ iff $A_kv=0$, so a null direction is a kernel vector. On
$\mathbb{C}^2$ a nonzero positive semidefinite matrix with a kernel vector has
rank exactly $1$; conversely a rank-$1$ matrix on $\mathbb{C}^2$ has a
one-dimensional kernel. On $\mathbb{C}^d$ a kernel vector only gives
$\rk A_k\le d-1$.
\end{proof}

% claim receiver/rem:rankone
\claimsummary{receiver/rem:rankone}{No result of this contract assumes one optical mode per detector port: a port collecting several modes can have a Gram of rank above one. The one-mode-per-port description holds exactly under the rank criterion of \hyperlink{claim-receiver/thm:representation-independence}{Theorem~\ref*{thm:representation-independence}} (iv).}
\claimbookmark{receiver/rem:rankone}{Remark}{rem:rankone}{Ports of rank greater than one}
\begin{remark}[Ports of rank greater than one]\label{rem:rankone}
Example~\ref{ex:bb84} ends with four rank-one Grams, which is why the textbook
$4\times2$ BB84 instrument is legitimate \emph{there}. Nothing in this
contract assumes one mode per port anywhere else: the click-pattern POVM,
the theorems built on it and the decomposition of the \termref{term-loss-deficiency}{loss deficiency},
among them Theorems~\ref{thm:click-povm}, \ref{thm:coarse-graining},
\ref{thm:multiclick-bound}, \ref{thm:emitted-bound} and
\ref{thm:exact-and-evaluated} and Lemma~\ref{lem:deficiency} of the core
and Theorems~\ref{thm:efficiency} and~\ref{thm:dark-counts} of its
submodules, are proved for arbitrary $\{A_k\}$, and a port that collects
several modes, such as a multimode fiber onto one active area or two spatial
modes onto one detector, can have $\rk A_k>1$, and every result above holds
when it does. The criterion for a
one-mode-per-port description is Theorem~\ref{thm:representation-independence}\,(iv); it is decidable
from the emitted Grams, and Remark~\ref{rem:effrank} exhibits the one
operation in this contract, apart from declaring a zero efficiency, that can
change its answer.
Requirement~\ref{req:one-mode-per-port} states the rule the engine follows.
\end{remark}
\claimmeta{\metalabel{Uses}\mref{Thm.}{thm:representation-independence}
\metasep \metalabel{Used by}\mref{Rem.}{rem:effrank}
\metasep newly derived
\metasep accepted}

% claim receiver/def:coarse
\claimsummary{receiver/def:coarse}{A declared surjection from raw click patterns onto the protocol's outcome alphabet defines the emitted measurement. Each emitted element is the sum of the click-pattern elements in its fibre.}
\claimbookmark{receiver/def:coarse}{Definition}{def:coarse}{Declared coarse-graining and the emitted POVM}
\claimleadin{%
The protocol reads an outcome alphabet that is usually coarser than the raw
click pattern, for instance one that merges every pattern with two or more
clicks into a single multiclick outcome.
}
\begin{definition}[Declared coarse-graining and the emitted POVM]\label{def:coarse}
A \termdef{term-declared-coarse-graining}{declared coarse-graining} is a surjection
$\cg:2^{[K]}\to\Vset$ onto the declared emitted alphabet, supplied as data by
the protocol declaration, $[K]$ indexing the \termref{term-detector-ports}{detector ports} of a
\termref{term-receiver-instrument}{receiver instrument} (Definition~\ref{def:instrument}). With $M_S$ the click-pattern elements of
Theorem~\ref{thm:click-povm} (\eqref{eq:MSdef}), the \termdef{term-emitted-povm}{emitted POVM} is
\begin{equation}\label{eq:coarse}
  F_v\coloneqq\sum_{S\in\cg^{-1}(v)}M_S ,\qquad v\in\Vset .
\end{equation}
\end{definition}
\claimmeta{\metalabel{Uses}\mref{Def.}{def:instrument}, \mref{Thm.}{thm:click-povm}
\metasep \metalabel{Used by}\mref{Thm.}{thm:enclosed-blocks}, \mref{Thm.}{thm:factor-blocks}, \mref{Def.}{def:dark-composed}, \mref{Prop.}{prop:noise-floor}, \mref{Prop.}{prop:efficiency-normalization}
\metasep newly derived
\metasep implemented (effective: accepted)
\metasep \metalabel{Code}\coderef{https://github.com/fraunhofersg/certiqs-key/blob/v0.1.0/src/certiqs_key/receiver/coarse.py}{certiqs_key.receiver.coarse:CoarseGraining}}

% claim receiver/thm:coarse-graining
\claimsummary{receiver/thm:coarse-graining}{The emitted POVM of a declared coarse-graining is a POVM, block diagonal in the photon number, and lower bounds on disjoint raw outcomes add. A bound on a merged outcome gives nothing about the click patterns it merges.}
\claimbookmark{receiver/thm:coarse-graining}{Theorem}{thm:coarse-graining}{Coarse-graining of the click-pattern POVM}
\claimleadin{%
Merging outcomes keeps everything the POVM structure needs and loses
information in one direction only. Part (iii) is why the
\termref{term-strengthening-slot}{strengthening slot} of Definition~\ref{def:slot} is indexed by
outcome sets rather than by emitted letters: by (ii) a bound on a set of raw
outcomes passes to every emitted outcome whose fibre contains the set, while
by (iii) a bound on an emitted outcome cannot be split back.
}
\begin{mainclaim}
\begin{theorem}[Coarse-graining of the click-pattern POVM]\label{thm:coarse-graining}
Let $\cg$ be a \termref{term-declared-coarse-graining}{declared coarse-graining}.
\begin{enumerate}[label=(\roman*),leftmargin=3.0em]
\item $\{F_v\}_{v\in\Vset}$ is a POVM on $\mathcal{F}(\mathbb{C}^{\Inp})$, and
  every $F_v$ is block diagonal in the photon number with
  $\blk{F}{\eqn}_v=\sum_{S\in\cg^{-1}(v)}\blk{M}{\eqn}_S$.
\item \emph{Bounds add.} If $\mathcal{O}\subseteq2^{[K]}$ and
  $\blk{M}{\gtN}_S\succeq p_S\,\id^{(>N)}$ for every $S\in\mathcal{O}$ with
  $p_S\ge0$, then
  $\sum_{S\in\mathcal{O}}\blk{M}{\gtN}_S\succeq
  \big(\sum_{S\in\mathcal{O}}p_S\big)\id^{(>N)}$. In particular, if
  $\cg^{-1}(v)\supseteq\mathcal{O}$ then
  $\blk{F}{\gtN}_v\succeq\big(\sum_{S\in\mathcal{O}}p_S\big)\id^{(>N)}$.
\item \emph{Coarsening is lossy in one direction only.} A bound on a coarse
  outcome does not descend to its fibres: $F_v\succeq p\,\id$ implies nothing
  about the individual $M_S$ with $\cg(S)=v$.
\end{enumerate}
\end{theorem}
\end{mainclaim}
\claimmeta{\metalabel{Uses}\mref{Lem.}{lem:fock-block-spectrum}, \mref{Thm.}{thm:click-povm}
\metasep \metalabel{Used by}\mref{Thm.}{thm:enclosed-blocks}, \mref{Thm.}{thm:exact-blocks}, \mref{Def.}{def:residuals}, \mref{Prop.}{prop:res5}, \extref{../reductions/contract.pdf}{reductions/req:recv-row-provenance}, \extref{../reductions/contract.pdf}{reductions/wpfss/cor:normalized-flat}, \extref{../reductions/contract.pdf}{reductions/wpfss/prop:route}, \extref{../reductions/contract.pdf}{reductions/wpfss/rem:receiver-imports}
\metasep newly derived
\metasep implemented (effective: accepted)
\metasep \metalabel{Code}\coderef{https://github.com/fraunhofersg/certiqs-key/blob/v0.1.0/src/certiqs_key/receiver/coarse.py}{certiqs_key.receiver.coarse:coarse_block}
\metasep \metalabel{Oracles}\coderef{https://github.com/fraunhofersg/certiqs-key/blob/v0.1.0/tests/oracles/registry_receiver_povm.py}{registry_receiver_povm.py}}

\begin{proof}
(i) Each $F_v$ is a finite sum of positive operators, hence positive; the
fibres $\cg^{-1}(v)$ partition $2^{[K]}$ because $\cg$ is a map, so
$\sum_vF_v=\sum_SM_S=\id$ by Theorem~\ref{thm:click-povm}\,(i). Block diagonality is
inherited from each summand (Theorem~\ref{thm:click-povm}\,(iii)), and the block of a
finite sum is the sum of the blocks.

(ii) The Löwner order is preserved by sums:
$\sum_{S\in\mathcal{O}}\blk{M}{\gtN}_S-\big(\sum_Sp_S\big)\id^{(>N)}
=\sum_{S\in\mathcal{O}}\big(\blk{M}{\gtN}_S-p_S\id^{(>N)}\big)\succeq0$. For
the second sentence add the remaining, positive, elements of the fibre.

(iii) Take $K=2$ and $\Vset=\{v\}$ with $\cg$ constant, so $F_v=\id$ and the
bound $F_v\succeq1\cdot\id$ is available. Choose port Grams with
$\lmin(A_\Sigma)=0$, i.e.\ $A_\Sigma$ singular: any lossy receiver with an
input direction no port sees. By Lemma~\ref{lem:fock-block-spectrum},
$\lmax\big(\blk{M}{\eqn}_{\varnothing}\big)=(1-\lmin(A_\Sigma))^n=1$, so on
every block there is a unit vector $\ket{\psi}$ with
$\bra{\psi}\blk{M}{\eqn}_{\varnothing}\ket{\psi}=1$; completeness and
positivity then force
$\bra{\psi}\blk{M}{\eqn}_{S}\ket{\psi}=0$ for every $S\ne\varnothing$, and in
particular $\lmin\big(\blk{M}{\eqn}_{\{1,2\}}\big)=0$. So a bound on the
coarse outcome gives nothing about its fibres.
\end{proof}

\section{The analytic high-photon bound}\label{group-receiver/def:bounddata}
\sectionintro{A squashing argument needs a lower bound on the \termref{term-multiclick-element}{multiclick element} valid on every photon-number block above the cutoff. Extreme eigenvalues of the no-click operators give a scalar \termref{term-bound-function}{bound function} on each block (\hyperlink{claim-receiver/thm:multiclick-bound}{Theorem~\ref*{thm:multiclick-bound}}); it is not monotone in general. Where its value at zero photons is at most zero, the bound cannot decrease from any photon number n \ensuremath{\ge} 1 at which it is non-negative, so the non-negative part of the infimum over the tail is that of its value at the first photon number above the cutoff (\hyperlink{claim-receiver/lem:collapse}{Lemma~\ref*{lem:collapse}}); otherwise a finite window with an analytic tail gives a valid bound that converges to the infimum when c \textless{} 1, with a stated residual gap (\hyperlink{claim-receiver/lem:windowtail}{Lemma~\ref*{lem:windowtail}}).}

% claim receiver/def:bounddata
\claimsummary{receiver/def:bounddata}{Ordered scalars between 0 and 1, one pair shared and one pair per port, define an exponential bound function of the photon number for the multiclick element. Admissibility is exact, and distinct from the tolerance rules that admit a floating-point evaluation.}
\claimbookmark{receiver/def:bounddata}{Definition}{def:bounddata}{Admissible bound data and the bound function}
\claimleadin{%
What a squashing argument consumes is, for every declared probe setting, one
number: a lower bound on the smallest eigenvalue of that setting's
\termref{term-multiclick-element}{multiclick element}, valid on \emph{every} photon-number block above the
declared cutoff. The bound itself is three lines from
Lemma~\ref{lem:fock-block-spectrum}, which gives the spectrum of a lifted no-click
operator on each block. Making the quantifier ``every block above the
cutoff'' into a finite computation is the work of this section. The
definition closes with how its one computed condition is decided in the
floating-point \termref{term-evaluation-model}{evaluation model} of Definition~\ref{def:evaluation-model}.
}
\begin{definition}[Admissible bound data and the bound function]\label{def:bounddata}
A tuple $\big(\gamma,\{\beta_k\},c,\{d_k\}\big)$ with $K\ge1$ entries is
\termdef{term-admissible-bound-data}{admissible bound data} when
\begin{equation}\label{eq:bounddata}
  0\le\gamma\le\beta_k\le1,\qquad 0\le c\le d_k\le1
  \qquad(k=1,\dots,K),
\end{equation}
and the associated \termdef{term-bound-function}{bound function} is
\begin{equation}\label{eq:beta}
  \Bnd(n)\coloneqq 1+(K-1)\,\gamma\,c^{\,n}-\sum_{k=1}^{K}\beta_k\,d_k^{\,n},
  \qquad n\ge0 ,
\end{equation}
with the convention $0^0=1$. Admissibility \eqref{eq:bounddata} is exact. The
one part of it that this contract decides on computed data is $0\le c$; the
other parts are not separately tested (Definition~\ref{def:tolmap}\,(iv)).
In floating point that part is decided by
the widened test $c\ge-\varepsilon_2$ of Definition~\ref{def:tolmap}\,(iv),
which decides it on \emph{every} input route and is not displaced by the
route-(ii) admission of Definition~\ref{def:tolmap}\,(v); the computed $c$
enters \eqref{eq:beta} unclamped.
\end{definition}
\claimmeta{\metalabel{Uses}\mref{Lem.}{lem:fock-block-spectrum}
\metasep \metalabel{Used by}\mref{Thm.}{thm:multiclick-bound}, \mref{Lem.}{lem:collapse}, \mref{Rem.}{rem:collapsesharp}, \mref{Lem.}{lem:windowtail}, \mref{Def.}{def:setting-family}, \mref{Def.}{def:pmc}, \mref{Rem.}{rem:noslack}, \mref{Rem.}{rem:exacteig}, \mref{Def.}{def:enclosed-weight}, \mref{Thm.}{thm:enclosed-weight}, \mref{Def.}{def:tolmap}, \mref{Rem.}{rem:bound-data-band}, \mref{Rem.}{rem:cdirection}, \mref{Thm.}{thm:exact-and-evaluated}, \mref{Rem.}{rem:ia}, \mref{Rem.}{rem:d2slack}
\metasep newly derived
\metasep implemented (effective: accepted)
\metasep \metalabel{Code}\coderef{https://github.com/fraunhofersg/certiqs-key/blob/v0.1.0/src/certiqs_key/receiver/bound.py}{certiqs_key.receiver.bound:BoundData}}

% claim receiver/thm:multiclick-bound
\claimsummary{receiver/thm:multiclick-bound}{For no-click operators of prefactor form, every n-photon block of the multiclick combination is bounded below by the bound function at n, built from extreme eigenvalues of the mode operators B and B\ensuremath{_{k}}. The bound is at most one.}
\claimbookmark{receiver/thm:multiclick-bound}{Theorem}{thm:multiclick-bound}{Lower bound on every multiclick block}
\claimleadin{%
The multiclick element is the identity plus a multiple of the all-dark
element, the \termref{term-joint-no-click-elements}{joint no-click element} of all
$K$ ports, minus the joint no-click elements that leave out one port. Bounding
each lifted operator by an extreme eigenvalue of its mode-space operator,
from below where it enters with a non-negative coefficient and from above
where it is subtracted, gives a scalar bound on each photon-number block.
}
\begin{theorem}[Lower bound on every multiclick block]\label{thm:multiclick-bound}
Let $B,B_1,\dots,B_K$ be Hermitian operators on a finite-dimensional space with
\begin{equation}\label{eq:bbkhyp}
  0\ \preceq\ B_k\ \preceq\ B\ \preceq\ \id\qquad(k=1,\dots,K),
\end{equation}
let $\gamma,\beta_1,\dots,\beta_K\in(0,1]$ with $\gamma\le\beta_k$ for every
$k$, let $\Pvac{[K]}$ and $\Pvac{[K]\setminus\{k\}}$ ($k=1,\dots,K$) be
operators on the \termref{term-fock-space}{Fock space} of that space of the form
\begin{equation}\label{eq:prefactorform}
  \Pvac{[K]}=\gamma\,\Fock(\id-B),\qquad
  \Pvac{[K]\setminus\{k\}}=\beta_k\,\Fock\big(\id-B_k\big),
\end{equation}
and put $M_{\mc}\coloneqq\id+(K-1)\Pvac{[K]}-\sum_k\Pvac{[K]\setminus\{k\}}$. Put
\begin{equation}\label{eq:cd}
  c\coloneqq1-\lmax(B),\qquad
  d_k\coloneqq1-\lmin\big(B_k\big).
\end{equation}
Then $\big(\gamma,\{\beta_k\},c,\{d_k\}\big)$ is admissible bound data, and for
every $n\ge0$
\begin{equation}\label{eq:mcbound}
  \lmin\big(\blk{M}{\eqn}_{\mc}\big)\ \ge\ \Bnd(n) ,
  \qquad\text{equivalently}\qquad
  \blk{M}{\eqn}_{\mc}\ \succeq\ \Bnd(n)\,\id^{(=n)} .
\end{equation}
Moreover $\Bnd(n)\le1$ for every $n$.
\end{theorem}
\claimmeta{\metalabel{Uses}\mref{Lem.}{lem:fock-block-spectrum}, \mref{Def.}{def:bounddata}
\metasep \metalabel{Used by}\mref{Rem.}{rem:kamin}, \mref{Def.}{def:exactbnd}, \mref{Lem.}{lem:exactbasic}, \mref{Rem.}{rem:exactpositivity}, \mref{Lem.}{lem:exactwindow}, \mref{Def.}{def:setting-family}, \mref{Cor.}{cor:unconditioned-family}, \mref{Rem.}{rem:exacteig}, \mref{Thm.}{thm:enclosed-weight}, \mref{Thm.}{thm:emitted-bound}, \mref{Def.}{def:tolmap}, \mref{Rem.}{rem:bound-data-band}, \mref{Rem.}{rem:tolerancedomain}, \mref{Thm.}{thm:exact-and-evaluated}, \mref{Rem.}{rem:ia}, \mref{Req.}{req:rigorous-weight}, \mref{Ex.}{ex:bb84exact}, \mref{Prop.}{prop:noise-floor}
\metasep newly derived
\metasep implemented (effective: accepted)
\metasep \metalabel{Code}\coderef{https://github.com/fraunhofersg/certiqs-key/blob/v0.1.0/src/certiqs_key/receiver/bound.py}{certiqs_key.receiver.bound:setting_bound}
\metasep \metalabel{Oracles}\coderef{https://github.com/fraunhofersg/certiqs-key/blob/v0.1.0/tests/oracles/registry_receiver_povm.py}{registry_receiver_povm.py}}

\begin{proof}
\emph{Admissibility of the data.} $0\preceq B\preceq\id$ gives
$\lmax(B)\in[0,1]$, hence $c\in[0,1]$; likewise $0\preceq B_k\preceq\id$ gives
$d_k\in[0,1]$. For $c\le d_k$: from $B_k\preceq B$ we get
$\lmin(B_k)\le\lmin(B)\le\lmax(B)$, and subtracting from $1$ reverses the
inequality, so $d_k=1-\lmin(B_k)\ge1-\lmax(B)=c$. The relation
$\gamma\le\beta_k$ is a hypothesis.

\emph{The bound.} By Lemma~\ref{lem:fock-block-spectrum} applied to $B$ and to each
$B_k$ --- legitimate because \eqref{eq:bbkhyp} puts all of them in
$[0,\id]$ ---
\begin{equation}\label{eq:blockopbounds}
  \Fock(\id-B)^{(n)}\succeq c^{\,n}\id^{(=n)},
  \qquad
  \Fock\big(\id-B_k\big)^{(n)}\preceq d_k^{\,n}\id^{(=n)} .
\end{equation}
The Löwner order is preserved by sums and by multiplication with non-negative
scalars, and $K-1\ge0$, $\gamma>0$, $\beta_k>0$. Taking the $n$-photon block of
$M_{\mc}=\id+(K-1)\Pvac{[K]}-\sum_k\Pvac{[K]\setminus\{k\}}$, bounding the
term with the non-negative coefficient from below and each subtracted term from
above, gives
$\blk{M}{\eqn}_{\mc}\succeq\big(1+(K-1)\gamma c^n-\sum_k\beta_kd_k^n\big)
\id^{(=n)}$, and $X\succeq t\,\id$ says exactly $\lmin(X)\ge t$.

\emph{$\Bnd\le1$.} Since $d_k\ge c$ and $\beta_k\ge\gamma$, every term of
$\sum_k\beta_kd_k^n$ is at least $\gamma c^n$, so the sum is at least
$K\gamma c^n\ge(K-1)\gamma c^n$, whence $\Bnd(n)\le1$.
\end{proof}

% claim receiver/rem:k-minus-1
\claimsummary{receiver/rem:k-minus-1}{The all-dark element (no port clicks) enters the multiclick element with coefficient K\ensuremath{-}1: each of the K single-click elements contributes one, and the no-click element, which is this same operator, takes one away. An error by one would show only where c \textgreater{} 0.}
\claimbookmark{receiver/rem:k-minus-1}{Remark}{rem:k-minus-1}{The coefficient K-1 of the all-dark element}
\begin{remark}[The coefficient $K-1$ of the all-dark element]\label{rem:k-minus-1}
The multiclick element is the complement of the no-click and the $K$
single-click elements, and each single-click element already contains one
subtraction of the all-dark element \eqref{eq:Msc}. The $K$ subtractions and
the one from $M_{\varnothing}$ combine into the coefficient $K-1$
(Theorem~\ref{thm:click-povm}\,(iv)). An error by one in this coefficient changes
the bound by $\gamma c^{n}$, a term that vanishes on every block $n\ge1$ when
$c=0$ and survives when $c>0$: the error would be invisible in the first case
and material in the second.
\end{remark}
\claimmeta{\metalabel{Uses}\mref{Thm.}{thm:click-povm}
\metasep newly derived
\metasep accepted}

% claim receiver/rem:kamin
\claimsummary{receiver/rem:kamin}{A published construction bounds a photon-number subspace weight from an observed multiclick rate, for one row per detector. \hyperlink{claim-receiver/thm:multiclick-bound}{Theorem~\ref*{thm:multiclick-bound}} is instead an operator inequality for every input state and ports of any rank; no step here depends on the citation.}
\claimbookmark{receiver/rem:kamin}{Remark}{rem:kamin}{Relation to a published multiclick weight bound}
\begin{remark}[Relation to a published multiclick weight bound]\label{rem:kamin}
A published route to a related quantity is
\cite[Thm.~1, Eq.~(58)]{kamin2024improved}, which writes a multiclick POVM
element as $M^{\mathrm{mult}}(G)=\id-\dyad{0}-\sum_iF_{\mathrm{sic}}$ for a
passive linear-optical setup with mode transformation $G$, and then bounds the
weight of a photon-number subspace using an \emph{observed} multiclick rate
together with constants obtained by partitioning the rows of $G$. That is a
different object from \eqref{eq:mcbound}: it is a statement about a measured
statistic and a subspace weight, whereas Theorem~\ref{thm:multiclick-bound} is an
operator inequality on each photon-number block that holds for every input
state and is used as a device fact rather than as an observation. The
construction is also less general in the direction that matters here: it is
written for a one-row-per-detector transformation, whereas
Theorem~\ref{thm:multiclick-bound} is proved for arbitrary port Grams, hence for ports
of arbitrary rank. The citation is for orientation, and no step of this
contract depends on it.
\end{remark}
\claimmeta{\metalabel{Uses}\mref{Thm.}{thm:multiclick-bound}
\metasep newly derived
\metasep accepted}

% claim receiver/lem:collapse
\claimsummary{receiver/lem:collapse}{If the bound function is at most zero at zero photons, it cannot decrease from any n \ensuremath{\ge} 1 at which it is non-negative. The non-negative part of its infimum above a cutoff is then that of its value at the next photon number.}
\claimbookmark{receiver/lem:collapse}{Lemma}{lem:collapse}{Collapse of the tail infimum to the first photon number}
\claimleadin{%
A squashing argument consumes not $\Bnd(N+1)$ but the infimum
$\inf_{n>N}\lmin\big(\blk{M}{\eqn}_{\mc}\big)$, and a squashing request whose
bounds are not uniform over \emph{all} photon numbers above the cutoff
fails (\href{../../../spec/architecture.md\#54-failure-policy}{\nolinkurl{spec/architecture.md}\texttt{\#54-failure-policy}}). The
\termref{term-bound-function}{bound function} is not monotone in general: it is the sum of a
non-increasing term $(K-1)\gamma c^n$ and a non-decreasing term
$-\sum_k\beta_kd_k^n$. The lemma shows that the non-monotone regime is
confined to where the bound is useless anyway.
}
\begin{lemma}[Collapse of the tail infimum to the first photon number]\label{lem:collapse}
Let $\big(\gamma,\{\beta_k\},c,\{d_k\}\big)$ be admissible bound data
satisfying
\begin{equation}\label{eq:collapsehyp}
  \sum_{k=1}^{K}\beta_k\ \ge\ 1+(K-1)\gamma
  \qquad\text{(equivalently, }\Bnd(0)\le0\text{)}.
\end{equation}
Then for every $n\ge1$,
\begin{equation}\label{eq:collapsestep}
  \Bnd(n)\ge0\ \Longrightarrow\ \Bnd(n+1)\ \ge\ \Bnd(n) .
\end{equation}
Consequently, for every cutoff $N\ge0$,
\begin{equation}\label{eq:collapseinf}
  \max\Big\{0,\ \inf_{n>N}\Bnd(n)\Big\}\;=\;\max\big\{0,\ \Bnd(N+1)\big\} ,
\end{equation}
and if $\Bnd(N+1)\ge0$ then $\Bnd$ is non-decreasing on $\{N+1,N+2,\dots\}$ and
$\inf_{n>N}\Bnd(n)=\Bnd(N+1)$. Hypothesis \eqref{eq:collapsehyp} holds
automatically when $\gamma=\beta_k=1$, where it reads $K\ge K$.
\end{lemma}
\claimmeta{\metalabel{Uses}\mref{Def.}{def:bounddata}
\metasep \metalabel{Used by}\mref{Rem.}{rem:n0}, \mref{Rem.}{rem:collapsesharp}, \mref{Lem.}{lem:windowtail}, \mref{Lem.}{lem:exactwindow}, \mref{Rem.}{rem:exactscope}, \mref{Def.}{def:setting-family}, \mref{Cor.}{cor:unconditioned-family}, \mref{Def.}{def:pmc}, \mref{Rem.}{rem:pmc-branch-choice}, \mref{Rem.}{rem:noslack}, \mref{Def.}{def:pmcex}, \mref{Rem.}{rem:exact-window-remainders}, \mref{Prop.}{prop:exactnoworse}, \mref{Rem.}{rem:noworsefl}, \mref{Thm.}{thm:enclosed-weight}, \mref{Thm.}{thm:emitted-bound}, \mref{Rem.}{rem:notsharp}, \mref{Thm.}{thm:exact-and-evaluated}, \mref{Rem.}{rem:ia}, \mref{Ex.}{ex:bb84bound}
\metasep newly derived
\metasep implemented (effective: accepted)
\metasep \metalabel{Code}\coderef{https://github.com/fraunhofersg/certiqs-key/blob/v0.1.0/src/certiqs_key/receiver/bound.py}{certiqs_key.receiver.bound:collapse_slack}}

\begin{proof}
\emph{Idea.} The increment $\Bnd(n+1)-\Bnd(n)$ is non-negative exactly when
a weighted sum over the ports dominates the all-dark term. After the
substitution $u_k=\beta_kd_k^{\,n}$ each port's contribution is a concave
function of $u_k$, so it lies above its chord; the constraint
$\Bnd(n)\ge0$ bounds $\sum_ku_k$, and the hypothesis at zero photons closes
the gap. An induction then turns the step into the statement about the
infimum.

\emph{Step 1 (reduction to one inequality).} Fix $n\ge1$ and assume
$\Bnd(n)\ge0$, that is
\begin{equation}\label{eq:constraint}
  \sum_{k=1}^{K}\beta_kd_k^{\,n}\ \le\ T\coloneqq1+(K-1)\gamma c^{\,n} .
\end{equation}
Since
$\Bnd(n+1)-\Bnd(n)=\sum_k\beta_kd_k^{\,n}(1-d_k)-(K-1)\gamma c^{\,n}(1-c)$,
the implication \eqref{eq:collapsestep} is exactly the statement
\begin{equation}\label{eq:target}
  \sum_{k=1}^{K}\beta_kd_k^{\,n}(1-d_k)\ \ge\ (K-1)\gamma\,c^{\,n}(1-c)
\end{equation}
under the constraint \eqref{eq:constraint} and $c\le d_k\le1$.

\emph{Step 2 (the case $c=1$).} If $c=1$ then $d_k=1$ for every $k$ and both
sides of \eqref{eq:target} vanish. Assume therefore $c<1$, so that
$c^{\,n}<1$ for $n\ge1$.

\emph{Step 3 (the chord bound).} Substitute $u_k\coloneqq\beta_kd_k^{\,n}$, a strictly increasing bijection from
$d_k\in[c,1]$ onto $u_k\in[\beta_kc^{\,n},\beta_k]$ (for $\beta_k>0$; a
$\beta_k=0$ contributes $0$ to both sides of the chord bound below and may be
dropped). Under it
\begin{equation}\label{eq:psi}
  \beta_kd_k^{\,n}(1-d_k)=\psi_k(u_k),\qquad
  \psi_k(u)\coloneqq u-\beta_k^{-1/n}\,u^{1+1/n},
\end{equation}
because $\beta_k^{-1/n}u^{1+1/n}=\beta_k^{-1/n}
(\beta_kd_k^{\,n})^{1+1/n}=\beta_kd_k^{\,n+1}$. Two values and one convexity
fact are all that is needed:
\begin{equation}\label{eq:psivalues}
  \psi_k(\beta_k)=0,\qquad
  \psi_k(\beta_kc^{\,n})=\beta_kc^{\,n}(1-c),\qquad
  \psi_k''(u)=-\beta_k^{-1/n}\Big(1+\tfrac1n\Big)\tfrac1n\,u^{1/n-1}\le0 ,
\end{equation}
so $\psi_k$ is concave on $[0,\beta_k]$ and therefore lies \emph{above} the
chord joining its values at the two endpoints of
$[\beta_kc^{\,n},\beta_k]$. Writing a point of that interval as
$u=\lambda\beta_kc^{\,n}+(1-\lambda)\beta_k$ with
$\lambda=(\beta_k-u)/\big(\beta_k(1-c^{\,n})\big)\in[0,1]$, the chord bound is
\begin{equation}\label{eq:chord}
  \psi_k(u)\ \ge\ \lambda\,\psi_k(\beta_kc^{\,n})+(1-\lambda)\,\psi_k(\beta_k)
  \ =\ \frac{c^{\,n}(1-c)}{1-c^{\,n}}\,\big(\beta_k-u\big).
\end{equation}
Summing \eqref{eq:chord} over $k$ and using $\sum_ku_k\le T$ from
\eqref{eq:constraint} together with $\frac{c^{\,n}(1-c)}{1-c^{\,n}}\ge0$,
\begin{equation}\label{eq:sumchord}
  \sum_{k=1}^{K}\beta_kd_k^{\,n}(1-d_k)
  =\sum_{k=1}^{K}\psi_k(u_k)
  \ \ge\ \frac{c^{\,n}(1-c)}{1-c^{\,n}}
    \Big(\sum_{k}\beta_k-\sum_ku_k\Big)
  \ \ge\ \frac{c^{\,n}(1-c)}{1-c^{\,n}}
    \Big(\sum_{k}\beta_k-T\Big).
\end{equation}

\emph{Step 4 (the hypothesis at zero photons).} Hypothesis
\eqref{eq:collapsehyp} gives
\begin{equation}\label{eq:gap}
  \sum_k\beta_k-T\ \ge\ \big(1+(K-1)\gamma\big)-1-(K-1)\gamma c^{\,n}
  \ =\ (K-1)\gamma\big(1-c^{\,n}\big)\ \ge\ 0 ,
\end{equation}
and substituting \eqref{eq:gap} into \eqref{eq:sumchord} cancels the factor
$1-c^{\,n}>0$ and yields exactly \eqref{eq:target}.

\emph{Step 5 (from \eqref{eq:collapsestep} to \eqref{eq:collapseinf}).} The induction
runs over $n\ge N+1\ge1$, which is where \eqref{eq:collapsestep} was proved. If
$\Bnd(N+1)\ge0$ then, inductively, $\Bnd(n+1)\ge\Bnd(n)\ge\Bnd(N+1)\ge0$ for
every $n\ge N+1$, so $\Bnd$ is non-decreasing on $\{N+1,N+2,\dots\}$ and
$\inf_{n>N}\Bnd(n)=\Bnd(N+1)$, which gives \eqref{eq:collapseinf}. If
$\Bnd(N+1)<0$ then $\inf_{n>N}\Bnd(n)\le\Bnd(N+1)<0$ and both sides of
\eqref{eq:collapseinf} equal $0$.
\end{proof}

% claim receiver/rem:n0
\claimsummary{receiver/rem:n0}{The monotonicity step of \hyperlink{claim-receiver/lem:collapse}{Lemma~\ref*{lem:collapse}} fails at n = 0, with an explicit two-port receiver as witness. Nothing is lost, because the infimum runs over photon numbers above the cutoff, all at least one.}
\claimbookmark{receiver/rem:n0}{Remark}{rem:n0}{The monotonicity step at zero photons}
\begin{remark}[The monotonicity step at zero photons]
\label{rem:n0}
Lemma~\ref{lem:collapse} is stated from $n\ge1$ and cannot be extended to
$n=0$. Take $K\ge2$, $\gamma=\beta_k=1$ and $d_k=1$ for every $k$, and $c<1$: then
$\Bnd(0)=1+(K-1)-K=0\ge0$, while
$\Bnd(1)=1+(K-1)c-K=(K-1)(c-1)<0$, so \eqref{eq:collapsestep} fails at $n=0$.
The data are realised, for instance, by $K=2$ on two input modes with
$A_1=a\dyad{0}$, $A_2=a\dyad{1}$, $a\in(0,1)$: there $c=1-a$ and
$d_1=d_2=1$. The substitution $u_k=\beta_kd_k^{\,n}$ degenerates at $n=0$,
which is precisely where the proof breaks. Nothing is lost: the infimum in
\eqref{eq:collapseinf} runs over $n>N\ge0$, so the induction never starts below
$n=1$.
\end{remark}
\claimmeta{\metalabel{Uses}\mref{Lem.}{lem:collapse}
\metasep newly derived
\metasep accepted}

% claim receiver/rem:collapsesharp
\claimsummary{receiver/rem:collapsesharp}{At the data \ensuremath{\gamma} = \ensuremath{\beta}\ensuremath{_{k}} = 1, d\ensuremath{_{1}} = 1 and d\ensuremath{_{k}} = c for k \ensuremath{\ge} 2, the bound function is identically zero, so the monotonicity step of \hyperlink{claim-receiver/lem:collapse}{Lemma~\ref*{lem:collapse}} cannot be strengthened to strict increase.}
\claimbookmark{receiver/rem:collapsesharp}{Remark}{rem:collapsesharp}{Sharpness of the collapse lemma}
\begin{remark}[Sharpness of the collapse lemma]\label{rem:collapsesharp}
Equality holds in \eqref{eq:target} at the data $\gamma=\beta_k=1$,
$d_1=1$, $d_2=\dots=d_K=c$: there $\Bnd(n)=1+(K-1)c^{\,n}-1-(K-1)c^{\,n}=0$ for
every $n$, so $\Bnd$ is constant and the implication
\eqref{eq:collapsestep} cannot be strengthened to strict increase. More
generally, at $\gamma=\beta_k=1$ every $d_1=1$ gives $\Bnd(n)\le0$ for every
$n$, because $d_k\ge c$ by \eqref{eq:bounddata} makes
$\Bnd(n)\le1+(K-1)c^{\,n}-1-(K-1)c^{\,n}$. Physically, $d_1=1$ is a singular
$B_1$, and at a receiver without a probe, where $B_1=A_{[K]\setminus\{1\}}$
(Corollary~\ref{cor:unconditioned-family}), that is a receiver in which some
input polarization reaches no port other than port~$1$: no multiclick weight
can be guaranteed at any photon number, and the bound correctly returns $0$.
The data above attain equality; at a receiver without a probe they force
ports $2$ to $K$ to receive nothing, since $d_k=c$ means
$\lmin(B_k)=\lmax(B)$, which with $B_k\preceq B$ forces $B_k=B$, that is
$A_k=0$.
\end{remark}
\claimmeta{\metalabel{Uses}\mref{Def.}{def:bounddata}, \mref{Lem.}{lem:collapse}, \mref{Cor.}{cor:unconditioned-family}
\metasep \metalabel{Used by}\mref{Rem.}{rem:ia}
\metasep newly derived
\metasep accepted}

% claim receiver/lem:windowtail
\claimsummary{receiver/lem:windowtail}{The minimum of the bound function over a finite window, capped by an analytic tail term, is a valid uniform bound above the cutoff. It improves as the window widens, converges to the infimum when c \textless{} 1, and a residual gap bounds what remains.}
\claimbookmark{receiver/lem:windowtail}{Lemma}{lem:windowtail}{Window-and-tail bound on the tail infimum}
\claimleadin{%
Where the hypothesis of Lemma~\ref{lem:collapse} fails, the infimum of the
\termref{term-bound-function}{bound function} over all photon numbers above the cutoff is still
bounded from below by a finite computation, which converges to it when
$c<1$: evaluate $\Bnd$ on a finite window and bound the rest by a term that
drops the non-negative all-dark contribution.
}
\begin{lemma}[Window-and-tail bound on the tail infimum]\label{lem:windowtail}
Let $\big(\gamma,\{\beta_k\},c,\{d_k\}\big)$ be admissible bound data and
$N\ge0$. Define, for an integer $m>N$,
\begin{equation}\label{eq:windowtail}
  \underline{\Bnd}_{N,m}\coloneqq
  \min\big\{\,W_{N,m},\ T_m\,\big\},\qquad
  W_{N,m}\coloneqq\min_{N<n\le m}\Bnd(n),\qquad
  T_m\coloneqq1-\sum_{k=1}^{K}\beta_kd_k^{\,m} .
\end{equation}
The window $\{n:N<n\le m\}$ is non-empty for every $m>N$; at $m=N+1$ it is the
single index $N+1$. Then:
\begin{enumerate}[label=(\roman*),leftmargin=3.0em]
\item $\underline{\Bnd}_{N,m}\le\inf_{n>N}\Bnd(n)$ for every $m>N$: the
  quantity is a valid uniform bound;
\item $\underline{\Bnd}_{N,m}$ is non-decreasing in $m$, and
  $\lim_{m\to\infty}\underline{\Bnd}_{N,m}=\inf_{n>N}\Bnd(n)$ whenever
  $c<1$; moreover the \emph{residual gap}
  \begin{equation}\label{eq:resgap}
    g_{N,m}\coloneqq\max\big\{0,\ W_{N,m}-T_m\big\}
  \end{equation}
  is a finite non-negative real for every $m>N$ --- at $m=N+1$ it equals
  $(K-1)\gamma c^{\,N+1}$ --- and satisfies
  $\inf_{n>N}\Bnd(n)-\underline{\Bnd}_{N,m}\le g_{N,m}$, so the evaluation is
  convergent and $g_{N,m}$ bounds what widening the window could still recover;
\item if $c=0$ then $\Bnd(n)=1-\sum_k\beta_kd_k^{\,n}$ for $n\ge1$ is
  non-decreasing, $\underline{\Bnd}_{N,N+1}=\Bnd(N+1)$ is exact, and
  $g_{N,N+1}=0$;
\item the two lemmas are consistent: under the hypothesis of
  Lemma~\ref{lem:collapse} with $\Bnd(N+1)\ge0$ \emph{and $c<1$}, (i) and (ii)
  give $\underline{\Bnd}_{N,m}\le\Bnd(N+1)$ for every $m$ with equality in the
  limit $m\to\infty$, so the first branch of Definition~\ref{def:pmc} is the
  exact value the window-and-tail evaluation converges to, and using it loses
  nothing. The hypothesis $c<1$ is needed: at $c=1$ admissibility forces
  $d_k=1$ for every $k$, $\Bnd$ is constant, and under
  Lemma~\ref{lem:collapse}'s hypothesis with $\Bnd(N+1)\ge0$ that constant is
  $0$ while $T_m=1-\sum_k\beta_k=-(K-1)\gamma$ for every $m$, so the limit
  clause fails whenever $\gamma>0$ and $K\ge2$.
\end{enumerate}
\end{lemma}
\claimmeta{\metalabel{Uses}\mref{Def.}{def:bounddata}, \mref{Lem.}{lem:collapse}
\metasep \metalabel{Used by}\mref{Def.}{def:exactbnd}, \mref{Lem.}{lem:exactwindow}, \mref{Rem.}{rem:exactscope}, \mref{Def.}{def:pmc}, \mref{Rem.}{rem:pmc-branch-choice}, \mref{Rem.}{rem:noslack}, \mref{Def.}{def:pmcex}, \mref{Rem.}{rem:exact-window-remainders}, \mref{Prop.}{prop:exactnoworse}, \mref{Rem.}{rem:exacteig}, \mref{Def.}{def:enclosed-weight}, \mref{Thm.}{thm:enclosed-weight}, \mref{Thm.}{thm:emitted-bound}, \mref{Rem.}{rem:notsharp}, \mref{Def.}{def:tolmap}, \mref{Thm.}{thm:exact-and-evaluated}, \mref{Rem.}{rem:ia}, \mref{Req.}{req:residual-refusal}
\metasep newly derived
\metasep implemented (effective: accepted)
\metasep \metalabel{Code}\coderef{https://github.com/fraunhofersg/certiqs-key/blob/v0.1.0/src/certiqs_key/receiver/bound.py}{certiqs_key.receiver.bound:window_floor}}

\begin{proof}
(i) For $N<n\le m$ the value $\Bnd(n)$ appears in $W_{N,m}$. For $n>m$,
$d_k\le1$ gives $d_k^{\,n}\le d_k^{\,m}$ and $\gamma c^{\,n}\ge0$, so
$\Bnd(n)\ge1-\sum_k\beta_kd_k^{\,n}\ge T_m$. Hence every $\Bnd(n)$ with $n>N$
dominates $\min\{W_{N,m},T_m\}$.

(ii) \emph{Monotonicity.} $W_{N,m+1}=\min\{W_{N,m},\Bnd(m+1)\}$ and
$T_{m+1}\ge T_m$ because $d_k^{\,m+1}\le d_k^{\,m}$; and $\Bnd(m+1)\ge T_m$ by
the computation in (i) at $n=m+1>m$. Each of $W_{N,m}$, $\Bnd(m+1)$ and
$T_{m+1}$ is therefore at least $\min\{W_{N,m},T_m\}=\underline{\Bnd}_{N,m}$,
and $\underline{\Bnd}_{N,m+1}$ is the minimum of those three.

\emph{Limit.} $W_{N,m}$ is non-increasing in $m$ and the windows exhaust
$\{n:n>N\}$, so $W_{N,m}\downarrow\inf_{n>N}\Bnd(n)$. Put
$D\coloneqq\{k:d_k=1\}$; then $T_m\uparrow1-\sum_{k\in D}\beta_k$, and since
$c<1$ also $\Bnd(n)\to1-\sum_{k\in D}\beta_k$, so
$\lim_mT_m\ge\inf_{n>N}\Bnd(n)$. Both sequences converge, hence
$\lim_m\min\{W_{N,m},T_m\}=\min\{\inf_{n>N}\Bnd(n),\,\lim_mT_m\}
=\inf_{n>N}\Bnd(n)$.

\emph{The residual gap.} Both $W_{N,m}$ and $T_m$ are minima of finitely many
reals, so $g_{N,m}$ is a finite non-negative real; at $m=N+1$ the window is
$\{N+1\}$, so $W_{N,N+1}=\Bnd(N+1)=T_{N+1}+(K-1)\gamma c^{\,N+1}$ and
$g_{N,N+1}=(K-1)\gamma c^{\,N+1}$. For the bound, write $W=W_{N,m}$, $T=T_m$.
If $T\le W$ then $\underline{\Bnd}_{N,m}=T$ and $\inf_{n>N}\Bnd(n)\le W$ (the
window values are among those infimized), so the difference is at most
$W-T=g_{N,m}$. If $W<T$ then $\underline{\Bnd}_{N,m}=W=\Bnd(n_0)$ for some
$N<n_0\le m$, so $\inf_{n>N}\Bnd(n)\le\underline{\Bnd}_{N,m}$ and the
difference is at most $0=g_{N,m}$.

(iii) $c=0$ makes $\gamma c^{\,n}=0$ for $n\ge1$, so
$\Bnd(n)=1-\sum_k\beta_kd_k^{\,n}$ there, and each $d_k^{\,n}$ is
non-increasing in $n$; in particular $W_{N,N+1}=\Bnd(N+1)=T_{N+1}$, so
$\underline{\Bnd}_{N,N+1}=\Bnd(N+1)$ and $g_{N,N+1}=0$.

(iv) Under the hypothesis of Lemma~\ref{lem:collapse} with $\Bnd(N+1)\ge0$ and
$c<1$, that lemma gives $\inf_{n>N}\Bnd(n)=\Bnd(N+1)$; then (i) reads
$\underline{\Bnd}_{N,m}\le\Bnd(N+1)$ and (ii) gives
$\underline{\Bnd}_{N,m}\to\Bnd(N+1)$.
\end{proof}

\section{Exact block eigenvalues}\label{group-receiver/def:exactbnd}
\sectionintro{The analytic bound bounds three operators separately and loses slack that no wider window recovers. Computing the smallest eigenvalue of a block itself recovers it: that \termref{term-exact-block-bound}{exact block bound} dominates the analytic bound, is at most one, and is non-negative wherever the \termref{term-multiclick-element}{multiclick element} is positive (\hyperlink{claim-receiver/lem:exactbasic}{Lemma~\ref*{lem:exactbasic}}). Lower estimates on finitely many exact blocks combine with an admissible analytic remainder into a uniform bound over the whole tail (\hyperlink{claim-receiver/lem:exactwindow}{Lemma~\ref*{lem:exactwindow}}).}

% claim receiver/def:exactbnd
\claimsummary{receiver/def:exactbnd}{The exact block bound \ensuremath{\beta}*(n) is the smallest eigenvalue of the n-photon block of the multiclick element. Computing it on finitely many blocks recovers the slack the analytic bound loses by bounding three operators separately.}
\claimbookmark{receiver/def:exactbnd}{Definition}{def:exactbnd}{The exact block bound function}
\claimleadin{%
Both evaluations of the uniform bound so far, the collapse of
Lemma~\ref{lem:collapse} and the window and tail of
Lemma~\ref{lem:windowtail}, bound $\lmin(\blk{M}{\eqn}_{\mc})$ by the
analytic $\Bnd(n)$, and \eqref{eq:mcbound} is where all of the slack enters:
Theorem~\ref{thm:multiclick-bound} bounds three operators separately. Widening the
analytic window does not recover it: at $c=0$, for instance at a receiver
with no probe whose instrument sends every input entirely into
\termref{term-detector-ports}{detector ports} ($A_\Sigma=\id$), Lemma~\ref{lem:windowtail}\,(iii) makes the residual
gap exactly zero. What recovers it is computing the block itself, on finitely
many indices. Lemma~\ref{lem:exactbasic} and Lemma~\ref{lem:exactwindow} say
what may be computed, what must still be argued, and how the two compose.
}
\begin{definition}[The exact block bound function]\label{def:exactbnd}
Under the hypotheses of Theorem~\ref{thm:multiclick-bound}, the
\termdef{term-exact-block-bound}{exact block bound}
is
\begin{equation}\label{eq:exactbnd}
  \Bnd^{\ast}(n)\ \coloneqq\ \lmin\big(\blk{M}{\eqn}_{\mc}\big)
  \ =\ \lmin\Big(\id^{(=n)}+(K-1)\gamma\,\Fock(\id-B)^{(n)}
   -\sum_{k=1}^{K}\beta_k\,\Fock\big(\id-B_k\big)^{(n)}\Big),
\end{equation}
for $n\ge0$, the second equality being \eqref{eq:prefactorform} restricted to
the $n$-photon block. The operator inside is a Hermitian matrix of dimension
$\binom{D+n-1}{n}$, where $D$ is the dimension of the space $B$ acts on:
the \termref{term-channel-input-space}{channel input space} of the \termref{term-setting-family}{setting family}
(Definition~\ref{def:setting-family}), which is $|\InpC|$ at both instances
this module constructs. No new object is introduced: each
$\Fock(\cdot)^{(n)}=\big((\cdot)^{\otimes n}\big)\big|_{\Symn{n}}$ is the
construction \eqref{eq:blockMS} already uses for the blocks the artifact
emits at $n\le N$, in the occupation basis and order that
Definition~\ref{def:fock} and the interface conventions of this module fix,
evaluated at an index above the cutoff instead of below it.
\end{definition}
\claimmeta{\metalabel{Uses}\mref{Def.}{def:fock}, \mref{Thm.}{thm:click-povm}, \mref{Thm.}{thm:multiclick-bound}, \mref{Lem.}{lem:windowtail}
\metasep \metalabel{Used by}\mref{Lem.}{lem:exactbasic}, \mref{Def.}{def:pmcex}, \mref{Def.}{def:enclosed-weight}, \mref{Thm.}{thm:enclosed-weight}, \mref{Thm.}{thm:emitted-bound}, \mref{Rem.}{rem:notsharp}, \mref{Ass.}{ass:eig}, \mref{Req.}{req:rigorous-weight}, \mref{Ex.}{ex:bb84exact}
\metasep newly derived
\metasep implemented (effective: accepted)
\metasep \metalabel{Code}\coderef{https://github.com/fraunhofersg/certiqs-key/blob/v0.1.0/src/certiqs_key/receiver/bound.py}{certiqs_key.receiver.bound:exact_block_bound}}

% claim receiver/lem:exactbasic
\claimsummary{receiver/lem:exactbasic}{The exact block bound lies between the analytic bound and one on every photon-number block, and it is non-negative when the multiclick element is positive. That positivity holds at every setting of a setting family but does not follow from the multiclick bound's hypotheses.}
\claimbookmark{receiver/lem:exactbasic}{Lemma}{lem:exactbasic}{Range of the exact block bound}
\claimleadin{%
The three bounds of the lemma need different hypotheses. The comparison with
the analytic bound and the upper bound hold under the hypotheses of
Theorem~\ref{thm:multiclick-bound} alone. Non-negativity needs the
\termref{term-multiclick-element}{multiclick element} itself to be positive, \eqref{eq:mcpos}, which holds at every setting
of a \termref{term-setting-family}{setting family} (Definition~\ref{def:setting-family}\,(a)), because
there $M_{\mc}$ is a sum of elements of a POVM; those are the only instances
to which this module applies Theorem~\ref{thm:multiclick-bound}. It does not follow
from \eqref{eq:bbkhyp} and \eqref{eq:prefactorform}, and
Remark~\ref{rem:exactpositivity} is the instance that shows so.
}
\begin{lemma}[Range of the exact block bound]
\label{lem:exactbasic}
Let the hypotheses of Theorem~\ref{thm:multiclick-bound} hold. Then, for every
$n\ge0$,
\begin{equation}\label{eq:exactbasic}
  \Bnd(n)\ \le\ \Bnd^{\ast}(n)\ \le\ 1 .
\end{equation}
If in addition
\begin{equation}\label{eq:mcpos}
  M_{\mc}\ \succeq\ 0 ,
\end{equation}
then $\Bnd^{\ast}(n)\ge0$ as well, so that $\Bnd^{\ast}(n)\in[0,1]$.
\end{lemma}
\claimmeta{\metalabel{Uses}\mref{Lem.}{lem:fock-functor}, \mref{Lem.}{lem:fock-block-spectrum}, \mref{Lem.}{lem:fock-monotone}, \mref{Thm.}{thm:multiclick-bound}, \mref{Def.}{def:exactbnd}
\metasep \metalabel{Used by}\mref{Rem.}{rem:exactpositivity}, \mref{Lem.}{lem:exactwindow}, \mref{Rem.}{rem:exactscope}, \mref{Def.}{def:setting-family}, \mref{Rem.}{rem:exacteig}, \mref{Thm.}{thm:enclosed-weight}, \mref{Thm.}{thm:emitted-bound}, \mref{Def.}{def:residuals}, \mref{Rem.}{rem:tolerancedomain}, \mref{Thm.}{thm:exact-and-evaluated}, \mref{Rem.}{rem:ia}
\metasep newly derived
\metasep accepted}

\begin{proof}
\emph{Lower comparison.} $\Bnd(n)\le\Bnd^{\ast}(n)$ is \eqref{eq:mcbound}.

\emph{Upper bound.} For $\Bnd^{\ast}(n)\le1$ it suffices that
$M_{\mc}\preceq\id$, which is the operator form of the argument that closes
the proof of Theorem~\ref{thm:multiclick-bound}. From $B_k\preceq B$,
Lemma~\ref{lem:fock-monotone} gives $\Fock(\id-B_k)\succeq\Fock(\id-B)$;
multiplying by $\beta_k>0$, summing, and then using $\beta_k\ge\gamma$
together with $\Fock(\id-B)\succeq0$ (Lemma~\ref{lem:fock-block-spectrum}),
\begin{equation}\label{eq:mcupper}
  \sum_{k=1}^{K}\beta_k\Fock\big(\id-B_k\big)
  \ \succeq\ \Big(\sum_{k=1}^{K}\beta_k\Big)\Fock(\id-B)
  \ \succeq\ K\gamma\,\Fock(\id-B)
  \ \succeq\ (K-1)\gamma\,\Fock(\id-B) ,
\end{equation}
so \eqref{eq:prefactorform} gives $M_{\mc}=\id+(K-1)\gamma\Fock(\id-B)
-\sum_k\beta_k\Fock(\id-B_k)\preceq\id$. Every summand is photon-number block
diagonal (Lemma~\ref{lem:fock-functor}\,(i)), so the $n$-photon block of
$M_{\mc}$ is its compression by $\Pi_n$ and inherits the bound:
$\blk{M}{\eqn}_{\mc}\preceq\id^{(=n)}$, whence $\Bnd^{\ast}(n)\le1$.

\emph{Non-negativity.} Under \eqref{eq:mcpos} the same compression inherits
positivity, since
$\bra{\psi}\Pi_nM_{\mc}\Pi_n\ket{\psi}=\bra{\Pi_n\psi}M_{\mc}\ket{\Pi_n\psi}
\ge0$ for every $\ket\psi$; so $\blk{M}{\eqn}_{\mc}\succeq0$ and
$\Bnd^{\ast}(n)\ge0$.
\end{proof}

% claim receiver/rem:exactpositivity
\claimsummary{receiver/rem:exactpositivity}{Data satisfying every hypothesis of the multiclick bound can display a multiclick element with exact block bound \ensuremath{-}3, so non-negativity of the exact block bound needs positivity as its own hypothesis. No receiver of this contract produces such an operator.}
\claimbookmark{receiver/rem:exactpositivity}{Remark}{rem:exactpositivity}{Positivity of the multiclick element as a separate hypothesis}
\begin{remark}[Positivity of the multiclick element as a separate hypothesis]\label{rem:exactpositivity}
Hypothesis \eqref{eq:mcpos} does not follow from \eqref{eq:bbkhyp} and
\eqref{eq:prefactorform}, and the failure is not delicate. Take $K=4$,
$\gamma=\beta_k=1$, $B=\id$ and $B_k=0$: then $0\preceq B_k\preceq B\preceq\id$
and $\gamma\le\beta_k\le1$, so every hypothesis of
Theorem~\ref{thm:multiclick-bound} holds, while the operator $M_{\mc}$ that
the theorem forms from the operators \eqref{eq:prefactorform} is
$M_{\mc}=\id+3\,\Fock(0)-4\,\Fock(\id)$, which on every block $n\ge1$ is
$-3\,\id^{(=n)}$ because $\Fock(0)^{(n)}=0$ there. So $\Bnd^{\ast}(n)=-3<0$.

Theorem~\ref{thm:multiclick-bound} is untouched by the instance: the same data give $c=1-\lmax(B)=0$ and $d_k=1-\lmin(B_k)=1$, hence
$\Bnd(n)=1+0-4=-3$, and \eqref{eq:mcbound} holds with equality. A
\emph{lower} bound on $\lmin$ needs no positivity; the statement
$\Bnd^{\ast}\ge0$ is about the operator itself and does. What the instance
excludes is the operator, not the numbers: no receiver of this contract
produces it, because there the left-hand sides of \eqref{eq:prefactorform}
are the \termref{term-joint-no-click-elements}{joint no-click elements} of a POVM and $M_{\mc}$ is a sum of
that POVM's elements. Stating \eqref{eq:mcpos} explicitly, rather than
reading it off the notation \eqref{eq:prefactorform} uses, is what keeps
Lemma~\ref{lem:exactbasic} and Lemma~\ref{lem:exactwindow} true as printed;
Theorem~\ref{thm:emitted-bound} discharges it at each of the two instances, and no
emitted value depends on the difference.
\end{remark}
\claimmeta{\metalabel{Uses}\mref{Thm.}{thm:multiclick-bound}, \mref{Lem.}{lem:exactbasic}
\metasep newly derived
\metasep accepted}

% claim receiver/lem:exactwindow
\claimsummary{receiver/lem:exactwindow}{Lower estimates of the exact block bound on finitely many blocks above the cutoff, combined with an admissible remainder for the blocks beyond them, give a weight in [0,1] that bounds the multiclick element on every block above the cutoff. Three admissible choices are named.}
\claimbookmark{receiver/lem:exactwindow}{Lemma}{lem:exactwindow}{Exact window with an analytic remainder}
\claimleadin{%
The lemma turns information about finitely many blocks into a bound uniform
over all blocks above the cutoff: a lower estimate on each exact block in a
window, and a remainder that covers every block beyond it. The three choices
listed at its end are all this contract uses.
}
\begin{lemma}[Exact window with an analytic remainder]\label{lem:exactwindow}
Let the hypotheses of Theorem~\ref{thm:multiclick-bound} hold together with
\eqref{eq:mcpos}, let $N\ge0$, and let
$m^{\ast}>N$ be an integer. Let $\ell(n)$, $N<n\le m^{\ast}$, be reals with
\begin{equation}\label{eq:ellhyp}
  \ell(n)\ \le\ \Bnd^{\ast}(n)\qquad\text{for }N<n\le m^{\ast},
\end{equation}
and let $\tau$ be a real with
\begin{equation}\label{eq:tailadm}
  \tau\le0
  \qquad\text{or}\qquad
  \tau\le\inf_{n>m^{\ast}}\Bnd^{\ast}(n).
\end{equation}
Then
\begin{equation}\label{eq:exactcomposite}
  P\ \coloneqq\ \max\Big\{0,\ \min\Big\{\ \min_{N<n\le m^{\ast}}\ell(n),
   \ \ \tau\ \Big\}\Big\}
\end{equation}
satisfies $0\le P\le1$ and
\begin{equation}\label{eq:exactstatement}
  \blk{M}{\gtN}_{\mc}\ \succeq\ P\,\id^{(>N)} .
\end{equation}
Three choices are admissible:
\begin{enumerate}[label=(\roman*),leftmargin=3.0em]
\item $\tau=T_{m^{\ast}}$ of \eqref{eq:windowtail}, formed from the bound
  data $\big(\gamma,\{\beta_k\},c,\{d_k\}\big)$ that
  Theorem~\ref{thm:multiclick-bound} forms by \eqref{eq:cd}, satisfies
  \eqref{eq:tailadm} with no further hypothesis;
\item if the bound data satisfy \eqref{eq:collapsehyp}, then
  $\tau=\Bnd(m^{\ast}+1)$ satisfies \eqref{eq:tailadm};
\item $\ell(n)=\Bnd(n)$ satisfies \eqref{eq:ellhyp}, and the pointwise maximum
  of finitely many admissible choices of $\ell$ is admissible.
\end{enumerate}
\end{lemma}
\claimmeta{\metalabel{Uses}\mref{Thm.}{thm:multiclick-bound}, \mref{Lem.}{lem:collapse}, \mref{Lem.}{lem:windowtail}, \mref{Lem.}{lem:exactbasic}
\metasep \metalabel{Used by}\mref{Rem.}{rem:exactscope}, \mref{Def.}{def:pmcex}, \mref{Rem.}{rem:exact-window-remainders}, \mref{Rem.}{rem:exacteig}, \mref{Thm.}{thm:enclosed-weight}, \mref{Thm.}{thm:emitted-bound}, \mref{Rem.}{rem:notsharp}, \mref{Thm.}{thm:exact-and-evaluated}
\metasep newly derived
\metasep implemented (effective: accepted)
\metasep \metalabel{Code}\coderef{https://github.com/fraunhofersg/certiqs-key/blob/v0.1.0/src/certiqs_key/receiver/bound.py}{certiqs_key.receiver.bound:exact_composite}}

\begin{proof}
\emph{Step 1 (range of $P$).} $0\le P$ is the outer maximum. For $P\le1$:
either $P=0$, or $P\le\ell(N+1)\le\Bnd^{\ast}(N+1)\le1$ by \eqref{eq:ellhyp}
and \eqref{eq:exactbasic}, which needs no more than the hypotheses of
Theorem~\ref{thm:multiclick-bound}, the window being non-empty because
$m^{\ast}>N$.

\emph{Step 2 (the bound \eqref{eq:exactstatement}).} If $P=0$ the assertion is
$\blk{M}{\gtN}_{\mc}\succeq0$, and this is the one step that uses
\eqref{eq:mcpos}: that hypothesis gives $\blk{M}{\eqn}_{\mc}\succeq0$ on every
block (Lemma~\ref{lem:exactbasic}), hence on their direct sum. So let $P>0$.
Then
$P=\min\{\min_{N<n\le m^{\ast}}\ell(n),\ \tau\}$, and in particular
$\tau\ge P>0$, so the first disjunct of \eqref{eq:tailadm} fails and the second
holds. Take any $n>N$. If $n\le m^{\ast}$ then $\Bnd^{\ast}(n)\ge\ell(n)\ge P$
by \eqref{eq:ellhyp}; if $n>m^{\ast}$ then
$\Bnd^{\ast}(n)\ge\inf_{j>m^{\ast}}\Bnd^{\ast}(j)\ge\tau\ge P$. Either way
$\lmin(\blk{M}{\eqn}_{\mc})\ge P$, that is
$\blk{M}{\eqn}_{\mc}\succeq P\,\id^{(=n)}$; the blocks with $n>N$ are
orthogonal direct summands, and a lower bound holding on each holds on the
direct sum.

\emph{Step 3 (choice (i)).} For $n>m^{\ast}$, $d_k\le1$ gives
$d_k^{\,n}\le d_k^{\,m^{\ast}}$ and $\gamma c^{\,n}\ge0$, so
$\Bnd(n)\ge1-\sum_k\beta_kd_k^{\,n}\ge
1-\sum_k\beta_kd_k^{\,m^{\ast}}=T_{m^{\ast}}$, the tail step of
Lemma~\ref{lem:windowtail}\,(i); and $\Bnd^{\ast}(n)\ge\Bnd(n)$ by
Lemma~\ref{lem:exactbasic}. Hence
$\inf_{n>m^{\ast}}\Bnd^{\ast}(n)\ge T_{m^{\ast}}$.

\emph{Step 4 (choice (ii)).} If $\Bnd(m^{\ast}+1)\le0$, the first disjunct of
\eqref{eq:tailadm} holds. Otherwise $\Bnd(m^{\ast}+1)>0$ and
Lemma~\ref{lem:collapse}, applied with cutoff $m^{\ast}$ in place of $N$,
gives $\inf_{n>m^{\ast}}\Bnd(n)=\Bnd(m^{\ast}+1)$; with
$\Bnd^{\ast}\ge\Bnd$ pointwise this is
$\inf_{n>m^{\ast}}\Bnd^{\ast}(n)\ge\Bnd(m^{\ast}+1)$, the second disjunct.

\emph{Step 5 (choice (iii)).} The first statement is
Lemma~\ref{lem:exactbasic}. For the second, a maximum of finitely many
numbers each at most $\Bnd^{\ast}(n)$ is at most $\Bnd^{\ast}(n)$.
\end{proof}

% claim receiver/rem:exactscope
\claimsummary{receiver/rem:exactscope}{The exact-window lemma makes finitely many blocks yield a bound uniform above the cutoff, from remainders already proved, and uses positivity at one step only. It is stated on the exact block bound and makes no computed eigenvalue available by itself.}
\claimbookmark{receiver/rem:exactscope}{Remark}{rem:exactscope}{Scope of the exact-window lemma}
\claimleadin{%
A squashing request fails when a bound does not hold at every photon number
above the cutoff (\href{../../../spec/architecture.md}{\nolinkurl{spec/architecture.md}}~\S5.4), so a bound on
finitely many blocks is of no use there by itself.
}
\begin{remark}[Scope of the exact-window lemma]\label{rem:exactscope}
\emph{What the lemma buys.} Lemma~\ref{lem:exactwindow} converts information
about \emph{finitely many} blocks into a bound uniform over the infinitely
many above the cutoff, which is what a squashing request requires and
what Remark~\ref{rem:notsharp} names as the obstacle to using block eigenvalues at all. The conversion needs no new
analysis: both admissible remainders are results already proved, (i) being
the tail step inside Lemma~\ref{lem:windowtail}\,(i) and (ii) being
Lemma~\ref{lem:collapse} read at a larger cutoff. Some hypotheses are
\emph{not} used: (i) is free of \eqref{eq:collapsehyp}, neither part needs
\termref{term-predicate-x}{predicate~(X)}, which enters only in choosing between
them, and \eqref{eq:mcpos} is used at exactly one step, the $P=0$ branch that
the outer $\max\{0,\cdot\}$ of \eqref{eq:exactcomposite} creates.

\emph{What it does not do.} The lemma does not make an exact eigenvalue
available. It is stated on $\Bnd^{\ast}$, an exact quantity, and
\eqref{eq:ellhyp} is a hypothesis on whatever numbers are supplied for it.
Definition~\ref{def:pmcex} is where those numbers come from, and
Remark~\ref{rem:exacteig} is where the price of computing them is stated.
\end{remark}
\claimmeta{\metalabel{Uses}\mref{Lem.}{lem:collapse}, \mref{Lem.}{lem:windowtail}, \mref{Lem.}{lem:exactbasic}, \mref{Lem.}{lem:exactwindow}
\metasep newly derived
\metasep accepted}

\section{Setting families and the emitted weight}\label{group-receiver/def:setting-family}
\sectionintro{A \termref{term-setting-family}{setting family} packages, for each declared setting, a photon-number block-diagonal POVM whose \termref{term-joint-no-click-elements}{joint no-click elements} have the prefactor form of the multiclick bound; a probe-free receiver is the one-setting family, and a submodule that constructs another family proves there that it is one. Where a depth is declared, a setting's weight is an exact window; elsewhere it is the collapse value where the declaration marks the setting exact and a declared window and tail otherwise; their minimum is the \termref{term-uniform-weight}{uniform weight}. Formed on bounds that lie on the conservative side of the exact data, the same formulas give the \termref{term-enclosed-weight}{enclosed weight}, a bound with no assumption (\hyperlink{claim-receiver/thm:enclosed-weight}{Theorem~\ref*{thm:enclosed-weight}}). The first main result, \hyperlink{claim-receiver/thm:emitted-bound}{Theorem~\ref*{thm:emitted-bound}}, is that each \termref{term-per-setting-weight}{per-setting weight} bounds its setting's \termref{term-multiclick-element}{multiclick element} on every block above the cutoff, and the uniform weight bounds it at every setting, under the eigenvalue-slack assumption only where a computed block eigenvalue enters. The second, \hyperlink{claim-receiver/prop:slotvalid}{Proposition~\ref*{prop:slotvalid}}, is that the \termref{term-strengthening-slot}{strengthening slot} in which the artifact carries the weights has weights that are non-negative and sum to at most one at every setting, as a squashing request requires.}

% claim receiver/def:setting-family
\claimsummary{receiver/def:setting-family}{A setting family gives, per setting, a photon-number block-diagonal POVM whose joint no-click elements have the prefactor form of the multiclick bound, and a flag that reads exact only where every prefactor is one by construction. The high-photon weights are stated on it.}
\claimbookmark{receiver/def:setting-family}{Definition}{def:setting-family}{Setting family}
\claimleadin{%
The weights built on the multiclick bound
(Theorem~\ref{thm:multiclick-bound} onwards) and the residuals of
Definition~\ref{def:residuals} are stated per \emph{setting}: a receiver
without a probe has one setting, a receiver with a probe has one per declared
probe setting, and in both cases the object they are stated on is the same. This definition names that object once, so that the emitted value and
its guarantee are proved for it and each instance is established separately.
The core constructs the unconditioned receiver as an instance
(Corollary~\ref{cor:unconditioned-family}, one setting); a submodule that
constructs another, such as a receiver conditioned on its own probe with one
setting per declared probe setting, proves there that it is one.
}
\begin{definition}[Setting family]\label{def:setting-family}
Let $D\ge1$ and write $\mathcal{F}_D\coloneqq\mathcal{F}(\mathbb{C}^{D})$ for
the \termref{term-fock-space}{Fock space} of Definition~\ref{def:fock} over a $D$-dimensional
\termdef{term-channel-input-space}{channel input space},
with photon-number projectors $\Pi_n$ onto $\Symn{n}(\mathbb{C}^{D})$. A
\termdef{term-setting-family}{setting family}
of $K\ge1$ ports on $\mathcal{F}_D$ is a finite non-empty set $\Tset$ of
\emph{settings} together with, for every $t\in\Tset$:
\begin{enumerate}[label=(\alph*),leftmargin=2.6em]
\item a POVM $\{M_{S|t}\}_{S\subseteq[K]}$ on $\mathcal{F}_D$ every
  element of which commutes with the photon number, so that
  $M_{S|t}=\bigoplus_{n\ge0}\blk{M}{\eqn}_{S|t}$ with $\blk{M}{\eqn}_{S|t}$
  acting on $\Symn{n}(\mathbb{C}^{D})$. Its
  \termdef{term-joint-no-click-elements}{joint no-click elements}
  and its
  \termdef{term-multiclick-element}{multiclick element}
  are
  \begin{equation}\label{eq:family-elements}
    \Pvac{J}^{(t)}\coloneqq\sum_{S\subseteq[K]\setminus J}M_{S|t}
    \quad(J\subseteq[K]),\qquad
    M_{\mc|t}\coloneqq\sum_{S\in\MCK}M_{S|t},
  \end{equation}
  with $\MCK$ the \termref{term-multiclick-set}{multiclick set} of Definition~\ref{def:clickalphabet};
\item Hermitian operators $B^{(t)},B^{(t)}_1,\dots,B^{(t)}_K$ on
  $\mathbb{C}^{D}$ and scalars $\gamma^{(t)},\beta^{(t)}_1,\dots,
  \beta^{(t)}_K\in(0,1]$ satisfying the hypotheses of
  Theorem~\ref{thm:multiclick-bound}, namely \eqref{eq:bbkhyp} and
  $\gamma^{(t)}\le\beta^{(t)}_k$ for every $k$, such that the joint
  no-click elements have the prefactor form \eqref{eq:prefactorform},
  \begin{equation}\label{eq:family-prefactor}
    \Pvac{[K]}^{(t)}=\gamma^{(t)}\,\Fock\big(\id-B^{(t)}\big),\qquad
    \Pvac{[K]\setminus\{k\}}^{(t)}=\beta^{(t)}_k\,\Fock\big(\id-B^{(t)}_k\big);
  \end{equation}
\item a flag $\chi^{(t)}\in\{\mathrm{exact},\mathrm{general}\}$, the
  setting's
  \termdef{term-exact-flag}{exact flag},
  which may be $\mathrm{exact}$ only if $\gamma^{(t)}=\beta^{(t)}_k=1$ for
  every $k$ holds \emph{exactly, by the construction that produced the
  setting}, read off the declaration with no arithmetic performed, so that
  at an exact setting \eqref{eq:collapsehyp} reads $K\ge K$ and holds with
  equality by construction.
\end{enumerate}
The
\termdef{term-bound-data-of-the-setting}{bound data of the setting}
are $\big(\gamma^{(t)},\{\beta^{(t)}_k\},c^{(t)},\{d^{(t)}_k\}\big)$ with
$c^{(t)}\coloneqq1-\lmax(B^{(t)})$ and $d^{(t)}_k\coloneqq1-\lmin(B^{(t)}_k)$
as in \eqref{eq:cd}; they are \termref{term-admissible-bound-data}{admissible bound data}
(Definition~\ref{def:bounddata}) by Theorem~\ref{thm:multiclick-bound}, and
$\Bnd^{(t)}$ denotes their \termref{term-bound-function}{bound function} \eqref{eq:beta}.

Two consequences are used throughout. First, the multiclick element has the
inclusion--exclusion form
\begin{equation}\label{eq:family-mc}
  M_{\mc|t}=\id+(K-1)\,\Pvac{[K]}^{(t)}-\sum_{k=1}^{K}\Pvac{[K]\setminus\{k\}}^{(t)} ,
\end{equation}
because completeness gives $M_{\mc|t}=\id-M_{\varnothing|t}-\sum_kM_{\{k\}|t}$
while \eqref{eq:family-elements} gives $M_{\varnothing|t}=\Pvac{[K]}^{(t)}$
and $M_{\{k\}|t}=\Pvac{[K]\setminus\{k\}}^{(t)}-\Pvac{[K]}^{(t)}$; so
Theorem~\ref{thm:multiclick-bound} applies to every setting, with
$\lmin(\blk{M}{\eqn}_{\mc|t})\ge\Bnd^{(t)}(n)$ for every $n\ge0$. Second,
$M_{\mc|t}$ is a sum of elements of a POVM, so hypothesis \eqref{eq:mcpos}
of Lemma~\ref{lem:exactbasic}, $M_{\mc|t}\succeq0$, holds at every setting.
\end{definition}
\claimmeta{\metalabel{Uses}\mref{Def.}{def:fock}, \mref{Def.}{def:clickalphabet}, \mref{Def.}{def:bounddata}, \mref{Thm.}{thm:multiclick-bound}, \mref{Lem.}{lem:collapse}, \mref{Lem.}{lem:exactbasic}
\metasep \metalabel{Used by}\mref{Cor.}{cor:unconditioned-family}, \mref{Def.}{def:pmc}, \mref{Rem.}{rem:flag-instances}, \mref{Rem.}{rem:noslack}, \mref{Def.}{def:pmcex}, \mref{Def.}{def:enclosed-weight}, \mref{Thm.}{thm:enclosed-weight}, \mref{Thm.}{thm:emitted-bound}, \mref{Def.}{def:slot}, \mref{Prop.}{prop:slotvalid}, \mref{Def.}{def:residuals}, \mref{Prop.}{prop:res5}, \mref{Def.}{def:tolmap}, \mref{Def.}{def:evaluation-model}, \mref{Thm.}{thm:exact-and-evaluated}, \mref{Req.}{req:rigorous-weight}, \mref{Thm.}{thm:dark-counts}, \mref{Prop.}{prop:dark-slack-bound}, \mref{Prop.}{prop:noise-floor}, \extref{../reductions/contract.pdf}{reductions/req:recv-row-provenance}, \extref{../reductions/contract.pdf}{reductions/wpfss/prop:route}, \extref{../reductions/contract.pdf}{reductions/wpfss/rem:receiver-imports}
\metasep newly derived
\metasep accepted}

% claim receiver/cor:unconditioned-family
\claimsummary{receiver/cor:unconditioned-family}{A receiver without probe modes is a one-setting family whose prefactors are all exactly one, with bound data read off its port Grams. Its flag is exact by construction, so the collapse branch applies to it.}
\claimbookmark{receiver/cor:unconditioned-family}{Corollary}{cor:unconditioned-family}{The unconditioned receiver is a setting family}
\begin{corollary}[The unconditioned receiver is a setting family]
\label{cor:unconditioned-family}
Let $\Recv$ be a \termref{term-receiver-instrument}{receiver instrument} (Definition~\ref{def:instrument})
with $\InpS=\varnothing$, so that $\Inp=\InpC$, and let $\{A_k\}$ be its
\termref{term-port-gram-matrices}{port Grams} and $\{M_S\}$ its click-pattern POVM
(Theorem~\ref{thm:click-povm}). Then $\Tset=\{\off\}$ with
\begin{equation}\label{eq:unconditioned-data}
  M_{S|\off}\coloneqq M_S,\qquad
  B^{(\off)}\coloneqq A_\Sigma,\qquad
  B^{(\off)}_k\coloneqq A_{[K]\setminus\{k\}},\qquad
  \gamma^{(\off)}=\beta^{(\off)}_k\coloneqq1,\qquad
  \chi^{(\off)}\coloneqq\mathrm{exact},
\end{equation}
is a setting family (Definition~\ref{def:setting-family}) on
$\mathcal{F}(\mathbb{C}^{\Inp})$, with $D=|\Inp|$. Its bound data are
$c=1-\lmax(A_\Sigma)$ and $d_k=1-\lmin(A_{[K]\setminus\{k\}})$, and
\eqref{eq:collapsehyp} holds with equality.
\end{corollary}
\claimmeta{\metalabel{Uses}\mref{Def.}{def:instrument}, \mref{Def.}{def:gram}, \mref{Thm.}{thm:port-grams}, \mref{Thm.}{thm:click-povm}, \mref{Thm.}{thm:multiclick-bound}, \mref{Lem.}{lem:collapse}, \mref{Def.}{def:setting-family}
\metasep \metalabel{Assumes}\aref{A-TRUST}{ass:trust}, \aref{A-PASSIVE}{ass:passive}, \aref{A-VAC}{ass:vac}, \aref{A-MODES}{ass:modes}, \aref{A-THRESH}{ass:thresh}, \aref{A-PORT}{ass:port}, \aref{A-EXACT}{ass:exact}, \aref{A-NOSEL}{ass:nosel}, \aref{A-ROUND}{ass:round}
\metasep \metalabel{Used by}\mref{Rem.}{rem:collapsesharp}, \mref{Rem.}{rem:flag-instances}, \mref{Cor.}{cor:emitted-unconditioned}, \mref{Def.}{def:residuals}, \mref{Rem.}{rem:bound-data-band}, \mref{Thm.}{thm:exact-and-evaluated}, \mref{Req.}{req:rigorous-weight}, \extref{../reductions/contract.pdf}{reductions/req:recv-row-provenance}, \extref{../reductions/contract.pdf}{reductions/wpfss/prop:route}, \extref{../reductions/contract.pdf}{reductions/wpfss/rem:receiver-imports}
\metasep newly derived
\metasep accepted}

\begin{proof}
The dimension is $D=|\Inp|=|\InpC|\ge1$, since $\InpC\neq\varnothing$ by
Definition~\ref{def:instrument}\,(ii), as Definition~\ref{def:setting-family}
requires. We check the three items of Definition~\ref{def:setting-family}.

(a) $\{M_S\}$ is a POVM by Theorem~\ref{thm:click-povm}\,(i), block diagonal in
the photon number of $\Inp=\InpC$ by Theorem~\ref{thm:click-povm}\,(iii), and by
Theorem~\ref{thm:click-povm}\,(i) its joint no-click elements in the sense of
\eqref{eq:family-elements} are $\Pvac{J}=\sum_{S\subseteq[K]\setminus J}M_S$.

(b) Theorem~\ref{thm:click-povm}\,(ii) gives $\Pvac{J}=\Fock(\id-A_J)$ for every
$J$, which is \eqref{eq:family-prefactor} with $\gamma=\beta_k=1$,
$B=A_\Sigma$ and $B_k=A_{[K]\setminus\{k\}}$. Hypothesis \eqref{eq:bbkhyp}
is $0\preceq A_{[K]\setminus\{k\}}\preceq A_\Sigma\preceq\id$: the first
two inequalities because every $A_k\succeq0$ and $A_\Sigma-A_{[K]\setminus
\{k\}}=A_k$ (Theorem~\ref{thm:port-grams}\,(i), Definition~\ref{def:gram}), the
last by Theorem~\ref{thm:port-grams}\,(i). $\gamma\le\beta_k$ is $1\le1$.

(c) $\gamma=\beta_k=1$ are identities of Theorem~\ref{thm:click-povm}\,(ii), read
off the single declared fact $\InpS=\varnothing$ with no arithmetic
performed, so the flag $\mathrm{exact}$ is admissible; at
$\gamma=\beta_k=1$ hypothesis \eqref{eq:collapsehyp} of
Lemma~\ref{lem:collapse} reads $K\ge K$, which is equality. The bound data
are \eqref{eq:cd} at these operators.
\end{proof}

% claim receiver/def:pmc
\claimsummary{receiver/def:pmc}{At a setting without an exact-window depth, the per-setting weight is the collapse value where the flag is exact and a declared window with its tail otherwise. The uniform weight is the minimum of the per-setting weights over all settings.}
\claimbookmark{receiver/def:pmc}{Definition}{def:pmc}{The uniform high-photon multiclick weight}
\claimleadin{%
Each setting without an exact-window depth gets its own weight, by the
collapse of Lemma~\ref{lem:collapse} where the setting's flag guarantees its hypothesis
and by the window and tail of Lemma~\ref{lem:windowtail} otherwise.
Remark~\ref{rem:flag-instances} says where the flag is set at the two
instances this module constructs.
}
\begin{definition}[The uniform high-photon multiclick weight]\label{def:pmc}
Fix a cutoff $N$ and a \termref{term-setting-family}{setting family}
(Definition~\ref{def:setting-family}); for each setting $t\in\Tset$ let
$\Bnd^{(t)}$ be the \termref{term-bound-function}{bound function} \eqref{eq:beta} of the
\termref{term-bound-data-of-the-setting}{bound data of the setting}. At every setting $t$ whose declaration
carries \emph{no} \termref{term-exact-window-depth}{exact-window depth} (Definition~\ref{def:pmcex}),
the
\termdef{term-per-setting-weight}{per-setting weight}
is
\begin{equation}\label{eq:pmcdef}
  \pmc^{(t)}\coloneqq
  \begin{cases}
    \max\big\{0,\ \Bnd^{(t)}(N+1)\big\}
      &\text{if (X) holds at }t,\\[4pt]
    \max\big\{0,\ \underline{\Bnd}^{(t)}_{N,m_t}\big\}
      &\text{otherwise, for the declared window }m_t>N,
  \end{cases}
\end{equation}
and the
\termdef{term-uniform-weight}{uniform weight}
is $\pmc\coloneqq\min_{t\in\Tset}\pmc^{(t)}$, where the branch-1
\termdef{term-predicate-x}{predicate (X)}
is
\begin{description}[leftmargin=3.2em,style=nextline]
\item[(X) \emph{exact by construction, decidable from the declaration:}] the
  setting's \termref{term-exact-flag}{exact flag} is $\chi^{(t)}=\mathrm{exact}$
  (Definition~\ref{def:setting-family}\,(c)). Then
  $\gamma^{(t)}=\beta^{(t)}_k=1$ exactly, so \eqref{eq:collapsehyp} reads
  $K\ge K$ and holds with \emph{equality}, no arithmetic having been
  performed.
\end{description}
\emph{Where the exact window replaces this.} At a setting whose declaration
carries an exact-window depth $m^{\ast}_t$, \eqref{eq:pmcdef} does not apply
and $\pmc^{(t)}$ is \eqref{eq:pmcex} instead
(Definition~\ref{def:pmcex}); $\pmc=\min_{t\in\Tset}\pmc^{(t)}$ is unchanged,
and the two kinds of setting may be mixed in one declaration. A declaration
carrying no depth at any setting is governed by \eqref{eq:pmcdef} alone: no
$\Bnd^{\ast}$ is formed, no eigenvalue of any block above the cutoff is
computed, and Assumption~\ref{ass:eig} is not invoked.
\end{definition}
\claimmeta{\metalabel{Uses}\mref{Def.}{def:bounddata}, \mref{Lem.}{lem:collapse}, \mref{Lem.}{lem:windowtail}, \mref{Def.}{def:setting-family}
\metasep \metalabel{Used by}\mref{Rem.}{rem:pmc-branch-choice}, \mref{Rem.}{rem:flag-instances}, \mref{Rem.}{rem:noslack}, \mref{Def.}{def:pmcex}, \mref{Rem.}{rem:exact-window-remainders}, \mref{Prop.}{prop:exactnoworse}, \mref{Rem.}{rem:noworsefl}, \mref{Def.}{def:enclosed-weight}, \mref{Thm.}{thm:enclosed-weight}, \mref{Thm.}{thm:emitted-bound}, \mref{Cor.}{cor:emitted-unconditioned}, \mref{Def.}{def:tolmap}, \mref{Thm.}{thm:exact-and-evaluated}, \mref{Req.}{req:window-declaration}, \mref{Req.}{req:exact-window-declaration}, \mref{Req.}{req:rigorous-weight}, \mref{Req.}{req:residual-refusal}, \mref{Ex.}{ex:bb84bound}, \mref{Ex.}{ex:bb84exact}
\metasep newly derived
\metasep implemented (effective: accepted)
\metasep \metalabel{Code}\coderef{https://github.com/fraunhofersg/certiqs-key/blob/v0.1.0/src/certiqs_key/receiver/bound.py}{certiqs_key.receiver.bound:SettingBound}}

% claim receiver/rem:pmc-branch-choice
\claimsummary{receiver/rem:pmc-branch-choice}{The collapse branch is taken only where predicate (X) holds, because it can overstate the bound where the collapse hypothesis fails, while the window branch is valid always. A computed margin test would never fire, since the collapse slack is never positive.}
\claimbookmark{receiver/rem:pmc-branch-choice}{Remark}{rem:pmc-branch-choice}{Choice of the collapse predicate}
\begin{remark}[Choice of the collapse predicate]\label{rem:pmc-branch-choice}
\emph{Why the predicate is exactly (X).} Branch 2 of \eqref{eq:pmcdef} is
valid for \emph{every} admissible bound data with no hypothesis
(Lemma~\ref{lem:windowtail}\,(i)), whereas branch 1 needs
\eqref{eq:collapsehyp} and can, where that fails, return a value strictly
above $\inf_{n>N}\Bnd^{(t)}(n)$: an over-large lower bound, which is the
unsafe direction. An instance on admissible bound data is $K=4$,
$\gamma^{(t)}=\beta^{(t)}_k=\tfrac{1}{100}$ and $d^{(t)}_k=1$ for every $k$,
and $0<c^{(t)}<1$, for which
$\Bnd^{(t)}(n)=\tfrac{96}{100}+\tfrac{3}{100}\,(c^{(t)})^{n}$ is positive and
strictly decreasing in $n$, so that branch~1 returns
$\Bnd^{(t)}(N+1)>\tfrac{96}{100}=\inf_{n>N}\Bnd^{(t)}(n)$.
Predicate~(X) is decided on the declaration, before any
arithmetic, and admits the equality case, which is the generic one:
\eqref{eq:collapsehyp} is satisfied with slack exactly zero by every
probe-free receiver and at every probe-off setting, which is precisely the
configuration Lemma~\ref{lem:collapse} was proved for. (X) is
\emph{sufficient, not necessary}: a receiver whose slack happens to be zero
for some other reason, such as a probe that is on but decoupled from every
port, where $\Gamma^{1/2}A_{\Prb\Prb}\Gamma^{1/2}=0$ makes
$\gamma^{(t)}=\beta^{(t)}_k=1$ as well, fails (X) and is served by branch 2.
That is conservative, and Lemma~\ref{lem:windowtail}\,(i) makes it sound.

\emph{Why there is no computed companion to (X).} A test of the form ``take
branch 1 when the evaluated slack $\sum_k\beta^{(t)}_k-1-(K-1)\gamma^{(t)}$
clears a declared margin'' would never fire, and Remark~\ref{rem:noslack}
proves it: that slack equals $-\Bnd^{(t)}(0)$, which is $\le0$ for every
receiver, at every setting, at every parameter value. A margin test on it
could therefore be met only by floating-point noise, which is the one case a
margin exists to exclude. There is no such predicate here, and there is no
declared margin among the inputs.
\end{remark}
\claimmeta{\metalabel{Uses}\mref{Lem.}{lem:collapse}, \mref{Lem.}{lem:windowtail}, \mref{Def.}{def:pmc}, \mref{Rem.}{rem:noslack}
\metasep newly derived
\metasep accepted}

% claim receiver/rem:flag-instances
\claimsummary{receiver/rem:flag-instances}{The construction of each setting sets its exact flag; at the one setting of a probe-free receiver it is exact. Every setting whose flag is general takes the general branch of whichever weight definition governs it, read off the declaration alone.}
\claimbookmark{receiver/rem:flag-instances}{Remark}{rem:flag-instances}{Where the exact flag is set, and what a general flag selects}
\begin{remark}[Where the exact flag is set, and what a general flag selects]
\label{rem:flag-instances}
Predicate~(X) of Definition~\ref{def:pmc} is the setting's exact flag
(Definition~\ref{def:setting-family}\,(c)), and the flag is set by the
construction that produced the setting. At the instance the core
constructs, the one setting of a probe-free receiver, it is set
(Corollary~\ref{cor:unconditioned-family}); a submodule that constructs a
setting family states at which of its settings it is set. At every setting
whose flag is $\mathrm{general}$ the general branch of whichever definition
governs the setting is taken: branch~2 of \eqref{eq:pmcdef} where the declaration
carries no exact-window depth, and the second branch of
\eqref{eq:pmcex} where it carries one (Definition~\ref{def:pmcex}). Which of
the two flag values applies is read off the declaration alone, before any
arithmetic.
\end{remark}
\claimmeta{\metalabel{Uses}\mref{Def.}{def:setting-family}, \mref{Cor.}{cor:unconditioned-family}, \mref{Def.}{def:pmc}, \mref{Def.}{def:pmcex}
\metasep newly derived
\metasep accepted}

% claim receiver/rem:noslack
\claimsummary{receiver/rem:noslack}{At every setting of a setting family the collapse slack is \ensuremath{-}\ensuremath{\beta}(0), where \ensuremath{\beta}(0) equals the vacuum block of the multiclick element, so it is never positive and the collapse hypothesis holds only with equality. A computed margin test for branch 1 could therefore never fire.}
\claimbookmark{receiver/rem:noslack}{Remark}{rem:noslack}{Sign of the collapse slack}
\begin{remark}[Sign of the collapse slack]\label{rem:noslack}
\emph{The slack is the negative vacuum block.} Fix a setting $t$ and write
$\sigma\coloneqq\sum_k\beta^{(t)}_k-1-(K-1)\gamma^{(t)}$ for the slack in
\eqref{eq:collapsehyp}. Evaluating \eqref{eq:beta} at $n=0$ under the
convention $0^0=1$ gives
$\Bnd^{(t)}(0)=1+(K-1)\gamma^{(t)}-\sum_k\beta^{(t)}_k=-\sigma$, so the two are
negatives of each other. For a setting of a setting family
(Definition~\ref{def:setting-family}) the vacuum block of the multiclick
element is $\blk{M}{(=0)}_{\mc|t}=\Bnd^{(t)}(0)$ exactly: that block is one
dimensional and every $\Fock(\cdot)^{(0)}=1$, so \eqref{eq:family-prefactor}
reads $\blk{\Pvac{[K]}}{(=0)}=\gamma^{(t)}$ and
$\blk{\Pvac{[K]\setminus\{k\}}}{(=0)}=\beta^{(t)}_k$, and the $n=0$ block of
\eqref{eq:family-mc} is
$1+(K-1)\gamma^{(t)}-\sum_k\beta^{(t)}_k=\Bnd^{(t)}(0)$. That
block is a diagonal block of a positive operator
(Definition~\ref{def:setting-family}\,(a)), hence a non-negative real. So
$\Bnd^{(t)}(0)\ge0$ and $\sigma\le0$ always, and \eqref{eq:collapsehyp} holds
\emph{if and only if} $\Bnd^{(t)}(0)=0$, that is $\sigma=0$ exactly.

\emph{No computed margin test can select branch 1.} A test demanding
$\sigma\ge\varepsilon>0$ is therefore unsatisfiable, and a test demanding
$\sigma\ge0$ is satisfied only in the exact-equality case, which is what (X)
already guarantees from the declaration without computing anything. So
branch 1 of \eqref{eq:pmcdef} is taken only through (X), and it is by design
that no arithmetic decides it. The declared $\varepsilon_{\mathrm{eig}}$ of
Definition~\ref{def:pmcex} is not a margin of this kind and decides no
predicate: it is a slack subtracted from a computed eigenvalue, and the exact
window is selected by the presence of a declared depth, which is again a
property of the declaration.

\emph{Nothing is lost.} Branch 2 is unconditionally valid
(Lemma~\ref{lem:windowtail}\,(i)). At a setting where \eqref{eq:collapsehyp}
holds but (X) does not, it recovers branch 1's value in the limit of a
widening window whenever $c^{(t)}<1$: by Lemma~\ref{lem:windowtail}\,(ii) its
value $\max\{0,\underline{\Bnd}^{(t)}_{N,m}\}$ tends to
$\max\{0,\inf_{n>N}\Bnd^{(t)}(n)\}$, which is $\max\{0,\Bnd^{(t)}(N+1)\}$ by
\eqref{eq:collapseinf}.
\end{remark}
\claimmeta{\metalabel{Uses}\mref{Def.}{def:bounddata}, \mref{Lem.}{lem:collapse}, \mref{Lem.}{lem:windowtail}, \mref{Def.}{def:setting-family}, \mref{Def.}{def:pmc}
\metasep \metalabel{Used by}\mref{Rem.}{rem:pmc-branch-choice}
\metasep newly derived
\metasep accepted}

% claim receiver/def:pmcex
\claimsummary{receiver/def:pmcex}{At a setting with a declared depth, the weight uses on each block above the cutoff up to the depth the larger of the analytic bound and the computed block eigenvalue less the eigenvalue slack, and keeps the setting's analytic remainder above. A depth changes only its own setting's weight.}
\claimbookmark{receiver/def:pmcex}{Definition}{def:pmcex}{Exact-window weight of a setting with a declared depth}
\claimleadin{%
The analytic weight of Definition~\ref{def:pmc} bounds three operators
separately and so loses slack (Remark~\ref{rem:notsharp}). Where a
declaration asks for it, the exact window uses, on finitely many blocks
above the cutoff, the larger of the analytic value and a computed block
eigenvalue less a declared slack, and keeps the setting's analytic
remainder above those blocks; Lemma~\ref{lem:exactwindow} makes the
combination a bound under Assumption~\ref{ass:eig}, and
Remark~\ref{rem:exact-window-remainders} explains the remainder each branch
keeps. What a declaration must supply, and which depths are refused, is
Requirement~\ref{req:exact-window-declaration}.
}
\begin{definition}[Exact-window weight of a setting with a declared depth]
\label{def:pmcex}
Let $t\in\Tset$ be a setting whose declaration carries an
\termdef{term-exact-window-depth}{exact-window depth}
$m^{\ast}_t$: an integer with $m^{\ast}_t>N$ and, where predicate~(X)
fails at $t$, also $m^{\ast}_t\le m_t$. Write $\widehat\Bnd^{\ast(t)}(n)$ for
the value an implementation computes for $\Bnd^{\ast(t)}(n)$ of
\eqref{eq:exactbnd} at that setting, let $\varepsilon_{\mathrm{eig}}$ be the
declared
\termdef{term-eigenvalue-slack}{eigenvalue slack},
and put
\begin{equation}\label{eq:elldef}
  \ell^{(t)}(n)\ \coloneqq\
  \max\big\{\ \Bnd^{(t)}(n),\ \ \widehat\Bnd^{\ast(t)}(n)
   -\varepsilon_{\mathrm{eig}}\ \big\},
  \qquad N<n\le m^{\ast}_t .
\end{equation}
Then
\begin{equation}\label{eq:pmcex}
  \pmc^{(t)}\coloneqq
  \begin{cases}
    \max\Big\{0,\ \min\big\{\ \min\limits_{N<n\le m^{\ast}_t}\ell^{(t)}(n),
      \ \ \Bnd^{(t)}\big(m^{\ast}_t+1\big)\ \big\}\Big\}
      &\text{if (X) holds at }t,\\[10pt]
    \max\Big\{0,\ \min\big\{\ \min\limits_{N<n\le m^{\ast}_t}\ell^{(t)}(n),
      \ \ \min\limits_{m^{\ast}_t<n\le m_t}\Bnd^{(t)}(n),
      \ \ T^{(t)}_{m_t}\ \big\}\Big\}
      &\text{otherwise,}
  \end{cases}
\end{equation}
the middle minimum in the second case running over the empty index set ---
and so equal to $+\infty$ and inactive --- when $m^{\ast}_t=m_t$.

\emph{Per setting.} The depth is optional and per setting, and each
per-setting weight $\pmc^{(t)}$ is its own setting's: a depth changes the
weight of the setting that declares it and of no other. A declaration that
carries a depth at some settings and not at others therefore obtains the
exact-window value as the weight of each setting with one, while the uniform
weight $\pmc=\min_{t\in\Tset}\pmc^{(t)}$ is bounded by the settings without
one.
\end{definition}
\claimmeta{\metalabel{Uses}\mref{Lem.}{lem:collapse}, \mref{Lem.}{lem:windowtail}, \mref{Def.}{def:exactbnd}, \mref{Lem.}{lem:exactwindow}, \mref{Def.}{def:setting-family}, \mref{Def.}{def:pmc}
\metasep \metalabel{Used by}\mref{Rem.}{rem:flag-instances}, \mref{Rem.}{rem:exact-window-remainders}, \mref{Prop.}{prop:exactnoworse}, \mref{Rem.}{rem:exacteig}, \mref{Def.}{def:enclosed-weight}, \mref{Thm.}{thm:enclosed-weight}, \mref{Thm.}{thm:emitted-bound}, \mref{Cor.}{cor:emitted-unconditioned}, \mref{Rem.}{rem:notsharp}, \mref{Def.}{def:residuals}, \mref{Def.}{def:tolmap}, \mref{Rem.}{rem:tolerancedomain}, \mref{Thm.}{thm:exact-and-evaluated}, \mref{Rem.}{rem:ia}, \mref{Ass.}{ass:eig}, \mref{Req.}{req:exact-window-declaration}, \mref{Req.}{req:rigorous-weight}, \mref{Ex.}{ex:bb84exact}
\metasep newly derived
\metasep implemented (effective: accepted)
\metasep \metalabel{Code}\coderef{https://github.com/fraunhofersg/certiqs-key/blob/v0.1.0/src/certiqs_key/receiver/bound.py}{certiqs_key.receiver.bound:ExactWindowRecord}}

% claim receiver/rem:exact-window-remainders
\claimsummary{receiver/rem:exact-window-remainders}{The exact window covers finitely many blocks, so each branch keeps an admissible remainder above it: branch 1 the bound function one past the depth, branch 2 its own analytic window and tail. Neither remainder is new or weakened.}
\claimbookmark{receiver/rem:exact-window-remainders}{Remark}{rem:exact-window-remainders}{The analytic remainder of each exact-window branch}
\begin{remark}[The analytic remainder of each exact-window branch]
\label{rem:exact-window-remainders}
The exact window of Definition~\ref{def:pmcex} covers $N<n\le m^{\ast}_t$
and nothing above it, so a remainder admissible in the sense of
\eqref{eq:tailadm} is still required, and each branch of \eqref{eq:pmcex}
already has one. Branch~1 has \eqref{eq:collapsehyp} exactly, by
predicate~(X) of Definition~\ref{def:pmc}, and takes
$\Bnd^{(t)}(m^{\ast}_t+1)$ by Lemma~\ref{lem:exactwindow}\,(ii). Branch~2
has no hypothesis available and keeps its own analytic window and tail
\eqref{eq:windowtail} above the exact one; this is why
$m^{\ast}_t\le m_t$ is required there, so that the two windows compose into
the single window $N<n\le m_t$ to which Lemma~\ref{lem:exactwindow} is
applied. Neither remainder is new and neither is weakened.
\end{remark}
\claimmeta{\metalabel{Uses}\mref{Lem.}{lem:collapse}, \mref{Lem.}{lem:windowtail}, \mref{Lem.}{lem:exactwindow}, \mref{Def.}{def:pmc}, \mref{Def.}{def:pmcex}
\metasep newly derived
\metasep accepted}

% claim receiver/prop:exactnoworse
\claimsummary{receiver/prop:exactnoworse}{At every setting with a declared depth, the exact-window weight is at least the analytic weight on the same bound data. The comparison is between exact expressions; \hyperlink{claim-receiver/rem:noworsefl}{Remark~\ref*{rem:noworsefl}} says which of its steps survive an evaluation.}
\claimbookmark{receiver/prop:exactnoworse}{Proposition}{prop:exactnoworse}{Dominance of the exact-window weight over the analytic weight}
\begin{proposition}[Dominance of the exact-window weight over the analytic weight]\label{prop:exactnoworse}
At every setting carrying a declared \termref{term-exact-window-depth}{exact-window depth}, the value
\eqref{eq:pmcex} is at least the value \eqref{eq:pmcdef} gives at that
setting on the same bound data.
\end{proposition}
\claimmeta{\metalabel{Uses}\mref{Lem.}{lem:collapse}, \mref{Lem.}{lem:windowtail}, \mref{Def.}{def:pmc}, \mref{Def.}{def:pmcex}
\metasep \metalabel{Used by}\mref{Rem.}{rem:noworsefl}, \mref{Rem.}{rem:exacteig}
\metasep newly derived
\metasep accepted}

\begin{proof}
Write $P^{\ast}$ for \eqref{eq:pmcex} and $P$ for \eqref{eq:pmcdef} and drop
the superscript $(t)$. From the maximum in \eqref{eq:elldef},
$\ell(n)\ge\Bnd(n)$ for every $n$ in the exact window.

\emph{Branch 2.} By \eqref{eq:windowtail},
$P=\max\{0,\min\{\min_{N<n\le m_t}\Bnd(n),\,T_{m_t}\}\}$. Splitting the window
at $m^{\ast}_t$,
$\min_{N<n\le m_t}\Bnd(n)=\min\{\min_{N<n\le m^{\ast}_t}\Bnd(n),\,
\min_{m^{\ast}_t<n\le m_t}\Bnd(n)\}$, and \eqref{eq:pmcex} is that same
expression with the first of the two replaced by
$\min_{N<n\le m^{\ast}_t}\ell(n)$, which is at least it. A minimum is
non-decreasing in each argument and $\max\{0,\cdot\}$ is non-decreasing, so
$P^{\ast}\ge P$.

\emph{Branch 1.} $P=\max\{0,\Bnd(N+1)\}$. If $\Bnd(N+1)<0$ then $P=0\le
P^{\ast}$. If $\Bnd(N+1)\ge0$ then, predicate~(X) making \eqref{eq:collapsehyp}
hold with equality, Lemma~\ref{lem:collapse} makes $\Bnd$ non-decreasing on
$\{N+1,N+2,\dots\}$; hence $\ell(n)\ge\Bnd(n)\ge\Bnd(N+1)$ for
$N<n\le m^{\ast}_t$, and $\Bnd(m^{\ast}_t+1)\ge\Bnd(N+1)$. Both arguments of
the inner minimum in \eqref{eq:pmcex} are therefore at least $\Bnd(N+1)$, so
$P^{\ast}\ge\max\{0,\Bnd(N+1)\}=P$.
\end{proof}

% claim receiver/rem:noworsefl
\claimsummary{receiver/rem:noworsefl}{Two steps of \hyperlink{claim-receiver/prop:exactnoworse}{Proposition~\ref*{prop:exactnoworse}} compare evaluated numbers and survive an evaluation; branch 1's use of the collapse lemma does not. There the exact-window weight can sit below the analytic one by evaluation errors of the bound function between the first block above the cutoff and one past the depth.}
\claimbookmark{receiver/rem:noworsefl}{Remark}{rem:noworsefl}{The exact-window comparison under floating-point evaluation}
\begin{remark}[The exact-window comparison under floating-point evaluation]
\label{rem:noworsefl}
Proposition~\ref{prop:exactnoworse} compares two \emph{expressions} in exact
arithmetic. Two of its three steps survive an evaluation term for term,
because they compare evaluated numbers with each other: $\ell(n)\ge\Bnd(n)$
is an evaluated maximum, and branch~2 minimises over exactly the evaluated
$\Bnd(n)$ that \eqref{eq:pmcdef} minimises over. The step that does not is
branch~1's use of Lemma~\ref{lem:collapse},
$\Bnd(n)\ge\Bnd(N+1)$ for $N+1<n\le m^{\ast}_t+1$: a statement about exact
values, and those further $\Bnd(n)$ are quantities \eqref{eq:pmcdef} never
evaluates, whose evaluated values can fall below the evaluated $\Bnd(N+1)$.
So on branch~1 the emitted value can sit below the one \eqref{eq:pmcdef}
would produce by the evaluation errors of those further values together with
that of $\Bnd(N+1)$, and this contract claims no bound on either
(Theorem~\ref{thm:exact-and-evaluated}\,(ii)). Unaffected either way are the
result that underwrites the emitted number (Theorem~\ref{thm:emitted-bound}) and
the fact that a declaration carrying no depth performs none of this
arithmetic.
\end{remark}
\claimmeta{\metalabel{Uses}\mref{Lem.}{lem:collapse}, \mref{Def.}{def:pmc}, \mref{Prop.}{prop:exactnoworse}
\metasep newly derived
\metasep accepted}

% claim receiver/rem:exacteig
\claimsummary{receiver/rem:exacteig}{In the exact-window maximum the analytic entry is a proven lower bound, while the computed eigenvalue less the slack is one only under \hyperlink{claim-receiver/ass:eig}{A-EIG}; the rigorous evaluation replaces it by a proven bound. No test catches an overshoot, r\_6 detects only an undershoot, and each weight stays in [0, 1].}
\claimbookmark{receiver/rem:exacteig}{Remark}{rem:exacteig}{The eigenvalue-slack assumption behind the exact window}
\claimleadin{%
The exact window trades the analytic bound for a sharper value on finitely
many blocks, and in the floating-point evaluation the price is one assumption
on the computed eigenvalues. The rigorous evaluation pays no such price: it
puts a rigorous lower bound on the block eigenvalue in the computed value's
place (Definition~\ref{def:enclosed-weight},
Requirement~\ref{req:rigorous-weight}), and the weight is then a bound by
Theorem~\ref{thm:enclosed-weight}.
}
\begin{remark}[The eigenvalue-slack assumption behind the exact window]
\label{rem:exacteig}
\eqref{eq:elldef} is a maximum of two lower bounds on $\Bnd^{\ast(t)}(n)$,
and they are of two different \emph{kinds}. The first,
$\Bnd^{(t)}(n)\le\Bnd^{\ast(t)}(n)$, is Theorem~\ref{thm:multiclick-bound}: an
inequality between exact quantities, holding with no hypothesis about any
computation. The second,
$\widehat\Bnd^{\ast(t)}(n)-\varepsilon_{\mathrm{eig}}\le\Bnd^{\ast(t)}(n)$,
is \emph{not a theorem of this contract}:
the \termref{term-evaluation-model}{evaluation model} of Definition~\ref{def:evaluation-model} asserts no bound on
the difference between a quantity computed in floating point and the exact
one, and $\widehat\Bnd^{\ast}$ is the output
of an extreme-eigenvalue computation on a matrix of dimension
$\binom{|\InpC|+n-1}{n}$. It is Assumption~\ref{ass:eig}, and that
assumption is invoked \emph{only} at the settings and indices where a depth
is declared and the weight is that of Definition~\ref{def:pmcex}. What is a
theorem is the same maximum with a \termref{term-rigorous-window-bounds}{rigorous
window bound} in the second entry: in the \termref{term-enclosed-weight}{enclosed weight} of
Definition~\ref{def:enclosed-weight} the entry is a real proved to lie below
$\Bnd^{\ast(t)}(n)$, not an evaluation of it, no slack is subtracted, and
Theorem~\ref{thm:enclosed-weight} makes the weight a bound with no assumption.
Everything below concerns the floating-point evaluation.

\emph{The direction.} A computed eigenvalue landing \emph{below} the exact
one costs tightness and nothing else, and \eqref{eq:elldef} floors the value
at $\Bnd^{(t)}(n)$, which is sound with no assumption at all; so that
direction cannot take the emitted weight below what the analytic evaluation
emits (Proposition~\ref{prop:exactnoworse}). A computed eigenvalue landing
\emph{above} the exact one by more than $\varepsilon_{\mathrm{eig}}$ can make
the emitted weight $\pmc^{(t)}$ of that setting exceed
$\inf_{n>N}\lmin(\blk{M}{\eqn}_{\mc|t})$: an over-large lower bound, the
unsafe direction, and one that reaches a key length through
Definition~\ref{def:slot}. No range check catches it, whatever the size of
the overshoot. Every branch of \eqref{eq:pmcex} is capped by an analytic
remainder that is at most $1$ and does not depend on
$\widehat\Bnd^{\ast(t)}$, so $\pmc^{(t)}$ stays in $[0,1]$ (the
\emph{Range} paragraph of this remark), and neither the ``$p_y<0$'' nor the
``$\sum_yp_y>1$'' failure condition that Proposition~\ref{prop:slotvalid}
names can fire on it. An overshoot of any size is caught by nothing. That
asymmetry is why Assumption~\ref{ass:eig} is
stated one-sidedly, on the side that matters, and why
$\varepsilon_{\mathrm{eig}}$ is \emph{subtracted} rather than used as a
comparison threshold: it decides nothing, it lowers a number.

\emph{A slack, not a repair.} The computed $\widehat\Bnd^{\ast(t)}(n)$
enters \eqref{eq:elldef} unmodified, and $\varepsilon_{\mathrm{eig}}$ enters
only as a subtracted declared slack, so that \eqref{eq:elldef} can be
recomputed from the recorded values
(Requirement~\ref{req:exact-window-record}). This is the distinction
Remark~\ref{rem:d2slack} draws for $\bar r$ and
Definition~\ref{def:tolmap}\,(iv) draws for $c$: a coefficient or a slack
inside a stated inequality is not a repair of a datum.

\emph{What $r_6$ tests, and the direction it is blind to.} Since
$\Bnd(n)\le\Bnd^{\ast}(n)$ is a theorem (Lemma~\ref{lem:exactbasic}), the
window consistency defect $r_6$ of \eqref{eq:res6} has exact value $0$, and
it is a residual in exactly the sense of the other five: what is emitted is
its floating-point evaluation, and a value above $\varepsilon_6$ refuses. It
catches an evaluation that contradicts an inequality this contract proves
--- a block formed at the wrong index, a wrongly ordered occupation basis, a
conditioned family taken at the wrong setting, an iteration that has not
converged --- which is the failure mode most likely to be gross. It does
\emph{not} verify Assumption~\ref{ass:eig}: $r_6$ fires when
$\widehat\Bnd^{\ast}$ lands too \emph{low}, and Assumption~\ref{ass:eig}
constrains it from landing too \emph{high}. No test in this contract bounds
the unsafe direction; the assumption carries it.

\emph{Range.} $0\le\pmc^{(t)}\le1$ holds exactly, whatever the computed
values $\widehat\Bnd^{\ast(t)}(n)$ and whether or not
Assumption~\ref{ass:eig} holds: the outer maximum of \eqref{eq:pmcex} gives
$\pmc^{(t)}\ge0$, and the inner minimum of each branch contains an analytic
remainder at most $1$ --- $\Bnd^{(t)}(m^{\ast}_t+1)\le1$ on branch~1 by
Theorem~\ref{thm:multiclick-bound}, and the tail value
$T^{(t)}_{m_t}=1-\sum_k\beta^{(t)}_k\big(d^{(t)}_k\big)^{m_t}\le1$ of
\eqref{eq:windowtail} on branch~2 because $\beta^{(t)}_k,d^{(t)}_k\ge0$ by
\eqref{eq:bounddata}. Lemma~\ref{lem:exactwindow} gives the same range only
through \eqref{eq:ellhyp}, which at a setting carrying a depth is
Assumption~\ref{ass:eig}. So Theorem~\ref{thm:emitted-bound}'s range statement
and Proposition~\ref{prop:slotvalid} hold unchanged; their floating-point
evaluation is governed by Theorem~\ref{thm:exact-and-evaluated}\,(ii) here as everywhere
else, and no computed value is clamped into $[0,1]$ to make them look true.
\end{remark}
\claimmeta{\metalabel{Uses}\mref{Def.}{def:bounddata}, \mref{Thm.}{thm:multiclick-bound}, \mref{Lem.}{lem:windowtail}, \mref{Lem.}{lem:exactbasic}, \mref{Lem.}{lem:exactwindow}, \mref{Def.}{def:pmcex}, \mref{Prop.}{prop:exactnoworse}
\metasep \metalabel{Assumes}\aref{A-EIG}{ass:eig}
\metasep \metalabel{Used by}\mref{Rem.}{rem:tolerancedomain}, \mref{Req.}{req:exact-window-record}
\metasep newly derived
\metasep accepted}

% claim receiver/def:enclosed-weight
\claimsummary{receiver/def:enclosed-weight}{Enclosed bound data take c and \ensuremath{\gamma} from below and each d and \ensuremath{\beta} from above, within [0, 1], and rigorous window bounds lie below the exact block eigenvalues. The enclosed weight is the per-setting weight formed on them, with no computed eigenvalue and no eigenvalue slack.}
\claimbookmark{receiver/def:enclosed-weight}{Definition}{def:enclosed-weight}{The weight evaluated on enclosed bound data}
\claimleadin{%
The weights of Definitions~\ref{def:pmc} and~\ref{def:pmcex} are formulas in
the bound data of a setting and, where a depth is declared, in computed block
eigenvalues. The same formulas, evaluated on numbers known to lie on the
conservative side of the exact data, give a weight that rests on no premise
about how it was computed (Theorem~\ref{thm:enclosed-weight}). The four
directions are those Remark~\ref{rem:ia} reads off the bound function, and in
the window a rigorous lower bound on the exact block eigenvalue takes the
place of the computed eigenvalue less the slack, which is why the slack does
not enter (Remark~\ref{rem:exacteig}).
}
\begin{definition}[The weight evaluated on enclosed bound data]\label{def:enclosed-weight}
Fix a cutoff $N$, a \termref{term-setting-family}{setting family}
(Definition~\ref{def:setting-family}) and a setting $t\in\Tset$, and write
$\big(\gamma,\{\beta_k\},c,\{d_k\}\big)$ for the
\termref{term-bound-data-of-the-setting}{bound data of the setting}, the index $t$ omitted.
\begin{enumerate}[label=(\roman*),leftmargin=2.4em]
\item Reals $\big(\check\gamma,\{\check\beta_k\},\check c,\{\check d_k\}\big)$
  are
  \termdef{term-enclosed-bound-data}{enclosed bound data} of $t$ when
  \begin{equation}\label{eq:enclosed-directions}
    0\le\check\gamma\le\gamma,\qquad \beta_k\le\check\beta_k\le1,\qquad
    0\le\check c\le c,\qquad d_k\le\check d_k\le1
    \qquad(k=1,\dots,K).
  \end{equation}
  $\check\Bnd$ denotes the \termref{term-bound-function}{bound function} \eqref{eq:beta} of this
  tuple, and $\check T_m$ and $\underline{\check\Bnd}_{N,m}$ the quantities
  $T_m$ and $\underline{\Bnd}_{N,m}$ of \eqref{eq:windowtail} formed from it.
\item Where $t$ carries an \termref{term-exact-window-depth}{exact-window depth} $m^{\ast}_t$
  (Definition~\ref{def:pmcex}), reals $w(n)$, $N<n\le m^{\ast}_t$, are
  \termdef{term-rigorous-window-bounds}{rigorous window bounds} of $t$ when
  $w(n)\le\Bnd^{\ast(t)}(n)$ for each such $n$, $\Bnd^{\ast(t)}$ being the
  \termref{term-exact-block-bound}{exact block bound} \eqref{eq:exactbnd} at $t$.
\item The
  \termdef{term-enclosed-weight}{enclosed weight} $\check p^{(t)}_{\mc}$ of $t$ formed
  from them is \eqref{eq:pmcdef} where $t$ carries no exact-window depth and
  \eqref{eq:pmcex} where it carries one, its branch chosen by
  \termref{term-predicate-x}{predicate (X)} at $t$ as there, each read with $\check\Bnd$,
  $\underline{\check\Bnd}_{N,m_t}$ and $\check T_{m_t}$ in place of
  $\Bnd^{(t)}$, $\underline{\Bnd}^{(t)}_{N,m_t}$ and $T^{(t)}_{m_t}$, and with
  \begin{equation}\label{eq:enclosed-ell}
    \check\ell(n)\ \coloneqq\ \max\big\{\check\Bnd(n),\ w(n)\big\},
    \qquad N<n\le m^{\ast}_t ,
  \end{equation}
  in place of $\ell^{(t)}(n)$ of \eqref{eq:elldef}. No computed eigenvalue of
  a block and no \termref{term-eigenvalue-slack}{eigenvalue slack} enters it.
\end{enumerate}
\end{definition}
\claimmeta{\metalabel{Uses}\mref{Def.}{def:bounddata}, \mref{Lem.}{lem:windowtail}, \mref{Def.}{def:exactbnd}, \mref{Def.}{def:setting-family}, \mref{Def.}{def:pmc}, \mref{Def.}{def:pmcex}
\metasep \metalabel{Used by}\mref{Thm.}{thm:enclosed-weight}, \mref{Thm.}{thm:emitted-bound}, \mref{Cor.}{cor:emitted-unconditioned}, \mref{Def.}{def:evaluation-model}, \mref{Thm.}{thm:exact-and-evaluated}, \mref{Req.}{req:rigorous-weight}
\metasep newly derived
\metasep implemented (effective: accepted)
\metasep \metalabel{Code}\coderef{https://github.com/fraunhofersg/certiqs-key/blob/v0.1.0/src/certiqs_key/receiver/rigorous.py}{certiqs_key.receiver.rigorous:enclosed_weight}}

% claim receiver/thm:enclosed-weight
\claimsummary{receiver/thm:enclosed-weight}{Bounds on the four inputs and on the window blocks, taken on the conservative side and clipped to [0, 1], give enclosed bound data, and the enclosed weight bounds its setting's multiclick element on every block above the cutoff. No premise about arithmetic enters, so no eigenvalue slack.}
\claimbookmark{receiver/thm:enclosed-weight}{Theorem}{thm:enclosed-weight}{The weight on enclosed bound data bounds the multiclick element}
\claimleadin{%
Each formula of the weight is non-decreasing in $c$ and in $\gamma$ and
non-increasing in each $d_k$ and $\beta_k$ on $[0,1]$, and each minimum and
maximum in it is non-decreasing in its arguments, so the weight formed on
enclosed bound data lies below the weight formed on the exact data, entry by
entry; this is the argument of Remark~\ref{rem:ia}. In the exact window the
lemma that combines the window with its remainder admits any lower bounds on
the block eigenvalues and any lower remainder, so a rigorous lower bound
serves where a computed eigenvalue needed an assumption. The theorem is the
rigorous counterpart of Theorem~\ref{thm:emitted-bound} at one setting, and it
has no assumption at all.
}
\begin{theorem}[The weight on enclosed bound data bounds the multiclick element]\label{thm:enclosed-weight}
Let a \termref{term-setting-family}{setting family} (Definition~\ref{def:setting-family}) be given, let
$N\ge0$, and let $t\in\Tset$, with operators $B=B^{(t)}$ and $B_k=B^{(t)}_k$,
\termref{term-bound-data-of-the-setting}{bound data of the setting} $\big(\gamma,\{\beta_k\},c,\{d_k\}\big)$,
\termref{term-bound-function}{bound function} $\Bnd=\Bnd^{(t)}$ and $T_m$ of \eqref{eq:windowtail}
formed from those data. Write
$\operatorname{clip}(x)\coloneqq\min\{1,\max\{0,x\}\}$ for real $x$.
\begin{enumerate}[label=(\roman*),leftmargin=2.4em]
\item \emph{Construction.} Let $u$, $\ell_1,\dots,\ell_K$, $g$ and
  $b_1,\dots,b_K$ be reals with $u\ge\lmax(B)$, $\ell_k\le\lmin(B_k)$,
  $g\le\gamma$ and $b_k\ge\beta_k$. Then
  \begin{equation}\label{eq:enclosed-construction}
    \big(\operatorname{clip}(g),\ \{\operatorname{clip}(b_k)\},\
    \operatorname{clip}(1-u),\ \{\operatorname{clip}(1-\ell_k)\}\big)
  \end{equation}
  are \termref{term-enclosed-bound-data}{enclosed bound data} of $t$. Where $t$ carries an
  \termref{term-exact-window-depth}{exact-window depth} $m^{\ast}_t$ and reals $\lambda_n$ satisfy
  $\lambda_n\le\lmin\big(\blk{M}{\eqn}_{\mc|t}\big)$ for
  $N<n\le m^{\ast}_t$, the reals $\max\{0,\lambda_n\}$ are
  \termref{term-rigorous-window-bounds}{rigorous window bounds} of $t$.
\item \emph{Comparison.} Enclosed bound data
  $\big(\check\gamma,\{\check\beta_k\},\check c,\{\check d_k\}\big)$ of $t$
  are \termref{term-admissible-bound-data}{admissible bound data} (Definition~\ref{def:bounddata}), and their
  bound function $\check\Bnd$ and $\check T_m$ satisfy
  \begin{equation}\label{eq:enclosed-comparison}
    \check\Bnd(n)\ \le\ \Bnd(n)\quad(n\ge0),\qquad
    \check T_m\ \le\ T_m\quad(m\ge0).
  \end{equation}
\item \emph{The bound.} For enclosed bound data of $t$ and, where $t$
  carries an exact-window depth, rigorous window bounds of $t$, the
  \termref{term-enclosed-weight}{enclosed weight} $\check p^{(t)}_{\mc}$ formed from them
  (Definition~\ref{def:enclosed-weight}) satisfies
  $0\le\check p^{(t)}_{\mc}\le1$ and
  \begin{equation}\label{eq:enclosed-bound}
    \blk{M}{\gtN}_{\mc|t}\ \succeq\ \check p^{(t)}_{\mc}\,\id^{(>N)} ,
  \end{equation}
  and every real $p$ with $0\le p\le\check p^{(t)}_{\mc}$ satisfies
  $0\le p\le1$ and $\blk{M}{\gtN}_{\mc|t}\succeq p\,\id^{(>N)}$.
\end{enumerate}
\end{theorem}
\claimreading{Where the receiver's model is known exactly, its high-photon multiclick weight can be computed from numbers that are proven to lie on the safe side of the exact ones, and the weight so computed, or any smaller number, is a guaranteed lower bound on the multiclick probability of every block above the cutoff. It assumes nothing about the accuracy of floating-point arithmetic or of an eigenvalue routine, and it needs no eigenvalue slack. It still rests on the receiver's device assumptions, which say that the declared model is your receiver.}
\claimmeta{\metalabel{Uses}\mref{Def.}{def:bounddata}, \mref{Thm.}{thm:multiclick-bound}, \mref{Lem.}{lem:collapse}, \mref{Lem.}{lem:windowtail}, \mref{Def.}{def:exactbnd}, \mref{Lem.}{lem:exactbasic}, \mref{Lem.}{lem:exactwindow}, \mref{Def.}{def:setting-family}, \mref{Def.}{def:pmc}, \mref{Def.}{def:pmcex}, \mref{Def.}{def:enclosed-weight}
\metasep \metalabel{Used by}\mref{Thm.}{thm:emitted-bound}, \mref{Thm.}{thm:exact-and-evaluated}, \mref{Req.}{req:rigorous-weight}, \extref{../reductions/contract.pdf}{reductions/req:recv-row-provenance}
\metasep newly derived
\metasep implemented (effective: accepted)
\metasep \metalabel{Code}\coderef{https://github.com/fraunhofersg/certiqs-key/blob/v0.1.0/src/certiqs_key/receiver/rigorous.py}{certiqs_key.receiver.rigorous:rigorous_weights}}

\begin{proof}
Two facts about the setting are used throughout. By
Definition~\ref{def:setting-family}\,(b) the hypotheses of
Theorem~\ref{thm:multiclick-bound} hold at $t$, so its bound data are
admissible bound data, $\gamma,\beta_k\in(0,1]$, and by \eqref{eq:family-mc}
\begin{equation*}
  \blk{M}{\eqn}_{\mc|t}\ \succeq\ \Bnd(n)\,\id^{(=n)}\qquad\text{for every }n\ge0
  \tag{F1}
\end{equation*}
(Theorem~\ref{thm:multiclick-bound}, \eqref{eq:mcbound}). And
$M_{\mc|t}\succeq0$, the consequence \eqref{eq:mcpos} recorded in
Definition~\ref{def:setting-family}, so Lemma~\ref{lem:exactbasic} gives
$\Bnd^{\ast(t)}(n)\in[0,1]$ for every $n$, that is,
\begin{equation*}
  \blk{M}{\eqn}_{\mc|t}\ \succeq\ 0\qquad\text{for every }n\ge0 .
  \tag{F2}
\end{equation*}
The blocks $n>N$ are orthogonal direct summands of $\blk{M}{\gtN}_{\mc|t}$, so
a lower bound $s\,\id^{(=n)}$ holding on each of them holds as
$s\,\id^{(>N)}$ on their sum.

(i) The map $\operatorname{clip}$ is non-decreasing, takes values in
$[0,1]$, and is the identity on $[0,1]$. Since $\gamma\in[0,1]$ and
$g\le\gamma$, $0\le\operatorname{clip}(g)\le\operatorname{clip}(\gamma)=\gamma$;
since $\beta_k\in[0,1]$ and $b_k\ge\beta_k$,
$\beta_k=\operatorname{clip}(\beta_k)\le\operatorname{clip}(b_k)\le1$. By
\eqref{eq:cd}, $c=1-\lmax(B)\in[0,1]$ and $u\ge\lmax(B)$ give $1-u\le c$, so
$0\le\operatorname{clip}(1-u)\le c$; and $d_k=1-\lmin(B_k)\in[0,1]$ and
$\ell_k\le\lmin(B_k)$ give $1-\ell_k\ge d_k$, so
$d_k\le\operatorname{clip}(1-\ell_k)\le1$. These are the four conditions of
\eqref{eq:enclosed-directions}. For the window, $\lmin(\blk{M}{\eqn}_{\mc|t})$
is $\Bnd^{\ast(t)}(n)$ by \eqref{eq:exactbnd}, which is at least $0$ by (F2),
so $\max\{0,\lambda_n\}\le\Bnd^{\ast(t)}(n)$.

(ii) By \eqref{eq:enclosed-directions} and the admissibility of the exact data,
$0\le\check\gamma\le\gamma\le\beta_k\le\check\beta_k\le1$ and
$0\le\check c\le c\le d_k\le\check d_k\le1$ for every $k$, which is
\eqref{eq:bounddata} for the enclosed tuple. Fix $n\ge0$. With the convention
$0^0=1$ the map $x\mapsto x^{\,n}$ is non-decreasing on $[0,\infty)$ (at $n=0$
it is the constant $1$). From $0\le\check c\le c$ follows
$0\le\check c^{\,n}\le c^{\,n}$, and with $0\le\check\gamma\le\gamma$,
$\check\gamma\,\check c^{\,n}\le\gamma\,\check c^{\,n}\le\gamma\,c^{\,n}$; as
$K-1\ge0$, $(K-1)\check\gamma\check c^{\,n}\le(K-1)\gamma c^{\,n}$. From
$0\le d_k\le\check d_k$ follows $d_k^{\,n}\le\check d_k^{\,n}$, and with
$0\le\beta_k\le\check\beta_k$,
$\beta_kd_k^{\,n}\le\beta_k\check d_k^{\,n}\le\check\beta_k\check d_k^{\,n}$.
Adding, $\check\Bnd(n)=1+(K-1)\check\gamma\check c^{\,n}
-\sum_k\check\beta_k\check d_k^{\,n}\le1+(K-1)\gamma c^{\,n}
-\sum_k\beta_kd_k^{\,n}=\Bnd(n)$. The second of these comparisons, with $m$ in
place of $n$, gives $\check T_m=1-\sum_k\check\beta_k\check d_k^{\,m}
\le1-\sum_k\beta_kd_k^{\,m}=T_m$.

(iii) Write $\check p\coloneqq\check p^{(t)}_{\mc}$. Minima and maxima of
finitely many reals are non-decreasing in each argument, which is all that is
used of them below.

\emph{Case 1: $t$ carries no exact-window depth.} If predicate~(X) holds at
$t$, $\check p=\max\{0,\check\Bnd(N+1)\}\le\max\{0,\Bnd(N+1)\}$ by (ii); the
flag gives $\gamma=\beta_k=1$ (Definition~\ref{def:setting-family}\,(c)), so
\eqref{eq:collapsehyp} holds and \eqref{eq:collapseinf} of
Lemma~\ref{lem:collapse} gives $\max\{0,\Bnd(N+1)\}=
\max\{0,\inf_{n>N}\Bnd(n)\}$. Otherwise
$\check p=\max\{0,\underline{\check\Bnd}_{N,m_t}\}$, and
$\underline{\check\Bnd}_{N,m_t}=\min\{\min_{N<n\le m_t}\check\Bnd(n),\check
T_{m_t}\}\le\min\{\min_{N<n\le m_t}\Bnd(n),T_{m_t}\}
=\underline{\Bnd}_{N,m_t}\le\inf_{n>N}\Bnd(n)$ by (ii) and
Lemma~\ref{lem:windowtail}\,(i). In both cases
$\check p\le q\coloneqq\max\{0,\inf_{n>N}\Bnd(n)\}$. For every $n>N$, (F1)
gives $\blk{M}{\eqn}_{\mc|t}\succeq\Bnd(n)\,\id^{(=n)}\succeq
\big(\inf_{j>N}\Bnd(j)\big)\id^{(=n)}$ and (F2) gives
$\blk{M}{\eqn}_{\mc|t}\succeq0$, so $\blk{M}{\eqn}_{\mc|t}\succeq q\,\id^{(=n)}
\succeq\check p\,\id^{(=n)}$, and \eqref{eq:enclosed-bound} follows on the
direct sum. $\check p\ge0$ by the outer maximum, and
$\check p\le q\le\max\{0,\Bnd(N+1)\}\le1$, since $\Bnd\le1$
(Theorem~\ref{thm:multiclick-bound}).

\emph{Case 2: a depth $m^{\ast}_t$, and predicate~(X) holds at $t$.} Apply
Lemma~\ref{lem:exactwindow} at $t$, whose hypotheses are those of
Theorem~\ref{thm:multiclick-bound} and \eqref{eq:mcpos}, with
$m^{\ast}=m^{\ast}_t$, $\ell=\check\ell$ of \eqref{eq:enclosed-ell} and
$\tau=\check\Bnd(m^{\ast}_t+1)$. Hypothesis \eqref{eq:ellhyp}: for
$N<n\le m^{\ast}_t$, $\check\Bnd(n)\le\Bnd(n)\le\Bnd^{\ast(t)}(n)$ by (ii) and
\eqref{eq:exactbasic}, and $w(n)\le\Bnd^{\ast(t)}(n)$ by
Definition~\ref{def:enclosed-weight}\,(ii), so their maximum $\check\ell(n)$
is at most $\Bnd^{\ast(t)}(n)$. Hypothesis \eqref{eq:tailadm}: the flag makes
\eqref{eq:collapsehyp} hold, as in Case~1, so
$\tau_0\coloneqq\Bnd(m^{\ast}_t+1)$ satisfies \eqref{eq:tailadm} by
Lemma~\ref{lem:exactwindow}\,(ii), and $\tau\le\tau_0$ by (ii). A real below
one that satisfies \eqref{eq:tailadm} satisfies it too: $\tau_0\le0$ gives
$\tau\le0$, and $\tau_0\le\inf_{n>m^{\ast}_t}\Bnd^{\ast(t)}(n)$ gives
$\tau\le\inf_{n>m^{\ast}_t}\Bnd^{\ast(t)}(n)$. With these choices
\eqref{eq:exactcomposite} is the first branch of \eqref{eq:pmcex} read as
Definition~\ref{def:enclosed-weight}\,(iii) reads it, that is $P=\check p$,
and the lemma gives $0\le\check p\le1$ and \eqref{eq:exactstatement}, which is
\eqref{eq:enclosed-bound}.

\emph{Case 3: a depth $m^{\ast}_t$, and predicate~(X) fails at $t$.} Then
$m^{\ast}_t\le m_t$ (Definition~\ref{def:pmcex}). Apply
Lemma~\ref{lem:exactwindow} at $t$ with $m^{\ast}=m_t$, with
$\ell(n)=\check\ell(n)$ for $N<n\le m^{\ast}_t$ and $\ell(n)=\check\Bnd(n)$ for
$m^{\ast}_t<n\le m_t$, and with $\tau=\check T_{m_t}$. Hypothesis
\eqref{eq:ellhyp} holds on the first range as in Case~2 and on the second by
$\check\Bnd(n)\le\Bnd(n)\le\Bnd^{\ast(t)}(n)$. Hypothesis \eqref{eq:tailadm}:
$T_{m_t}$ satisfies it by Lemma~\ref{lem:exactwindow}\,(i), and
$\check T_{m_t}\le T_{m_t}$ by (ii), lowered as in Case~2. With these choices
\eqref{eq:exactcomposite} reads
$\max\{0,\min\{\min_{N<n\le m^{\ast}_t}\check\ell(n),
\min_{m^{\ast}_t<n\le m_t}\check\Bnd(n),\check T_{m_t}\}\}$, the middle minimum
being over the empty set and inactive when $m^{\ast}_t=m_t$; this is the
second branch of \eqref{eq:pmcex} read as
Definition~\ref{def:enclosed-weight}\,(iii) reads it, so $P=\check p$, and the
lemma concludes as in Case~2.

\emph{Any smaller weight.} For $0\le p\le\check p$,
$\blk{M}{\gtN}_{\mc|t}\succeq\check p\,\id^{(>N)}\succeq p\,\id^{(>N)}$, since
$(\check p-p)\id^{(>N)}\succeq0$, and $0\le p\le\check p\le1$.
\end{proof}

% claim receiver/thm:emitted-bound
\claimsummary{receiver/thm:emitted-bound}{Each setting's emitted weight bounds its multiclick element on every photon-number block above the cutoff, and their minimum does so at every setting, under the eigenvalue-slack assumption \hyperlink{claim-receiver/ass:eig}{A-EIG} only where computed block eigenvalues enter. Each weight lies in [0, 1] and is finitely computable.}
\claimbookmark{receiver/thm:emitted-bound}{Theorem}{thm:emitted-bound}{Multiclick bounds carried by the emitted weights}
\claimleadin{%
This is the theorem the emitted weights rest on. Each setting's weight is
the analytic weight of Definition~\ref{def:pmc} or, where a depth is
declared, the exact-window weight of Definition~\ref{def:pmcex}; where the
engine evaluates the weight rigorously (Requirement~\ref{req:rigorous-weight})
it is the enclosed weight of Definition~\ref{def:enclosed-weight}, formed by
the same formulas from bounds on the conservative side. The theorem says that
each bounds the multiclick element of its own setting on every
photon-number block above the cutoff, and that the minimum over the settings
bounds it at every setting, under Assumption~\ref{ass:eig} only where a
computed block eigenvalue enters, at a setting governed by
Definition~\ref{def:pmcex}. Remark~\ref{rem:exacteig} explains why that entry
rests on the assumption rather than on a theorem, and
Theorem~\ref{thm:enclosed-weight} is the theorem that replaces it.
}
\begin{mainclaim}
\begin{theorem}[Multiclick bounds carried by the emitted weights]\label{thm:emitted-bound}
Let a setting family (Definition~\ref{def:setting-family}) be given. For
every setting $t\in\Tset$ let $\pmc^{(t)}$ be as in
Definition~\ref{def:pmc} if $t$ carries no exact-window depth and as in
Definition~\ref{def:pmcex} if it carries one, or, at any setting, let it be
the \termref{term-enclosed-weight}{enclosed weight} $\check p^{(t)}_{\mc}$ of
Definition~\ref{def:enclosed-weight} formed from \termref{term-enclosed-bound-data}{enclosed bound data} of
$t$ and, where $t$ carries a depth, from \termref{term-rigorous-window-bounds}{rigorous window bounds} of $t$,
or a real $p$ with $0\le p\le\check p^{(t)}_{\mc}$;
and let $\pmc=\min_{t\in\Tset}\pmc^{(t)}$ as in Definition~\ref{def:pmc}. Then
\begin{equation}\label{eq:pmcsetting}
  \blk{M}{\gtN}_{\mc|t}\ \succeq\ \pmc^{(t)}\,\id^{(>N)}
  \qquad\text{for every }t\in\Tset ,
\end{equation}
where $\blk{M}{\gtN}_{\mc|t}$ is the restriction of the multiclick element
$M_{\mc|t}$ of \eqref{eq:family-elements} at setting $t$ to the photon-number
blocks $n>N$; and consequently, with the one weight $\pmc$ at every setting,
\begin{equation}\label{eq:pmcstatement}
  \blk{M}{\gtN}_{\mc|t}\ \succeq\ \pmc\,\id^{(>N)}
  \qquad\text{for every }t\in\Tset .
\end{equation}
At a setting governed by Definition~\ref{def:pmc} or by
Definition~\ref{def:enclosed-weight} both inequalities hold at
that setting with no hypothesis beyond this contract's ledger; at a setting
governed by Definition~\ref{def:pmcex} they hold there \emph{under
Assumption~\ref{ass:eig}}, which is what makes the second entry of
\eqref{eq:elldef} a lower bound. At a
one-setting family, $\Tset=\{\off\}$, one has $\pmc=\pmc^{(\off)}$ and the
two inequalities are one. Moreover $0\le\pmc^{(t)}\le1$ at every setting,
hence $0\le\pmc\le1$, and each $\pmc^{(t)}$ of Definition~\ref{def:pmc}
or~\ref{def:pmcex} is computable by finitely many
arithmetic operations on
$\{\lmax(B^{(t)}),\lmin(B^{(t)}_k),\gamma^{(t)},\beta^{(t)}_k\}$ and the
declared $N$, $m_t$ and $m^{\ast}_t$ --- where $B^{(t)},B^{(t)}_k$ are the
operators of the setting (Definition~\ref{def:setting-family}\,(b)) ---
together with, where the setting carries a declared depth, the declared
$\varepsilon_{\mathrm{eig}}$ and the $m^{\ast}_t-N$ extreme eigenvalues
$\widehat\Bnd^{\ast(t)}(n)$ of the Hermitian matrices \eqref{eq:exactbnd},
each of dimension $\binom{D+n-1}{n}$; each enclosed weight by finitely many
on the enclosed bound data, the rigorous window bounds and the declared $N$,
$m_t$ and $m^{\ast}_t$; and $\pmc$ by finitely many comparisons among them.
\end{theorem}
\end{mainclaim}
\claimmeta{\metalabel{Uses}\mref{Thm.}{thm:multiclick-bound}, \mref{Lem.}{lem:collapse}, \mref{Lem.}{lem:windowtail}, \mref{Def.}{def:exactbnd}, \mref{Lem.}{lem:exactbasic}, \mref{Lem.}{lem:exactwindow}, \mref{Def.}{def:setting-family}, \mref{Def.}{def:pmc}, \mref{Def.}{def:pmcex}, \mref{Def.}{def:enclosed-weight}, \mref{Thm.}{thm:enclosed-weight}
\metasep \metalabel{Assumes}\aref{A-EIG}{ass:eig}
\metasep \metalabel{Used by}\mref{Cor.}{cor:emitted-unconditioned}, \mref{Prop.}{prop:slotvalid}, \mref{Rem.}{rem:tolerancedomain}, \mref{Thm.}{thm:exact-and-evaluated}, \extref{../reductions/contract.pdf}{reductions/wpfss/prop:route}, \extref{../reductions/contract.pdf}{reductions/wpfss/rem:minweight}, \extref{../reductions/contract.pdf}{reductions/wpfss/rem:receiver-imports}
\metasep newly derived
\metasep implemented (effective: accepted)
\metasep \metalabel{Code}\coderef{https://github.com/fraunhofersg/certiqs-key/blob/v0.1.0/src/certiqs_key/receiver/bound.py}{certiqs_key.receiver.bound:emitted_weights}
\metasep \metalabel{Oracles}\coderef{https://github.com/fraunhofersg/certiqs-key/blob/v0.1.0/tests/oracles/registry_receiver_povm.py}{registry_receiver_povm.py}}

\begin{proof}
\emph{Idea.} At a fixed setting, Theorem~\ref{thm:multiclick-bound} bounds the
multiclick element block by block. Lemma~\ref{lem:collapse} or
Lemma~\ref{lem:windowtail} turns the block bounds into the analytic weight,
and Lemma~\ref{lem:exactwindow} turns them into the exact-window weight;
Theorem~\ref{thm:enclosed-weight} does both for the enclosed weight. The
uniform statement follows because a minimum lies below each of its terms.

\emph{Step 1 (the analytic weight at a fixed setting).} Fix $t$. By
Theorem~\ref{thm:multiclick-bound}, applicable at every setting by
\eqref{eq:family-mc}, $\blk{M}{\eqn}_{\mc|t}\succeq\Bnd^{(t)}(n)\id$ for
every $n$, so
$\blk{M}{\gtN}_{\mc|t}\succeq\big(\inf_{n>N}\Bnd^{(t)}(n)\big)\id^{(>N)}$
--- the blocks are orthogonal summands and a lower bound holding on each
holds on the direct sum. Since $\blk{M}{\gtN}_{\mc|t}\succeq0$ as well
($M_{\mc|t}$ is a sum of POVM elements,
Definition~\ref{def:setting-family}\,(a)), the bound may be replaced by its
positive part. In the first branch of \eqref{eq:pmcdef} --- where the flag
is $\mathrm{exact}$, so that \eqref{eq:collapsehyp} holds with equality by
construction (Definition~\ref{def:setting-family}\,(c)) ---
Lemma~\ref{lem:collapse} identifies that positive part with
$\max\{0,\Bnd^{(t)}(N+1)\}$; in the second,
Lemma~\ref{lem:windowtail}\,(i) gives
$\underline{\Bnd}^{(t)}_{N,m_t}\le\inf_{n>N}\Bnd^{(t)}(n)$, so its positive
part is a valid lower bound as well.

\emph{Step 2 (the exact-window weight at a fixed setting).} At a setting
carrying a declared depth, \eqref{eq:exactstatement} of
Lemma~\ref{lem:exactwindow} gives the conclusion directly for
\eqref{eq:pmcex}; its hypothesis \eqref{eq:mcpos} holds at every setting,
because $M_{\mc|t}=\sum_{S\in\MCK}M_{S|t}$ is a sum of elements of a POVM
(Definition~\ref{def:setting-family}\,(a)), which is the discharge
Remark~\ref{rem:exactpositivity} says is needed. On branch~1, apply it with
$m^{\ast}=m^{\ast}_t$, $\ell=\ell^{(t)}$ and
$\tau=\Bnd^{(t)}(m^{\ast}_t+1)$, the remainder being admissible by
part~(ii) because predicate~(X) makes \eqref{eq:collapsehyp} hold. On
branch~2, apply it with $m^{\ast}=m_t$, $\tau=T^{(t)}_{m_t}$ admissible by
part~(i), and the window function $\ell(n)=\ell^{(t)}(n)$ for
$N<n\le m^{\ast}_t$, $\ell(n)=\Bnd^{(t)}(n)$ for $m^{\ast}_t<n\le m_t$ ---
the two ranges together being $N<n\le m_t$ because $m^{\ast}_t\le m_t$ is
required there. In both cases hypothesis \eqref{eq:ellhyp} holds entry by
entry: for $\Bnd^{(t)}(n)$ by Lemma~\ref{lem:exactbasic}, and for
$\widehat\Bnd^{\ast(t)}(n)-\varepsilon_{\mathrm{eig}}$ by
Assumption~\ref{ass:eig}; a maximum of the two is admissible by
part~(iii).

\emph{Step 2a (the enclosed weight at a fixed setting).} At a setting whose
weight is the enclosed weight $\check p^{(t)}_{\mc}$ or a real $p$ with
$0\le p\le\check p^{(t)}_{\mc}$,
Theorem~\ref{thm:enclosed-weight}\,(iii) gives \eqref{eq:pmcsetting} at $t$
and $0\le\pmc^{(t)}\le1$, for the enclosed weight and for every real between
$0$ and it, with no hypothesis beyond Definition~\ref{def:setting-family},
whether or not $t$ carries a depth.

\emph{Step 3 (the uniform weight).} So \eqref{eq:pmcsetting} holds at $t$,
under Assumption~\ref{ass:eig} exactly where $t$ is governed by
Definition~\ref{def:pmcex}. Since
$\pmc$ is the minimum of the $\pmc^{(s)}$ over the finite set $\Tset$,
$\pmc\le\pmc^{(t)}$, and $\pmc^{(t)}\id^{(>N)}\succeq\pmc\,\id^{(>N)}$ turns
\eqref{eq:pmcsetting} at $t$ into \eqref{eq:pmcstatement} at $t$, with no
further hypothesis; at $\Tset=\{\off\}$ the minimum is over one number.

\emph{Step 4 (range and computability).} The range $0\le\pmc^{(t)}\le1$
follows at every setting from the positive part and from an entry at most $1$
in the inner minimum. At a setting governed by Definition~\ref{def:pmc} every
entry is a value of $\Bnd^{(t)}$ or the tail value $T^{(t)}_{m_t}$, and
$\Bnd\le1$ (Theorem~\ref{thm:multiclick-bound}). At a setting governed by
Definition~\ref{def:pmcex} the window entries $\ell^{(t)}(n)$ may exceed $1$,
but the inner minimum of each branch contains an analytic remainder at most
$1$: $\Bnd^{(t)}(m^{\ast}_t+1)\le1$ on branch~1, and
$T^{(t)}_{m_t}=1-\sum_k\beta^{(t)}_k\big(d^{(t)}_k\big)^{m_t}\le1$ on
branch~2 because $\beta^{(t)}_k,d^{(t)}_k\ge0$ (the \emph{Range} paragraph of
Remark~\ref{rem:exacteig}). At a setting governed by
Definition~\ref{def:enclosed-weight} the range is Step~2a. And $\pmc$, a
minimum of finitely many
numbers in $[0,1]$, lies there with them. Computability of each $\pmc^{(t)}$
is the finiteness of its window, and of $\pmc$ the finiteness of $\Tset$.
\end{proof}

% claim receiver/cor:emitted-unconditioned
\claimsummary{receiver/cor:emitted-unconditioned}{For a receiver without probe modes the one emitted weight bounds the multiclick element on every block above the cutoff and lies in [0, 1], under \hyperlink{claim-receiver/ass:eig}{A-EIG} only where a computed block eigenvalue enters it. The slot entry of the multiclick set meets the guarantee of \hyperlink{claim-receiver/def:slot}{Definition~\ref*{def:slot}}.}
\claimbookmark{receiver/cor:emitted-unconditioned}{Corollary}{cor:emitted-unconditioned}{Emitted multiclick weight of a probe-free receiver}
\begin{corollary}[Emitted multiclick weight of a probe-free receiver]
\label{cor:emitted-unconditioned}
Let $\Recv$ be a \termref{term-receiver-instrument}{receiver instrument} with $\InpS=\varnothing$
(Definition~\ref{def:instrument}), let $N$ be the declared cutoff, and let
$\pmc$ be computed by Definition~\ref{def:pmc}, or by
Definition~\ref{def:pmcex} where an exact-window depth is declared, from the
one-setting family of Corollary~\ref{cor:unconditioned-family}, or let it be
a real with $0\le\pmc\le\check p_{\mc}$, where $\check p_{\mc}$ is an
\termref{term-enclosed-weight}{enclosed weight} of that family's one setting
(Definition~\ref{def:enclosed-weight}), $\check p_{\mc}$ itself included. Then
\begin{equation}\label{eq:emitted-unconditioned}
  \blk{M}{\gtN}_{\mc}\ \succeq\ \pmc\,\id^{(>N)},\qquad 0\le\pmc\le1,
\end{equation}
with no hypothesis beyond the ledger where $\pmc$ is given by
Definition~\ref{def:pmc} or~\ref{def:enclosed-weight} and under
Assumption~\ref{ass:eig} where it is given by Definition~\ref{def:pmcex}; and
the strengthening slot
$\{(\MCK,\pmc)\}$ satisfies the guarantee of Definition~\ref{def:slot}.
\end{corollary}
\claimmeta{\metalabel{Uses}\mref{Def.}{def:instrument}, \mref{Cor.}{cor:unconditioned-family}, \mref{Def.}{def:pmc}, \mref{Def.}{def:pmcex}, \mref{Def.}{def:enclosed-weight}, \mref{Thm.}{thm:emitted-bound}, \mref{Def.}{def:slot}
\metasep \metalabel{Assumes}\aref{A-TRUST}{ass:trust}, \aref{A-PASSIVE}{ass:passive}, \aref{A-VAC}{ass:vac}, \aref{A-MODES}{ass:modes}, \aref{A-THRESH}{ass:thresh}, \aref{A-PORT}{ass:port}, \aref{A-EXACT}{ass:exact}, \aref{A-NOSEL}{ass:nosel}, \aref{A-ROUND}{ass:round}, \aref{A-EIG}{ass:eig}
\metasep \metalabel{Used by}\mref{Rem.}{rem:normalization-gain}
\metasep newly derived
\metasep implemented (effective: accepted)
\metasep \metalabel{Code}\coderef{https://github.com/fraunhofersg/certiqs-key/blob/v0.1.0/src/certiqs_key/receiver/bound.py}{certiqs_key.receiver.bound:emitted_weight}}

\begin{proof}
By Corollary~\ref{cor:unconditioned-family} the receiver is a setting
family with $\Tset=\{\off\}$ and $M_{\mc|\off}=M_{\mc}$, so
Theorem~\ref{thm:emitted-bound} applies and \eqref{eq:pmcstatement} at the one
setting is \eqref{eq:emitted-unconditioned}; the range and the hypothesis
clause are the theorem's. The slot guarantee is
\eqref{eq:emitted-unconditioned} read as $\sum_{S\in\MCK}\blk{M}{\gtN}_S
\succeq\pmc\,\id^{(>N)}$ at the one setting.
\end{proof}

% claim receiver/def:slot
\claimsummary{receiver/def:slot}{The artifact carries high-photon lower bounds on disjoint sets of raw outcomes, one weight per set and setting; a set's uniform weight is its minimum over the settings. The receiver carries one entry, the multiclick set with the per-setting multiclick weights.}
\claimbookmark{receiver/def:slot}{Definition}{def:slot}{The strengthening slot}
\claimleadin{%
The slot is indexed by sets of raw outcomes rather than by emitted letters:
lower bounds on disjoint raw outcomes add, while a bound on a merged outcome
implies nothing about its parts (Theorem~\ref{thm:coarse-graining}).
}
\begin{definition}[The strengthening slot]\label{def:slot}
The
\termdef{term-strengthening-slot}{strengthening slot}
is the finite list of pairs in which the artifact carries the high-photon
evidence,
\begin{equation}\label{eq:slot}
  \Big\{\big(\mathcal{O}_i,\ (p^{(t)}_i)_{t\in\Tset}\big)\Big\}_{i=1}^{r},
  \qquad \mathcal{O}_i\subseteq2^{[K]}\ \text{pairwise disjoint},\quad
  p^{(t)}_i\ge0,
\end{equation}
one weight per entry and per setting $t\in\Tset$ of the setting family
(Definition~\ref{def:setting-family}), with the guarantee
$\sum_{S\in\mathcal{O}_i}\blk{M}{\gtN}_{S|t}\succeq p^{(t)}_i\,\id^{(>N)}$
for every setting $t\in\Tset$ and every $i$. The
\termdef{term-uniform-weight-of-an-entry}{uniform weight of an entry}
$i$ is $p_i\coloneqq\min_{t\in\Tset}p^{(t)}_i$; since
$0\le p_i\le p^{(t)}_i$, it carries the guarantee
$\sum_{S\in\mathcal{O}_i}\blk{M}{\gtN}_{S|t}\succeq p_i\,\id^{(>N)}$ for
every setting $t\in\Tset$ and every $i$, one weight at every setting. At a
one-setting family, $\Tset=\{\off\}$, an entry has the single weight
$p^{(\off)}_i=p_i$ and is written $(\mathcal{O}_i,p_i)$. The artifact
carries $r=1$ with
$(\mathcal{O}_1,(p^{(t)}_1)_{t\in\Tset})=(\MCK,(\pmc^{(t)})_{t\in\Tset})$,
the per-setting weights of Theorem~\ref{thm:emitted-bound}, each given at its
setting by Definition~\ref{def:pmc} or~\ref{def:pmcex} or as a non-negative
real at most the enclosed weight of Definition~\ref{def:enclosed-weight},
whose uniform weight is $\pmc$; at a
one-setting family that entry is $(\MCK,\pmc)$.
\end{definition}
\claimmeta{\metalabel{Uses}\mref{Def.}{def:setting-family}
\metasep \metalabel{Used by}\mref{Cor.}{cor:emitted-unconditioned}, \mref{Prop.}{prop:slotvalid}, \mref{Ass.}{ass:eig}, \mref{Req.}{req:rigorous-weight}, \extref{../reductions/contract.pdf}{reductions/req:recv-row-provenance}, \extref{../reductions/contract.pdf}{reductions/wpfss/prop:route}, \extref{../reductions/contract.pdf}{reductions/wpfss/rem:absorbneeded}, \extref{../reductions/contract.pdf}{reductions/wpfss/rem:minweight}, \extref{../reductions/contract.pdf}{reductions/wpfss/rem:receiver-imports}
\metasep newly derived
\metasep implemented (effective: accepted)
\metasep \metalabel{Code}\coderef{https://github.com/fraunhofersg/certiqs-key/blob/v0.1.0/src/certiqs_key/receiver/bound.py}{certiqs_key.receiver.bound:StrengtheningSlot}}

% claim receiver/prop:slotvalid
\claimsummary{receiver/prop:slotvalid}{At every setting the strengthening slot's weights (high-photon lower bounds on disjoint sets of raw outcomes) are non-negative and sum to at most one, as are their minima over the settings, and each emitted multiclick weight lies in [0, 1]. That the emitted weights are valid bounds is \hyperlink{claim-receiver/thm:emitted-bound}{Theorem~\ref*{thm:emitted-bound}}.}
\claimbookmark{receiver/prop:slotvalid}{Proposition}{prop:slotvalid}{Range of the strengthening-slot weights}
\claimleadin{%
Among the conditions under which a squashing request fails are a negative
weight, weights summing to more than one, and a bound that does not hold
where it is stated (\href{../../../spec/architecture.md\#54-failure-policy}{\nolinkurl{spec/architecture.md}\texttt{\#54-failure-policy}}). The proposition rules
out the first two for every slot and names the result that rules out the third for the slot this module emits,
under Assumption~\ref{ass:eig} only where a computed block eigenvalue enters a
weight.
}
\begin{mainclaim}
\begin{proposition}[Range of the strengthening-slot weights]
\label{prop:slotvalid}
For any list \eqref{eq:slot} satisfying the stated guarantee with pairwise
disjoint $\mathcal{O}_i$, one has at every setting $t\in\Tset$ that
$p^{(t)}_i\ge0$ for every $i$ and $\sum_{i}p^{(t)}_i\le1$; hence the uniform
weights satisfy $p_i\ge0$ for every $i$ and $\sum_{i}p_i\le1$. In particular
the list this module emits --- by Corollary~\ref{cor:emitted-unconditioned}
and, for a setting family a submodule constructs, by that submodule's
instance of Theorem~\ref{thm:emitted-bound} --- satisfies
$\pmc^{(t)}\in[0,1]$ at every setting, so the failure conditions
``$p_y<0$'' and ``$\sum_yp_y>1$'' of a squashing request cannot be triggered
by this module's output; and what the further failure condition of a
squashing request, ``a bound does not hold at every declared setting it is
stated for and every photon number above the cutoff'', requires to hold --- each $\pmc^{(t)}$ at
its own setting $t$ and $\pmc$ at every setting, for every photon number
above the cutoff --- is Theorem~\ref{thm:emitted-bound}.
\end{proposition}
\end{mainclaim}
\claimmeta{\metalabel{Uses}\mref{Def.}{def:setting-family}, \mref{Thm.}{thm:emitted-bound}, \mref{Def.}{def:slot}
\metasep \metalabel{Used by}\extref{../reductions/contract.pdf}{reductions/wpfss/prop:route}, \extref{../reductions/contract.pdf}{reductions/wpfss/rem:receiver-imports}
\metasep newly derived
\metasep accepted}

\begin{proof}
Non-negativity is part of \eqref{eq:slot}. For the sum, fix any $t$ and any
unit vector $\ket{\psi}$ in a block $n>N$. Since the $\mathcal{O}_i$ are
pairwise disjoint subsets of $2^{[K]}$ and $\{M_{S|t}\}$ is a POVM
(Definition~\ref{def:setting-family}\,(a)),
$\sum_i\sum_{S\in\mathcal{O}_i}\blk{M}{\eqn}_{S|t}\preceq
\sum_{S\subseteq[K]}\blk{M}{\eqn}_{S|t}=\id^{(=n)}$, so by the guarantee at
$t$,
$\sum_ip^{(t)}_i\le\sum_i\bra{\psi}\sum_{S\in\mathcal{O}_i}\blk{M}{\eqn}_{S|t}\ket{\psi}
\le\bra{\psi}\id\ket{\psi}=1$. The uniform weights satisfy
$0\le p_i\le p^{(t)}_i$ (Definition~\ref{def:slot}), so $p_i\ge0$ and
$\sum_ip_i\le\sum_ip^{(t)}_i\le1$ at any $t$.
\end{proof}

% claim receiver/rem:notsharp
\claimsummary{receiver/rem:notsharp}{The analytic bound bounds three operators separately and need not be tight, while the exact block eigenvalue covers one block at a time. The exact window combines them: computed block eigenvalues less the slack, floored at the analytic bound, on finitely many blocks, an analytic remainder above.}
\claimbookmark{receiver/rem:notsharp}{Remark}{rem:notsharp}{Slack of the analytic bound against the block eigenvalue}
\begin{remark}[Slack of the analytic bound against the block eigenvalue]\label{rem:notsharp}
$\Bnd(n)$ bounds three operators separately and so need not be tight
(Example~\ref{ex:bb84exact} exhibits a strict gap at $n=2$): on
a given block, $\lmin(\blk{M}{\eqn}_{\mc})$ can be computed exactly from
\eqref{eq:blockMS} by diagonalizing a matrix of dimension
$\binom{|\InpC|+n-1}{n}$, which is what Definition~\ref{def:exactbnd} names
$\Bnd^{\ast}(n)$. The analytic bound is what the module emits where no depth
is declared, because it is uniform in $n$ --- by Lemma~\ref{lem:collapse}
where predicate~(X) of Definition~\ref{def:pmc} holds and by
Lemma~\ref{lem:windowtail}\,(i) where it does not --- while a per-block
computation covers finitely many blocks. The exact window does not answer
that by computing more blocks: on the finitely many blocks up to a declared
depth it uses the computed block eigenvalues less the eigenvalue slack,
floored at the analytic bound, and it keeps an analytic remainder above them
(Lemma~\ref{lem:exactwindow}, Definition~\ref{def:pmcex}). Which of the two
evaluations enters the artifact, and what it records beside each weight, is
Requirement~\ref{req:weight-provenance}.
\end{remark}
\claimmeta{\metalabel{Uses}\mref{Thm.}{thm:click-povm}, \mref{Lem.}{lem:collapse}, \mref{Lem.}{lem:windowtail}, \mref{Def.}{def:exactbnd}, \mref{Lem.}{lem:exactwindow}, \mref{Def.}{def:pmcex}
\metasep \metalabel{Used by}\mref{Req.}{req:weight-provenance}
\metasep newly derived
\metasep accepted}

\section{The reference, enclosed and exact}\label{group-receiver/def:complex-enclosure}
\sectionintro{The reference a run hands beside the snap of the emitted blocks, and evaluates its first-order bound on, must be known through rational boxes, and the binary64 evaluation is not one. \termref{term-complex-enclosure}{Complex enclosures}, a Gaussian-rational centre with a rational radius, are carried through sums, products and reciprocals by explicit rules (\hyperlink{claim-receiver/lem:enclosure-arithmetic}{Lemma~\ref*{lem:enclosure-arithmetic}}). An entry of a lifted operator in the occupation basis is a permanent divided by square roots of factorials (\hyperlink{claim-receiver/lem:fock-permanent}{Lemma~\ref*{lem:fock-permanent}}), so enclosures of the port Grams, themselves obtained from enclosures of square roots, cosines and sines, enclose every block, and the emitted blocks up to the cutoff are a reference for any declared parameters (\hyperlink{claim-receiver/thm:enclosed-blocks}{Theorem~\ref*{thm:enclosed-blocks}}). Where the port Grams are Gaussian-rational, the \termref{term-exact-route}{exact route} (\hyperlink{claim-receiver/def:exact-route}{Definition~\ref*{def:exact-route}}), every block of at most one photon is known exactly, which a snap needs in order to keep a block unmoved (\hyperlink{claim-receiver/thm:exact-blocks}{Theorem~\ref*{thm:exact-blocks}}). Read through a \termref{term-port-dilation}{port dilation}, an isometry from the inputs that reproduces the port Grams (\hyperlink{claim-receiver/def:port-dilation}{Definition~\ref*{def:port-dilation}}), every block is also a weighted sum of rank-one terms whose columns are enclosed the same way, a factor enclosure, which fixes the block's rank where those columns are certified independent (\hyperlink{claim-receiver/thm:factor-blocks}{Theorem~\ref*{thm:factor-blocks}}).}
\sectionimports{\extref{../foundations/contract.pdf}{foundations/def:binary64}, \extref{../foundations/contract.pdf}{foundations/def:exact-reference}, \extref{../foundations/contract.pdf}{foundations/def:factor-enclosure}, \extref{../foundations/contract.pdf}{foundations/def:sector-map}}

% claim receiver/def:complex-enclosure
\claimsummary{receiver/def:complex-enclosure}{A complex enclosure of a number is a Gaussian-rational centre with a rational radius that bounds its distance from the number; a matrix is enclosed entry by entry. Square roots, cosines and sines are enclosed by any such pair with a rational centre, exactly where the value is plainly rational.}
\claimbookmark{receiver/def:complex-enclosure}{Definition}{def:complex-enclosure}{Complex enclosures, and enclosures of square roots and angles}
\claimleadin{%
A reference of a measurement family is read through rational boxes around its
entries (\xref{foundations/def:exact-reference}). The receiver's
photon-number blocks are built from the declared parameters by finitely many
arithmetic operations, square roots, cosines and sines, so a box around each
block follows from boxes around those square roots and trigonometric values,
carried through the arithmetic (Theorem~\ref{thm:enclosed-blocks}). The
definition fixes what such a box must satisfy and leaves open how it is
computed; a series with a bounded remainder, or bisection on the squares for a
square root, are ways to obtain one.
}
\begin{definition}[Complex enclosures, and enclosures of square roots and angles]\label{def:complex-enclosure}
Write $\mathbb{Q}(i)=\{x+iy:x,y\in\mathbb{Q}\}$ and, for $z\in\mathbb{Q}(i)$,
$\lvert z\rvert_{1}\coloneqq\lvert\mathrm{Re}\,z\rvert+\lvert\mathrm{Im}\,z\rvert$,
a rational with $\lvert z\rvert\le\lvert z\rvert_{1}$. A
\termdef{term-complex-enclosure}{complex enclosure} of a complex number $x$ is a pair $(x_c,r)$ of
$x_c\in\mathbb{Q}(i)$, its \emph{centre}, and a rational $r\ge0$, its
\emph{radius}, with $\lvert x-x_c\rvert\le r$. It is \emph{exact} when $r=0$,
and then $x=x_c$. A complex enclosure of a complex matrix $X$ is a pair
$(X_c,r)$ of a matrix $X_c$ of the same shape with entries in
$\mathbb{Q}(i)$ and a matrix $r$ of non-negative rationals such that
$\big((X_c)_{il},r_{il}\big)$ is a complex enclosure of $X_{il}$ for all
$i,l$.

A \termdef{term-square-root-enclosure}{square-root enclosure} of a rational $p\ge0$ is a complex enclosure
of $\sqrt p$ with a rational centre, exact where $p$ is the square of a
rational. An
\termdef{term-angle-enclosure}{angle enclosure} of a real number $\theta$ that is
given either as a rational, such as a \exttermref{../foundations/contract.pdf}{term-declared-value}{declared value}
(\xref{foundations/def:binary64}\,(iii)), or as $c\pi$ with $c$ a rational,
is a pair of complex enclosures with rational centres of $\cos\theta$ and of
$\sin\theta$, both exact where $\theta=c\pi$ with $2c$ an integer.
\end{definition}
\claimmeta{\metalabel{Uses}\extref{../foundations/contract.pdf}{foundations/def:binary64}
\metasep \metalabel{Used by}\mref{Lem.}{lem:enclosure-arithmetic}, \mref{Thm.}{thm:enclosed-blocks}, \mref{Thm.}{thm:factor-blocks}, \mref{Thm.}{thm:exact-blocks}, \mref{Def.}{def:evaluation-model}, \mref{Thm.}{thm:exact-and-evaluated}, \mref{Req.}{req:rigorous-weight}, \mref{Req.}{req:exact-route-decision}, \mref{Cor.}{cor:exact-dark-composed}, \mref{Rem.}{rem:noise-floor-receiver}, \mref{Req.}{req:dark-composed-reference}, \mref{Req.}{req:snap-fix}, \extref{../reductions/contract.pdf}{reductions/cor:reference-companion}
\metasep newly derived
\metasep implemented
\metasep \metalabel{Code}\coderef{https://github.com/fraunhofersg/certiqs-key/blob/v0.1.0/src/certiqs_key/receiver/enclosure.py}{certiqs_key.receiver.enclosure:Ball}}

% claim receiver/lem:enclosure-arithmetic
\claimsummary{receiver/lem:enclosure-arithmetic}{Sums, products, conjugates and reciprocals of enclosed numbers are enclosed by explicit rational rules, so any finite expression in enclosed numbers is enclosed and stays exact where its inputs are. Averaging a Hermitian matrix's enclosure with its adjoint makes the centre Hermitian.}
\claimbookmark{receiver/lem:enclosure-arithmetic}{Lemma}{lem:enclosure-arithmetic}{Arithmetic on complex enclosures}
\claimleadin{%
This is interval arithmetic on complex discs with Gaussian-rational centres,
written so that every rule is computed in exact rational arithmetic. A rule
never shrinks a radius, and an enclosure obtained by it can be far wider than
the true uncertainty when the same quantity enters an expression twice;
soundness is all the lemma asserts.
}
\begin{technicalclaim}
\begin{lemma}[Arithmetic on complex enclosures]\label{lem:enclosure-arithmetic}
Let $(x_c,r_x)$ and $(y_c,r_y)$ be \termref{term-complex-enclosure}{complex
enclosures} of complex numbers $x$ and $y$
(Definition~\ref{def:complex-enclosure}).
\begin{enumerate}[label=(\roman*),leftmargin=2.4em]
\item \emph{Sums and conjugates.} $(x_c+y_c,\,r_x+r_y)$,
  $(-x_c,\,r_x)$ and $(\overline{x_c},\,r_x)$ are complex enclosures of
  $x+y$, $-x$ and $\overline{x}$.
\item \emph{Products.}
  $\big(x_cy_c,\ \lvert x_c\rvert_{1}\,r_y+r_x\,\lvert y_c\rvert_{1}+r_xr_y\big)$
  is a complex enclosure of $xy$.
\item \emph{Reciprocals.} If $\ell$ is a rational with $\ell>0$,
  $\ell^{2}\le\lvert x_c\rvert^{2}$ and $r_x<\ell$, then $x\neq0$, $x_c\neq0$,
  $1/x_c=\overline{x_c}/\lvert x_c\rvert^{2}$ lies in $\mathbb{Q}(i)$, and
  $\big(1/x_c,\ r_x/(\ell(\ell-r_x))\big)$ is a complex enclosure of $1/x$.
\item \emph{Expressions.} Let $E$ be a finite expression whose leaves are
  complex numbers and whose nodes are sums, differences, products,
  conjugates and reciprocals, and let a complex enclosure of every leaf be
  given. Assign to every node, in an order in which its operands come first,
  the enclosure that (i)--(iii) give from the enclosures of its operands, a
  reciprocal node receiving one only where some rational $\ell$ meets the
  condition of (iii). If every reciprocal node receives one, the enclosure of
  every node is a complex enclosure of its value, and a node whose operands
  have exact enclosures has an exact one. This applies entry by entry to sums
  and products of matrices, to permanents and determinants, which are sums of
  products of entries, and to the inverse of an invertible matrix, whose
  entries are signed minors divided by its determinant.
\item \emph{Hermitian centres.} If $X$ is a Hermitian $d\times d$ matrix and
  $(X_c,r)$ a complex enclosure of $X$, then
  $\big(\tfrac12(X_c+X_c^{\dagger}),\,\tfrac12(r+r^{\transpose})\big)$ is a
  complex enclosure of $X$ whose centre is Hermitian with entries in
  $\mathbb{Q}(i)$ and whose radii are symmetric; it is therefore an
  \exttermref{../foundations/contract.pdf}{term-entrywise-enclosure}{entrywise enclosure} of $X$
  (\xref{foundations/def:exact-reference}), and its entry $(i,l)$ is exact
  where $r_{il}=r_{li}=0$.
\end{enumerate}
\end{lemma}
\end{technicalclaim}
\claimmeta{\metalabel{Uses}\mref{Def.}{def:complex-enclosure}, \extref{../foundations/contract.pdf}{foundations/def:exact-reference}
\metasep \metalabel{Used by}\mref{Thm.}{thm:enclosed-blocks}, \mref{Thm.}{thm:factor-blocks}, \mref{Thm.}{thm:exact-blocks}, \mref{Thm.}{thm:exact-and-evaluated}, \mref{Req.}{req:rigorous-weight}, \mref{Req.}{req:exact-route-decision}, \mref{Cor.}{cor:exact-dark-composed}, \mref{Req.}{req:dark-composed-reference}, \mref{Req.}{req:snap-fix}, \extref{../reductions/contract.pdf}{reductions/cor:reference-companion}
\metasep newly derived
\metasep implemented
\metasep \metalabel{Code}\coderef{https://github.com/fraunhofersg/certiqs-key/blob/v0.1.0/src/certiqs_key/receiver/enclosure.py}{certiqs_key.receiver.enclosure:Ball}}

\begin{proof}
(i) $\lvert(x+y)-(x_c+y_c)\rvert\le\lvert x-x_c\rvert+\lvert y-y_c\rvert
\le r_x+r_y$; negation and conjugation preserve moduli, and
$\mathbb{Q}(i)$ is closed under both.

(ii) $xy-x_cy_c=x_c(y-y_c)+(x-x_c)y_c+(x-x_c)(y-y_c)$, so
$\lvert xy-x_cy_c\rvert\le\lvert x_c\rvert r_y+r_x\lvert y_c\rvert+r_xr_y$,
and $\lvert z\rvert\le\lvert z\rvert_{1}$ for $z\in\mathbb{Q}(i)$. The centre
lies in $\mathbb{Q}(i)$ and the radius is a non-negative rational.

(iii) $\lvert x_c\rvert\ge\ell>0$, so $x_c\neq0$ and
$1/x_c=\overline{x_c}/\lvert x_c\rvert^{2}$, with $\lvert x_c\rvert^{2}$ a
positive rational. Further $\lvert x\rvert\ge\lvert x_c\rvert-r_x\ge\ell-r_x>0$,
so $x\neq0$ and
\begin{equation*}
  \Big\lvert\frac1x-\frac1{x_c}\Big\rvert
  =\frac{\lvert x_c-x\rvert}{\lvert x\rvert\,\lvert x_c\rvert}
  \le\frac{r_x}{(\ell-r_x)\,\ell} ,
\end{equation*}
a non-negative rational.

(iv) By induction along the order of assignment: a leaf has a complex
enclosure by hypothesis, and a node whose operands have complex enclosures
receives one by (i)--(iii). If every operand radius is zero, the radii of (i)
and (ii) are zero, and so is that of (iii), its numerator being $r_x=0$. A
sum or product of matrices has entries that are sums of products of entries; a
permanent or determinant is a sum of products of entries, with signs for the
determinant; and Cramer's rule writes every entry of an inverse as a signed
minor, a determinant, times the reciprocal of the determinant of the matrix.

(v) The centre
$Z\coloneqq\tfrac12(X_c+X_c^{\dagger})$ is Hermitian with entries in
$\mathbb{Q}(i)$. Since $X_{il}=\overline{X_{li}}$,
\begin{equation*}
  \lvert X_{il}-Z_{il}\rvert
  =\tfrac12\big\lvert(X_{il}-(X_c)_{il})
     +\overline{(X_{li}-(X_c)_{li})}\big\rvert
  \le\tfrac12\,(r_{il}+r_{li}) ,
\end{equation*}
and $\tfrac12(r+r^{\transpose})$ is a symmetric matrix of non-negative
rationals. These are
the data of an entrywise enclosure of the Hermitian matrix $X$, and the radius
of entry $(i,l)$ is zero where $r_{il}=r_{li}=0$.
\end{proof}

% claim receiver/lem:fock-permanent
\claimsummary{receiver/lem:fock-permanent}{An entry of the n-photon block of a lifted operator is a permanent of a submatrix with repeated rows and columns, divided by the square roots of two products of factorials. For a Gaussian-rational operator the entry is Gaussian-rational on the diagonal and on every block of at most one photon.}
\claimbookmark{receiver/lem:fock-permanent}{Lemma}{lem:fock-permanent}{Entries of a lifted operator in the occupation basis}
\begin{technicalclaim}
\begin{lemma}[Entries of a lifted operator in the occupation basis]\label{lem:fock-permanent}
Let $I$ be a finite mode set with orthonormal basis $\{\ket{e_i}\}_{i\in I}$,
fix an order on $I$, and let $X$ be an operator on $\mathbb{C}^{I}$ and
$n\ge0$. For an occupation vector $\vec m\in\mathbb{N}_0^{I}$ with
$|\vec m|=n$ let $w_{\vec m}=(w_{\vec m}(1),\dots,w_{\vec m}(n))$ be the word
listing each $i\in I$ exactly $m_i$ times in the order of $I$, write
$\vec m!\coloneqq\prod_{i\in I}m_i!$, and for two such vectors let
$X[\vec m|\vec m']$ be the $n\times n$ matrix with entries
$X_{w_{\vec m}(r),\,w_{\vec m'}(s)}$. Write
$\operatorname{per}Y\coloneqq\sum_{\pi}\prod_{r=1}^{n}Y_{r,\pi(r)}$, the sum
over the permutations of $\{1,\dots,n\}$, with $\operatorname{per}Y=1$ for
the empty matrix.
\begin{enumerate}[label=(\roman*),leftmargin=3.0em]
\item In the occupation basis of Definition~\ref{def:fock},
  \begin{equation}\label{eq:fock-permanent}
    \bra{\vec m}\Fock(X)^{(n)}\ket{\vec m'}
    =\frac{\operatorname{per}X[\vec m|\vec m']}{\sqrt{\vec m!\,\vec m'!}} .
  \end{equation}
\item If $X$ has entries in $\mathbb{Q}(i)=\{x+iy:x,y\in\mathbb{Q}\}$, then
  $\operatorname{per}X[\vec m|\vec m']$ lies in $\mathbb{Q}(i)$, and the entry
  \eqref{eq:fock-permanent} lies in $\mathbb{Q}(i)$ wherever
  $\vec m!\,\vec m'!$ is the square of an integer; in particular on the
  diagonal, $\vec m=\vec m'$, and for every $\vec m,\vec m'$ when $n\le1$.
\end{enumerate}
\end{lemma}
\end{technicalclaim}
\claimmeta{\metalabel{Uses}\mref{Def.}{def:fock}, \mref{Lem.}{lem:fock-functor}
\metasep \metalabel{Used by}\mref{Lem.}{lem:fock-permanent-maps}, \mref{Thm.}{thm:enclosed-blocks}, \mref{Thm.}{thm:exact-blocks}
\metasep newly derived
\metasep implemented (effective: accepted)
\metasep \metalabel{Code}\coderef{https://github.com/fraunhofersg/certiqs-key/blob/v0.1.0/src/certiqs_key/receiver/generic.py}{certiqs_key.receiver.generic:fock_block}}

\begin{proof}
(i) Write $P_{\Symn{n}}=\frac{1}{n!}\sum_\pi\Pi_\pi$ for the orthogonal
projector onto $\Symn{n}(\mathbb{C}^{I})$, $\Pi_\pi$ permuting the tensor
factors, and $\ket{w_{\vec m}}\coloneqq\ket{e_{w_{\vec m}(1)}}\otimes\cdots
\otimes\ket{e_{w_{\vec m}(n)}}$. The permutations with
$\Pi_\pi\ket{w_{\vec m}}=\ket{w_{\vec m}}$ are those that permute equal
letters among themselves, of which there are $\vec m!$; every other
$\Pi_\pi\ket{w_{\vec m}}$ is a product basis vector orthogonal to
$\ket{w_{\vec m}}$. Hence
$\lVert P_{\Symn{n}}\ket{w_{\vec m}}\rVert^2
=\bra{w_{\vec m}}P_{\Symn{n}}\ket{w_{\vec m}}=\vec m!/n!$, and the
symmetrized and normalized vector of Definition~\ref{def:fock} is
$\ket{\vec m}=\sqrt{n!/\vec m!}\,P_{\Symn{n}}\ket{w_{\vec m}}$. The block
$\Fock(X)^{(n)}$ is the restriction of $X^{\otimes n}$ to
$\Symn{n}(\mathbb{C}^{I})$, and $X^{\otimes n}$ commutes with
$P_{\Symn{n}}$ (proof of Lemma~\ref{lem:fock-functor}\,(i)), so
\begin{equation*}
  \bra{\vec m}X^{\otimes n}\ket{\vec m'}
  =\frac{n!}{\sqrt{\vec m!\,\vec m'!}}
   \bra{w_{\vec m}}X^{\otimes n}P_{\Symn{n}}\ket{w_{\vec m'}}
  =\frac{1}{\sqrt{\vec m!\,\vec m'!}}\sum_\pi
   \bra{w_{\vec m}}X^{\otimes n}\Pi_\pi\ket{w_{\vec m'}} ,
\end{equation*}
using $P_{\Symn{n}}^{\dagger}=P_{\Symn{n}}=P_{\Symn{n}}^{2}$. The vector
$\Pi_\pi\ket{w_{\vec m'}}$ is the product vector whose $r$-th factor is
$\ket{e_{w_{\vec m'}(\sigma(r))}}$ for a permutation $\sigma$ determined by
$\pi$, and $\pi\mapsto\sigma$ is a bijection of the permutation group. The
summand is therefore $\prod_r\bra{e_{w_{\vec m}(r)}}X\ket{e_{w_{\vec m'}(\sigma(r))}}
=\prod_rX_{w_{\vec m}(r),\,w_{\vec m'}(\sigma(r))}$, and the sum over
$\sigma$ is $\operatorname{per}X[\vec m|\vec m']$. At $n=0$ both sides are
$1$.

(ii) The permanent is a sum of products of entries of $X$ with coefficient
$1$, so it lies in $\mathbb{Q}(i)$, which is closed under sums and products.
If $\vec m!\,\vec m'!=s^{2}$ with $s$ a positive integer, the entry
\eqref{eq:fock-permanent} is $\operatorname{per}X[\vec m|\vec m']/s$, which
lies in $\mathbb{Q}(i)$. On the diagonal $\vec m!\,\vec m'!=(\vec m!)^{2}$,
and for $n\le1$ every factorial in $\vec m!$ and $\vec m'!$ is $0!$ or $1!$,
so $\vec m!\,\vec m'!=1$.
\end{proof}

% claim receiver/lem:fock-permanent-maps
\claimsummary{receiver/lem:fock-permanent-maps}{The Fock lift of a linear map between two mode sets respects composition and adjoints and keeps an isometry an isometry, and its entries in the occupation bases are permanents divided by square roots of factorials. Reading a block through a port dilation needs both.}
\claimbookmark{receiver/lem:fock-permanent-maps}{Lemma}{lem:fock-permanent-maps}{The lift of a map between two mode sets}
\claimleadin{%
Lemma~\ref{lem:fock-functor} and Lemma~\ref{lem:fock-permanent} are stated for
an operator on one mode space. Theorem~\ref{thm:factor-blocks} lifts an
isometry from the declared inputs into a different set of modes, and needs the
same two facts for it: that the lift composes and takes adjoints, and that its
entries are permanents.
}
\begin{technicalclaim}
\begin{lemma}[The lift of a map between two mode sets]\label{lem:fock-permanent-maps}
Let $I$, $O$ and $O'$ be finite mode sets, each with an orthonormal basis
listed in a fixed order, let $X:\mathbb{C}^{I}\to\mathbb{C}^{O}$ and
$X':\mathbb{C}^{O}\to\mathbb{C}^{O'}$ be linear maps, and let $\Fock(X)$ be
the lift of Definition~\ref{def:fock}.
\begin{enumerate}[label=(\roman*),leftmargin=3.0em]
\item $\Fock(X)$ maps $\Symn{n}(\mathbb{C}^{I})$ into
  $\Symn{n}(\mathbb{C}^{O})$ for every $n\ge0$;
  $\Fock(X'X)=\Fock(X')\Fock(X)$ and $\Fock(X)^{\dagger}=\Fock(X^{\dagger})$;
  and if $X^{\dagger}X=\id_{I}$, then $\Fock(X)^{\dagger}\Fock(X)$ is the
  identity of $\mathcal{F}(\mathbb{C}^{I})$, so that $\Fock(X)$ is an
  isometry.
\item Let $\vec m\in\mathbb{N}_0^{O}$ and $\vec m'\in\mathbb{N}_0^{I}$ be
  occupation vectors with $|\vec m|=|\vec m'|=n$, let $w_{\vec m}$,
  $w_{\vec m'}$, $\vec m!$ and $\vec m'!$ be the words and the factorial
  products of Lemma~\ref{lem:fock-permanent}, each formed in its own mode set,
  and let $X[\vec m|\vec m']$ be the $n\times n$ matrix with entries
  $X_{w_{\vec m}(r),\,w_{\vec m'}(s)}$. Then
  \begin{equation}\label{eq:fock-permanent-maps}
    \bra{\vec m}\Fock(X)^{(n)}\ket{\vec m'}
    =\frac{\operatorname{per}X[\vec m|\vec m']}{\sqrt{\vec m!\,\vec m'!}} ,
  \end{equation}
  with $\operatorname{per}$ as in Lemma~\ref{lem:fock-permanent}; at $n=0$
  both sides are $1$.
\end{enumerate}
\end{lemma}
\end{technicalclaim}
\claimmeta{\metalabel{Uses}\mref{Def.}{def:fock}, \mref{Lem.}{lem:fock-permanent}
\metasep \metalabel{Used by}\mref{Thm.}{thm:factor-blocks}
\metasep newly derived
\metasep implemented (effective: accepted)
\metasep \metalabel{Code}\coderef{https://github.com/fraunhofersg/certiqs-key/blob/v0.1.0/src/certiqs_key/receiver/factor.py}{certiqs_key.receiver.factor:_column}}

\begin{proof}
(i) On $\Symn{n}(\mathbb{C}^{I})$ the lift $\Fock(X)$ is the restriction of
$X^{\otimes n}$, which carries $\Symn{n}(\mathbb{C}^{I})$ into
$\Symn{n}(\mathbb{C}^{O})$ (Definition~\ref{def:fock}); so $\Fock(X)$ maps
each photon-number summand into the summand of the same $n$. Factor by
factor, $(X'X)^{\otimes n}=X'^{\otimes n}X^{\otimes n}$, and since
$X^{\otimes n}$ lands in $\Symn{n}(\mathbb{C}^{O})$, the restriction of the
product to $\Symn{n}(\mathbb{C}^{I})$ is the product of the restrictions;
summing over $n$ gives $\Fock(X'X)=\Fock(X')\Fock(X)$. On the full tensor
powers $(X^{\otimes n})^{\dagger}=(X^{\dagger})^{\otimes n}$, and
$(X^{\dagger})^{\otimes n}$ carries $\Symn{n}(\mathbb{C}^{O})$ into
$\Symn{n}(\mathbb{C}^{I})$ (Definition~\ref{def:fock} for the map
$X^{\dagger}$). So for $\psi\in\Symn{n}(\mathbb{C}^{O})$ and
$\phi\in\Symn{n}(\mathbb{C}^{I})$,
$\langle\psi,X^{\otimes n}\phi\rangle=\langle(X^{\dagger})^{\otimes n}\psi,\phi\rangle$
with both vectors on the right in their symmetric subspaces, and the
restriction of $(X^{\dagger})^{\otimes n}$ to $\Symn{n}(\mathbb{C}^{O})$ is
the adjoint of the restriction of $X^{\otimes n}$ to
$\Symn{n}(\mathbb{C}^{I})$. Distinct $n$ give orthogonal summands, so
$\Fock(X)^{\dagger}=\Fock(X^{\dagger})$. If $X^{\dagger}X=\id_{I}$, these give
$\Fock(X)^{\dagger}\Fock(X)=\Fock(X^{\dagger}X)=\Fock(\id_{I})$, and
$\id^{\otimes n}$ restricts to the identity of every
$\Symn{n}(\mathbb{C}^{I})$.

(ii) Let $P_I=\frac1{n!}\sum_{\pi}\Pi^{I}_{\pi}$ and
$P_O=\frac1{n!}\sum_{\pi}\Pi^{O}_{\pi}$ be the orthogonal projectors onto
$\Symn{n}(\mathbb{C}^{I})$ and $\Symn{n}(\mathbb{C}^{O})$, $\Pi^{I}_{\pi}$ and
$\Pi^{O}_{\pi}$ permuting the tensor factors, and let $\ket{w_{\vec m}}$ and
$\ket{w_{\vec m'}}$ be the product basis vectors whose $r$-th factors are the
basis vectors of the letters $w_{\vec m}(r)$ and $w_{\vec m'}(r)$. The
permutations that fix $\ket{w_{\vec m}}$ are those that permute equal letters
among themselves, $\vec m!$ of them, and every other $\Pi^{O}_{\pi}$ sends it
to an orthogonal product basis vector; so
$\lVert P_O\ket{w_{\vec m}}\rVert^{2}=\bra{w_{\vec m}}P_O\ket{w_{\vec m}}
=\vec m!/n!$, and the symmetrized and normalized vector of
Definition~\ref{def:fock} is
$\ket{\vec m}=\sqrt{n!/\vec m!}\,P_O\ket{w_{\vec m}}$, and likewise
$\ket{\vec m'}=\sqrt{n!/\vec m'!}\,P_I\ket{w_{\vec m'}}$. For every
permutation $\pi$, $X^{\otimes n}\Pi^{I}_{\pi}=\Pi^{O}_{\pi}X^{\otimes n}$,
both sending $v_1\otimes\cdots\otimes v_n$ to the tensor product of the
$Xv_r$ permuted by $\pi$; so $X^{\otimes n}P_I=P_OX^{\otimes n}$. With
$P_O^{\dagger}=P_O$ and $P_I^{2}=P_I$,
\begin{equation*}
  \bra{\vec m}X^{\otimes n}\ket{\vec m'}
  =\frac{n!}{\sqrt{\vec m!\,\vec m'!}}
   \bra{w_{\vec m}}P_OX^{\otimes n}P_I\ket{w_{\vec m'}}
  =\frac{n!}{\sqrt{\vec m!\,\vec m'!}}
   \bra{w_{\vec m}}X^{\otimes n}P_I\ket{w_{\vec m'}}
  =\frac{1}{\sqrt{\vec m!\,\vec m'!}}\sum_{\pi}
   \bra{w_{\vec m}}X^{\otimes n}\Pi^{I}_{\pi}\ket{w_{\vec m'}} .
\end{equation*}
The vector $\Pi^{I}_{\pi}\ket{w_{\vec m'}}$ is the product vector whose $r$-th
factor is the basis vector of $w_{\vec m'}(\sigma(r))$ for a permutation
$\sigma$ determined by $\pi$, and $\pi\mapsto\sigma$ is a bijection of the
permutation group. The summand is therefore
$\prod_{r}X_{w_{\vec m}(r),\,w_{\vec m'}(\sigma(r))}$, and the sum over
$\sigma$ is $\operatorname{per}X[\vec m|\vec m']$. The left-hand side is the
entry of $\Fock(X)^{(n)}$, $\ket{\vec m'}$ lying in
$\Symn{n}(\mathbb{C}^{I})$. At $n=0$ the lift is $1$ on the vacuum, the
permanent of the empty matrix is $1$ and both factorial products are $1$.
\end{proof}

% claim receiver/thm:enclosed-blocks
\claimsummary{receiver/thm:enclosed-blocks}{Enclosures of the port Grams, obtained from enclosures of square roots, cosines and sines, carry through permanents, inclusion and exclusion and the coarse-graining to enclosures of every block. With Hermitian centres they make the emitted blocks up to any cutoff a reference.}
\claimbookmark{receiver/thm:enclosed-blocks}{Theorem}{thm:enclosed-blocks}{Rigorous enclosures of the photon-number blocks}
\claimleadin{%
Every block of the receiver's measurement is a finite expression in the
entries of the port Grams: a permanent divided by square roots of factorials
(Lemma~\ref{lem:fock-permanent}), combined by inclusion and exclusion and
summed over the fibres of the coarse-graining. A compiled receiver's port
Grams are in turn finite expressions in the entries of its elements, which are
square roots, cosines and sines of the declared parameters. Enclosures of
these, carried through the expressions by
Lemma~\ref{lem:enclosure-arithmetic}, therefore enclose the blocks of the
model at the declared parameters, whatever the parameters are, with no
floating-point step. The enclosures depend on the order in which the products
are formed; the value they enclose does not
(Theorem~\ref{thm:compilation-well-defined}\,(i)). Where the enclosures are
exact, the reference is exact (Theorem~\ref{thm:exact-blocks}).
}
\begin{theorem}[Rigorous enclosures of the photon-number blocks]\label{thm:enclosed-blocks}
Let $\Recv$ be a \termref{term-receiver-instrument}{receiver instrument} with
\termref{term-port-gram-matrices}{port Grams} $\{A_k\}_{k=1}^{K}$ on
$\mathbb{C}^{\Inp}$ (Definitions~\ref{def:instrument} and~\ref{def:gram}),
$\{M_S\}_{S\subseteq[K]}$ its click-pattern POVM
(Theorem~\ref{thm:click-povm}), $\cg$ a \termref{term-declared-coarse-graining}{declared coarse-graining} onto
$\Vset$ and $\{F_v\}_{v\in\Vset}$ the \termref{term-emitted-povm}{emitted POVM}
(Definition~\ref{def:coarse}). Fix $N\ge0$, and read every block
$\Symn{n}(\mathbb{C}^{\Inp})$, $n\le N$, in its occupation basis listed in an
order fixed for that block. Enclosures are
\termref{term-complex-enclosure}{complex enclosures}
(Definition~\ref{def:complex-enclosure}), and ``the rules'' are those of
Lemma~\ref{lem:enclosure-arithmetic}\,(i)--(iv).
\begin{enumerate}[label=(\roman*),leftmargin=3.0em]
\item \emph{Port Grams of a compiled receiver.} Suppose that the
  \termref{term-instrument-matrix}{instrument matrix} of $\Recv$ is the
  \termref{term-compiled-instrument-matrix}{compiled instrument matrix} $\Gmat=D_\eta\Gmat^{\sigma}_{\Grph}$ of a
  well-formed graph for a \termref{term-linear-extension}{linear extension} $\sigma$
  (Definition~\ref{def:compile}), and that enclosures are given of the matrix
  $\hat\Gmat_e$ of every element $e$ and of $\sqrt{\eta_k}$ for every port
  $k$. Enclose every other entry of $D_\eta$, and every entry of each
  $\Emd{\{k\}}$, exactly by itself. Then the rules, applied to the products
  \eqref{eq:compile} in the order of $\sigma$, to the restriction of
  $\Gmat$ to the columns $\Inp$, and to \eqref{eq:gramdef}, give enclosures
  of the $A_k$.
\item \emph{Blocks.} Let $(C_k,r_k)$ be enclosures of the $A_k$, and for
  every $n\le N$ and occupation vectors $\vec m,\vec m'$ of weight $n$ let
  $(a_{\vec m\vec m'},e_{\vec m\vec m'})$ be a \termref{term-square-root-enclosure}{square-root enclosure}
  of $1/(\vec m!\,\vec m'!)$. For $J\subseteq[K]$ enclose $\id-A_J$ by the
  rules from the $(C_k,r_k)$, the identity being exact; enclose the entry
  $(\vec m,\vec m')$ of $\Pvac{J}^{(n)}$ by the rules applied to the product
  of $\operatorname{per}\big((\id-A_J)[\vec m|\vec m']\big)$ with
  $\sqrt{1/(\vec m!\,\vec m'!)}$; enclose $\blk{M}{\eqn}_S$ by the rules
  applied to \eqref{eq:blockMS}, the signs being exact; and enclose
  $\blk{F}{\eqn}_v$ by the rules applied to
  $\sum_{S\in\cg^{-1}(v)}\blk{M}{\eqn}_S$. Each enclosure so obtained is a
  complex enclosure of the block it names.
\item \emph{The reference.} Replace each enclosure of a block
  $\blk{F}{\eqn}_v$ of (ii) by the one with a Hermitian centre of
  Lemma~\ref{lem:enclosure-arithmetic}\,(v). Let $\mathfrak{m}_N$ be a
  \exttermref{../foundations/contract.pdf}{term-sector-map}{sector map} (\xref{foundations/def:sector-map}) whose blocks are the
  photon numbers $n=0,\dots,N$ of the declared inputs, in increasing order,
  of dimensions $\binom{|\Inp|+n-1}{n}$, with any block kind and block
  identifiers, each block read in its fixed order. Then the family
  $\{F_v\}$ restricted to the blocks $n\le N$, read with these enclosures, is
  a \exttermref{../foundations/contract.pdf}{term-reference}{reference} with constant $1$ against $\mathfrak{m}_N$ over $\Vset$
  (\xref{foundations/def:exact-reference}); the same holds for $\{M_S\}$ over
  the click patterns.
\item \emph{Exactness.} An entry of an enclosure of (i)--(iii) is exact where
  every enclosure it is computed from is exact; the reference of (iii) is an
  \exttermref{../foundations/contract.pdf}{term-exact-reference}{exact reference} when every one of its radii is zero.
\end{enumerate}
\end{theorem}
\claimmeta{\metalabel{Uses}\mref{Def.}{def:instrument}, \mref{Def.}{def:gram}, \mref{Def.}{def:compile}, \mref{Thm.}{thm:click-povm}, \mref{Def.}{def:coarse}, \mref{Thm.}{thm:coarse-graining}, \mref{Def.}{def:complex-enclosure}, \mref{Lem.}{lem:enclosure-arithmetic}, \mref{Lem.}{lem:fock-permanent}, \extref{../foundations/contract.pdf}{foundations/def:exact-reference}, \extref{../foundations/contract.pdf}{foundations/def:sector-map}
\metasep \metalabel{Used by}\mref{Thm.}{thm:exact-blocks}, \mref{Thm.}{thm:exact-and-evaluated}, \mref{Req.}{req:rigorous-weight}, \mref{Req.}{req:exact-route-decision}, \mref{Cor.}{cor:exact-dark-composed}, \mref{Req.}{req:dark-composed-reference}, \mref{Req.}{req:snap-fix}, \extref{../reductions/contract.pdf}{reductions/req:recv-row-provenance}
\metasep newly derived
\metasep implemented (effective: accepted)
\metasep \metalabel{Code}\coderef{https://github.com/fraunhofersg/certiqs-key/blob/v0.1.0/src/certiqs_key/receiver/generic.py}{certiqs_key.receiver.generic:element_blocks}}

\begin{proof}
(i) Every entry of a product or a sum of matrices is a sum of products of
entries, the adjoint conjugates and transposes entries, and restricting to
the columns $\Inp$ selects entries. The matrix $D_\eta$ is diagonal with the
entry $\sqrt{\eta_k}$ on the modes of port $k$ and $1$ elsewhere
(Definition~\ref{def:compile}), and $\Emd{\{k\}}$ is diagonal with entries
$0$ and $1$, so all their entries are enclosed. Each entry of
$A_k=\Gmat_{\Inp}^{\dagger}\Emd{\{k\}}\Gmat_{\Inp}$ is thus a finite
expression in enclosed numbers, and
Lemma~\ref{lem:enclosure-arithmetic}\,(iv), which needs no reciprocal here,
encloses it.

(ii) By Theorem~\ref{thm:click-povm}\,(ii) and its block structure (iii),
$\Pvac{J}=\Fock(\id-A_J)$ is block diagonal, and by
\eqref{eq:fock-permanent} of Lemma~\ref{lem:fock-permanent}\,(i) the entry
$(\vec m,\vec m')$ of its block $n$ is
$\operatorname{per}\big((\id-A_J)[\vec m|\vec m']\big)\cdot
\sqrt{1/(\vec m!\,\vec m'!)}$. The permanent is a sum of products of entries
of $\id-A_J=\id-\sum_{k\in J}A_k$, each of which is enclosed by
Lemma~\ref{lem:enclosure-arithmetic}\,(i); the square root is enclosed by
hypothesis; so Lemma~\ref{lem:enclosure-arithmetic}\,(iv) encloses the entry.
Equation \eqref{eq:blockMS} writes $\blk{M}{\eqn}_S$ as
$\sum_{T\subseteq S}(-1)^{|T|}\Pvac{([K]\setminus S)\cup T}^{(n)}$, an
integer combination taken entry by entry, and
Theorem~\ref{thm:coarse-graining}\,(i) gives
$\blk{F}{\eqn}_v=\sum_{S\in\cg^{-1}(v)}\blk{M}{\eqn}_S$; both are finite
expressions in the enclosed entries, enclosed by the same lemma.

(iii) Each $\blk{F}{\eqn}_v$ is Hermitian, being a block of the
block-diagonal POVM element $F_v$, so
Lemma~\ref{lem:enclosure-arithmetic}\,(v) gives an
\exttermref{../foundations/contract.pdf}{term-entrywise-enclosure}{entrywise enclosure} with a Hermitian Gaussian-rational centre and
rational radii. The blocks are positive semidefinite, since
$M_S\succeq0$ (Theorem~\ref{thm:click-povm}\,(i)) and a block of a
block-diagonal positive semidefinite operator is positive semidefinite, and
$F_v$ is a sum of such. Completeness on every block holds because
$\sum_SM_S=\id$ restricts to $\sum_S\blk{M}{\eqn}_S=\id$ on
$\Symn{n}(\mathbb{C}^{\Inp})$ and the fibres of $\cg$ partition $2^{[K]}$, so
$\sum_v\blk{F}{\eqn}_v=\id$ there. The dimension of
$\Symn{n}(\mathbb{C}^{\Inp})$ is the number of occupation vectors of weight
$n$ on $|\Inp|$ modes, $\binom{|\Inp|+n-1}{n}$, which the sector map
declares, and the constant $1$ is a positive declared value. These are the
conditions of a reference. Nothing depends on the block kind or the
identifiers; the enclosures were formed in each block's fixed order, which is
the order in which the sector map reads the block. The statement for
$\{M_S\}$ is the same argument without the fibre sums.

(iv) Lemma~\ref{lem:enclosure-arithmetic}\,(iv) gives an exact enclosure at
every node whose operands are exactly enclosed, and (v) of the same lemma
keeps an entry exact where it and its transposed entry are. A reference all
of whose radii are zero is an exact reference by definition.
\end{proof}

% claim receiver/def:port-dilation
\claimsummary{receiver/def:port-dilation}{A port dilation is an isometry from the declared inputs into a finite mode set with disjoint ports that reproduces every port Gram. Any such map carries the click statistics as the unitary dilation of the physical model does, and it may have far fewer modes.}
\claimbookmark{receiver/def:port-dilation}{Definition}{def:port-dilation}{Port dilation of a receiver instrument}
\claimleadin{%
The \termref{term-measurement-isometry}{measurement isometry} of Definition~\ref{def:physmodel} applies a unitary
on the modes and an environment to the declared inputs, and
Theorem~\ref{thm:port-grams}\,(iii) and~(v) show that the click statistics
read of it only the amplitudes it sends from the inputs to each port. The
definition names the least that a map from the inputs must satisfy to carry
the same statistics, so that a smaller isometry, into the port modes and as
many loss modes as the instrument's loss needs, can stand in for the unitary
dilation: Theorem~\ref{thm:factor-blocks} reads every photon-number block
through such a map. The restriction $U\iota$ of a unitary dilation $U$ to the
declared inputs, with the ports of the instrument, is one, since
$(U\iota)^{\dagger}U\iota=\id_{\Inp}$ and the rows of $U\iota$ on a port are
those of $\Gmat_{\Inp}$.
}
\begin{definition}[Port dilation of a receiver instrument]\label{def:port-dilation}
Let $\Recv$ be a \termref{term-receiver-instrument}{receiver instrument} with
\termref{term-port-gram-matrices}{port Grams} $\{A_k\}_{k=1}^{K}$ on
$\mathbb{C}^{\Inp}$ (Definitions~\ref{def:instrument} and~\ref{def:gram}). A
\termdef{term-port-dilation}{port dilation} of $\Recv$ is a triple
$(\mathbb{O};\mathbb{O}_1,\dots,\mathbb{O}_K;\Omega)$ of a finite mode set
$\mathbb{O}$ with a fixed order, pairwise disjoint subsets
$\mathbb{O}_1,\dots,\mathbb{O}_K$ of $\mathbb{O}$, its \emph{ports}, and a
linear map $\Omega:\mathbb{C}^{\Inp}\to\mathbb{C}^{\mathbb{O}}$ with
\begin{equation}\label{eq:port-dilation}
  \Omega^{\dagger}\Omega=\id_{\Inp}
  \qquad\text{and}\qquad
  \Omega^{\dagger}\Emd{\mathbb{O}_k}\Omega=A_k\quad(k=1,\dots,K),
\end{equation}
where $\Emd{\mathbb{O}_k}$ is the coordinate projector of
$\mathbb{C}^{\mathbb{O}}$ onto the modes of $\mathbb{O}_k$. The modes of
$\mathbb{O}$ in no port are its \emph{loss modes}, and the \emph{row} of a mode
$i\in\mathbb{O}$ is the row of $\Omega$ that $i$ indexes. The \emph{click
pattern} of an occupation vector $\vec m\in\mathbb{N}_0^{\mathbb{O}}$
(Definition~\ref{def:fock}) is
\begin{equation}\label{eq:click-pattern-of}
  S_{\vec m}\coloneqq\big\{k\in[K]:\ m_i\ge1\text{ for some }i\in\mathbb{O}_k\big\}.
\end{equation}
\end{definition}
\claimmeta{\metalabel{Uses}\mref{Def.}{def:fock}, \mref{Def.}{def:instrument}, \mref{Def.}{def:gram}
\metasep \metalabel{Used by}\mref{Thm.}{thm:factor-blocks}, \mref{Lem.}{lem:catalogue-dilations}, \mref{Req.}{req:dark-composed-reference}
\metasep newly derived
\metasep implemented (effective: accepted)
\metasep \metalabel{Code}\coderef{https://github.com/fraunhofersg/certiqs-key/blob/v0.1.0/src/certiqs_key/receiver/factor.py}{certiqs_key.receiver.factor:PortDilation}}

% claim receiver/thm:factor-blocks
\claimsummary{receiver/thm:factor-blocks}{Read through a port dilation, a non-negative combination of click-pattern blocks is a weighted sum of rank-one terms whose columns are permanents of the dilation, enclosed by the same rules as the blocks. With rational weights, columns and weights are a factor enclosure of the block.}
\claimbookmark{receiver/thm:factor-blocks}{Theorem}{thm:factor-blocks}{Photon-number blocks in enclosed factor form}
\claimleadin{%
Theorem~\ref{thm:enclosed-blocks} encloses every block entry by entry. A block
whose rank is below its dimension and whose support no coordinate vectors
span is then known only through a box that also contains matrices of higher
rank, and no lower bound on a least eigenvalue over that box fixes its rank.
The physical model gives such a block in another form. A click pattern is an
event on the output modes that is diagonal in their occupation basis, so
pulling it back through the measurement writes its block as a sum of
rank-one terms, one for each output occupation vector with that pattern,
whose columns are the pulled-back occupation vectors. A port dilation
(Definition~\ref{def:port-dilation}) serves in place of the unitary dilation,
and the columns are permanents of its entries
(Lemma~\ref{lem:fock-permanent-maps}), enclosed by the rules that enclose the
blocks. Read through its columns, the block keeps the rank they fix
(\xref{foundations/cor:rank-over-factor-enclosure}). A coarse-graining, and
a classical channel of dark counts on the recorded patterns, combine click
patterns with non-negative weights, which keeps the form.
}
\begin{theorem}[Photon-number blocks in enclosed factor form]\label{thm:factor-blocks}
Let $\Recv$ be a \termref{term-receiver-instrument}{receiver instrument} with
\termref{term-port-gram-matrices}{port Grams} $\{A_k\}_{k=1}^{K}$ on
$\mathbb{C}^{\Inp}$ and click-pattern POVM $\{M_S\}_{S\subseteq[K]}$
(Theorem~\ref{thm:click-povm}), and let
$(\mathbb{O};\mathbb{O}_1,\dots,\mathbb{O}_K;\Omega)$ be a
\termref{term-port-dilation}{port dilation} of $\Recv$ (Definition~\ref{def:port-dilation}), with
$S_{\vec m}$ the click pattern of an occupation vector $\vec m$ of
$\mathbb{O}$. Read every block $\Symn{n}(\mathbb{C}^{\Inp})$ in its
occupation basis listed in an order fixed for that block, and for an
occupation vector $\vec m\in\mathbb{N}_0^{\mathbb{O}}$ of weight $n$ put
\begin{equation}\label{eq:factor-column}
  c_{\vec m}\coloneqq\Fock(\Omega)^{\dagger}\ket{\vec m}
  \ \in\ \Symn{n}(\mathbb{C}^{\Inp}).
\end{equation}
Enclosures are \termref{term-complex-enclosure}{complex enclosures}
(Definition~\ref{def:complex-enclosure}), and ``the rules'' are those of
Lemma~\ref{lem:enclosure-arithmetic}\,(i)--(iv).
\begin{enumerate}[label=(\roman*),leftmargin=3.0em]
\item \emph{The factor form.} For every $n\ge0$ and all reals $w_S\ge0$,
  $S\subseteq[K]$,
  \begin{equation}\label{eq:factor-form}
    \sum_{S\subseteq[K]}w_S\,\blk{M}{\eqn}_S
    \;=\;\sum_{\vec m:\,|\vec m|=n}w_{S_{\vec m}}\,c_{\vec m}c_{\vec m}^{\dagger},
  \end{equation}
  the sum on the right over the occupation vectors of $\mathbb{O}$ of weight
  $n$. In particular $\blk{M}{\eqn}_S$ is the sum of the
  $c_{\vec m}c_{\vec m}^{\dagger}$ with $S_{\vec m}=S$, and the block
  $\blk{F}{\eqn}_v$ of the \termref{term-emitted-povm}{emitted POVM} of a
  \termref{term-declared-coarse-graining}{declared coarse-graining} $\cg$ (Definition~\ref{def:coarse}) is
  \eqref{eq:factor-form} with $w_S=1$ for $S\in\cg^{-1}(v)$ and $w_S=0$
  otherwise. The entry of $c_{\vec m}$ at an occupation vector $\vec m'$ of
  $\Inp$ is
  $\overline{\operatorname{per}\Omega[\vec m|\vec m']}\big/\sqrt{\vec m!\,\vec m'!}$,
  with $\Omega[\vec m|\vec m']$ as in
  Lemma~\ref{lem:fock-permanent-maps}\,(ii).
\item \emph{Enclosed columns.} Let enclosures of the entries of $\Omega$ be
  given, and, for all occupation vectors $\vec m$ of $\mathbb{O}$ and
  $\vec m'$ of $\Inp$ of weight $n$, a \termref{term-square-root-enclosure}{square-root enclosure} of
  $1/(\vec m!\,\vec m'!)$. The rules, applied to the conjugate of
  $\operatorname{per}\Omega[\vec m|\vec m']$ times that square root, enclose
  every entry of every $c_{\vec m}$. An entry so enclosed is exact where every
  enclosure it is computed from is exact, and it is enclosed exactly as $0$
  where every term of the permanent, a product of entries of $\Omega$, has a
  factor enclosed exactly as $0$; in particular every entry of $c_{\vec m}$
  is, where $\vec m$ has a photon in a mode whose row is enclosed exactly as
  zero.
\item \emph{The enclosure of the factor form.} Let $n\ge0$, let the weights $w_S\ge0$ be
  rationals, let enclosures of the columns $c_{\vec m}$ be given, such as
  those of (ii), and let $\mathcal{C}$ be the set of the occupation vectors
  $\vec m$ of $\mathbb{O}$ of weight $n$ with $w_{S_{\vec m}}>0$ whose column
  is not enclosed exactly as zero, listed in a fixed order. Let $(L_c,\varrho)$
  be the \exttermref{../foundations/contract.pdf}{term-matrix-box}{matrix box} of complex
  $\dim\Symn{n}(\mathbb{C}^{\Inp})\times|\mathcal{C}|$ matrices whose column
  $\vec m$ has the centres and the radii of the enclosure of $c_{\vec m}$,
  and $W\coloneqq\operatorname{diag}(w_{S_{\vec m}})_{\vec m\in\mathcal{C}}$.
  Then $(L_c,\varrho,W)$ is a \exttermref{../foundations/contract.pdf}{term-factor-enclosure}{factor enclosure}
  (\xref{foundations/def:factor-enclosure}) of the block
  \eqref{eq:factor-form}, with $|\mathcal{C}|$ columns; where $\mathcal{C}$
  is empty, that block is zero.
\item \emph{Port dilations from factors.} Let
  $\mathbb{O}_0,\mathbb{O}_1,\dots,\mathbb{O}_K$ be pairwise disjoint finite
  mode sets and $Y_k:\mathbb{C}^{\Inp}\to\mathbb{C}^{\mathbb{O}_k}$ linear
  maps with $Y_k^{\dagger}Y_k=A_k$ for $k=1,\dots,K$ and
  $Y_0^{\dagger}Y_0=\defic$, the \termref{term-loss-deficiency}{loss deficiency}
  (Definition~\ref{def:gram}). Then the map $\Omega$ into
  $\mathbb{O}\coloneqq\mathbb{O}_0\sqcup\mathbb{O}_1\sqcup\dots\sqcup\mathbb{O}_K$,
  in any fixed order, whose rows on $\mathbb{O}_k$ are those of $Y_k$ for
  $k=0,\dots,K$, is with the ports $\mathbb{O}_1,\dots,\mathbb{O}_K$ a port
  dilation of $\Recv$. If a matrix $X$ is $E^{\dagger}DE$ with $D$ diagonal with
  positive rational diagonal entries $d_i$ and $E$ with entries in
  $\mathbb{Q}(i)$, then $Y\coloneqq D^{1/2}E$ satisfies $Y^{\dagger}Y=X$, and
  the rules applied to the products of square-root enclosures of the $d_i$
  with the entries of $E$ enclose $Y$, exactly in every row $i$ in which
  $d_i$ is the square of a rational.
\end{enumerate}
\end{theorem}
\claimmeta{\metalabel{Uses}\mref{Def.}{def:fock}, \mref{Lem.}{lem:fock-functor}, \mref{Def.}{def:gram}, \mref{Thm.}{thm:click-povm}, \mref{Def.}{def:coarse}, \mref{Def.}{def:complex-enclosure}, \mref{Lem.}{lem:enclosure-arithmetic}, \mref{Lem.}{lem:fock-permanent-maps}, \mref{Def.}{def:port-dilation}, \extref{../foundations/contract.pdf}{foundations/def:factor-enclosure}
\metasep \metalabel{Used by}\mref{Req.}{req:dark-composed-reference}
\metasep newly derived
\metasep implemented (effective: accepted)
\metasep \metalabel{Code}\coderef{https://github.com/fraunhofersg/certiqs-key/blob/v0.1.0/src/certiqs_key/receiver/factor.py}{certiqs_key.receiver.factor:factor_enclosures}}

\begin{proof}
\emph{(i).} By Lemma~\ref{lem:fock-permanent-maps}\,(i), $\Fock(\Omega)$ maps
$\Symn{n}(\mathbb{C}^{\Inp})$ into $\Symn{n}(\mathbb{C}^{\mathbb{O}})$, and its
adjoint $\Fock(\Omega^{\dagger})$ maps back, so $c_{\vec m}$ lies in
$\Symn{n}(\mathbb{C}^{\Inp})$. For $J\subseteq[K]$ put
$E_J\coloneqq\sum_{k\in J}\Emd{\mathbb{O}_k}$, a coordinate projector of
$\mathbb{C}^{\mathbb{O}}$ since the ports are pairwise disjoint, and
$\Pi_J\coloneqq\Fock(\id_{\mathbb{O}}-E_J)$. By the same lemma and
\eqref{eq:port-dilation},
\begin{equation*}
  \Fock(\Omega)^{\dagger}\,\Pi_J\,\Fock(\Omega)
  =\Fock\big(\Omega^{\dagger}(\id_{\mathbb{O}}-E_J)\Omega\big)
  =\Fock\Big(\id_{\Inp}-\sum_{k\in J}A_k\Big)
  =\Fock(\id_{\Inp}-A_J).
\end{equation*}
By Lemma~\ref{lem:fock-functor}\,(iv) on $\mathbb{C}^{\mathbb{O}}$, whose
coordinate basis is adapted to the range and kernel of the coordinate
projector $\id_{\mathbb{O}}-E_J$, $\Pi_J$ is the orthogonal projector onto the
span of the occupation vectors with no photon in a mode of
$\bigcup_{k\in J}\mathbb{O}_k$; the occupation vectors being an orthonormal
basis, $\Pi_J\ket{\vec m}=\ket{\vec m}$ where $S_{\vec m}\cap J=\varnothing$
and $\Pi_J\ket{\vec m}=0$ otherwise. Fix $S\subseteq[K]$. By
Theorem~\ref{thm:click-povm}\,(ii) and the identity above,
\begin{equation*}
  M_S=\sum_{T\subseteq S}(-1)^{|T|}\,
      \Fock\big(\id_{\Inp}-A_{([K]\setminus S)\cup T}\big)
  =\Fock(\Omega)^{\dagger}P_S\,\Fock(\Omega),
  \qquad
  P_S\coloneqq\sum_{T\subseteq S}(-1)^{|T|}\,\Pi_{([K]\setminus S)\cup T}.
\end{equation*}
Let $\vec m$ be an occupation vector of $\mathbb{O}$. If $S_{\vec m}$ meets
$[K]\setminus S$, every term of $P_S$ annihilates $\ket{\vec m}$. Otherwise
$S_{\vec m}\subseteq S$, and the term of $T$ keeps $\ket{\vec m}$ exactly
when $T\cap S_{\vec m}=\varnothing$, so
$P_S\ket{\vec m}=\sum_{T\subseteq S\setminus S_{\vec m}}(-1)^{|T|}\ket{\vec m}
=(1-1)^{|S\setminus S_{\vec m}|}\ket{\vec m}$, which is $\ket{\vec m}$ where
$S_{\vec m}=S$ and $0$ otherwise. So $P_S$ restricted to
$\Symn{n}(\mathbb{C}^{\mathbb{O}})$ is
$\sum_{|\vec m|=n,\,S_{\vec m}=S}\ket{\vec m}\bra{\vec m}$, and since
$\Fock(\Omega)$ maps $\Symn{n}(\mathbb{C}^{\Inp})$ into
$\Symn{n}(\mathbb{C}^{\mathbb{O}})$, the block of $M_S$ on
$\Symn{n}(\mathbb{C}^{\Inp})$ is
\begin{equation*}
  \blk{M}{\eqn}_S=\sum_{|\vec m|=n,\,S_{\vec m}=S}
  \Fock(\Omega)^{\dagger}\ket{\vec m}\bra{\vec m}\Fock(\Omega)
  =\sum_{|\vec m|=n,\,S_{\vec m}=S}c_{\vec m}c_{\vec m}^{\dagger}.
\end{equation*}
Every occupation vector has exactly one click pattern, so the sum of the
$w_S\blk{M}{\eqn}_S$ is \eqref{eq:factor-form}. The emitted POVM is
$F_v=\sum_{S\in\cg^{-1}(v)}M_S$ (Definition~\ref{def:coarse}), and every $M_S$
is block diagonal in the photon number (Theorem~\ref{thm:click-povm}\,(iii)),
so the block of $F_v$ is the sum of the blocks of its $M_S$. Finally,
$\bra{\vec m'}c_{\vec m}=\bra{\vec m'}\Fock(\Omega)^{\dagger}\ket{\vec m}
=\overline{\bra{\vec m}\Fock(\Omega)\ket{\vec m'}}$, which
Lemma~\ref{lem:fock-permanent-maps}\,(ii) evaluates, the square root being a
positive real.

\emph{(ii).} The permanent of $\Omega[\vec m|\vec m']$ is a sum of products of
entries of $\Omega$, each of them enclosed, and the entry of $c_{\vec m}$ is
its conjugate times the square root; Lemma~\ref{lem:enclosure-arithmetic}\,(iv)
encloses the value of this expression and keeps it exact where its inputs
are. A product one of whose operands is enclosed exactly as $0$, with centre
and radius $0$, receives by part (ii) of that lemma the centre $0$ and the
radius $0\cdot r_y+0\cdot\lvert y_c\rvert_{1}+0\cdot r_y=0$, so a term of the
permanent with such a factor is enclosed exactly as $0$, whatever order its
factors are multiplied in; a sum of such terms is enclosed exactly as $0$ by
part (i), and so are its conjugate and its product with the square root.
Where $\vec m$ has a photon in a mode $i$, the letter $i$ occurs in the word
$w_{\vec m}$, so every term of the permanent has a factor from the row of
$i$; if that row is enclosed exactly as zero, every term has a factor
enclosed exactly as $0$.

\emph{(iii).} Let $L$ be the matrix whose columns are the $c_{\vec m}$,
$\vec m\in\mathcal{C}$, in the fixed order. Each of its entries lies within
the radius of the centre of its enclosure, so the box $(L_c,\varrho)$, whose
centres have entries in $\mathbb{Q}(i)$ and whose radii are non-negative
rationals (Definition~\ref{def:complex-enclosure}), contains $L$; and $W$ is
diagonal with positive rational diagonal entries. A column enclosed exactly
as zero is zero, an exact enclosure being equal to its centre. In
\eqref{eq:factor-form} the terms with $w_{S_{\vec m}}=0$ and the terms whose
column is zero vanish, and the remaining terms are those of $\mathcal{C}$, so
the block is
$\sum_{\vec m\in\mathcal{C}}w_{S_{\vec m}}c_{\vec m}c_{\vec m}^{\dagger}=LWL^{\dagger}$
with $L$ in the box, which is what a factor enclosure asks of the matrix it
encloses (\xref{foundations/def:factor-enclosure}). Where $\mathcal{C}$ is
empty the sum is empty and the block is zero.

\emph{(iv).} The rows of $\Omega$ on $\mathbb{O}_k$ are those of $Y_k$, so
$\Omega^{\dagger}\Emd{\mathbb{O}_k}\Omega=Y_k^{\dagger}Y_k=A_k$ for
$k=1,\dots,K$, and, the $\mathbb{O}_k$ partitioning $\mathbb{O}$,
$\Omega^{\dagger}\Omega=\sum_{k=0}^{K}Y_k^{\dagger}Y_k
=\defic+\sum_{k=1}^{K}A_k=\id_{\Inp}$ by \eqref{eq:deficiency}. The ports
$\mathbb{O}_1,\dots,\mathbb{O}_K$ are pairwise disjoint, so this is
\eqref{eq:port-dilation}. For the second sentence, $D^{1/2}$ is the real
diagonal matrix of the positive square roots $\sqrt{d_i}$, so
$Y^{\dagger}Y=E^{\dagger}D^{1/2}D^{1/2}E=E^{\dagger}DE=X$. The entry $(i,l)$
of $Y$ is $\sqrt{d_i}\,E_{il}$, the product of $\sqrt{d_i}$, which a
square-root enclosure of $d_i$ encloses, and of $E_{il}$, enclosed exactly
by itself; part (ii) of Lemma~\ref{lem:enclosure-arithmetic} encloses it,
exactly where the square-root enclosure is exact, which
Definition~\ref{def:complex-enclosure} makes it where $d_i$ is the square of
a rational.
\end{proof}

% claim receiver/def:exact-route
\claimsummary{receiver/def:exact-route}{A declaration is on the exact route when the port Grams its declared data determine, computed exactly, have Gaussian-rational entries and are an admissible Gram family exactly, and its probe is of the Gaussian tier. A declaration of Grams or of an instrument matrix admissible exactly is one.}
\claimbookmark{receiver/def:exact-route}{Definition}{def:exact-route}{Gaussian-rational port Grams and the exact route}
\claimleadin{%
The receiver's blocks are enclosed for every declaration
(Theorem~\ref{thm:enclosed-blocks}). Where the port Grams are
Gaussian-rational, many entries of the blocks, every entry of a block of at
most one photon among them, are Gaussian rationals that can be known exactly
(Theorem~\ref{thm:exact-blocks}), and a snap can keep such a block at weight
zero only if it is known exactly. The route is decided on the Grams and not
on the instrument matrix: a half-wave plate at $\pi/8$ has entries
$\pm1/\sqrt2$, yet behind an ideal polarizing beamsplitter the port Grams it
produces are rational.
}
\begin{definition}[Gaussian-rational port Grams and the exact route]\label{def:exact-route}
Let $\Recv$ be a \termref{term-receiver-instrument}{receiver instrument} with
\termref{term-port-gram-matrices}{port Grams} $\{A_k\}_{k=1}^{K}$ on
$\mathbb{C}^{\Inp}$ (Definitions~\ref{def:instrument} and~\ref{def:gram}),
and write $\mathbb{Q}(i)=\{x+iy:x,y\in\mathbb{Q}\}$. A declaration of a
receiver is on the
\termdef{term-exact-route}{exact route} when (i) the port Grams $\{A_k\}$ that its
\exttermref{../foundations/contract.pdf}{term-declared-data}{declared data} determine, computed in exact arithmetic, have entries in
$\mathbb{Q}(i)$, are Hermitian, and are an \termref{term-admissible-gram-family}{admissible Gram family}
exactly (Definition~\ref{def:gram}), that is,
\begin{equation}\label{eq:exact-admissible}
  A_k\succeq0\quad(k=1,\dots,K)
  \qquad\text{and}\qquad
  \id-A_{[K]}\succeq0 ,
\end{equation}
conditions on Gaussian-rational matrices; and (ii) its
\termref{term-probe-declaration}{probe declaration}, where it has one, is of the \termref{term-gaussian-tier}{Gaussian tier}
(Definition~\ref{def:probe}), whose mean photon numbers $\bar n^{(t)}_m$ are
declared values and so rationals (\xref{foundations/def:binary64}\,(iii)).
Under (i) the port Grams are those of a receiver instrument
(Theorem~\ref{thm:representation-independence}\,(ii)), and every receiver
instrument with these port Grams induces the same click-pattern POVM
(Theorem~\ref{thm:representation-independence}\,(i)); where the
\termref{term-instrument-matrix}{instrument matrix} $\Gmat$, declared or compiled, satisfies
$\Gmat^\dagger\Gmat\preceq\id$ exactly, the declared instrument is one of
them. A declared family admitted only at the \termref{term-receiver-tolerances}{receiver tolerances}
(Definition~\ref{def:tolmap}) and not admissible exactly is not on the exact
route. A port Gram family or an instrument matrix declared directly has
entries that are dyadic rationals (\xref{foundations/def:binary64}\,(iii));
its port Grams then have entries in $\mathbb{Q}(i)$, by \eqref{eq:gramdef}
for a declared instrument matrix, and such a declaration satisfies (i)
exactly when its port Grams are Hermitian, as they always are for a declared
instrument matrix, and admissible exactly.
\end{definition}
\claimmeta{\metalabel{Uses}\mref{Def.}{def:instrument}, \mref{Def.}{def:gram}, \mref{Def.}{def:probe}, \mref{Thm.}{thm:representation-independence}, \mref{Def.}{def:tolmap}, \extref{../foundations/contract.pdf}{foundations/def:binary64}
\metasep \metalabel{Used by}\mref{Thm.}{thm:exact-blocks}, \mref{Def.}{def:evaluation-model}, \mref{Thm.}{thm:exact-and-evaluated}, \mref{Req.}{req:exact-route-decision}, \mref{Rem.}{rem:noise-floor-receiver}, \mref{Req.}{req:snap-fix}
\metasep newly derived
\metasep implemented (effective: accepted)
\metasep \metalabel{Code}\coderef{https://github.com/fraunhofersg/certiqs-key/blob/v0.1.0/src/certiqs_key/receiver/reference.py}{certiqs_key.receiver.reference:ExactRouteDecision}}

% claim receiver/thm:exact-blocks
\claimsummary{receiver/thm:exact-blocks}{With Gaussian-rational port Grams, every block entry whose two occupation vectors have factorial products forming a square, every entry of a block of at most one photon among them, is Gaussian-rational and enclosed exactly. Up to one photon the reference is then exact.}
\claimbookmark{receiver/thm:exact-blocks}{Theorem}{thm:exact-blocks}{Exact entries of the blocks on the exact route}
\claimleadin{%
A snap keeps a block only where its reference is exact on the block and
every element is already a \exttermref{../foundations/contract.pdf}{term-lattice-value}{lattice value}
(\xref{foundations/req:snap-reference}\,(ii)). The theorem says where the
enclosures of Theorem~\ref{thm:enclosed-blocks} are exact: on every entry
that the factorials of the occupation basis do not make irrational. Above one
photon the off-diagonal entries generally carry a factor $1/\sqrt2$ or
similar, so those blocks have non-zero radii even at Gaussian-rational
Grams.
}
\begin{theorem}[Exact entries of the blocks on the exact route]\label{thm:exact-blocks}
Let $\Recv$ be a \termref{term-receiver-instrument}{receiver instrument} whose
\termref{term-port-gram-matrices}{port Grams} $\{A_k\}$ have entries in
$\mathbb{Q}(i)$, as on the \termref{term-exact-route}{exact route}
(Definition~\ref{def:exact-route}), with click-pattern POVM $\{M_S\}$, a
\termref{term-declared-coarse-graining}{declared coarse-graining} $\cg$ and emitted POVM $\{F_v\}$. Let $n\ge0$
and let $\vec m,\vec m'$ be occupation vectors of weight $n$ such that
$\vec m!\,\vec m'!$ is the square of an integer; this holds on the diagonal,
$\vec m=\vec m'$, and for every $\vec m,\vec m'$ when $n\le1$.
\begin{enumerate}[label=(\roman*),leftmargin=3.0em]
\item The entry $(\vec m,\vec m')$ of every block $\Pvac{J}^{(n)}$,
  $\blk{M}{\eqn}_S$ and $\blk{F}{\eqn}_v$ lies in $\mathbb{Q}(i)$.
\item If the enclosures of the $A_k$ in Theorem~\ref{thm:enclosed-blocks}\,(ii)
  are exact, that is, each $A_k$ is enclosed by itself with radius zero,
  then the enclosure of the entry $(\vec m,\vec m')$ of each of these blocks
  in Theorem~\ref{thm:enclosed-blocks}\,(ii), and of the blocks
  $\blk{F}{\eqn}_v$ in Theorem~\ref{thm:enclosed-blocks}\,(iii), is exact,
  its centre being the entry.
\item Under the hypothesis of (ii), every block $n\le1$ of the reference of
  Theorem~\ref{thm:enclosed-blocks}\,(iii) has radius zero everywhere, and for
  a cutoff $N\le1$ that reference is an \exttermref{../foundations/contract.pdf}{term-exact-reference}{exact reference}
  (\xref{foundations/def:exact-reference}).
\end{enumerate}
\end{theorem}
\claimmeta{\metalabel{Uses}\mref{Thm.}{thm:click-povm}, \mref{Thm.}{thm:coarse-graining}, \mref{Def.}{def:complex-enclosure}, \mref{Lem.}{lem:enclosure-arithmetic}, \mref{Lem.}{lem:fock-permanent}, \mref{Thm.}{thm:enclosed-blocks}, \mref{Def.}{def:exact-route}, \extref{../foundations/contract.pdf}{foundations/def:exact-reference}
\metasep \metalabel{Used by}\mref{Thm.}{thm:exact-and-evaluated}, \mref{Req.}{req:exact-route-decision}, \mref{Rem.}{rem:noise-floor-receiver}, \mref{Req.}{req:dark-composed-reference}, \mref{Req.}{req:snap-fix}
\metasep newly derived
\metasep accepted}

\begin{proof}
(i) Every entry of $\id-A_J$ lies in $\mathbb{Q}(i)$, so
Lemma~\ref{lem:fock-permanent}\,(ii) puts the entry $(\vec m,\vec m')$ of
$\Pvac{J}^{(n)}=\Fock(\id-A_J)^{(n)}$ (Theorem~\ref{thm:click-povm}\,(ii)) in
$\mathbb{Q}(i)$. By \eqref{eq:blockMS} and
Theorem~\ref{thm:coarse-graining}\,(i), the entry $(\vec m,\vec m')$ of
$\blk{M}{\eqn}_S$ and of $\blk{F}{\eqn}_v$ is an integer combination of the
entries $(\vec m,\vec m')$ of blocks $\Pvac{J}^{(n)}$, and $\mathbb{Q}(i)$ is
closed under such combinations.

(ii) With exact enclosures of the $A_k$, every entry of $\id-A_J$ and every
permanent formed from them has an exact enclosure
(Lemma~\ref{lem:enclosure-arithmetic}\,(iv)). If
$\vec m!\,\vec m'!=s^{2}$ with $s$ a positive integer, then
$1/(\vec m!\,\vec m'!)=(1/s)^{2}$ is the square of a rational, so its
\termref{term-square-root-enclosure}{square-root enclosure} is exact (Definition~\ref{def:complex-enclosure}),
and the product of two exactly enclosed numbers is exactly enclosed. The
entries $(\vec m,\vec m')$ of $\blk{M}{\eqn}_S$ and $\blk{F}{\eqn}_v$ are
enclosed by integer combinations of these exact enclosures, hence exactly. The
condition on $\vec m!\,\vec m'!$ is symmetric in $\vec m$ and $\vec m'$, so
the entry $(\vec m',\vec m)$ is exact as well, and the Hermitian centre of
Lemma~\ref{lem:enclosure-arithmetic}\,(v) keeps the entry exact. An exact
enclosure has as its centre the number it encloses.

(iii) For $n\le1$ every pair $\vec m,\vec m'$ satisfies the condition, so by
(ii) every radius of the block is zero. For $N\le1$ every block of the
reference is of this kind, and a reference all of whose radii are zero is
exact.
\end{proof}

\section{Residuals, tolerances and floating-point evaluation}\label{group-receiver/def:residuals}
\sectionintro{Every statement of this contract is exact unless it names a receiver tolerance or the eigenvalue-slack assumption; its evaluation is binary64 floating point, beside a rigorous evaluation of the weights in exact rational arithmetic wherever the port Grams are enclosed (\hyperlink{claim-receiver/def:evaluation-model}{Definition~\ref*{def:evaluation-model}}). Six residuals check the floating-point evaluation's consistency, each with an exact value of zero, and the \termref{term-receiver-tolerances}{receiver tolerances}, numbers an instance declares to the receiver as a supplier, decide each exact condition on declared data (\hyperlink{claim-receiver/def:tolmap}{Definition~\ref*{def:tolmap}}), with the direction of the unclamped bound-data constant stated on both sides. None of this encloses a floating-point number: the one relation between a floating-point and an exact quantity is the opt-in eigenvalue-slack assumption (\hyperlink{claim-receiver/thm:exact-and-evaluated}{Theorem~\ref*{thm:exact-and-evaluated}}), and the floating-point blocks are evaluations of the model's blocks, which the rational enclosures of the previous section contain, while a rigorously evaluated weight is a bound by theorem. The last remark records in which direction each input of the bound is enclosed for the rigorous evaluation.}
\sectionimports{\extref{../foundations/contract.pdf}{foundations/cor:least-eigenvalue-over-enclosure}, \extref{../foundations/contract.pdf}{foundations/def:binary64}, \extref{../foundations/contract.pdf}{foundations/def:exact-reference}, \extref{../foundations/contract.pdf}{foundations/lem:enclosure-psd}, \extref{../foundations/contract.pdf}{foundations/lem:least-eigenvalue-bound}}

% claim receiver/def:residuals
\claimsummary{receiver/def:residuals}{Six residuals, each exactly zero in exact arithmetic, check contractivity, Gram positivity, the deficiency identity, POVM completeness, block Hermiticity and positivity, and window consistency. Which exist depends on the input route and on whether an exact window is requested.}
\claimbookmark{receiver/def:residuals}{Definition}{def:residuals}{The six residuals}
\claimleadin{%
Route (iii) changes the nature of the arithmetic. Where the Grams are
\exttermref{../foundations/contract.pdf}{term-declared-data}{declared data}, the only numerical step is an eigenvalue computation. Where they are
\emph{compiled}, $\Gmat$ is a product of many matrices, and $A_k$, $\lmin$ and
$\lmax$ are derived quantities whose floating-point error compounds through the
product. The residuals below check the consistency of that evaluation; which of
them the artifact carries, and what a residual above its tolerance causes, are
Requirements~\ref{req:residual-emission} and~\ref{req:residual-refusal}.
Remark~\ref{rem:exacteig} states which direction of failure the window
residual $r_6$ is blind to.
}
\begin{definition}[The six residuals]\label{def:residuals}
The \termdef{term-residuals}{residuals} of an evaluation of this contract are those of the six quantities
below that exist on the input route taken and for the evaluation requested
(which exist is stated after \eqref{eq:rhodef}):
\begin{align}
  r_1&\coloneqq\big\lVert\big(\Gmat^\dagger\Gmat-\id\big)_{+}\big\rVert
    &&\text{(contractivity defect: the positive part must be }0),
    \label{eq:res1}\\
  r_2&\coloneqq\max_k\ \max\big\{0,\,-\lmin(A_k)\big\}
    &&\text{(Gram positivity)},\label{eq:res2}\\
  r_3&\coloneqq\big\lVert\,\id_{\Inp}-A_\Sigma
     -\Gmat_{\Inp}^\dagger\Emd{\mathrm{un}}\Gmat_{\Inp}
     -\big(\id_{\Inp}-\Gmat_{\Inp}^\dagger\Gmat_{\Inp}\big)\big\rVert
    &&\text{(deficiency identity \eqref{eq:deficiencyid})},\label{eq:res3}\\
  r_4&\coloneqq\Big\lVert\,\id-\sum_{S\subseteq[K]}
     \blk{M}{\eqn}_{S|t}\Big\rVert\ \ \text{for each emitted }n
     \text{ and each setting }t
    &&\text{(POVM completeness, blockwise)},\label{eq:res4}\\
  r_5&\coloneqq\max_{v,\,n,\,t}\ \varrho\big(\blk{M}{\eqn}_{v|t}\big)
    &&\text{(block Hermiticity and positivity)},\label{eq:res5}\\
  r_6&\coloneqq\max_{t,\ N<n\le m^{\ast}_t}
     \big(\Bnd^{(t)}(n)-\widehat\Bnd^{\ast(t)}(n)\big)_{+}
    &&\text{(window consistency, Lem.~\ref{lem:exactbasic})},\label{eq:res6}
\end{align}
in which the \termdef{term-hermiticity-and-positivity-defect}{Hermiticity-and-positivity defect} of a square matrix $X$ is
\begin{equation}\label{eq:rhodef}
  \varrho(X)\ \coloneqq\ \max\Big\{\,\big\lVert X-X^\dagger\big\rVert,\ \
    \big(-\lmin\big(\tfrac12(X+X^\dagger)\big)\big)_{+}\,\Big\} ,
\end{equation}
where $\Emd{\mathrm{un}}$ projects onto the declared modes in no port,
$(\cdot)_+$ is the positive part, and $\lVert\cdot\rVert$ is the spectral norm
throughout \eqref{eq:res1}--\eqref{eq:rhodef}. The maximum in \eqref{eq:res5}
runs over exactly the blocks $\blk{M}{\eqn}_{v|t}$ of the coarse-grained
elements, before any \termref{term-dark-count-channel}{dark-count channel}, that the artifact emits: every
emitted outcome letter
$v\in\Vset$, every photon number $0\le n\le N$, and every setting $t\in\Tset$
of the \termref{term-setting-family}{setting family} (Definition~\ref{def:setting-family}; one setting when
$\InpS=\varnothing$, Corollary~\ref{cor:unconditioned-family}); the emitted
bound of the one-setting family is Corollary~\ref{cor:emitted-unconditioned},
and that of a setting family a submodule constructs is that submodule's
instance of Theorem~\ref{thm:emitted-bound}. Here $\blk{M}{\eqn}_{v|t}$ is the
emitted block: the $n$-photon block of the coarse-grained conditional element
$F_{v|t}$, written $\blk{F}{\eqn}_v$ in Theorem~\ref{thm:coarse-graining}\,(i) at a single
setting.

\emph{Which residuals exist on which route.} $r_2$, $r_4$ and $r_5$ are
functions of the port Grams and of the emitted blocks, which every route
produces, so all three exist on all three routes. $r_1$ \eqref{eq:res1} and
$r_3$ \eqref{eq:res3} are written on $\Gmat$ and $\Gmat_{\Inp}$, and route (i)
supplies no \termref{term-instrument-matrix}{instrument matrix}: there the Grams are the declared datum and no
$\Gmat$ exists, so $r_1$ and $r_3$ are \emph{not defined}. Synthesizing a
$\Gmat$ from Theorem~\ref{thm:representation-independence}\,(ii) to evaluate them on would test the
synthesis rather than the declaration. This agrees with
Theorem~\ref{thm:exact-and-evaluated}\,(iv), where route (i) is the case in which ``there is
no compilation''. On routes (ii) and (iii) all five of $r_1,\dots,r_5$ exist.

\emph{$r_6$ depends on the request, not on the route.} It compares the two
evaluations of a single proved inequality, $\Bnd^{(t)}(n)\le
\Bnd^{\ast(t)}(n)$ (Lemma~\ref{lem:exactbasic}), so its exact value is $0$
like the other five, and it exists on all three routes but \emph{only at the
settings whose declaration carries an \termref{term-exact-window-depth}{exact-window depth}}
(Definition~\ref{def:pmcex}), over exactly the indices $N<n\le m^{\ast}_t$ that
depth names. A declaration that requests no exact window computes no
$\widehat\Bnd^{\ast}$ and has nothing to compare, so $r_6$ is not defined
there, as $r_1$ and $r_3$ are not on route (i).

No residual bounds the distance between a computed quantity and the exact
one (Theorem~\ref{thm:exact-and-evaluated}\,(ii)).
\end{definition}
\claimmeta{\metalabel{Uses}\mref{Lem.}{lem:deficiency}, \mref{Thm.}{thm:representation-independence}, \mref{Thm.}{thm:coarse-graining}, \mref{Lem.}{lem:exactbasic}, \mref{Def.}{def:setting-family}, \mref{Cor.}{cor:unconditioned-family}, \mref{Def.}{def:pmcex}
\metasep \metalabel{Used by}\mref{Prop.}{prop:res5}, \mref{Rem.}{rem:res5why}, \mref{Def.}{def:tolmap}, \mref{Rem.}{rem:tolerancedomain}, \mref{Thm.}{thm:exact-and-evaluated}, \mref{Req.}{req:residual-emission}, \mref{Req.}{req:residual-refusal}, \mref{Rem.}{rem:wrongcheck}, \mref{Def.}{def:darktol}
\metasep newly derived
\metasep implemented (effective: accepted)
\metasep \metalabel{Code}\coderef{https://github.com/fraunhofersg/certiqs-key/blob/v0.1.0/src/certiqs_key/receiver/residuals.py}{certiqs_key.receiver.residuals:evaluate_residuals}}

% claim receiver/prop:res5
\claimsummary{receiver/prop:res5}{The defect of a matrix vanishes exactly when it is Hermitian and positive semidefinite, and bounds how far its expectation values stray below zero or off the real axis. Every emitted block before the dark counts has defect zero exactly, so the fifth residual measures the evaluation alone.}
\claimbookmark{receiver/prop:res5}{Proposition}{prop:res5}{Properties of the Hermiticity-and-positivity defect}
\begin{proposition}[Properties of the Hermiticity-and-positivity defect]\label{prop:res5}
Let $X$ be a complex $D\times D$ matrix and $\varrho(X)$ its
\termref{term-hermiticity-and-positivity-defect}{Hermiticity-and-positivity defect} \eqref{eq:rhodef}.
\begin{enumerate}[label=(\roman*),leftmargin=3.0em]
\item $\varrho(X)=0$ if and only if $X=X^\dagger$ and $X\succeq0$.
\item If $X=X^\dagger$ then $\varrho(X)=\big(-\lmin(X)\big)_+$, which is the
  form \eqref{eq:res2} takes on a Gram.
\item For every unit vector $\ket{\psi}$,
  $\Re\bra{\psi}X\ket{\psi}\ge-\varrho(X)$ and
  $\big|\Im\bra{\psi}X\ket{\psi}\big|\le\tfrac12\varrho(X)$. Hence
  $\varrho(X)\le\varepsilon$ bounds every diagonal expectation of $X$ away from
  the negative reals by $\varepsilon$ and off the real axis by
  $\varepsilon/2$.
\item Every block $\blk{M}{\eqn}_{v|t}$ of the coarse-grained elements, before
  any dark-count channel, that the artifact emits satisfies $\varrho=0$
  \emph{exactly}: for
  each $v$, $n$ and $t$, $\blk{M}{\eqn}_{v|t}$ is a diagonal block of the
  positive operator $F_{v|t}$ (Definition~\ref{def:setting-family}\,(a) and
  Theorem~\ref{thm:coarse-graining}), and a diagonal block --- the compression $\Pi X\Pi$ of a positive
  operator $X$ by an orthogonal projector $\Pi$, read on the range of $\Pi$ ---
  is Hermitian and positive semidefinite. So $r_5$ is a residual in exactly the
  sense of the other five: the exact value is $0$, and what is emitted is the
  floating-point evaluation.
\end{enumerate}
\end{proposition}
\claimmeta{\metalabel{Uses}\mref{Thm.}{thm:coarse-graining}, \mref{Def.}{def:setting-family}, \mref{Def.}{def:residuals}
\metasep \metalabel{Used by}\mref{Req.}{req:residual-emission}
\metasep newly derived
\metasep implemented (effective: accepted)
\metasep \metalabel{Code}\coderef{https://github.com/fraunhofersg/certiqs-key/blob/v0.1.0/src/certiqs_key/receiver/numerics.py}{certiqs_key.receiver.numerics:hermiticity_positivity_defect}}

\begin{proof}
Write $H\coloneqq\tfrac12(X+X^\dagger)$ and $S\coloneqq X-X^\dagger$, so that
$X=H+\tfrac12S$ with $H$ Hermitian and $S$ skew-Hermitian.

(i) If $\varrho(X)=0$ then $\lVert S\rVert=0$, so $X=X^\dagger$ and $H=X$; and
$(-\lmin(X))_+=0$ gives $\lmin(X)\ge0$, i.e.\ $X\succeq0$. Conversely if
$X=X^\dagger\succeq0$ then $S=0$ and $\lmin(X)\ge0$, so both entries of the
maximum vanish.

(ii) $S=0$ gives $\lVert S\rVert=0$ and $H=X$, so the maximum is its second
entry.

(iii) $\bra{\psi}X\ket{\psi}=\bra{\psi}H\ket{\psi}
+\tfrac12\bra{\psi}S\ket{\psi}$. The first term is real, with
$\bra{\psi}H\ket{\psi}\ge\lmin(H)\ge-\varrho(X)$ by the definition of
$\varrho$; the second is purely imaginary, because $iS$ is Hermitian and
therefore $\bra\psi S\ket\psi\in i\mathbb{R}$, with modulus at most
$\tfrac12\lVert S\rVert\le\tfrac12\varrho(X)$. Splitting into real and
imaginary parts gives both inequalities.

(iv) A POVM element is positive semidefinite, and so is any sum of them, so
$F_{v|t}=\sum_{S\in\cg^{-1}(v)}M_{S|t}\succeq0$. Each $F_{v|t}$ is block
diagonal in the photon number (Definition~\ref{def:setting-family}\,(a) and
Theorem~\ref{thm:coarse-graining}), so $\blk{M}{\eqn}_{v|t}$ is the compression
$\Pi_nF_{v|t}\Pi_n$ read on the range of $\Pi_n$. For any $Y\succeq0$ and any
orthogonal projector $\Pi$, $\Pi Y\Pi$ is Hermitian and
$\bra{\psi}\Pi Y\Pi\ket{\psi}=\bra{\Pi\psi}Y\ket{\Pi\psi}\ge0$, so the
compression is Hermitian positive semidefinite. Now apply (i).
\end{proof}

% claim receiver/rem:res5why
\claimsummary{receiver/rem:res5why}{The fifth residual is the only one that looks at the emitted photon-number blocks, which can fail Hermiticity or positivity at the last step while the other residuals pass. It is neither a repair nor an enclosure.}
\claimbookmark{receiver/rem:res5why}{Remark}{rem:res5why}{What the fifth residual checks that the others do not}
\begin{remark}[What the fifth residual checks that the others do not]\label{rem:res5why}
Four of the other five residuals of Definition~\ref{def:residuals} check the
\emph{optics}: that the compiled matrix is a contraction, that the Grams are
positive, that the loss identity closes, that the POVM sums to the identity.
None of them looks at the object the artifact actually emits, the
photon-number block, and a block can fail Hermiticity or positivity at the last
floating-point step even when all four pass, because $r_4$ constrains only the
\emph{sum} over $S$ and $r_2$ lives on the Grams, one compilation stage
earlier. $r_5$ makes that defect visible where the receiver is in hand, with
the port, the setting and the photon number that produced it.

Two things $r_5$ is not. It is not a repair: a block that fails is refused and
never repaired (Requirement~\ref{req:residual-refusal}), because a repaired
block is a different operator and no claim of this contract holds for it. And
it is not an enclosure: $r_5=0$ says that the \emph{computed} block is
Hermitian and positive, not that it is close to the exact one, which is
Theorem~\ref{thm:exact-and-evaluated}\,(ii) exactly as for $r_1,\dots,r_4$.
\end{remark}
\claimmeta{\metalabel{Uses}\mref{Def.}{def:residuals}
\metasep \metalabel{Used by}\mref{Rem.}{rem:d2slack}
\metasep newly derived
\metasep accepted}

% claim receiver/def:tolmap
\claimsummary{receiver/def:tolmap}{Each exact condition on declared data that the receiver decides in floating point has one receiver tolerance and one test, and every other validation is a domain comparison, a finite inspection or a residual test. Admitting an instrument matrix does not admit the bound data derived from it.}
\claimbookmark{receiver/def:tolmap}{Definition}{def:tolmap}{Receiver tolerances and the exact conditions they decide}
\claimleadin{%
The receiver is a \exttermref{../foundations/contract.pdf}{term-supplier}{supplier} in the sense of
\xref{foundations/def:declared-data}\,(ii): it decides its own conditions on the
data it reads and emits in floating point, at slacks an instance declares to
it. Those slacks are declared numbers of their own, never the instance's
\exttermref{../foundations/contract.pdf}{term-declared-tolerance}{declared tolerance} (\xref{foundations/req:tolerance-scope}\,(iii)); this
definition names them and says which condition each one decides. The
floating-point evaluation model in which they decide is
Definition~\ref{def:evaluation-model}. Among the engine's obligations behind the
validations of kinds (a) to (c) are the reserved tolerance field
(Requirement~\ref{req:reserved-tolerance-field}), the exclusive dark-count
routes (Requirement~\ref{req:dark-route-exclusive}), the exact-window
declaration (Requirement~\ref{req:exact-window-declaration}) and the residual
refusal (Requirement~\ref{req:residual-refusal}); Remark~\ref{rem:exacteig}
says why $\varepsilon_{\mathrm{eig}}$ is no test.
}
\begin{definition}[Receiver tolerances and the exact conditions they decide]
\label{def:tolmap}
Several conditions this contract states are \emph{exact} statements about
\exttermref{../foundations/contract.pdf}{term-declared-data}{declared data} that this contract nevertheless decides in floating point. Each is decided by a declared number, fixed here,
and by no other threshold. The \termdef{term-receiver-tolerances}{receiver tolerances} are the declared
$\varepsilon_1,\dots,\varepsilon_6$, $\varepsilon_{\mathrm{gap}}$ and
$\varepsilon_{\mathrm{eig}}$, each in the declared domain $[0,1)$; every quantitative statement below uses
that upper bound and says where. Items (i) to (v) spend $\varepsilon_1$,
$\varepsilon_2$ and $\varepsilon_4$; $\varepsilon_3$, $\varepsilon_5$,
$\varepsilon_6$ and $\varepsilon_{\mathrm{gap}}$ decide no exact condition on
declared data and bound computed quantities only (kind (c) below); and
$\varepsilon_{\mathrm{eig}}$ decides nothing at all. It is subtracted from a
computed eigenvalue in \eqref{eq:elldef}, so it is the one declared number this
contract spends on a \emph{hypothesis} (Assumption~\ref{ass:eig}) rather than
on a test.
\begin{enumerate}[label=(\roman*),leftmargin=3.0em]
\item \emph{Route-(i) Gram admissibility}, $A_k\succeq0$ and
  $\sum_kA_k\preceq\id$ (Definition~\ref{def:gram}), is decided at
  $\varepsilon_2$: the family is admitted when $\lmin(A_k)\ge-\varepsilon_2$
  for every $k$ and $\lmax(A_\Sigma)\le1+\varepsilon_2$. $\varepsilon_2$ is the
  tolerance the contract already attaches to Gram positivity through
  \eqref{eq:res2}, so the same declared number bounds the residual and decides
  the admission, in the same units.
\item \emph{The route-(i) rank consistency check} $\rk A_k\le|\Port{k}|$ is
  decided by the $\varepsilon_2$-rank
  \begin{equation}\label{eq:eps2rank}
    \rk_{\varepsilon_2}(X)\ \coloneqq\
    \#\big\{\,i:\ \lambda_i\big(\tfrac12(X+X^\dagger)\big)>\varepsilon_2\,\big\}
    \qquad\text{(eigenvalues with multiplicity)},
  \end{equation}
  the check being $\rk_{\varepsilon_2}(A_k)\le|\Port{k}|$ for every $k$. The
  threshold is \emph{absolute} and it is $\varepsilon_2$; no separate tolerance
  is added.
\item \emph{Dark-channel admissibility} for a \emph{declared} $R$ is decided
  in the detector-imperfections submodule, Definition~\ref{def:darktol}, which
  extends this list; the tolerances that definition spends are the
  $\varepsilon_4$ and $\varepsilon_2$ of this list.
\item \emph{Bound-data admissibility}, $0\le c$ in \eqref{eq:bounddata}, is
  decided at $\varepsilon_2$: the data are admitted when
  $c=1-\lmax(B)\ge-\varepsilon_2$. \emph{Nothing is clamped}: the computed
  $c$ enters \eqref{eq:beta} exactly as it stands. The remaining inequalities of
  \eqref{eq:bounddata}, $\gamma\le\beta_k\le1$ and $c\le d_k\le1$, are
  \emph{not} separately tested and no threshold decides them: they are
  redundant by theorem, discharged by the admissibility paragraph of
  Theorem~\ref{thm:multiclick-bound} from \eqref{eq:bbkhyp} and from
  $\gamma\le\beta_k$, both hypotheses of the setting family
  (Definition~\ref{def:setting-family}\,(b)), and their floating-point
  evaluation is governed by Theorem~\ref{thm:exact-and-evaluated}\,(ii) like every other
  computed quantity.

  \emph{Item (iv) decides this admission alone, on every input route, and
  item~(v) does not decide it.} Whenever $\varepsilon_1>\varepsilon_2$ the band
  \begin{equation}\label{eq:tolmapband}
    c\in\big[-\varepsilon_1,\,-\varepsilon_2\big)
  \end{equation}
  is non-empty, admitted by item~(v)'s test and refused by item (iv)'s, and
  item (iv) wins on it: the band is refused.
\item \emph{Route-(ii) contractivity, and (E1) at catalogue entry C10.}
  $\Gmat^\dagger\Gmat\preceq\id$ (Definition~\ref{def:instrument}) is decided
  at $\varepsilon_1$: the matrix is admitted when
  $\lmax(\Gmat^\dagger\Gmat)\le1+\varepsilon_1$. The same test decides the (E1)
  predicate of the catalogue's \termref{term-generic-block}{generic block} entry C10
  (Definition~\ref{def:c10-tolerance} of the component-catalogue submodule).

  The tolerance that bounds the residual $r_1$ also decides this admission.
  $\Gmat^\dagger\Gmat$ is Hermitian positive semidefinite, so
  $\big(\Gmat^\dagger\Gmat-\id\big)_+$ has eigenvalues
  $\big(\lambda_i(\Gmat^\dagger\Gmat)-1\big)_+$ and therefore
  \begin{equation}\label{eq:r1islmax}
    r_1=\big\lVert\big(\Gmat^\dagger\Gmat-\id\big)_{+}\big\rVert
      =\big(\lmax(\Gmat^\dagger\Gmat)-1\big)_{+} ,
  \end{equation}
  so ``$r_1\le\varepsilon_1$'' and ``$\lmax(\Gmat^\dagger\Gmat)\le
  1+\varepsilon_1$'' are the \emph{same test on the same number}, and no
  separate tolerance is added.
\end{enumerate}
\emph{The five items, together with the two extensions the submodules declare
in the same form, Definition~\ref{def:darktol}, which is item (iii)'s content,
and Definition~\ref{def:c10-tolerance}, are the complete set of exact conditions
about \emph{declared} data that this module and its submodules decide under a
threshold.} Every other validation they state is one of three kinds, and none
of them adds a threshold this definition has not named:
\begin{enumerate}[label=(\alph*),leftmargin=2.2em]
\item a comparison of declared data against a declared domain: every parameter
  domain of Table~\ref{tab:catalogue} except C10's, the detector-layer ranges,
  $m_t>N$, $m^{\ast}_t>N$, $m^{\ast}_t\le m_t$, $N\ge0$, the domain of the
  receiver tolerances itself, the test of
  the reserved tolerance field of Assumption~\ref{ass:exact}, and the diagonal
  non-negativity of $\Gamma^{(t)}$;
\item a combinatorial predicate on the declaration decided by finite
  inspection: (E2) and (E3) of Definition~\ref{def:graph}, route exclusivity,
  the port map's disjointness, totality and surjectivity of $\cg$, the
  single-valued $\tau$ of Definition~\ref{def:modes}, the exclusivity of the
  two dark-count routes of Definition~\ref{def:darkadm}, predicate~(X) of
  Definition~\ref{def:pmc}, and the conditional presence of
  $\varepsilon_{\mathrm{eig}}$ and $\varepsilon_6$ at a declaration carrying an
  exact-window depth;
\item the residual tests themselves: $r_3$, $r_4$, $r_5$, $r_6$ against
  $\varepsilon_3$, $\varepsilon_4$, $\varepsilon_5$, $\varepsilon_6$, and the
  residual gap $g_{N,m_t}$ against $\varepsilon_{\mathrm{gap}}$. These are
  threshold decisions, but on \emph{computed} quantities rather than on
  declared data, and each is a declared-threshold decision, for the residuals of
  Definition~\ref{def:residuals} and for the residual gap \eqref{eq:resgap},
  so nothing undeclared decides them either. The tests $r_1\le\varepsilon_1$
  and $r_2\le\varepsilon_2$ are of this kind too, except where they coincide
  with an item: on route (ii), where $\Gmat$ is declared, the test
  $r_1\le\varepsilon_1$ \emph{is} item (v) by \eqref{eq:r1islmax}, and on
  route (i), where the Grams are declared, $r_2\le\varepsilon_2$ is the
  positivity half of item (i) written as a residual; on route (iii) $r_1$
  tests the compiled $\Gmat$, and on routes (ii) and (iii) $r_2$ tests
  computed Grams. Neither is
  $\varepsilon_{\mathrm{eig}}$, which is not a test of any kind: it is
  subtracted from $\widehat\Bnd^{\ast}$ in \eqref{eq:elldef} and changes an
  emitted value, not an admission.
\end{enumerate}
\end{definition}
\claimmeta{\metalabel{Uses}\mref{Def.}{def:instrument}, \mref{Def.}{def:gram}, \mref{Def.}{def:graph}, \mref{Def.}{def:bounddata}, \mref{Thm.}{thm:multiclick-bound}, \mref{Lem.}{lem:windowtail}, \mref{Def.}{def:setting-family}, \mref{Def.}{def:pmc}, \mref{Def.}{def:pmcex}, \mref{Def.}{def:residuals}
\metasep \metalabel{Used by}\mref{Def.}{def:exact-route}, \mref{Rem.}{rem:bound-data-band}, \mref{Rem.}{rem:tolerancedomain}, \mref{Rem.}{rem:ranktwosided}, \mref{Rem.}{rem:cdirection}, \mref{Thm.}{thm:exact-and-evaluated}, \mref{Rem.}{rem:ia}, \mref{Req.}{req:tolerance-tests}, \mref{Req.}{req:exact-route-decision}, \mref{Def.}{def:c10-tolerance}, \mref{Def.}{def:darkadm}, \mref{Def.}{def:darktol}, \mref{Prop.}{prop:dark-slack-bound}, \mref{Rem.}{rem:rmintrap}, \mref{Rem.}{rem:d2slack}, \mref{Rem.}{rem:corrdark}, \mref{Req.}{req:dark-composed-reference}
\metasep newly derived
\metasep implemented (effective: accepted)
\metasep \metalabel{Code}\coderef{https://github.com/fraunhofersg/certiqs-key/blob/v0.1.0/src/certiqs_key/receiver/numerics.py}{certiqs_key.receiver.numerics:tolerances_from_declaration}}

% claim receiver/rem:bound-data-band
\claimsummary{receiver/rem:bound-data-band}{Where the contractivity test at \ensuremath{\varepsilon}\ensuremath{_{1}} admits and the bound-data test at \ensuremath{\varepsilon}\ensuremath{_{2}} refuses, the bound-data test decides: an exact implication between the two conditions transports no admission between computed numbers. Admitting the band would make a stated error term false on admitted values.}
\claimbookmark{receiver/rem:bound-data-band}{Remark}{rem:bound-data-band}{The bound-data test and the contractivity test on the band}
\claimleadin{%
Items (iv) and (v) of Definition~\ref{def:tolmap} both widen an exact
condition by a receiver tolerance, and in exact arithmetic the condition of
item (v) implies that of item (iv). This remark says why the implication
nevertheless admits nothing, and what each widened test does and does not
guard against.
}
\begin{remark}[The bound-data test and the contractivity test on the band]\label{rem:bound-data-band}
\emph{Why item (iv) is widened at all.} Admissible receivers with
$\lmax(B)=1$ exactly are ordinary: at a receiver without a probe $B=A_\Sigma$
(Corollary~\ref{cor:unconditioned-family}), and the channel-only receiver of
Example~\ref{ex:bb84} has $A_\Sigma=\id_2$. There a computed $\lmax(B)$ one
rounding unit above $1$ makes $c$ negative on a receiver that is admissible,
and refusing it would refuse the simplest legitimate device.

\emph{The two items decide different objects.} A declaration can satisfy one
and fail the other: item~(v) admits a route-(ii) $\Gmat$ at $\varepsilon_1$,
item (iv) decides the derived $c$ at $\varepsilon_2$, and whenever
$\varepsilon_1>\varepsilon_2$, the ordinary declaration, the band
\eqref{eq:tolmapband} is non-empty, admitted by item~(v)'s test and refused by
item (iv)'s. On route (i) the band has no analogue, item~(i) and item (iv)
carrying the same $\varepsilon_2$; but item (iv) is a test there too and not a
formality, for the reason step~(2) below gives. Item (iv) wins on the band, in
four steps.

(1) \emph{In exact arithmetic item~(v)'s condition implies item (iv)'s, so
item (iv) is redundant, exactly.} $\Gmat^\dagger\Gmat\preceq\id$ gives
$A_\Sigma\preceq\Gmat_{\Inp}^\dagger\Gmat_{\Inp}\preceq\id$ by
Theorem~\ref{thm:port-grams}\,(i), hence $\lmax(A_\Sigma)\le1$ and $c\ge0$; and at
every setting $B^{(t)}\preceq\id$ by \eqref{eq:bbkhyp}, hence $c\ge0$ again.

(2) \emph{But item~(v) never establishes its exact condition.} It admits under
a widened test on a \emph{computed} number, and Theorem~\ref{thm:exact-and-evaluated}\,(ii)
bounds the distance from that number to the exact one by nothing at all. So
what step~(1) transports is an implication between exact quantities and
\emph{no implication whatever between the two computed quantities}, which are
separated by at least one matrix product
($A_k=\Gmat_{\Inp}^\dagger\Emd{\{k\}}\Gmat_{\Inp}$,
Definition~\ref{def:gram}) and, at a setting of a family a submodule
constructs from the Grams, by the further operations of that construction as
well. Reading
item~(v) as admitting anything under item (iv) is reading an exact implication
as a bound on an evaluation, which is the one substitution
Theorem~\ref{thm:exact-and-evaluated}\,(ii) exists to forbid. This step is also why item (iv)
is not redundant on route (i), where step~(1) holds just as well from
$\sum_kA_k\preceq\id$: redundancy of the \emph{exact} condition is not
admission of the \emph{computed} one. And $0\le c$ is tested where the rest of
\eqref{eq:bounddata} is not because the error term of
Remark~\ref{rem:cdirection}, on which step~(3) turns, is stated for a negative
computed $c$ and bounded through this test alone.

(3) \emph{The band is not empty of consequence, and the consequence is
quantitative and unsafe.} Remark~\ref{rem:cdirection} bounds the even-$n$
overshoot of the unclamped evaluation by
$(K-1)|c|^{\,n}\le(K-1)\varepsilon_2^{\,2}$, and that bound is stated in
$\varepsilon_2$ \emph{because item (iv) is what bounds $|c|$}. Admitting
\eqref{eq:tolmapband} would make $|c|$ as large as $\varepsilon_1$ and the
true bound $(K-1)\varepsilon_1^{\,2}$, larger by the factor
$(\varepsilon_1/\varepsilon_2)^2$, which at an ordinary pair of declared
numbers is many orders of magnitude. A stated error term of this contract
would then be false on values this contract admits, in the direction where
the evaluated $\Bnd(n)$ \emph{exceeds} the exact one, which is the direction
that reaches a key length. Remark~\ref{rem:tolerancedomain} says what a
declared domain buys, ``that every error term this contract states holds on
every value the contract admits'', and that is the property item~(v) winning
would destroy.

(4) \emph{Nothing is lost by item (iv) deciding.} The case the widening exists
for, a receiver with $\lmax(B)=1$ exactly whose computed $\lmax(B)$ sits a
rounding unit above $1$, is admitted by item (iv)'s own $-\varepsilon_2$; and item (iv) tests the
very quantity that enters \eqref{eq:beta}, where item~(v) tests one two
computations upstream. A declarer who wants the wider band declares a wider
$\varepsilon_2$, which is what a receiver tolerance is for and which carries
the stated error term with it, rather than leaving it behind.

\emph{Item (v) is one-sided as a test and two-sided in its consequences.}
$\Gmat^\dagger\Gmat\preceq\id$ is an upper bound on a spectrum, so the test
relaxes in one direction only: it admits
$\lmax(\Gmat^\dagger\Gmat)\in(1,1+\varepsilon_1]$, which the exact condition
excludes, and it refuses nothing that the exact condition, applied to the same
computed number, would admit. The \emph{errors} are two-sided in the sense of
Remark~\ref{rem:ranktwosided}, and deliberately so. Too tight refuses a
legitimate device: a receiver with a unitary $\Gmat$, such as the worked
instrument \eqref{eq:bb84G} of Example~\ref{ex:bb84}, has
$\lmax(\Gmat^\dagger\Gmat)=1$ \emph{exactly}, so a computed value one rounding unit above
$1$ would refuse it, which is the argument for widening item~(iv) verbatim.
Too loose admits a genuinely non-contractive $\Gmat$, which violates
Assumption~\ref{ass:passive}. Route (ii) carries no separate Gram
admissibility check, because Theorem~\ref{thm:port-grams}\,(i) \emph{derives}
$\sum_kA_k\preceq\Gmat_{\Inp}^\dagger\Gmat_{\Inp}\preceq\id$ from
contractivity, so the \emph{exact} condition item~(i) tests at route (i) is
redundant here rather than dropped.

\emph{What that derivation transports, and what it does not.} It is an
implication between \emph{exact} quantities. It therefore makes item~(i)'s
exact condition redundant on route (ii), and, through $B\preceq\id$ in
\eqref{eq:bbkhyp}, item~(iv)'s exact condition $c\ge0$ redundant too. It
transports \emph{no admission}: the computed $\lmax(A_\Sigma)$ and the computed
$c$ are separate evaluations, related to the exact ones by nothing this
contract asserts (Theorem~\ref{thm:exact-and-evaluated}\,(ii)), so an $\varepsilon_1$-admitted
$\Gmat$ does not thereby satisfy any test stated at $\varepsilon_2$. Item~(iv)
accordingly decides the bound data on route (ii) as on every other, and the
band \eqref{eq:tolmapband} that opens when $\varepsilon_1>\varepsilon_2$ is
refused. \emph{Each route therefore tests its own declared datum under its own
declared number}, $\varepsilon_1$ on the supplied $\Gmat$ at route (ii) and
$\varepsilon_2$ on the supplied Grams at route (i), because each is the
tolerance attached to \emph{that} datum's residual, and neither is thereby a
tolerance on the derived $c$. A declarer who wants one slack on
$\lmax(A_\Sigma)$ regardless of route declares
$\varepsilon_1=\varepsilon_2$, which also empties \eqref{eq:tolmapband}.
\end{remark}
\claimmeta{\metalabel{Uses}\mref{Def.}{def:gram}, \mref{Thm.}{thm:port-grams}, \mref{Def.}{def:bounddata}, \mref{Thm.}{thm:multiclick-bound}, \mref{Cor.}{cor:unconditioned-family}, \mref{Def.}{def:tolmap}, \mref{Rem.}{rem:tolerancedomain}, \mref{Rem.}{rem:cdirection}, \mref{Thm.}{thm:exact-and-evaluated}
\metasep newly derived
\metasep accepted}

% claim receiver/rem:tolerancedomain
\claimsummary{receiver/rem:tolerancedomain}{Every quantity a receiver tolerance bounds, decides or reduces lies in [0, 1], so at a tolerance of one or more its test is vacuous or degenerate. The declared domain [0, 1) makes every error term the contract states hold on every value it admits.}
\claimbookmark{receiver/rem:tolerancedomain}{Remark}{rem:tolerancedomain}{The domain of the receiver tolerances}
\begin{remark}[The domain of the receiver tolerances]
\label{rem:tolerancedomain}
Each of $\varepsilon_1,\dots,\varepsilon_6$, $\varepsilon_{\mathrm{gap}}$ and
$\varepsilon_{\mathrm{eig}}$ is declared in $[0,1)$ and validated at input
(Requirement~\ref{req:tolerance-domain}). A declared number needs a declared
domain for the same reason a declared parameter does, and here the domain is
not a convention: \emph{a \termref{term-receiver-tolerances}{receiver tolerance} is a
slack on a normalized quantity, never a scale.} Every quantity the eight bound,
decide or reduce already lies in $[0,1]$: eigenvalues of operators between $0$
and $\id$ (\eqref{eq:res2}, \eqref{eq:eps2rank},
Definition~\ref{def:tolmap}\,(i),\,(v)), the spectral-norm defects
\eqref{eq:res1}, \eqref{eq:res3}, \eqref{eq:res4} and \eqref{eq:rhodef} of
operators bounded by $\id$, entries and column defects of a stochastic matrix
(Definition~\ref{def:tolmap}\,(iii)), and $c=1-\lmax(B)\in[0,1]$
(Definition~\ref{def:tolmap}\,(iv)). At $\varepsilon\ge1$ the corresponding
test is not merely loose, it is \emph{vacuous or degenerate}, and namably so:
Definition~\ref{def:tolmap}\,(i) would admit every Hermitian family with
$\lmin\ge-1$; \eqref{eq:eps2rank} would return $\rk_{\varepsilon_2}(A_k)=0$ for
\emph{every} port Gram, so the rank check could never refuse (the
under-reporting direction of Remark~\ref{rem:ranktwosided} taken to its
limit); (iii) would call every declared matrix monotone; (iv) would admit every
computed $c$; and (v) would admit a $\Gmat$ of operator norm $\sqrt2$.
$\varepsilon_{\mathrm{gap}}$ is bounded for the same reason at one remove: it
bounds a slack in $\pmc\in[0,1]$ (Theorem~\ref{thm:emitted-bound}), so a value
$\ge1$ permits the emitted bound to sit below the attainable one by the whole
range of the quantity, which is no constraint. The two members the exact
window adds are bounded on the same argument, up to a computed
$\widehat\Bnd^{\ast}$ that rounding puts outside $[0,1]$, which is not clamped
into that interval and is not claimed to lie in it
(Remark~\ref{rem:exacteig}). $\varepsilon_6$ bounds the positive part of
$\Bnd(n)-\widehat\Bnd^{\ast}(n)$, where $\Bnd(n)\le1$
(Theorem~\ref{thm:multiclick-bound}) and the exact value that
$\widehat\Bnd^{\ast}(n)$ evaluates lies in $[0,1]$
(Lemma~\ref{lem:exactbasic} with Definition~\ref{def:setting-family}); so at
$\varepsilon_6\ge1$ the test $r_6\le\varepsilon_6$ can refuse only an
evaluation that puts $\widehat\Bnd^{\ast}(n)$ below the evaluated $\Bnd(n)$
less $1$, which is at most $0$ up to rounding, outside the interval where the
exact value lies. $\varepsilon_{\mathrm{eig}}$ is subtracted from a quantity
whose exact value $\Bnd^{\ast}$ lies in $[0,1]$ at every receiver
(Lemma~\ref{lem:exactbasic} with \eqref{eq:mcpos}, which
Definition~\ref{def:setting-family} gives). So at
$\varepsilon_{\mathrm{eig}}\ge1$ the second entry of \eqref{eq:elldef} is
positive only where a computed $\widehat\Bnd^{\ast}$ exceeds $1$; elsewhere it
is never positive and the exact window is never the binding term: the
declaration pays the full cost of the computation and buys nothing, which is
degenerate in the same namable sense as the rest of this list.

Two arguments use the upper bound. Remark~\ref{rem:cdirection}'s even-$n$
estimate uses $\varepsilon_2^{\,n}\le\varepsilon_2^{\,2}$ for every even
$n\ge2$, which is exactly $\varepsilon_2\le1$; and the justification of
\eqref{eq:eps2rank} in Remark~\ref{rem:ranktwosided}, that the spectrum it
thresholds is normalized to $[0,1]$, is an argument about a threshold
\emph{inside} that interval.

The domain is a floor, not sizing advice. Nothing here says a tolerance near
$1$ is sensible; Remark~\ref{rem:ranktwosided} records that the choice of
$\varepsilon_2$ decides which declarations are accepted, and
Proposition~\ref{prop:dark-slack-bound} gives the inequality a declarer sizes $\varepsilon_2$
and $\varepsilon_4$ against. What the declared domain buys is that every error
term this contract states holds on every value the contract admits.
\end{remark}
\claimmeta{\metalabel{Uses}\mref{Thm.}{thm:multiclick-bound}, \mref{Lem.}{lem:exactbasic}, \mref{Def.}{def:pmcex}, \mref{Rem.}{rem:exacteig}, \mref{Thm.}{thm:emitted-bound}, \mref{Def.}{def:residuals}, \mref{Def.}{def:tolmap}
\metasep \metalabel{Used by}\mref{Rem.}{rem:bound-data-band}, \mref{Rem.}{rem:cdirection}, \mref{Req.}{req:tolerance-domain}
\metasep newly derived
\metasep accepted}

% claim receiver/rem:ranktwosided
\claimsummary{receiver/rem:ranktwosided}{Under any positive threshold a counted rank can exceed or fall short of the exact one, so the rank check is not refuse-only. The contract therefore fixes the threshold at the declared \ensuremath{\varepsilon}\ensuremath{_{2}}, absolute rather than relative, instead of an undeclared library default.}
\claimbookmark{receiver/rem:ranktwosided}{Remark}{rem:ranktwosided}{Two-sided errors of a rank cutoff, and the choice of threshold}
\begin{remark}[Two-sided errors of a rank cutoff, and the choice of threshold]\label{rem:ranktwosided}
Under any positive threshold $\tau$, the count \eqref{eq:eps2rank} can differ
from the exact rank in \emph{both} directions.
\begin{itemize}[leftmargin=1.6em]
\item \emph{Over-reported.} An exactly singular declared Gram whose zero
  eigenvalue evaluates at $+2\tau$ (rounding noise in the declared entries,
  or in the eigenvalue routine) is counted as one rank higher, and a
  declaration consistent with its port map is \emph{refused}.
\item \emph{Under-reported.} A Gram whose exact rank is $2$ but whose second
  eigenvalue is genuinely small, below $\tau$, is counted as rank $1$, and a
  declaration \emph{inconsistent} with its port map is \emph{admitted} on a
  single-mode port.
\end{itemize}
So the check is not refuse-only, and the choice of $\tau$ is a choice about
which declarations are accepted, not a numerical detail. That is why this
contract fixes $\tau=\varepsilon_2$, a number the declarer supplies and can
reason about, rather than inheriting whatever cutoff a linear-algebra library
applies by default; the usual library default is
$\lmax\cdot\max(D_1,D_2)\cdot\epsilon_{\mathrm{mach}}$, which is relative,
undeclared, and invisible in the artifact. Under \eqref{eq:eps2rank} both
directions remain possible (no threshold can remove them), but both are now
consequences of a declared number, and both are recorded: the rank the check
used is a function of $\{A_k\}$ and $\varepsilon_2$ alone.

The threshold of \eqref{eq:eps2rank} is \emph{absolute}, because a port Gram is
already normalized by construction: $0\preceq A_k\preceq A_\Sigma\preceq\id$
puts every eigenvalue in $[0,1]$, and $\vec\alpha^\dagger A_k\vec\alpha$ is a
mean photon number per unit input intensity, so there is no scale to normalize
against. A threshold taken relative to $\lmax(A_k)$ would make the decision at
a nearly dark port (all eigenvalues $\sim10^{-6}$) a million times tighter than
at a lossless one, which is a property of the port's throughput and not of the
declarer's numerics. And the threshold is $\varepsilon_2$, because
\eqref{eq:eps2rank} and Definition~\ref{def:tolmap}\,(i) are the same test on
the same spectrum from the two sides: item (i) admits eigenvalues down to
$-\varepsilon_2$ as ``numerically zero'', and \eqref{eq:eps2rank} counts as
nonzero exactly those above $+\varepsilon_2$, so a declared Gram admitted by
item (i) has its numerically-zero eigenvalues in
$[-\varepsilon_2,\varepsilon_2]$ and the two thresholds cannot disagree about
which they are.
\end{remark}
\claimmeta{\metalabel{Uses}\mref{Def.}{def:tolmap}
\metasep \metalabel{Used by}\mref{Rem.}{rem:d2slack}
\metasep newly derived
\metasep accepted}

% claim receiver/rem:cdirection
\claimsummary{receiver/rem:cdirection}{An admitted negative c lowers the analytic bound at odd photon numbers and can raise it at even ones, by at most (K\ensuremath{-}1)\ensuremath{\varepsilon}\ensuremath{_{2}}\ensuremath{^{2}}, second order on the declared domain. Neither comparison is an enclosure of the exact bound.}
\claimbookmark{receiver/rem:cdirection}{Remark}{rem:cdirection}{The direction of the unclamped c at odd and even photon numbers}
\begin{remark}[The direction of the unclamped $c$ at odd and even photon numbers]
\label{rem:cdirection}
Definition~\ref{def:tolmap}\,(iv) admits $c\in[-\varepsilon_2,0)$ and feeds
that value to \eqref{eq:beta} unchanged. Since
$\Bnd(n)=1+(K-1)\gamma c^{\,n}-\sum_k\beta_kd_k^{\,n}$ is increasing in
$c^{\,n}$, and since the exact $c$ satisfies $c_{\mathrm{ex}}\ge0$, the effect
splits by the parity of $n$, and both sides are stated.
\begin{itemize}[leftmargin=1.6em]
\item \emph{Odd $n$.} $c^{\,n}<0\le c_{\mathrm{ex}}^{\,n}$, so the evaluated
  $\Bnd(n)$ is \emph{below} the exact one. Conservative, and by a first-order
  amount: measured against the value at $c=0$ the deficit is
  $(K-1)\gamma|c|^{\,n}\le(K-1)\varepsilon_2$ at $n=1$, which is a lower bound
  on the deficit against $\Bnd(n)$ at $c_{\mathrm{ex}}$, that one being
  $(K-1)\gamma\big(c_{\mathrm{ex}}^{\,n}+|c|^{\,n}\big)$.
\item \emph{Even $n\ge2$.} $c^{\,n}=|c|^{\,n}>0$, so the evaluated $\Bnd(n)$
  can \emph{exceed} the exact one, by at most
  $(K-1)\gamma\big(|c|^{\,n}-c_{\mathrm{ex}}^{\,n}\big)
  \le(K-1)|c|^{\,n}\le(K-1)\varepsilon_2^{\,2}$, using $\gamma\le1$,
  $|c|\le\varepsilon_2$ and $\varepsilon_2\le1$. The bound
  $|c|\le\varepsilon_2$ is Definition~\ref{def:tolmap}\,(iv) and no other rule,
  which is why that item decides the admission of $c$ on \emph{every} route,
  item~(v)'s route-(ii) admission at $\varepsilon_1$ not being a bound on $|c|$
  at all. The last step is where the declared upper bound of
  Remark~\ref{rem:tolerancedomain} is used, since
  $\varepsilon_2^{\,n}\le\varepsilon_2^{\,2}$ for even $n\ge2$ holds exactly on
  $\varepsilon_2\le1$ and fails above it. The overshoot is second order in
  $\varepsilon_2$ on the declared domain, and this is the one direction in
  which clamping $c$ to $0$ would have been the safer choice.
\end{itemize}
The contract fixes the unclamped rule because a clamped and an unclamped
evaluation are both defensible and they differ, so leaving the choice to an
implementation would leave the emitted number undetermined by the contract.
Neither statement above is an enclosure: they compare two evaluations of
\eqref{eq:beta}, and Theorem~\ref{thm:exact-and-evaluated}\,(ii) still governs the distance
from either to the exact quantity. The case arises in the floating-point
evaluation only: the rigorous evaluation takes its lower end of $c$ in
$[0,1]$ (Remark~\ref{rem:ia}, Definition~\ref{def:enclosed-weight}), so no
negative $c$ enters a weight it forms.
\end{remark}
\claimmeta{\metalabel{Uses}\mref{Def.}{def:bounddata}, \mref{Def.}{def:tolmap}, \mref{Rem.}{rem:tolerancedomain}
\metasep \metalabel{Used by}\mref{Rem.}{rem:bound-data-band}, \mref{Rem.}{rem:ia}
\metasep newly derived
\metasep accepted}

% claim receiver/def:evaluation-model
\claimsummary{receiver/def:evaluation-model}{The claims are exact statements about exactly declared data, evaluated in binary64 with diagnostic residuals, while the blocks are enclosed and, where the port Grams are enclosed, the weights evaluated rigorously, in exact rational arithmetic. The restriction asserts nothing a run could violate.}
\claimbookmark{receiver/def:evaluation-model}{Definition}{def:evaluation-model}{The evaluation model: floating point, enclosures and rigour}
\begin{definition}[The evaluation model: floating point, enclosures and rigour]
\label{def:evaluation-model}
The claims of this contract are exact statements about exactly declared data
(\xref{foundations/def:binary64}\,(iii)). This contract \emph{evaluates} them
in binary64 floating-point arithmetic (\xref{foundations/def:binary64}),
reports those of the six \termref{term-residuals}{residuals} of Definition~\ref{def:residuals}
that exist on the input route taken and for the evaluation requested, and
asserts no rigorous enclosure of any quantity it computes in floating point
(Assumption~\ref{ass:eig} is the one exception): this is its
\emph{floating-point evaluation}. For a declaration whose
model is an admissible receiver exactly it moreover computes, in exact
rational arithmetic, \termref{term-complex-enclosure}{complex enclosures}
(Definition~\ref{def:complex-enclosure}) of the photon-number blocks of the
emitted family: those of Theorem~\ref{thm:enclosed-blocks}, and in the
submodules those of the setting families a submodule constructs, each
enclosed by that submodule, and of their composition with the dark counts
(Corollary~\ref{cor:exact-dark-composed}). They are exact where
the declaration is on the \termref{term-exact-route}{exact route}
(Definition~\ref{def:exact-route}) and the entry is known exactly
(Theorem~\ref{thm:exact-blocks}). From the same enclosures of the port Grams
it computes, in exact rational arithmetic, at every setting at which the rules
of Lemma~\ref{lem:enclosure-arithmetic} give them, complex enclosures of the
setting's operators $B^{(t)}$, $B^{(t)}_k$ and prefactors $\gamma^{(t)}$,
$\beta^{(t)}_k$ (Definition~\ref{def:setting-family}) and, at a setting
carrying an exact-window depth, of its blocks $\blk{M}{\eqn}_{\mc|t}$ in the
window; from these, an upper bound $u$ on $\lmax(B^{(t)})$ and lower bounds
$\ell_k$ on $\lmin(B^{(t)}_k)$ and $\lambda_n$ on the least eigenvalue of
each window block, by the least-eigenvalue bounds over an enclosure
(\xref{foundations/lem:least-eigenvalue-bound},
\xref{foundations/cor:least-eigenvalue-over-enclosure}), and a lower bound $g$
on $\gamma^{(t)}$ and upper bounds $b_k$ on $\beta^{(t)}_k$, each the real
part of a centre less or plus its radius; from these, the reals
$g,b_k,1-u,1-\ell_k$ each intersected with $[0,1]$ and
$\max\{0,\lambda_n\}$; and from those, by the formulas of
Definition~\ref{def:enclosed-weight}\,(iii), a weight: this is its
\emph{rigorous evaluation} of the weights. A weight emitted at a setting is
the weight so formed where the rigorous evaluation is formed there, and
otherwise the floating-point evaluation of Definition~\ref{def:pmc}
or~\ref{def:pmcex}.
These enclosures and bounds are not evaluations, and a decision taken on
them is exact; their interval arithmetic is proved in this contract
(Lemma~\ref{lem:enclosure-arithmetic}) and their eigenvalue bounds in the
foundations contract, and no interval-arithmetic assumption is imported. This
\termdef{term-evaluation-model}{evaluation model}
is a \emph{restriction} this contract places on itself and asserts nothing a
run could violate.
\end{definition}
\claimmeta{\metalabel{Uses}\mref{Def.}{def:setting-family}, \mref{Def.}{def:enclosed-weight}, \mref{Def.}{def:complex-enclosure}, \mref{Def.}{def:exact-route}
\metasep \metalabel{Used by}\mref{Thm.}{thm:exact-and-evaluated}, \mref{Ass.}{ass:eig}, \mref{Req.}{req:rigorous-weight}, \mref{Req.}{req:exact-route-decision}
\metasep newly derived
\metasep accepted}

% claim receiver/thm:exact-and-evaluated
\claimsummary{receiver/thm:exact-and-evaluated}{The receiver's identities and inequalities are exact unless a claim names a receiver tolerance or \hyperlink{claim-receiver/ass:eig}{A-EIG}; its floating-point numbers carry no error bound except through \hyperlink{claim-receiver/ass:eig}{A-EIG}, while rational enclosures contain the model's blocks and rigorously evaluated weights are bounds outright.}
\claimbookmark{receiver/thm:exact-and-evaluated}{Theorem}{thm:exact-and-evaluated}{Exact statements, floating-point and rigorous evaluation}
\claimleadin{%
Remark~\ref{rem:exactness-scope} says what item~(i) does not assert, and
Remark~\ref{rem:exacteig} says why the window residual $r_6$ of item~(v) is
blind to the direction Assumption~\ref{ass:eig} constrains.
Requirement~\ref{req:rigorous-weight} is where the engine forms the rigorous
evaluation of item~(vii), and Requirement~\ref{req:evaluation-label} what its
label means for the tier of a run. Items~(vi) and~(vii) speak of the core's
own family; the submodules state the counterparts for theirs, such as the
enclosures of the dark-composed blocks
(Corollary~\ref{cor:exact-dark-composed}), and a submodule that constructs a
setting family states the enclosures of its blocks and the rigorously
evaluated weight at its settings.
}
\begin{theorem}[Exact statements, floating-point and rigorous evaluation]\label{thm:exact-and-evaluated}
For the evaluation model of Definition~\ref{def:evaluation-model}:
\begin{enumerate}[label=(\roman*),leftmargin=3.0em]
\item \emph{The asserted mathematics is exact, except where a claim says so in
  its own statement.} \emph{Every asserted identity and inequality} of this
  module and its submodules is between exactly specified quantities and
  contains no error term, no tolerance and no numerical hypothesis
  (Theorem~\ref{thm:multiclick-bound}, Lemma~\ref{lem:collapse},
  Lemma~\ref{lem:windowtail}, Lemma~\ref{lem:exactbasic} and
  Lemma~\ref{lem:exactwindow} included, the last two being stated on the exact
  $\Bnd^{\ast}$ and carrying no tolerance), \emph{except where a claim's own
  statement names a \termref{term-receiver-tolerances}{receiver tolerance} or
  Assumption~\ref{ass:eig}}.
\item \emph{The floating-point evaluation is not.} The emitted $\{A_k\}$,
  block matrices and those bound values that the floating-point evaluation
  produces are floating-point evaluations of those quantities; the enclosures
  of item (vi) and the rigorous evaluation of item (vii) are not. This
  contract asserts no bound on the difference, and the residuals of
  Definition~\ref{def:residuals} are \emph{diagnostics}: passing them is
  necessary for the artifact to be self-consistent and is not sufficient for
  any enclosure statement.
\item \emph{The evaluation label.} An artifact carries the rigorous evaluation
  label exactly where every \termref{term-per-setting-weight}{per-setting weight} it emits is formed by the
  rigorous evaluation of Definition~\ref{def:evaluation-model}, as the
  \termref{term-enclosed-weight}{enclosed weight} of Definition~\ref{def:enclosed-weight} or a
  non-negative real below it; every other artifact, on every input route, the
  route-(iii) compiler included, carries the floating-point evaluation label.
  The rigorous label concerns the weights alone: the blocks are evaluated in
  floating point under (ii) whatever the label. No path in this module
  produces a label stronger than these two.
\item \emph{Where route (i) stands.} If the Grams are \emph{declared}, as
  binary64 values, which are dyadic rationals
  (\xref{foundations/def:binary64}\,(iii)), the compilation error of (ii) does
  not exist: there is no compilation. The eigenvalue quantities of
  Theorem~\ref{thm:multiclick-bound} are still evaluated in floating point by
  the floating-point evaluation, so (ii) and (iii) apply to a weight it forms;
  what changes is that the rigorous evaluation of (vii) then encloses no
  matrix product, the enclosures it starts from being the declared Grams
  themselves. This is also why $r_1$ and $r_3$ do not exist on this route: they
  are written on a $\Gmat$ that is not supplied
  (Definition~\ref{def:residuals}). Nothing further is claimed.
\item \emph{The one place a floating-point number is asserted to bound an
  exact one, and it is opt-in.} Where a setting $t$ carries an exact-window
  depth and its emitted weight $\pmc^{(t)}$ is that of
  Definition~\ref{def:pmcex}, the weight depends on the
  computed $\widehat\Bnd^{\ast}$ through \eqref{eq:elldef}, and through the
  minimum so may the uniform $\pmc$; the inequality
  $\widehat\Bnd^{\ast(t)}(n)-\varepsilon_{\mathrm{eig}}\le\Bnd^{\ast(t)}(n)$
  that makes each such emitted value a lower bound is Assumption~\ref{ass:eig},
  not a theorem of this contract: (ii) governs $\widehat\Bnd^{\ast}$ exactly as
  it governs every other computed quantity, and $r_6$ is a diagnostic in
  exactly the sense (ii) fixes and is blind to the direction the assumption
  constrains. Where a setting's emitted weight is an enclosed weight, no
  computed $\widehat\Bnd^{\ast}$ and no $\varepsilon_{\mathrm{eig}}$ enters
  it and Assumption~\ref{ass:eig} is not invoked there, whatever depth is
  declared. Where no depth is declared, no
  $\Bnd^{\ast}$ is computed, Assumption~\ref{ass:eig} is not invoked, and the
  emitted weights are the values \eqref{eq:pmcdef} defines, expression for
  expression, evaluated in floating point or on \termref{term-enclosed-bound-data}{enclosed bound data}. Item
  (iii) covers all of these: an artifact with a weight resting on
  Assumption~\ref{ass:eig} carries the floating-point label, so no artifact
  resting on Assumption~\ref{ass:eig} can claim more.
\item \emph{The model's blocks are enclosed, and the evaluation approximates
  them.} For a declaration whose model is an admissible receiver exactly, the
  photon-number blocks of the click-pattern and
  \termref{term-emitted-povm}{emitted POVMs} of that model are quantities the
  identities and inequalities of (i) are about. Their
  \termref{term-complex-enclosure}{complex enclosures}
  (Definition~\ref{def:complex-enclosure}) of
  Theorem~\ref{thm:enclosed-blocks}, where they are formed, are computed in
  exact rational arithmetic and contain them: every entry of a
  block lies within its radius of its centre, an entry whose radius is zero
  is its centre, and on the \termref{term-exact-route}{exact route} the entries that
  Theorem~\ref{thm:exact-blocks} names have radius zero when the port Grams
  are enclosed exactly. A positive semidefiniteness or a positive
  definiteness established through
  \xref{foundations/lem:enclosure-psd}\,(i) on such an enclosure with a
  Hermitian centre holds of the block itself. The floating-point block
  matrices are evaluations of the same blocks and fall under (ii): this
  contract asserts no bound on their difference from the blocks, so no
  decision about the blocks is taken on them.
\item \emph{The rigorously evaluated weights are bounds.} Let $t$ be a
  setting of a \termref{term-setting-family}{setting family} (Definition~\ref{def:setting-family}) at
  which the rigorous evaluation is formed from complex enclosures of
  $B^{(t)}$, of every $B^{(t)}_k$, of $\gamma^{(t)}$, of every
  $\beta^{(t)}_k$ and, at a setting carrying a depth, of the window blocks
  $\blk{M}{\eqn}_{\mc|t}$; for a receiver without a probe, whose setting
  family is that of Corollary~\ref{cor:unconditioned-family}, the enclosures
  it forms from the enclosures of the port Grams are such. Then the four
  clipped reals it forms
  are \termref{term-enclosed-bound-data}{enclosed bound data} of the setting and, at a setting carrying a
  depth, the window values it forms are \termref{term-rigorous-window-bounds}{rigorous window bounds}; the
  weight it forms from them is their enclosed weight, a rational computed
  exactly; and it and every real between $0$
  and it bound the \termref{term-multiclick-element}{multiclick element} of the setting from below on every
  photon-number block above the cutoff and lie in $[0,1]$. Its validity rests
  on no number computed in floating point: such a number enters it, if at
  all, only as the exact rational it is, of which nothing is assumed, as the
  approximate eigenbasis of \xref{foundations/lem:least-eigenvalue-bound}
  does. So (ii) does not govern it, and it rests on no assumption about
  arithmetic.
\end{enumerate}
\end{theorem}
\claimmeta{\metalabel{Uses}\mref{Def.}{def:bounddata}, \mref{Thm.}{thm:multiclick-bound}, \mref{Lem.}{lem:collapse}, \mref{Lem.}{lem:windowtail}, \mref{Lem.}{lem:exactbasic}, \mref{Lem.}{lem:exactwindow}, \mref{Def.}{def:setting-family}, \mref{Cor.}{cor:unconditioned-family}, \mref{Def.}{def:pmc}, \mref{Def.}{def:pmcex}, \mref{Def.}{def:enclosed-weight}, \mref{Thm.}{thm:enclosed-weight}, \mref{Thm.}{thm:emitted-bound}, \mref{Def.}{def:complex-enclosure}, \mref{Lem.}{lem:enclosure-arithmetic}, \mref{Thm.}{thm:enclosed-blocks}, \mref{Def.}{def:exact-route}, \mref{Thm.}{thm:exact-blocks}, \mref{Def.}{def:residuals}, \mref{Def.}{def:tolmap}, \mref{Def.}{def:evaluation-model}, \extref{../foundations/contract.pdf}{foundations/cor:least-eigenvalue-over-enclosure}, \extref{../foundations/contract.pdf}{foundations/def:binary64}, \extref{../foundations/contract.pdf}{foundations/def:exact-reference}, \extref{../foundations/contract.pdf}{foundations/lem:enclosure-psd}, \extref{../foundations/contract.pdf}{foundations/lem:least-eigenvalue-bound}
\metasep \metalabel{Assumes}\aref{A-EIG}{ass:eig}
\metasep \metalabel{Used by}\mref{Rem.}{rem:bound-data-band}, \mref{Req.}{req:evaluation-label}
\metasep newly derived
\metasep accepted
\metasep \metalabel{Oracles}\coderef{https://github.com/fraunhofersg/certiqs-key/blob/v0.1.0/tests/oracles/registry_receiver_povm.py}{registry_receiver_povm.py}}

\begin{proof}
(i) By inspection of the statements: each is a Löwner-order or scalar
inequality between quantities defined from $\{A_k\}$, $\Gmat$, $\Tset$,
$\{\rho^{(t)}\}$, $N$, $\{m_t\}$ and $\{m^{\ast}_t\}$, with no epsilon
anywhere. A claim falling under the exception is recognised from its own
statement: \eqref{eq:elldef} introduces $\varepsilon_{\mathrm{eig}}$ into an
emitted value; Theorem~\ref{thm:emitted-bound} and the two corollaries that
instantiate it name Assumption~\ref{ass:eig}; and the submodules' tolerance
claims carry their declared slacks in the inequalities they assert. The
receiver tolerances occurring elsewhere in the claims this item quantifies over
occur in cross-references to Definition~\ref{def:tolmap}, in remarks
quantifying the effect of a rule stated here, and in
Example~\ref{ex:bb84exact}'s evaluation of \eqref{eq:elldef}, none of them an
asserted identity or inequality of the kind this item quantifies over.
Theorem~\ref{thm:enclosed-weight} and Definition~\ref{def:enclosed-weight}
state inequalities between reals and the exact quantities of the setting and
name neither.

(ii) is the content of Definition~\ref{def:evaluation-model}. That the residuals are not
sufficient is immediate: $r_1=\dots=r_6=0$ is a statement about the
\emph{computed} matrices, which are consistent with each other but are not
shown to be close to the exact ones; a compilation that lost precision
uniformly would satisfy every residual.

(iii) is a stipulation of the emitted behaviour, together with the refusal
convention of the module's interface
(\href{../../../spec/modules/receiver/interface.yaml}{\nolinkurl{spec/modules/receiver/interface.yaml}}); it proves nothing, and
Requirement~\ref{req:evaluation-label} is its engine form. It is consistent
because the rigorous label is given only where every emitted weight is a bound
by (vii) at every setting whose enclosures meet its hypothesis, which (vii)
shows for the unconditioned family and each submodule shows for its own, and
that is the only thing the label says; the floating-point label claims
nothing; and it is sound because a label withheld is never unsafe.

(iv) The compiled matrix is absent by hypothesis, so no product is formed; the
remaining floating-point steps of the floating-point evaluation are the
eigenvalue extremes in \eqref{eq:cd} and the arithmetic of \eqref{eq:beta}. On
this route the enclosures of the port Grams that the rigorous evaluation reads
are the declared Grams, exactly (Definition~\ref{def:exact-route}).

(v) The first sentence is Definition~\ref{def:pmcex} together with
Assumption~\ref{ass:eig}, which is the hypothesis
Theorem~\ref{thm:emitted-bound}'s proof discharges \eqref{eq:ellhyp} with on that
route. The last conjunct of the first sentence, that $r_6$ is blind to the direction
the assumption constrains, is read off \eqref{eq:res6}.
Assumption~\ref{ass:eig} bounds each computed $\widehat\Bnd^{\ast(t)}(n)$
from above at the settings carrying a depth whose emitted weight is that of
Definition~\ref{def:pmcex} in floating point, which are among those over
which the maximum in \eqref{eq:res6} runs, and at the indices
$N<n\le m^{\ast}_t$; at the other settings of that maximum it constrains
nothing, and the argument below is term by term. The
direction it constrains is a computed value above
$\Bnd^{\ast(t)}(n)+\varepsilon_{\mathrm{eig}}$. At such a setting the
hypotheses of Theorem~\ref{thm:multiclick-bound} hold
(Definition~\ref{def:setting-family}\,(b)), so \eqref{eq:exactbasic} of
Lemma~\ref{lem:exactbasic} gives $\Bnd^{(t)}(n)\le\Bnd^{\ast(t)}(n)$, and
$\varepsilon_{\mathrm{eig}}\ge0$ (Definition~\ref{def:tolmap}). A term
$\big(\Bnd^{(t)}(n)-\widehat\Bnd^{\ast(t)}(n)\big)_{+}$ is therefore positive
only where $\widehat\Bnd^{\ast(t)}(n)<\Bnd^{(t)}(n)\le\Bnd^{\ast(t)}(n)$, an
undershoot. A computed value with
$\widehat\Bnd^{\ast(t)}(n)>\Bnd^{\ast(t)}(n)+\varepsilon_{\mathrm{eig}}
\ge\Bnd^{(t)}(n)$ makes its term $0$, which is also the term the exact value
$\Bnd^{\ast(t)}(n)$ gives in its place. Each term depends on its own computed
value alone, so $r_6$ takes the same value whether the computed values that
violate Assumption~\ref{ass:eig} stand or are replaced by the exact ones: no
value of $r_6$, and so no comparison with $\varepsilon_6$, tells the two
apart. This is a property of $r_6$ as \eqref{eq:res6} defines it on the
computed values; the number reported is its floating-point evaluation, which
(ii) governs as it governs every residual.
The second sentence is Definition~\ref{def:enclosed-weight}\,(iii): the window
entry \eqref{eq:enclosed-ell} is formed from $\check\Bnd(n)$ and the rigorous
window bound $w(n)$, and neither is a computed eigenvalue or contains
$\varepsilon_{\mathrm{eig}}$; Theorem~\ref{thm:enclosed-weight}\,(iii), which
assumes nothing, is then what makes the weight a bound.
The third sentence is a property of the definitions rather than an argument:
$\Bnd^{\ast}$ occurs in \eqref{eq:pmcdef} nowhere, $r_6$ exists only where a
depth is declared, and $\varepsilon_6$ and $\varepsilon_{\mathrm{eig}}$ are
required only there (Definition~\ref{def:residuals},
Requirements~\ref{req:residual-emission}
and~\ref{req:exact-window-declaration}; the inputs table of
\href{../../../spec/modules/receiver/interface.yaml}{\nolinkurl{spec/modules/receiver/interface.yaml}}); and an enclosed weight at a
setting without a depth is \eqref{eq:pmcdef} read on enclosed bound data
(Definition~\ref{def:enclosed-weight}\,(iii)). The fourth is (iii), since a
weight of Definition~\ref{def:pmcex} is not formed by the rigorous evaluation.

(vi) Theorem~\ref{thm:enclosed-blocks}\,(ii) and~(iii) state that the
enclosures it forms are complex enclosures of the blocks of the click-pattern
and emitted POVMs; Definition~\ref{def:evaluation-model} computes them in exact
rational arithmetic, every rule of Lemma~\ref{lem:enclosure-arithmetic} being
a finite computation on Gaussian rationals and rationals. By
Definition~\ref{def:complex-enclosure}, an entry $x$ with centre $x_c$ and
radius $r$ satisfies $\lvert x-x_c\rvert\le r$, so $x=x_c$ when $r=0$; the
entries named in Theorem~\ref{thm:exact-blocks} have radius zero under its
hypothesis by its item (ii). A block is Hermitian, and an enclosure with a
Hermitian centre supplies the hypotheses of
\xref{foundations/lem:enclosure-psd}\,(i) with $X$ the block, so the
conclusion of that lemma is a property of the block itself. The last two
sentences are (ii) applied to the floating-point block matrices.

(vii) \emph{The enclosures.} For a receiver without a probe the setting family
is that of Corollary~\ref{cor:unconditioned-family}: $B=A_\Sigma$ and
$B_k=A_{[K]\setminus\{k\}}$ are sums of port Grams, enclosed from the
enclosures of the port Grams by Lemma~\ref{lem:enclosure-arithmetic}\,(i);
$\gamma=\beta_k=1$ exactly; and each window block
$\blk{M}{\eqn}_{\mc}=\id+(K-1)\Pvac{[K]}^{(n)}-\sum_k\Pvac{[K]\setminus\{k\}}^{(n)}$
(\eqref{eq:family-mc}) is enclosed by Theorem~\ref{thm:enclosed-blocks}\,(ii),
taken with the depth as its cutoff, and the same rules. Each is a complex
enclosure of the quantity it names, which is the hypothesis of (vii) for such
a receiver. At every setting the hypothesis covers, an enclosure of a
Hermitian matrix made Hermitian-centred by
Lemma~\ref{lem:enclosure-arithmetic}\,(v) is an \exttermref{../foundations/contract.pdf}{term-entrywise-enclosure}{entrywise enclosure}.

\emph{The four bounds.} For an entrywise enclosure $(C,r)$ of a Hermitian $X$
and a rational $\rho\ge0$ with $\rho^{2}\ge\sum_{i,l}r_{il}^{2}$,
\xref{foundations/cor:least-eigenvalue-over-enclosure}\,(i) and~(ii) give
$\lmin(X)\ge\lmin(C)-\rho$ and $\lmax(X)\le\lmax(C)+\rho$, and
\xref{foundations/lem:least-eigenvalue-bound}\,(ii) and~(iv) bound
$\lmin(C)$ from below and $\lmax(C)$ from above at every real $V$ that
admits a rational $\theta$ with $0\le\theta<1$ and
$\lVert V^{\transpose}V-\id\rVert_{F}\le\theta$, the hypothesis of that lemma,
the identity with $\theta=0$ among them, however $V$ was obtained, and as
rationals by its item (iii). For a complex
enclosure $(x_c,r)$ of a real $x$,
$\lvert x-\mathrm{Re}\,x_c\rvert=\lvert\mathrm{Re}(x-x_c)\rvert\le\lvert
x-x_c\rvert\le r$, so $\mathrm{Re}\,x_c-r\le x\le\mathrm{Re}\,x_c+r$. Hence the
reals of Definition~\ref{def:evaluation-model} satisfy $u\ge\lmax(B)$,
$\ell_k\le\lmin(B_k)$, $g\le\gamma$, $b_k\ge\beta_k$ and
$\lambda_n\le\lmin(\blk{M}{\eqn}_{\mc|t})$.

\emph{The weight.} Theorem~\ref{thm:enclosed-weight}\,(i) makes the clipped
reals enclosed bound data and the window values rigorous window bounds, so the
weight formed from them by Definition~\ref{def:enclosed-weight}\,(iii) is their
enclosed weight; it is a finite expression in rationals built from sums,
products, minima and maxima by \eqref{eq:pmcdef} and \eqref{eq:pmcex}, so a
rational computed exactly. Theorem~\ref{thm:enclosed-weight}\,(iii) gives the
bound and the range, for the enclosed weight and for every real between $0$
and it, under the hypotheses of Definition~\ref{def:setting-family} alone.
A floating-point number enters, if at all, only as the basis $V$ read as the
rational matrix it is, which decides tightness and never validity, so the
difference (ii) speaks of does not arise.
\end{proof}

% claim receiver/rem:exactness-scope
\claimsummary{receiver/rem:exactness-scope}{The eigenvalue-slack assumption is the one relation between a floating-point and an exact number the receiver asserts; enclosures and rigorous weights are proved. A claim outside the exception of \hyperlink{claim-receiver/thm:exact-and-evaluated}{Theorem~\ref*{thm:exact-and-evaluated}} may still mention a receiver tolerance, in a cross-reference, a remark or an example.}
\claimbookmark{receiver/rem:exactness-scope}{Remark}{rem:exactness-scope}{Scope of the exactness rule}
\begin{remark}[Scope of the exactness rule]\label{rem:exactness-scope}
The one relation between a floating-point number and an exact one that this
contract does assert is Assumption~\ref{ass:eig}; it is a separate assumption
because it is invoked only on the opt-in route of Definition~\ref{def:pmcex},
evaluated in floating point.
The enclosures of the photon-number blocks
(Theorem~\ref{thm:enclosed-blocks}) relate rationals computed in exact
arithmetic to the model, and so do the rigorously evaluated weights
(Theorem~\ref{thm:enclosed-weight}); they are proved, not assumed.
An artifact carries the rigorous evaluation label exactly where every weight
it emits was evaluated rigorously, and the floating-point label otherwise; no
path in this module produces another label
(Theorem~\ref{thm:exact-and-evaluated}\,(iii),
Requirement~\ref{req:evaluation-label}).

Item~(i) of Theorem~\ref{thm:exact-and-evaluated} does not say that a claim
outside its exception never \emph{mentions} a receiver tolerance: such claims
do, in cross-references to Definition~\ref{def:tolmap} recording which
declared number decides an exact condition in the evaluation model of
Definition~\ref{def:evaluation-model} (Definition~\ref{def:bounddata},
Definition~\ref{def:darkadm}, Remark~\ref{rem:corrdark}), and in remarks and
an example that quantify the effect of a rule stated in that theorem
(Remark~\ref{rem:noslack}, Remark~\ref{rem:exacteig},
Example~\ref{ex:bb84exact}); a cross-reference asserts nothing. Nor does the
item \emph{enumerate} the exceptions: the rule decides each claim by reading
that claim's own statement.
\end{remark}
\claimmeta{newly derived
\metasep accepted}

% claim receiver/rem:ia
\claimsummary{receiver/rem:ia}{A rigorous evaluation of the bound needs lower enclosures of c and \ensuremath{\gamma} and upper ones of the per-port d and \ensuremath{\beta}, each intersected with [0, 1], which excludes no exact value. In the window a rigorous lower bound on the block eigenvalue replaces the computed one, and \hyperlink{claim-receiver/ass:eig}{A-EIG} is not invoked.}
\claimbookmark{receiver/rem:ia}{Remark}{rem:ia}{Directions of rigorous enclosures of the bound inputs}
\claimleadin{%
The floating-point evaluation of Theorem~\ref{thm:exact-and-evaluated} leaves a gap between
the emitted and the exact bound. This remark records in which direction each
input has to be enclosed to close it; Requirement~\ref{req:rigorous-weight}
closes it in these directions, and Theorem~\ref{thm:enclosed-weight} proves
that the result is a bound.
}
\begin{remark}[Directions of rigorous enclosures of the bound inputs]\label{rem:ia}
Closing the gap means replacing the four inputs of \eqref{eq:beta} by
rigorous \emph{outer} enclosures taken in the conservative direction, formed
from enclosures of $B$, the $B_k$ and the prefactors rather than from their
floating-point evaluation. This engine forms them in exact rational
arithmetic, from the enclosures of the port Grams behind the reference of the
blocks (Theorem~\ref{thm:enclosed-blocks}), which are proved and need no
interval-arithmetic assumption, and from the least-eigenvalue bounds of the
foundations contract (Requirement~\ref{req:rigorous-weight}). A weight so
evaluated can carry a run to the \exttermref{../foundations/contract.pdf}{term-certified}{certified} tier
(\xref{foundations/def:tiers}) where the run's own certified evaluation
succeeds; a weight evaluated in floating point cannot. There are four
directions, not two. Reading
the signs off $\Bnd(n)=1+(K-1)\gamma c^{\,n}-\sum_k\beta_kd_k^{\,n}$ with
$n\ge1$ and $c,d_k\in[0,1]$:
\begin{itemize}[leftmargin=1.6em]
\item $c$: a \emph{lower} enclosure, i.e.\ an \emph{upper} enclosure of
  $\lmax(B)$ in \eqref{eq:cd};
\item $d_k$: an \emph{upper} enclosure, i.e.\ a \emph{lower} enclosure of
  $\lmin(B_k)$, since the coefficient $-\beta_k$ is non-positive;
\item $\gamma$: a \emph{lower} enclosure, since the coefficient
  $(K-1)c^{\,n}$ is non-negative;
\item each $\beta_k$: an \emph{upper} enclosure, since the coefficient
  $-d_k^{\,n}$ is non-positive.
\end{itemize}
The last two also preserve the hypothesis $\gamma\le\beta_k$ of
Theorem~\ref{thm:multiclick-bound}, which an enclosure in the opposite direction could
break.

\emph{Every one of the four is taken in $[0,1]$, and that is not optional.}
Each of $c,d_k,\gamma,\beta_k\in[0,1]$ is a \emph{theorem}: the
admissibility paragraph of Theorem~\ref{thm:multiclick-bound} derives
\eqref{eq:bounddata} from $0\preceq B_k\preceq B\preceq\id$ and
$\gamma\le\beta_k\le1$. So intersecting each rigorous enclosure with $[0,1]$
excludes no exact value, keeps it a valid enclosure, and keeps it on the
conservative side of \eqref{eq:beta}. Without the intersection a rigorous outer
enclosure can leave the interval in each of the four cases, and it can do so
at ordinary devices, not only at contrived ones: $\overline{\lmax(B)}>1$ where
$\lmax(B)=1$ exactly puts $\underline c<0$; $\underline{\lmin(B_k)}<0$ at a
singular $B_k$ (at a receiver without a probe, one with an input polarization
that reaches no port other than $k$, the configuration
Remark~\ref{rem:collapsesharp} discusses) puts $\overline{d_k}>1$; and
$\overline{\beta_k}>1$ is possible wherever $\beta_k=1$ exactly, as at every
probe-off setting. Data outside \eqref{eq:bounddata} is not
\termref{term-admissible-bound-data}{admissible bound data}, and \emph{both} lemmas that read the bound data use the
excluded part: Lemma~\ref{lem:windowtail}\,(i) obtains its tail step from
$d_k\le1$, and Lemma~\ref{lem:collapse}'s chord argument is taken on
$u_k\in[\beta_kc^{\,n},\beta_k]$, which is $d_k\in[c,1]$. Taken in $[0,1]$, the
enclosed tuple is admissible bound data, both lemmas apply to it unchanged, and
$\Bnd$ evaluated on it is a lower bound on $\Bnd$ evaluated on the exact
tuple; these are the \termref{term-enclosed-bound-data}{enclosed bound data} of
Definition~\ref{def:enclosed-weight}, and Theorem~\ref{thm:enclosed-weight}
states both facts.

\emph{The domain on which the enclosure of $c$ is taken, and why this is not
the clamping Definition~\ref{def:tolmap}\,(iv) forbids.} ``$\Bnd$ is
non-decreasing in $c$'' holds on $c\ge0$ and \emph{fails} below zero at even
$n$, which is exactly what Remark~\ref{rem:cdirection} states for the
floating-point evaluation; and a rigorous outer enclosure at a receiver
where $\lmax(B)=1$ exactly can put its lower end below zero. The
enclosure of $c$ is therefore taken \emph{in $[0,1]$}: the evaluation uses
$\underline c\coloneqq\max\{0,\,1-\overline{\lmax(B)}\}$. That is sound, and
for a reason available inside this contract: $c_{\mathrm{ex}}\in[0,1]$ is a
\emph{theorem} (the admissibility paragraph of Theorem~\ref{thm:multiclick-bound},
from $0\preceq B\preceq\id$), so intersecting a rigorous enclosure with $[0,1]$
loses no exact value, $\underline c\le c_{\mathrm{ex}}$ still holds, and on
$[0,1]$ the monotonicity does hold, so
$\Bnd(\underline c)\le\Bnd(c_{\mathrm{ex}})$ at every $n$. This is the
opposite of the situation of Definition~\ref{def:tolmap}\,(iv): there, no
enclosure exists, the computed $c$ is not known to bound $c_{\mathrm{ex}}$ on
either side, and clamping it to $0$ would assert a relation to the exact value
that nothing supports, so nothing is clamped and
Remark~\ref{rem:cdirection} states the resulting two-sided direction. Here, the
enclosure \emph{is} a rigorous bound and the intersection is with a domain the
contract proves; nothing is asserted that is not already a theorem.

\emph{A fifth direction, where an exact window is declared.} The four above
are the four inputs of \eqref{eq:beta}. Definition~\ref{def:pmcex} adds one
more computed quantity to the emitted bound, $\Bnd^{\ast}(n)$, and it takes a
\emph{lower} enclosure: it enters \eqref{eq:elldef} as a lower bound on
itself. As with the other four the enclosure is taken in $[0,1]$, which
excludes no exact value because $\Bnd^{\ast}(n)\in[0,1]$ is a theorem
(Lemma~\ref{lem:exactbasic} with Definition~\ref{def:setting-family}), and
as with $c$ this is an intersection with a
domain the contract proves rather than a clamp. A rigorous lower enclosure of
$\Bnd^{\ast}(n)$ is a \termref{term-rigorous-window-bounds}{rigorous window bound} of
Definition~\ref{def:enclosed-weight}: it takes the place of the computed
eigenvalue less the slack, so $\varepsilon_{\mathrm{eig}}$ does not enter the
weight and Assumption~\ref{ass:eig} is not invoked at that setting, the
weight being a bound by Theorem~\ref{thm:enclosed-weight}. This is the only
route in this remark that removes an assumption from the ledger rather than
tightening a number; the declared $\varepsilon_{\mathrm{eig}}$ stays for a
setting the engine evaluates in floating point.
\end{remark}
\claimmeta{\metalabel{Uses}\mref{Def.}{def:bounddata}, \mref{Thm.}{thm:multiclick-bound}, \mref{Lem.}{lem:collapse}, \mref{Rem.}{rem:collapsesharp}, \mref{Lem.}{lem:windowtail}, \mref{Lem.}{lem:exactbasic}, \mref{Def.}{def:pmcex}, \mref{Def.}{def:tolmap}, \mref{Rem.}{rem:cdirection}
\metasep \metalabel{Used by}\mref{Req.}{req:rigorous-weight}
\metasep newly derived
\metasep accepted}

\section{Assumptions}\label{group-receiver/ass:trust}
\sectionintro{This section states every assumption the derivations use, the implicit ones included. Most device assumptions spell out the trusted, exactly characterised receiver (\hyperlink{claim-receiver/ass:trust}{Assumption~\ref*{ass:trust}}): the declared instrument is exactly the physical receiver, with every optical path declared, passive, fed by vacuum on every undeclared mode, read by threshold detectors that respond to any photon in their ports, and fixed independently of the adversary. Round homogeneity of the receiver (\hyperlink{claim-receiver/ass:round}{Assumption~\ref*{ass:round}}) refines the protocol compiler's round-homogeneity premise instead. The eigenvalue-slack assumption is the one premise on the arithmetic, and only a setting with an exact window whose weight is evaluated in floating point invokes it. The engine verifies none of them; conditions decided on the declared data alone are definitions with requirements, not assumptions.}
\sectionimports{\extref{../foundations/contract.pdf}{foundations/def:round-homogeneity}}

% claim receiver/ass:trust
\claimsummary{receiver/ass:trust}{The declared instrument describes the physical receiver exactly, and the receiver is under the legitimate party's control. It is the parent of most of the receiver's device assumptions, which spell out its parts.}
\claimbookmark{receiver/ass:trust}{Assumption}{ass:trust}{Trusted and exactly characterised receiver}
\claimleadin{%
The receiver is trusted and known exactly. No uncertainty set, no outer Gram
family, no device family and no characterisation protocol is represented
anywhere in this contract; a receiver known only up to a stated uncertainty is
future scope (\href{../../../spec/future-scope.md}{\nolinkurl{spec/future-scope.md}}).
}
\SetShortId{A-TRUST}
\begin{assumption}[Trusted and exactly characterised receiver]
\label{ass:trust}
(i) The declared instrument describes the physical receiver \emph{exactly}:
its mode space, its \termref{term-declared-element-graph}{declared element graph} or the supplied
\termref{term-instrument-matrix}{instrument matrix} or \termref{term-port-gram-matrices}{port Gram family}, the
port map, the efficiencies, the dark-count probabilities and the probe states;
(ii) the receiver is under the legitimate party's control: it is not
selected, tuned or otherwise influenced by the adversary.
\end{assumption}
\claimverification{The engine does not verify this assumption; a run that rests on it records it as unverified.}
\claimreading{The receiver you declared is the receiver you run. Its mode space, its optical layout (or the instrument or Gram matrices you supplied), the assignment of outputs to detectors, the efficiencies, the dark-count probabilities and the probe states are those of the physical device, with no uncertainty to account for. The device is under your control, and nobody else can select, tune or influence it. A characterisation uncertainty is not modelled: the result applies to the device exactly as declared.}
\claimmeta{\metalabel{Category}receiver
\metasep \metalabel{Affects}secrecy
\metasep \metalabel{Assumed by}\mref{Lem.}{lem:deficiency}, \mref{Cor.}{cor:unconditioned-family}, \mref{Cor.}{cor:emitted-unconditioned}, \mref{Rem.}{rem:powervsclick}, \mref{Prop.}{prop:efficiency-normalization}
\metasep \metalabel{Machine key}\iface{receiver_trusted_exact}
\metasep newly derived
\metasep accepted}

% claim receiver/rem:unchecked-assumptions
\claimsummary{receiver/rem:unchecked-assumptions}{The assumptions on the physical receiver and the run, and \hyperlink{claim-receiver/ass:eig}{A-EIG} on the arithmetic, are never verified. The conditions decided on the declared data are definitions with requirements, and the floating-point evaluation is a restriction the contract places on itself.}
\claimbookmark{receiver/rem:unchecked-assumptions}{Remark}{rem:unchecked-assumptions}{Assumptions this module cannot check}
\begin{remark}[Assumptions this module cannot check]\label{rem:unchecked-assumptions}
Every assumption of this contract that is a statement about the physical
receiver or the run is one the engine never verifies, as the contract says
of each: no function of the \exttermref{../foundations/contract.pdf}{term-declared-data}{declared data} decides it, and validation
decides only those of its clauses that constrain the declared data.
Assumption~\ref{ass:eig} is a statement about the \emph{arithmetic}; it is
likewise never verified, and it is invoked only at a setting where a caller
declares an \termref{term-exact-window-depth}{exact-window depth} and the weight is evaluated in floating
point (Requirement~\ref{req:eig-ledger-row}), so a declaration that declares
none, or whose weights are all evaluated rigorously, is unaffected by it. The conditions that constrain the declared data alone are not
assumptions of this contract: the well-formedness of a declared element graph
(Definition~\ref{def:graph}) and the occupation-diagonality of the declared
probe states (Definition~\ref{def:probe}) are decided on the declared data
(Requirements~\ref{req:graph-refusal} and~\ref{req:probe-diagonal-admission}).
The floating-point evaluation of Definition~\ref{def:evaluation-model} is a
restriction this contract places on itself and asserts nothing that a run
could violate.
\end{remark}
\claimmeta{newly derived
\metasep accepted}

% claim receiver/ass:passive
\claimsummary{receiver/ass:passive}{Every optical element is passive linear optics, with no gain and no mixing of creation and annihilation operators. Total photon number is then a conserved grading, on which every block statement of the module rests.}
\claimbookmark{receiver/ass:passive}{Assumption}{ass:passive}{Passive linear optics}
\claimleadin{%
Under this assumption the total photon number is a conserved grading of the
receiver's modes.
}
\SetShortId{A-PASSIVE}
\begin{assumption}[Passive linear optics]\label{ass:passive}
(i) Every optical element is passive linear optics: it maps annihilation
operators to combinations of annihilation operators,
$b_i=\sum_j\Gmat_{ij}a_j$, with $\Gmat^\dagger\Gmat\preceq\id$;
(ii) no element mixes $a$ with $a^\dagger$, and none provides gain.
\end{assumption}
\claimverification{The engine does not verify this assumption; a run that rests on it records it as unverified.}
\claimreading{Every optical component of your receiver is passive and linear: beam splitters, wave plates, polarisers, attenuators and the like. No component amplifies light, and none is nonlinear (no frequency conversion, no squeezing). The total photon number is then never increased inside the receiver. If a component can add photons, the photon-number blocks the analysis rests on do not exist.}
\claimmeta{\metalabel{Category}receiver
\metasep \metalabel{Affects}secrecy
\metasep \metalabel{Refines}\aref{A-TRUST}{ass:trust}
\metasep \metalabel{Assumed by}\mref{Lem.}{lem:deficiency}, \mref{Cor.}{cor:unconditioned-family}, \mref{Cor.}{cor:emitted-unconditioned}, \mref{Rem.}{rem:out}, \mref{Prop.}{prop:efficiency-normalization}
\metasep \metalabel{Machine key}\iface{passive_optics}
\metasep newly derived
\metasep accepted}

% claim receiver/ass:vac
\claimsummary{receiver/ass:vac}{Undeclared input modes and dilation modes start in vacuum. A mode that can carry light to a detector is declared, as a channel mode if the adversary can populate it and as a source mode if the legitimate party does.}
\claimbookmark{receiver/ass:vac}{Assumption}{ass:vac}{Vacuum on undeclared input modes}
\claimleadin{%
Vacuum on every undeclared input is what lets the receiver be described by
the columns of its declared inputs alone.
}
\SetShortId{A-VAC}
\begin{assumption}[Vacuum on undeclared input modes]\label{ass:vac}
(i) Every mode of $\Modes\setminus\Inp$, and every mode of the dilation
$\Env$, carries the vacuum state at the input;
(ii) every mode that can carry a non-vacuum state and can couple to a port is
declared in $\Inp$: in $\InpC$ if the adversary can populate it, in $\InpS$
if the legitimate party does.
\end{assumption}
\claimverification{The engine does not verify this assumption; a run that rests on it records it as unverified.}
\claimreading{Every optical input of your receiver that you did not declare as an input, and every internal loss or ancilla mode, receives no light. A mode that can carry light to a detector is declared: as a channel mode if an attacker could put light into it, and as a source mode if your own equipment, for example a probe laser, does.}
\claimmeta{\metalabel{Category}receiver
\metasep \metalabel{Affects}secrecy
\metasep \metalabel{Refines}\aref{A-TRUST}{ass:trust}
\metasep \metalabel{Assumed by}\mref{Lem.}{lem:vacelim}, \mref{Cor.}{cor:compiled-instrument}, \mref{Lem.}{lem:deficiency}, \mref{Cor.}{cor:unconditioned-family}, \mref{Cor.}{cor:emitted-unconditioned}, \mref{Rem.}{rem:layers}, \mref{Rem.}{rem:bsconv}, \mref{Prop.}{prop:efficiency-normalization}
\metasep \metalabel{Machine key}\iface{undeclared_inputs_vacuum}
\metasep newly derived
\metasep accepted}

% claim receiver/ass:modes
\claimsummary{receiver/ass:modes}{Every optical path from outside the receiver to a detector is a declared mode. A bright-illumination or wavelength-shifted attack targets exactly this premise.}
\claimbookmark{receiver/ass:modes}{Assumption}{ass:modes}{Completeness of the declared modes}
\claimleadin{%
A photon arriving in a spatial, spectral, temporal or polarization mode that
the model does not carry does not exist for this module. The assumption is
therefore not free: it is exactly the one a bright-illumination or
wavelength-shifted attack violates, and violating it makes the \termref{term-emitted-povm}{emitted POVM} a
description of a device the adversary is not attacking.
}
\SetShortId{A-MODES}
\begin{assumption}[Completeness of the declared modes]\label{ass:modes}
(i) Every optical path from outside the receiver to a detector is represented
by a declared mode of $\Modes$.
\end{assumption}
\claimverification{The engine does not verify this assumption; a run that rests on it records it as unverified.}
\claimreading{Every way light from outside can reach one of your detectors (every spatial, spectral, temporal and polarisation mode it can arrive in) is represented by one of the declared modes. Light the model does not carry, for example at another wavelength or along a path straight to the detectors as in bright-illumination attacks, violates this. Declare the mode, or declare a model of the attack. This premise is about the completeness of the receiver's optical description; that nothing reaches the station except through its declared inputs, and that nothing leaks out of it, is the laboratory premise of this station's party, \hyperref{../meat-core/contract.pdf}{}{claim-meat-core/ass:lab}{A-LAB}.}
\claimmeta{\metalabel{Category}receiver
\metasep \metalabel{Affects}secrecy
\metasep \metalabel{Refines}\aref{A-TRUST}{ass:trust}
\metasep \metalabel{Assumed by}\mref{Lem.}{lem:deficiency}, \mref{Cor.}{cor:unconditioned-family}, \mref{Cor.}{cor:emitted-unconditioned}, \mref{Rem.}{rem:out}, \mref{Prop.}{prop:efficiency-normalization}
\metasep \metalabel{Machine key}\iface{no_undeclared_paths}
\metasep newly derived
\metasep accepted}

% claim receiver/ass:thresh
\claimsummary{receiver/ass:thresh}{Each detector reports one bit per gate, with no photon-number resolution, no memory or timing effects, and no correlations beyond the declared dark-count channel. It fixes the click-or-no-click outcome alphabet the receiver's measurement is built on.}
\claimbookmark{receiver/ass:thresh}{Assumption}{ass:thresh}{Threshold detectors without memory}
\SetShortId{A-THRESH}
\begin{assumption}[Threshold detectors without memory]\label{ass:thresh}
(i) Each detector reports one bit per detection gate: click or no click;
(ii) there is no photon-number resolution, no dead time, no afterpulsing and
no timing jitter;
(iii) there is no correlation of any kind between gates or between detectors
beyond the declared \termref{term-dark-count-channel}{dark-count channel};
(iv) dark counts are independent across rounds, and across ports unless a
correlated channel is declared explicitly (Definition~\ref{def:darkadm}).
\end{assumption}
\claimverification{The engine does not verify this assumption; a run that rests on it records it as unverified.}
\claimreading{Your detectors are threshold (click or no-click) detectors read once per gate. They have no photon-number resolution, no dead time, no afterpulsing and no timing jitter. Gates and detectors are correlated only through the dark-count model you declared. Dark counts are independent from round to round, and between detectors unless you declared a correlated dark-count channel. Detectors with dead time or afterpulsing at your gate rate violate this.}
\claimmeta{\metalabel{Uses}\mref{Def.}{def:modes}
\metasep \metalabel{Category}receiver
\metasep \metalabel{Affects}secrecy
\metasep \metalabel{Refines}\aref{A-TRUST}{ass:trust}
\metasep \metalabel{Assumed by}\mref{Cor.}{cor:unconditioned-family}, \mref{Cor.}{cor:emitted-unconditioned}, \mref{Rem.}{rem:out}, \mref{Rem.}{rem:corrdark}, \mref{Prop.}{prop:efficiency-normalization}
\metasep \metalabel{Machine key}\iface{threshold_memoryless_detectors}
\metasep newly derived
\metasep accepted}

% claim receiver/ass:port
\claimsummary{receiver/ass:port}{A port clicks exactly when the modes it collects carry at least one photon, whatever their distribution over those modes. It is the multimode form of \hyperlink{claim-receiver/ass:thresh}{A-THRESH}, and it gives a detector that collects several modes a meaning.}
\claimbookmark{receiver/ass:port}{Assumption}{ass:port}{Mode-insensitive port response}
\claimleadin{%
This is the multimode form of Assumption~\ref{ass:thresh}. It is easily left
implicit, but without it ``a detector that collects several modes'' is
undefined. Under it the click decision is invariant under any passive unitary
acting inside the port, since such a unitary preserves the total photon number
of the port's modes.
}
\SetShortId{A-PORT}
\begin{assumption}[Mode-insensitive port response]\label{ass:port}
(i) Port $k$ clicks if and only if the total photon number in the modes
$\Port{k}$ is nonzero; equivalently, its no-click element is the joint vacuum
projector $\Pi^0_k=\Fock(\id-\Emd{\{k\}})$ on those modes.
\end{assumption}
\claimverification{The engine does not verify this assumption; a run that rests on it records it as unverified.}
\claimreading{Each detector port clicks exactly when at least one photon arrives in any of the modes it collects, whatever their distribution over those modes. A detector whose click probability depends on the spatial, spectral or temporal mode of the arriving photon (efficiency mismatch) violates this. Such a mismatch has to be declared as mode-dependent loss in front of the port.}
\claimmeta{\metalabel{Category}receiver
\metasep \metalabel{Affects}secrecy
\metasep \metalabel{Refines}\aref{A-THRESH}{ass:thresh}
\metasep \metalabel{Assumed by}\mref{Lem.}{lem:vacelim}, \mref{Cor.}{cor:unconditioned-family}, \mref{Cor.}{cor:emitted-unconditioned}, \mref{Prop.}{prop:efficiency-normalization}
\metasep \metalabel{Machine key}\iface{mode_insensitive_ports}
\metasep newly derived
\metasep accepted}

% claim receiver/ass:exact
\claimsummary{receiver/ass:exact}{Declared component parameters, efficiencies, dark-count probabilities and probe mean photon numbers are exact. A non-zero parameter tolerance would be an uncertainty set, which is out of scope and is refused rather than ignored.}
\claimbookmark{receiver/ass:exact}{Assumption}{ass:exact}{Exact declared parameters}
\claimleadin{%
This is the parameter part of Assumption~\ref{ass:trust}. A parameter
tolerance is an uncertainty set and is out of scope: uncertainty sets on the
device parameters are future scope (\href{../../../spec/future-scope.md}{\nolinkurl{spec/future-scope.md}}~\S11), and
the reserved tolerance field that would carry one is refused unless absent or
exactly zero.
}
\SetShortId{A-EXACT}
\begin{assumption}[Exact declared parameters]\label{ass:exact}
(i) Every declared component parameter, efficiency, dark-count probability and
probe mean photon number is exact.
\end{assumption}
\claimverification{The engine does not verify this assumption; a run that rests on it records it as unverified.}
\claimreading{Every number you declared for a receiver component (each efficiency, dark-count probability, probe mean photon number and other component parameter) is the device's true value, not an estimate with an error bar. The engine refuses a non-zero tolerance field, and an uncertainty set is not modelled. A result computed from a measured value therefore applies to a device with exactly that value.}
\claimmeta{\metalabel{Category}receiver
\metasep \metalabel{Affects}secrecy
\metasep \metalabel{Refines}\aref{A-TRUST}{ass:trust}
\metasep \metalabel{Assumed by}\mref{Lem.}{lem:deficiency}, \mref{Cor.}{cor:unconditioned-family}, \mref{Cor.}{cor:emitted-unconditioned}, \mref{Prop.}{prop:efficiency-normalization}
\metasep \metalabel{Machine key}\iface{exact_parameters}
\metasep newly derived
\metasep accepted}

% claim receiver/ass:nosel
\claimsummary{receiver/ass:nosel}{The receiver configuration and the probe declaration are fixed before the run and cannot be chosen or influenced by the adversary. With \hyperlink{claim-receiver/ass:trust}{A-TRUST} this makes one declared POVM, rather than a family, the right object.}
\claimbookmark{receiver/ass:nosel}{Assumption}{ass:nosel}{Device fixed independently of the adversary}
\claimleadin{%
Together with Assumption~\ref{ass:trust}, this assumption is what makes a
single POVM, rather than a family, the right object.
}
\SetShortId{A-NOSEL}
\begin{assumption}[Device fixed independently of the adversary]
\label{ass:nosel}
(i) The component list, its parameters, the port map and the
\termref{term-probe-declaration}{probe declaration} are fixed before the protocol run;
(ii) none of them is chosen by the adversary, and none can be influenced by
her, in particular not through a degree of freedom the model does not carry,
such as wavelength.
\end{assumption}
\claimverification{The engine does not verify this assumption; a run that rests on it records it as unverified.}
\claimreading{Your receiver's configuration (its components and their settings, which output goes to which detector, and the probe settings) is fixed before the run. An attacker cannot choose or change it, including through properties the model does not describe, such as the wavelength of injected light shifting a beam splitter's ratio. If an attacker can steer the receiver this way, the declared measurement is the wrong one and the result does not apply. This premise is about the configuration of the receiver itself; what may leave or enter the station by other routes is the laboratory premise of this station's party, \hyperref{../meat-core/contract.pdf}{}{claim-meat-core/ass:lab}{A-LAB}.}
\claimmeta{\metalabel{Category}receiver
\metasep \metalabel{Affects}secrecy
\metasep \metalabel{Refines}\aref{A-TRUST}{ass:trust}
\metasep \metalabel{Assumed by}\mref{Lem.}{lem:deficiency}, \mref{Cor.}{cor:unconditioned-family}, \mref{Cor.}{cor:emitted-unconditioned}, \mref{Prop.}{prop:efficiency-normalization}
\metasep \metalabel{Machine key}\iface{device_not_adversarial}
\metasep newly derived
\metasep accepted}

% claim receiver/ass:round
\claimsummary{receiver/ass:round}{The same instrument, probe declaration and dark-count channel describe every round, and the receiver has no memory across rounds. One receiver artifact then describes the whole run.}
\claimbookmark{receiver/ass:round}{Assumption}{ass:round}{Round homogeneity of the instrument, probe and dark counts}
\claimleadin{%
Each round is described by the same three objects, independently of every
other round. Rounds are not averaged, and per-round variation is not modelled.
}
\SetShortId{A-ROUND}
\begin{assumption}[Round homogeneity of the instrument, probe and dark counts]\label{ass:round}
The family of the instrument, the probe declaration and the dark-count channel
is \exttermref{../foundations/contract.pdf}{term-round-homogeneous}{round-homogeneous} (\xref{foundations/def:round-homogeneity}) in both
clauses of that definition:
(i) the same instrument, the same probe declaration and the same dark-count
channel describe every round of the protocol;
(ii) the receiver acts in every round as they describe, whatever the inputs
and outputs of the other rounds, with no memory across rounds.
\end{assumption}
\claimverification{The engine does not verify this assumption; a run that rests on it records it as unverified.}
\claimreading{Your receiver behaves the same way in every round: the same optics, the same probe settings and the same dark-count behaviour. Nothing it does in one round depends on the inputs or outputs of other rounds: no detector memory, and no drift that the declaration averages over. A receiver whose efficiency drifts during the run violates this. Only parts of a run fixed before it starts, each processed as its own execution and composed under \hyperref{../meat-core/contract.pdf}{}{claim-meat-core/ass:separate-executions}{A-SEPARATE}, can be analysed separately; a split chosen after seeing the data is not covered. A receiver whose detectors remember earlier clicks, through afterpulsing, a dead time or any other dependence on earlier rounds, violates it too, and nothing in this engine repairs that: no input describes more than one round.}
\claimmeta{\metalabel{Uses}\extref{../foundations/contract.pdf}{foundations/def:round-homogeneity}
\metasep \metalabel{Category}receiver
\metasep \metalabel{Affects}secrecy
\metasep \metalabel{Refines}\extaref{../protocol-compiler/contract.pdf}{protocol-compiler/ass:rhom}{A-RHOM}
\metasep \metalabel{Assumed by}\mref{Lem.}{lem:deficiency}, \mref{Cor.}{cor:unconditioned-family}, \mref{Cor.}{cor:emitted-unconditioned}, \mref{Prop.}{prop:efficiency-normalization}
\metasep \metalabel{Machine key}\iface{round_homogeneity}
\metasep newly derived
\metasep accepted}

% claim receiver/ass:eig
\claimsummary{receiver/ass:eig}{Where an exact window is evaluated in floating point, the declared slack bounds any overestimate of the exact block minimum eigenvalue by the computed one. This one-sided premise on the arithmetic is the price of that window; a setting without one, or evaluated rigorously, does not invoke it.}
\claimbookmark{receiver/ass:eig}{Assumption}{ass:eig}{Declared slack on the computed block eigenvalues}
\claimleadin{%
On finitely many blocks the exact window takes the larger of the analytic
bound and a computed smallest eigenvalue less a declared slack. That the computed extreme
eigenvalue of a matrix of dimension $\binom{D+n-1}{n}$, $D$ being the
dimension of the \termref{term-channel-input-space}{channel input space}, does not overshoot the exact one
by more than the slack is a premise about the arithmetic, not a theorem. It is one-sided: since the analytic bound stays as a floor, an undershoot
costs tightness only. A declaration
that requests no exact window does not invoke it, and neither does a setting
whose weight is the \termref{term-enclosed-weight}{enclosed weight} of Definition~\ref{def:enclosed-weight},
where a proved lower bound on the block eigenvalue stands in place of the
computed one.
}
\SetShortId{A-EIG}
\begin{assumption}[Declared slack on the computed block eigenvalues]
\label{ass:eig}
(i) At every setting $t$ carrying a declared exact-window depth $m^{\ast}_t$
whose emitted weight is that of Definition~\ref{def:pmcex} formed by the
floating-point evaluation of Definition~\ref{def:evaluation-model}, and every
index $N<n\le m^{\ast}_t$, the value
$\widehat\Bnd^{\ast(t)}(n)$ an implementation computes for the exact block
minimum eigenvalue $\Bnd^{\ast(t)}(n)$ of \eqref{eq:exactbnd} satisfies
\begin{equation}\label{eq:eigassumption}
  \widehat\Bnd^{\ast(t)}(n)-\varepsilon_{\mathrm{eig}}
  \ \le\ \Bnd^{\ast(t)}(n) .
\end{equation}
\end{assumption}
\claimverification{The engine does not verify this assumption; a run that rests on it records it as unverified.}
\claimreading{Applies only if you request an exact-window depth at a setting whose weight the receiver evaluates in floating point. Wherever your declaration describes a valid receiver exactly, the receiver evaluates the weight rigorously instead, from a proved lower bound on each block eigenvalue, and does not use this premise; the record says which evaluation each setting got. Where it falls back to floating point, the eigensolver that computes the smallest eigenvalue of each requested high-photon block may overestimate it by at most the slack \ensuremath{\varepsilon}\_eig you declared. The engine cannot check this and trusts its floating-point eigensolver up to \ensuremath{\varepsilon}\_eig. Without an exact-window request this premise is not used, and the bound comes from a closed form that rests on no premise about the eigensolver.}
\claimmeta{\metalabel{Uses}\mref{Def.}{def:exactbnd}, \mref{Def.}{def:pmcex}, \mref{Def.}{def:slot}, \mref{Def.}{def:evaluation-model}
\metasep \metalabel{Category}numerical trust
\metasep \metalabel{Affects}secrecy
\metasep \metalabel{Assumed by}\mref{Rem.}{rem:exacteig}, \mref{Thm.}{thm:emitted-bound}, \mref{Cor.}{cor:emitted-unconditioned}, \mref{Thm.}{thm:exact-and-evaluated}
\metasep \metalabel{Machine key}\iface{exact_window_eigenvalue_slack}
\metasep newly derived
\metasep accepted}

\clearpage
\part{Engine requirements}
\sectionintro{What an implementation of this module must do, record or refuse; each requirement names the result that motivates it.}

\section{Requirements: compilation, representation and outcomes}\label{group-receiver/req:graph-refusal}
\sectionintro{These requirements bind the engine where the mathematics leaves a choice to the implementation. A component graph that is not \termref{term-well-formed}{well formed}, and a probe state outside the \termref{term-gaussian-tier}{Gaussian tier}, are refused rather than repaired, and the port map, the one-mode-per-port reading and the emitted alphabet are declared or decided, never inferred.}

% claim receiver/req:graph-refusal
\claimsummary{receiver/req:graph-refusal}{A declared element graph failing any of (E1) to (E3) is refused with the failing condition named, and the offending element or edge for (E1) or element pair for (E3). No matrix is emitted and no ordering is invented.}
\claimbookmark{receiver/req:graph-refusal}{Requirement}{req:graph-refusal}{Refusal of a graph that is not well formed}
\claimleadin{%
Theorem~\ref{thm:compilation-well-defined}\,(iv) makes each condition of a well-formed graph
decidable by finite inspection; this requirement says what the engine does
when one fails.
}
\begin{requirement}[Refusal of a graph that is not well formed]\label{req:graph-refusal}
The engine refuses a \termref{term-declared-element-graph}{declared element graph} (input \iface{element_graph})
that fails any of conditions (E1)--(E3) of Definition~\ref{def:graph}, with
the failing condition named, together with the offending element or edge
for (E1) and the offending element pair for (E3); no matrix is
emitted, and no ordering is invented to repair it. The refusal follows the
\iface{fail_closed} convention.
\end{requirement}
\claimmeta{\metalabel{Motivated by}\mref{Def.}{def:graph}, \mref{Thm.}{thm:compilation-well-defined}
\metasep implemented (effective: accepted)
\metasep \metalabel{Code}\coderef{https://github.com/fraunhofersg/certiqs-key/blob/v0.1.0/src/certiqs_key/receiver/graph.py}{certiqs_key.receiver.graph:check_element_graph}}

% claim receiver/req:probe-diagonal-admission
\claimsummary{receiver/req:probe-diagonal-admission}{The engine admits declared probe states of the Gaussian tier only, which are occupation-diagonal by construction, and refuses any other. Conditioning on an admitted probe therefore keeps the channel photon-number blocks.}
\claimbookmark{receiver/req:probe-diagonal-admission}{Requirement}{req:probe-diagonal-admission}{Admission of Gaussian-tier probe states only}
\claimleadin{%
The channel grading of the blocks is recovered from the total one only under
a condition on the probe state (Remark~\ref{rem:blockmeaning});
occupation-diagonality is the stronger form of that condition which
validation can check.
}
\begin{requirement}[Admission of Gaussian-tier probe states only]\label{req:probe-diagonal-admission}
The engine admits a \termref{term-probe-declaration}{probe declaration} (input \iface{probe_declaration})
only in the \termref{term-gaussian-tier}{Gaussian tier} of Definition~\ref{def:probe}, whose states
are \termref{term-occupation-diagonal}{occupation-diagonal} by construction, and refuses a declared probe
state outside it under the \iface{fail_closed} convention.
\end{requirement}
\claimmeta{\metalabel{Motivated by}\mref{Def.}{def:probe}, \mref{Rem.}{rem:blockmeaning}
\metasep implemented (effective: accepted)
\metasep \metalabel{Code}\coderef{https://github.com/fraunhofersg/certiqs-key/blob/v0.1.0/src/certiqs_key/receiver/settings.py}{certiqs_key.receiver.settings:check_probe_declaration}}

% claim receiver/req:declared-port-map
\claimsummary{receiver/req:declared-port-map}{The port map is an input of the declaration and is never inferred from the instrument matrix. Which output modes lie in a port has no observable content beyond the port Grams it produces.}
\claimbookmark{receiver/req:declared-port-map}{Requirement}{req:declared-port-map}{Declared port map}
\begin{requirement}[Declared port map]\label{req:declared-port-map}
The engine takes the \termref{term-detector-ports}{detector ports} as a declared input
(\iface{port_map}) and never infers them from the \termref{term-instrument-matrix}{instrument matrix}:
by Remark~\ref{rem:unmonitored}, the partition of the output modes into
``in some port'' and ``in no port'' has no observable content beyond the port
Grams it produces.
\end{requirement}
\claimmeta{\metalabel{Motivated by}\mref{Rem.}{rem:unmonitored}
\metasep accepted}

% claim receiver/req:one-mode-per-port
\claimsummary{receiver/req:one-mode-per-port}{The engine uses a one-mode-per-port representation of a receiver only where every emitted port Gram has rank at most one, the criterion of \hyperlink{claim-receiver/thm:representation-independence}{Theorem~\ref*{thm:representation-independence}} (iv), and never by default. Ranks are decided on the Grams after every detector efficiency is applied.}
\claimbookmark{receiver/req:one-mode-per-port}{Requirement}{req:one-mode-per-port}{One mode per port only where the rank criterion holds}
\claimleadin{%
A receiver architecture must not assume one \termref{term-optical-mode}{optical mode} per detector
unless rank one is established
(\href{../../../spec/architecture.md\#32-the-seven-cores}{\nolinkurl{spec/architecture.md}\texttt{\#32-the-seven-cores}}, the receiver core).
Remark~\ref{rem:rankone} explains why nothing else in this contract needs
the assumption.
}
\begin{requirement}[One mode per port only where the rank criterion holds]\label{req:one-mode-per-port}
The engine uses a one-mode-per-port representation of a receiver only where
the criterion of Theorem~\ref{thm:representation-independence}\,(iv), $\rk A_k\le1$ for every $k$,
holds, and never assumes it by default. The criterion is decided on the
emitted port Grams, that is after every detector efficiency is applied
(Remark~\ref{rem:effrank}), with each rank decided as
Requirement~\ref{req:tolerance-tests} fixes.
\end{requirement}
\claimmeta{\metalabel{Motivated by}\mref{Thm.}{thm:representation-independence}, \mref{Rem.}{rem:rankone}, \mref{Req.}{req:tolerance-tests}, \mref{Rem.}{rem:effrank}
\metasep implemented (effective: accepted)
\metasep \metalabel{Code}\coderef{https://github.com/fraunhofersg/certiqs-key/blob/v0.1.0/src/certiqs_key/receiver/grams.py}{certiqs_key.receiver.grams:one_mode_per_port}}

% claim receiver/req:declared-alphabet
\claimsummary{receiver/req:declared-alphabet}{The engine takes the emitted outcome alphabet and the coarse-graining onto it as declared inputs and never fixes an alphabet internally. The optical profile's six outcomes at four detectors are one declared instance, not the definition.}
\claimbookmark{receiver/req:declared-alphabet}{Requirement}{req:declared-alphabet}{Declared emitted alphabet and coarse-graining}
\claimleadin{%
The optical profile's six-outcome alphabet
$\Yset=\{\nc,\singlec_H,\singlec_V,\singlec_+,\singlec_-,\mc\}$
(\href{../../../spec/notation.md\#37-optical-profile-symbols}{\nolinkurl{spec/notation.md}\texttt{\#37-optical-profile-symbols}}) is the instance
$K=4$, $\cg(\varnothing)=\nc$, $\cg(\{k\})=\singlec_k$, and $\cg(S)=\mc$ for
$S\in\MCK$: one declared instance among many.
}
\begin{requirement}[Declared emitted alphabet and coarse-graining]\label{req:declared-alphabet}
The engine takes the emitted alphabet $\Vset$ as an input
(\iface{emitted_outcome_alphabet}) and the \termref{term-declared-coarse-graining}{declared coarse-graining}
$\cg$ of Definition~\ref{def:coarse} with it
(\iface{outcome_coarse_graining}); it never fixes an outcome alphabet
internally. A receiver that fixed six outcomes over four detectors
internally would make the outcome alphabet an implementation detail of the
receiver rather than a declared property of the protocol, and every
downstream index would then depend on an undeclared convention.
\end{requirement}
\claimmeta{\metalabel{Motivated by}\mref{Def.}{def:coarse}
\metasep accepted}

\section{Requirements: the emitted weight}\label{group-receiver/req:window-declaration}
\sectionintro{These requirements fix what a declaration must supply before any weight is computed and which declarations are refused. They also fix that the weight is evaluated rigorously wherever the port Grams are enclosed, with no \termref{term-eigenvalue-slack}{eigenvalue slack}, and what the artifact records beside each emitted weight, so that a reader can recompute it.}
\sectionimports{\extref{../foundations/contract.pdf}{foundations/cor:least-eigenvalue-over-enclosure}, \extref{../foundations/contract.pdf}{foundations/def:exact-reference}, \extref{../foundations/contract.pdf}{foundations/lem:least-eigenvalue-bound}, \extref{../foundations/contract.pdf}{foundations/req:rounding-direction}}

% claim receiver/req:window-declaration
\claimsummary{receiver/req:window-declaration}{A declaration must supply a window at every setting where predicate (X) fails, and one missing there is refused rather than defaulted. Which settings need one is read off the declaration; a window at or below the cutoff is refused everywhere, and one above it at an exact setting is unused.}
\claimbookmark{receiver/req:window-declaration}{Requirement}{req:window-declaration}{Declared window at every setting whose exact flag is general}
\claimleadin{%
Branch 2 of \eqref{eq:pmcdef} is evaluated on a declared window $m_t$, and
the contract fixes no default for it. Whether a setting takes that branch is a
property of the declaration: by Remark~\ref{rem:flag-instances} it is
read off the declaration alone, before any arithmetic.
}
\begin{requirement}[Declared window at every setting whose exact flag is general]\label{req:window-declaration}
A declaration supplying no window $m_t$ (\iface{tail_window}) at a setting
where \termref{term-predicate-x}{predicate (X)} of Definition~\ref{def:pmc} fails is refused
rather than defaulted; by Remark~\ref{rem:flag-instances} that condition is
decidable at input validation from \exttermref{../foundations/contract.pdf}{term-declared-data}{declared data} alone. The engine
therefore requires $m_t$ exactly at the settings whose \termref{term-exact-flag}{exact flag} is
$\mathrm{general}$, and at no others. A declared window satisfies $m_t>N$ (Definition~\ref{def:pmc}); one
that does not is refused at every setting, whatever its flag. A window
$m_t>N$ supplied at a setting whose flag is $\mathrm{exact}$ is unused, not
an error.
\end{requirement}
\claimmeta{\metalabel{Motivated by}\mref{Def.}{def:pmc}, \mref{Rem.}{rem:flag-instances}
\metasep implemented (effective: accepted)
\metasep \metalabel{Code}\coderef{https://github.com/fraunhofersg/certiqs-key/blob/v0.1.0/src/certiqs_key/receiver/validation.py}{certiqs_key.receiver.validation:check_tail_window}}

% claim receiver/req:exact-window-declaration
\claimsummary{receiver/req:exact-window-declaration}{A declaration carrying an exact-window depth must also declare the eigenvalue slack and the window-consistency threshold, and a missing one is refused, never defaulted. A depth at or below the cutoff, or above the declared window where predicate (X) fails, is refused.}
\claimbookmark{receiver/req:exact-window-declaration}{Requirement}{req:exact-window-declaration}{Declaration of an exact-window depth and its tolerances}
\begin{requirement}[Declaration of an exact-window depth and its tolerances]
\label{req:exact-window-declaration}
Declaring an \termref{term-exact-window-depth}{exact-window depth} (Definition~\ref{def:pmcex}) at any
setting requires $\varepsilon_{\mathrm{eig}}$ and $\varepsilon_6$ (the bound
on $r_6$ of Definition~\ref{def:residuals}); the engine refuses a
declaration missing either, rather than defaulting it, exactly as a missing
window $m_t$ is refused (Requirement~\ref{req:window-declaration}), with
\iface{exact_window_declared_without_required_tolerance}. The engine refuses
a depth violating $m^{\ast}_t>N$, or violating $m^{\ast}_t\le m_t$ at a
setting where \termref{term-predicate-x}{predicate (X)} of Definition~\ref{def:pmc} fails, with
\iface{exact_window_depth_outside_declared_range}.
\end{requirement}
\claimmeta{\metalabel{Motivated by}\mref{Def.}{def:pmc}, \mref{Def.}{def:pmcex}, \mref{Def.}{def:residuals}, \mref{Req.}{req:window-declaration}
\metasep implemented (effective: accepted)
\metasep \metalabel{Code}\coderef{https://github.com/fraunhofersg/certiqs-key/blob/v0.1.0/src/certiqs_key/receiver/validation.py}{certiqs_key.receiver.validation:check_exact_window}}

% claim receiver/req:exact-window-record
\claimsummary{receiver/req:exact-window-record}{The artifact carries every computed block eigenvalue of an exact window exactly as computed, and the eigenvalue slack enters only as the slack the window value subtracts. Nothing is clamped, so a reader can recompute each window value and see what the slack cost.}
\claimbookmark{receiver/req:exact-window-record}{Requirement}{req:exact-window-record}{Computed block eigenvalues carried as computed}
\begin{requirement}[Computed block eigenvalues carried as computed]
\label{req:exact-window-record}
Nothing is clamped, and the computed value is emitted: at every setting
carrying an \termref{term-exact-window-depth}{exact-window depth}, the engine carries each
$\widehat\Bnd^{\ast(t)}(n)$, $N<n\le m^{\ast}_t$, in the artifact exactly as
computed, in \iface{bound_quantifier_domain}, and
$\varepsilon_{\mathrm{eig}}$ enters only as the declared slack subtracted in
\eqref{eq:elldef} (Definition~\ref{def:pmcex}), so that a reader can
recompute \eqref{eq:elldef} and see what the slack cost.
\end{requirement}
\claimmeta{\metalabel{Motivated by}\mref{Def.}{def:pmcex}, \mref{Rem.}{rem:exacteig}
\metasep accepted}

% claim receiver/req:weight-provenance
\claimsummary{receiver/req:weight-provenance}{An exact window is evaluated only on request, never by default. The slot carries the per-setting weights, and the artifact says for each which evaluation and which definition produced it and at what depth; block eigenvalues may also be reported as diagnostics.}
\claimbookmark{receiver/req:weight-provenance}{Requirement}{req:weight-provenance}{Provenance of each emitted multiclick weight}
\begin{requirement}[Provenance of each emitted multiclick weight]
\label{req:weight-provenance}
The engine evaluates an exact window, that of Definition~\ref{def:pmcex} or
the \termref{term-rigorous-window-bounds}{rigorous window bounds} of Definition~\ref{def:enclosed-weight}, only
where a declaration requests it by an exact-window depth, and never by
default. Where it is not requested, the engine may also report block
eigenvalues as a diagnostic; the values entering the \termref{term-strengthening-slot}{strengthening slot}
(Definition~\ref{def:slot}), \iface{high_photon_bounds}, are the
\termref{term-per-setting-weight}{per-setting weights} $\pmc^{(t)}$ of
Theorem~\ref{thm:emitted-bound} --- at a receiver without a probe the single
value $\pmc$ --- each the \termref{term-enclosed-weight}{enclosed weight} of
Definition~\ref{def:enclosed-weight}, rounded down, where
Requirement~\ref{req:rigorous-weight} forms it, and otherwise the
floating-point evaluation of Definition~\ref{def:pmc} or~\ref{def:pmcex}; and
the artifact says, for each, which evaluation produced it, which of
Definitions~\ref{def:pmc} and~\ref{def:pmcex} gives its formula and at what
depth, in \iface{bound_quantifier_domain}.
\end{requirement}
\claimmeta{\metalabel{Motivated by}\mref{Def.}{def:pmc}, \mref{Def.}{def:pmcex}, \mref{Def.}{def:enclosed-weight}, \mref{Thm.}{thm:emitted-bound}, \mref{Def.}{def:slot}, \mref{Rem.}{rem:notsharp}, \mref{Req.}{req:rigorous-weight}
\metasep accepted}

% claim receiver/req:rigorous-weight
\claimsummary{receiver/req:rigorous-weight}{Wherever the port Grams are enclosed, the engine bounds the four inputs and the window blocks from rational enclosures, forms the enclosed weight exactly with no eigenvalue slack and emits it rounded down. Elsewhere it emits the floating-point weight and records why.}
\claimbookmark{receiver/req:rigorous-weight}{Requirement}{req:rigorous-weight}{Rigorous evaluation of the multiclick weight}
\claimleadin{%
Remark~\ref{rem:ia} reads off the \termref{term-bound-function}{bound function} the direction in which each
input of the weight must be bounded, and Theorem~\ref{thm:enclosed-weight}
shows that a weight formed from bounds taken in those directions is a bound
outright. The receiver holds what that needs: rational enclosures of its port
Grams for every declaration with a model admissible exactly
(Requirement~\ref{req:exact-route-decision}), from which the operators of the
bound, the prefactors and the window blocks are enclosed by the rules that
enclose its photon-number blocks, and the least-eigenvalue bounds of
\xref{foundations/lem:least-eigenvalue-bound} and
\xref{foundations/cor:least-eigenvalue-over-enclosure}, which turn an
enclosure of a matrix into bounds on its extreme eigenvalues.
}
\begin{requirement}[Rigorous evaluation of the multiclick weight]\label{req:rigorous-weight}
\begin{enumerate}[label=(\roman*),leftmargin=2.4em]
\item \emph{Where it applies.} On every declaration whose port Grams the
  engine encloses by Requirement~\ref{req:exact-route-decision}\,(iii), the
  engine evaluates the weight of every setting by (ii) to (iv), in exact
  rational arithmetic: the rigorous evaluation of
  Definition~\ref{def:evaluation-model}. At a setting where it cannot, it
  emits the weight of Definition~\ref{def:pmc} or~\ref{def:pmcex} evaluated in
  floating point, with the reason \iface{port_grams_not_enclosed} where the
  declaration has no model admissible exactly and
  \iface{bound_inputs_not_enclosed} where a reciprocal that (ii) forms
  receives no enclosure.
\item \emph{The enclosures.} At each setting $t$ the engine forms, from the
  enclosures of the port Grams by the rules of
  Lemma~\ref{lem:enclosure-arithmetic},
  \termref{term-complex-enclosure}{complex enclosures} of $B^{(t)}$ and of every
  $B^{(t)}_k$, of $\gamma^{(t)}$ and of every $\beta^{(t)}_k$, and, where $t$
  carries an \termref{term-exact-window-depth}{exact-window depth} $m^{\ast}_t$, of
  $\blk{M}{\eqn}_{\mc|t}$ for $N<n\le m^{\ast}_t$, each enclosure of one of
  these Hermitian matrices made Hermitian-centred by part~(v) of that lemma.
  For a receiver without a probe these are, by
  Corollary~\ref{cor:unconditioned-family}, $B=A_\Sigma$ and
  $B_k=A_{[K]\setminus\{k\}}$, sums of enclosed Grams, $\gamma=\beta_k=1$
  exactly, and the blocks of \eqref{eq:family-mc}, formed from the enclosures
  of Theorem~\ref{thm:enclosed-blocks}\,(ii) taken with $m^{\ast}_t$ as its
  cutoff. For a \termref{term-setting-family}{setting family} a submodule constructs they are the
  enclosures of its prefactors, its operators and its window blocks that the
  submodule states, formed by the same rules.
\item \emph{The four inputs and the window bounds.} For an
  \exttermref{../foundations/contract.pdf}{term-entrywise-enclosure}{entrywise enclosure} $(C,r)$ of a Hermitian matrix $X$ and a rational
  $\rho\ge0$ with $\rho^{2}\ge\sum_{i,l}r_{il}^{2}$, an upper bound on
  $\lmax(C)$ from \xref{foundations/lem:least-eigenvalue-bound}\,(iv) plus
  $\rho$ bounds $\lmax(X)$ from above, and a lower bound on $\lmin(C)$ from
  part~(ii) of that lemma less $\rho$ bounds $\lmin(X)$ from below
  (\xref{foundations/cor:least-eigenvalue-over-enclosure}\,(i),~(ii)); for a
  complex enclosure $(x_c,r)$ of a real $x$,
  $\mathrm{Re}\,x_c-r\le x\le\mathrm{Re}\,x_c+r$. By these rules the engine
  forms an upper bound $u$ on $\lmax(B^{(t)})$, lower bounds $\ell_k$ on
  $\lmin(B^{(t)}_k)$, a lower bound $g$ on $\gamma^{(t)}$, upper bounds $b_k$
  on $\beta^{(t)}_k$ and lower bounds $\lambda_n$ on
  $\lmin(\blk{M}{\eqn}_{\mc|t})$. Its \termref{term-enclosed-bound-data}{enclosed bound data} are
  \eqref{eq:enclosed-construction} of Theorem~\ref{thm:enclosed-weight}\,(i),
  each of the four intersected with $[0,1]$, and its
  \termref{term-rigorous-window-bounds}{rigorous window bounds} are $\max\{0,\lambda_n\}$. The output of a
  floating-point eigensolver may enter only as the basis $V$ of
  \xref{foundations/lem:least-eigenvalue-bound}, read as the exact rational
  matrix it is; a basis for which the engine forms no rational $\theta<1$
  satisfying the hypothesis of that lemma is replaced by the identity, which
  admits $\theta=0$. No computed eigenvalue is used as a bound, and
  $\varepsilon_{\mathrm{eig}}$ is not subtracted.
\item \emph{The weight.} The engine computes the \termref{term-enclosed-weight}{enclosed weight}
  $\check p^{(t)}_{\mc}$ (Definition~\ref{def:enclosed-weight}) as an exact
  rational and emits, as the setting's weight in the
  \termref{term-strengthening-slot}{strengthening slot} \iface{high_photon_bounds}
  (Definition~\ref{def:slot}), the largest binary64 number at or below it,
  rounded down as \xref{foundations/req:rounding-direction} rounds an emitted
  lower bound; by Theorem~\ref{thm:enclosed-weight}\,(iii) the emitted number
  bounds the \termref{term-multiclick-element}{multiclick element} of its setting on every block above the
  cutoff. The \termref{term-uniform-weight}{uniform weight} is the minimum of the emitted per-setting
  weights.
\item \emph{What stays.} The floating-point evaluation of the blocks, the
  residuals and their tests, and, at a setting carrying a depth, the computed
  values of Requirement~\ref{req:exact-window-record}, are formed as they are
  without this requirement. At a setting evaluated by (ii) to (iv), the
  floating-point weight of Definition~\ref{def:pmc} or~\ref{def:pmcex} enters
  no emitted number, and it may be recorded beside the emitted one as a
  diagnostic.
\item \emph{The record.} Per setting, in \iface{bound_quantifier_domain}: the
  evaluation that produced the emitted weight, \iface{rigorous} or
  \iface{floating_point}, with the reason of (i) where it is
  \iface{floating_point}; and, where it is \iface{rigorous}, the enclosed
  bound data, the rigorous window bounds and the enclosed weight as exact
  rationals, so that a reader can recompute the weight. A setting whose weight
  is \iface{rigorous} is not one at which Assumption~\ref{ass:eig} is invoked
  (Requirement~\ref{req:eig-ledger-row}), and the artifact's evaluation label
  follows Requirement~\ref{req:evaluation-label}.
\end{enumerate}
\end{requirement}
\claimreading{Wherever your receiver declaration describes a valid receiver exactly, the engine computes each high-photon multiclick weight it can from rational enclosures of the receiver's model, every bound taken on the safe side, and emits that weight rounded down: no floating-point eigenvalue is trusted and no eigenvalue slack is subtracted. The artifact records, setting by setting, whether it did. Where the declaration is valid only within the receiver tolerances, or a needed enclosure cannot be formed at a setting, the engine emits the floating-point weight there instead and records why; a key rate that reads such a weight is not certified.}
\claimmeta{\metalabel{Motivated by}\mref{Thm.}{thm:multiclick-bound}, \mref{Def.}{def:exactbnd}, \mref{Def.}{def:setting-family}, \mref{Cor.}{cor:unconditioned-family}, \mref{Def.}{def:pmc}, \mref{Def.}{def:pmcex}, \mref{Def.}{def:enclosed-weight}, \mref{Thm.}{thm:enclosed-weight}, \mref{Def.}{def:slot}, \mref{Def.}{def:complex-enclosure}, \mref{Lem.}{lem:enclosure-arithmetic}, \mref{Thm.}{thm:enclosed-blocks}, \mref{Def.}{def:evaluation-model}, \mref{Rem.}{rem:ia}, \mref{Ass.}{ass:eig}, \mref{Req.}{req:exact-window-record}, \mref{Req.}{req:evaluation-label}, \mref{Req.}{req:exact-route-decision}, \mref{Req.}{req:eig-ledger-row}, \extref{../foundations/contract.pdf}{foundations/cor:least-eigenvalue-over-enclosure}, \extref{../foundations/contract.pdf}{foundations/def:exact-reference}, \extref{../foundations/contract.pdf}{foundations/lem:least-eigenvalue-bound}, \extref{../foundations/contract.pdf}{foundations/req:rounding-direction}
\metasep \metalabel{Used by}\mref{Req.}{req:weight-provenance}, \mref{Req.}{req:evaluation-label}, \mref{Req.}{req:eig-ledger-row}
\metasep implemented (effective: accepted)
\metasep \metalabel{Code}\coderef{https://github.com/fraunhofersg/certiqs-key/blob/v0.1.0/src/certiqs_key/receiver/rigorous.py}{certiqs_key.receiver.rigorous:rigorous_weights}}

\section{Requirements: residuals, tolerances and the evaluation label}\label{group-receiver/req:residual-emission}
\sectionintro{These requirements say which residuals the engine emits, and that a residual or an admission test that fails refuses while nothing is repaired or clamped. The \termref{term-receiver-tolerances}{receiver tolerances} have a declared domain, and every artifact carries an evaluation label, the rigorous one only where every weight was evaluated rigorously. Whether a declaration defines an admissible receiver exactly, and whether it is on the \termref{term-exact-route}{exact route}, is decided in exact arithmetic and recorded with its cause, and the port Grams of every such declaration are enclosed.}
\sectionimports{\extref{../foundations/contract.pdf}{foundations/def:binary64}}

% claim receiver/req:residual-emission
\claimsummary{receiver/req:residual-emission}{Every artifact carries exactly the residuals that exist on its input route and for its request; the key of any other residual is absent, never zero. The tolerances of the first and third residual are declared on every route, that of the sixth is required only with an exact-window depth.}
\claimbookmark{receiver/req:residual-emission}{Requirement}{req:residual-emission}{Emission of the residuals that exist on the route and request}
\claimleadin{%
The artifact already follows this rule for the record of the receiver
construction, which is present exactly on route (iii); a residual that does
not exist on the route or for the request taken is likewise absent, because
emitting $0$ would report a check that was not made.
}
\begin{requirement}[Emission of the residuals that exist on the route and request]\label{req:residual-emission}
The engine emits with every artifact, in \iface{validation_results}, those
residuals of Definition~\ref{def:residuals} that exist on the input route taken
and for the evaluation requested, and no other: the keys of $r_1$ and $r_3$ are
absent, not zero, on route (i), and the key of $r_6$ is absent, not zero, at a
declaration that requests no exact window. The blocks over which
\eqref{eq:res5} maximises are the floating-point blocks
$\blk{M}{\eqn}_{v|t}$ of the coarse-grained elements that
Definition~\ref{def:residuals} names, before any \termref{term-dark-count-channel}{dark-count channel}, which
the \iface{photon_blocks} entry carries beside its dark-composed blocks
(Requirement~\ref{req:dark-composed-reference}); $r_5$ does not run over the
dark-composed blocks, on which the exact value $0$ of
Proposition~\ref{prop:res5}\,(iv) is not asserted.
$\varepsilon_1$ and $\varepsilon_3$ are declared on every route, route (i)
included, where they are unused (\iface{numerical_tolerances});
$\varepsilon_6$ is required exactly where an \termref{term-exact-window-depth}{exact-window depth} is
declared, on the same condition as $r_6$, and where it is declared elsewhere
it is unused and not an error.
\end{requirement}
\claimmeta{\metalabel{Motivated by}\mref{Def.}{def:residuals}, \mref{Prop.}{prop:res5}, \mref{Req.}{req:dark-composed-reference}
\metasep implemented (effective: accepted)
\metasep \metalabel{Code}\coderef{https://github.com/fraunhofersg/certiqs-key/blob/v0.1.0/src/certiqs_key/receiver/construct.py}{certiqs_key.receiver.construct:emitted_residuals}}

% claim receiver/req:residual-refusal
\claimsummary{receiver/req:residual-refusal}{Each emitted residual, and the residual gap of a declared window, is compared with its receiver tolerance, and one above it refuses the request. No block is Hermitized, symmetrized, projected, clamped or otherwise repaired: the artifact carries the computed block and its own defect.}
\claimbookmark{receiver/req:residual-refusal}{Requirement}{req:residual-refusal}{Refusal on a residual above its tolerance, and no repair}
\claimleadin{%
A repaired block is a different operator, and no claim of this contract holds
for it (Remark~\ref{rem:res5why}); so a residual only ever refuses.
}
\begin{requirement}[Refusal on a residual above its tolerance, and no repair]\label{req:residual-refusal}
The engine compares each residual that exists (Definition~\ref{def:residuals})
with its \termref{term-receiver-tolerances}{receiver tolerance} $\varepsilon_i$, and a
residual above its tolerance is a \emph{refusal}, not a warning; in particular
$r_5>\varepsilon_5$ refuses with the reason
\iface{photon_block_positivity_residual_above_tolerance}, and
$r_6>\varepsilon_6$ with the reason
\iface{window_eigenvalue_consistency_residual_above_tolerance}. In the same
way, at every setting where \termref{term-predicate-x}{predicate (X)} of
Definition~\ref{def:pmc} fails, the engine compares the residual gap
$g_{N,m_t}$ of \eqref{eq:resgap} on the declared window $m_t$ with
$\varepsilon_{\mathrm{gap}}$, and a gap above it refuses. The direction
is refuse-only: no block is Hermitized, symmetrized, projected onto the
positive cone, clamped or repaired in any way; the artifact carries the
computed block itself, and $r_5$ reports that block's own defect. The same
holds of $r_6$: it refuses and repairs nothing (Remark~\ref{rem:exacteig}
states the direction of failure it is blind to).
\end{requirement}
\claimmeta{\metalabel{Motivated by}\mref{Lem.}{lem:windowtail}, \mref{Def.}{def:pmc}, \mref{Rem.}{rem:exacteig}, \mref{Def.}{def:residuals}, \mref{Rem.}{rem:res5why}
\metasep implemented (effective: accepted)
\metasep \metalabel{Code}\coderef{https://github.com/fraunhofersg/certiqs-key/blob/v0.1.0/src/certiqs_key/receiver/residuals.py}{certiqs_key.receiver.residuals:check_residuals}}

% claim receiver/req:tolerance-tests
\claimsummary{receiver/req:tolerance-tests}{The engine decides each exact condition of \hyperlink{claim-receiver/def:tolmap}{Definition~\ref*{def:tolmap}} by the test given there, every rank by the \ensuremath{\varepsilon}\ensuremath{_{2}}-rank and by no other rule. Bound data failing their admission test are refused and no artifact is emitted.}
\claimbookmark{receiver/req:tolerance-tests}{Requirement}{req:tolerance-tests}{Decisions at the receiver tolerances, and refusal of failed bound data}
\claimleadin{%
The refusal rules of the module's interface
(\href{../../../spec/modules/receiver/interface.yaml}{\nolinkurl{spec/modules/receiver/interface.yaml}}) tabulate no reason string
for a bound-data failure, and this
contract adds none: fixing a contract-verbatim key here would settle an
interface detail that no claim of this contract consumes.
}
\begin{requirement}[Decisions at the receiver tolerances, and refusal of failed bound data]\label{req:tolerance-tests}
The engine decides each exact condition that Definition~\ref{def:tolmap} lists
by the test given there and by no other threshold, and every rank statement
this module decides on floating-point data by \eqref{eq:eps2rank}; no other
rule is licensed, there or at Theorem~\ref{thm:representation-independence}\,(iv). Bound data failing
item (iv) of Definition~\ref{def:tolmap} are a failed validation, so
\iface{fail_closed} applies at this site: no artifact is emitted. The reason
string is the implementer's to name, under the rule the refusal rules of the
module's interface state for the validations they do not tabulate, none of
the tabulated strings
covering this check.
\end{requirement}
\claimmeta{\metalabel{Motivated by}\mref{Thm.}{thm:representation-independence}, \mref{Def.}{def:tolmap}
\metasep implemented (effective: accepted)
\metasep \metalabel{Code}\coderef{https://github.com/fraunhofersg/certiqs-key/blob/v0.1.0/src/certiqs_key/receiver/construct.py}{certiqs_key.receiver.construct:construct_receiver}}

% claim receiver/req:tolerance-domain
\claimsummary{receiver/req:tolerance-domain}{Each receiver tolerance is validated at input against its declared domain [0, 1), and a value outside it is refused. At one or more every test the tolerance decides would be vacuous or degenerate (\hyperlink{claim-receiver/rem:tolerancedomain}{Remark~\ref*{rem:tolerancedomain}}).}
\claimbookmark{receiver/req:tolerance-domain}{Requirement}{req:tolerance-domain}{Validation of the receiver tolerances against their domain}
\begin{requirement}[Validation of the receiver tolerances against their domain]\label{req:tolerance-domain}
The engine validates each of the \termref{term-receiver-tolerances}{receiver tolerances}
$\varepsilon_1,\dots,\varepsilon_6$, $\varepsilon_{\mathrm{gap}}$ and
$\varepsilon_{\mathrm{eig}}$ of Definition~\ref{def:tolmap} at input against
the declared domain $[0,1)$, and refuses a value outside that interval with the
reason \iface{tolerance_outside_declared_domain}.
\end{requirement}
\claimmeta{\metalabel{Motivated by}\mref{Def.}{def:tolmap}, \mref{Rem.}{rem:tolerancedomain}
\metasep accepted}

% claim receiver/req:evaluation-label
\claimsummary{receiver/req:evaluation-label}{Every receiver artifact says whether its weights were all evaluated rigorously or not, and no path produces another label. A run whose declaration reads a per-setting weight evaluated in floating point, directly or through a uniform weight, reports at most a first-order tier.}
\claimbookmark{receiver/req:evaluation-label}{Requirement}{req:evaluation-label}{The evaluation label on every artifact}
\claimleadin{%
This contract claims no enclosure of any quantity it computes in floating
point (Definition~\ref{def:evaluation-model}), while the weights it evaluates
rigorously are bounds by theorem (Theorem~\ref{thm:enclosed-weight}). The label
says which of the two an artifact's weights are
(Theorem~\ref{thm:exact-and-evaluated}\,(iii)), and the requirement's last
sentence states the consequence for the tier of a run.
}
\begin{requirement}[The evaluation label on every artifact]\label{req:evaluation-label}
The engine writes \iface{bound_evaluation} on every artifact of this module,
on every input route: \iface{rigorous} where the weight of every setting was
formed by the rigorous evaluation of Requirement~\ref{req:rigorous-weight},
and \iface{floating_point} otherwise, an artifact with a weight resting on
Assumption~\ref{ass:eig} among them; which evaluation formed each setting's
weight is recorded as Requirement~\ref{req:rigorous-weight}\,(vi) says. No
path in this module produces another label
(Theorem~\ref{thm:exact-and-evaluated}\,(iii)): the label \iface{certified},
which would name a tier and not an evaluation, and the retired fallback reason
\iface{receiver_grams_not_enclosed} are never produced. A run whose
declaration carries a weight into which a \termref{term-per-setting-weight}{per-setting weight} enters
that the record of Requirement~\ref{req:rigorous-weight}\,(vi) names as
evaluated in floating point, the
\termref{term-uniform-weight-of-an-entry}{uniform weight} of a slot entry counting
the weight of every setting as entering it, is not reported at the certified
tier
(\xref{meat-core/frank-wolfe/req:certified-evaluation}\,(vii)); in every other
respect the tier a run reports from these artifacts follows
\xref{foundations/req:tiers}.
\end{requirement}
\claimmeta{\metalabel{Motivated by}\mref{Thm.}{thm:enclosed-weight}, \mref{Def.}{def:evaluation-model}, \mref{Thm.}{thm:exact-and-evaluated}, \mref{Ass.}{ass:eig}, \mref{Req.}{req:rigorous-weight}, \extref{../foundations/contract.pdf}{foundations/req:tiers}, \extref{../meat-core/contract.pdf}{meat-core/frank-wolfe/req:certified-evaluation}
\metasep accepted}

% claim receiver/req:exact-route-decision
\claimsummary{receiver/req:exact-route-decision}{The engine decides exactly whether a declaration has a model admissible exactly and, where its arithmetic represents the parameters, whether it is on the exact route. For every such model it encloses the port Grams in rational arithmetic, exactly on that route, and records which.}
\claimbookmark{receiver/req:exact-route-decision}{Requirement}{req:exact-route-decision}{Deciding the exact route and enclosing the port Grams}
\claimleadin{%
The engine encloses the port Grams of every declaration that defines an
admissible receiver exactly, so the receiver's blocks have a reference
whatever the parameters are (Theorem~\ref{thm:enclosed-blocks}); the
reference the receiver offers is formed from these enclosures
(Requirement~\ref{req:dark-composed-reference}). The \termref{term-exact-route}{exact route}
is decided in addition, because only exact entries let a snap keep a
block at weight zero (Theorem~\ref{thm:exact-blocks}). The entries of the
catalogue (Definition~\ref{def:catalogue-entries}) carry the square roots of
their transmittances, reflectivities, efficiencies and pass and leak
parameters, and of the complements $1-x$, $1-t_H$, $1-t_V$ and $1-y$ of the
splitting ratios and the losses; at axis angles that are multiples of
$\pi/8$ they involve otherwise only $\cos$ and $\sin$ of multiples of
$\pi/4$ and phases that are multiples of $\pi/4$, which lie in
$\mathbb{Q}(i,\sqrt2)$. That is why the floor of item (ii) names those square
roots and those angles. A compiled graph whose \termref{term-generic-block}{generic
blocks} are contractions exactly is admissible exactly at every other
parameter in its domain (Theorem~\ref{thm:catalogue-admissibility}), so on
that route the model needs no decision beyond those blocks and the domain
comparisons.
}
\begin{requirement}[Deciding the exact route and enclosing the port Grams]\label{req:exact-route-decision}
\begin{enumerate}[label=(\roman*),leftmargin=2.4em]
\item \emph{Declared Grams or matrix.} On a route that
  declares the port Grams or the \termref{term-instrument-matrix}{instrument matrix}, the engine computes
  the port Grams in exact rational arithmetic from the
  \exttermref{../foundations/contract.pdf}{term-declared-value}{declared values}, by \eqref{eq:gramdef} where the
  instrument matrix is declared, and decides in exact rational arithmetic
  whether every $A_k$ is Hermitian and whether the family satisfies
  \eqref{eq:exact-admissible}, that is, whether it is an
  \termref{term-admissible-gram-family}{admissible Gram family} exactly. The declaration has a model
  admissible exactly, and is then on the \termref{term-exact-route}{exact route}
  (Definition~\ref{def:exact-route}), exactly when both hold.
\item \emph{Component graph.} On the route that declares a component graph,
  the engine decides in exact rational arithmetic whether the matrix of every
  generic block, whose entries are declared values, is a contraction,
  $\Gmat_\Sigma^{\dagger}\Gmat_\Sigma\preceq\id$. Where every one is, the
  declaration has a model admissible exactly: the \termref{term-compiled-instrument-matrix}{compiled instrument matrix}
  is a contraction (Theorem~\ref{thm:compilation-well-defined}\,(ii)) and its
  port Grams are an admissible Gram family (Theorem~\ref{thm:port-grams}\,(i)),
  at the exact values of the declared parameters. For a declaration that
  has one and whose parameters its exact arithmetic represents, the engine
  moreover compiles the instrument matrix of Definition~\ref{def:compile} and
  the port Grams exactly and decides whether every port Gram has entries in
  $\mathbb{Q}(i)$, which places the declaration on the exact route. Every
  transmittance, loss, reflectivity, extinction and efficiency is the
  rational number the artifact carries for it: the declared value, or, for a
  loss declared in decibels, the converted efficiency that
  \iface{presentation_rule} makes the declared datum, a binary64 number and
  so a dyadic rational (\xref{foundations/def:binary64}\,(iii)). An angle
  parameter given in the presentation \iface{angle_presentation}, as a
  rational coefficient $c$ of $\pi$, is $c\pi$ exactly; an angle given in
  radians is a \exttermref{../foundations/contract.pdf}{term-declared-value}{declared value}, the dyadic rational it denotes. The
  arithmetic represents at least the field generated over $\mathbb{Q}(i)$ by
  $\sqrt2$ and by $\sqrt p$ and $\sqrt{1-p}$ for every declared rational
  parameter $p$ in $[0,1]$, which contains every entry of every catalogue
  entry whose angles are given as coefficients of $\pi$, whose axis angles
  are multiples of $\pi/8$, whose retardances are multiples of $\pi/2$ and
  whose rotations and phases are multiples of $\pi/4$. A declaration whose
  parameters the arithmetic does not represent, a non-zero angle given in
  radians among them, is treated as off the exact route, with its cause
  recorded by (iv), even where its port Grams happen to have entries in
  $\mathbb{Q}(i)$ (two half-wave plates in sequence at one such angle, for
  instance, compose to $-\id$).
\item \emph{Enclosure of the port Grams.} For every declaration with a
  model admissible exactly, the engine computes, in exact rational
  arithmetic, \termref{term-complex-enclosure}{complex enclosures}
  (Definition~\ref{def:complex-enclosure}) of its port Grams at the exact
  values of the declared parameters. On the exact route as decided by (i)
  or (ii) each port Gram is enclosed exactly, by itself. Otherwise, on the
  route of (ii), the
  enclosures are those of Theorem~\ref{thm:enclosed-blocks}\,(i), with the
  matrix of every element enclosed by the rules of
  Lemma~\ref{lem:enclosure-arithmetic} applied to the formula of its
  catalogue entry from \termref{term-square-root-enclosure}{square-root enclosures}
  and \termref{term-angle-enclosure}{angle enclosures} of its parameters, and
  $\sqrt{\eta_k}$ by a square-root enclosure; where the exact arithmetic of
  (ii) has computed a port Gram, its entries, which are expressions in square
  roots of rationals, may be enclosed from those expressions instead. A
  phase $e^{i\varphi}$ is $\cos\varphi+i\sin\varphi$, and the angle $\delta/2$
  of a retarder is again a rational or a rational multiple of $\pi$. No
  floating-point value enters an enclosure.
\item \emph{Record.} The \iface{photon_blocks} entry of the artifact records
  whether (i) or (ii) places the declaration on the exact route and, where
  neither does, why:
  a parameter or an angle that the exact arithmetic of (ii) does not
  represent (a non-zero angle given in radians among them), port Grams that
  are not Gaussian-rational, or no model admissible exactly, which is
  declared port Grams that are not Hermitian, a declared family that is not
  an admissible Gram family exactly (a declaration admitted only at the
  \termref{term-receiver-tolerances}{receiver tolerances}), or a generic block that is not a contraction
  exactly. Being off the exact route is not a refusal of the receiver and
  takes nothing away from the floating-point blocks (Definition~\ref{def:evaluation-model});
  where the declaration has a model admissible exactly, its port Grams are
  enclosed by (iii).
\end{enumerate}
\end{requirement}
\claimmeta{\metalabel{Motivated by}\mref{Def.}{def:gram}, \mref{Def.}{def:compile}, \mref{Thm.}{thm:compilation-well-defined}, \mref{Thm.}{thm:port-grams}, \mref{Def.}{def:complex-enclosure}, \mref{Lem.}{lem:enclosure-arithmetic}, \mref{Thm.}{thm:enclosed-blocks}, \mref{Def.}{def:exact-route}, \mref{Thm.}{thm:exact-blocks}, \mref{Def.}{def:tolmap}, \mref{Def.}{def:evaluation-model}, \mref{Def.}{def:catalogue-entries}, \mref{Thm.}{thm:catalogue-admissibility}, \mref{Req.}{req:dark-composed-reference}, \extref{../foundations/contract.pdf}{foundations/def:binary64}
\metasep \metalabel{Used by}\mref{Req.}{req:rigorous-weight}, \mref{Req.}{req:dark-composed-reference}, \mref{Req.}{req:snap-fix}
\metasep implemented (effective: accepted)
\metasep \metalabel{Code}\coderef{https://github.com/fraunhofersg/certiqs-key/blob/v0.1.0/src/certiqs_key/receiver/reference.py}{certiqs_key.receiver.reference:enclose_port_grams}}

\section{Requirements: the assumptions}\label{group-receiver/req:decidable-assumptions}
\sectionintro{The engine validates, where it is decidable, every clause of an assumption that constrains the declared data, and it refuses a non-zero parameter tolerance. It carries the eigenvalue-slack assumption beside the settings at which it was invoked, those whose weight uses a computed block eigenvalue.}

% claim receiver/req:decidable-assumptions
\claimsummary{receiver/req:decidable-assumptions}{Every clause of an assumption that constrains the declared data is enforced by validation wherever it is decidable. The assumptions the engine never verifies form the never-verified set, and no code path decides more of them than their clauses on the declared data.}
\claimbookmark{receiver/req:decidable-assumptions}{Requirement}{req:decidable-assumptions}{Validation of the assumptions on the declared data}
\claimleadin{%
Which assumptions the engine never verifies is recorded in each one's
verification field
(\href{../../../spec/claim-format.md\#32-whether-the-engine-verifies-an-assumption}{\nolinkurl{spec/claim-format.md}\texttt{\#32-whether-the-engine-verifies-an-assumption}}),
from which the contract generates the sentence that says so.
}
\begin{requirement}[Validation of the assumptions on the declared data]\label{req:decidable-assumptions}
The engine enforces by validation, wherever it is decidable, every clause of
an assumption of this contract that constrains the \exttermref{../foundations/contract.pdf}{term-declared-data}{declared data}. The
assumptions that the engine never verifies, as the contract says of each
(Remark~\ref{rem:unchecked-assumptions}), form \iface{never_verified_set},
and no code path decides what they say beyond those clauses.
\end{requirement}
\claimmeta{\metalabel{Motivated by}\mref{Rem.}{rem:unchecked-assumptions}
\metasep implemented
\metasep \metalabel{Code}\coderef{https://github.com/fraunhofersg/certiqs-key/blob/v0.1.0/src/certiqs_key/receiver/construct.py}{certiqs_key.receiver.construct:construct_receiver}}

% claim receiver/req:reserved-tolerance-field
\claimsummary{receiver/req:reserved-tolerance-field}{The reserved tolerance field of a declared parameter is admitted only when absent or exactly zero and is refused otherwise. An ignored tolerance would be a silent claim that the device is exact (\hyperlink{claim-receiver/ass:exact}{A-EXACT}).}
\claimbookmark{receiver/req:reserved-tolerance-field}{Requirement}{req:reserved-tolerance-field}{Refusal of a non-zero parameter tolerance}
\claimleadin{%
Assumption~\ref{ass:exact} makes every declared parameter exact, and a
tolerance silently ignored would be a claim of exactness that was not made;
the field is therefore validated and never ignored. Its place in the
declaration is the parameters table of the interface
(\href{../../../spec/modules/receiver/interface.yaml}{\nolinkurl{spec/modules/receiver/interface.yaml}}).
}
\begin{requirement}[Refusal of a non-zero parameter tolerance]\label{req:reserved-tolerance-field}
The engine validates the reserved \iface{tolerance} field of every declared
parameter as absent or exactly zero, and refuses it otherwise.
\end{requirement}
\claimmeta{\metalabel{Motivated by}\mref{Ass.}{ass:exact}
\metasep implemented
\metasep \metalabel{Code}\coderef{https://github.com/fraunhofersg/certiqs-key/blob/v0.1.0/src/certiqs_key/receiver/validation.py}{certiqs_key.receiver.validation:check_parameter_tolerances}}

% claim receiver/req:eig-ledger-row
\claimsummary{receiver/req:eig-ledger-row}{The assumption ledger carries \hyperlink{claim-receiver/ass:eig}{A-EIG} beside the settings at which it was invoked: those whose emitted weight uses a computed block eigenvalue. A reader of the artifact sees which emitted weights rest on the declared eigenvalue slack.}
\claimbookmark{receiver/req:eig-ledger-row}{Requirement}{req:eig-ledger-row}{Invoking settings on the eigenvalue-slack row of the ledger}
\begin{technicalclaim}
\begin{requirement}[Invoking settings on the eigenvalue-slack row of the ledger]\label{req:eig-ledger-row}
The engine carries Assumption~\ref{ass:eig} in the artifact's assumption
ledger \iface{assumption_ledger} beside the settings at which it was invoked.
A setting is one at which it was invoked exactly when its emitted weight is
that of Definition~\ref{def:pmcex}, evaluated in floating point; a setting
whose weight is an \termref{term-enclosed-weight}{enclosed weight} of Definition~\ref{def:enclosed-weight},
formed by Requirement~\ref{req:rigorous-weight}, is not one, whatever depth it
declares.
\end{requirement}
\end{technicalclaim}
\claimmeta{\metalabel{Motivated by}\mref{Def.}{def:pmcex}, \mref{Def.}{def:enclosed-weight}, \mref{Ass.}{ass:eig}, \mref{Req.}{req:rigorous-weight}
\metasep accepted}

\clearpage
\part{Component catalogue}
\sectionintro{The component catalogue lets a receiver be declared as an optical layout instead of a single matrix. Each of its ten entries fixes a support, a parameter domain and a passive matrix with every phase and every loss the part has written out (\hyperlink{claim-receiver/component-catalogue/def:catalogue-entries}{Definition~\ref*{def:catalogue-entries}}). The main result is that every entry maps its whole declared domain into local elements, lossless exactly on a named subfamily (\hyperlink{claim-receiver/component-catalogue/thm:catalogue-admissibility}{Theorem~\ref*{thm:catalogue-admissibility}}); one attenuation lemma carries most of the proof. The amplitude each entry removes has rows of its own, none on the lossless subfamily, and along a compiled graph they make a port dilation of the receiver (\hyperlink{claim-receiver/component-catalogue/lem:catalogue-dilations}{Lemma~\ref*{lem:catalogue-dilations}}). Detector efficiency and the probe source are declared layers outside the catalogue because they act at different points of the measurement. A worked BB84 receiver follows one declaration through compilation to its port Grams, its analytic high-photon weight and the improvement an exact window buys at a smaller cutoff.}

\section{Catalogue entries and their admissibility}\label{group-receiver/component-catalogue/lem:attenuated}
\sectionintro{The main result, \hyperlink{claim-receiver/component-catalogue/thm:catalogue-admissibility}{Theorem~\ref*{thm:catalogue-admissibility}}, is that every entry maps its whole declared domain into \termref{term-local-element}{local elements}, lossless exactly on the subfamily its table row names. Every entry but the \termref{term-generic-block}{generic block} is written as a unitary core with a positive contractive attenuation on one or both sides, so one lemma gives this for each of them. The generic block is the one entry whose admission is a spectral test, decided at the route-(ii) receiver tolerance (\hyperlink{claim-receiver/component-catalogue/def:c10-tolerance}{Definition~\ref*{def:c10-tolerance}}). The remarks place the rotator outside the retarder family, give the physical roles of the entries, separate the detector and source layers from the chain, and say which convention factors are free. The rows that carry what each entry removes, stacked along a compiled graph with those of the detector efficiencies, make a \termref{term-port-dilation}{port dilation} of the receiver (\hyperlink{claim-receiver/component-catalogue/lem:catalogue-dilations}{Lemma~\ref*{lem:catalogue-dilations}}).}

% claim receiver/component-catalogue/lem:attenuated
\claimsummary{receiver/component-catalogue/lem:attenuated}{A unitary with a positive semidefinite contraction on either side is a passive local element, lossless exactly when both contractions are the identity. Every catalogue entry has this form, so the lemma carries most of the catalogue's admissibility proof.}
\claimbookmark{receiver/component-catalogue/lem:attenuated}{Lemma}{lem:attenuated}{Admissibility of an attenuated unitary}
\claimleadin{%
The catalogue has ten \termref{term-entry-type}{entry types}, each a short matrix with
a short verification. Every one of them is a unitary core with a positive
semidefinite attenuation on one or both sides, so the following lemma does
almost all of the work.
}
\begin{technicalclaim}
\begin{lemma}[Admissibility of an attenuated unitary]\label{lem:attenuated}
Let $V$ be unitary on $\mathbb{C}^\Sigma$ and let $D_1,D_2$ be positive
semidefinite contractions on $\mathbb{C}^\Sigma$. Then
$\Gmat_\Sigma\coloneqq D_2VD_1$ satisfies
$\Gmat_\Sigma^\dagger\Gmat_\Sigma\preceq\id_\Sigma$, and
$\Gmat_\Sigma^\dagger\Gmat_\Sigma=\id_\Sigma$ if and only if
$D_1=D_2=\id_\Sigma$.
\end{lemma}
\end{technicalclaim}
\claimmeta{\metalabel{Used by}\mref{Thm.}{thm:catalogue-admissibility}
\metasep newly derived
\metasep accepted}

\begin{proof}
\emph{The bound.} Conjugation by a unitary preserves the Löwner order, so
$V^\dagger D_2^2V\preceq V^\dagger\id V=\id$. Since $X\preceq Y$ implies
$Z^\dagger XZ\preceq Z^\dagger YZ$ for any $Z$,
$\Gmat_\Sigma^\dagger\Gmat_\Sigma=D_1(V^\dagger D_2^2V)D_1\preceq D_1\id
D_1=D_1^2\preceq\id$.

\emph{The equality case.} If the composite equals $\id$, the chain forces
$D_1^2=\id$. A positive semidefinite contraction with $D_1^2=\id$ has all its
eigenvalues in $[0,1]$ squaring to $1$, hence equal to $1$, so $D_1=\id$. Then
$V^\dagger D_2^2V=\id$ gives $D_2^2=\id$, and the same argument gives
$D_2=\id$. Conversely, $D_1=D_2=\id$ gives $\Gmat_\Sigma=V$, which is
unitary.
\end{proof}

% claim receiver/component-catalogue/def:catalogue-entries
\claimsummary{receiver/component-catalogue/def:catalogue-entries}{Ten entry types, from the beamsplitter C1 to the generic block C10, each fix a support, a parameter domain and a passive matrix with every phase and every loss the part has written out. A table collects the supports, the domains and the lossless subfamilies.}
\claimbookmark{receiver/component-catalogue/def:catalogue-entries}{Definition}{def:catalogue-entries}{Entries of the component catalogue}
\claimleadin{%
A receiver declared as an optical layout is built from the entries below,
which form the \termref{term-component-catalogue}{component catalogue} of
Definition~\ref{def:catalogue}. Each
entry is written the way the part is specified: every phase and every loss
the part has appears as a parameter, and none is dropped because it could
have been absorbed elsewhere.
}
\begin{definition}[Entries of the component catalogue]\label{def:catalogue-entries}
The \termdef{term-catalogue-entries}{catalogue entries} are the entry types C1--C10 below.
Throughout, a \termdef{term-path}{path} is a spatial index $p$ and carries the two
polarization modes $(p,H)$ and $(p,V)$; a matrix written on two paths acts as
$\cdot\otimes\id_2$ on polarization unless stated otherwise, and a matrix
written on polarization acts on one path.

\emph{C1 --- Beamsplitter.} Support: two paths $(p_1,p_2)$, four modes.
Parameters: transmittance $x\in[0,1]$, reflection phase
$\varphi\in[0,2\pi)$, per-output losses $y_1,y_2\in[0,1]$. Core, on the ordered
path pair with rows $=$ outputs:
\begin{equation}\label{eq:bs}
  \Bsp(x,\varphi)=
  \begin{pmatrix}
    \sqrt{x} & \sqrt{1-x}\,e^{i\varphi}\\[2pt]
    \sqrt{1-x}\,e^{-i\varphi} & -\sqrt{x}
  \end{pmatrix},
  \qquad
  \Gmat_{\mathrm{C1}}=
  \Big[\operatorname{diag}\big(\sqrt{1-y_1},\sqrt{1-y_2}\big)\,
       \Bsp(x,\varphi)\Big]\otimes\id_2 .
\end{equation}
A \termdef{term-polarization-dependent-beamsplitter}{polarization-dependent beamsplitter} declares $(x_H,\varphi_H)$ and
$(x_V,\varphi_V)$ separately, with the per-output losses $y_1,y_2$ of C1, and
is
$\operatorname{diag}\big(\sqrt{1-y_1},\sqrt{1-y_2},\sqrt{1-y_1},\sqrt{1-y_2}\big)
\big[\Bsp(x_H,\varphi_H)\oplus\Bsp(x_V,\varphi_V)\big]$ in the reordered basis
$(p_1H,p_2H,p_1V,p_2V)$; this is the matrix of entry C2 at
$(t_H,\varphi_H,t_V,\varphi_V)=(x_H,\varphi_H,x_V,\varphi_V)$ written in that
basis.

\emph{C2 --- Polarizing beamsplitter.} Support: two paths, four modes.
Parameters: $t_H,t_V\in[0,1]$ (ideal $t_H=1$, $t_V=0$), phases
$\varphi_H,\varphi_V\in[0,2\pi)$, losses $y_1,y_2\in[0,1]$. In the ordering
$(p_1H,p_1V,p_2H,p_2V)$,
\begin{equation}\label{eq:pbs}
  \Gmat_{\mathrm{C2}}=
  \operatorname{diag}\big(\sqrt{1-y_1},\sqrt{1-y_1},
                          \sqrt{1-y_2},\sqrt{1-y_2}\big)\;
  \Xi^\dagger\big[\Bsp(t_H,\varphi_H)\oplus\Bsp(t_V,\varphi_V)\big]\,\Xi ,
\end{equation}
where $\Xi$ is the permutation from $(p_1H,p_1V,p_2H,p_2V)$ to
$(p_1H,p_2H,p_1V,p_2V)$.

\emph{C3 --- \termdef{term-retarder}{Retarder}, with half- and quarter-wave presets.} Support: the two
polarization modes of one path. Parameters: retardance $\delta\in[0,2\pi)$,
fast-axis angle $\theta\in[0,\pi)$, scalar loss $y\in[0,1]$. With
$\Rot(\theta)=\left(\begin{smallmatrix}\cos\theta&-\sin\theta\\
\sin\theta&\cos\theta\end{smallmatrix}\right)$,
\begin{equation}\label{eq:retarder}
  \Wplate(\delta,\theta)
  \coloneqq e^{-i\delta/2}\,\Rot(\theta)\operatorname{diag}(1,e^{i\delta})
    \Rot(\theta)^{\transpose}
  = e^{-i\delta/2}
  \begin{pmatrix}
    \cos^2\theta+e^{i\delta}\sin^2\theta &
    (1-e^{i\delta})\sin\theta\cos\theta\\[2pt]
    (1-e^{i\delta})\sin\theta\cos\theta &
    \sin^2\theta+e^{i\delta}\cos^2\theta
  \end{pmatrix},
\end{equation}
and $\Gmat_{\mathrm{C3}}=\sqrt{1-y}\,\Wplate(\delta,\theta)$.
The presets are
\begin{equation}\label{eq:hwpqwp}
  \mathrm{HWP}(\theta)\coloneqq\Wplate(\pi,\theta)
  =-i\begin{pmatrix}\cos2\theta&\sin2\theta\\
                    \sin2\theta&-\cos2\theta\end{pmatrix},
  \qquad
  \mathrm{QWP}(\theta)\coloneqq\Wplate(\tfrac{\pi}{2},\theta),
\end{equation}
so that $\mathrm{HWP}(\pi/8)=-\tfrac{i}{\sqrt2}
\left(\begin{smallmatrix}1&1\\1&-1\end{smallmatrix}\right)$.

\emph{C4 --- Circular retarder, or polarization \termdef{term-rotator}{rotator}.} Support: the two
polarization modes of one path. Parameter: rotation $\rho\in[0,2\pi)$.
$\Gmat_{\mathrm{C4}}=\Rot(\rho)$.

\emph{C5 --- Mirror or fold.} Support: one path. Parameters: reflectivities
$R_H,R_V\in[0,1]$, phases $\varphi_H,\varphi_V\in[0,2\pi)$, handedness flip
$\mathrm{fl}\in\{\mathrm{false},\mathrm{true}\}$ (default $\mathrm{false}$).
\begin{equation}\label{eq:mirror}
  \Gmat_{\mathrm{C5}}=\operatorname{diag}\big(\sqrt{R_H},\sqrt{R_V}\big)\,
  \operatorname{diag}\big(e^{i\varphi_H},e^{i\varphi_V}\big)\,
  F,\qquad
  F=\begin{cases}\operatorname{diag}(1,-1)&\mathrm{fl}=\mathrm{true}\\
                 \id_2&\text{otherwise.}\end{cases}
\end{equation}

\emph{C6 --- Phase shifter.} Support: any declared mode subset $\Sigma$.
Parameters: $\varphi_m\in[0,2\pi)$ per $m\in\Sigma$.
$\Gmat_{\mathrm{C6}}=\operatorname{diag}(e^{i\varphi_m})_{m\in\Sigma}$.

\emph{C7 --- Loss.} Support: any declared mode subset $\Sigma$. Parameters:
$\eta_m\in[0,1]$ per $m\in\Sigma$, optionally declared as
$L_{\mathrm{dB}}\ge0$ with $\eta=10^{-L_{\mathrm{dB}}/10}$.
$\Gmat_{\mathrm{C7}}=\operatorname{diag}(\sqrt{\eta_m})_{m\in\Sigma}$.

\emph{C8 --- Polarizer.} Support: the two polarization modes of one path.
Parameters: transmission axis $\theta\in[0,\pi)$, pass amplitude $\sqrt{t}$ and
leak amplitude $\sqrt{t_\perp}$ with $t,t_\perp\in[0,1]$. With
$\ket{e_\theta}=\cos\theta\ket{H}+\sin\theta\ket{V}$,
\begin{equation}\label{eq:polarizer}
  \Gmat_{\mathrm{C8}}
  =\sqrt{t}\,\dyad{e_\theta}+\sqrt{t_\perp}\,\dyad{e_\theta^{\perp}}
  =\Rot(\theta)\operatorname{diag}\big(\sqrt{t},\sqrt{t_\perp}\big)
   \Rot(\theta)^{\transpose}.
\end{equation}

\emph{C9 --- Mode routing.} Support: any declared mode subset.
Parameter: a permutation $\pi$ of the modes of $\Sigma$;
$(\Gmat_{\mathrm{C9}})_{ij}=\delta_{i,\pi(j)}$.

\emph{C10 --- \termdef{term-generic-block}{Generic block}.} Support: any declared mode subset $\Sigma$.
Parameter: a matrix $\Gmat_\Sigma\in\mathbb{C}^{\Sigma\times\Sigma}$, admitted
after the explicit check $\Gmat_\Sigma^\dagger\Gmat_\Sigma\preceq\id_\Sigma$.

Table~\ref{tab:catalogue} collects the supports and the parameter domains.

\begin{table}[htbp]
\centering
{\small
\begin{tabular}{@{}l >{\raggedright\arraybackslash}p{0.18\textwidth}
                  >{\raggedright\arraybackslash}p{0.30\textwidth}
                  >{\raggedright\arraybackslash}p{0.32\textwidth}@{}}
\hline
ID & Support & Parameters (domain) &
Lossless exactly when \\
\hline
C1 & 2 paths (4 modes) & $x\in[0,1]$, $\varphi\in[0,2\pi)$,
$y_1,y_2\in[0,1]$ & $y_1=y_2=0$ \\
C2 & 2 paths (4 modes) & $t_H,t_V\in[0,1]$, $\varphi_H,\varphi_V\in[0,2\pi)$,
$y_1,y_2\in[0,1]$ & $y_1=y_2=0$ \\
C3 & 1 path, polarization & $\delta\in[0,2\pi)$, $\theta\in[0,\pi)$,
$y\in[0,1]$ & $y=0$ \\
C4 & 1 path, polarization & $\rho\in[0,2\pi)$ & always \\
C5 & 1 path, polarization & $R_H,R_V\in[0,1]$,
$\varphi_H,\varphi_V\in[0,2\pi)$, $\mathrm{fl}$ & $R_H=R_V=1$ \\
C6 & any $\Sigma$ & $\varphi_m\in[0,2\pi)$ & always \\
C7 & any $\Sigma$ & $\eta_m\in[0,1]$ & every $\eta_m=1$ \\
C8 & 1 path, polarization & $\theta\in[0,\pi)$, $t,t_\perp\in[0,1]$ &
$t=t_\perp=1$ (then $\Gmat_\Sigma=\id$) \\
C9 & any $\Sigma$ & a permutation of $\Sigma$ & always \\
C10 & any $\Sigma$ & $\Gmat_\Sigma$ with
$\Gmat_\Sigma^\dagger\Gmat_\Sigma\preceq\id$ (checked) &
$\Gmat_\Sigma^\dagger\Gmat_\Sigma=\id$ (checked) \\
\hline
\end{tabular}}
\caption{The catalogue entries. Support shapes are stated in modes; the
parameter domains are the declared domains of
Definition~\ref{def:graph}\,(E1); the last column names the subfamily on
which Theorem~\ref{thm:catalogue-admissibility} proves the entry lossless.}\label{tab:catalogue}
\end{table}
\end{definition}
\claimmeta{\metalabel{Used by}\mref{Thm.}{thm:catalogue-admissibility}, \mref{Lem.}{lem:catalogue-dilations}
\metasep newly derived
\metasep implemented
\metasep \metalabel{Code}\coderef{https://github.com/fraunhofersg/certiqs-key/blob/v0.1.0/src/certiqs_key/receiver/catalogue.py}{certiqs_key.receiver.catalogue:Component}}

% claim receiver/component-catalogue/thm:catalogue-admissibility
\claimsummary{receiver/component-catalogue/thm:catalogue-admissibility}{Every catalogue entry maps its whole declared parameter domain to contractions on its modes, lossless exactly on the subfamily the catalogue's table lists for it. Only the generic block, an arbitrary matrix on declared modes, needs its contractivity checked at admission.}
\claimbookmark{receiver/component-catalogue/thm:catalogue-admissibility}{Theorem}{thm:catalogue-admissibility}{Admissibility of the catalogue entries}
\claimleadin{%
A \termref{term-declared-element-graph}{declared element graph} is \termref{term-well-formed}{well formed} only if every element
is a contraction on its support. For entries C1--C9 this holds by theorem on
the whole declared domain, so for each of them condition (E1) of
Definition~\ref{def:graph} reduces to checking that its parameters lie in
that domain.
}
\begin{mainclaim}
\begin{theorem}[Admissibility of the catalogue entries]\label{thm:catalogue-admissibility}
Let $\Gmat_\Sigma$ be the matrix of an entry type C1--C10 of
Definition~\ref{def:catalogue-entries} at a parameter tuple in its declared
domain. Then:
\begin{enumerate}[label=(\roman*),leftmargin=2.2em]
\item $\Gmat_\Sigma^\dagger\Gmat_\Sigma\preceq\id_\Sigma$, so $\Gmat_\Sigma$
  is a \termref{term-local-element}{local element} in the sense of Definition~\ref{def:embed};
\item $\Gmat_\Sigma$ is lossless,
  $\Gmat_\Sigma^\dagger\Gmat_\Sigma=\id_\Sigma$, exactly on the subfamily
  named in Table~\ref{tab:catalogue};
\item for C1--C9, a matrix whose parameter tuple passes the domain comparison
  of (E1) (Definition~\ref{def:graph}) is contractive by (i); C10 is the only
  entry whose declared domain is the contractivity condition itself, so that
  its check of Theorem~\ref{thm:compilation-well-defined}\,(iv)(E1) is a
  spectral test on the supplied matrix.
\end{enumerate}
\end{theorem}
\end{mainclaim}
\claimmeta{\metalabel{Uses}\mref{Def.}{def:embed}, \mref{Def.}{def:graph}, \mref{Def.}{def:compile}, \mref{Thm.}{thm:compilation-well-defined}, \mref{Lem.}{lem:attenuated}, \mref{Def.}{def:catalogue-entries}
\metasep \metalabel{Used by}\mref{Def.}{def:c10-tolerance}, \mref{Rem.}{rem:rotator}, \mref{Rem.}{rem:entry-roles}, \mref{Ex.}{ex:bb84}
\metasep newly derived
\metasep implemented (effective: accepted)
\metasep \metalabel{Code}\coderef{https://github.com/fraunhofersg/certiqs-key/blob/v0.1.0/src/certiqs_key/receiver/catalogue.py}{certiqs_key.receiver.catalogue:Component}
\metasep \metalabel{Oracles}\coderef{https://github.com/fraunhofersg/certiqs-key/blob/v0.1.0/tests/oracles/registry_receiver_povm.py}{registry_receiver_povm.py}}

\begin{proof}
\emph{Idea.} Each of the entries C1--C9 is exhibited in the form $D_2VD_1$ of
Lemma~\ref{lem:attenuated}, after which only the unitarity of the core has to
be checked, and the lemma then gives (i) and the lossless case (ii) together.
Entry C10 is decided by its hypothesis, checked on the supplied matrix.

\emph{C1.} $\Bsp(x,\varphi)$ is unitary: its first column has squared norm
$x+(1-x)=1$, its second $(1-x)+x=1$, and their inner product is
$\sqrt{x}\sqrt{1-x}e^{i\varphi}+\overline{\sqrt{1-x}e^{-i\varphi}}(-\sqrt{x})
=\sqrt{x(1-x)}e^{i\varphi}-\sqrt{x(1-x)}e^{i\varphi}=0$. Tensoring with
$\id_2$ preserves unitarity, and
$\operatorname{diag}(\sqrt{1-y_1},\sqrt{1-y_2})\otimes\id_2$ is a positive
semidefinite contraction because $1-y_i\in[0,1]$. Lemma~\ref{lem:attenuated}
with $D_1=\id$ gives admissibility and losslessness exactly at $y_1=y_2=0$.
(At $\varphi=0$ the core is real with $\det\Bsp=-1$; that sign is an artefact
of the convention, see Remark~\ref{rem:bsconv}.)

\emph{C2.} $\Bsp(t_H,\varphi_H)\oplus\Bsp(t_V,\varphi_V)$ is a direct sum of
unitaries, hence unitary; conjugating by the permutation $\Xi$ preserves
unitarity; the prefactor is again a positive semidefinite diagonal contraction.
The polarization-dependent beamsplitter of entry C1 is this matrix written in
the reordered basis $(p_1H,p_2H,p_1V,p_2V)$, so the same argument covers it.
At $t_H=1$, $t_V=0$, $\varphi_H=\varphi_V=0$, $y_1=y_2=0$ the element is the
signed permutation
$p_1H\mapsto p_1H$, $p_1V\mapsto p_2V$, $p_2H\mapsto -p_2H$,
$p_2V\mapsto p_1V$, which is the intended routing up to per-mode phases.

\emph{C3.} $\Rot(\theta)$ is real orthogonal ($\Rot(\theta)^{\transpose}
\Rot(\theta)=\id$ by $\cos^2+\sin^2=1$) and
$\operatorname{diag}(1,e^{i\delta})$ is unitary, so
$\Wplate(\delta,\theta)$ is a product of unitaries times the unimodular scalar
$e^{-i\delta/2}$, hence unitary; the expanded form in \eqref{eq:retarder} is
the product carried out, and \eqref{eq:hwpqwp} is $\delta=\pi$, using
$e^{i\pi}=-1$, $\cos^2\theta-\sin^2\theta=\cos2\theta$ and
$2\sin\theta\cos\theta=\sin2\theta$, with the scalar $e^{-i\pi/2}=-i$. The
prefactor $\sqrt{1-y}$ is a scalar contraction.

\emph{C4.} $\Rot(\rho)$ is real orthogonal, hence unitary, with no attenuation
factor.

\emph{C5.} $\operatorname{diag}(e^{i\varphi_H},e^{i\varphi_V})F$ is a product
of unitaries ($F$ is diagonal with entries of modulus one in both cases), and
$\operatorname{diag}(\sqrt{R_H},\sqrt{R_V})$ is a positive semidefinite
contraction, equal to $\id$ exactly when $R_H=R_V=1$.

\emph{C6, C9.} Diagonal with unimodular entries, respectively a permutation
matrix: both unitary, both with $D_1=D_2=\id$.

\emph{C7.} $V=\id$, $D_2=\operatorname{diag}(\sqrt{\eta_m})$, a positive
semidefinite contraction equal to $\id$ exactly when every $\eta_m=1$.

\emph{C8.} The two displayed forms agree because
$\Rot(\theta)\ket{H}=\ket{e_\theta}$ and $\Rot(\theta)\ket{V}=
\ket{e_\theta^{\perp}}$, so conjugating $\operatorname{diag}(\sqrt t,
\sqrt{t_\perp})$ by $\Rot(\theta)$ is the spectral decomposition written in
the rotated basis. The result is Hermitian positive semidefinite with
eigenvalues $\sqrt t,\sqrt{t_\perp}\in[0,1]$, so it is itself a positive
semidefinite contraction: Lemma~\ref{lem:attenuated} applies with
$V=\id$, $D_1=\id$, and equality holds exactly at $t=t_\perp=1$.

\emph{C10.} The condition is the hypothesis, checked on the supplied matrix;
losslessness is the same check with equality.

\emph{Item (iii).} For C1--C9 it is (i), which holds on the whole declared
domain; for C10 the declared domain is, by
Definition~\ref{def:catalogue-entries}, the set of blocks $\Gmat_\Sigma$ with
$\Gmat_\Sigma^\dagger\Gmat_\Sigma\preceq\id_\Sigma$.
\end{proof}

% claim receiver/component-catalogue/lem:catalogue-dilations
\claimsummary{receiver/component-catalogue/lem:catalogue-dilations}{Every catalogue entry has a loss factor, one row per amplitude it can remove and no row on its lossless subfamily; a Gaussian-rational generic block has one from an exact decomposition, which exists. Stacked along a compiled graph with the efficiency layer's loss rows, they make a port dilation.}
\claimbookmark{receiver/component-catalogue/lem:catalogue-dilations}{Lemma}{lem:catalogue-dilations}{Loss factors of the catalogue entries and the dilation of a graph}
\claimleadin{%
Theorem~\ref{thm:catalogue-admissibility} writes every entry but the generic
block as a unitary core with positive semidefinite attenuations on its sides
(Lemma~\ref{lem:attenuated}). The amplitude an attenuation removes can be
sent to new modes, one per attenuated mode: the core and the attenuation
together with those rows are then an isometry, and the rows vanish exactly
where nothing is attenuated. Done element by element along a compiled graph,
and once more for the detector efficiencies, this gives a
\termref{term-port-dilation}{port dilation} of the receiver whose loss modes are only those its
declared losses need, which Theorem~\ref{thm:factor-blocks} reads the
photon-number blocks through.
}
\begin{lemma}[Loss factors of the catalogue entries and the dilation of a graph]\label{lem:catalogue-dilations}
For a diagonal matrix $D$ write $D^{\flat}$ for the matrix of the rows of $D$
whose diagonal entry is not zero, so that
$(D^{\flat})^{\dagger}D^{\flat}=D^{\dagger}D$.
\begin{enumerate}[label=(\roman*),leftmargin=2.2em]
\item \emph{Loss factors.} Let $\Gmat_\Sigma$ be the matrix of an entry type
  of Definition~\ref{def:catalogue-entries} at a parameter tuple in its
  declared domain, in the notation of that definition, and let $\Theta$ be
  the matrix given for its type below, each of its rows indexed by a mode of
  its own:
  \begin{itemize}
  \item C1: $\Theta=\big(\operatorname{diag}(\sqrt{y_1},\sqrt{y_2})^{\flat}\,\Bsp(x,\varphi)\big)\otimes\id_2$;
  \item C2: $\Theta=\operatorname{diag}(\sqrt{y_1},\sqrt{y_1},\sqrt{y_2},\sqrt{y_2})^{\flat}\,
    \Xi^{\dagger}\big[\Bsp(t_H,\varphi_H)\oplus\Bsp(t_V,\varphi_V)\big]\Xi$;
  \item C3: $\Theta=\operatorname{diag}(\sqrt{y},\sqrt{y})^{\flat}$;
  \item C4, C6 and C9: $\Theta$ with no row;
  \item C5: $\Theta=\operatorname{diag}\big(\sqrt{1-R_H},\sqrt{1-R_V}\big)^{\flat}$;
  \item C7: $\Theta=\operatorname{diag}\big(\sqrt{1-\eta_m}\big)_{m\in\Sigma}^{\flat}$;
  \item C8: $\Theta=\operatorname{diag}\big(\sqrt{1-t},\sqrt{1-t_\perp}\big)^{\flat}\,\Rot(\theta)^{\transpose}$;
  \item C10, with the entries of $\Gmat_\Sigma$ in $\mathbb{Q}(i)$:
    $\Theta=D^{1/2}E$, where
    $\id_\Sigma-\Gmat_\Sigma^{\dagger}\Gmat_\Sigma=E^{\dagger}DE$ with $D$
    diagonal with positive rational diagonal entries and $E$ with entries in
    $\mathbb{Q}(i)$; such a decomposition exists, with at most $|\Sigma|$
    rows.
  \end{itemize}
  Then $\Gmat_\Sigma^{\dagger}\Gmat_\Sigma+\Theta^{\dagger}\Theta=\id_\Sigma$,
  and $\Theta$ is a \emph{loss factor} of the entry. For C1--C9, $\Theta$ has
  no row exactly on the lossless subfamily named in
  Table~\ref{tab:catalogue}.
\item \emph{The dilation of a compiled graph.} Let $\Recv$ be a
  \termref{term-receiver-instrument}{receiver instrument} whose \termref{term-instrument-matrix}{instrument matrix} is the
  \termref{term-compiled-instrument-matrix}{compiled instrument matrix}
  $\Gmat=D_\eta\Gmat^{\sigma}_{\Grph}$ of a well-formed graph $\Grph$ of
  catalogue elements for a \termref{term-linear-extension}{linear extension}
  $\sigma=(e_1,\dots,e_L)$ (Definition~\ref{def:compile}), with ports
  $\Port{1},\dots,\Port{K}$ and efficiencies $\eta_k$. For each $j$ let
  $\Theta_j$ be a loss factor of (i) of the element $e_j$, its rows indexed
  by a mode set $\Env_j$, and let
  $\hat\Theta_j:\mathbb{C}^{\Modes}\to\mathbb{C}^{\Env_j}$ act as $\Theta_j$
  on the modes of the support $\Sigma_{e_j}$ and as zero on the others. Let
  $\Theta_\eta\coloneqq\big(\sum_{k}\sqrt{1-\eta_k}\,\Emd{\{k\}}\big)^{\flat}$,
  its rows indexed by a mode set $\Env_\eta$, one for each mode of a port
  with $\eta_k<1$, the sets $\Modes$, $\Env_\eta$ and the $\Env_j$ being
  pairwise disjoint. With $\iota:\mathbb{C}^{\Inp}\to\mathbb{C}^{\Modes}$ the
  inclusion of the \termref{term-declared-input-modes}{declared input modes}, put $X_0\coloneqq\iota$ and
  $X_j\coloneqq\hat\Gmat_{e_j}X_{j-1}$ for $j=1,\dots,L$, and let
  \begin{equation}\label{eq:graph-dilation}
    \Omega:\mathbb{C}^{\Inp}\to\mathbb{C}^{\Modes\sqcup\Env_\eta\sqcup\Env_L\sqcup\dots\sqcup\Env_1}
  \end{equation}
  have the rows $D_\eta X_L$ on $\Modes$, $\Theta_\eta X_L$ on $\Env_\eta$
  and $\hat\Theta_jX_{j-1}$ on $\Env_j$. Then $\Omega$, with the ports
  $\Port{1},\dots,\Port{K}$ of $\Recv$, is a port dilation of $\Recv$
  (Definition~\ref{def:port-dilation}), and its rows on $\Modes$ are those of
  $\Gmat_{\Inp}$.
\end{enumerate}
\end{lemma}
\claimmeta{\metalabel{Uses}\mref{Def.}{def:embed}, \mref{Def.}{def:instrument}, \mref{Def.}{def:gram}, \mref{Def.}{def:graph}, \mref{Def.}{def:compile}, \mref{Def.}{def:port-dilation}, \mref{Def.}{def:catalogue-entries}
\metasep newly derived
\metasep implemented (effective: accepted)
\metasep \metalabel{Code}\coderef{https://github.com/fraunhofersg/certiqs-key/blob/v0.1.0/src/certiqs_key/receiver/factor.py}{certiqs_key.receiver.factor:port_dilation}}

\begin{proof}
\emph{The flat matrix.} For $D=\operatorname{diag}(a_1,\dots,a_r)$,
$(D^{\flat})^{\dagger}D^{\flat}=\sum_{i:\,a_i\ne0}\lvert a_i\rvert^{2}e_ie_i^{\dagger}
=D^{\dagger}D$, the rows omitted being zero. So in every case of (i) but C10,
$\Theta^{\dagger}\Theta$ is the same product with the full diagonal matrix,
and below $\Theta^{\dagger}\Theta$ is computed with it.

\emph{(i).} \emph{C1.} $\Bsp(x,\varphi)$ is unitary: its columns have
squared norms $x+(1-x)=1$ and $(1-x)+x=1$, and their inner product is
$\sqrt{x}\sqrt{1-x}\,e^{i\varphi}-\sqrt{1-x}\,e^{i\varphi}\sqrt{x}=0$. With
$D_2\coloneqq\operatorname{diag}(\sqrt{1-y_1},\sqrt{1-y_2})$ and
$(Z\otimes\id_2)^{\dagger}(Z\otimes\id_2)=Z^{\dagger}Z\otimes\id_2$,
$\Gmat_\Sigma^{\dagger}\Gmat_\Sigma+\Theta^{\dagger}\Theta
=\big(\Bsp^{\dagger}(D_2^{2}+\operatorname{diag}(y_1,y_2))\Bsp\big)\otimes\id_2
=\Bsp^{\dagger}\Bsp\otimes\id_2=\id$.
\emph{C2.} $V\coloneqq\Xi^{\dagger}[\Bsp(t_H,\varphi_H)\oplus\Bsp(t_V,\varphi_V)]\Xi$
is unitary, a direct sum of unitaries conjugated by a permutation; with
$D_2$ the diagonal prefactor of \eqref{eq:pbs},
$\Gmat_\Sigma^{\dagger}\Gmat_\Sigma+\Theta^{\dagger}\Theta
=V^{\dagger}\big(D_2^{2}+\operatorname{diag}(y_1,y_1,y_2,y_2)\big)V=V^{\dagger}V=\id$.
\emph{C3.} $\Wplate(\delta,\theta)$ is unitary, a product of the real
orthogonal $\Rot(\theta)$, the unitary $\operatorname{diag}(1,e^{i\delta})$,
$\Rot(\theta)^{\transpose}$ and the unimodular $e^{-i\delta/2}$; so
$\Gmat_\Sigma^{\dagger}\Gmat_\Sigma=(1-y)\id_2$ and
$\Theta^{\dagger}\Theta=y\,\id_2$.
\emph{C4, C6, C9.} $\Rot(\rho)$ is real orthogonal, a diagonal matrix with
unimodular entries is unitary, and so is a permutation matrix; so
$\Gmat_\Sigma^{\dagger}\Gmat_\Sigma=\id$, and $\Theta^{\dagger}\Theta=0$, an
empty sum.
\emph{C5.} The three factors of \eqref{eq:mirror} are diagonal, so they
commute, $F^{\dagger}F=\id$ and
$\lvert e^{i\varphi_H}\rvert=\lvert e^{i\varphi_V}\rvert=1$; so
$\Gmat_\Sigma^{\dagger}\Gmat_\Sigma=\operatorname{diag}(R_H,R_V)$ and
$\Theta^{\dagger}\Theta=\operatorname{diag}(1-R_H,1-R_V)$.
\emph{C7.} $\Gmat_\Sigma^{\dagger}\Gmat_\Sigma=\operatorname{diag}(\eta_m)$
and $\Theta^{\dagger}\Theta=\operatorname{diag}(1-\eta_m)$.
\emph{C8.} With $\Rot\coloneqq\Rot(\theta)$ real orthogonal, the second form
of \eqref{eq:polarizer} gives
$\Gmat_\Sigma^{\dagger}\Gmat_\Sigma=\Rot\operatorname{diag}(t,t_\perp)\Rot^{\transpose}$,
and $\Theta^{\dagger}\Theta=\Rot\operatorname{diag}(1-t,1-t_\perp)\Rot^{\transpose}$;
their sum is $\Rot\Rot^{\transpose}=\id$.
\emph{C10.} $\Theta^{\dagger}\Theta=E^{\dagger}(D^{1/2})^{2}E=E^{\dagger}DE
=\id_\Sigma-\Gmat_\Sigma^{\dagger}\Gmat_\Sigma$, $D^{1/2}$ being the real
diagonal matrix of the positive square roots. Such a decomposition exists. Put
$H\coloneqq\id_\Sigma-\Gmat_\Sigma^{\dagger}\Gmat_\Sigma$: it is Hermitian,
positive semidefinite, since $\Gmat_\Sigma^{\dagger}\Gmat_\Sigma\preceq\id$ on
the declared domain (Definition~\ref{def:catalogue-entries}), and has entries
in $\mathbb{Q}(i)$. A diagonal entry $H_{jj}=0$ forces $H_{jl}=0$ for every
$l$, since $\lvert H_{jl}\rvert^{2}\le H_{jj}H_{ll}$ for a positive
semidefinite $H$. If $H=0$, take $E$ and $D$ with no row. Otherwise some
$h\coloneqq H_{pp}\neq0$; it is real and lies in $\mathbb{Q}(i)$, so it is a
positive rational. Let $e$ be the row with entries $e_l\coloneqq H_{pl}/h$,
in $\mathbb{Q}(i)$, and put $H'\coloneqq H-h\,e^{\dagger}e$, whose entries
$H_{il}-H_{ip}H_{pl}/h$ lie in $\mathbb{Q}(i)$. For a vector $x$ put
$c\coloneqq ex$ and $y\coloneqq x-c\,\delta_p$, $\delta_p$ the $p$-th
coordinate vector; as $(Hx)_p=hc$,
$y^{\dagger}Hy=x^{\dagger}Hx-2h\lvert c\rvert^{2}+h\lvert c\rvert^{2}
=x^{\dagger}H'x$, so $H'$ is positive semidefinite. Row $p$ of $H'$ is zero,
and so is every row of $H'$ whose row in $H$ is zero, its entry $H_{jp}$
being zero; so $H'$ has fewer non-zero diagonal entries than $H$. Repeating
the step on $H'$ ends after at most $|\Sigma|$ steps at the zero matrix, and
$H$ is the sum of the terms $h\,e^{\dagger}e$ formed, which is $E^{\dagger}DE$
with the rows $e$ stacked into $E$ and the pivots $h$ on the diagonal of $D$.
In every case of C1--C9 the sum is $\id_\Sigma$. A matrix $D^{\flat}$ has no
row exactly when every diagonal entry of $D$ is zero, and a product with it
on the left then has none either; so $\Theta$ has no row exactly when
$y_1=y_2=0$ for C1 and C2, $y=0$ for C3, $R_H=R_V=1$ for C5, every
$\eta_m=1$ for C7 and $t=t_\perp=1$ for C8, and always for C4, C6 and C9:
the subfamilies of Table~\ref{tab:catalogue}.

\emph{(ii).} \emph{Each layer keeps the norm.} Fix $j$ and split
$\mathbb{C}^{\Modes}=\mathbb{C}^{\Sigma}\oplus\mathbb{C}^{\Modes\setminus\Sigma}$
with $\Sigma\coloneqq\Sigma_{e_j}$. The \termref{term-global-embedding}{global embedding} is
$\hat\Gmat_{e_j}=\Gmat_\Sigma\oplus\id$ (Definition~\ref{def:embed}) and
$\hat\Theta_j=\big(\Theta_j\ \ 0\big)$, so
$\hat\Gmat_{e_j}^{\dagger}\hat\Gmat_{e_j}+\hat\Theta_j^{\dagger}\hat\Theta_j
=\big(\Gmat_\Sigma^{\dagger}\Gmat_\Sigma+\Theta_j^{\dagger}\Theta_j\big)\oplus\id
=\id_{\Modes}$ by (i). The terminal layer
$D_\eta=\sum_k\sqrt{\eta_k}\,\Emd{\{k\}}+\Emd{\Modes\setminus\bigcup_k\Port{k}}$
(Definition~\ref{def:compile}) is real and diagonal, the ports being pairwise
disjoint (Definition~\ref{def:instrument}\,(iii)), and
$D_\eta^{2}+\Theta_\eta^{\dagger}\Theta_\eta
=\sum_k\big(\eta_k+(1-\eta_k)\big)\Emd{\{k\}}+\Emd{\Modes\setminus\bigcup_k\Port{k}}
=\id_{\Modes}$.
\emph{Telescoping.} The first identity gives
$X_j^{\dagger}X_j=X_{j-1}^{\dagger}\hat\Gmat_{e_j}^{\dagger}\hat\Gmat_{e_j}X_{j-1}
=X_{j-1}^{\dagger}X_{j-1}-(\hat\Theta_jX_{j-1})^{\dagger}(\hat\Theta_jX_{j-1})$,
and the blocks of rows of $\Omega$ lying on pairwise disjoint mode sets,
\begin{equation*}
  \Omega^{\dagger}\Omega
  =X_L^{\dagger}\big(D_\eta^{2}+\Theta_\eta^{\dagger}\Theta_\eta\big)X_L
   +\sum_{j=1}^{L}(\hat\Theta_jX_{j-1})^{\dagger}(\hat\Theta_jX_{j-1})
  =X_L^{\dagger}X_L+\sum_{j=1}^{L}\big(X_{j-1}^{\dagger}X_{j-1}-X_j^{\dagger}X_j\big)
  =X_0^{\dagger}X_0=\iota^{\dagger}\iota=\id_{\Inp}.
\end{equation*}
\emph{The ports.} $X_L=\hat\Gmat_{e_L}\cdots\hat\Gmat_{e_1}\iota
=\Gmat^{\sigma}_{\Grph}\iota$ by \eqref{eq:compile}, so the rows of $\Omega$
on $\Modes$ are $D_\eta\Gmat^{\sigma}_{\Grph}\iota=\Gmat\iota=\Gmat_{\Inp}$,
the columns $\Inp$ of $\Gmat$ (Definition~\ref{def:instrument}). Every port
$\Port{k}$ lies in $\Modes$, so the rows of $\Omega$ on it are those of
$\Emd{\{k\}}\Gmat_{\Inp}$, and
$\Omega^{\dagger}\Emd{\Port{k}}\Omega=\Gmat_{\Inp}^{\dagger}\Emd{\{k\}}\Gmat_{\Inp}=A_k$
by \eqref{eq:gramdef}. The ports are pairwise disjoint subsets of the mode
set of $\Omega$, so \eqref{eq:port-dilation} holds.
\end{proof}

% claim receiver/component-catalogue/def:c10-tolerance
\claimsummary{receiver/component-catalogue/def:c10-tolerance}{The contractivity of the generic block C10 is decided by the route-(ii) test of \hyperlink{claim-receiver/def:tolmap}{Definition~\ref*{def:tolmap}} (v), at the same widened threshold, and its lossless case, an equality, is decided two-sidedly within \ensuremath{\varepsilon}\ensuremath{_{1}}. Without the widening, rounding one unit above one would refuse a unitary block.}
\claimbookmark{receiver/component-catalogue/def:c10-tolerance}{Definition}{def:c10-tolerance}{Admission test of the generic block}
\claimleadin{%
The generic block is the one catalogue entry whose admission is a spectral
test on a supplied matrix, and that test is evaluated in floating point. It
is decided at $\varepsilon_1$, one of the \termref{term-receiver-tolerances}{receiver tolerances}, for the
same reason as at route (ii): a lossless block has
$\lmax(\Gmat_\Sigma^\dagger\Gmat_\Sigma)=1$ \emph{exactly} (the worked
instrument \eqref{eq:bb84G} of Example~\ref{ex:bb84} is unitary), so a
computed value one rounding unit above $1$ would refuse it.
}
\begin{definition}[Admission test of the generic block]
\label{def:c10-tolerance}
The (E1) predicate of catalogue entry C10 of
Definition~\ref{def:catalogue-entries} is decided by the route-(ii)
contractivity test of Definition~\ref{def:tolmap}\,(v) applied to the supplied
block $\Gmat_\Sigma$: the block is admitted when
$\lmax(\Gmat_\Sigma^\dagger\Gmat_\Sigma)\le1+\varepsilon_1$. C10 is the one
entry whose (E1) is a spectral condition rather than a parameter-domain
comparison (Theorem~\ref{thm:catalogue-admissibility}\,(iii)). Its lossless
sub-case $\Gmat_\Sigma^\dagger\Gmat_\Sigma=\id$ is an \emph{equality} and is
decided two-sidedly, as
$\lVert\Gmat_\Sigma^\dagger\Gmat_\Sigma-\id\rVert\le\varepsilon_1$.
\end{definition}
\claimmeta{\metalabel{Uses}\mref{Def.}{def:tolmap}, \mref{Thm.}{thm:catalogue-admissibility}
\metasep newly derived
\metasep accepted}

% claim receiver/component-catalogue/rem:rotator
\claimsummary{receiver/component-catalogue/rem:rotator}{Every retarder of the catalogue is generated in the \ensuremath{\sigma}z\textendash{}\ensuremath{\sigma}x plane, while the rotator is generated by \ensuremath{\sigma}y, so no single retarder is a non-trivial rotation. A chain of two half-wave plates is one, and the rotator is the one-element primitive.}
\claimbookmark{receiver/component-catalogue/rem:rotator}{Remark}{rem:rotator}{The rotator outside the retarder family}
\begin{remark}[The rotator outside the retarder family]\label{rem:rotator}
The rotator C4 of Definition~\ref{def:catalogue-entries} is \emph{not} a
special case of the retarder C3. Write $\sigma_x,\sigma_y,\sigma_z$ for the
Pauli matrices in the $(H,V)$ basis. The scalar of \eqref{eq:retarder}
cancels exactly,
$e^{-i\delta/2}\operatorname{diag}(1,e^{i\delta})=e^{-i(\delta/2)\sigma_z}$,
so that, using
$\Rot(\theta)\sigma_z\Rot(\theta)^{\transpose}
=\cos2\theta\,\sigma_z+\sin2\theta\,\sigma_x$,
\begin{equation}\label{eq:retgen}
  \Wplate(\delta,\theta)
  =\Rot(\theta)\,e^{-i(\delta/2)\sigma_z}\,\Rot(\theta)^{\transpose}
  =\exp\!\Big(-\tfrac{i\delta}{2}
    \big(\cos2\theta\,\sigma_z+\sin2\theta\,\sigma_x\big)\Big),
\end{equation}
an identity, not an equality up to a phase. Its generator lies in the
$\sigma_z$--$\sigma_x$ plane, while $\Rot(\rho)=\exp(-i\rho\sigma_y)$ generates
about the missing axis. So no single C3 element equals $\Rot(\rho)$ for
$\rho\in(0,2\pi)$: a lossy one is not unitary, and
$\exp\big(-\tfrac{i\delta}{2}\vec n\cdot\vec\sigma\big)$ with $\vec n$ in the
$z$--$x$ plane equals $\exp(-i\rho\sigma_y)$ only where both are real
multiples of $\id$, that is $\sin\rho=0$ and $\cos\tfrac\delta2=\cos\rho$,
which leaves $\rho=\pi$ with $\delta=2\pi$, outside the domain of C3. A chain
of two half-wave plates does equal a rotation: with
$M(a)\coloneqq\big(\begin{smallmatrix}\cos a&\sin a\\
\sin a&-\cos a\end{smallmatrix}\big)$, \eqref{eq:hwpqwp} reads
$\mathrm{HWP}(\theta)=-i\,M(2\theta)$, and $M(a)M(b)=\Rot(a-b)$, so
$\mathrm{HWP}(t_2)\,\mathrm{HWP}(t_1)=-\Rot\big(2(t_2-t_1)\big)
=\Rot\big(2(t_2-t_1)+\pi\big)$, which reaches every $\Rot(\rho)$. C4 is the
one-element primitive.
\end{remark}
\claimmeta{\metalabel{Uses}\mref{Thm.}{thm:catalogue-admissibility}
\metasep newly derived
\metasep accepted}

% claim receiver/component-catalogue/rem:entry-roles
\claimsummary{receiver/component-catalogue/rem:entry-roles}{The mirror's handedness flip is a choice of frame, the phase shifter is the interferometer's control knob, and fibre, coupling, filter and connector losses are loss entries at different sites. Route (ii) is the special case of route (iii) with one generic block spanning the receiver.}
\claimbookmark{receiver/component-catalogue/rem:entry-roles}{Remark}{rem:entry-roles}{Physical roles of the catalogue entries}
\begin{remark}[Physical roles of the catalogue entries]\label{rem:entry-roles}
Read physically, the entries of Definition~\ref{def:catalogue-entries} play
the following roles.
\begin{itemize}[leftmargin=1.6em]
\item \emph{C2.} Entry C2 is the polarization-dependent beamsplitter of
  entry C1 with its per-output losses, written in the mode order
  $(p_1H,p_1V,p_2H,p_2V)$; its ideal preset $t_H=1$, $t_V=0$ is the extreme
  case, the polarizing beamsplitter.
\item \emph{C3.} The preset $\mathrm{HWP}(\pi/8)$ is the $22.5^{\circ}$
  plate that maps the rectilinear basis to the diagonal one.
\item \emph{C5.} The handedness flip is a declared choice of coordinate frame,
  not a physical claim.
\item \emph{C6.} This is the interferometric receiver's control knob.
\item \emph{C7.} This is the loss primitive: fibre attenuation, coupling
  efficiency, spectral-filter insertion loss and connector loss are all C7 at
  different \emph{sites}, and the sites are what Theorem~\ref{thm:efficiency} is
  about.
\item \emph{C8.} It is contractive and not unitary unless $t=t_\perp=1$: the
  blocked polarization is loss.
\item \emph{C10.} This entry makes route (ii) a special case of route (iii).
  A user with a measured \termref{term-instrument-matrix}{instrument matrix}, or with a component the
  catalogue does not have, declares one block, and a route-(ii) declaration
  is recovered as the route-(iii) declaration of exactly one C10 block
  spanning the whole receiver.
\end{itemize}
\end{remark}
\claimmeta{\metalabel{Uses}\mref{Thm.}{thm:catalogue-admissibility}
\metasep newly derived
\metasep accepted}

% claim receiver/component-catalogue/rem:layers
\claimsummary{receiver/component-catalogue/rem:layers}{Detector efficiency with dark counts, and the probe source, are declared layers rather than catalogue entries. Efficiency enters as the terminal factor, dark counts act on the outcome label, and source modes, unlike vacuum inputs, are never eliminated.}
\claimbookmark{receiver/component-catalogue/rem:layers}{Remark}{rem:layers}{Detector and source layers outside the catalogue}
\begin{remark}[Detector and source layers outside the catalogue]\label{rem:layers}
Two further declared objects are \emph{layers}, not chain elements, and are
therefore not rows of Table~\ref{tab:catalogue} of
Definition~\ref{def:catalogue-entries}.

\emph{L1 --- Detector layer.} A port map $\{\Port{k}\}_{k=1}^{K}$ with
per-port efficiency $\eta_k\in[0,1]$ and dark-count probability
$p_{\dc|k}\in[0,1]$. The efficiency enters as the
\termref{term-terminal-detector-efficiency-layer}{terminal factor} $D_\eta$ of
Definition~\ref{def:compile}; the dark counts never enter $\Gmat$ at all and
act on the outcome label (Theorem~\ref{thm:dark-counts}). \emph{Detector} efficiency
is \emph{not} C7: it is the terminal layer, and the schema keeps the two apart
(Theorem~\ref{thm:efficiency}\,(iii) and Requirement~\ref{req:efficiency-schema}).

\emph{L2 --- Source (probe) layer.} The input split
$\Inp=\InpC\sqcup\InpS$ and the \termref{term-probe-declaration}{probe declaration} of
Definition~\ref{def:probe}. Undeclared input modes carry vacuum and may be
eliminated (Theorem~\ref{thm:port-grams}\,(iv)); \termref{term-declared-source-modes}{declared source modes} carry a
declared state and \emph{may not}. That difference is the entire structural
distinction between ``loss'' and ``probe'', and it is a definition rather than
a convention (Remark~\ref{rem:lossvsprobe}).
\end{remark}
\claimmeta{\metalabel{Uses}\mref{Def.}{def:probe}, \mref{Def.}{def:compile}
\metasep \metalabel{Assumes}\aref{A-VAC}{ass:vac}
\metasep newly derived
\metasep accepted}

% claim receiver/component-catalogue/rem:bsconv
\claimsummary{receiver/component-catalogue/rem:bsconv}{The symmetric beamsplitter differs from the catalogue's real one by a phase on each side. The two conventions give the same port Grams whenever the element is terminal with each output path inside one port and the input-side phase is a scalar on the declared input space.}
\claimbookmark{receiver/component-catalogue/rem:bsconv}{Remark}{rem:bsconv}{Beamsplitter conventions and the port Grams}
\claimleadin{%
A reader who imports a beamsplitter matrix written in another convention needs
to know which of its phase factors is free. The catalogue fixes one
convention, the core $\Bsp(x,\varphi)$ of \eqref{eq:bs}, and the interface
transcribes it (\href{../../../spec/modules/receiver/interface.yaml}{\nolinkurl{spec/modules/receiver/interface.yaml}}).
}
\begin{remark}[Beamsplitter conventions and the port Grams]\label{rem:bsconv}
The ``symmetric'' beamsplitter convention
$\Bsp_{\mathrm{sym}}(x)\coloneqq
\left(\begin{smallmatrix}\sqrt x&i\sqrt{1-x}\\
i\sqrt{1-x}&\sqrt x\end{smallmatrix}\right)$ differs from $\Bsp(x,0)$ by
phases on \emph{both} sides, and it matters which side each one is on. Writing
$\Upsilon\coloneqq\operatorname{diag}(1,i)$ on the ordered path pair, multiplication
gives
\begin{equation}\label{eq:bssym}
  \Upsilon\,\Bsp(x,0)\,\Upsilon
  =\begin{pmatrix}1&0\\0&i\end{pmatrix}
   \begin{pmatrix}\sqrt x&\sqrt{1-x}\\ \sqrt{1-x}&-\sqrt x\end{pmatrix}
   \begin{pmatrix}1&0\\0&i\end{pmatrix}
  =\begin{pmatrix}\sqrt x& i\sqrt{1-x}\\ i\sqrt{1-x}&\sqrt x\end{pmatrix}
  =\Bsp_{\mathrm{sym}}(x) ,
\end{equation}
the sign of the $(2,2)$ entry being restored by $i\cdot i=-1$. The two factors
have different status.
\begin{itemize}[leftmargin=1.6em]
\item The \emph{left} factor is an output-side phase. If the element is
  terminal and each of its two output paths lies inside one detector port,
  Corollary~\ref{cor:phases} applies and it changes no port Gram.
\item The \emph{right} factor is an \emph{input}-side phase, which
  Corollary~\ref{cor:phases} does not cover. Consider the case in which the
  element acts directly on the inputs, each of its two input paths being
  either \termref{term-declared-input-modes}{declared input modes} or an undeclared vacuum port
  (Assumption~\ref{ass:vac}), so that the factor acts on
  $\mathbb{C}^{\Inp}$ as a fixed diagonal unitary $\widetilde\Upsilon$,
  the vacuum columns being eliminated exactly
  (Theorem~\ref{thm:port-grams}\,(iv)). (With optics upstream of the element the
  pull-back of a diagonal phase need be neither diagonal nor unitary on
  $\mathbb{C}^{\Inp}$, and the comparison that follows is not the right one to
  make; the conclusion, to convert before declaring, still holds, and more
  strongly.) The two conventions then give
  $A_k^{\mathrm{sym}}=\widetilde\Upsilon^{\dagger}A_k\widetilde\Upsilon$ for
  every $k$: the \emph{same} family only up to that conjugation. The
  conjugation is the identity map exactly when $\widetilde\Upsilon$ is a
  scalar on $\mathbb{C}^{\Inp}$, that is, when every declared input mode
  receives the same phase from the factor, as when the element's one declared
  input path is the whole declared input space and the other input path is
  the vacuum port: a global phase on all of $\mathbb{C}^{\Inp}$ cancels in
  $A_k=\Gmat_{\Inp}^\dagger\Emd{\{k\}}\Gmat_{\Inp}$. When declared input modes
  receive different phases, as when both input paths are declared,
  $\widetilde\Upsilon$ is a genuine relative phase between them, and the two
  conventions can describe \emph{different} devices, whose Grams differ by
  that conjugation.
\end{itemize}
So the two conventions agree on the port Grams, and hence on the POVM,
whenever both conditions hold: the element is terminal with each of its
output paths inside one detector port, and the input-side factor is a scalar
on the declared input space.
Neither is automatic, which is why the catalogue fixes the one convention
$\Bsp(x,\varphi)$ of \eqref{eq:bs}, and why no catalogue entry drops a phase
it could have absorbed.
\end{remark}
\claimmeta{\metalabel{Assumes}\aref{A-VAC}{ass:vac}
\metasep newly derived
\metasep accepted}

% claim receiver/component-catalogue/rem:out
\claimsummary{receiver/component-catalogue/rem:out}{Active elements, modulators and switches driven by a party's random choice, cavities and feedback paths, and mode-shape-dependent effects are outside the catalogue, each for a stated reason. A spectral filter reduced to a scalar transmittance is a loss entry and is in scope.}
\claimbookmark{receiver/component-catalogue/rem:out}{Remark}{rem:out}{Elements outside the catalogue}
\claimleadin{%
Each exclusion from the catalogue follows from an assumption of the receiver
or from where the protocol declaration puts classical data. The interface records that placement in its
test-statistics model (\href{../../../spec/modules/receiver/interface.yaml}{\nolinkurl{spec/modules/receiver/interface.yaml}}).
}
\begin{remark}[Elements outside the catalogue]\label{rem:out}
Squeezers, parametric amplifiers and any active element are excluded by
Assumption~\ref{ass:passive}: their Bogoliubov map mixes $a$ with
$a^\dagger$, so $\Gmat^\dagger\Gmat\preceq\id$ fails, total photon number is
not conserved, and the photon-number block structure on which the receiver's
analysis rests does not exist. Modulators and optical switches driven by a
party's random choice are excluded because a classical choice belongs in the
protocol declaration as a weight, not in the mode algebra: the receiver
states its measurement per setting (Definition~\ref{def:setting-family}), and
the weights of the settings enter only downstream, as the class weights of
\xref{reductions/def:joint-family}. Cavities and feedback paths are excluded by
condition (E2) of Definition~\ref{def:graph}. Anything mode-shape-dependent is outside the
declared-mode model, in which each label of Definition~\ref{def:modes} is one
mode: dispersion and frequency conversion would act on the reserved temporal
and spectral index, which Definition~\ref{def:modes} fixes to a single value
that no element reads, and a mode-profile mismatch would carry light into
modes the declaration does not contain, which Assumption~\ref{ass:modes}
excludes; a spectral filter reduced to a scalar transmittance is C7 and is in
scope.
\end{remark}
\claimmeta{\metalabel{Uses}\mref{Def.}{def:graph}
\metasep \metalabel{Assumes}\aref{A-PASSIVE}{ass:passive}, \aref{A-THRESH}{ass:thresh}, \aref{A-MODES}{ass:modes}
\metasep newly derived
\metasep accepted}

\section{A worked BB84 receiver}\label{group-receiver/component-catalogue/ex:bb84}
\sectionintro{A four-detector BB84 receiver with a probe path is declared from three \termref{term-catalogue-entries}{catalogue entries}, compiled, and reduced to its channel port Grams (\hyperlink{claim-receiver/component-catalogue/ex:bb84}{Example~\ref*{ex:bb84}}). Its analytic multiclick weight is computed in closed form, and an exact window of depth two, evaluated in floating point, reaches at cutoff one, less the declared \termref{term-eigenvalue-slack}{eigenvalue slack}, the weight the analytic bound reaches only at cutoff two.}

% claim receiver/component-catalogue/ex:bb84
\claimsummary{receiver/component-catalogue/ex:bb84}{A four-detector BB84 receiver with a probe path, declared from three catalogue entries, compiles to a unitary instrument matrix with rank-one channel port Grams. Its dropped wave-plate phase is invisible at the ports, and eliminating the vacuum probe inputs is exact.}
\claimbookmark{receiver/component-catalogue/ex:bb84}{Example}{ex:bb84}{A BB84 receiver compiled from the catalogue}
\claimleadin{%
A four-detector BB84 receiver with a probe path can be declared from three
\termref{term-catalogue-entries}{catalogue entries}. The example shows that the catalogue and its conventions
express such a receiver, exhibits a case where a dropped wave-plate phase is
harmless \emph{for the reason Corollary~\ref{cor:phases} gives}, and supplies
the numbers used in Example~\ref{ex:bb84bound}.
}
\begin{example}[A BB84 receiver compiled from the catalogue]\label{ex:bb84}
Declare eight modes: a channel path $\Chan$ with $(\Chan_H,\Chan_V)$, a probe
path $\Prb$ with $(\Prb_H,\Prb_V)$, a $Z$ arm with $(Z_H,Z_V)$ and an $X$ arm
with $(X_H,X_V)$. Declare $\InpC=\{\Chan_H,\Chan_V\}$,
$\InpS=\{\Prb_H,\Prb_V\}$, and four single-mode ports
$\Port{H}=\{Z_H\}$, $\Port{V}=\{Z_V\}$, $\Port{+}=\{X_H\}$,
$\Port{-}=\{X_V\}$, with unit efficiencies
$\eta_H=\eta_V=\eta_+=\eta_-=1$ on the four ports, so that the
\termref{term-terminal-detector-efficiency-layer}{terminal detector-efficiency layer} $D_\eta$ of Definition~\ref{def:compile} is the
identity and the \termref{term-compiled-instrument-matrix}{compiled instrument matrix}
$\Gmat=D_\eta\Gmat_\Grph$ is $\Gmat_\Grph$ itself. Declare the chain
\begin{enumerate}[label=(\arabic*),leftmargin=2.2em]
\item C1 with $x=p_Z$, $\varphi=0$, $y_1=y_2=0$, on the ordered path pair
  $(\Chan,\Prb)$, so that (per polarization) the first output, on the
  $\Chan$ path, carries $\sqrt{p_Z}\,\Chan+\sqrt{p_X}\,\Prb$ and the second,
  on the $\Prb$ path, carries $\sqrt{p_X}\,\Chan-\sqrt{p_Z}\,\Prb$, with
  $p_X=1-p_Z$;
\item C9 on all eight modes, the permutation exchanging $\Chan_H,\Chan_V$
  with $Z_H,Z_V$ and $\Prb_H,\Prb_V$ with $X_H,X_V$, so that
  $Z\leftarrow\sqrt{p_Z}\,\Chan+\sqrt{p_X}\,\Prb$ and
  $X\leftarrow\sqrt{p_X}\,\Chan-\sqrt{p_Z}\,\Prb$;
\item C3 preset $\mathrm{HWP}(\pi/8)$ on the $X$ arm only.
\end{enumerate}
Every two of these elements whose supports intersect are ordered by the
declared chain ((1) and (2) share the $\Chan$ and $\Prb$ modes, (2) and (3)
the $X$ arm), so (E3) holds; (E1) holds because the three declared
entries carry in-domain parameters, and (E2) because the declaration is a
chain and a chain has no cycle. With (E1), (E2) and (E3) the declaration is
\termref{term-well-formed}{well formed}, so by Corollary~\ref{cor:compiled-instrument} this declaration is
a \termref{term-receiver-instrument}{receiver instrument}, and the results stated of a receiver
instrument apply to it. Compiling by \eqref{eq:compile} and reading
off $\Gmat_{\Inp}$ with rows ordered $(H,V,+,-)$ and columns ordered
$(\Chan_H,\Chan_V,\Prb_H,\Prb_V)$:
\begin{equation}\label{eq:bb84G}
  \Gmat_{\Inp}=
  \begin{pmatrix}
    \sqrt{p_Z} & 0 & \sqrt{p_X} & 0\\
    0 & \sqrt{p_Z} & 0 & \sqrt{p_X}\\
    -i\sqrt{\tfrac{p_X}{2}} & -i\sqrt{\tfrac{p_X}{2}} &
      \phantom{-}i\sqrt{\tfrac{p_Z}{2}} & \phantom{-}i\sqrt{\tfrac{p_Z}{2}}\\
    -i\sqrt{\tfrac{p_X}{2}} & \phantom{-}i\sqrt{\tfrac{p_X}{2}} &
      \phantom{-}i\sqrt{\tfrac{p_Z}{2}} & -i\sqrt{\tfrac{p_Z}{2}}
  \end{pmatrix},
\end{equation}
which is unitary: each column has squared norm $p_Z+2\cdot\frac{p_X}{2}=1$ or
$p_X+2\cdot\frac{p_Z}{2}=1$, and the six pairwise inner products vanish
(for instance
$\langle\Chan_H,\Chan_V\rangle=\frac{p_X}{2}-\frac{p_X}{2}=0$ and
$\langle\Chan_H,\Prb_H\rangle=\sqrt{p_Zp_X}-2\cdot\frac{\sqrt{p_Zp_X}}{2}=0$).

\emph{The phase.} Rows $3$ and $4$ of \eqref{eq:bb84G} carry the overall factor
$-i$: it is the physical retarder phase $e^{-i\delta/2}$ of
\eqref{eq:retarder} at $\delta=\pi$, which an implementation that writes
$\mathrm{HWP}$ as the real matrix
$\left(\begin{smallmatrix}\cos2\theta&\sin2\theta\\
\sin2\theta&-\cos2\theta\end{smallmatrix}\right)$ silently drops. Here that
factor is a \emph{per-port} phase --- both affected rows belong to ports that
are never recombined --- so by Theorem~\ref{thm:representation-independence}\,(iii) the port Grams, and
hence the POVM, are identical either way. That is a statement about this
receiver and this factor, not about phases in general.

\emph{The channel-only receiver.} With the probe off (vacuum on $\InpS$),
Lemma~\ref{lem:vacelim} lets the $\Prb$ columns be dropped: the channel-only
instrument is the left $4\times2$ block of \eqref{eq:bb84G}, and the
channel port Grams are
\begin{equation}\label{eq:bb84grams}
  A_H=p_Z\dyad{H},\quad A_V=p_Z\dyad{V},\quad
  A_+=p_X\dyad{+},\quad A_-=p_X\dyad{-},
\end{equation}
each of rank one --- so by Theorem~\ref{thm:representation-independence}\,(iv) this receiver \emph{is}
exactly a one-mode-per-port device, and
$A_\Sigma=p_Z(\dyad{H}+\dyad{V})+p_X(\dyad{+}+\dyad{-})=\id_2$: lossless, as
the construction promises. Dropping the $\Prb$ columns is the vacuum
elimination of Theorem~\ref{thm:port-grams}\,(iv), not a second model.
\end{example}
\claimmeta{\metalabel{Uses}\mref{Lem.}{lem:vacelim}, \mref{Def.}{def:compile}, \mref{Cor.}{cor:compiled-instrument}, \mref{Thm.}{thm:port-grams}, \mref{Thm.}{thm:representation-independence}, \mref{Cor.}{cor:phases}, \mref{Thm.}{thm:catalogue-admissibility}
\metasep newly derived
\metasep accepted}

% claim receiver/component-catalogue/ex:bb84bound
\claimsummary{receiver/component-catalogue/ex:bb84bound}{The channel-only BB84 receiver has a closed-form analytic multiclick bound, and its one setting is exact, so the emitted weight is the positive part of the bound at the first block above the cutoff, with no tail window. At balanced bases the cutoffs N = 2 and N = 3 give 1/2 and 3/4.}
\claimbookmark{receiver/component-catalogue/ex:bb84bound}{Example}{ex:bb84bound}{Analytic multiclick weight of the BB84 receiver}
\claimleadin{%
The rank-one Grams of the BB84 receiver make every quantity of the analytic
multiclick bound explicit, and its single setting shows the collapse branch at
work.
}
\begin{example}[Analytic multiclick weight of the BB84 receiver]\label{ex:bb84bound}
Take the \emph{channel-only} receiver of the last paragraph of
Example~\ref{ex:bb84} (the probe off and its $\Prb$ columns eliminated,
Lemma~\ref{lem:vacelim}), which declares $\InpS=\varnothing$ and therefore
has $\Tset=\{\off\}$ by Definition~\ref{def:probe}. Then $\gamma=\beta_k=1$
and the Grams are \eqref{eq:bb84grams} with $A_\Sigma=\id_2$, so
$c=1-\lmax(A_\Sigma)=0$, and
\begin{equation}
  A_\Sigma-A_H=\id-p_Z\dyad{H}\ \Rightarrow\ \lmin=1-p_Z=p_X\ \Rightarrow\
  d_H=p_Z,
\end{equation}
and likewise $d_V=p_Z$, $d_+=d_-=p_X$. Hence, for $n\ge1$,
\begin{equation}\label{eq:bb84beta}
  \Bnd(n)=1-2p_Z^{\,n}-2p_X^{\,n}.
\end{equation}
This receiver declares no probe modes, so its one setting is flagged
$\mathrm{exact}$ (Corollary~\ref{cor:unconditioned-family}) and predicate~(X)
of Definition~\ref{def:pmc} holds, \eqref{eq:collapsehyp}
holds with equality ($\gamma=\beta_k=1$), branch~1 is taken, and no tail
window is declared or required (Remark~\ref{rem:flag-instances},
Requirement~\ref{req:window-declaration}). Since $\Tset=\{\off\}$, the minimum over
$\Tset$ in Definition~\ref{def:pmc} is over that one setting, so what is
computed here is the emitted $\pmc$ and not one of several per-setting values.
Hence $\pmc=\max\{0,\Bnd(N+1)\}$ exactly. At $p_Z=p_X=\tfrac12$ this is
$\Bnd(n)=1-2^{2-n}$: $\Bnd(1)=-1$, $\Bnd(2)=0$, $\Bnd(3)=\tfrac12$,
$\Bnd(4)=\tfrac34$, non-decreasing from the first non-negative index, as
Lemma~\ref{lem:collapse} requires. A cutoff $N=2$ therefore yields
$\pmc=\tfrac12$, and $N=3$ yields $\pmc=\tfrac34$. Note that $\pmc\to1$
exponentially exactly when $\max(p_Z,p_X)<1$, i.e.\ when no single basis
absorbs the whole input.
\end{example}
\claimmeta{\metalabel{Uses}\mref{Def.}{def:probe}, \mref{Lem.}{lem:vacelim}, \mref{Def.}{def:compile}, \mref{Lem.}{lem:collapse}, \mref{Def.}{def:pmc}
\metasep newly derived
\metasep accepted}

% claim receiver/component-catalogue/ex:bb84exact
\claimsummary{receiver/component-catalogue/ex:bb84exact}{The exact two-photon multiclick block of the channel-only BB84 receiver shows the slack in the analytic bound. An exact window of depth two reaches at cutoff one, less the declared eigenvalue slack, the weight the analytic bound reaches only at cutoff two.}
\claimbookmark{receiver/component-catalogue/ex:bb84exact}{Example}{ex:bb84exact}{Exact window on the BB84 receiver}
\claimleadin{%
An exact window takes, on the first blocks above the cutoff, the larger of
the analytic bound and a computed block eigenvalue less the
\termref{term-eigenvalue-slack}{eigenvalue slack}. On the BB84 receiver the two-photon block can be
diagonalised by hand, which shows what the window buys.
}
\begin{example}[Exact window on the BB84 receiver]\label{ex:bb84exact}
Continue with the channel-only receiver of Example~\ref{ex:bb84bound}: $K=4$,
$\gamma=\beta_k=1$, $\InpC=\{\Chan_H,\Chan_V\}$ so $D=|\InpC|=2$, and the Grams
\eqref{eq:bb84grams} with $A_\Sigma=\id_2$. Here $B=A_\Sigma=\id_2$ and
$B_k=A_{[K]\setminus\{k\}}=\id-A_k$, so $\id-B=0$ and $\id-B_k=A_k$, and
\eqref{eq:exactbnd} reads, for $n\ge1$,
\begin{equation}\label{eq:bb84exactblock}
  \blk{M}{\eqn}_{\mc}=\id^{(=n)}-\sum_{k}\Fock(A_k)^{(n)}
  =\id^{(=n)}-p_Z^{\,n}\big(\dyad{n_H}+\dyad{n_V}\big)
   -p_X^{\,n}\big(\dyad{n_+}+\dyad{n_-}\big).
\end{equation}
Two steps: $\Fock(0)^{(n)}=0$ for $n\ge1$ kills the $\gamma$ term; and for a
rank-one $A_k=p\dyad{\psi}$, $(p\dyad{\psi})^{\otimes n}$ is $p^{\,n}$ times
the orthogonal projector onto $(\operatorname{span}\ket{\psi})^{\otimes n}$,
whose restriction to $\Symn{n}$ is $\dyad{n_\psi}$ with
$\ket{n_\psi}=\ket{\psi}^{\otimes n}$ the occupation vector carrying all $n$
photons in $\ket{\psi}$ (Lemma~\ref{lem:fock-functor}\,(iv) with
$P=\dyad{\psi}$).

\emph{The two-photon block.} The space $\Symn{2}(\mathbb{C}^2)$ has
dimension $\binom{2+2-1}{2}=3$ and occupation basis
$\{\ket{2_H},\ket{1_H1_V},\ket{2_V}\}$, so
$\dyad{2_H}+\dyad{2_V}=\id-\dyad{1_H1_V}$; likewise
$\dyad{2_+}+\dyad{2_-}=\id-\dyad{1_+1_-}$. Expanding
$\ket{\pm}=(\ket{H}\pm\ket{V})/\sqrt2$,
\begin{equation}\label{eq:bb84exactvec}
  \ket{1_+1_-}=\tfrac{1}{\sqrt2}\big(\ket{+}\otimes\ket{-}
   +\ket{-}\otimes\ket{+}\big)
  =\tfrac{1}{\sqrt2}\big(\ket{2_H}-\ket{2_V}\big),
\end{equation}
which is orthogonal to $\ket{1_H1_V}$. Substituting into
\eqref{eq:bb84exactblock} at $n=2$,
\begin{equation}\label{eq:bb84exact2}
  \blk{M}{(=2)}_{\mc}=\big(1-p_Z^2-p_X^2\big)\,\id
   +p_Z^2\,\dyad{1_H1_V}+p_X^2\,\dyad{1_+1_-},
\end{equation}
whose spectrum is therefore $\{1-p_X^2,\ 1-p_Z^2,\ 1-p_Z^2-p_X^2\}$, the
first two on the two displayed orthonormal vectors, the third on
$(\ket{2_H}+\ket{2_V})/\sqrt2$. So
\begin{equation}\label{eq:bb84exactvalue}
  \Bnd^{\ast}(2)=1-p_Z^2-p_X^2
  \qquad\text{against}\qquad
  \Bnd(2)=1-2p_Z^2-2p_X^2 ,
\end{equation}
a gap of $p_Z^2+p_X^2>0$ at every $(p_Z,p_X)$.

\emph{At the cutoff that matters.} Take $p_Z=p_X=\tfrac12$, where
$\Bnd(n)=1-2^{2-n}$ and $\Bnd^{\ast}(2)=\tfrac12$. At the cutoff $N=1$,
Definition~\ref{def:pmc} emits $\max\{0,\Bnd(2)\}=0$: the two-photon block is
discarded entirely, although half of its weight is guaranteed. Declare the
minimal \termref{term-exact-window-depth}{depth} $m^{\ast}=N+1=2$. Predicate~(X) holds (this receiver declares
$\InpS=\varnothing$), so \eqref{eq:pmcex} takes branch~1 with the remainder
$\Bnd(m^{\ast}+1)=\Bnd(3)=\tfrac12$, and, at an evaluation
$\widehat\Bnd^{\ast}(2)=\Bnd^{\ast}(2)$,
\begin{equation}\label{eq:bb84exactemit}
  \pmc=\max\big\{0,\ \min\big\{\ell(2),\ \Bnd(3)\big\}\big\}
   =\max\big\{0,\ \tfrac12-\varepsilon_{\mathrm{eig}}\big\} ,
\end{equation}
since $\ell(2)=\max\{\Bnd(2),\widehat\Bnd^{\ast}(2)
-\varepsilon_{\mathrm{eig}}\}=\max\{0,\tfrac12-\varepsilon_{\mathrm{eig}}\}$,
which is $\le\tfrac12=\Bnd(3)$ on the whole declared domain
$\varepsilon_{\mathrm{eig}}\in[0,1)$, so the minimum is $\ell(2)$ throughout.
The value is $\tfrac12-\varepsilon_{\mathrm{eig}}$ at
$\varepsilon_{\mathrm{eig}}\le\tfrac12$ and $0$ above it; the outer
$\max\{0,\cdot\}$ is carried through the chain rather than dropped, because
the declared domain reaches past $\tfrac12$.
The declared slack is subtracted once, from the computed block eigenvalue
$\widehat\Bnd^{\ast}(2)=\tfrac12$.

\emph{What that buys, and what it does not change.}
Example~\ref{ex:bb84bound} reaches $\pmc=\tfrac12$ only at the cutoff $N=2$;
the exact window reaches $\max\{0,\tfrac12-\varepsilon_{\mathrm{eig}}\}$ at
$N=1$, from a single $3\times3$ eigenvalue
problem, and $N$ is the parameter every downstream target dimension grows with.
Theorem~\ref{thm:multiclick-bound} is untouched: $\Bnd(3)=\tfrac12$ is still
the remainder covering every block above the second, and the analytic bound is
still what makes the quantifier over $n>N$ finite. Note also that the
declaration must supply $\varepsilon_{\mathrm{eig}}$ and $\varepsilon_6$ to
declare the depth at all, and that this receiver still declares no tail
window, predicate~(X) holding
(Remark~\ref{rem:flag-instances}).
\end{example}
\claimmeta{\metalabel{Uses}\mref{Lem.}{lem:fock-functor}, \mref{Def.}{def:compile}, \mref{Thm.}{thm:multiclick-bound}, \mref{Def.}{def:exactbnd}, \mref{Def.}{def:pmc}, \mref{Def.}{def:pmcex}
\metasep newly derived
\metasep accepted}

\clearpage
\part{Detector imperfections}
\sectionintro{This submodule adds detector efficiency and dark counts to the passive receiver, and its two main results say how each acts on the measurement. Efficiency uniform within a port scales that port's Gram and enters at the single terminal layer of the compiled matrix (\hyperlink{claim-receiver/detector-imperfections/thm:efficiency}{Theorem~\ref*{thm:efficiency}}); efficiency that varies within a port is not a scalar and can lower a Gram's rank. Dark counts are a classical channel on the recorded click pattern that adds clicks and never removes them. Such a channel preserves the POVM structure and improves the multiclick bound p\_mc to (1 \ensuremath{-} r\_min)p\_mc + r\_min, r\_min being the least favourable transition from every non-multiclick outcome, the no-click outcome included (\hyperlink{claim-receiver/detector-imperfections/thm:dark-counts}{Theorem~\ref*{thm:dark-counts}}). A declared channel is admitted at two receiver tolerances, and the bound a downstream module composes with carries their cost explicitly. The emitted blocks are those of the dark-composed family, enclosed rigorously for any declared parameters and exact where the enclosures have radius zero; loss and a dark count in every port make each of its elements positive definite on every block, and neither can be dropped (\hyperlink{claim-receiver/detector-imperfections/prop:noise-floor}{Proposition~\ref*{prop:noise-floor}}), which is the floor a model needs if no snap of its blocks is to revive a zero row. A receiver without a probe is a uniform loss of transmittance t followed by a normalized station whose efficiencies are divided by t, for any t from its largest efficiency \ensuremath{\eta}\_max to one, so that loss can be handed to the adversary with the observed statistics unchanged (\hyperlink{claim-receiver/detector-imperfections/prop:efficiency-normalization}{Proposition~\ref*{prop:efficiency-normalization}}); by default the engine compiles the normalized station at t = \ensuremath{\eta}\_max/(99/100), and the declared receiver (t = 1) where \ensuremath{\eta}\_max \ensuremath{\ge} 99/100.}

\section{Detector efficiency}\label{group-receiver/detector-imperfections/thm:efficiency}
\sectionintro{The main result, \hyperlink{claim-receiver/detector-imperfections/thm:efficiency}{Theorem~\ref*{thm:efficiency}}, is that detector efficiency is a parameter of the terminal layer, applied after every chain element, that multiplies each port Gram by its efficiency when uniform within the port. The remarks show a within-port attenuation that lowers a rank, why a product of transmittances is the wrong loss audit in an interferometer, and why power data and click data differ by exactly the efficiencies.}

% claim receiver/detector-imperfections/thm:efficiency
\claimsummary{receiver/detector-imperfections/thm:efficiency}{Efficiency \ensuremath{\eta}\_k uniform on detector port k scales that port's Gram by \ensuremath{\eta}\_k, keeping its rank if \ensuremath{\eta}\_k \textgreater{} 0; varying within a port, it is no scalar and can lower the rank. It enters only as the compiled instrument matrix's last layer, and the loss decomposition \hyperlink{claim-receiver/lem:deficiency}{Lemma~\ref*{lem:deficiency}} checks that loss is counted once.}
\claimbookmark{receiver/detector-imperfections/thm:efficiency}{Theorem}{thm:efficiency}{Detector efficiency in the port Grams}
\claimleadin{%
``Apply the efficiency once'' is easy to state for a fixed device and does not
survive verbatim once a receiver is compiled from components: a component graph
has many legitimate loss sites, among them every beamsplitter's loss, every C7
element, every polarizer's extinction, every mirror's reflectivity and the
detector efficiency itself. The theorem fixes where the detector efficiency
enters, what it does to the port Grams and what the emitted record of the loss
sites lists. What the engine keeps apart in its input and records in its output
is stated in Requirements~\ref{req:efficiency-schema}
and~\ref{req:efficiency-record}.
}
\begin{mainclaim}
\begin{theorem}[Detector efficiency in the port Grams]\label{thm:efficiency}
Let $\Recv$ be a \termref{term-receiver-instrument}{receiver instrument} with \termref{term-port-gram-matrices}{port
Grams} $\{A_k\}$ and an \termref{term-instrument-matrix}{instrument matrix} $\Gmat$ that does not yet
contain a detector-efficiency layer (for a \termref{term-declared-element-graph}{declared element graph}, the
product $\Gmat^{\sigma}_{\Grph}$ of Definition~\ref{def:compile}), let
$\eta_1,\dots,\eta_K\in[0,1]$, and let
$D_\eta=\sum_k\sqrt{\eta_k}\,\Emd{\{k\}}+\Emd{\Modes\setminus\bigcup_k\Port{k}}$
be the \termref{term-terminal-detector-efficiency-layer}{terminal detector-efficiency layer} of
Definition~\ref{def:compile}.
\begin{enumerate}[label=(\roman*),leftmargin=3.0em]
\item \emph{Uniform within a port: a scalar on the Gram.} $D_\eta\Gmat$ is a
  contraction, and its port Grams are
  \begin{equation}\label{eq:etagram}
    A_k\ \longmapsto\ \eta_k\,A_k .
  \end{equation}
  For $\eta_k>0$ this preserves $\rk A_k$, so it leaves the one-mode-per-port
  criterion of Theorem~\ref{thm:representation-independence}\,(iv) untouched. Every formula of
  Theorems~\ref{thm:click-povm}, \ref{thm:coarse-graining}, \ref{thm:multiclick-bound}
  and~\ref{thm:emitted-bound} therefore covers inefficient detectors under the
  single substitution \eqref{eq:etagram}; in particular
  $A_\Sigma\mapsto\sum_k\eta_kA_k$, which is in general not a multiple of
  $A_\Sigma$ unless the $\eta_k$ agree.
\item \emph{Non-uniform within a port: not a scalar, and the rank can drop.} If
  attenuation varies \emph{within} a port, that is, for a diagonal contraction
  $D$ that is not constant on $\Port{k}$, the effective Gram is
  \begin{equation}\label{eq:noneffgram}
    A_k^{\mathrm{eff}}=\Gmat_{\Inp}^\dagger D^\dagger\Emd{\{k\}}D\,
    \Gmat_{\Inp},
  \end{equation}
  which is in general \emph{not} a multiple of $A_k$ and satisfies only
  $\rk A_k^{\mathrm{eff}}\le\rk A_k$, with strict inequality possible. All
  formulas of the theorems named in (i) still apply, to $A_k^{\mathrm{eff}}$.
\item \emph{One structural site.} Detector efficiency is a parameter of the
  terminal layer L1 and enters the \termref{term-compiled-instrument-matrix}{compiled instrument matrix}
  $D_\eta\Gmat^{\sigma}_{\Grph}$ only through the single left factor
  $D_\eta$ of Definition~\ref{def:compile}, which is applied after every chain
  element. It is therefore not an audit but an invariant of the compiler: there
  is exactly one site at which it can enter.
\item \emph{The composition record.} The emitted efficiency record
  lists every declared loss parameter $q_i$, under its name, with the element
  that declares it and its position in the compiled product: the
  transmittance $x$ of a beamsplitter (C1) and $t_H,t_V$ of a polarizing
  beamsplitter (C2), which is also the entry of a
  \termref{term-polarization-dependent-beamsplitter}{polarization-dependent beamsplitter}, its $x_H,x_V$ recorded as
  $t_H,t_V$; the output losses $y_1,y_2$ of C1 and C2 and the loss $y$ of a
  retarder (C3), whose transmittances are $1-y_1$, $1-y_2$ and $1-y$; the
  reflectivities $R_H,R_V$ of a mirror (C5); the transmittances $\eta_m$ of a
  loss element (C7); the pass and leak transmittances $t,t_\perp$ of a
  polarizer (C8); and the detector efficiencies $\eta_k$. Two elements may
  carry equal numerical values; they are still two losses and appear as two
  rows.
\item \emph{The deficiency identity.} The identity \eqref{eq:deficiencyid} of
  Lemma~\ref{lem:deficiency} holds for every receiver instrument, before and
  after the layer $D_\eta$, and is the correct check that loss was accounted
  for exactly once: its two summands are the amplitude that left through the
  dilation and the amplitude that reached an unmonitored declared mode, and
  neither carries a separate efficiency term; after the layer, the efficiency
  loss is part of the dilation summand.
\end{enumerate}
\end{theorem}
\end{mainclaim}
\claimmeta{\metalabel{Uses}\mref{Prop.}{prop:mode-algebra}, \mref{Def.}{def:compile}, \mref{Lem.}{lem:deficiency}, \mref{Thm.}{thm:representation-independence}
\metasep \metalabel{Used by}\mref{Rem.}{rem:powervsclick}, \mref{Rem.}{rem:sperlingeta}, \mref{Prop.}{prop:noise-floor}, \mref{Def.}{def:normalized-station}, \mref{Rem.}{rem:normalization-scope}, \mref{Req.}{req:efficiency-schema}, \mref{Req.}{req:efficiency-record}
\metasep newly derived
\metasep implemented (effective: accepted)
\metasep \metalabel{Code}\coderef{https://github.com/fraunhofersg/certiqs-key/blob/v0.1.0/src/certiqs_key/receiver/efficiency.py}{certiqs_key.receiver.efficiency:efficiency_record}
\metasep \metalabel{Oracles}\coderef{https://github.com/fraunhofersg/certiqs-key/blob/v0.1.0/tests/oracles/registry_receiver_povm.py}{registry_receiver_povm.py}}

\begin{proof}
(i) $D_\eta$ is diagonal with entries in $[0,1]$, hence a positive semidefinite
contraction, and a product of contractions is a contraction
(Proposition~\ref{prop:mode-algebra}\,(ii)). The ports are pairwise disjoint, so
$\Emd{\{k\}}D_\eta=\sqrt{\eta_k}\,\Emd{\{k\}}$, whence
\begin{equation}
  (D_\eta\Gmat)^\dagger_{\Inp}\Emd{\{k\}}(D_\eta\Gmat)_{\Inp}
  =\Gmat_{\Inp}^\dagger D_\eta\Emd{\{k\}}D_\eta\Gmat_{\Inp}
  =\eta_k\,\Gmat_{\Inp}^\dagger\Emd{\{k\}}\Gmat_{\Inp}=\eta_kA_k .
\end{equation}
A strictly positive scalar multiple of a matrix has the same rank.

(ii) The displayed formula is the definition of the port Gram for the
instrument $D\Gmat$. For the rank, $\Emd{\{k\}}D\Gmat_{\Inp}
=D_k\big(\Emd{\{k\}}\Gmat_{\Inp}\big)$ where $D_k$ is the diagonal
restriction of $D$ to $\Port{k}$, and $\rk(D_kY)\le\rk Y$ for any $Y$; the
Gram of a matrix has the rank of that matrix. Strictness is exhibited in
Remark~\ref{rem:effrank}.

(iii) is read off Definition~\ref{def:compile}, in which $D_\eta$ is the one
factor that carries the detector efficiencies.

(iv) is a statement about the emitted artifact.

(v) is Lemma~\ref{lem:deficiency}, which uses only that the
\termref{term-instrument-matrix}{instrument matrix} is a contraction; $D_\eta\Gmat$ is one by (i).
\end{proof}

% claim receiver/detector-imperfections/rem:effrank
\claimsummary{receiver/detector-imperfections/rem:effrank}{A port of two modes with Gram diag(1, \ensuremath{\varepsilon}\ensuremath{^{2}}) has rank two; attenuating one of its modes to zero leaves rank one. Apart from a zero efficiency, within-port attenuation is the one operation in this contract that can change the answer to the one-mode-per-port rank test.}
\claimbookmark{receiver/detector-imperfections/rem:effrank}{Remark}{rem:effrank}{A port Gram whose rank drops under within-port attenuation}
\begin{remark}[A port Gram whose rank drops under within-port attenuation]\label{rem:effrank}
Take $|\Inp|=2$ and a port collecting two modes, with the two rows of
$\Emd{\{k\}}\Gmat_{\Inp}$ equal to $(1,0)$ and $(0,\varepsilon)$ for some
$\varepsilon\in(0,1]$. Then $A_k=\operatorname{diag}(1,\varepsilon^2)$ has rank
$2$. Attenuating the second of those modes to zero, by a diagonal $D$ with
entries $1$ and $0$ inside the same port, gives
$A_k^{\mathrm{eff}}=\operatorname{diag}(1,0)$, of rank $1$. So a device that
does \emph{not} satisfy the one-mode-per-port criterion of
Theorem~\ref{thm:representation-independence}\,(iv) satisfies it after a within-port attenuation, and
the rank test gives its answer only on the Grams after every efficiency is
applied, never before (Requirement~\ref{req:one-mode-per-port}). For
single-mode ports the situation never arises, which is why the layer L1
declares one $\eta_k$ per port (the uniform case of
Theorem~\ref{thm:efficiency}\,(i)) by construction. Apart from declaring some
$\eta_k=0$, which sends $A_k$ to $0$ by Theorem~\ref{thm:efficiency}\,(i), this
is the one operation in this contract that can change the answer to the rank
test (Remark~\ref{rem:rankone}).
\end{remark}
\claimmeta{\metalabel{Uses}\mref{Thm.}{thm:representation-independence}, \mref{Rem.}{rem:rankone}
\metasep newly derived
\metasep accepted}

% claim receiver/detector-imperfections/rem:wrongcheck
\claimsummary{receiver/detector-imperfections/rem:wrongcheck}{In a lossless balanced Mach\textendash{}Zehnder interferometer at zero phase, a product of transmittances along paths reports a spurious loss of one half at the first port, as it cannot see the interference cross term. The deficiency identity is the audit that holds, and the artifact reports it as a residual.}
\claimbookmark{receiver/detector-imperfections/rem:wrongcheck}{Remark}{rem:wrongcheck}{Failure of the path-product loss audit under interference}
\begin{remark}[Failure of the path-product loss audit under interference]
\label{rem:wrongcheck}
The obvious audit, multiplying the declared transmittances along each path and
comparing with the port throughput, is \emph{wrong} in any interferometric
topology, and the failure is not small. Consider a balanced Mach--Zehnder
interferometer built from the \termref{term-component-catalogue}{component catalogue}: C1 with
$x=\tfrac12,\varphi=0$ and no loss, then C6 with phase $\phi$ on the second
path, then C1 with $x=\tfrac12,\varphi=0$ again, with the two output paths
monitored by two ports, one per path. Every declared loss parameter is $0$
($y_1=y_2=0$ at both beamsplitters, and the phase shifter is unitary), and
\eqref{eq:deficiencyid} correctly returns $\defic=0$. The compiled matrix is,
on the two paths (each entry acting as that multiple of $\id_2$ on
polarization),
\begin{equation}\label{eq:mz}
  \Gmat=\Bsp(\tfrac12,0)\operatorname{diag}(1,e^{i\phi})\Bsp(\tfrac12,0)
  =\tfrac12\begin{pmatrix}1+e^{i\phi}&1-e^{i\phi}\\
                          1-e^{i\phi}&1+e^{i\phi}\end{pmatrix},
\end{equation}
so at $\phi=0$ it is the identity and \emph{all} the light entering path $1$
reaches port $1$. A path-product audit sees two paths from input $1$ to output
$1$, transmit--transmit and reflect--reflect, each of power weight
$\tfrac12\cdot\tfrac12=\tfrac14$, and concludes a throughput of $\tfrac12$,
reporting a spurious $50\%$ loss that does not exist. The missing $\tfrac12$ is
the interference cross term, which no product of transmittances can see. The
checkable invariant is \eqref{eq:deficiencyid}, and it is what the artifact
reports (residual $r_3$ of Definition~\ref{def:residuals}).
\end{remark}
\claimmeta{\metalabel{Uses}\mref{Lem.}{lem:deficiency}, \mref{Def.}{def:residuals}
\metasep newly derived
\metasep accepted}

% claim receiver/detector-imperfections/rem:powervsclick
\claimsummary{receiver/detector-imperfections/rem:powervsclick}{A power characterisation of the optics returns the port Grams A\_k, click statistics return \ensuremath{\eta}\_k A\_k. Double counting is avoided exactly when the efficiencies enter once, and which convention the declared numbers follow cannot be read off the numbers.}
\claimbookmark{receiver/detector-imperfections/rem:powervsclick}{Remark}{rem:powervsclick}{Power data and click data measure different objects}
\begin{remark}[Power data and click data measure different objects]
\label{rem:powervsclick}
A classical power characterisation of the optics measures intensities
independently of the single-photon detectors and returns $A_k$; detector click
statistics return $\eta_kA_k$. Double counting is avoided exactly when the
efficiencies enter through \eqref{eq:etagram} once: either the Grams are
declared from power data and then rescaled by $\eta_k$, or they are declared
from click data and the $\eta_k$ are already inside them and must be declared as
$1$. Which of the two a declaration means is not inferable from the numbers
(Requirement~\ref{req:efficiency-convention} says where it is recorded).
\emph{Characterising} the device is out of scope
(Assumption~\ref{ass:trust}); recording which convention the declared numbers
follow is not.
\end{remark}
\claimmeta{\metalabel{Uses}\mref{Thm.}{thm:efficiency}
\metasep \metalabel{Assumes}\aref{A-TRUST}{ass:trust}
\metasep \metalabel{Used by}\mref{Req.}{req:efficiency-convention}
\metasep newly derived
\metasep accepted}

% claim receiver/detector-imperfections/rem:sperlingeta
\claimsummary{receiver/detector-imperfections/rem:sperlingeta}{For one mode the efficiency substitution gives an on/off detector of efficiency \ensuremath{\eta} the normally ordered no-click element exp(\ensuremath{-}\ensuremath{\eta}n), the zero-count, noise-free case of a published photo-counting formula. The agreement corroborates and backs no step.}
\claimbookmark{receiver/detector-imperfections/rem:sperlingeta}{Remark}{rem:sperlingeta}{The single-mode efficiency substitution in the photo-counting formula}
\begin{remark}[The single-mode efficiency substitution in the photo-counting formula]\label{rem:sperlingeta}
For a single mode, \eqref{eq:etagram} together with
Proposition~\ref{prop:normord} says that an on/off detector of efficiency
$\eta$ has no-click element $\normord{e^{-\eta\hat n}}$ rather than
$\normord{e^{-\hat n}}$. The same substitution appears in
\cite[Eq.~(3)]{sperling2012}, whose photo-counting formula with finite quantum
efficiency $\eta$ and noise $\nu$ reads
$p_k=\langle\normord{(\eta\hat n+\nu)^ke^{-(\eta\hat n+\nu)}}\rangle/k!$ and
whose $k=0$, $\nu=0$ case is $\langle\normord{e^{-\eta\hat n}}\rangle$.
This is a corroboration only; no proof relies on it.
\end{remark}
\claimmeta{\metalabel{Uses}\mref{Prop.}{prop:normord}, \mref{Thm.}{thm:efficiency}
\metasep newly derived
\metasep accepted}

\section{Dark counts}\label{group-receiver/detector-imperfections/def:darkadm}
\sectionintro{A \termref{term-dark-count-channel}{dark-count channel} is column-stochastic and monotone. The main result, \hyperlink{claim-receiver/detector-imperfections/thm:dark-counts}{Theorem~\ref*{thm:dark-counts}}, is that at every setting it keeps the composed family a photon-number block-diagonal POVM, never moves weight out of the \termref{term-multiclick-set}{multiclick set}, and improves the multiclick bound; the minimum defining the improvement must run over the no-click column too. A declared channel is admitted at two \termref{term-receiver-tolerances}{receiver tolerances}, and the composed inequality that holds for every admitted channel carries both slacks (\hyperlink{claim-receiver/detector-imperfections/prop:dark-slack-bound}{Proposition~\ref*{prop:dark-slack-bound}}).}

% claim receiver/detector-imperfections/def:darkadm
\claimsummary{receiver/detector-imperfections/def:darkadm}{A dark-count channel is a column-stochastic matrix on click patterns that can add clicks and never removes one. It is generated from per-port dark-count probabilities in product form, or declared as a matrix and then checked at the receiver tolerances.}
\claimbookmark{receiver/detector-imperfections/def:darkadm}{Definition}{def:darkadm}{Dark-count channel and its independent product form}
\claimleadin{%
With the efficiencies inside the Grams, the remaining detector imperfection is
a click produced with no incident light. Dark counts change the recorded
\termref{term-outcome}{outcome}, not the \termref{term-optical-mode}{optical modes}, so they are a
classical channel on the outcome label, and a channel satisfying the
definition's conditions (D1) and (D2) exactly can only \emph{help} the
multiclick bound. The care is in
\emph{how much} it helps. A declared channel is checked in the floating-point
\termref{term-evaluation-model}{evaluation model} of Definition~\ref{def:evaluation-model}.
}
\begin{definition}[Dark-count channel and its independent product form]\label{def:darkadm}
A \termdef{term-dark-count-channel}{dark-count channel}
is a matrix $R=(R_{S|S'})_{S,S'\subseteq[K]}$ that is
\begin{enumerate}[label=(D\arabic*),leftmargin=3.0em]
\item \emph{column-stochastic}: $R_{S|S'}\ge0$ and $\sum_SR_{S|S'}=1$ for every
  $S'$;
\item \emph{monotone}: $R_{S|S'}>0$ implies $S'\subseteq S$; a dark count
  adds a click and never removes one.
\end{enumerate}
The \termdef{term-independent-product-form}{independent product form}
generated by per-port dark-count probabilities $p_{\dc|k}\in[0,1]$ is
\begin{equation}\label{eq:darkR}
  R_{S|S'}\coloneqq
  \mathds{1}\big[S'\subseteq S\big]
  \prod_{k\in S\setminus S'}p_{\dc|k}
  \prod_{k\notin S}\big(1-p_{\dc|k}\big).
\end{equation}
For $\mathcal{O}\subseteq2^{[K]}$ we write
$R_{\mathcal{O}|S'}\coloneqq\sum_{S\in\mathcal{O}}R_{S|S'}$.
A dark-count channel is supplied on one of two routes: on the
\termdef{term-generated-dark-count-route}{generated dark-count route}
it is \eqref{eq:darkR}, and on the
\termdef{term-declared-dark-count-route}{declared dark-count route}
it is a declared matrix $R$.
(D1) and (D2) hold by construction for \eqref{eq:darkR}
(Theorem~\ref{thm:dark-counts}\,(i)); for a \emph{declared} $R$ they are conditions to
be checked, and in floating point they are decided at the
\termref{term-receiver-tolerances}{receiver tolerances} $\varepsilon_4$ and $\varepsilon_2$ respectively, by
item (iii) of Definition~\ref{def:tolmap}, whose content is
Definition~\ref{def:darktol}.
\end{definition}
\claimmeta{\metalabel{Uses}\mref{Def.}{def:tolmap}
\metasep \metalabel{Used by}\mref{Thm.}{thm:dark-counts}, \mref{Def.}{def:darktol}, \mref{Prop.}{prop:dark-slack-bound}, \mref{Def.}{def:dark-composed}, \mref{Cor.}{cor:exact-dark-composed}, \mref{Prop.}{prop:noise-floor}, \mref{Prop.}{prop:efficiency-normalization}, \mref{Req.}{req:dark-route-exclusive}, \mref{Req.}{req:dark-composed-reference}, \mref{Req.}{req:snap-fix}, \extref{../reductions/contract.pdf}{reductions/wpfss/prop:route}, \extref{../reductions/contract.pdf}{reductions/wpfss/rem:darkcompose}, \extref{../reductions/contract.pdf}{reductions/wpfss/rem:receiver-imports}
\metasep newly derived
\metasep implemented (effective: accepted)
\metasep \metalabel{Code}\coderef{https://github.com/fraunhofersg/certiqs-key/blob/v0.1.0/src/certiqs_key/receiver/dark.py}{certiqs_key.receiver.dark:check_dark_channel}}

% claim receiver/detector-imperfections/thm:dark-counts
\claimsummary{receiver/detector-imperfections/thm:dark-counts}{A dark-count channel keeps the photon-number block-diagonal POVM structure and never moves weight out of the multiclick set. It raises the multiclick bound p\_mc to (1 \ensuremath{-} r\_min)p\_mc + r\_min, where r\_min is the least multiclick mass over every non-multiclick column, the no-click column included.}
\claimbookmark{receiver/detector-imperfections/thm:dark-counts}{Theorem}{thm:dark-counts}{Dark counts on a photon-number block-diagonal POVM}
\claimleadin{%
The theorem is stated for a channel that satisfies (D1) and (D2) exactly. For
a channel admitted at the receiver tolerances,
Proposition~\ref{prop:dark-slack-bound} states what holds in place of items
(iii) and (iv), and Remark~\ref{rem:d2slack}\,(f) states what remains of
item (ii). Of the POVM it uses only that it is a POVM whose elements commute
with the photon number, so it holds at every setting of a \termref{term-setting-family}{setting family}, the
click-pattern POVM being the one setting of a receiver without a probe.
}
\begin{mainclaim}
\begin{theorem}[Dark counts on a photon-number block-diagonal POVM]\label{thm:dark-counts}
Let $R$ be a dark-count channel, let $\{M_S\}_{S\subseteq[K]}$ be the
click-pattern POVM of Theorem~\ref{thm:click-povm} or, at a setting $t$ of a
\termref{term-setting-family}{setting family} of $K$ ports, its POVM $\{M_{S|t}\}_{S\subseteq[K]}$
(Definition~\ref{def:setting-family}\,(a)), written $\{M_S\}$, and put
$\widetilde{M}_S\coloneqq\sum_{S'\subseteq[K]}R_{S|S'}M_{S'}$.
\begin{enumerate}[label=(\roman*),leftmargin=3.0em]
\item \emph{The product form is admissible.} \eqref{eq:darkR} satisfies (D1)
  and (D2).
\item \emph{Structure is preserved.} $\{\widetilde M_S\}_{S\subseteq[K]}$ is a
  POVM, every $\widetilde M_S$ is block diagonal in the photon number with
  $\blk{\widetilde M}{\eqn}_S=\sum_{S'}R_{S|S'}\blk{M}{\eqn}_{S'}$, and $R$ acts
  on the outcome label alone: it commutes with every operation on the optical
  modes.
\item \emph{Multiclick absorption.} (D2) implies
  $R_{\MCK|S'}=1$ for every $S'\in\MCK$: dark counts never move weight out of
  the \termref{term-multiclick-set}{multiclick set}.
\item \emph{The bound improves.} With
  \begin{equation}\label{eq:rmin}
    r_{\min}\coloneqq\min_{S'\in2^{[K]}\setminus\MCK}R_{\MCK|S'} ,
  \end{equation}
  the minimum running over \emph{every} non-multiclick column, the no-click
  column $S'=\varnothing$ included,
  one has, writing $\widetilde M_{\MCK}=\sum_{S\in\MCK}\widetilde M_S$,
  \begin{equation}\label{eq:darkimprove}
    \widetilde M_{\MCK}\ \succeq\ (1-r_{\min})\,M_{\MCK}+r_{\min}\,\id ,
  \end{equation}
  and consequently, if $\blk{M}{\gtN}_{\MCK}\succeq\pmc\,\id^{(>N)}$, then
  \begin{equation}\label{eq:darkpmc}
    \blk{\widetilde M}{\gtN}_{\MCK}\ \succeq\
    \big((1-r_{\min})\,\pmc+r_{\min}\big)\,\id^{(>N)} ,
  \end{equation}
  which is $\ge\pmc$ since $\pmc\le1$.
\end{enumerate}
\end{theorem}
\end{mainclaim}
\claimmeta{\metalabel{Uses}\mref{Thm.}{thm:click-povm}, \mref{Def.}{def:setting-family}, \mref{Def.}{def:darkadm}
\metasep \metalabel{Used by}\mref{Prop.}{prop:dark-slack-bound}, \mref{Rem.}{rem:rmintrap}, \extref{../reductions/contract.pdf}{reductions/wpfss/prop:route}, \extref{../reductions/contract.pdf}{reductions/wpfss/rem:darkcompose}, \extref{../reductions/contract.pdf}{reductions/wpfss/rem:receiver-imports}
\metasep newly derived
\metasep implemented (effective: accepted)
\metasep \metalabel{Code}\coderef{https://github.com/fraunhofersg/certiqs-key/blob/v0.1.0/src/certiqs_key/receiver/dark.py}{certiqs_key.receiver.dark:compose_blocks}
\metasep \metalabel{Oracles}\coderef{https://github.com/fraunhofersg/certiqs-key/blob/v0.1.0/tests/oracles/registry_receiver_povm.py}{registry_receiver_povm.py}}

\begin{proof}
(i) Non-negativity is clear. Fix $S'$ and write $S=S'\sqcup W$ with
$W\subseteq[K]\setminus S'$; then
$\sum_SR_{S|S'}=\sum_{W\subseteq[K]\setminus S'}
\prod_{k\in W}p_{\dc|k}\prod_{k\in[K]\setminus S'\setminus W}(1-p_{\dc|k})
=\prod_{k\notin S'}\big(p_{\dc|k}+(1-p_{\dc|k})\big)=1$, the middle equality
being the expansion of the product. (D2) is the indicator in
\eqref{eq:darkR}.

Of $\{M_S\}$ the proof uses only that it is a POVM and that every element
commutes with the photon number: for the click-pattern POVM these are
Theorem~\ref{thm:click-povm}\,(i) and~(iii), and at a setting of a setting
family they are Definition~\ref{def:setting-family}\,(a).

(ii) $\widetilde M_S\succeq0$ as a non-negative combination of positive
operators, and, by (D1) and the completeness of $\{M_S\}$,
\begin{equation}
  \sum_S\widetilde M_S=\sum_{S'}\Big(\sum_SR_{S|S'}\Big)M_{S'}
  =\sum_{S'}M_{S'}=\id .
\end{equation}
Block diagonality is inherited
from each $M_{S'}$; the coefficients are
scalars, so they commute with everything.

(iii) If $S'\in\MCK$ and $R_{S|S'}>0$ then $S\supseteq S'$ by (D2), so
$|S|\ge|S'|\ge2$ and $S\in\MCK$. Hence every non-zero entry of the column $S'$
lies in $\MCK$, and column-stochasticity gives $R_{\MCK|S'}=1$.

(iv) Split the sum over columns:
\begin{align}
  \widetilde M_{\MCK}
  &=\sum_{S'\in\MCK}R_{\MCK|S'}M_{S'}
   +\sum_{S'\notin\MCK}R_{\MCK|S'}M_{S'}
  \ \ \overset{\text{(iii)}}{=}\ \
   M_{\MCK}+\sum_{S'\notin\MCK}R_{\MCK|S'}M_{S'}\nonumber\\
  &\succeq M_{\MCK}+r_{\min}\sum_{S'\notin\MCK}M_{S'}
   =M_{\MCK}+r_{\min}\big(\id-M_{\MCK}\big)
   =(1-r_{\min})M_{\MCK}+r_{\min}\id ,
\end{align}
using $M_{S'}\succeq0$ and $R_{\MCK|S'}\ge r_{\min}$ in the inequality and
the completeness of $\{M_S\}$ in the last line. All operators are block diagonal
in the photon number, so the inequality restricts to the $>N$ blocks; since
$1-r_{\min}\ge0$, substituting the assumed bound gives \eqref{eq:darkpmc}. The
final comparison is
$(1-r_{\min})\pmc+r_{\min}-\pmc=r_{\min}(1-\pmc)\ge0$.
\end{proof}

% claim receiver/detector-imperfections/def:darktol
\claimsummary{receiver/detector-imperfections/def:darktol}{For a declared dark-count matrix, column-stochasticity is decided at \ensuremath{\varepsilon}\_4 and monotonicity at \ensuremath{\varepsilon}\_2. This is item (iii) of \hyperlink{claim-receiver/def:tolmap}{Definition~\ref*{def:tolmap}}; the slack the two tests admit is carried into the bound by \hyperlink{claim-receiver/detector-imperfections/prop:dark-slack-bound}{Proposition~\ref*{prop:dark-slack-bound}}.}
\claimbookmark{receiver/detector-imperfections/def:darktol}{Definition}{def:darktol}{Receiver tolerances deciding dark-channel admissibility}
\claimleadin{%
Proposition~\ref{prop:dark-slack-bound} states what both slacks cost and which
side of each test each step uses: the lower side of (D1)'s column sum together
with the upper side of (D2)'s entry bound gives \eqref{eq:d2slack}, the upper
side of (D1)'s column sum permits $r_{\min}>1$ \eqref{eq:rminupper}, and what a
downstream module composes with is \eqref{eq:darkpmcslack} rather than
\eqref{eq:darkpmc}. The floating-point evaluation model in which the two tests
decide is Definition~\ref{def:evaluation-model}.
}
\begin{definition}[Receiver tolerances deciding dark-channel admissibility]
\label{def:darktol}
This definition extends the list of Definition~\ref{def:tolmap} by its
item~(iii), under the same rule: an exact condition about \exttermref{../foundations/contract.pdf}{term-declared-data}{declared data}
is decided in floating point by a receiver
tolerance, fixed here, and by no other threshold.

\emph{Dark-channel admissibility} (Definition~\ref{def:darkadm}) for a
\emph{declared} $R$: (D1) column-stochasticity is decided at
$\varepsilon_4$, admitted when $R_{S|S'}\ge-\varepsilon_4$ for every
entry and $\big|\sum_SR_{S|S'}-1\big|\le\varepsilon_4$ for every column;
and (D2) monotonicity is decided at $\varepsilon_2$: an entry counts as
positive when $R_{S|S'}>\varepsilon_2$, so the channel is refused exactly
when some entry with $S'\not\subseteq S$ exceeds $\varepsilon_2$.
$\varepsilon_4$ bounds the completeness residual \eqref{eq:res4}, which is
the same kind of quantity as a column sum's defect from one; $\varepsilon_2$
is the threshold at which this contract already calls a spectral quantity
zero.
\end{definition}
\claimmeta{\metalabel{Uses}\mref{Def.}{def:residuals}, \mref{Def.}{def:tolmap}, \mref{Def.}{def:darkadm}
\metasep \metalabel{Used by}\mref{Prop.}{prop:dark-slack-bound}, \mref{Req.}{req:dark-composed-reference}, \extref{../reductions/contract.pdf}{reductions/wpfss/rem:darkcompose}
\metasep newly derived
\metasep accepted}

% claim receiver/detector-imperfections/prop:dark-slack-bound
\claimsummary{receiver/detector-imperfections/prop:dark-slack-bound}{For a declared dark-count matrix admitted at \ensuremath{\varepsilon}\_2 and \ensuremath{\varepsilon}\_4, each multiclick column keeps mass at least 1 \ensuremath{-} \ensuremath{\delta}\_dc, r\_min can exceed one by at most (K+2)\ensuremath{\varepsilon}\_4, and the dark-count improvement holds with rbar = min\{r\_min, 1\} and \ensuremath{\delta}\_dc subtracted. This is the inequality a downstream module composes with.}
\claimbookmark{receiver/detector-imperfections/prop:dark-slack-bound}{Proposition}{prop:dark-slack-bound}{Multiclick bound under a dark-count channel admitted with slack}
\claimleadin{%
Theorem~\ref{thm:dark-counts} needs (D1) and (D2) exactly, and the tests of
Definition~\ref{def:darktol} admit a declared channel that satisfies them only
up to $\varepsilon_4$ and $\varepsilon_2$. The proposition carries both slacks
into the inequality that is composed with the multiclick bound: the lower side
of (D1)'s column sum and the upper side of (D2)'s entry bound cost a constant
$\delta_{\dc}$. The upper side of (D1), column sums above one and entries
down to $-\varepsilon_4$, lets $r_{\min}$ exceed one; the coefficient is
capped at one, at a cost bounded in (ii).
}
\begin{proposition}[Multiclick bound under a dark-count channel admitted with slack]\label{prop:dark-slack-bound}
Let $\{M_S\}_{S\subseteq[K]}$ be the click-pattern POVM of
Theorem~\ref{thm:click-povm} or, at a setting $t$ of a \termref{term-setting-family}{setting family}
of $K$ ports, its POVM $\{M_{S|t}\}_{S\subseteq[K]}$
(Definition~\ref{def:setting-family}\,(a)), written $\{M_S\}$, and for a matrix $R=(R_{S|S'})_{S,S'\subseteq[K]}$ let
$\widetilde M_S$, $\widetilde M_{\MCK}$ and $r_{\min}$ be formed as in
Theorem~\ref{thm:dark-counts} and \eqref{eq:rmin}. Put
\begin{equation}\label{eq:darkdelta}
  \delta_{\dc}\ \coloneqq\
  \begin{cases}
    0
      & \text{generated route: }R\text{ is the product form }\eqref{eq:darkR},
      \\[3pt]
    \varepsilon_4+(K+1)\,\varepsilon_2
      & \text{declared route: }R\text{ admitted under
        Def.~\ref{def:tolmap}\,(iii),}
  \end{cases}
\end{equation}
the two cases being the two routes of Definition~\ref{def:darkadm}, and
\begin{equation}\label{eq:rbar}
  \bar r\ \coloneqq\ \min\{\,r_{\min},\,1\,\} .
\end{equation}
Let $R$ be declared and admitted under item (iii) of
Definition~\ref{def:tolmap} (Definition~\ref{def:darktol}), on either side of
either test and whatever the sign of $r_{\min}$. Then:
\begin{enumerate}[label=(\roman*),leftmargin=3.0em]
\item \emph{(the multiclick columns)} for every $S'\in\MCK$,
  \begin{equation}\label{eq:d2slack}
    R_{\MCK|S'}\ \ge\ (1-\varepsilon_4)-(K+1)\varepsilon_2\ =\ 1-\delta_{\dc} ;
  \end{equation}
\item \emph{(the upper side)} for every column $S'$,
  \begin{equation}\label{eq:rminupper}
    R_{\MCK|S'}\ \le\ (1+\varepsilon_4)+(K+1)\varepsilon_4\ =\ 1+(K+2)\varepsilon_4 ,
  \end{equation}
  and consequently
  \begin{equation}\label{eq:rbarcost}
    0\ \le\ r_{\min}-\bar r\ \le\ (K+2)\,\varepsilon_4 ;
  \end{equation}
\item \emph{(the composed inequality)}
  \begin{equation}\label{eq:darkimproveslack}
    \widetilde M_{\MCK}\ \succeq\
    (1-\delta_{\dc})\,M_{\MCK}+\bar r\,\big(\id-M_{\MCK}\big)
    \ \succeq\ \Big[(1-\bar r)\,M_{\MCK}+\bar r\,\id\Big]-\delta_{\dc}\,\id ,
  \end{equation}
  and consequently, whenever $\blk{M}{\gtN}_{\MCK}\succeq\pmc\,\id^{(>N)}$,
  \begin{equation}\label{eq:darkpmcslack}
    \blk{\widetilde M}{\gtN}_{\MCK}\ \succeq\
    \big((1-\bar r)\,\pmc+\bar r-\delta_{\dc}\big)\,\id^{(>N)} .
  \end{equation}
\end{enumerate}
\end{proposition}
\claimmeta{\metalabel{Uses}\mref{Thm.}{thm:click-povm}, \mref{Def.}{def:setting-family}, \mref{Def.}{def:tolmap}, \mref{Def.}{def:darkadm}, \mref{Thm.}{thm:dark-counts}, \mref{Def.}{def:darktol}
\metasep \metalabel{Used by}\mref{Rem.}{rem:rmintrap}, \mref{Rem.}{rem:d2slack}, \mref{Rem.}{rem:corrdark}, \mref{Req.}{req:dark-channel-record}
\metasep newly derived
\metasep implemented (effective: accepted)
\metasep \metalabel{Code}\coderef{https://github.com/fraunhofersg/certiqs-key/blob/v0.1.0/src/certiqs_key/receiver/dark.py}{certiqs_key.receiver.dark:composed_slack}}

\begin{proof}
\emph{Idea.} A multiclick column mass is the column sum minus the entries in
the $K+1$ rows with at most one click; the two tests bound each of these from
both sides, which gives (i) and (ii). The composed operator is a combination of
the positive operators $M_{S'}$ with the column masses as coefficients, and
lowering every coefficient to its bound gives (iii).

(i) Let $S'\in\MCK$. The rows $S\notin\MCK$ are exactly those with $|S|\le1$,
of which there are $K+1$, and none of them contains $S'$ (as
$|S|\le1<2\le|S'|$), so every entry $R_{S|S'}$ with $S\notin\MCK$ sits in a
(D2)-forbidden position and is at most $\varepsilon_2$. The (D1) test bounds
the column sum below by $1-\varepsilon_4$. Hence
\begin{equation*}
  R_{\MCK|S'}\ =\ \sum_{S}R_{S|S'}-\sum_{S\notin\MCK}R_{S|S'}
  \ \ge\ (1-\varepsilon_4)-(K+1)\varepsilon_2 ,
\end{equation*}
which is $1-\delta_{\dc}$ by the declared case of \eqref{eq:darkdelta}.

(ii) The widened (D1) admits column sums up to $1+\varepsilon_4$ and entries
down to $-\varepsilon_4$, so for every column
\begin{equation*}
  R_{\MCK|S'}\ =\ \sum_SR_{S|S'}-\sum_{S\notin\MCK}R_{S|S'}
  \ \le\ (1+\varepsilon_4)+(K+1)\varepsilon_4 ,
\end{equation*}
which is \eqref{eq:rminupper}. The cost \eqref{eq:rbarcost} is bounded by
\eqref{eq:rminupper}: $\bar r\le r_{\min}$ by \eqref{eq:rbar}, and
$r_{\min}-\bar r$ is $0$ when $r_{\min}\le1$ and $r_{\min}-1$ otherwise, where
$r_{\min}$, a minimum of column masses, is at most $1+(K+2)\varepsilon_4$.

(iii) $\widetilde M_{\MCK}=\sum_{S'}R_{\MCK|S'}M_{S'}$ with every
$M_{S'}\succeq0$, $\{M_S\}$ being a POVM (Theorem~\ref{thm:click-povm}\,(i),
or Definition~\ref{def:setting-family}\,(a) at a setting), so replacing each coefficient by
a smaller real gives a L\"owner-smaller operator. On the columns $S'\in\MCK$
use $R_{\MCK|S'}\ge1-\delta_{\dc}$, which is \eqref{eq:d2slack}; on the columns
$S'\notin\MCK$ use $R_{\MCK|S'}\ge r_{\min}\ge\bar r$, the first by the
definition \eqref{eq:rmin} of $r_{\min}$ as their minimum and the second by
\eqref{eq:rbar}. Summing, with $\sum_{S'\in\MCK}M_{S'}=M_{\MCK}$ and
$\sum_{S'\notin\MCK}M_{S'}=\id-M_{\MCK}$ by completeness, gives
the first inequality of \eqref{eq:darkimproveslack}; regrouping and applying
$0\preceq M_{\MCK}\preceq\id$ with $\delta_{\dc}\ge0$ gives the second. Every
operator here is block diagonal in the photon number
(Theorem~\ref{thm:click-povm}\,(iii), or
Definition~\ref{def:setting-family}\,(a) at a setting), so the inequality
restricts to the blocks $n>N$. There $1-\bar r\ge0$ by construction of
\eqref{eq:rbar}, so
$(1-\bar r)\blk{M}{\gtN}_{\MCK}\succeq(1-\bar r)\pmc\,\id^{(>N)}$, and adding
$(\bar r-\delta_{\dc})\,\id^{(>N)}$ gives \eqref{eq:darkpmcslack}. No step uses
a bound on $r_{\min}$ from either side, and no step uses the sign of $r_{\min}$
or of $\bar r$.
\end{proof}

% claim receiver/detector-imperfections/rem:rmintrap
\claimsummary{receiver/detector-imperfections/rem:rmintrap}{Taking r\_min over the single-click columns only gives an invalid bound whenever the no-click column has less multiclick mass, the generic case: at K = 4 and p = 0.00001 the two differ by four orders of magnitude. A lossy receiver keeps no-click weight in every high-photon block.}
\claimbookmark{receiver/detector-imperfections/rem:rmintrap}{Remark}{rem:rmintrap}{The no-click column in the minimum defining r\_min}
\begin{remark}[The no-click column in the minimum defining $r_{\min}$]\label{rem:rmintrap}
The minimum in \eqref{eq:rmin} must run over \emph{all} non-multiclick columns,
the no-click column $S'=\varnothing$ included. Minimising only over the
single-click columns gives an \emph{invalid} bound whenever
$R_{\MCK|\varnothing}<R_{\MCK|\{k\}}$, and that is the generic case rather than
a corner one. With equal per-port probabilities $p_{\dc|k}=p$ and $K=4$,
\begin{equation}\label{eq:rmintrapnum}
  R_{\MCK|\varnothing}=1-(1-p)^4-4p(1-p)^3=6p^2+O(p^3),
  \qquad
  R_{\MCK|\{k\}}=1-(1-p)^3=3p+O(p^2),
\end{equation}
so at a realistic $p=10^{-5}$ the two differ by more than four orders of
magnitude: $6\cdot10^{-10}$ against $3\cdot10^{-5}$. Using the larger number
would claim a multiclick weight the device does not have.

The no-click column cannot be excluded because the no-click element does not
vanish in the high-photon sectors of a lossy receiver: by
Lemma~\ref{lem:fock-block-spectrum},
$\lmin\big(\blk{M}{\eqn}_{\varnothing}\big)
=\big(1-\lmax(A_\Sigma)\big)^{n}>0$ whenever $\lmax(A_\Sigma)<1$. So
in exactly the regime where the high-photon bound matters, the column that the
shortcut discards carries weight. The correct statement is
Theorem~\ref{thm:dark-counts}\,(iv) as written, and \eqref{eq:darkpmcslack} of
Proposition~\ref{prop:dark-slack-bound} for a declared $R$ admitted under
item (iii) of Definition~\ref{def:tolmap} (Definition~\ref{def:darktol}).
Which columns the minimum runs over is a question about the index set and is
untouched by either widened test; the numbers above are those of the
\termref{term-generated-dark-count-route}{generated dark-count route}, where both conditions are theorems.
\end{remark}
\claimmeta{\metalabel{Uses}\mref{Lem.}{lem:fock-block-spectrum}, \mref{Def.}{def:tolmap}, \mref{Thm.}{thm:dark-counts}, \mref{Prop.}{prop:dark-slack-bound}
\metasep newly derived
\metasep accepted}

% claim receiver/detector-imperfections/rem:d2slack
\claimsummary{receiver/detector-imperfections/rem:d2slack}{Admitting a declared channel at \ensuremath{\varepsilon}\_2 and \ensuremath{\varepsilon}\_4 relaxes a hypothesis, not a precision; r\_min above one is reachable, and ignoring it or \ensuremath{\delta}\_dc overstates the multiclick bound in the unsafe direction. The composed family is a POVM only up to \ensuremath{\varepsilon}\_4, a defect left visible and not repaired.}
\claimbookmark{receiver/detector-imperfections/rem:d2slack}{Remark}{rem:d2slack}{Cost of admitting a dark-count channel with slack}
\begin{remark}[Cost of admitting a dark-count channel with slack]\label{rem:d2slack}
\emph{(a) What is relaxed.} Definition~\ref{def:tolmap}\,(iii) admits a
declared $R$ that violates (D2) by up to $\varepsilon_2$ in every forbidden
position and (D1) by up to $\varepsilon_4$ in every column.
Theorem~\ref{thm:dark-counts} is stated and proved for an $R$ satisfying (D1) and (D2)
\emph{exactly}, and admitting one that satisfies them only up to
$(\varepsilon_2,\varepsilon_4)$ does not weaken the theorem; but the module then
emits an object the theorem's hypothesis does not cover, and what that costs
has to be stated at the inequality a downstream module actually uses, not only
at the identity it is proved through. Proposition~\ref{prop:dark-slack-bound}
states it. This is not a matter of the evaluation model of
Definition~\ref{def:evaluation-model}: the datum is
exactly as declared, and what is relaxed is a hypothesis, not the precision of
an evaluation, so Theorem~\ref{thm:exact-and-evaluated}\,(ii) does not reach it.

\emph{(b) The two values of $\delta_{\dc}$.} On the
\termref{term-generated-dark-count-route}{generated dark-count route} (D1) and (D2) are \emph{theorems}
(Theorem~\ref{thm:dark-counts}\,(i)) and not tests, so there is no slack to carry and
the correction in \eqref{eq:darkdelta} is exactly $0$; on the
\termref{term-declared-dark-count-route}{declared dark-count route} they are tested, and
Proposition~\ref{prop:dark-slack-bound} is what the test admits. In
\eqref{eq:d2slack}, against the exact value $1$ of
Theorem~\ref{thm:dark-counts}\,(iii), \emph{both} slacks appear: reading
\eqref{eq:d2slack} off an exact column sum would drop the $\varepsilon_4$ and
overstate the bound for every channel this contract admits.

\emph{(c) The upper side of the same tests, and where it bites.}
\eqref{eq:rmin} defines $r_{\min}$ as a minimum of multiclick column masses, and
for an $R$ satisfying (D1) and (D2) exactly it satisfies $r_{\min}\le1$ for
free: entries are non-negative and columns sum to $1$, so
$R_{\MCK|S'}\le\sum_SR_{S|S'}=1$. The widened (D1) does not supply that, and
$r_{\min}>1$ is reachable within \eqref{eq:rminupper}: a declared $R$ placing
mass $1+\varepsilon_4$ on the row $S=[K]$ in every column $S'\notin\MCK$ and
$1-\varepsilon_4$ on that same row in every column $S'\in\MCK$ occupies no
(D2)-forbidden position at all, is admitted at
$(\varepsilon_2,\varepsilon_4)=(0,\varepsilon_4)$, and has
$r_{\min}=1+\varepsilon_4$. Using the lower side of (D1) in \eqref{eq:d2slack}
and passing over the upper side here would be exactly the one-sided reading
that Remarks~\ref{rem:ranktwosided} and~\ref{rem:cdirection} guard against,
and it would not be harmless: the coefficient $1-r_{\min}$ that
Theorem~\ref{thm:dark-counts}\,(iv)'s substitution step requires to be non-negative
would not be, and the composed bound would fail, in the unsafe direction, by
$(r_{\min}-1)(1-\pmc)$.

\emph{(d) $\bar r$ is not a repair, and it costs a bounded amount.} Nothing is
clamped (Requirement~\ref{req:dark-channel-record}): $\bar r$ is a separate
named scalar appearing as a \emph{coefficient in a stated inequality}.
Lowering that coefficient \emph{weakens} the bound, and needs nothing beyond
$\id-M_{\MCK}\succeq0$; no claim is made that $\bar r$ approximates $r_{\min}$
or any exact quantity. That is the distinction
Definition~\ref{def:tolmap}\,(iv) draws in the other direction, where clamping
a computed $c$ before it enters \eqref{eq:beta} would assert a relation to an
exact value that nothing supports. The cost \eqref{eq:rbarcost} is first order
in $\varepsilon_4$ and no larger.

\emph{(e) Relation to the exact composed inequality.} On the generated route
$\delta_{\dc}=0$ and $r_{\min}\le1$, so $\bar r=r_{\min}$ and
\eqref{eq:darkpmcslack} \emph{is} \eqref{eq:darkpmc}. On the declared route a
downstream module that used \eqref{eq:darkpmc} instead would overstate by
\begin{equation}\label{eq:darkoverstate}
  \big(r_{\min}-\bar r\big)\big(1-\pmc\big)+\delta_{\dc} ,
\end{equation}
first order in the \termref{term-receiver-tolerances}{receiver tolerances} and in the \emph{unsafe} direction: a
(D2) violation is precisely weight moved out of $\MCK$, and $r_{\min}>1$ is
multiclick mass the column arithmetic does not have. And
\eqref{eq:darkpmcslack} improves on $\pmc$, which is the point of
Theorem~\ref{thm:dark-counts}\,(iv), only where $\bar r\,(1-\pmc)\ge\delta_{\dc}$;
below that threshold the declared channel's admitted slack costs more than its
dark counts gain, and the two emitted numbers say so.

\emph{(f) Theorem~\ref{thm:dark-counts}\,(ii) is outside the widened hypotheses too,
and is not repaired either.} For an $R$ admitted at
$(\varepsilon_2,\varepsilon_4)$ the composed family $\{\widetilde M_S\}$ need
not be a POVM. Writing $\sigma_{S'}\coloneqq\sum_SR_{S|S'}$, one has
$\sum_S\widetilde M_S=\sum_{S'}\sigma_{S'}M_{S'}$ with
$\sigma_{S'}\in[1-\varepsilon_4,1+\varepsilon_4]$, so
$(1-\varepsilon_4)\id\preceq\sum_S\widetilde M_S\preceq(1+\varepsilon_4)\id$
in place of $=\id$; and a single entry at $-\varepsilon_4$ gives only
$\widetilde M_S\succeq-\varepsilon_4\,\id$ in place of $\succeq0$. Both defects
are bounded by $\varepsilon_4$, both are properties of the declared datum, and
neither is repaired. They are not silent: every consumer decides the
admission predicates it states of the data it reads exactly and refuses on a
failure (\xref{foundations/req:consumer-decides-exactly}), so a consumer that
requires the composed family to be a POVM, or the channel to satisfy (D1) and
(D2), refuses it; and on this side a positivity defect of an emitted block is
what $r_5$ makes visible (Remark~\ref{rem:res5why}).
\end{remark}
\claimmeta{\metalabel{Uses}\mref{Thm.}{thm:click-povm}, \mref{Def.}{def:bounddata}, \mref{Rem.}{rem:res5why}, \mref{Def.}{def:tolmap}, \mref{Rem.}{rem:ranktwosided}, \mref{Prop.}{prop:dark-slack-bound}
\metasep \metalabel{Used by}\mref{Rem.}{rem:corrdark}
\metasep newly derived
\metasep accepted}

% claim receiver/detector-imperfections/rem:corrdark
\claimsummary{receiver/detector-imperfections/rem:corrdark}{A declared dark-count channel need not have product form: if it satisfies (D1) and (D2) exactly, the dark-count theorem holds for it verbatim, and if it passes the widened tests it holds in the weakened forms of \hyperlink{claim-receiver/detector-imperfections/prop:dark-slack-bound}{Proposition~\ref*{prop:dark-slack-bound}} and \hyperlink{claim-receiver/detector-imperfections/rem:d2slack}{Remark~\ref*{rem:d2slack}}. Correlations across rounds are excluded.}
\claimbookmark{receiver/detector-imperfections/rem:corrdark}{Remark}{rem:corrdark}{Correlated dark counts within a round}
\begin{remark}[Correlated dark counts within a round]\label{rem:corrdark}
Nothing in Theorem~\ref{thm:dark-counts} uses the product form beyond part~(i): a
declared, experimentally certified $R$ satisfying (D1) and (D2)
\emph{exactly} is admissible, and (ii)--(iv) hold for it verbatim. What must
then be \emph{checked} rather than derived is (D1) and (D2), and the artifact
records which of the two routes the emitted $R$ came from
(Requirement~\ref{req:dark-channel-record}). Where the check is the one this
contract performs, item (iii) of Definition~\ref{def:tolmap}
(Definition~\ref{def:darktol}) at $\varepsilon_4$ and $\varepsilon_2$, the
declared $R$ satisfies the two conditions only up to those tolerances, and
then (ii), (iii) and (iv) hold only as Proposition~\ref{prop:dark-slack-bound}
and Remark~\ref{rem:d2slack} restate them: the composed inequality is
\eqref{eq:darkpmcslack} in $\bar r$ and $\delta_{\dc}$, and the composed family
is a POVM only up to $\varepsilon_4$. Correlations \emph{across rounds} are a
different matter and are excluded by Assumption~\ref{ass:thresh}: they would
violate Assumption~\ref{ass:round} downstream, and no channel on a single
round's label can repair that.
\end{remark}
\claimmeta{\metalabel{Uses}\mref{Def.}{def:tolmap}, \mref{Prop.}{prop:dark-slack-bound}, \mref{Rem.}{rem:d2slack}
\metasep \metalabel{Assumes}\aref{A-THRESH}{ass:thresh}
\metasep newly derived
\metasep accepted}

\section{The dark-composed family and the noise floor}\label{group-receiver/detector-imperfections/def:dark-composed}
\sectionintro{The \termref{term-dark-composed-family}{dark-composed family} applies the \termref{term-dark-count-channel}{dark-count channel} to every setting's click patterns and then the coarse-graining (\hyperlink{claim-receiver/detector-imperfections/def:dark-composed}{Definition~\ref*{def:dark-composed}}). A dark-count channel has rational entries and carries enclosures of the click-pattern blocks to enclosures of the dark-composed ones, which form the reference handed beside the snap (\hyperlink{claim-receiver/detector-imperfections/cor:exact-dark-composed}{Corollary~\ref*{cor:exact-dark-composed}}). The main point of the section is \hyperlink{claim-receiver/detector-imperfections/prop:noise-floor}{Proposition~\ref*{prop:noise-floor}}: loss and a dark-count probability in (0, 1) in every port make every dark-composed element positive definite on every photon-number block, with an explicit floor, and without either some element is singular. The closing remark says what a declaration should therefore model, why an ideal device measuring in a rotated basis dominates no lattice family beyond one photon at a small weight (for the two-photon block of the passive BB84 receiver at a balanced basis choice, 1/4), and why a run at a cutoff of two or more photons that uses no enclosed reference typically rests on its snap.}
\sectionimports{\extref{../foundations/contract.pdf}{foundations/def:binary64}, \extref{../foundations/contract.pdf}{foundations/def:evaluated-off-snap}, \extref{../foundations/contract.pdf}{foundations/def:exact-reference}, \extref{../foundations/contract.pdf}{foundations/def:sector-map}, \extref{../foundations/contract.pdf}{foundations/lem:rational-obstruction}, \extref{../foundations/contract.pdf}{foundations/rem:noise-floor}}

% claim receiver/detector-imperfections/def:dark-composed
\claimsummary{receiver/detector-imperfections/def:dark-composed}{The dark-composed family applies a dark-count channel to the click patterns of each setting's POVM and then the declared coarse-graining. It is the measurement the device performs, and the family a snap is taken from.}
\claimbookmark{receiver/detector-imperfections/def:dark-composed}{Definition}{def:dark-composed}{The dark-composed family}
\claimleadin{%
Theorem~\ref{thm:dark-counts} composes a \termref{term-dark-count-channel}{dark-count channel} with the
click-pattern POVM of one receiver. The definition names the composed object
for every setting at once, and after the coarse-graining, because that is the
measurement a declaration describes when its device has dark counts. The
families it is applied to are the unconditioned receiver
(Corollary~\ref{cor:unconditioned-family}) and every \termref{term-setting-family}{setting family} a
submodule constructs.
}
\begin{definition}[The dark-composed family]\label{def:dark-composed}
Let $\Tset$ be a finite non-empty set of settings and, for every $t\in\Tset$,
let $\{M_{S|t}\}_{S\subseteq[K]}$ be a POVM on a \termref{term-fock-space}{Fock space}
$\mathcal{F}(\mathbb{C}^{I})$ (Definition~\ref{def:fock}) every element of
which commutes with the photon number. Let $R$ be a
\termref{term-dark-count-channel}{dark-count channel} (Definition~\ref{def:darkadm}) and $\cg$ a
\termref{term-declared-coarse-graining}{declared coarse-graining} onto $\Vset$ (Definition~\ref{def:coarse}).
The
\termdef{term-dark-composed-family}{dark-composed
family} is
\begin{equation}\label{eq:dark-composed}
  \widetilde M_{S|t}\coloneqq\sum_{S'\subseteq[K]}R_{S|S'}\,M_{S'|t},
  \qquad
  \widetilde F_{v|t}\coloneqq\sum_{S\in\cg^{-1}(v)}\widetilde M_{S|t}
  =\sum_{S'\subseteq[K]}R_{\cg^{-1}(v)|S'}\,M_{S'|t},
\end{equation}
for $S\subseteq[K]$, $v\in\Vset$ and $t\in\Tset$. Each
$\{\widetilde M_{S|t}\}_S$ and each $\{\widetilde F_{v|t}\}_v$ is a POVM
every element of which commutes with the photon number, with
$\blk{\widetilde F}{\eqn}_{v|t}=\sum_{S'}R_{\cg^{-1}(v)|S'}\blk{M}{\eqn}_{S'|t}$:
the coefficients are non-negative scalars, they sum to one over $S$ in every
column by (D1), and the fibres of $\cg$ partition $2^{[K]}$.
\end{definition}
\claimmeta{\metalabel{Uses}\mref{Def.}{def:fock}, \mref{Def.}{def:coarse}, \mref{Def.}{def:darkadm}
\metasep \metalabel{Used by}\mref{Cor.}{cor:exact-dark-composed}, \mref{Prop.}{prop:noise-floor}, \mref{Rem.}{rem:noise-floor-receiver}, \mref{Prop.}{prop:efficiency-normalization}, \mref{Req.}{req:dark-composed-reference}, \extref{../reductions/contract.pdf}{reductions/wpfss/cor:normalized-flat}
\metasep newly derived
\metasep implemented (effective: accepted)
\metasep \metalabel{Code}\coderef{https://github.com/fraunhofersg/certiqs-key/blob/v0.1.0/src/certiqs_key/receiver/generic.py}{certiqs_key.receiver.generic:dark_composed_blocks}}

% claim receiver/detector-imperfections/cor:exact-dark-composed
\claimsummary{receiver/detector-imperfections/cor:exact-dark-composed}{A dark-count channel has rational entries and mixes whole elements with non-negative weights, so enclosures of the click-pattern blocks give enclosures of the dark-composed blocks, exact where theirs are. With Hermitian centres they make the dark-composed blocks up to the cutoff a reference.}
\claimbookmark{receiver/detector-imperfections/cor:exact-dark-composed}{Corollary}{cor:exact-dark-composed}{The dark-composed reference, enclosed or exact}
\claimleadin{%
A dark count acts on the recorded click pattern, not on the light, so it
mixes whole elements with scalar weights. An enclosure therefore passes
through it with the same weights, and an exact one stays exact. Item~(iii)
applies the corollary to the unconditioned receiver; for a setting family a
submodule constructs, the hypothesis of item~(i) is the enclosure of its
blocks that the submodule states. Where the click-pattern blocks
are enclosed exactly up to one photon (Theorem~\ref{thm:exact-blocks}), so are
the dark-composed ones.
}
\begin{corollary}[The dark-composed reference, enclosed or exact]\label{cor:exact-dark-composed}
Let $\{\widetilde M_{S|t}\}$ and $\{\widetilde F_{v|t}\}$ be the
\termref{term-dark-composed-family}{dark-composed family} of POVMs $\{M_{S|t}\}$ on
$\mathcal{F}(\mathbb{C}^{I})$ (Definition~\ref{def:dark-composed}), with a
\termref{term-dark-count-channel}{dark-count channel} $R$ whose entries are rational.
\begin{enumerate}[label=(\roman*),leftmargin=3.0em]
\item Suppose that for some $N\ge0$ every block $\blk{M}{\eqn}_{S|t}$ with
  $n\le N$, read in the occupation basis of $\Symn{n}(\mathbb{C}^{I})$ in an
  order fixed for that block, has a \termref{term-complex-enclosure}{complex enclosure}
  $\big(C^{(n)}_{S|t},r^{(n)}_{S|t}\big)$
  (Definition~\ref{def:complex-enclosure}). Then
  \begin{equation}\label{eq:exact-dark}
    \Big(\sum_{S'\subseteq[K]}R_{\cg^{-1}(v)|S'}\,C^{(n)}_{S'|t},\
         \sum_{S'\subseteq[K]}R_{\cg^{-1}(v)|S'}\,r^{(n)}_{S'|t}\Big)
  \end{equation}
  is a complex enclosure of $\blk{\widetilde F}{\eqn}_{v|t}$, exact on every
  entry on which every $r^{(n)}_{S'|t}$ with $R_{\cg^{-1}(v)|S'}\neq0$ is
  zero; the same holds for $\blk{\widetilde M}{\eqn}_{S|t}$, with $R_{S|S'}$
  in place of $R_{\cg^{-1}(v)|S'}$.
\item Made Hermitian-centred by
  Lemma~\ref{lem:enclosure-arithmetic}\,(v), the enclosures of (i) make
  $\{\widetilde F_{v|t}\}_v$, restricted to the blocks $n\le N$, a
  \exttermref{../foundations/contract.pdf}{term-reference}{reference} with constant $1$
  (\xref{foundations/def:exact-reference}) at every setting, against every
  \exttermref{../foundations/contract.pdf}{term-sector-map}{sector map} whose blocks are the photon numbers $n=0,\dots,N$ in
  increasing order, with any block kind and identifiers, each block read in
  its fixed order; it is an \exttermref{../foundations/contract.pdf}{term-exact-reference}{exact reference} when every radius is zero.
  The same holds for $\{\widetilde M_{S|t}\}_S$.
\item For a \termref{term-receiver-instrument}{receiver instrument} with $\InpS=\varnothing$, the
  hypothesis of (i) holds with $\Tset=\{\off\}$, $M_{S|\off}=M_S$ and the
  enclosures of Theorem~\ref{thm:enclosed-blocks}\,(ii), at every $N$.
\item $R$ has rational entries on the \termref{term-generated-dark-count-route}{generated dark-count route},
  whose $p_{\dc|k}$ are declared values, and on the declared route, whose
  entries are declared values; a declared value is the dyadic rational it
  denotes (\xref{foundations/def:binary64}\,(iii)).
\end{enumerate}
\end{corollary}
\claimmeta{\metalabel{Uses}\mref{Def.}{def:complex-enclosure}, \mref{Lem.}{lem:enclosure-arithmetic}, \mref{Thm.}{thm:enclosed-blocks}, \mref{Def.}{def:darkadm}, \mref{Def.}{def:dark-composed}, \extref{../foundations/contract.pdf}{foundations/def:binary64}, \extref{../foundations/contract.pdf}{foundations/def:exact-reference}, \extref{../foundations/contract.pdf}{foundations/def:sector-map}
\metasep \metalabel{Used by}\mref{Rem.}{rem:noise-floor-receiver}, \mref{Req.}{req:dark-composed-reference}, \mref{Req.}{req:snap-fix}
\metasep newly derived
\metasep accepted}

\begin{proof}
(i) By Definition~\ref{def:dark-composed},
$\blk{\widetilde F}{\eqn}_{v|t}=\sum_{S'}R_{\cg^{-1}(v)|S'}\blk{M}{\eqn}_{S'|t}$
with non-negative rational coefficients $R_{\cg^{-1}(v)|S'}$, a sum of
entries of $R$. Entry by entry,
\begin{equation*}
  \Big\lvert\big(\blk{\widetilde F}{\eqn}_{v|t}\big)_{il}
   -\sum_{S'}R_{\cg^{-1}(v)|S'}\big(C^{(n)}_{S'|t}\big)_{il}\Big\rvert
  \le\sum_{S'}R_{\cg^{-1}(v)|S'}\,\big(r^{(n)}_{S'|t}\big)_{il} ,
\end{equation*}
the centre has entries in $\mathbb{Q}(i)$, and the radius is a non-negative
rational, zero exactly where every term with a non-zero coefficient has radius
zero. The statement for $\widetilde M_{S|t}$ is the same computation with
$\cg$ the identity.

(ii) $\{\widetilde F_{v|t}\}_v$ is a POVM every element of which commutes with
the photon number (Definition~\ref{def:dark-composed}), so each block is
Hermitian and positive semidefinite and the blocks of one $n$ sum to the
identity. Lemma~\ref{lem:enclosure-arithmetic}\,(v) turns the enclosures of
(i) into entrywise enclosures with Hermitian Gaussian-rational centres, exact
where both an entry and its transposed entry were. The block of photon number
$n$ has the dimension $\binom{|I|+n-1}{n}$ of $\Symn{n}(\mathbb{C}^{I})$, and
the constant $1$ is a positive declared value. These are the conditions of
\xref{foundations/def:exact-reference}; the block kind and identifiers enter
nowhere.

(iii) Theorem~\ref{thm:enclosed-blocks}\,(ii) gives complex enclosures of
every $\blk{M}{\eqn}_S$, formed in each block's fixed order.

(iv) Each entry of \eqref{eq:darkR} is a product of the numbers $p_{\dc|k}$
and $1-p_{\dc|k}$, rational when the $p_{\dc|k}$ are, and a declared value
denotes the dyadic rational it is.
\end{proof}

% claim receiver/detector-imperfections/prop:noise-floor
\claimsummary{receiver/detector-imperfections/prop:noise-floor}{With loss and a dark-count probability in (0,1) in every port, every dark-composed element is positive definite on every photon-number block, with a least eigenvalue bounded below through the no-click element. Without loss, or without a dark count in some port, some element is singular.}
\claimbookmark{receiver/detector-imperfections/prop:noise-floor}{Proposition}{prop:noise-floor}{Loss and dark counts make every dark-composed element definite}
\claimleadin{%
The \termref{term-no-click-outcome}{no-click outcome} is the only one that the vacuum and a lossy device
both reach on every photon-number block, and a dark count can turn it into
any other outcome. That single column of the dark-count channel therefore
spreads a positive multiple of the no-click element into every element. Items
(i)--(iii) say when this makes every element positive definite; items (iv)
and (v) show that neither ingredient can be dropped. Item (iii) applies at
every setting whose operator satisfies $\lmax(B^{(t)})\le\lmax(A_\Sigma)$;
the unconditioned receiver has $B=A_\Sigma$ and $\gamma=1$
(Corollary~\ref{cor:unconditioned-family}), and a submodule that constructs
a setting family states whether its operators do.
}
\begin{proposition}[Loss and dark counts make every dark-composed element definite]\label{prop:noise-floor}
Let $\{\widetilde M_{S|t}\}$ and $\{\widetilde F_{v|t}\}$ be the
\termref{term-dark-composed-family}{dark-composed family} (Definition~\ref{def:dark-composed}) of a
\termref{term-setting-family}{setting family} $\big(\Tset,\{M_{S|t}\}\big)$ with operators $B^{(t)}$
and prefactors $\gamma^{(t)}$ (Definition~\ref{def:setting-family}), a
\termref{term-dark-count-channel}{dark-count channel} $R$ and a declared coarse-graining $\cg$, and write
$r_S\coloneqq R_{S|\varnothing}$.
\begin{enumerate}[label=(\roman*),leftmargin=3.0em]
\item \emph{The floor.} For every $t$, $S$, $v$ and $n\ge0$,
  \begin{equation}\label{eq:noise-floor}
    \blk{\widetilde M}{\eqn}_{S|t}\succeq r_S\,\gamma^{(t)}
      \big(1-\lmax(B^{(t)})\big)^{n}\,\id,
    \qquad
    \blk{\widetilde F}{\eqn}_{v|t}\succeq R_{\cg^{-1}(v)|\varnothing}\,
      \gamma^{(t)}\big(1-\lmax(B^{(t)})\big)^{n}\,\id .
  \end{equation}
  Hence, if $r_S>0$ for every $S\subseteq[K]$ and $\lmax(B^{(t)})<1$, every
  $\widetilde M_{S|t}$ and every $\widetilde F_{v|t}$ is positive definite on
  every photon-number block.
\item \emph{Dark counts in product form.} On the generated dark-count route,
  $r_S=\prod_{k\in S}p_{\dc|k}\prod_{k\notin S}(1-p_{\dc|k})$, and $r_S>0$
  for every $S$ if and only if $0<p_{\dc|k}<1$ for every $k$.
\item \emph{Loss through the efficiencies.} If the \termref{term-instrument-matrix}{instrument matrix} of a
  receiver instrument is $D_\eta\Gmat'$, with $\Gmat'$ a contraction and
  $D_\eta$ the \termref{term-terminal-detector-efficiency-layer}{terminal detector-efficiency layer}, then
  $\lmax(A_\Sigma)\le\max_k\eta_k$. For the unconditioned receiver, with
  $\InpS=\varnothing$, $B=A_\Sigma$ and $\gamma=1$, so $\eta_k<1$ and
  $0<p_{\dc|k}<1$ for every $k$ give
  $\lmin\big(\blk{\widetilde M}{\eqn}_S\big)\ge r_S\,(1-\max_k\eta_k)^{n}>0$.
\item \emph{Loss is necessary.} If $\lmax(B^{(t)})=1$, then for every
  dark-count channel and every $n\ge1$ the element
  $\blk{\widetilde M}{\eqn}_{\varnothing|t}=r_\varnothing\,\blk{M}{\eqn}_{\varnothing|t}$
  is singular; if $B^{(t)}=\id$ it is zero.
\item \emph{Dark counts in every port are necessary.} For the unconditioned
  receiver on the generated route, let
  $Z\coloneqq\{k:p_{\dc|k}=0\}$. Then $\blk{\widetilde M}{\eqn}_S=0$ for
  every $S$ with $|S\cap Z|>n$; in particular
  $\blk{\widetilde M}{(=0)}_{\{k\}}=0$ for every $k\in Z$. If $p_{\dc|k}=1$,
  then $\widetilde M_S=0$ for every $S\not\ni k$.
\item \emph{Equivalence.} For the unconditioned receiver on the generated
  route, every $\widetilde M_S$ is positive definite on every
  photon-number block if and only if $\lmax(A_\Sigma)<1$ and
  $0<p_{\dc|k}<1$ for every $k$.
\end{enumerate}
\end{proposition}
\claimmeta{\metalabel{Uses}\mref{Def.}{def:fock}, \mref{Lem.}{lem:fock-functor}, \mref{Lem.}{lem:fock-block-spectrum}, \mref{Def.}{def:instrument}, \mref{Def.}{def:physmodel}, \mref{Thm.}{thm:port-grams}, \mref{Thm.}{thm:click-povm}, \mref{Def.}{def:coarse}, \mref{Thm.}{thm:multiclick-bound}, \mref{Def.}{def:setting-family}, \mref{Thm.}{thm:efficiency}, \mref{Def.}{def:darkadm}, \mref{Def.}{def:dark-composed}
\metasep \metalabel{Used by}\mref{Rem.}{rem:noise-floor-receiver}, \mref{Rem.}{rem:normalization-gain}, \mref{Rem.}{rem:normalization-scope}, \mref{Req.}{req:dark-composed-reference}, \mref{Req.}{req:snap-fix}
\metasep newly derived
\metasep accepted}

\begin{proof}
(i) By \eqref{eq:family-elements}, $\Pvac{[K]}^{(t)}=M_{\varnothing|t}$, so
\eqref{eq:family-prefactor} gives
$M_{\varnothing|t}=\gamma^{(t)}\Fock(\id-B^{(t)})$, and
$0\preceq B^{(t)}\preceq\id$ by \eqref{eq:bbkhyp}. Lemma~\ref{lem:fock-block-spectrum}
gives $\blk{M}{\eqn}_{\varnothing|t}\succeq\gamma^{(t)}(1-\lmax(B^{(t)}))^{n}\id$.
Every term of $\widetilde M_{S|t}=\sum_{S'}R_{S|S'}M_{S'|t}$ is positive
semidefinite, so dropping all but $S'=\varnothing$ gives
$\widetilde M_{S|t}\succeq r_SM_{\varnothing|t}$, and the blocks inherit the
inequality because every term is block diagonal. Summing over the fibre of
$v$ gives the second inequality, with
$\sum_{S\in\cg^{-1}(v)}r_S=R_{\cg^{-1}(v)|\varnothing}$. If every $r_S>0$ and
$\lmax(B^{(t)})<1$, the right-hand sides are positive multiples of $\id$:
$\gamma^{(t)}>0$ by Definition~\ref{def:setting-family}, and every fibre is
non-empty because $\cg$ is onto (Definition~\ref{def:coarse}).

(ii) \eqref{eq:darkR} at $S'=\varnothing$ is the stated product. It is positive
for $S=[K]$ only if every $p_{\dc|k}>0$, and for $S=\varnothing$ only if every
$p_{\dc|k}<1$; conversely, under both, every factor is positive.

(iii) By Theorem~\ref{thm:efficiency}\,(i), $A_k=\eta_kA'_k$ with $A'_k$ the
port Grams of $\Gmat'$, and $\sum_kA'_k\preceq\id$ by
Theorem~\ref{thm:port-grams}\,(i) applied to $\Gmat'$. Hence
$A_\Sigma=\sum_k\eta_kA'_k\preceq(\max_k\eta_k)\sum_kA'_k\preceq(\max_k\eta_k)\id$,
each $A'_k$ being positive semidefinite. For $\InpS=\varnothing$,
$M_\varnothing=\Fock(\id-A_\Sigma)$ by \eqref{eq:Mnc} of
Theorem~\ref{thm:click-povm}, which is \eqref{eq:family-prefactor} with
$\gamma=1$ and $B=A_\Sigma$; the argument of (i), which uses only this form of
$M_\varnothing$ and $0\preceq A_\Sigma\preceq\id$
(Theorem~\ref{thm:port-grams}\,(i)), gives the bound, and (ii) makes $r_S$
positive.

(iv) By (D2), $R_{\varnothing|S'}>0$ forces $S'\subseteq\varnothing$, so
$\widetilde M_{\varnothing|t}=r_\varnothing M_{\varnothing|t}
=r_\varnothing\gamma^{(t)}\Fock(\id-B^{(t)})$. By
Lemma~\ref{lem:fock-block-spectrum} the least eigenvalue of its $n$-photon block
is $r_\varnothing\gamma^{(t)}(1-\lmax(B^{(t)}))^{n}=0$ for $n\ge1$. If
$B^{(t)}=\id$, the block is $r_\varnothing\gamma^{(t)}(0^{\otimes n})|_{\Symn{n}}=0$
for $n\ge1$.

(v) \emph{Step 1 (no more clicks than photons).} For $\InpS=\varnothing$,
$\blk{M}{\eqn}_{S'}=0$ whenever $|S'|>n$. Indeed
$M_{S'}=V^\dagger X V$ with
$X=\prod_{k\in S'}(\id-\Pi^0_k)\prod_{k\notin S'}\Pi^0_k$
(Theorem~\ref{thm:click-povm}, \eqref{eq:MSdef}). The \termref{term-measurement-isometry}{measurement isometry}
$V=\Fock(U)\Fock(\iota)$ of Definition~\ref{def:physmodel} maps
$\Symn{n}(\mathbb{C}^{\Inp})$ into the $n$-photon sector of
$\mathcal{F}(\mathbb{C}^{\Modes\sqcup\Env})$, because both lifts preserve the
photon number (Lemma~\ref{lem:fock-functor}\,(i) and Definition~\ref{def:fock}).
Each $\Pi^0_k$ is the projector onto the span of the occupation vectors with
no photon in $\Port{k}$ (Lemma~\ref{lem:fock-functor}\,(iv)), so all of them
are diagonal in the occupation basis and $\prod_{k\in S'}(\id-\Pi^0_k)$ is the
projector onto the span of the occupation vectors with at least one photon in
every port of $S'$. The ports are pairwise disjoint
(Definition~\ref{def:instrument}\,(iii)), so such a vector carries at least
$|S'|$ photons, and $X$ vanishes on the $n$-photon sector when $|S'|>n$. Hence
$M_{S'}$ vanishes on $\Symn{n}(\mathbb{C}^{\Inp})$.

\emph{Step 2.} By \eqref{eq:darkR}, $R_{S|S'}\neq0$ requires $S'\subseteq S$
and $p_{\dc|k}\neq0$ for every $k\in S\setminus S'$, that is,
$S\cap Z\subseteq S'$. If $|S\cap Z|>n$, every $S'$ with $R_{S|S'}\neq0$ has
$|S'|>n$, and Step 1 gives
$\blk{\widetilde M}{\eqn}_S=\sum_{S'}R_{S|S'}\blk{M}{\eqn}_{S'}=0$. For
$k\in Z$ and $S=\{k\}$, $|S\cap Z|=1>0$. If $p_{\dc|k}=1$, every entry
$R_{S|S'}$ with $k\notin S$ carries the factor $1-p_{\dc|k}=0$.

(vi) Sufficiency is (i) with (ii), at $\gamma=1$ and $B=A_\Sigma$ as in (iii).
Necessity: if $\lmax(A_\Sigma)=1$, (iv) gives a singular block of
$\widetilde M_\varnothing$ at $n=1$; if some $p_{\dc|k}=0$, (v) gives
$\blk{\widetilde M}{(=0)}_{\{k\}}=0$; if some $p_{\dc|k}=1$, (v) gives
$\widetilde M_\varnothing=0$.
\end{proof}

% claim receiver/detector-imperfections/rem:noise-floor-receiver
\claimsummary{receiver/detector-imperfections/rem:noise-floor-receiver}{Declare the loss and dark counts the device surely has, in every port; then every element is positive definite and no snap revives a zero row. A run that uses no enclosed reference rests, at a cutoff of two or more, on its snap; an angle decides whether the reference can be exact.}
\claimbookmark{receiver/detector-imperfections/rem:noise-floor-receiver}{Remark}{rem:noise-floor-receiver}{What a receiver declaration needs for a rate that rests on its model}
\claimleadin{%
The remark ends with reductions that apply a channel to the model before
comparing it with a lattice family; they are planned beyond these contracts
(\href{../../../spec/future-scope.md}{\nolinkurl{spec/future-scope.md}}).
}
\begin{remark}[What a receiver declaration needs for a rate that rests on its model]\label{rem:noise-floor-receiver}
The emitted blocks of a receiver are snapped onto a lattice before a key
rate is computed, and the dark-composed reference
(Corollary~\ref{cor:exact-dark-composed}), enclosed for any declared
parameters and exact where its radii are zero, is handed beside the snap.
Where the run's first-order bound is evaluated on that reference and its key
length is formed from a number on it, the blocks carrying it are
\exttermref{../foundations/contract.pdf}{term-evaluated-off-its-snap}{evaluated off its snap} (\xref{foundations/def:evaluated-off-snap}) and
the key rate rests on the declared model, not on the snap. Where the key
length is formed instead from a number certified on the snapped declaration
(\xref{meat-core/frank-wolfe/req:certified-evaluation}\,(i)), the moved
ones among those blocks rest on the snapped-measurement premise, as every
other moved block does
(\xref{foundations/req:snap-reference}\,(iii)).
This remark says what that means for a receiver declaration.

\emph{What to declare.} Model the loss and the dark counts the device has,
in every port: an efficiency $\eta_k<1$ (or any loss that leaves
$\lmax(A_\Sigma)<1$) and a dark-count probability $0<p_{\dc|k}<1$. By
Proposition~\ref{prop:noise-floor} every element of the dark-composed family
is then positive definite on every photon-number block. At a setting $t$
whose operator and prefactor in the setting family are $B^{(t)}$ and
$\gamma^{(t)}$, the least eigenvalue on block $n$ is at least
$r_S\gamma^{(t)}\big(1-\lmax(B^{(t)})\big)^{n}$ for a click pattern $S$, and
at least the sum of these over a fibre for an emitted outcome
(Proposition~\ref{prop:noise-floor}\,(i)). For the unconditioned receiver,
$B=A_\Sigma$ and $\gamma=1$, so the bound is
$r_S\big(1-\lmax(A_\Sigma)\big)^{n}$, which is at least
$r_S(1-\max_k\eta_k)^{n}$ when the loss sits in the terminal efficiency layer
(Proposition~\ref{prop:noise-floor}\,(iii)); a setting whose operator has
$\lmax(B^{(t)})\le\lmax(A_\Sigma)$ and whose prefactor $\gamma^{(t)}$ is below
one has at least $\gamma^{(t)}$ times that floor. No row of such a block is
structurally zero, so even the last resort of
\xref{foundations/req:support-safe-snap}\,(iv) gives no element of it
support that the data deny. The same floor is what makes domination by such
a model possible at a small weight
(\xref{foundations/rem:noise-floor}); no run relates a snap to its model
that way. The rates to declare are the smallest ones the
user is sure the device has. A larger rate gives a larger floor, but a model
noisier than the device is a wrong model, which nothing in the engine can
detect.

\emph{The enclosure.} An enclosed reference is evaluated at its centres,
which displaces the measured family by up to the enclosure's radius
(\xref{meat-core/frank-wolfe/rem:fo-arithmetic}). The radii of the
\termref{term-square-root-enclosure}{square-root enclosures} and
\termref{term-angle-enclosure}{angle enclosures}
(Definition~\ref{def:complex-enclosure}) from which the blocks are computed
are amplified on the way to the radius of a block: through the permanents of
the Fock blocks, the inclusion--exclusion sums over click patterns and the
sums over a fibre of the coarse-graining. The radius of a block bears on a
run in two places. The companion's entries are relaxed outward by it, so a
larger radius loosens them, which never raises the bound formed on the
companion (\xref{reductions/lem:row-relaxation}); and the evaluated family
lies up to the radius from the model, a perturbation that no claim bounds
in its effect on the key rate (\xref{meat-core/frank-wolfe/rem:fo-arithmetic}).
No size of the radius is a threshold for either: the enclosures are chosen
as tight as their exact computation allows, and the record shows the
largest radius of every handed block and of the companion's datum.

\emph{A run that uses no enclosed reference.} A run may declare that enclosed
references are not used (\xref{foundations/req:snap-reference}\,(i)). Then
only the blocks whose reference is exact are handed, and every other moved
block rests on the snapped-measurement premise, in the group no usable
reference, unless the run declares the strict option
(\xref{foundations/req:snap-reference}\,(iv)), which refuses. The receiver's
reference is in general exact only on the \termref{term-exact-route}{exact route}, and there with
certainty only on the blocks of at most one photon
(Theorem~\ref{thm:exact-blocks}\,(iii)): on a block of two or more photons,
an entry between occupation vectors $\vec m,\vec m'$ with
$\vec m!\,\vec m'!$ not a square is in general a Gaussian rational times an
irrational factor such as $1/\sqrt2$, and is enclosed with a positive radius.
In general, therefore, every run with a cutoff $N\ge2$ that uses no enclosed
reference rests on its snap there, and so does every such run whose
declaration is off the exact route; its record and a report on it say so.

\emph{An ideal device beyond one photon.} Without loss the no-click element
vanishes on every block $n\ge1$, and without dark counts every click pattern
vanishes on the vacuum block (Proposition~\ref{prop:noise-floor}\,(iv),
(v)); so an ideal device has rank-deficient elements on every block, and a
snap of it may need its last resort, which the record flags. A receiver that
measures in a rotated basis can have some of them rank-deficient along
irrational directions: on the one-photon block where a direction of the
basis has no Gaussian-rational representative, and on blocks of more photons
even where every direction has one. In an ideal passive BB84 receiver, the
two-photon element in which only the detector of the diagonal polarization
clicks, behind a half-wave plate at $\pi/8$, is a multiple of the projector
onto $\tfrac12\ket{2,0}+\tfrac1{\sqrt2}\ket{1,1}+\tfrac12\ket{0,2}$ in the
occupation basis of the channel's two polarization modes, and that ray
contains no Gaussian-rational vector, although the one-photon directions of
the diagonal basis are rational. By
\xref{foundations/lem:rational-obstruction}, every weight at which the model
dominates a lattice family declared in this basis is then at least the
largest eigenvalue of that element, at every lattice exponent, which is $1/4$
for a balanced basis choice; that is why a declared device weight, and not a
snap, is where domination could serve. At a cutoff $N\le1$ on the
\termref{term-exact-route}{exact route} the blocks are known exactly
(Theorem~\ref{thm:exact-blocks}\,(iii)); where they are moreover dyadic
rationals of bounded denominator, as at a balanced basis choice with dyadic
efficiencies, they are lattice values and the snap keeps them. The
obstruction belongs to domination, which compares the model with a lattice
family element by element in the occupation basis. A reduction that first
maps the $m$-photon block of the two polarization modes isometrically into
the $m$-fold tensor power of $\mathbb{C}^{2}$ carries the ray above to a
vector with rational coordinates (\xref{foundations/rem:noise-floor}); for
an ideal device on these two modes whose port Grams are Gaussian-rational it
would give a target with entries in $\mathbb{Q}(i)$ at every cutoff, at
dimension $2^{m}$ instead of $m+1$ (with $d$ input modes, $d^{m}$ instead of
$\binom{d+m-1}{m}$). The images of the elements sum to the projector onto
the symmetric subspace, not to the identity, so such a target is complete
only after an outcome is added on the complement of that subspace. No claim
of these contracts makes that reduction.

\emph{Angles.} A half-wave plate at axis angle $\theta$ has entries
$\pm\cos2\theta$ and $\pm\sin2\theta$
(\eqref{eq:hwpqwp} of Definition~\ref{def:catalogue-entries}). Every angle
is enclosed, so the angle decides nothing about whether a reference exists;
it decides whether the reference can be exact. Behind an ideal polarizing
beamsplitter the plate's port Gram on the two polarization modes has the
diagonal entries $(1\pm\cos4\theta)/2$, up to a rational factor, and the exact
route needs Gaussian-rational Grams (Definition~\ref{def:exact-route}), so it
needs $\cos4\theta$ rational. At multiples of $\pi/8$,
$\cos4\theta\in\{0,\pm1\}$; at $\theta=\pi/24$, $\cos4\theta=\sqrt3/2$ and
the reference is enclosed, which the run hands and evaluates as it does an
exact one. The \termref{term-exact-route}{exact route} reads an angle given as a rational
coefficient of $\pi$ (Requirement~\ref{req:exact-route-decision}\,(ii)); an
angle given in radians is a dyadic rational, and a non-zero one is no
rational multiple of $\pi$, so a plate declared at the binary64 value nearest
$\pi/8$ is enclosed, while one declared as $\tfrac18\pi$ can be on the
exact route.
\end{remark}
\claimmeta{\metalabel{Uses}\mref{Def.}{def:complex-enclosure}, \mref{Def.}{def:exact-route}, \mref{Thm.}{thm:exact-blocks}, \mref{Def.}{def:dark-composed}, \mref{Cor.}{cor:exact-dark-composed}, \mref{Prop.}{prop:noise-floor}, \extref{../foundations/contract.pdf}{foundations/def:evaluated-off-snap}, \extref{../foundations/contract.pdf}{foundations/lem:rational-obstruction}, \extref{../foundations/contract.pdf}{foundations/rem:noise-floor}
\metasep newly derived
\metasep accepted}

\section{Efficiency normalization}\label{group-receiver/detector-imperfections/def:pure-loss}
\sectionintro{A common factor of the port Grams is a uniform loss that commutes with every passive element, so it can be moved in front of the receiver and counted as part of the channel. The section defines the \termref{term-pure-loss-channel}{pure-loss channel} and the \termref{term-normalized-station}{normalized station}, whose efficiencies are the receiver's divided by a \termref{term-factored-loss}{factored loss} t, and its main point is \hyperlink{claim-receiver/detector-imperfections/prop:efficiency-normalization}{Proposition~\ref*{prop:efficiency-normalization}}: every click pattern of the receiver, and every dark-count composition of them, is the pull-back of the normalized station's through the pure-loss channel of transmittance t. The remarks show that at low efficiency the normalization removes a squared-efficiency factor from the two-photon multiclick block, and say what it leaves in place, efficiency mismatch included.}

% claim receiver/detector-imperfections/def:pure-loss
\claimsummary{receiver/detector-imperfections/def:pure-loss}{The pure-loss channel of transmittance t sends every channel mode through a beamsplitter of transmittance t whose other input is vacuum and discards the reflected light. It is the uniform loss a normalized station hands to the channel.}
\claimbookmark{receiver/detector-imperfections/def:pure-loss}{Definition}{def:pure-loss}{Pure-loss channel on the channel modes}
\claimleadin{%
Uniform loss in front of a receiver acts on the light the adversary sends,
before any optical element. Writing it as a channel on the channel system lets
Proposition~\ref{prop:efficiency-normalization} hand it to the adversary.
}
\begin{definition}[Pure-loss channel on the channel modes]\label{def:pure-loss}
Let $\InpC$ be the \termref{term-channel-modes}{channel modes} of a \termref{term-receiver-instrument}{receiver instrument}, so
that the channel system is $\Hs{\Chan}=\mathcal{F}(\mathbb{C}^{\InpC})$
(Definition~\ref{def:instrument}), let $\InpC'$ be a disjoint copy of
$\InpC$, with $m\mapsto m'$ the copying bijection, and let $t\in[0,1]$. Put
\begin{equation}\label{eq:lossiso}
  w_t\colon\mathbb{C}^{\InpC}\to\mathbb{C}^{\InpC\sqcup\InpC'},\qquad
  w_t\ket{m}\coloneqq\sqrt{t}\,\ket{m}+\sqrt{1-t}\,\ket{m'} ,
\end{equation}
and identify $\mathcal{F}(\mathbb{C}^{\InpC\sqcup\InpC'})$ with
$\mathcal{F}(\mathbb{C}^{\InpC})\otimes\mathcal{F}(\mathbb{C}^{\InpC'})$ by
the occupation identification of Lemma~\ref{lem:fock-maps}\,(ii),
which is unitary. The
\termdef{term-pure-loss-channel}{pure-loss channel}
of transmittance $t$ is
\begin{equation}\label{eq:pureloss}
  \Lambda^{\mathrm{loss}}_t(\rho)\;\coloneqq\;
  \tr_{\mathcal{F}(\mathbb{C}^{\InpC'})}\!\big[\Fock(w_t)\,\rho\,\Fock(w_t)^{\dagger}\big],
  \qquad \rho\in\Tcl{\Chan},
\end{equation}
and its adjoint is
\begin{equation}\label{eq:purelossadj}
  (\Lambda^{\mathrm{loss}}_t)^{\dagger}(Y)\;\coloneqq\;
  \Fock(w_t)^{\dagger}\big(Y\otimes\idon{\mathcal{F}(\mathbb{C}^{\InpC'})}\big)\Fock(w_t),
  \qquad Y\in\Lin{\Chan}.
\end{equation}
Since $w_t^{\dagger}w_t=t\,\id+(1-t)\,\id=\id$, the map $\Fock(w_t)$ is an
isometry (Lemma~\ref{lem:fock-maps}\,(i)); so $\Lambda^{\mathrm{loss}}_t$ is completely positive and trace
preserving, $(\Lambda^{\mathrm{loss}}_t)^{\dagger}$ is completely positive and
unital, and
$\tr[Y\,\Lambda^{\mathrm{loss}}_t(\rho)]=\tr[(\Lambda^{\mathrm{loss}}_t)^{\dagger}(Y)\,\rho]$
for every bounded $Y$ and trace-class $\rho$, by the defining property of the
partial trace and cyclicity of the trace. At $t=1$ the channel is the
identity.
\end{definition}
\claimmeta{\metalabel{Uses}\mref{Def.}{def:fock}, \mref{Lem.}{lem:fock-maps}, \mref{Def.}{def:instrument}
\metasep \metalabel{Used by}\mref{Lem.}{lem:pure-loss-lift}, \mref{Prop.}{prop:efficiency-normalization}, \mref{Rem.}{rem:normalization-gain}, \mref{Req.}{req:loss-before-evaluation}, \extref{../reductions/contract.pdf}{reductions/wpfss/cor:normalized-flat}
\metasep newly derived
\metasep accepted}

% claim receiver/detector-imperfections/def:normalized-station
\claimsummary{receiver/detector-imperfections/def:normalized-station}{A normalized station at factor t divides every detector efficiency of a probe-free receiver by t, for t between the largest efficiency and one; t is the factored loss. Its port Grams are those of the receiver divided by t.}
\claimbookmark{receiver/detector-imperfections/def:normalized-station}{Definition}{def:normalized-station}{Normalized station and factored loss}
\claimleadin{%
The efficiencies enter a receiver's port Grams as one scalar per port
(Theorem~\ref{thm:efficiency}\,(i)), on each of the three input routes: as
the \termref{term-terminal-detector-efficiency-layer}{terminal detector-efficiency layer} of an \termref{term-instrument-matrix}{instrument matrix}, or as
the scalars a declared Gram family is multiplied by. The definition divides
those scalars by a common factor and keeps everything else.
}
\begin{definition}[Normalized station and factored loss]\label{def:normalized-station}
Let $\Recv$ be a \termref{term-receiver-instrument}{receiver instrument} without
\termref{term-declared-source-modes}{declared source modes}, $\InpS=\varnothing$, with $K$ ports, whose
\termref{term-port-gram-matrices}{port Grams} are
\begin{equation}\label{eq:efficiencygrams}
  A_k=\eta_k\,A^{\circ}_k,\qquad k=1,\dots,K,
\end{equation}
for an \termref{term-admissible-gram-family}{admissible Gram family} $\{A^{\circ}_k\}$ on
$\mathbb{C}^{\InpC}$ and efficiencies $\eta_k\in[0,1]$ with
$\eta_{\max}\coloneqq\max_k\eta_k>0$. This is the case for an instrument
matrix $D_\eta\Gmat_0$, $\Gmat_0$ a contraction, with $A^{\circ}_k$ the port
Grams of $\Gmat_0$ (Theorem~\ref{thm:efficiency}\,(i)), and for a declared
Gram family $\{A^{\circ}_k\}$ with declared efficiencies. For
$t\in[\eta_{\max},1]$, a
\termdef{term-normalized-station}{normalized station}
of $\Recv$ at the
\termdef{term-factored-loss}{factored loss}
$t$ is a receiver instrument $\Recv'$ without declared source modes, with the
channel modes and the $K$ ports of $\Recv$, whose port Grams are
\begin{equation}\label{eq:normalizedgrams}
  A'_k=\frac{\eta_k}{t}\,A^{\circ}_k ,\qquad\text{so that}\qquad A_k=t\,A'_k .
\end{equation}
One exists: $\eta_k/t\in[0,1]$, so $0\preceq A'_k\preceq A^{\circ}_k$ and
$\sum_kA'_k\preceq\sum_kA^{\circ}_k\preceq\id$, and an admissible family is
the port-Gram family of some receiver instrument
(Theorem~\ref{thm:representation-independence}\,(ii)); for the instrument
matrix $D_\eta\Gmat_0$ one is $D_{\eta/t}\Gmat_0$, a product of two
contractions (Proposition~\ref{prop:mode-algebra}\,(ii)), with
$D_{\eta/t}$ the terminal layer of Definition~\ref{def:compile} at the
efficiencies $\eta_k/t$. Its click-pattern POVM depends on the port Grams
alone (Theorem~\ref{thm:representation-independence}\,(i)), so it is one
POVM for every choice. The choice $t=\eta_{\max}$ is the \emph{full}
normalization, at which the largest efficiency becomes one, and $t=1$ gives
$\Recv'$ the port Grams of $\Recv$.
\end{definition}
\claimmeta{\metalabel{Uses}\mref{Prop.}{prop:mode-algebra}, \mref{Def.}{def:instrument}, \mref{Def.}{def:gram}, \mref{Def.}{def:compile}, \mref{Thm.}{thm:representation-independence}, \mref{Thm.}{thm:efficiency}
\metasep \metalabel{Used by}\mref{Prop.}{prop:efficiency-normalization}, \mref{Rem.}{rem:normalization-gain}, \mref{Rem.}{rem:normalization-scope}, \mref{Req.}{req:efficiency-normalization}, \mref{Req.}{req:normalization-record}, \mref{Req.}{req:loss-before-evaluation}, \extref{../reductions/contract.pdf}{reductions/wpfss/cor:normalized-flat}
\metasep newly derived
\metasep implemented (effective: accepted)
\metasep \metalabel{Code}\coderef{https://github.com/fraunhofersg/certiqs-key/blob/v0.1.0/src/certiqs_key/receiver/request.py}{certiqs_key.receiver.request:normalized_station}}

% claim receiver/detector-imperfections/lem:pure-loss-lift
\claimsummary{receiver/detector-imperfections/lem:pure-loss-lift}{The adjoint of the pure-loss channel of transmittance t takes the lift of a contraction X to the lift of tX + (1 \ensuremath{-} t)\ensuremath{\cdot}1. Every joint no-click element is such a lift, so loss in front of a receiver acts on its Grams as a scalar.}
\claimbookmark{receiver/detector-imperfections/lem:pure-loss-lift}{Lemma}{lem:pure-loss-lift}{The pure-loss channel on the lift of a contraction}
\claimleadin{%
The \termref{term-joint-no-click-elements}{joint no-click elements} of a receiver are lifts $\Fock(\id-A_J)$ of
contractions (Theorem~\ref{thm:click-povm}\,(ii)), so the following identity
is what Proposition~\ref{prop:efficiency-normalization} needs from the
pure-loss channel.
}
\begin{technicalclaim}
\begin{lemma}[The pure-loss channel on the lift of a contraction]\label{lem:pure-loss-lift}
Let $t\in[0,1]$ and let $X$ be an operator on $\mathbb{C}^{\InpC}$ with
$X^{\dagger}X\preceq\id$. Then
\begin{equation}\label{eq:purelosslift}
  (\Lambda^{\mathrm{loss}}_t)^{\dagger}\big(\Fock(X)\big)
  \;=\;\Fock\big(t\,X+(1-t)\,\id\big).
\end{equation}
\end{lemma}
\end{technicalclaim}
\claimmeta{\metalabel{Uses}\mref{Lem.}{lem:fock-maps}, \mref{Def.}{def:pure-loss}
\metasep \metalabel{Used by}\mref{Prop.}{prop:efficiency-normalization}
\metasep newly derived
\metasep accepted}

\begin{proof}
\emph{Step 1 (the adjoint).} By Lemma~\ref{lem:fock-maps}\,(iii) with
$X_2=\id$, under the identification of Definition~\ref{def:pure-loss},
$\Fock(X)\otimes\idon{\mathcal{F}(\mathbb{C}^{\InpC'})}=\Fock(X\oplus\id)$,
an operator of norm at most one (Lemma~\ref{lem:fock-maps}\,(i)). So
\eqref{eq:purelossadj} and Lemma~\ref{lem:fock-maps}\,(i) give
\begin{equation*}
  (\Lambda^{\mathrm{loss}}_t)^{\dagger}\big(\Fock(X)\big)
  =\Fock(w_t)^{\dagger}\,\Fock(X\oplus\id)\,\Fock(w_t)
  =\Fock\big(w_t^{\dagger}(X\oplus\id)w_t\big).
\end{equation*}

\emph{Step 2 (the compression).} For $v\in\mathbb{C}^{\InpC}$ write $v'$ for
its copy in $\mathbb{C}^{\InpC'}$. Then $w_tv=\sqrt t\,v+\sqrt{1-t}\,v'$,
$(X\oplus\id)w_tv=\sqrt t\,Xv+\sqrt{1-t}\,v'$, and $w_t^{\dagger}$ maps
$a+b'$ to $\sqrt t\,a+\sqrt{1-t}\,b$, so
$w_t^{\dagger}(X\oplus\id)w_t=t\,X+(1-t)\,\id$.
\end{proof}

% claim receiver/detector-imperfections/prop:efficiency-normalization
\claimsummary{receiver/detector-imperfections/prop:efficiency-normalization}{Every click pattern of a probe-free receiver is the pull-back, through the pure-loss channel of the factored loss t, of the same pattern of its normalized station, and so is every dark-count composition. The statistics agree, and the loss can be handed to the adversary.}
\claimbookmark{receiver/detector-imperfections/prop:efficiency-normalization}{Proposition}{prop:efficiency-normalization}{A receiver as uniform loss followed by its normalized station}
\claimleadin{%
The idea is that a scalar factor common to all port Grams is a loss that
commutes with every passive element, so it can be moved to the front of the
receiver, where it acts on the channel modes alone.
}
\begin{proposition}[A receiver as uniform loss followed by its normalized station]\label{prop:efficiency-normalization}
Let $\Recv$ be a \termref{term-receiver-instrument}{receiver instrument} without
\termref{term-declared-source-modes}{declared source modes}, with efficiencies $\eta_k$,
$\eta_{\max}>0$, as in Definition~\ref{def:normalized-station}, let
$t\in[\eta_{\max},1]$ and let $\Recv'$ be a \termref{term-normalized-station}{normalized station} of
$\Recv$ at the \termref{term-factored-loss}{factored loss} $t$. Let $\{M_S\}$ and $\{M'_S\}$,
$S\subseteq[K]$, be their click-pattern POVMs on
$\Hs{\Chan}=\mathcal{F}(\mathbb{C}^{\InpC})$ (Theorem~\ref{thm:click-povm}).
\begin{enumerate}[label=(\roman*),leftmargin=3.0em]
\item \emph{The click patterns.} For every $S\subseteq[K]$,
  \begin{equation}\label{eq:normalization}
    M_S\;=\;(\Lambda^{\mathrm{loss}}_t)^{\dagger}\big(M'_S\big),
    \qquad\text{equivalently}\qquad
    \tr[M_S\,\rho]=\tr\big[M'_S\,\Lambda^{\mathrm{loss}}_t(\rho)\big]
    \ \ \text{for every }\rho\in\Tcl{\Chan}.
  \end{equation}
\item \emph{Classical post-processing, dark counts included.} For every
  column-stochastic matrix $R$ on $2^{[K]}$, in particular every
  \termref{term-dark-count-channel}{dark-count channel}, and every \termref{term-declared-coarse-graining}{declared coarse-graining} $\cg$
  onto $\Vset$,
  \begin{equation}\label{eq:normalizationdark}
    \sum_{S'}R_{S|S'}M_{S'}=(\Lambda^{\mathrm{loss}}_t)^{\dagger}
    \Big(\sum_{S'}R_{S|S'}M'_{S'}\Big),
    \qquad
    \sum_{S'}R_{\cg^{-1}(v)|S'}M_{S'}=(\Lambda^{\mathrm{loss}}_t)^{\dagger}
    \Big(\sum_{S'}R_{\cg^{-1}(v)|S'}M'_{S'}\Big),
  \end{equation}
  so the \termref{term-dark-composed-family}{dark-composed family} of $\Recv$ with $R$ and $\cg$
  (Definition~\ref{def:dark-composed}) is the pull-back of that of $\Recv'$
  with the same $R$ and $\cg$: the dark counts act after the factored loss.
\item \emph{The device.} Under the receiver's device assumptions,
  Assumptions~\ref{ass:trust}, \ref{ass:passive}, \ref{ass:vac},
  \ref{ass:modes}, \ref{ass:thresh}, \ref{ass:port}, \ref{ass:exact},
  \ref{ass:nosel} and~\ref{ass:round}, $\{M_S\}$ is the measurement the
  physical receiver makes on the light it receives in its channel modes, and,
  by Assumptions~\ref{ass:modes} and~\ref{ass:port}, every port responds
  alike to every mode through which light reaches it; the physical receiver,
  followed by any dark-count channel, therefore acts on that light exactly as
  $\Lambda^{\mathrm{loss}}_t$ followed by $\Recv'$ and the same dark-count
  channel.
\end{enumerate}
\end{proposition}
\claimreading{If your receiver has no probe source, the largest of its detector efficiencies can be treated as uniform loss in front of the whole receiver, followed by a receiver whose best detector is perfect and whose other detectors keep their efficiencies relative to it. The dark counts stay at the detectors. Both descriptions give the same click probabilities for every input, so the loss in front can be counted as part of the channel the adversary controls. A detector efficiency that depends on time slot, wavelength or spatial mode is covered if declared as mode-dependent loss in front of the detector, which stays in place; a dependence on a degree of freedom the adversary can use but your declaration does not carry is efficiency mismatch, which this does not address.}
\claimmeta{\metalabel{Uses}\mref{Def.}{def:gram}, \mref{Def.}{def:physmodel}, \mref{Thm.}{thm:port-grams}, \mref{Thm.}{thm:click-povm}, \mref{Def.}{def:coarse}, \mref{Def.}{def:darkadm}, \mref{Def.}{def:dark-composed}, \mref{Def.}{def:pure-loss}, \mref{Def.}{def:normalized-station}, \mref{Lem.}{lem:pure-loss-lift}
\metasep \metalabel{Assumes}\aref{A-TRUST}{ass:trust}, \aref{A-PASSIVE}{ass:passive}, \aref{A-VAC}{ass:vac}, \aref{A-MODES}{ass:modes}, \aref{A-THRESH}{ass:thresh}, \aref{A-PORT}{ass:port}, \aref{A-EXACT}{ass:exact}, \aref{A-NOSEL}{ass:nosel}, \aref{A-ROUND}{ass:round}
\metasep \metalabel{Used by}\mref{Rem.}{rem:normalization-gain}, \mref{Rem.}{rem:normalization-scope}, \mref{Req.}{req:efficiency-normalization}, \mref{Req.}{req:normalization-record}, \mref{Req.}{req:loss-before-evaluation}, \extref{../reductions/contract.pdf}{reductions/wpfss/cor:normalized-flat}, \extref{../reductions/contract.pdf}{reductions/wpfss/rem:receiver-imports}
\metasep newly derived
\metasep accepted}

\begin{proof}
(i) By \eqref{eq:normalizedgrams}, $A_k=t\,A'_k$ for every $k$, so
$A_J=t\,A'_J$ for every $J\subseteq[K]$, in the notation of
Definition~\ref{def:gram}. The family $\{A'_k\}$ is the port-Gram family of a
receiver instrument, hence admissible (Theorem~\ref{thm:port-grams}\,(i)), so
$0\preceq A'_J\preceq\id$ and $X_J\coloneqq\id-A'_J$ satisfies
$0\preceq X_J\preceq\id$, in particular $X_J^{\dagger}X_J\preceq\id$. By
Theorem~\ref{thm:click-povm}\,(ii) and Lemma~\ref{lem:pure-loss-lift},
\begin{equation*}
  \Pvac{J}=\Fock(\id-A_J)=\Fock\big(t\,X_J+(1-t)\,\id\big)
  =(\Lambda^{\mathrm{loss}}_t)^{\dagger}\big(\Fock(X_J)\big)
  =(\Lambda^{\mathrm{loss}}_t)^{\dagger}\big(\Pvac{J}'\big),
\end{equation*}
with $\Pvac{J}'=\Fock(\id-A'_J)$ the joint no-click element of $\Recv'$.
Each $M_S$ is the same signed sum \eqref{eq:incexc} of the elements
$\Pvac{([K]\setminus S)\cup T}$ as $M'_S$ is of the elements
$\Pvac{([K]\setminus S)\cup T}'$, and $(\Lambda^{\mathrm{loss}}_t)^{\dagger}$
is linear, which gives the first identity of \eqref{eq:normalization}. The
second is the first paired with $\rho$, by the pairing of
Definition~\ref{def:pure-loss}.

(ii) Both sides of each identity are finite linear combinations with the same
scalar coefficients, of $\{M_{S'}\}$ on the left and of $\{M'_{S'}\}$ inside
the adjoint on the right, so (i) and linearity give them; the fibre sums are
the coefficients $R_{\cg^{-1}(v)|S'}=\sum_{S\in\cg^{-1}(v)}R_{S|S'}$ of
Definition~\ref{def:dark-composed}. Neither column-stochasticity nor any other
property of $R$ (Definition~\ref{def:darkadm}) is used.

(iii) These are the device assumptions under which the click-pattern POVM
of a probe-free receiver is the physical receiver's measurement, the set
Corollary~\ref{cor:unconditioned-family} reads its \termref{term-setting-family}{setting family} under:
Assumption~\ref{ass:trust}\,(i) makes the declared instrument, efficiencies
included, the physical one; Assumption~\ref{ass:passive} makes its optics a
passive linear map; Assumption~\ref{ass:vac} puts vacuum on every undeclared
input, the probe-free receiver having no other declared input than its
channel modes; Assumptions~\ref{ass:thresh} and~\ref{ass:port} make each
detector a threshold detector on its port's total photon number;
Assumption~\ref{ass:exact} makes the declared parameters exact; and
Assumptions~\ref{ass:nosel} and~\ref{ass:round} fix the device
independently of the adversary and identically in every round. Under them the
measurement on the channel modes is the click-pattern POVM of $\Recv$ of
Definition~\ref{def:physmodel} and Theorem~\ref{thm:click-povm}.
Assumption~\ref{ass:modes}\,(i) says that every optical path from outside to
a detector is a declared mode, and Assumption~\ref{ass:port}\,(i) that a port
clicks exactly when its modes carry a photon, whichever mode carries it; so no
mode outside the declaration reaches a detector with an efficiency of its
own, and each port's efficiency is the scalar $\eta_k$ of
\eqref{eq:efficiencygrams}. The rest is (i) and (ii).
\end{proof}

% claim receiver/detector-imperfections/rem:normalization-gain
\claimsummary{receiver/detector-imperfections/rem:normalization-gain}{At a common efficiency \ensuremath{\eta} and no dark counts, the two-photon multiclick block of a receiver is \ensuremath{\eta}\ensuremath{^{2}} times that of its normalized station, so every multiclick weight it can emit at a cutoff of one photon is at most \ensuremath{\eta}\ensuremath{^{2}}. Normalization removes that factor without enlarging the adversary's reach.}
\claimbookmark{receiver/detector-imperfections/rem:normalization-gain}{Remark}{rem:normalization-gain}{The two-photon multiclick block scales with the squared efficiency}
\begin{remark}[The two-photon multiclick block scales with the squared efficiency]\label{rem:normalization-gain}
Let every port of a \termref{term-receiver-instrument}{receiver} without
\termref{term-declared-source-modes}{declared source modes} have the same efficiency $\eta\in(0,1]$, and let
$\Recv'$ be its \termref{term-normalized-station}{normalized station} at the full normalization
$t=\eta$, whose efficiencies are all one. Write $M_{\mc}$ and $M'_{\mc}$ for
the \termref{term-multiclick-element}{multiclick elements} \eqref{eq:Mmc} of the receiver and of $\Recv'$.
Without dark counts, $M'_{\mc}$ vanishes on the photon-number blocks $n=0$ and
$n=1$, a pattern of two or more clicks needing two photons
(Proposition~\ref{prop:noise-floor}\,(v) with every port dark-free). For
$\psi\in\Symn{2}(\mathbb{C}^{\InpC})$ the component of
$\Fock(w_t)\psi=w_t^{\otimes2}\psi$ with no photon in the copy $\InpC'$
(Definition~\ref{def:pure-loss}) is $t\,\psi\otimes\ket{\vac}$, and its other
components carry one or two photons in the copy and so at most one in
$\InpC$. The operator $M'_{\mc}\otimes\id$ acts as the identity on the copy,
so it does not couple components with different photon numbers there, and it
vanishes on the components with at most one photon in $\InpC$. By
Proposition~\ref{prop:efficiency-normalization}\,(i) and
\eqref{eq:purelossadj}, therefore,
$\bra{\psi}M_{\mc}\ket{\psi}=t^{2}\bra{\psi}M'_{\mc}\ket{\psi}$ for every
such $\psi$, and since both blocks are Hermitian,
\begin{equation}\label{eq:twophotonscaling}
  \blk{M}{(=2)}_{\mc}\;=\;\eta^{2}\,\blk{M'}{(=2)}_{\mc} .
\end{equation}
A multiclick weight $\pmc$ emitted for the receiver at the cutoff $N=1$
satisfies $\blk{M}{(=2)}_{\mc}\succeq\pmc\,\id$, the block $n=2$ lying above
the cutoff (\eqref{eq:emitted-unconditioned} of
Corollary~\ref{cor:emitted-unconditioned}), so
$\pmc\le\eta^{2}\lmin\big(\blk{M'}{(=2)}_{\mc}\big)\le\eta^{2}$, whatever
window or enclosure produced it; the normalized station's own weight is
computed at unit efficiency and carries no such factor. The normalization
does not buy this by enlarging what the adversary can do: by
Proposition~\ref{prop:efficiency-normalization}\,(i), every attack on the
receiver followed by $\Lambda^{\mathrm{loss}}_t$ is an attack on $\Recv'$
with the same statistics, so an exact analysis of the receiver over all
attacks gives a rate at least that of $\Recv'$, and what the normalization
improves is the multiclick weight a reduction to finitely many photons reads.
\end{remark}
\claimmeta{\metalabel{Uses}\mref{Thm.}{thm:click-povm}, \mref{Cor.}{cor:emitted-unconditioned}, \mref{Prop.}{prop:noise-floor}, \mref{Def.}{def:pure-loss}, \mref{Def.}{def:normalized-station}, \mref{Prop.}{prop:efficiency-normalization}
\metasep \metalabel{Used by}\mref{Req.}{req:efficiency-normalization}
\metasep newly derived
\metasep accepted}

% claim receiver/detector-imperfections/rem:normalization-scope
\claimsummary{receiver/detector-imperfections/rem:normalization-scope}{Only the per-port efficiencies are divided: declared mode-dependent loss stays, undeclared efficiency mismatch is outside the assumptions either way, and a probe-fed receiver is not covered. The full normalization can remove the loss the noise floor needs, which a larger factor keeps.}
\claimbookmark{receiver/detector-imperfections/rem:normalization-scope}{Remark}{rem:normalization-scope}{What the normalization factors out and what it leaves in place}
\begin{remark}[What the normalization factors out and what it leaves in place]\label{rem:normalization-scope}
\emph{Only the terminal scalars are divided.} The family $\{A^{\circ}_k\}$ of
Definition~\ref{def:normalized-station} is kept: every loss declared in the
element graph, a polarizer's extinction, and an attenuation that varies
within a port (Theorem~\ref{thm:efficiency}\,(ii)) stay in the normalized
station unchanged. A detector whose efficiency depends on a spatial, spectral
or temporal mode is declared as such a mode-dependent loss in front of its
port, and the normalization leaves it where it is.

\emph{Efficiency mismatch is a separate problem.} The identity of
Proposition~\ref{prop:efficiency-normalization}\,(i) holds of the declared
instrument. That it describes the physical receiver, clause (iii), rests on
Assumptions~\ref{ass:modes} and~\ref{ass:port}: an efficiency that differs
between degrees of freedom the adversary can address but the declaration
does not carry violates one of them, with or without the normalization, which
neither creates such a mismatch nor removes one.

\emph{A probe-fed receiver is not covered.} With declared source modes the
probe light passes the same detectors, so a uniform loss moved to the front
would act on the probe modes too and replace each probe state by its image
under the loss; no claim here states that form.

\emph{The choice of the factor.} The identity holds at every
$t\in[\eta_{\max},1]$. At the full normalization $t=\eta_{\max}$ with lossless
optics the normalized station can have $\lmax(A'_\Sigma)=1$, and then some
element of its dark-composed family is singular on every block with a photon,
dark counts or not (Proposition~\ref{prop:noise-floor}\,(iv)). At
$t>\eta_{\max}$ the largest efficiency of the station is $\eta_{\max}/t<1$,
so $\lmax(A'_\Sigma)\le\eta_{\max}/t$ and, with a dark count in every port,
every such element is positive definite with the floor of
Proposition~\ref{prop:noise-floor}\,(iii) at $\eta_{\max}/t$. Where
$\eta_{\max}<99/100$, the default factor
$t=\eta_{\max}/(99/100)$ of Requirement~\ref{req:efficiency-normalization}
is of this kind, with $\lmax(A'_\Sigma)\le99/100$; where
$\eta_{\max}\ge99/100$ the default is $t=1$, the declared receiver, whose
largest efficiency can itself be one. Which factor in $[\eta_{\max},1]$ a run
is better served by is a matter of the engine's evaluation, not of the
identity.

\emph{Only declared efficiencies are factored.} The proof of
Proposition~\ref{prop:efficiency-normalization}\,(i) uses only $A_k=t\,A'_k$
with $\{A'_k\}$ admissible, which holds for every
$t\in[\lmax(A_\Sigma),1]$ with $t>0$ and $A'_k=A_k/t$; such a factor can be smaller
than $\eta_{\max}$, but it is an eigenvalue rather than a declared number.
A Gram family declared from click data, with every efficiency declared as one
(Remark~\ref{rem:powervsclick}), has $\eta_{\max}=1$, and nothing is factored
from it.
\end{remark}
\claimmeta{\metalabel{Uses}\mref{Thm.}{thm:efficiency}, \mref{Prop.}{prop:noise-floor}, \mref{Def.}{def:normalized-station}, \mref{Prop.}{prop:efficiency-normalization}
\metasep \metalabel{Used by}\mref{Req.}{req:efficiency-normalization}
\metasep newly derived
\metasep accepted}

\section{Requirements}\label{group-receiver/detector-imperfections/req:efficiency-schema}
\sectionintro{The engine keeps detector efficiency and component loss in separate declared fields, and it records every declared transmittance and the efficiency convention of the declared numbers. It compiles a receiver without a probe at its \termref{term-normalized-station}{normalized station}, by default at the largest declared efficiency divided by 99/100, just short of the full normalization, or at one where that efficiency is at least 99/100, records the \termref{term-factored-loss}{factored loss}, and has the loss applied to every input on which the emitted measurement stands for the receiver's statistics. It refuses a request that declares a dark-count matrix while some per-port dark-count probability is non-zero, and it emits the \termref{term-dark-count-channel}{dark-count channel} with the scalars of the composed bound. The emitted blocks are the dark-composed ones and carry their reference, recorded block by block as exact, enclosed or absent; the snap itself runs only on the exact ones. A moved block of their snap that rests on the snap premise is recorded with its cause, the group of the cause and the fix for it; a block snapped by the last resort is flagged; only a run that declares the strict option refuses.}
\sectionimports{\extref{../foundations/contract.pdf}{foundations/cor:rank-over-factor-enclosure}, \extref{../foundations/contract.pdf}{foundations/def:evaluated-off-snap}, \extref{../foundations/contract.pdf}{foundations/def:exact-island}, \extref{../foundations/contract.pdf}{foundations/def:exact-reference}, \extref{../foundations/contract.pdf}{foundations/def:factor-enclosure}, \extref{../foundations/contract.pdf}{foundations/def:sector-map}, \extref{../foundations/contract.pdf}{foundations/def:snap}, \extref{../foundations/contract.pdf}{foundations/def:tiers}, \extref{../foundations/contract.pdf}{foundations/req:snap-reference}, \extref{../foundations/contract.pdf}{foundations/req:support-safe-snap}, \extref{../foundations/contract.pdf}{foundations/req:who-snaps}}

% claim receiver/detector-imperfections/req:efficiency-schema
\claimsummary{receiver/detector-imperfections/req:efficiency-schema}{The input schema keeps a detector efficiency and a C7 element in different fields, so a C7 element is never read as a detector efficiency. The two enter the compiled matrix at different sites (\hyperlink{claim-receiver/detector-imperfections/thm:efficiency}{Theorem~\ref*{thm:efficiency}} (iii)).}
\claimbookmark{receiver/detector-imperfections/req:efficiency-schema}{Requirement}{req:efficiency-schema}{Detector efficiency and component loss in separate fields}
\claimleadin{%
A detector efficiency enters the compiled matrix only through the terminal
factor $D_\eta$ of Definition~\ref{def:compile}
(Theorem~\ref{thm:efficiency}\,(iii)); a C7 element is a component of the graph,
applied at its own position in the chain.
}
\begin{requirement}[Detector efficiency and component loss in separate fields]\label{req:efficiency-schema}
A C7 element is not a detector efficiency, and the engine's input schema keeps
the two in different fields.
\end{requirement}
\claimmeta{\metalabel{Motivated by}\mref{Def.}{def:compile}, \mref{Thm.}{thm:efficiency}
\metasep implemented (effective: accepted)
\metasep \metalabel{Code}\coderef{https://github.com/fraunhofersg/certiqs-key/blob/v0.1.0/src/certiqs_key/receiver/request.py}{certiqs_key.receiver.request:DetectorPort}}

% claim receiver/detector-imperfections/req:efficiency-record
\claimsummary{receiver/detector-imperfections/req:efficiency-record}{The emitted efficiency record lists every declared loss parameter under its name, with the element that declares it and its position in the compiled product. Two elements with equal values are two losses and two rows.}
\claimbookmark{receiver/detector-imperfections/req:efficiency-record}{Requirement}{req:efficiency-record}{The efficiency record lists every declared loss parameter}
\claimleadin{%
Theorem~\ref{thm:efficiency}\,(iii) gives the detector efficiency one site in the
compiled matrix, and item~(iv) of that theorem states the record this
requirement obliges the engine to emit: it lists every loss site, the detector
efficiency included.
}
\begin{requirement}[The efficiency record lists every declared loss parameter]\label{req:efficiency-record}
The emitted \iface{efficiency_record} lists every declared loss parameter
$q_i$ that Theorem~\ref{thm:efficiency}\,(iv) names, by catalogue entry and
parameter, with the element that declares it and its position in the compiled
product of Definition~\ref{def:compile}. Every row names its parameter beside
its value, which is read as that parameter: a transmittance, a reflectivity,
or, for $y$, $y_1$ and $y_2$, a loss whose transmittance is $1-q_i$. Two
elements may carry equal numerical values; they are still two losses and
appear as two rows. A polarizer's pass transmittance $t$ is one such row,
under its element; the \termref{term-factored-loss}{factored loss} $t$ that the same record carries by
Requirement~\ref{req:normalization-record} is not a row and is recorded under
its own name.
\end{requirement}
\claimmeta{\metalabel{Motivated by}\mref{Def.}{def:compile}, \mref{Thm.}{thm:efficiency}, \mref{Req.}{req:normalization-record}
\metasep implemented (effective: accepted)
\metasep \metalabel{Code}\coderef{https://github.com/fraunhofersg/certiqs-key/blob/v0.1.0/src/certiqs_key/receiver/efficiency.py}{certiqs_key.receiver.efficiency:efficiency_record}}

% claim receiver/detector-imperfections/req:efficiency-convention
\claimsummary{receiver/detector-imperfections/req:efficiency-convention}{Whether the declared Grams come from power data or from click data is a declared field, written down in the efficiency record. The numbers alone cannot say which, and the two readings differ by exactly the efficiencies.}
\claimbookmark{receiver/detector-imperfections/req:efficiency-convention}{Requirement}{req:efficiency-convention}{Declared efficiency convention of the port Grams}
\claimleadin{%
Power data return the \termref{term-port-gram-matrices}{port Grams} $A_k$ and click data return $\eta_kA_k$
(Remark~\ref{rem:powervsclick}); the efficiencies must enter through
\eqref{eq:etagram} exactly once.
}
\begin{requirement}[Declared efficiency convention of the port Grams]\label{req:efficiency-convention}
Which of the two conventions a declaration means, Grams from power data
rescaled by $\eta_k$ or Grams from click data with every $\eta_k$ declared as
$1$, is a declared field, and the \iface{efficiency_record}
(Requirement~\ref{req:efficiency-record}) is where the answer is written down.
\end{requirement}
\claimmeta{\metalabel{Motivated by}\mref{Thm.}{thm:efficiency}, \mref{Rem.}{rem:powervsclick}, \mref{Req.}{req:efficiency-record}
\metasep accepted}

% claim receiver/detector-imperfections/req:efficiency-normalization
\claimsummary{receiver/detector-imperfections/req:efficiency-normalization}{By default the engine compiles a probe-free receiver at the factor t = \ensuremath{\eta}\_max/(99/100), or t = 1 where \ensuremath{\eta}\_max \ensuremath{\ge} 99/100, and emits that station's blocks, weights and dark-composed family. A declared factor in [\ensuremath{\eta}\_max, 1] replaces it, and one outside is refused.}
\claimbookmark{receiver/detector-imperfections/req:efficiency-normalization}{Requirement}{req:efficiency-normalization}{Compilation of a probe-free receiver at its normalized station}
\claimleadin{%
A receiver's measurement is uniform loss followed by its \termref{term-normalized-station}{normalized station}
(Proposition~\ref{prop:efficiency-normalization}), and the weight a
consumer reads loses a factor of the squared efficiency without it
(Remark~\ref{rem:normalization-gain}); the factor itself is a choice, of which
Remark~\ref{rem:normalization-scope} gives the trade-off. The reason the
default stops short of the full normalization where the largest efficiency
is below $99/100$ is the engine's present handling of a singular block, not a
property of the identity, and is stated in
\href{../../../spec/modules/receiver/design.md}{\nolinkurl{spec/modules/receiver/design.md}}.
}
\begin{requirement}[Compilation of a probe-free receiver at its normalized station]\label{req:efficiency-normalization}
For a request whose receiver declares no source modes and some efficiency
$\eta_k>0$, the engine compiles, in place of the declared receiver, its
normalized station at the \termref{term-factored-loss}{factored loss} $t$, with the efficiency of port $k$
equal to $\eta_k/t$, and emits that station's blocks, multiclick weights and
\termref{term-dark-composed-family}{dark-composed family}; $t$ is the value of \iface{efficiency_normalization}
where the request declares one; otherwise, with $\eta_{\max}$ the largest
declared efficiency, $t=\eta_{\max}/(99/100)$, formed exactly, where
$\eta_{\max}<99/100$, and $t=1$ where $\eta_{\max}\ge 99/100$. On every route the station compiled is the one at
the exact rational efficiencies $\eta_k/t$: its exact statements, its
reference and its enclosures are formed from those rationals, its
floating-point evaluation evaluates them as the \termref{term-evaluation-model}{evaluation model}
evaluates every declared number, and no value rounded from $\eta_k/t$ is
declared in their place; the quotient is exactly one at a port with
$\eta_k=t$. A declared $t$ outside
$[\eta_{\max},1]$ is refused with
\iface{efficiency_normalization_out_of_range}. A request that declares source
modes, or whose efficiencies are all zero, is compiled at $t=1$, and a
declared $t<1$ in such a request is refused with the same reason.
\end{requirement}
\claimreading{Unless you declare otherwise, a receiver without a probe source is analysed with nearly all of its largest detector efficiency moved into the channel as loss: every efficiency is divided by t, the largest one divided by 99/100, so the best detector is analysed at 99/100; at a largest efficiency of 99/100 or more none is divided. This leaves your observed statistics unchanged and typically raises the certified key rate at low efficiencies. You may declare the factor, between the largest efficiency and one; declaring one turns the normalization off. Declaring the largest efficiency itself is allowed and sound, but with dark counts at a photon cutoff of two or more the engine can then report no key where your model gives one, which is why, below 99/100, the default stops short of it; a receiver whose largest efficiency is one or close to it can meet the same limitation.}
\claimmeta{\metalabel{Motivated by}\mref{Def.}{def:normalized-station}, \mref{Prop.}{prop:efficiency-normalization}, \mref{Rem.}{rem:normalization-gain}, \mref{Rem.}{rem:normalization-scope}
\metasep implemented (effective: accepted)
\metasep \metalabel{Code}\coderef{https://github.com/fraunhofersg/certiqs-key/blob/v0.1.0/src/certiqs_key/receiver/efficiency.py}{certiqs_key.receiver.efficiency:normalize_efficiencies}}

% claim receiver/detector-imperfections/req:normalization-record
\claimsummary{receiver/detector-imperfections/req:normalization-record}{The efficiency record carries the factored loss and whether it was declared, every declared efficiency beside the one compiled, and whether the normalization was applied, exactly where t \textless{} 1, and why not. Where t \textless{} 1 the artefact says that its measurement is the normalized station's.}
\claimbookmark{receiver/detector-imperfections/req:normalization-record}{Requirement}{req:normalization-record}{The efficiency record names the normalization}
\claimleadin{%
A consumer of the artefact reads a measurement that is the receiver's only
after the factored loss is put in front of it
(Proposition~\ref{prop:efficiency-normalization}), so the factor travels with
the blocks. Requirement~\ref{req:efficiency-normalization} fixes the factor,
and Requirement~\ref{req:efficiency-record} the rows of the declared
transmittances. For a probe-free receiver at $t=1$ the normalized station is
the declared receiver,
however the factor was reached, so whether the normalization was applied is
read off $t$ alone, and who chose $t$ is recorded beside it.
}
\begin{requirement}[The efficiency record names the normalization]\label{req:normalization-record}
The emitted \iface{efficiency_record} carries the factored loss $t$, whether
$t$ was declared in the request or set by
Requirement~\ref{req:efficiency-normalization}, every declared efficiency
$\eta_k$ beside the exact rational $\eta_k/t$ at which the receiver was
compiled, and whether the normalization was applied, which it was exactly
where $t<1$. Where it was not, it carries the reason, the first that holds of:
source modes declared, every efficiency zero, or the factor one, declared or
by default. Where $t<1$ the artefact states that its blocks, multiclick
weights and dark-composed family are those of the normalized station at $t$,
whose pull-back through the \termref{term-pure-loss-channel}{pure-loss channel} of transmittance $t$ is the
receiver's measurement.
\end{requirement}
\claimmeta{\metalabel{Motivated by}\mref{Def.}{def:normalized-station}, \mref{Prop.}{prop:efficiency-normalization}, \mref{Req.}{req:efficiency-record}, \mref{Req.}{req:efficiency-normalization}
\metasep implemented (effective: accepted)
\metasep \metalabel{Code}\coderef{https://github.com/fraunhofersg/certiqs-key/blob/v0.1.0/src/certiqs_key/receiver/efficiency.py}{certiqs_key.receiver.efficiency:EfficiencyNormalization}}

% claim receiver/detector-imperfections/req:loss-before-evaluation
\claimsummary{receiver/detector-imperfections/req:loss-before-evaluation}{Where the receiver was compiled at a factored loss t below one, every evaluation of its emitted measurement on a state of the channel system that stands for the receiver's statistics, the honest statistics among them, is made on the state after the pure-loss channel; one that cannot is refused.}
\claimbookmark{receiver/detector-imperfections/req:loss-before-evaluation}{Requirement}{req:loss-before-evaluation}{The factored loss is applied before the artefact is evaluated}
\claimleadin{%
The emitted measurement of a normalized receiver is not the receiver's: the
receiver's click statistics on an input $\rho$ are those of the emitted
measurement on $\Lambda^{\mathrm{loss}}_t(\rho)$
(Proposition~\ref{prop:efficiency-normalization}\,(i), second form, and (ii)).
A consumer that evaluated the emitted measurement on $\rho$ itself would form
the statistics of a more efficient receiver. The factor travels in the record
of Requirement~\ref{req:normalization-record}. The states the key-rate program
optimises over are another matter: the loss is handed to the adversary
(Proposition~\ref{prop:efficiency-normalization}\,(iii)), so it is already
part of the channel those states range over, and applying it to them again
would be sound but would give up the gain the normalization exists for.
}
\begin{requirement}[The factored loss is applied before the artefact is evaluated]\label{req:loss-before-evaluation}
Where a receiver artefact's \iface{efficiency_record} carries a factored loss
$t<1$, every evaluation of its emitted blocks or dark-composed
family, or of a target formed from them, on a state that stands
for the receiver's statistics on an input $\rho$ of the channel system (an
honest-statistics provider forming the honest law among them) is made on
$\Lambda^{\mathrm{loss}}_t(\rho)$, or on the state the same construction forms
from $\Lambda^{\mathrm{loss}}_t(\rho)$, in place of $\rho$, with the $t$ of that
receiver's record. Where the input is a joint state of which that receiver's
channel system is one factor, such as a state shared with another party or
with a second normalized receiver, $\Lambda^{\mathrm{loss}}_t$ stands here for
$\idmap\otimes\Lambda^{\mathrm{loss}}_t$ acting on that receiver's
channel-system factor, each receiver with the $t$ of its own record. The states the
key-rate program optimises over are not inputs in this sense, and the loss
is not applied to them. A consumer that cannot apply $\Lambda^{\mathrm{loss}}_t$ to
the input it is handed refuses with \iface{factored_loss_not_applied}.
\end{requirement}
\claimreading{When the receiver is analysed at its normalized station, the expected (honest) statistics the analysis is built around are those of your actual receiver: the efficiency moved into the channel is applied to the expected input before the normalized measurement is evaluated. Without this the acceptance test would be set for a receiver more efficient than yours, and the reported key would be one your device does not reach.}
\claimmeta{\metalabel{Motivated by}\mref{Def.}{def:pure-loss}, \mref{Def.}{def:normalized-station}, \mref{Prop.}{prop:efficiency-normalization}, \mref{Req.}{req:normalization-record}
\metasep implemented (effective: accepted)
\metasep \metalabel{Code}\coderef{https://github.com/fraunhofersg/certiqs-key/blob/v0.1.0/src/certiqs_key/runs/pipeline.py}{certiqs_key.runs.pipeline:LossAppliedProvider}}

% claim receiver/detector-imperfections/req:dark-route-exclusive
\claimsummary{receiver/detector-imperfections/req:dark-route-exclusive}{A request that declares a dark-count matrix R while some per-port dark-count probability is non-zero is refused. A declared R therefore comes with every per-port probability zero, and exactly one route is taken.}
\claimbookmark{receiver/detector-imperfections/req:dark-route-exclusive}{Requirement}{req:dark-route-exclusive}{Refusal of a dark-count channel supplied on both routes}
\claimleadin{%
A request carrying both would supply the channel twice with no rule that makes
the two agree, and the record of which route the emitted channel came from
(Remark~\ref{rem:corrdark}, Requirement~\ref{req:dark-channel-record})
presupposes that exactly one was taken.
}
\begin{requirement}[Refusal of a dark-count channel supplied on both routes]\label{req:dark-route-exclusive}
The two routes of Definition~\ref{def:darkadm} are exclusive: a request
declaring $R$ while some $p_{\dc|k}\neq0$ is refused with
\iface{dark_count_channel_declared_twice}.
\end{requirement}
\claimmeta{\metalabel{Motivated by}\mref{Def.}{def:darkadm}, \mref{Rem.}{rem:corrdark}, \mref{Req.}{req:dark-channel-record}
\metasep implemented (effective: accepted)
\metasep \metalabel{Code}\coderef{https://github.com/fraunhofersg/certiqs-key/blob/v0.1.0/src/certiqs_key/receiver/dark.py}{certiqs_key.receiver.dark:check_dark_count_source}}

% claim receiver/detector-imperfections/req:dark-channel-record
\claimsummary{receiver/detector-imperfections/req:dark-channel-record}{The artifact carries the dark-count matrix as declared or generated, its route, and r\_min, rbar and \ensuremath{\delta}\_dc under distinct names, none modified. The inequality a downstream module composes with is that of \hyperlink{claim-receiver/detector-imperfections/prop:dark-slack-bound}{Proposition~\ref*{prop:dark-slack-bound}}.}
\claimbookmark{receiver/detector-imperfections/req:dark-channel-record}{Requirement}{req:dark-channel-record}{The emitted dark-count channel and the scalars of its composed bound}
\claimleadin{%
Both $r_{\min}$ and $\bar r$ are emitted: $\bar r$ is the number the
composition uses, while $r_{\min}$ as computed is what \eqref{eq:rmin} defines
and is the only one from which a reader can see that a declared channel
produced $r_{\min}>1$ at all, itself a signal that the declaration is off by
more than rounding (Remark~\ref{rem:d2slack}\,(c)).
}
\begin{requirement}[The emitted dark-count channel and the scalars of its composed bound]\label{req:dark-channel-record}
The artifact's \iface{dark_count_channel} entry carries $R$, which of the two
routes of Definition~\ref{def:darkadm} it came from
(Remark~\ref{rem:corrdark}), and $r_{\min}$, $\bar r$ and $\delta_{\dc}$ under
distinct names. Nothing is clamped: $R$ is emitted as declared and $r_{\min}$
as \eqref{eq:rmin} computes it. \eqref{eq:darkpmcslack}, \emph{not}
\eqref{eq:darkpmc}, is the inequality a downstream module composes with.
\end{requirement}
\claimmeta{\metalabel{Motivated by}\mref{Def.}{def:darkadm}, \mref{Thm.}{thm:dark-counts}, \mref{Prop.}{prop:dark-slack-bound}, \mref{Rem.}{rem:d2slack}, \mref{Rem.}{rem:corrdark}
\metasep implemented (effective: accepted)
\metasep \metalabel{Code}\coderef{https://github.com/fraunhofersg/certiqs-key/blob/v0.1.0/src/certiqs_key/receiver/artifact.py}{certiqs_key.receiver.artifact:DarkCountRecord}}

% claim receiver/detector-imperfections/req:dark-composed-reference
\claimsummary{receiver/detector-imperfections/req:dark-composed-reference}{The emitted low-photon blocks are the dark-composed family's, the dark-count channel still emitted. For a model admissible exactly they carry sound enclosures, the handed reference, and a factor enclosure where a port dilation admits one; only exact ones where enclosed references are not used.}
\claimbookmark{receiver/detector-imperfections/req:dark-composed-reference}{Requirement}{req:dark-composed-reference}{The emitted blocks are dark-composed and carry their reference}
\claimleadin{%
The run evaluates its first-order bound on the reference handed beside the
snap, and the noise floor of the device (Proposition~\ref{prop:noise-floor})
belongs to the model the bound is evaluated on, not to a step after it. The
reference offered is therefore the \termref{term-dark-composed-family}{dark-composed family}, and the
blocks the receiver emits are its blocks. For a \termref{term-setting-family}{setting family} a
submodule constructs, the enclosures at each setting are those that submodule states.
The
enclosures cannot recover every zero of the model: an entry that vanishes
through cancellation, in an inclusion--exclusion sum or in the composition
with the \termref{term-dark-count-channel}{dark-count channel}, is enclosed with a positive radius wherever its
enclosure is not exact, and a snap then does not treat it as known to be
zero. That happens off the \termref{term-exact-route}{exact route}, and on it as well, at the
entries of a block $n\ge2$ between occupation vectors $\vec m,\vec m'$ with
$\vec m!\,\vec m'!$ not a square, which
Theorem~\ref{thm:exact-blocks} does not make exact. Item (ii) lists the zeros
that are kept exact. An enclosure with positive radii does not fix the rank
of a block whose rank is below its dimension and whose support no coordinate
vectors span, such as a block of an ideal analyser at an angle whose
projector has irrational entries; item (vii) offers such a block besides in
factor form, read through a \termref{term-port-dilation}{port dilation}
(Theorem~\ref{thm:factor-blocks}), wherever its columns are few enough to fix
its rank. The instance hands the reference beside its snap, and
the run evaluates its first-order bound on it, so that item (vi) asks for it
in every run that uses enclosed references. The emitted blocks remain a decomposition of the physical
measurement and are not a squashing argument (Remark~\ref{rem:blockmeaning}).
}
\begin{requirement}[The emitted blocks are dark-composed and carry their reference]\label{req:dark-composed-reference}
\begin{enumerate}[label=(\roman*),leftmargin=2.4em]
\item The \iface{photon_blocks} entry carries, for every probe setting $t$
  and every $n\le N$, the blocks
  $\blk{\widetilde F}{\eqn}_{v|t}=\sum_{S'\subseteq[K]}R_{\cg^{-1}(v)|S'}\,\blk{M}{\eqn}_{S'|t}$
  given by the formula \eqref{eq:dark-composed} for the emitted matrix $R$, in
  floating point. They are the blocks of the \termref{term-dark-composed-family}{dark-composed family}
  (Definition~\ref{def:dark-composed}) when $R$ is a \termref{term-dark-count-channel}{dark-count channel},
  satisfying (D1) and (D2) exactly; for a declared $R$ admitted only at the
  \termref{term-receiver-tolerances}{receiver tolerances} they are the same formula evaluated at that $R$.
\item For a declaration with a model admissible exactly
  (Requirement~\ref{req:exact-route-decision}) whose $R$ satisfies (D1) and
  (D2) exactly, decided in exact arithmetic on its declared entries on the
  \termref{term-declared-dark-count-route}{declared dark-count route}, the entry also carries, for every setting
  and every $n\le N$, a \termref{term-complex-enclosure}{complex enclosure} with a Hermitian centre of
  every block $\blk{\widetilde F}{\eqn}_{v|t}$
  (Definition~\ref{def:complex-enclosure}), computed in exact rational
  arithmetic on the \exttermref{../foundations/contract.pdf}{term-exact-island}{exact island}
  (\xref{foundations/def:exact-island}\,(i)) and with no binary64 value as an
  input: from the enclosures of the port Grams of
  Requirement~\ref{req:exact-route-decision}\,(iii), by
  Theorem~\ref{thm:enclosed-blocks}\,(ii) for an unconditioned receiver and
  by the enclosure of its blocks that a submodule constructing the setting
  family states for that family, then by Corollary~\ref{cor:exact-dark-composed}\,(i) and
  Lemma~\ref{lem:enclosure-arithmetic}\,(v). The
  \termref{term-square-root-enclosure}{square-root enclosures} and
  \termref{term-angle-enclosure}{angle enclosures} used may be any that
  Definition~\ref{def:complex-enclosure} admits. The following entries are
  enclosed exactly, as zero with radius zero:
  \begin{enumerate}[label=(\alph*),leftmargin=2.0em]
  \item every entry of a block $\blk{\widetilde F}{\eqn}_{v|t}$ such that
    every $\blk{\widetilde M}{\eqn}_{S|t}$ with $S\in\cg^{-1}(v)$, whose sum
    it is (Definition~\ref{def:dark-composed}), is shown to vanish: for the
    unconditioned receiver on the \termref{term-generated-dark-count-route}{generated dark-count route}, and
    for no other receiver or route, by
    Proposition~\ref{prop:noise-floor}\,(v), decided on its declared
    per-port dark-count probabilities; or, for $S=\varnothing$ and $n\ge1$, by
    Proposition~\ref{prop:noise-floor}\,(iv) at a setting whose operator
    $B^{(t)}$ the engine has computed in exact arithmetic and found to be the
    identity, which for the unconditioned receiver is $A_\Sigma$ computed
    from the port Grams of
    Requirement~\ref{req:exact-route-decision}\,(i) or (ii);
  \item every entry whose enclosure is exact by
    Theorem~\ref{thm:exact-blocks}\,(ii) or by the exactness of its blocks
    that a submodule constructing the setting family states, carried through
    Corollary~\ref{cor:exact-dark-composed}\,(i), and whose value is zero.
  \end{enumerate}
  No other entry is required to be enclosed as zero.
\item The enclosures of (ii), with the family they enclose, are the
  \exttermref{../foundations/contract.pdf}{term-reference}{reference} (\xref{foundations/def:exact-reference}) that the receiver
  offers for a snap of its emitted family
  (\xref{foundations/req:snap-reference}), one per setting, with constant $1$
  (Corollary~\ref{cor:exact-dark-composed}\,(ii)). The entry records the
  \emph{offered} reference kind of every block: \emph{exact} where every
  radius of the block is zero, and \emph{enclosed} otherwise, with the
  largest radius of the block and the largest radius of the enclosures of
  irrational inputs (square roots, trigonometric values) from which the
  centres are computed exactly. The record of every snap of the family
  carries besides the kind of the reference the snap ran on
  (\xref{foundations/req:who-snaps}): \emph{exact} on a block whose
  offered kind is exact, the snap running on its centres, and \emph{none}
  on every other block, the snap running on the floating-point blocks of
  (i). The reference is an
  \exttermref{../foundations/contract.pdf}{term-exact-reference}{exact reference} where every block is exact. It also names the claims that establish the
  enclosure, positive semidefiniteness and completeness of the reference,
  which \xref{foundations/req:snap-reference}\,(i) leaves to the
  constructing module: Corollary~\ref{cor:exact-dark-composed}\,(ii), with
  Theorem~\ref{thm:enclosed-blocks}\,(ii) through
  Corollary~\ref{cor:exact-dark-composed}\,(iii) for an unconditioned
  receiver and the enclosure claim of the constructing submodule for a
  setting family a submodule constructs. On the \termref{term-exact-route}{exact route}, with the
  port Grams enclosed exactly, every block $n\le1$ has radius zero
  (Theorem~\ref{thm:exact-blocks}\,(iii) for an unconditioned receiver, the
  exactness claim of the constructing submodule for a setting family a
  submodule constructs, and Corollary~\ref{cor:exact-dark-composed}\,(i)). The reference is read against
  the \exttermref{../foundations/contract.pdf}{term-sector-map}{sector map} of the floating-point blocks of (i) in the same entry
  (\xref{foundations/def:sector-map}): the same block kind and block
  identifiers, the photon numbers $n=0,\dots,N$ in increasing order, and
  within each block the occupation vectors in the one order in which the
  entry lists the floating-point blocks, which is also the order of the rows
  and columns of every centre and radius matrix. The entries of the reference
  that a snap treats as known to be exactly zero
  (\xref{foundations/def:snap}) are therefore exactly those enclosed as zero
  with radius zero.
\item Where the declaration has no model admissible exactly, where a declared
  $R$ fails (D1) or (D2) exactly, or where a reciprocal that the enclosure
  of a setting family a submodule constructs forms receives no enclosure, the
  entry carries no reference, records the offered reference kind of every
  block as \emph{none}, and records the cause.
\item $R$ remains emitted in \iface{dark_count_channel}
  (Requirement~\ref{req:dark-channel-record}), where it serves the
  high-photon weight. It is already composed into the blocks of (i)--(iii),
  and a consumer does not compose it into them again.
\item Items (ii) and (iii) bind every run, save one that declares that
  enclosed references are not used
  (\xref{foundations/req:snap-reference}\,(i)). Such a run need form no
  enclosure of (ii) with a radius that is not zero, beyond what deciding that
  a block is not enclosed exactly takes, and carries none: it carries the
  reference of (iii) only on the blocks every entry of which it encloses
  exactly, with radius zero, from port Grams enclosed exactly, each block so
  carried being an \exttermref{../foundations/contract.pdf}{term-exact-reference}{exact reference} of that block, and records their
  offered kind \emph{exact}; every other block it records as offered with
  the kind \emph{none}, with the cause that enclosed references are not
  used (\xref{foundations/req:snap-reference}\,(iii)). The reference the entry
  carries, its centres and radii as data, is the \exttermref{../foundations/contract.pdf}{term-handed-reference}{handed reference}
  that the instance hands beside its snap of the emitted family
  (\xref{foundations/def:evaluated-off-snap}\,(i)), whatever the snap itself
  uses of it: the snap uses only the centres of the exact blocks
  (\xref{foundations/req:snap-reference}\,(ii)), and the run's first-order
  bound reads the reference on every block it carries
  (\xref{meat-core/frank-wolfe/req:bracket-on-reference}). Item (iv) binds
  as stated.
\item \emph{Factor enclosures.} Where the entry carries the reference of
  (iii) for a receiver with $\InpS=\varnothing$, whose one setting is $\off$
  (Corollary~\ref{cor:exact-dark-composed}\,(iii)), the engine forms
  besides, on the blocks the count below admits and on no other, a
  \exttermref{../foundations/contract.pdf}{term-factor-enclosure}{factor enclosure} (\xref{foundations/def:factor-enclosure}) of the
  block $\blk{\widetilde F}{\eqn}_{v|\off}$: the one of
  Theorem~\ref{thm:factor-blocks}\,(iii) at the weights
  $w_{S'}\coloneqq R_{\cg^{-1}(v)|S'}$, rationals by
  Corollary~\ref{cor:exact-dark-composed}\,(iv), at which the block of that
  theorem is $\blk{\widetilde F}{\eqn}_{v|\off}$
  (Definition~\ref{def:dark-composed}), its rows in the order of (iii), for
  the port dilation (Definition~\ref{def:port-dilation}) formed as follows.
  \begin{enumerate}[label=(\alph*),leftmargin=2.0em]
  \item On the route that declares a component graph, the dilation of
    Lemma~\ref{lem:catalogue-dilations}\,(ii), with the loss factor of
    part (i) of that lemma for every element. A row of a loss factor, and of
    the efficiency layer's, is formed only where the declared value that
    scales it is decided in exact arithmetic to make it non-zero, so that
    whether an element is lossless is read off its declared parameters and
    never off a computed number. The loss factor of a \termref{term-generic-block}{generic block} is formed
    from a decomposition
    $\id_\Sigma-\Gmat_\Sigma^{\dagger}\Gmat_\Sigma=E^{\dagger}DE$ formed and
    decided in exact rational arithmetic; where none is formed, no dilation of
    the declaration is.
  \item On a route that declares the port Grams or the \termref{term-instrument-matrix}{instrument matrix}, the
    dilation of Theorem~\ref{thm:factor-blocks}\,(iv), where the engine forms
    decompositions $A_k=E_k^{\dagger}D_kE_k$ and $\defic=E_0^{\dagger}D_0E_0$
    of the kind named there, formed and decided in exact rational arithmetic.
  \end{enumerate}
  The entries of the dilation are enclosed by the rules of
  Lemma~\ref{lem:enclosure-arithmetic} from \termref{term-square-root-enclosure}{square-root enclosures} and \termref{term-angle-enclosure}{angle
  enclosures}, as Requirement~\ref{req:exact-route-decision}\,(iii) encloses
  the elements' matrices and the $\sqrt{\eta_k}$, and the columns as
  Theorem~\ref{thm:factor-blocks}\,(ii) encloses them, in exact rational
  arithmetic or by a \exttermref{../foundations/contract.pdf}{term-rigorously-enclosed}{rigorously enclosed} computation
  (\xref{foundations/def:tiers}), with no binary64 value as an input; a column
  enclosed exactly as zero is omitted, and the core is exact. The engine forms
  the factor enclosure of a block only where the number of occupation vectors
  $\vec m$ of weight $n$ on the modes of the dilation whose row is not
  enclosed exactly as zero, with $R_{\cg^{-1}(v)|S_{\vec m}}>0$, is at most
  the block's dimension $\binom{|\Inp|+n-1}{n}$, a count it decides before it
  computes any column: a factor enclosure with more columns than that
  dimension fixes no rank
  (\xref{foundations/cor:rank-over-factor-enclosure}\,(ii)). It forms none
  for a receiver with $\InpS\neq\varnothing$, whose blocks are not formed
  through a port dilation; none where no dilation of (a) or (b) is formed; none on a block
  the count does not admit; and none in a run that declares that enclosed
  references are not used (vi). A factor enclosure is carried beside the
  entrywise enclosure of its block, which it leaves unaltered, and is handed
  with the reference, altered by no step. The entry records per block whether
  one is carried, with its number of columns and its largest radius, or the
  cause where none is; and it names Theorem~\ref{thm:factor-blocks} and, on
  the route of (a), Lemma~\ref{lem:catalogue-dilations} as the claims that
  establish that the block is a member of its factor enclosure, which
  \xref{foundations/req:snap-reference}\,(i) leaves to the module that
  constructs the reference.
\end{enumerate}
\end{requirement}
\claimmeta{\metalabel{Motivated by}\mref{Rem.}{rem:blockmeaning}, \mref{Def.}{def:complex-enclosure}, \mref{Lem.}{lem:enclosure-arithmetic}, \mref{Thm.}{thm:enclosed-blocks}, \mref{Def.}{def:port-dilation}, \mref{Thm.}{thm:factor-blocks}, \mref{Thm.}{thm:exact-blocks}, \mref{Def.}{def:tolmap}, \mref{Req.}{req:exact-route-decision}, \mref{Lem.}{lem:catalogue-dilations}, \mref{Def.}{def:darkadm}, \mref{Def.}{def:darktol}, \mref{Def.}{def:dark-composed}, \mref{Cor.}{cor:exact-dark-composed}, \mref{Prop.}{prop:noise-floor}, \mref{Req.}{req:dark-channel-record}, \extref{../foundations/contract.pdf}{foundations/cor:rank-over-factor-enclosure}, \extref{../foundations/contract.pdf}{foundations/def:evaluated-off-snap}, \extref{../foundations/contract.pdf}{foundations/def:exact-island}, \extref{../foundations/contract.pdf}{foundations/def:exact-reference}, \extref{../foundations/contract.pdf}{foundations/def:factor-enclosure}, \extref{../foundations/contract.pdf}{foundations/def:sector-map}, \extref{../foundations/contract.pdf}{foundations/def:snap}, \extref{../foundations/contract.pdf}{foundations/def:tiers}, \extref{../foundations/contract.pdf}{foundations/req:snap-reference}, \extref{../foundations/contract.pdf}{foundations/req:who-snaps}, \extref{../meat-core/contract.pdf}{meat-core/frank-wolfe/req:bracket-on-reference}
\metasep \metalabel{Used by}\mref{Req.}{req:snap-fix}
\metasep implemented (effective: accepted)
\metasep \metalabel{Code}\coderef{https://github.com/fraunhofersg/certiqs-key/blob/v0.1.0/src/certiqs_key/receiver/reference.py}{certiqs_key.receiver.reference:dark_composed_reference}}

% claim receiver/detector-imperfections/req:snap-fix
\claimsummary{receiver/detector-imperfections/req:snap-fix}{A moved block of the receiver's snap that rests on the snap premise records the reference kind, its cause, its group and the fix the cause details: exact data, an enclosed reference used, tighter enclosures, or loss and dark counts. Only a strict run refuses, or one whose last resort fails.}
\claimbookmark{receiver/detector-imperfections/req:snap-fix}{Requirement}{req:snap-fix}{Cause and fix recorded for a receiver block resting on its snap}
\claimleadin{%
Under \xref{foundations/req:snap-reference}, a moved block of the receiver's
snap is not refused by default: unless the run evaluates it off its snap
(\xref{foundations/def:evaluated-off-snap}), the run rests there on the
premise that the device measures the snapped family
(\xref{foundations/ass:snapped-measurement}), and the record names the block
with its cause, its group and a fix. Where the support-safe rules cannot
restore the conditions of a snap on a block, a last resort snaps it anyway
and the record flags it (\xref{foundations/req:support-safe-snap}\,(iv)). A
user can act on the record only if it says what to change in the receiver
declaration, and for a receiver the cause of the group \emph{no usable
reference} and the last resort each have a definite remedy, which details
the fix of the group; Remark~\ref{rem:noise-floor-receiver} explains them.
The record also names, block by block, whether the snap ran on the centres
of an exact reference or on the floating-point blocks, and which reference
the receiver hands, so that a reader sees where an enclosure entered the
result. An angle at any value is no cause: every declaration with a model
admissible exactly has an enclosed reference
(Requirement~\ref{req:dark-composed-reference}).
}
\begin{requirement}[Cause and fix recorded for a receiver block resting on its snap]\label{req:snap-fix}
When a moved block of a snap of the receiver's emitted dark-composed blocks
(Requirement~\ref{req:dark-composed-reference}) rests on
\xref{foundations/ass:snapped-measurement} under
\xref{foundations/req:snap-reference}\,(iii), its record names, beside the
setting, the block with its cause and the group that requirement names for
it, the reference kind of the snap, as \xref{foundations/req:who-snaps}
defines it (\emph{exact}, where the snap ran on the centres of an exact
reference of the block, or \emph{none}), the kind of the reference handed on
the block, if any (\emph{exact}, every radius of the block zero, or
\emph{enclosed}, with the largest radius of the block and the largest radius
of the enclosures of irrational inputs (square roots, trigonometric values)
from which the centres were computed exactly), and the fix below for the
cause, which details the fix of its group. A run that declares the strict
option of \xref{foundations/req:snap-reference}\,(iv) refuses instead,
naming the same; no other run refuses for these causes, save where the last
resort of \xref{foundations/req:support-safe-snap}\,(iv) yields no snap.
\begin{enumerate}[label=(\roman*),leftmargin=2.4em]
\item \emph{No reference.} The cause recorded by
  Requirement~\ref{req:dark-composed-reference}\,(iv), in the group
  \emph{no usable reference}, and the fix for it:
  \begin{enumerate}[label=(\alph*),leftmargin=2.0em]
  \item for declared port Grams that are not Hermitian, a declared family that
    is not an \termref{term-admissible-gram-family}{admissible Gram family} exactly, or a \termref{term-generic-block}{generic block} that is not a
    contraction exactly (Requirement~\ref{req:exact-route-decision}\,(i),
    (ii)), declared data that are Hermitian and admissible exactly, in
    rational arithmetic and not only at the receiver tolerances;
  \item for a declared dark-count matrix that does not satisfy (D1) and (D2)
    exactly, one that is exactly column-stochastic and monotone, or the
    \termref{term-generated-dark-count-route}{generated dark-count route};
  \item for a reciprocal that the enclosure of a setting family a submodule
    constructs forms and that receives no enclosure, square-root and angle
    enclosures with smaller
    radii (Definition~\ref{def:complex-enclosure}).
  \end{enumerate}
\item \emph{Enclosed references not used.} Where the run declares that
  enclosed references are not used and the entry offers an enclosed one, the
  block is handed no reference, in the group \emph{no usable reference}, and
  the fix names the kind the entry offers for every block
  (Requirement~\ref{req:dark-composed-reference}\,(iii)), with the largest
  radius of every block offered as enclosed, and the alternative of an exact
  reference, which the receiver offers for a declaration on the exact route
  at a cutoff $N\le1$ when its port Grams are enclosed exactly
  (Theorem~\ref{thm:exact-blocks}\,(iii),
  Corollary~\ref{cor:exact-dark-composed}\,(ii)), or a run that uses
  enclosed references.
\end{enumerate}
A block that the snap makes by the last resort of
\xref{foundations/req:support-safe-snap}\,(iv) is flagged
\iface{support_not_preserved} in the record, beside its cause and reference
kind, and where it rests, in the group \emph{last resort}, the fix names the
ports whose declared efficiency is $1$ where $\lmax(A_\Sigma)=1$, and the
click patterns $S$ whose entry $R_{S|\varnothing}$ of the
\termref{term-dark-count-channel}{dark-count channel} is $0$, since Proposition~\ref{prop:noise-floor}\,(i)
needs $R_{S|\varnothing}>0$ for every $S$: on the
\termref{term-declared-dark-count-route}{declared dark-count route} it asks for a declared $R$ with
$R_{S|\varnothing}>0$ for every $S$, and on the
\termref{term-generated-dark-count-route}{generated dark-count route}, where every $R_{S|\varnothing}$ is positive
exactly when $0<p_{\dc|k}<1$ for every $k$
(Proposition~\ref{prop:noise-floor}\,(ii)), a dark-count probability in
$(0,1)$ in every port; in either case a model with loss, at the smallest loss
and dark-count rates the user is sure the device has
(Proposition~\ref{prop:noise-floor}); and, where the entry offers no
reference, the fix of (i) as well. Where every element of a block of the
reference is positive definite (Proposition~\ref{prop:noise-floor}\,(i)), no
row of it is structurally zero (\xref{foundations/def:snap}), so every index
lies in the \exttermref{../foundations/contract.pdf}{term-structural-support}{structural support} of every element and the uniform lift
of that item meets the rules (i) to (iii) of
\xref{foundations/req:support-safe-snap} on the block. For a block whose
cause is the run's, a handed reference on which the companion's bound is
withdrawn or not evaluated, or counts while the key length is formed from a
number certified on the snapped declaration in its place, the cause
\iface{companion_not_certified}
(\xref{meat-core/frank-wolfe/req:certified-evaluation}\,(i)), the record of
the run names the fix
\xref{foundations/req:snap-reference}\,(iii) names for that cause, and, for
such a block snapped by the last resort, which rests in the group
\emph{last resort}, that fix beside the one above; the run decides that cause
after its search, so the fix is completed there and not at the snap. A
handed block that the run evaluates off its snap rests on nothing and is
named with no fix, whether or not the last resort snapped it. A reference that fails
the decisions of \xref{foundations/req:snap-reference}\,(i) cannot arise for
the receiver: the centres are Hermitian with entries in $\mathbb{Q}(i)$ and
the radii symmetric non-negative rationals by construction
(Lemma~\ref{lem:enclosure-arithmetic}\,(v)), and the enclosure, positive
semidefiniteness and completeness that the constructing module establishes
are Corollary~\ref{cor:exact-dark-composed}\,(ii), whose hypothesis
Theorem~\ref{thm:enclosed-blocks}\,(ii), through
Corollary~\ref{cor:exact-dark-composed}\,(iii), and, for a setting family a
submodule constructs, that submodule's enclosure of its blocks supply.
\end{requirement}
\claimmeta{\metalabel{Motivated by}\mref{Def.}{def:complex-enclosure}, \mref{Lem.}{lem:enclosure-arithmetic}, \mref{Thm.}{thm:enclosed-blocks}, \mref{Def.}{def:exact-route}, \mref{Thm.}{thm:exact-blocks}, \mref{Req.}{req:exact-route-decision}, \mref{Def.}{def:darkadm}, \mref{Cor.}{cor:exact-dark-composed}, \mref{Prop.}{prop:noise-floor}, \mref{Rem.}{rem:noise-floor-receiver}, \mref{Req.}{req:dark-composed-reference}, \extref{../foundations/contract.pdf}{foundations/ass:snapped-measurement}, \extref{../foundations/contract.pdf}{foundations/def:evaluated-off-snap}, \extref{../foundations/contract.pdf}{foundations/def:exact-reference}, \extref{../foundations/contract.pdf}{foundations/def:snap}, \extref{../foundations/contract.pdf}{foundations/req:snap-reference}, \extref{../foundations/contract.pdf}{foundations/req:support-safe-snap}, \extref{../foundations/contract.pdf}{foundations/req:who-snaps}, \extref{../meat-core/contract.pdf}{meat-core/frank-wolfe/req:certified-evaluation}
\metasep implemented (effective: accepted)
\metasep \metalabel{Code}\coderef{https://github.com/fraunhofersg/certiqs-key/blob/v0.1.0/src/certiqs_key/receiver/snapfix.py}{certiqs_key.receiver.snapfix:receiver_fixes}}
\clearpage
\bookmarksetupnext{level=part}
\section*{Interface}\addcontentsline{toc}{section}{Interface}
\sectionintro{The declared inputs and outputs of the module and its submodules: each key with its symbol, when it is present and the claims that fix it. The domains, types and refusal rules are in the interface files themselves.}
\subsection*{Inputs: Receiver POVM (receiver)}
\noindent{\small From \href{interface.yaml}{\nolinkurl{interface.yaml}}.}\par\smallskip
{\footnotesize\begin{xltabular}{\linewidth}{@{}>{\raggedright\arraybackslash}p{4.4cm}>{\raggedright\arraybackslash}X>{\raggedright\arraybackslash}p{2.9cm}@{}}
\toprule
Key and symbol & Presence & Fixed by\\
\midrule
\endhead
\iface{mode_space}\newline M (the declared mode space) &  & ---\\
\addlinespace
\iface{input_modes}\newline Inp = Inp\_C (disjoint union) Inp\_S &  & ---\\
\addlinespace
\iface{port_map}\newline \{P\_k\}\_\{k=1..K\} &  & ---\\
\addlinespace
\iface{port_grams}\newline \{A\_k\} &  & ---\\
\addlinespace
\iface{instrument_matrix}\newline Gamma &  & ---\\
\addlinespace
\iface{element_graph}\newline G &  & ---\\
\addlinespace
\iface{detector_layer}\newline \{(P\_k, eta\_k, p\_dc\textbar{}k)\} &  & ---\\
\addlinespace
\iface{dark_count_channel}\newline R &  & ---\\
\addlinespace
\iface{probe_declaration}\newline (T, \{rho\textasciicircum{}(t)\}) &  & ---\\
\addlinespace
\iface{photon_cutoff}\newline N &  & ---\\
\addlinespace
\iface{tail_window}\newline \{m\_t\}\_\{t in T\} &  & ---\\
\addlinespace
\iface{exact_window}\newline \{m*\_t\} &  & ---\\
\addlinespace
\iface{numerical_tolerances}\newline \{eps\_1, ..., eps\_6\}, eps\_gap, eps\_eig &  & ---\\
\addlinespace
\iface{outcome_coarse_graining}\newline cg &  & ---\\
\addlinespace
\iface{requested_assurance} &  & ---\\
\addlinespace
\iface{efficiency_normalization}\newline t &  & ---\\
\bottomrule
\end{xltabular}}
\subsection*{Outputs: Receiver POVM (receiver)}
\noindent{\small From \href{interface.yaml}{\nolinkurl{interface.yaml}}.}\par\smallskip
{\footnotesize\begin{xltabular}{\linewidth}{@{}>{\raggedright\arraybackslash}p{4.4cm}>{\raggedright\arraybackslash}X>{\raggedright\arraybackslash}p{2.9cm}@{}}
\toprule
Key and symbol & Presence & Fixed by\\
\midrule
\endhead
\iface{input_mode_space}\newline Inp = Inp\_C (disjoint union) Inp\_S & always & \mref{Def.}{def:modes}, \mref{Def.}{def:instrument}\\
\addlinespace
\iface{outcome_alphabet}\newline V with cg & always & \mref{Def.}{def:coarse}, \mref{Thm.}{thm:coarse-graining}\\
\addlinespace
\iface{port_grams}\newline \{A\_k\} & always & \mref{Def.}{def:gram}, \mref{Thm.}{thm:efficiency}\\
\addlinespace
\iface{receiver_construction}\newline G, \{P\_k\}, probe declaration & route (iii) only & ---\\
\addlinespace
\iface{photon_blocks}\newline blk M\textasciicircum{}\{(=n)\}\_v, 0 \textless{}= n \textless{}= N & always & \mref{Thm.}{thm:click-povm}, \mref{Thm.}{thm:coarse-graining}, \mref{Thm.}{thm:dark-counts}, \mref{Def.}{def:dark-composed}, \mref{Thm.}{thm:enclosed-blocks}, \mref{Thm.}{thm:exact-blocks}, \mref{Thm.}{thm:factor-blocks}, \mref{Cor.}{cor:exact-dark-composed}, \mref{Req.}{req:dark-composed-reference}\\
\addlinespace
\iface{block_dimensions}\newline binom(\textbar{}Inp\_C\textbar{} + n - 1, n) & always & \mref{Def.}{def:fock}\\
\addlinespace
\iface{efficiency_record}\newline \{(q\_i, e\_i, s\_i)\}, t & always & \mref{Req.}{req:efficiency-record}, \mref{Req.}{req:normalization-record}\\
\addlinespace
\iface{dark_count_channel}\newline R, r\_min, rbar, delta\_dc & always & \mref{Thm.}{thm:dark-counts}, \mref{Prop.}{prop:dark-slack-bound}, \mref{Req.}{req:dark-channel-record}\\
\addlinespace
\iface{high_photon_bounds}\newline \{(O\_i, (p\_i\textasciicircum{}(t))\_\{t in T\})\} & always & \mref{Def.}{def:slot}, \mref{Def.}{def:pmc}, \mref{Def.}{def:pmcex}, \mref{Def.}{def:enclosed-weight}, \mref{Thm.}{thm:emitted-bound}, \mref{Thm.}{thm:enclosed-weight}, \mref{Req.}{req:rigorous-weight}, \mref{Thm.}{thm:dark-counts}, \mref{Prop.}{prop:dark-slack-bound}, \mref{Cor.}{cor:emitted-unconditioned}\\
\addlinespace
\iface{bound_quantifier_domain}\newline N, T, \{m\_t\}, \{m*\_t\}, \{beta-hat*\textasciicircum{}(t)(n)\} & always & \mref{Thm.}{thm:emitted-bound}, \mref{Def.}{def:pmc}, \mref{Def.}{def:pmcex}, \mref{Def.}{def:enclosed-weight}, \mref{Req.}{req:rigorous-weight}, \mref{Req.}{req:weight-provenance}\\
\addlinespace
\iface{validation_results}\newline r\_1..r\_6, g\_\{N,m\_t\} & always & \mref{Def.}{def:residuals}, \mref{Req.}{req:residual-emission}, \mref{Req.}{req:residual-refusal}, \mref{Prop.}{prop:res5}, \mref{Lem.}{lem:windowtail}, \mref{Thm.}{thm:exact-and-evaluated}\\
\addlinespace
\iface{bound_evaluation} & always & \mref{Thm.}{thm:exact-and-evaluated}, \mref{Req.}{req:evaluation-label}, \mref{Req.}{req:rigorous-weight}, \mref{Def.}{def:evaluation-model}\\
\addlinespace
\iface{assumption_ledger}\newline A-... & always & ---\\
\addlinespace
\iface{calibration_coverage} & always & \mref{Ass.}{ass:trust}\\
\addlinespace
\iface{mathematical_spec_version} & always & ---\\
\bottomrule
\end{xltabular}}
\clearpage
\bookmarksetupnext{level=part}
\section*{Claim index}\addcontentsline{toc}{section}{Claim index}
{\small\begin{xltabular}{\linewidth}{@{}>{\raggedright\arraybackslash}p{2.7cm}>{\raggedright\arraybackslash}Xr@{}}
\toprule
Claim & Title and identifier & Page\\
\midrule
\endhead
\mref{Definition}{def:modes} & Declared mode space\newline{\footnotesize\color{metaGrey}\claimid{receiver/def:modes}} & \pageref{def:modes}\\
\addlinespace[2pt]
\mref{Definition}{def:fock} & Fock space and the second-quantization functor\newline{\footnotesize\color{metaGrey}\claimid{receiver/def:fock}} & \pageref{def:fock}\\
\addlinespace[2pt]
\mref{Definition}{def:embed} & Local element and its embedding\newline{\footnotesize\color{metaGrey}\claimid{receiver/def:embed}} & \pageref{def:embed}\\
\addlinespace[2pt]
\mref{Proposition}{prop:mode-algebra} & Algebra of embedded local elements\newline{\footnotesize\color{metaGrey}\claimid{receiver/prop:mode-algebra}} & \pageref{prop:mode-algebra}\\
\addlinespace[2pt]
\mref{Remark}{rem:overlap} & Disjoint mode supports as the commutation hypothesis\newline{\footnotesize\color{metaGrey}\claimid{receiver/rem:overlap}} & \pageref{rem:overlap}\\
\addlinespace[2pt]
\mref{Lemma}{lem:fock-functor} & Functoriality of the Fock lift\newline{\footnotesize\color{metaGrey}\claimid{receiver/lem:fock-functor}} & \pageref{lem:fock-functor}\\
\addlinespace[2pt]
\mref{Lemma}{lem:fock-maps} & The lift of maps between mode spaces and of direct sums\newline{\footnotesize\color{metaGrey}\claimid{receiver/lem:fock-maps}} & \pageref{lem:fock-maps}\\
\addlinespace[2pt]
\mref{Remark}{rem:migdal} & Mode and particle descriptions in the literature\newline{\footnotesize\color{metaGrey}\claimid{receiver/rem:migdal}} & \pageref{rem:migdal}\\
\addlinespace[2pt]
\mref{Lemma}{lem:fock-block-spectrum} & Spectrum of a lifted no-click operator on a photon-number block\newline{\footnotesize\color{metaGrey}\claimid{receiver/lem:fock-block-spectrum}} & \pageref{lem:fock-block-spectrum}\\
\addlinespace[2pt]
\mref{Lemma}{lem:fock-monotone} & Order monotonicity of the Fock lift\newline{\footnotesize\color{metaGrey}\claimid{receiver/lem:fock-monotone}} & \pageref{lem:fock-monotone}\\
\addlinespace[2pt]
\mref{Proposition}{prop:normord} & Normal-ordered form of a lifted no-click operator\newline{\footnotesize\color{metaGrey}\claimid{receiver/prop:normord}} & \pageref{prop:normord}\\
\addlinespace[2pt]
\mref{Remark}{rem:konig} & The single-mode normal-ordering identity in the literature\newline{\footnotesize\color{metaGrey}\claimid{receiver/rem:konig}} & \pageref{rem:konig}\\
\addlinespace[2pt]
\mref{Remark}{rem:whyprojector} & Derivations without normal ordering\newline{\footnotesize\color{metaGrey}\claimid{receiver/rem:whyprojector}} & \pageref{rem:whyprojector}\\
\addlinespace[2pt]
\mref{Definition}{def:instrument} & Receiver instrument\newline{\footnotesize\color{metaGrey}\claimid{receiver/def:instrument}} & \pageref{def:instrument}\\
\addlinespace[2pt]
\mref{Definition}{def:gram} & Port Grams and the loss deficiency\newline{\footnotesize\color{metaGrey}\claimid{receiver/def:gram}} & \pageref{def:gram}\\
\addlinespace[2pt]
\mref{Definition}{def:probe} & Probe declaration\newline{\footnotesize\color{metaGrey}\claimid{receiver/def:probe}} & \pageref{def:probe}\\
\addlinespace[2pt]
\mref{Lemma}{lem:dilation} & Unitary dilation of a contraction\newline{\footnotesize\color{metaGrey}\claimid{receiver/lem:dilation}} & \pageref{lem:dilation}\\
\addlinespace[2pt]
\mref{Definition}{def:physmodel} & The physical measurement model\newline{\footnotesize\color{metaGrey}\claimid{receiver/def:physmodel}} & \pageref{def:physmodel}\\
\addlinespace[2pt]
\mref{Lemma}{lem:vacelim} & Exactness of vacuum elimination\newline{\footnotesize\color{metaGrey}\claimid{receiver/lem:vacelim}} & \pageref{lem:vacelim}\\
\addlinespace[2pt]
\mref{Remark}{rem:lossvsprobe} & Loss ancillas and probe inputs\newline{\footnotesize\color{metaGrey}\claimid{receiver/rem:lossvsprobe}} & \pageref{rem:lossvsprobe}\\
\addlinespace[2pt]
\mref{Definition}{def:catalogue} & Component catalogue\newline{\footnotesize\color{metaGrey}\claimid{receiver/def:catalogue}} & \pageref{def:catalogue}\\
\addlinespace[2pt]
\mref{Definition}{def:graph} & Declared element graph\newline{\footnotesize\color{metaGrey}\claimid{receiver/def:graph}} & \pageref{def:graph}\\
\addlinespace[2pt]
\mref{Definition}{def:compile} & Compilation of a declared element graph\newline{\footnotesize\color{metaGrey}\claimid{receiver/def:compile}} & \pageref{def:compile}\\
\addlinespace[2pt]
\mref{Theorem}{thm:compilation-well-defined} & Well-definedness of compilation\newline{\footnotesize\color{metaGrey}\claimid{receiver/thm:compilation-well-defined}} & \pageref{thm:compilation-well-defined}\\
\addlinespace[2pt]
\mref{Corollary}{cor:compiled-instrument} & Receiver instrument from a well-formed graph declaration\newline{\footnotesize\color{metaGrey}\claimid{receiver/cor:compiled-instrument}} & \pageref{cor:compiled-instrument}\\
\addlinespace[2pt]
\mref{Remark}{rem:e3} & Order dependence without support comparability\newline{\footnotesize\color{metaGrey}\claimid{receiver/rem:e3}} & \pageref{rem:e3}\\
\addlinespace[2pt]
\mref{Remark}{rem:edgemeaning} & Meaning of an edge\newline{\footnotesize\color{metaGrey}\claimid{receiver/rem:edgemeaning}} & \pageref{rem:edgemeaning}\\
\addlinespace[2pt]
\mref{Remark}{rem:noloops} & Feedback paths outside the acyclic scope\newline{\footnotesize\color{metaGrey}\claimid{receiver/rem:noloops}} & \pageref{rem:noloops}\\
\addlinespace[2pt]
\mref{Remark}{rem:acyc-efficiency} & Where the element graph enters the efficiency theorem\newline{\footnotesize\color{metaGrey}\claimid{receiver/rem:acyc-efficiency}} & \pageref{rem:acyc-efficiency}\\
\addlinespace[2pt]
\mref{Theorem}{thm:port-grams} & Port Grams and dilation independence\newline{\footnotesize\color{metaGrey}\claimid{receiver/thm:port-grams}} & \pageref{thm:port-grams}\\
\addlinespace[2pt]
\mref{Lemma}{lem:deficiency} & Decomposition of the loss deficiency\newline{\footnotesize\color{metaGrey}\claimid{receiver/lem:deficiency}} & \pageref{lem:deficiency}\\
\addlinespace[2pt]
\mref{Remark}{rem:unmonitored} & Unmonitored output modes\newline{\footnotesize\color{metaGrey}\claimid{receiver/rem:unmonitored}} & \pageref{rem:unmonitored}\\
\addlinespace[2pt]
\mref{Definition}{def:clickalphabet} & Click-pattern alphabet\newline{\footnotesize\color{metaGrey}\claimid{receiver/def:clickalphabet}} & \pageref{def:clickalphabet}\\
\addlinespace[2pt]
\mref{Theorem}{thm:click-povm} & Click-pattern POVM\newline{\footnotesize\color{metaGrey}\claimid{receiver/thm:click-povm}} & \pageref{thm:click-povm}\\
\addlinespace[2pt]
\mref{Remark}{rem:positivity} & Source of positivity of the click-pattern elements\newline{\footnotesize\color{metaGrey}\claimid{receiver/rem:positivity}} & \pageref{rem:positivity}\\
\addlinespace[2pt]
\mref{Remark}{rem:normordelements} & Click-pattern elements in normal-ordered form\newline{\footnotesize\color{metaGrey}\claimid{receiver/rem:normordelements}} & \pageref{rem:normordelements}\\
\addlinespace[2pt]
\mref{Remark}{rem:small-k} & Small numbers of ports\newline{\footnotesize\color{metaGrey}\claimid{receiver/rem:small-k}} & \pageref{rem:small-k}\\
\addlinespace[2pt]
\mref{Remark}{rem:blockmeaning} & Scope of the photon-number block structure\newline{\footnotesize\color{metaGrey}\claimid{receiver/rem:blockmeaning}} & \pageref{rem:blockmeaning}\\
\addlinespace[2pt]
\mref{Theorem}{thm:representation-independence} & Representation independence of the click-pattern POVM\newline{\footnotesize\color{metaGrey}\claimid{receiver/thm:representation-independence}} & \pageref{thm:representation-independence}\\
\addlinespace[2pt]
\mref{Corollary}{cor:phases} & Unobservability of per-port phases\newline{\footnotesize\color{metaGrey}\claimid{receiver/cor:phases}} & \pageref{cor:phases}\\
\addlinespace[2pt]
\mref{Corollary}{cor:extinction} & Extinction criterion for rank one at two input modes\newline{\footnotesize\color{metaGrey}\claimid{receiver/cor:extinction}} & \pageref{cor:extinction}\\
\addlinespace[2pt]
\mref{Remark}{rem:rankone} & Ports of rank greater than one\newline{\footnotesize\color{metaGrey}\claimid{receiver/rem:rankone}} & \pageref{rem:rankone}\\
\addlinespace[2pt]
\mref{Definition}{def:coarse} & Declared coarse-graining and the emitted POVM\newline{\footnotesize\color{metaGrey}\claimid{receiver/def:coarse}} & \pageref{def:coarse}\\
\addlinespace[2pt]
\mref{Theorem}{thm:coarse-graining} & Coarse-graining of the click-pattern POVM\newline{\footnotesize\color{metaGrey}\claimid{receiver/thm:coarse-graining}} & \pageref{thm:coarse-graining}\\
\addlinespace[2pt]
\mref{Definition}{def:bounddata} & Admissible bound data and the bound function\newline{\footnotesize\color{metaGrey}\claimid{receiver/def:bounddata}} & \pageref{def:bounddata}\\
\addlinespace[2pt]
\mref{Theorem}{thm:multiclick-bound} & Lower bound on every multiclick block\newline{\footnotesize\color{metaGrey}\claimid{receiver/thm:multiclick-bound}} & \pageref{thm:multiclick-bound}\\
\addlinespace[2pt]
\mref{Remark}{rem:k-minus-1} & The coefficient $K-1$ of the all-dark element\newline{\footnotesize\color{metaGrey}\claimid{receiver/rem:k-minus-1}} & \pageref{rem:k-minus-1}\\
\addlinespace[2pt]
\mref{Remark}{rem:kamin} & Relation to a published multiclick weight bound\newline{\footnotesize\color{metaGrey}\claimid{receiver/rem:kamin}} & \pageref{rem:kamin}\\
\addlinespace[2pt]
\mref{Lemma}{lem:collapse} & Collapse of the tail infimum to the first photon number\newline{\footnotesize\color{metaGrey}\claimid{receiver/lem:collapse}} & \pageref{lem:collapse}\\
\addlinespace[2pt]
\mref{Remark}{rem:n0} & The monotonicity step at zero photons\newline{\footnotesize\color{metaGrey}\claimid{receiver/rem:n0}} & \pageref{rem:n0}\\
\addlinespace[2pt]
\mref{Remark}{rem:collapsesharp} & Sharpness of the collapse lemma\newline{\footnotesize\color{metaGrey}\claimid{receiver/rem:collapsesharp}} & \pageref{rem:collapsesharp}\\
\addlinespace[2pt]
\mref{Lemma}{lem:windowtail} & Window-and-tail bound on the tail infimum\newline{\footnotesize\color{metaGrey}\claimid{receiver/lem:windowtail}} & \pageref{lem:windowtail}\\
\addlinespace[2pt]
\mref{Definition}{def:exactbnd} & The exact block bound function\newline{\footnotesize\color{metaGrey}\claimid{receiver/def:exactbnd}} & \pageref{def:exactbnd}\\
\addlinespace[2pt]
\mref{Lemma}{lem:exactbasic} & Range of the exact block bound\newline{\footnotesize\color{metaGrey}\claimid{receiver/lem:exactbasic}} & \pageref{lem:exactbasic}\\
\addlinespace[2pt]
\mref{Remark}{rem:exactpositivity} & Positivity of the multiclick element as a separate hypothesis\newline{\footnotesize\color{metaGrey}\claimid{receiver/rem:exactpositivity}} & \pageref{rem:exactpositivity}\\
\addlinespace[2pt]
\mref{Lemma}{lem:exactwindow} & Exact window with an analytic remainder\newline{\footnotesize\color{metaGrey}\claimid{receiver/lem:exactwindow}} & \pageref{lem:exactwindow}\\
\addlinespace[2pt]
\mref{Remark}{rem:exactscope} & Scope of the exact-window lemma\newline{\footnotesize\color{metaGrey}\claimid{receiver/rem:exactscope}} & \pageref{rem:exactscope}\\
\addlinespace[2pt]
\mref{Definition}{def:setting-family} & Setting family\newline{\footnotesize\color{metaGrey}\claimid{receiver/def:setting-family}} & \pageref{def:setting-family}\\
\addlinespace[2pt]
\mref{Corollary}{cor:unconditioned-family} & The unconditioned receiver is a setting family\newline{\footnotesize\color{metaGrey}\claimid{receiver/cor:unconditioned-family}} & \pageref{cor:unconditioned-family}\\
\addlinespace[2pt]
\mref{Definition}{def:pmc} & The uniform high-photon multiclick weight\newline{\footnotesize\color{metaGrey}\claimid{receiver/def:pmc}} & \pageref{def:pmc}\\
\addlinespace[2pt]
\mref{Remark}{rem:pmc-branch-choice} & Choice of the collapse predicate\newline{\footnotesize\color{metaGrey}\claimid{receiver/rem:pmc-branch-choice}} & \pageref{rem:pmc-branch-choice}\\
\addlinespace[2pt]
\mref{Remark}{rem:flag-instances} & Where the exact flag is set, and what a general flag selects\newline{\footnotesize\color{metaGrey}\claimid{receiver/rem:flag-instances}} & \pageref{rem:flag-instances}\\
\addlinespace[2pt]
\mref{Remark}{rem:noslack} & Sign of the collapse slack\newline{\footnotesize\color{metaGrey}\claimid{receiver/rem:noslack}} & \pageref{rem:noslack}\\
\addlinespace[2pt]
\mref{Definition}{def:pmcex} & Exact-window weight of a setting with a declared depth\newline{\footnotesize\color{metaGrey}\claimid{receiver/def:pmcex}} & \pageref{def:pmcex}\\
\addlinespace[2pt]
\mref{Remark}{rem:exact-window-remainders} & The analytic remainder of each exact-window branch\newline{\footnotesize\color{metaGrey}\claimid{receiver/rem:exact-window-remainders}} & \pageref{rem:exact-window-remainders}\\
\addlinespace[2pt]
\mref{Proposition}{prop:exactnoworse} & Dominance of the exact-window weight over the analytic weight\newline{\footnotesize\color{metaGrey}\claimid{receiver/prop:exactnoworse}} & \pageref{prop:exactnoworse}\\
\addlinespace[2pt]
\mref{Remark}{rem:noworsefl} & The exact-window comparison under floating-point evaluation\newline{\footnotesize\color{metaGrey}\claimid{receiver/rem:noworsefl}} & \pageref{rem:noworsefl}\\
\addlinespace[2pt]
\mref{Remark}{rem:exacteig} & The eigenvalue-slack assumption behind the exact window\newline{\footnotesize\color{metaGrey}\claimid{receiver/rem:exacteig}} & \pageref{rem:exacteig}\\
\addlinespace[2pt]
\mref{Definition}{def:enclosed-weight} & The weight evaluated on enclosed bound data\newline{\footnotesize\color{metaGrey}\claimid{receiver/def:enclosed-weight}} & \pageref{def:enclosed-weight}\\
\addlinespace[2pt]
\mref{Theorem}{thm:enclosed-weight} & The weight on enclosed bound data bounds the multiclick element\newline{\footnotesize\color{metaGrey}\claimid{receiver/thm:enclosed-weight}} & \pageref{thm:enclosed-weight}\\
\addlinespace[2pt]
\mref{Theorem}{thm:emitted-bound} & Multiclick bounds carried by the emitted weights\newline{\footnotesize\color{metaGrey}\claimid{receiver/thm:emitted-bound}} & \pageref{thm:emitted-bound}\\
\addlinespace[2pt]
\mref{Corollary}{cor:emitted-unconditioned} & Emitted multiclick weight of a probe-free receiver\newline{\footnotesize\color{metaGrey}\claimid{receiver/cor:emitted-unconditioned}} & \pageref{cor:emitted-unconditioned}\\
\addlinespace[2pt]
\mref{Definition}{def:slot} & The strengthening slot\newline{\footnotesize\color{metaGrey}\claimid{receiver/def:slot}} & \pageref{def:slot}\\
\addlinespace[2pt]
\mref{Proposition}{prop:slotvalid} & Range of the strengthening-slot weights\newline{\footnotesize\color{metaGrey}\claimid{receiver/prop:slotvalid}} & \pageref{prop:slotvalid}\\
\addlinespace[2pt]
\mref{Remark}{rem:notsharp} & Slack of the analytic bound against the block eigenvalue\newline{\footnotesize\color{metaGrey}\claimid{receiver/rem:notsharp}} & \pageref{rem:notsharp}\\
\addlinespace[2pt]
\mref{Definition}{def:complex-enclosure} & Complex enclosures, and enclosures of square roots and angles\newline{\footnotesize\color{metaGrey}\claimid{receiver/def:complex-enclosure}} & \pageref{def:complex-enclosure}\\
\addlinespace[2pt]
\mref{Lemma}{lem:enclosure-arithmetic} & Arithmetic on complex enclosures\newline{\footnotesize\color{metaGrey}\claimid{receiver/lem:enclosure-arithmetic}} & \pageref{lem:enclosure-arithmetic}\\
\addlinespace[2pt]
\mref{Lemma}{lem:fock-permanent} & Entries of a lifted operator in the occupation basis\newline{\footnotesize\color{metaGrey}\claimid{receiver/lem:fock-permanent}} & \pageref{lem:fock-permanent}\\
\addlinespace[2pt]
\mref{Lemma}{lem:fock-permanent-maps} & The lift of a map between two mode sets\newline{\footnotesize\color{metaGrey}\claimid{receiver/lem:fock-permanent-maps}} & \pageref{lem:fock-permanent-maps}\\
\addlinespace[2pt]
\mref{Theorem}{thm:enclosed-blocks} & Rigorous enclosures of the photon-number blocks\newline{\footnotesize\color{metaGrey}\claimid{receiver/thm:enclosed-blocks}} & \pageref{thm:enclosed-blocks}\\
\addlinespace[2pt]
\mref{Definition}{def:port-dilation} & Port dilation of a receiver instrument\newline{\footnotesize\color{metaGrey}\claimid{receiver/def:port-dilation}} & \pageref{def:port-dilation}\\
\addlinespace[2pt]
\mref{Theorem}{thm:factor-blocks} & Photon-number blocks in enclosed factor form\newline{\footnotesize\color{metaGrey}\claimid{receiver/thm:factor-blocks}} & \pageref{thm:factor-blocks}\\
\addlinespace[2pt]
\mref{Definition}{def:exact-route} & Gaussian-rational port Grams and the exact route\newline{\footnotesize\color{metaGrey}\claimid{receiver/def:exact-route}} & \pageref{def:exact-route}\\
\addlinespace[2pt]
\mref{Theorem}{thm:exact-blocks} & Exact entries of the blocks on the exact route\newline{\footnotesize\color{metaGrey}\claimid{receiver/thm:exact-blocks}} & \pageref{thm:exact-blocks}\\
\addlinespace[2pt]
\mref{Definition}{def:residuals} & The six residuals\newline{\footnotesize\color{metaGrey}\claimid{receiver/def:residuals}} & \pageref{def:residuals}\\
\addlinespace[2pt]
\mref{Proposition}{prop:res5} & Properties of the Hermiticity-and-positivity defect\newline{\footnotesize\color{metaGrey}\claimid{receiver/prop:res5}} & \pageref{prop:res5}\\
\addlinespace[2pt]
\mref{Remark}{rem:res5why} & What the fifth residual checks that the others do not\newline{\footnotesize\color{metaGrey}\claimid{receiver/rem:res5why}} & \pageref{rem:res5why}\\
\addlinespace[2pt]
\mref{Definition}{def:tolmap} & Receiver tolerances and the exact conditions they decide\newline{\footnotesize\color{metaGrey}\claimid{receiver/def:tolmap}} & \pageref{def:tolmap}\\
\addlinespace[2pt]
\mref{Remark}{rem:bound-data-band} & The bound-data test and the contractivity test on the band\newline{\footnotesize\color{metaGrey}\claimid{receiver/rem:bound-data-band}} & \pageref{rem:bound-data-band}\\
\addlinespace[2pt]
\mref{Remark}{rem:tolerancedomain} & The domain of the receiver tolerances\newline{\footnotesize\color{metaGrey}\claimid{receiver/rem:tolerancedomain}} & \pageref{rem:tolerancedomain}\\
\addlinespace[2pt]
\mref{Remark}{rem:ranktwosided} & Two-sided errors of a rank cutoff, and the choice of threshold\newline{\footnotesize\color{metaGrey}\claimid{receiver/rem:ranktwosided}} & \pageref{rem:ranktwosided}\\
\addlinespace[2pt]
\mref{Remark}{rem:cdirection} & The direction of the unclamped $c$ at odd and even photon numbers\newline{\footnotesize\color{metaGrey}\claimid{receiver/rem:cdirection}} & \pageref{rem:cdirection}\\
\addlinespace[2pt]
\mref{Definition}{def:evaluation-model} & The evaluation model: floating point, enclosures and rigour\newline{\footnotesize\color{metaGrey}\claimid{receiver/def:evaluation-model}} & \pageref{def:evaluation-model}\\
\addlinespace[2pt]
\mref{Theorem}{thm:exact-and-evaluated} & Exact statements, floating-point and rigorous evaluation\newline{\footnotesize\color{metaGrey}\claimid{receiver/thm:exact-and-evaluated}} & \pageref{thm:exact-and-evaluated}\\
\addlinespace[2pt]
\mref{Remark}{rem:exactness-scope} & Scope of the exactness rule\newline{\footnotesize\color{metaGrey}\claimid{receiver/rem:exactness-scope}} & \pageref{rem:exactness-scope}\\
\addlinespace[2pt]
\mref{Remark}{rem:ia} & Directions of rigorous enclosures of the bound inputs\newline{\footnotesize\color{metaGrey}\claimid{receiver/rem:ia}} & \pageref{rem:ia}\\
\addlinespace[2pt]
\mref{Assumption}{ass:trust} & Trusted and exactly characterised receiver\newline{\footnotesize\color{metaGrey}\claimid{receiver/ass:trust}} & \pageref{ass:trust}\\
\addlinespace[2pt]
\mref{Remark}{rem:unchecked-assumptions} & Assumptions this module cannot check\newline{\footnotesize\color{metaGrey}\claimid{receiver/rem:unchecked-assumptions}} & \pageref{rem:unchecked-assumptions}\\
\addlinespace[2pt]
\mref{Assumption}{ass:passive} & Passive linear optics\newline{\footnotesize\color{metaGrey}\claimid{receiver/ass:passive}} & \pageref{ass:passive}\\
\addlinespace[2pt]
\mref{Assumption}{ass:vac} & Vacuum on undeclared input modes\newline{\footnotesize\color{metaGrey}\claimid{receiver/ass:vac}} & \pageref{ass:vac}\\
\addlinespace[2pt]
\mref{Assumption}{ass:modes} & Completeness of the declared modes\newline{\footnotesize\color{metaGrey}\claimid{receiver/ass:modes}} & \pageref{ass:modes}\\
\addlinespace[2pt]
\mref{Assumption}{ass:thresh} & Threshold detectors without memory\newline{\footnotesize\color{metaGrey}\claimid{receiver/ass:thresh}} & \pageref{ass:thresh}\\
\addlinespace[2pt]
\mref{Assumption}{ass:port} & Mode-insensitive port response\newline{\footnotesize\color{metaGrey}\claimid{receiver/ass:port}} & \pageref{ass:port}\\
\addlinespace[2pt]
\mref{Assumption}{ass:exact} & Exact declared parameters\newline{\footnotesize\color{metaGrey}\claimid{receiver/ass:exact}} & \pageref{ass:exact}\\
\addlinespace[2pt]
\mref{Assumption}{ass:nosel} & Device fixed independently of the adversary\newline{\footnotesize\color{metaGrey}\claimid{receiver/ass:nosel}} & \pageref{ass:nosel}\\
\addlinespace[2pt]
\mref{Assumption}{ass:round} & Round homogeneity of the instrument, probe and dark counts\newline{\footnotesize\color{metaGrey}\claimid{receiver/ass:round}} & \pageref{ass:round}\\
\addlinespace[2pt]
\mref{Assumption}{ass:eig} & Declared slack on the computed block eigenvalues\newline{\footnotesize\color{metaGrey}\claimid{receiver/ass:eig}} & \pageref{ass:eig}\\
\addlinespace[2pt]
\mref{Requirement}{req:graph-refusal} & Refusal of a graph that is not well formed\newline{\footnotesize\color{metaGrey}\claimid{receiver/req:graph-refusal}} & \pageref{req:graph-refusal}\\
\addlinespace[2pt]
\mref{Requirement}{req:probe-diagonal-admission} & Admission of Gaussian-tier probe states only\newline{\footnotesize\color{metaGrey}\claimid{receiver/req:probe-diagonal-admission}} & \pageref{req:probe-diagonal-admission}\\
\addlinespace[2pt]
\mref{Requirement}{req:declared-port-map} & Declared port map\newline{\footnotesize\color{metaGrey}\claimid{receiver/req:declared-port-map}} & \pageref{req:declared-port-map}\\
\addlinespace[2pt]
\mref{Requirement}{req:one-mode-per-port} & One mode per port only where the rank criterion holds\newline{\footnotesize\color{metaGrey}\claimid{receiver/req:one-mode-per-port}} & \pageref{req:one-mode-per-port}\\
\addlinespace[2pt]
\mref{Requirement}{req:declared-alphabet} & Declared emitted alphabet and coarse-graining\newline{\footnotesize\color{metaGrey}\claimid{receiver/req:declared-alphabet}} & \pageref{req:declared-alphabet}\\
\addlinespace[2pt]
\mref{Requirement}{req:window-declaration} & Declared window at every setting whose exact flag is general\newline{\footnotesize\color{metaGrey}\claimid{receiver/req:window-declaration}} & \pageref{req:window-declaration}\\
\addlinespace[2pt]
\mref{Requirement}{req:exact-window-declaration} & Declaration of an exact-window depth and its tolerances\newline{\footnotesize\color{metaGrey}\claimid{receiver/req:exact-window-declaration}} & \pageref{req:exact-window-declaration}\\
\addlinespace[2pt]
\mref{Requirement}{req:exact-window-record} & Computed block eigenvalues carried as computed\newline{\footnotesize\color{metaGrey}\claimid{receiver/req:exact-window-record}} & \pageref{req:exact-window-record}\\
\addlinespace[2pt]
\mref{Requirement}{req:weight-provenance} & Provenance of each emitted multiclick weight\newline{\footnotesize\color{metaGrey}\claimid{receiver/req:weight-provenance}} & \pageref{req:weight-provenance}\\
\addlinespace[2pt]
\mref{Requirement}{req:rigorous-weight} & Rigorous evaluation of the multiclick weight\newline{\footnotesize\color{metaGrey}\claimid{receiver/req:rigorous-weight}} & \pageref{req:rigorous-weight}\\
\addlinespace[2pt]
\mref{Requirement}{req:residual-emission} & Emission of the residuals that exist on the route and request\newline{\footnotesize\color{metaGrey}\claimid{receiver/req:residual-emission}} & \pageref{req:residual-emission}\\
\addlinespace[2pt]
\mref{Requirement}{req:residual-refusal} & Refusal on a residual above its tolerance, and no repair\newline{\footnotesize\color{metaGrey}\claimid{receiver/req:residual-refusal}} & \pageref{req:residual-refusal}\\
\addlinespace[2pt]
\mref{Requirement}{req:tolerance-tests} & Decisions at the receiver tolerances, and refusal of failed bound data\newline{\footnotesize\color{metaGrey}\claimid{receiver/req:tolerance-tests}} & \pageref{req:tolerance-tests}\\
\addlinespace[2pt]
\mref{Requirement}{req:tolerance-domain} & Validation of the receiver tolerances against their domain\newline{\footnotesize\color{metaGrey}\claimid{receiver/req:tolerance-domain}} & \pageref{req:tolerance-domain}\\
\addlinespace[2pt]
\mref{Requirement}{req:evaluation-label} & The evaluation label on every artifact\newline{\footnotesize\color{metaGrey}\claimid{receiver/req:evaluation-label}} & \pageref{req:evaluation-label}\\
\addlinespace[2pt]
\mref{Requirement}{req:exact-route-decision} & Deciding the exact route and enclosing the port Grams\newline{\footnotesize\color{metaGrey}\claimid{receiver/req:exact-route-decision}} & \pageref{req:exact-route-decision}\\
\addlinespace[2pt]
\mref{Requirement}{req:decidable-assumptions} & Validation of the assumptions on the declared data\newline{\footnotesize\color{metaGrey}\claimid{receiver/req:decidable-assumptions}} & \pageref{req:decidable-assumptions}\\
\addlinespace[2pt]
\mref{Requirement}{req:reserved-tolerance-field} & Refusal of a non-zero parameter tolerance\newline{\footnotesize\color{metaGrey}\claimid{receiver/req:reserved-tolerance-field}} & \pageref{req:reserved-tolerance-field}\\
\addlinespace[2pt]
\mref{Requirement}{req:eig-ledger-row} & Invoking settings on the eigenvalue-slack row of the ledger\newline{\footnotesize\color{metaGrey}\claimid{receiver/req:eig-ledger-row}} & \pageref{req:eig-ledger-row}\\
\addlinespace[2pt]
\mref{Lemma}{lem:attenuated} & Admissibility of an attenuated unitary\newline{\footnotesize\color{metaGrey}\claimid{receiver/component-catalogue/lem:attenuated}} & \pageref{lem:attenuated}\\
\addlinespace[2pt]
\mref{Definition}{def:catalogue-entries} & Entries of the component catalogue\newline{\footnotesize\color{metaGrey}\claimid{receiver/component-catalogue/def:catalogue-entries}} & \pageref{def:catalogue-entries}\\
\addlinespace[2pt]
\mref{Theorem}{thm:catalogue-admissibility} & Admissibility of the catalogue entries\newline{\footnotesize\color{metaGrey}\claimid{receiver/component-catalogue/thm:catalogue-admissibility}} & \pageref{thm:catalogue-admissibility}\\
\addlinespace[2pt]
\mref{Lemma}{lem:catalogue-dilations} & Loss factors of the catalogue entries and the dilation of a graph\newline{\footnotesize\color{metaGrey}\claimid{receiver/component-catalogue/lem:catalogue-dilations}} & \pageref{lem:catalogue-dilations}\\
\addlinespace[2pt]
\mref{Definition}{def:c10-tolerance} & Admission test of the generic block\newline{\footnotesize\color{metaGrey}\claimid{receiver/component-catalogue/def:c10-tolerance}} & \pageref{def:c10-tolerance}\\
\addlinespace[2pt]
\mref{Remark}{rem:rotator} & The rotator outside the retarder family\newline{\footnotesize\color{metaGrey}\claimid{receiver/component-catalogue/rem:rotator}} & \pageref{rem:rotator}\\
\addlinespace[2pt]
\mref{Remark}{rem:entry-roles} & Physical roles of the catalogue entries\newline{\footnotesize\color{metaGrey}\claimid{receiver/component-catalogue/rem:entry-roles}} & \pageref{rem:entry-roles}\\
\addlinespace[2pt]
\mref{Remark}{rem:layers} & Detector and source layers outside the catalogue\newline{\footnotesize\color{metaGrey}\claimid{receiver/component-catalogue/rem:layers}} & \pageref{rem:layers}\\
\addlinespace[2pt]
\mref{Remark}{rem:bsconv} & Beamsplitter conventions and the port Grams\newline{\footnotesize\color{metaGrey}\claimid{receiver/component-catalogue/rem:bsconv}} & \pageref{rem:bsconv}\\
\addlinespace[2pt]
\mref{Remark}{rem:out} & Elements outside the catalogue\newline{\footnotesize\color{metaGrey}\claimid{receiver/component-catalogue/rem:out}} & \pageref{rem:out}\\
\addlinespace[2pt]
\mref{Example}{ex:bb84} & A BB84 receiver compiled from the catalogue\newline{\footnotesize\color{metaGrey}\claimid{receiver/component-catalogue/ex:bb84}} & \pageref{ex:bb84}\\
\addlinespace[2pt]
\mref{Example}{ex:bb84bound} & Analytic multiclick weight of the BB84 receiver\newline{\footnotesize\color{metaGrey}\claimid{receiver/component-catalogue/ex:bb84bound}} & \pageref{ex:bb84bound}\\
\addlinespace[2pt]
\mref{Example}{ex:bb84exact} & Exact window on the BB84 receiver\newline{\footnotesize\color{metaGrey}\claimid{receiver/component-catalogue/ex:bb84exact}} & \pageref{ex:bb84exact}\\
\addlinespace[2pt]
\mref{Theorem}{thm:efficiency} & Detector efficiency in the port Grams\newline{\footnotesize\color{metaGrey}\claimid{receiver/detector-imperfections/thm:efficiency}} & \pageref{thm:efficiency}\\
\addlinespace[2pt]
\mref{Remark}{rem:effrank} & A port Gram whose rank drops under within-port attenuation\newline{\footnotesize\color{metaGrey}\claimid{receiver/detector-imperfections/rem:effrank}} & \pageref{rem:effrank}\\
\addlinespace[2pt]
\mref{Remark}{rem:wrongcheck} & Failure of the path-product loss audit under interference\newline{\footnotesize\color{metaGrey}\claimid{receiver/detector-imperfections/rem:wrongcheck}} & \pageref{rem:wrongcheck}\\
\addlinespace[2pt]
\mref{Remark}{rem:powervsclick} & Power data and click data measure different objects\newline{\footnotesize\color{metaGrey}\claimid{receiver/detector-imperfections/rem:powervsclick}} & \pageref{rem:powervsclick}\\
\addlinespace[2pt]
\mref{Remark}{rem:sperlingeta} & The single-mode efficiency substitution in the photo-counting formula\newline{\footnotesize\color{metaGrey}\claimid{receiver/detector-imperfections/rem:sperlingeta}} & \pageref{rem:sperlingeta}\\
\addlinespace[2pt]
\mref{Definition}{def:darkadm} & Dark-count channel and its independent product form\newline{\footnotesize\color{metaGrey}\claimid{receiver/detector-imperfections/def:darkadm}} & \pageref{def:darkadm}\\
\addlinespace[2pt]
\mref{Theorem}{thm:dark-counts} & Dark counts on a photon-number block-diagonal POVM\newline{\footnotesize\color{metaGrey}\claimid{receiver/detector-imperfections/thm:dark-counts}} & \pageref{thm:dark-counts}\\
\addlinespace[2pt]
\mref{Definition}{def:darktol} & Receiver tolerances deciding dark-channel admissibility\newline{\footnotesize\color{metaGrey}\claimid{receiver/detector-imperfections/def:darktol}} & \pageref{def:darktol}\\
\addlinespace[2pt]
\mref{Proposition}{prop:dark-slack-bound} & Multiclick bound under a dark-count channel admitted with slack\newline{\footnotesize\color{metaGrey}\claimid{receiver/detector-imperfections/prop:dark-slack-bound}} & \pageref{prop:dark-slack-bound}\\
\addlinespace[2pt]
\mref{Remark}{rem:rmintrap} & The no-click column in the minimum defining $r_{\min}$\newline{\footnotesize\color{metaGrey}\claimid{receiver/detector-imperfections/rem:rmintrap}} & \pageref{rem:rmintrap}\\
\addlinespace[2pt]
\mref{Remark}{rem:d2slack} & Cost of admitting a dark-count channel with slack\newline{\footnotesize\color{metaGrey}\claimid{receiver/detector-imperfections/rem:d2slack}} & \pageref{rem:d2slack}\\
\addlinespace[2pt]
\mref{Remark}{rem:corrdark} & Correlated dark counts within a round\newline{\footnotesize\color{metaGrey}\claimid{receiver/detector-imperfections/rem:corrdark}} & \pageref{rem:corrdark}\\
\addlinespace[2pt]
\mref{Definition}{def:dark-composed} & The dark-composed family\newline{\footnotesize\color{metaGrey}\claimid{receiver/detector-imperfections/def:dark-composed}} & \pageref{def:dark-composed}\\
\addlinespace[2pt]
\mref{Corollary}{cor:exact-dark-composed} & The dark-composed reference, enclosed or exact\newline{\footnotesize\color{metaGrey}\claimid{receiver/detector-imperfections/cor:exact-dark-composed}} & \pageref{cor:exact-dark-composed}\\
\addlinespace[2pt]
\mref{Proposition}{prop:noise-floor} & Loss and dark counts make every dark-composed element definite\newline{\footnotesize\color{metaGrey}\claimid{receiver/detector-imperfections/prop:noise-floor}} & \pageref{prop:noise-floor}\\
\addlinespace[2pt]
\mref{Remark}{rem:noise-floor-receiver} & What a receiver declaration needs for a rate that rests on its model\newline{\footnotesize\color{metaGrey}\claimid{receiver/detector-imperfections/rem:noise-floor-receiver}} & \pageref{rem:noise-floor-receiver}\\
\addlinespace[2pt]
\mref{Definition}{def:pure-loss} & Pure-loss channel on the channel modes\newline{\footnotesize\color{metaGrey}\claimid{receiver/detector-imperfections/def:pure-loss}} & \pageref{def:pure-loss}\\
\addlinespace[2pt]
\mref{Definition}{def:normalized-station} & Normalized station and factored loss\newline{\footnotesize\color{metaGrey}\claimid{receiver/detector-imperfections/def:normalized-station}} & \pageref{def:normalized-station}\\
\addlinespace[2pt]
\mref{Lemma}{lem:pure-loss-lift} & The pure-loss channel on the lift of a contraction\newline{\footnotesize\color{metaGrey}\claimid{receiver/detector-imperfections/lem:pure-loss-lift}} & \pageref{lem:pure-loss-lift}\\
\addlinespace[2pt]
\mref{Proposition}{prop:efficiency-normalization} & A receiver as uniform loss followed by its normalized station\newline{\footnotesize\color{metaGrey}\claimid{receiver/detector-imperfections/prop:efficiency-normalization}} & \pageref{prop:efficiency-normalization}\\
\addlinespace[2pt]
\mref{Remark}{rem:normalization-gain} & The two-photon multiclick block scales with the squared efficiency\newline{\footnotesize\color{metaGrey}\claimid{receiver/detector-imperfections/rem:normalization-gain}} & \pageref{rem:normalization-gain}\\
\addlinespace[2pt]
\mref{Remark}{rem:normalization-scope} & What the normalization factors out and what it leaves in place\newline{\footnotesize\color{metaGrey}\claimid{receiver/detector-imperfections/rem:normalization-scope}} & \pageref{rem:normalization-scope}\\
\addlinespace[2pt]
\mref{Requirement}{req:efficiency-schema} & Detector efficiency and component loss in separate fields\newline{\footnotesize\color{metaGrey}\claimid{receiver/detector-imperfections/req:efficiency-schema}} & \pageref{req:efficiency-schema}\\
\addlinespace[2pt]
\mref{Requirement}{req:efficiency-record} & The efficiency record lists every declared loss parameter\newline{\footnotesize\color{metaGrey}\claimid{receiver/detector-imperfections/req:efficiency-record}} & \pageref{req:efficiency-record}\\
\addlinespace[2pt]
\mref{Requirement}{req:efficiency-convention} & Declared efficiency convention of the port Grams\newline{\footnotesize\color{metaGrey}\claimid{receiver/detector-imperfections/req:efficiency-convention}} & \pageref{req:efficiency-convention}\\
\addlinespace[2pt]
\mref{Requirement}{req:efficiency-normalization} & Compilation of a probe-free receiver at its normalized station\newline{\footnotesize\color{metaGrey}\claimid{receiver/detector-imperfections/req:efficiency-normalization}} & \pageref{req:efficiency-normalization}\\
\addlinespace[2pt]
\mref{Requirement}{req:normalization-record} & The efficiency record names the normalization\newline{\footnotesize\color{metaGrey}\claimid{receiver/detector-imperfections/req:normalization-record}} & \pageref{req:normalization-record}\\
\addlinespace[2pt]
\mref{Requirement}{req:loss-before-evaluation} & The factored loss is applied before the artefact is evaluated\newline{\footnotesize\color{metaGrey}\claimid{receiver/detector-imperfections/req:loss-before-evaluation}} & \pageref{req:loss-before-evaluation}\\
\addlinespace[2pt]
\mref{Requirement}{req:dark-route-exclusive} & Refusal of a dark-count channel supplied on both routes\newline{\footnotesize\color{metaGrey}\claimid{receiver/detector-imperfections/req:dark-route-exclusive}} & \pageref{req:dark-route-exclusive}\\
\addlinespace[2pt]
\mref{Requirement}{req:dark-channel-record} & The emitted dark-count channel and the scalars of its composed bound\newline{\footnotesize\color{metaGrey}\claimid{receiver/detector-imperfections/req:dark-channel-record}} & \pageref{req:dark-channel-record}\\
\addlinespace[2pt]
\mref{Requirement}{req:dark-composed-reference} & The emitted blocks are dark-composed and carry their reference\newline{\footnotesize\color{metaGrey}\claimid{receiver/detector-imperfections/req:dark-composed-reference}} & \pageref{req:dark-composed-reference}\\
\addlinespace[2pt]
\mref{Requirement}{req:snap-fix} & Cause and fix recorded for a receiver block resting on its snap\newline{\footnotesize\color{metaGrey}\claimid{receiver/detector-imperfections/req:snap-fix}} & \pageref{req:snap-fix}\\
\bottomrule
\end{xltabular}}

\clearpage
\bookmarksetupnext{level=part}
\printindex

\clearpage\phantomsection\bookmarksetupnext{level=part}%
\addcontentsline{toc}{section}{References}
\bibliographystyle{plain}
\bibliography{../../../refs/references}

\end{document}
